\documentclass[12pt,a5paper]{article}
\usepackage{xeCJK}
\usepackage{fontspec}
\usepackage[margin=18mm]{geometry}
\usepackage{parskip}
%\setCJKmainfont{Noto Serif CJK JP}


\setCJKsansfont[BoldFont=HaranoAjiGothic-Bold.otf]{HaranoAjiGothic-Regular.otf}
\setCJKmonofont[BoldFont=HaranoAjiGothic-Bold.otf]{HaranoAjiGothic-Regular.otf}
\setCJKmainfont[
  BoldFont=HaranoAjiMincho-Bold.otf,
  ItalicFont=HaranoAjiMincho-Regular.otf,
  BoldItalicFont=HaranoAjiMincho-Bold.otf
]{HaranoAjiMincho-Regular.otf}

%\setmainfont{TeX Gyre Termes}
\pagestyle{plain}
\begin{document}

\section*{741. Be bad at it}

\textbf{Date}: 2026.04.21 (Tue)\\
\textbf{Index}: udegawarui\\
\textbf{Expression (JA)}: (腕:うで)が(悪:わる)い\\
\textbf{Romanization}: Ude ga warui\\
\textbf{Expression (EN)}: Be bad at it

\textbf{Variants (JA)}: 腕が悪い / 腕がいい / 腕がない\\
\textbf{Variants (EN)}: be bad at it / be good at it / have no skill

\textbf{Intent (JA)}: 能力の評価, 技量の指摘, 軽い批判\\
\textbf{Intent (EN)}: evaluating ability, pointing out skill level, light criticism

\textbf{Tags (JA)}: 即時文法, 慣用表現, 口語, 能力評価, 技量\\
\textbf{Tags (EN)}: immediate grammar, idiomatic expression, colloquial, ability evaluation, skill

\textbf{Adjusted Expression (JA)}:\\
(彼:かれ)は、その(作業:さぎょう)を(十分:じゅうぶん)にこなす(技術:ぎじゅつ)がない。\\
\textbf{Adjusted Expression (EN)}:\\
He does not have the skill to perform that task adequately.

\textbf{Example (JA)}:\\
A「このパチンコ台、(全然:ぜんぜん)(入:はい)らないね」 B「(腕:うで)が(悪:わる)いんだよ」\\
\textbf{Example (EN)}:\\
A: This pachinko machine doesn't work at all. B: You don't have the skill for it.

\textbf{Notes (JA)}:\\
「腕が悪い」は、能力や技量が十分でないことを表現する言い方です。レオナルド熊のコントで、「どうしたんだ、このパチンコ、ぜんぜん入らないじゃないか」「それはパチンコ台が悪いんじゃなくて、あなたの腕が悪いんですよ」と言われた人が、「先生、どうにかしてください」と病院に来る場面があります。このように、評価が身体の問題として受け取られることで、意味のずれが生まれます。

\textbf{Notes (EN)}:\\
The phrase "ude ga warui" refers to a lack of skill or ability. In a comedy sketch by Leonardo Kuma, a man says, "What's wrong with this pachinko machine? It doesn't go in at all," and is told, "It's not the machine, it's your skill that's bad." He then goes to a doctor and says, "Doctor, please do something about it." This shows how an evaluation can be reinterpreted as a physical problem, creating a mismatch in meaning.

\newpage

\section*{740. Is this wrong?}

\textbf{Date}: 2026.04.20 (Mon)\\
\textbf{Index}: korettemachigatteru\\
\textbf{Expression (JA)}: これって(間違:まちが)ってる？\\
\textbf{Romanization}: Kore tte machigatteru?\\
\textbf{Expression (EN)}: Is this wrong?

\textbf{Variants (JA)}: これって間違ってる？ / これって間違い？\\
\textbf{Variants (EN)}: Is this wrong? / Is this a mistake?

\textbf{Intent (JA)}: 確認の表現, 疑念の共有, 口語的な表現\\
\textbf{Intent (EN)}: expression of confirmation, sharing of doubt, colloquial expression

\textbf{Tags (JA)}: 即時文法, 口語, 確認表現, 疑念共有表現, 口語的表現\\
\textbf{Tags (EN)}: immediate grammar, spoken, expression of confirmation, expression of doubt sharing, colloquial expression

\textbf{Adjusted Expression (JA)}:\\
これは(間違:まちが)っていますか？\\
\textbf{Adjusted Expression (EN)}:\\
Is this wrong?

\textbf{Example (JA)}:\\
これって(間違:まちが)ってる？(正解:せいかい)だと(思:おも)うんだけど、...\\
\textbf{Example (EN)}:\\
Is this wrong? I thought it was correct, but...

\textbf{Notes (JA)}:\\
「これって間違ってる？」は、目の前の答えや表現、やり方などについて、正しいかどうかをその場で気軽に確かめるときに使う口語的な表現です。単なる確認だけでなく、「自分は合っていると思ったのだけれど」という軽い疑念やためらいを含んで使われることもあります。

\textbf{Notes (EN)}:\\
'Kore tte machigatteru?' is a colloquial expression used when casually checking, on the spot, whether something such as an answer, expression, or method is wrong. It can express not only simple confirmation but also a slight sense of doubt or hesitation, as in a situation where the speaker thought it was correct.

\newpage

\section*{739. I hope it goes well.}

\textbf{Date}: 2026.04.19 (Sun)\\
\textbf{Index}: umakuikutoiine\\
\textbf{Expression (JA)}: うまく(行:い)くといいね。\\
\textbf{Romanization}: Umaku iku to ii ne.\\
\textbf{Expression (EN)}: I hope it goes well.

\textbf{Variants (JA)}: うまく行くといいね。 / うまくいくといいな。\\
\textbf{Variants (EN)}: I hope it goes well. / Hope it goes well.

\textbf{Intent (JA)}: 願望の表現, 応援の共有, 口語的な表現\\
\textbf{Intent (EN)}: expression of desire, sharing of support, colloquial expression

\textbf{Tags (JA)}: 即時文法, 口語, 願望表現, 応援共有表現, 口語的表現\\
\textbf{Tags (EN)}: immediate grammar, spoken, expression of desire, expression of support sharing, colloquial expression

\textbf{Adjusted Expression (JA)}:\\
うまく(行:い)くといいですね。\\
\textbf{Adjusted Expression (EN)}:\\
I hope it goes well.

\textbf{Example (JA)}:\\
A「(明日:あす)の(試験:しけん)、(大丈夫:だいじょうぶ)かな？」B「うまく(行:い)くといいね！」\\
\textbf{Example (EN)}:\\
A: Do you think I'll be okay on tomorrow's exam? B: I hope it goes well!

\textbf{Notes (JA)}:\\
「うまく行くといいね」は、何かが順調に進んだり、よい結果になったりすることを願うときに使う口語的な表現です。相手を軽く応援したり、気持ちを共有したりするときにも自然に使われます。

\textbf{Notes (EN)}:\\
'Umaku iku to ii ne.' is a colloquial expression used when you hope that something will go well or turn out well. It is also naturally used to lightly encourage someone or to share a supportive feeling.

\newpage

\section*{738. Can you hear me?}

\textbf{Date}: 2026.04.18 (Sat)\\
\textbf{Index}: kikoeru\\
\textbf{Expression (JA)}: (聞:き)こえる？\\
\textbf{Romanization}: Kikoeru?\\
\textbf{Expression (EN)}: Can you hear me?

\textbf{Variants (JA)}: 聞こえる？ / 聞こえてる？\\
\textbf{Variants (EN)}: Can you hear me? / Can you hear me? (casual)

\textbf{Intent (JA)}: 確認の表現, コミュニケーションの共有, 口語的な表現\\
\textbf{Intent (EN)}: expression of confirmation, sharing of communication, colloquial expression

\textbf{Tags (JA)}: 即時文法, 口語, 確認表現, コミュニケーション共有表現, 口語的表現\\
\textbf{Tags (EN)}: immediate grammar, spoken, expression of confirmation, expression of communication sharing, colloquial expression

\textbf{Adjusted Expression (JA)}:\\
(私:わたし)の(声:こえ)が(聞:き)こえますか？\\
\textbf{Adjusted Expression (EN)}:\\
Can you hear my voice?

\textbf{Example (JA)}:\\
(オンライン会議:おんらいんかいぎ)で、(相手:あいて)に(声:こえ)が(聞:き)こえるかどうかを(確認:かくにん)するときに(使:つか)います。\\
\textbf{Example (EN)}:\\
This is used when confirming whether the other person can hear your voice during an online meeting.

\textbf{Notes (JA)}:\\
「聞こえる？」は、相手が自分の声を聞いているかどうかを確認するための口語的な表現。「聞こえてる？」も同じ意味で使われます。

\textbf{Notes (EN)}:\\
'Kikoeru?' is a colloquial expression used to confirm whether the other person can hear your voice. 'Kikoeteru?' is also used with the same meaning.

\newpage

\section*{737. I have to go now.}

\textbf{Date}: 2026.04.17 (Fri)\\
\textbf{Index}: mouikanakucha\\
\textbf{Expression (JA)}: もう、(行:い)かなくっちゃ。\\
\textbf{Romanization}: Mou, ikanakucha.\\
\textbf{Expression (EN)}: I have to go now.

\textbf{Variants (JA)}: もう、行かなくっちゃ。 / もう、行かなきゃ。\\
\textbf{Variants (EN)}: I have to go now. / I gotta go now.

\textbf{Intent (JA)}: 必要性の表現, 時間の切迫の共有, 口語的な表現\\
\textbf{Intent (EN)}: expression of necessity, sharing urgency, colloquial expression

\textbf{Tags (JA)}: 即時文法, 口語, 必要性表現, 時間切迫表現, 口語的表現\\
\textbf{Tags (EN)}: immediate grammar, spoken, expression of necessity, urgency expression, colloquial expression

\textbf{Adjusted Expression (JA)}:\\
もう、(行:い)かなくてはなりません。\\
\textbf{Adjusted Expression (EN)}:\\
I have to go now.

\textbf{Example (JA)}:\\
もう、(行:い)かなくっちゃ。(次:つぎ)の(予定:よてい)があるからね。\\
\textbf{Example (EN)}:\\
I have to go now. I have other plans.

\textbf{Notes (JA)}:\\
「もう、行かなくっちゃ。」は、今すぐに行かなければならないことを表す口語的な表現。「もう」は、この時点でそうしなければならないという気持ちを表し、「行かなくっちゃ」は「I have to go」の意味です。

\textbf{Notes (EN)}:\\
'Mou, ikanakucha.' is a colloquial expression used to indicate that you have to leave now. 'Mou' expresses the feeling that it is already time to go, and 'ikanakucha' means 'I have to go'.

\newpage

\section*{736. Is there a good way to do it?}

\textbf{Date}: 2026.04.16 (Thu)\\
\textbf{Index}: nanika ii houhou wa\\
\textbf{Expression (JA)}: (何:なに)かいい(方法:ほうほう)は？\\
\textbf{Romanization}: Nanika ii hōhō wa?\\
\textbf{Expression (EN)}: Is there a good way to do it?

\textbf{Variants (JA)}: 何かいい方法は？ / 何かいい方法ある？ / 何かいい方法ない？ / 何かいいやり方は？\\
\textbf{Variants (EN)}: Is there a good way to do it? / Is there any good way to do it? / Do you have any good ideas?

\textbf{Intent (JA)}: 方法の質問, 解決策の共有, 口語的な表現\\
\textbf{Intent (EN)}: asking for a method, sharing of solutions, colloquial expression

\textbf{Tags (JA)}: 即時文法, 口語, 方法質問表現, 解決策共有表現, 口語的表現\\
\textbf{Tags (EN)}: immediate grammar, spoken, method-question expression, solution-sharing expression, colloquial expression

\textbf{Adjusted Expression (JA)}:\\
(何:なに)か(良:い)い(方法:ほうほう)はありますか？\\
\textbf{Adjusted Expression (EN)}:\\
Is there a good way to do it?

\textbf{Example (JA)}:\\
この(問題:もんだい)、どうにかしたいんだけど、(何:なに)かいい(方法:ほうほう)は？\\
\textbf{Example (EN)}:\\
I want to do something about this problem. Is there any good way to do it?

\textbf{Notes (JA)}:\\
(「何かいい方法は？」は、何かをするための良い方法や手段を尋ねる口語的な表現。「何か」は「something」の意味で、「いい方法」は「good way」の意味です。口語では、このように短く言うことが多いですが、正式な表現では「何か良い方法はありますか？」のように言うこともできます。しかし、長い言い方は丁寧ではあるかもしれませんが、会話の流れが悪くなったり、自然な感じが失われたりすることもあるので、状況に応じて使い分けることが大切です。

\textbf{Notes (EN)}:\\
'Nanika ii hōhō wa?' is a colloquial expression used to ask for a good method or way to do something. 'Nanika' means 'something', and 'ii houhou' means 'good way'. In casual conversation, it's common to say it this way, but in a more formal expression, you could say 'Nanika yoi houhou wa arimasu ka?'. However, while the longer form may be more polite, it can also disrupt the flow of conversation and feel less natural, so it's important to choose which form to use based on the situation.

\newpage

\section*{735. Related to that, ...}

\textbf{Date}: 2026.04.15 (Wed)\\
\textbf{Index}: chinaminichotto\\
\textbf{Expression (JA)}: ちなみに、...\\
\textbf{Romanization}: Chinami ni, ...\\
\textbf{Expression (EN)}: Related to that, ...

\textbf{Variants (JA)}: ちなみに、... / ちなみにね、... / ちなみにさ、...\\
\textbf{Variants (EN)}: Related to that, ... / This is a bit related, but...

\textbf{Intent (JA)}: 関連情報の追加, 補足の共有, 口語的な表現\\
\textbf{Intent (EN)}: adding related information, sharing additional information, colloquial expression

\textbf{Tags (JA)}: 即時文法, 口語, 関連情報の追加, 補足の追加, 口語的表現\\
\textbf{Tags (EN)}: immediate grammar, spoken, expression for adding related information, colloquial expression

\textbf{Adjusted Expression (JA)}:\\
(少:すこ)し、(関連:かんれん)する(話:はなし)なんですが、...\\
\textbf{Adjusted Expression (EN)}:\\
This is a bit related, but...

\textbf{Example (JA)}:\\
A「今日はバレンタインデーですね！」 B「そうですね。ちなみに、もうチョコ買いました？」\\
\textbf{Example (EN)}:\\
A: Today is Valentine's Day, isn't it? B: Yes, it is. Related to that, have you bought chocolate yet?

\textbf{Notes (JA)}:\\
「ちなみに」は、話題転換の言葉のように見えて、実際には完全に別の話題には使いにくい表現です。前の話と少しでもつながりがないと不自然に聞こえます。つまり「ちなみに」は、関係ない話を始める合図ではなく、前の話に関連する情報を添える合図です。

\textbf{Notes (EN)}:\\
The phrase 'chinami ni' may look like a way to change the topic, but it is actually difficult to use it for a completely unrelated topic. It sounds unnatural if there is no connection to the previous topic. In other words, 'chinami ni' is not a signal for starting an unrelated topic, but a signal for adding information related to the previous topic.

\newpage

\section*{734. Completely specialized in ...}

\textbf{Date}: 2026.04.14 (Tue)\\
\textbf{Index}: ni kanzen ni tokka shiteru\\
\textbf{Expression (JA)}: ...に(完全:かんぜん)に(特化:とっか)してる\\
\textbf{Romanization}: ...ni kanzen ni tokka shiteru\\
\textbf{Expression (EN)}: Completely specialized in ...

\textbf{Variants (JA)}: ...に完全に特化してる / ...のためだけに作られている\\
\textbf{Variants (EN)}: Completely specialized in ... / Made just for ... / Totally focused on ...

\textbf{Intent (JA)}: 専門性の表現, 強調の共有, 口語的な表現\\
\textbf{Intent (EN)}: expression of specialization, sharing of emphasis, colloquial expression

\textbf{Tags (JA)}: 即時文法, 口語, 専門性表現, 強調共有表現, 口語的表現\\
\textbf{Tags (EN)}: immediate grammar, spoken, expression of specialization, expression of emphasis sharing, colloquial expression

\textbf{Adjusted Expression (JA)}:\\
...に(完全:かんぜん)に(特化:とっか)しています。\\
\textbf{Adjusted Expression (EN)}:\\
...is completely specialized in ...

\textbf{Example (JA)}:\\
この(サイト:さいと)はこの(分野:ぶんや)に(完全:かんぜん)に(特化:とっか)してるから、(他:ほか)のことはあまり(説明:せつめい)していないんだよね。\\
\textbf{Example (EN)}:\\
This site is completely specialized in this field, so it doesn't explain much about other things.

\textbf{Notes (JA)}:\\
「...に完全に特化してる」は、ある分野や領域に対して非常に専門的であることを表す口語的な表現。「完全に」は「completely」の意味で、「特化してる」は「is specialized」の意味です。

\textbf{Notes (EN)}:\\
'...ni kanzen ni tokka shiteru' is a colloquial expression used to indicate that something is very specialized in a certain field or area. 'Kanzen ni' means 'completely', and 'tokka shiteru' means 'is specialized'.

\newpage

\section*{733. It would be nice if we could do that, but...}

\textbf{Date}: 2026.04.13 (Mon)\\
\textbf{Index}: soudekirebaiidakedo\\
\textbf{Expression (JA)}: そうできればいいんだけど...\\
\textbf{Romanization}: Sō dekireba ii n da kedo.\\
\textbf{Expression (EN)}: It would be nice if we could do that, but...

\textbf{Variants (JA)}: そうできればいいんだけど。 / そうできるといいんだけど。\\
\textbf{Variants (EN)}: It would be nice if we could do that, but... / It would be great if we could do that, but...

\textbf{Intent (JA)}: 願望の表現, 可能性の共有, 口語的な表現\\
\textbf{Intent (EN)}: expression of desire, sharing of possibility, colloquial expression

\textbf{Tags (JA)}: 即時文法, 口語, 願望表現, 可能性共有表現, 口語的表現\\
\textbf{Tags (EN)}: immediate grammar, spoken, expression of desire, expression of possibility sharing, colloquial expression

\textbf{Adjusted Expression (JA)}:\\
そうできれば(助:たす)かるのですが、\\
\textbf{Adjusted Expression (EN)}:\\
It would be helpful if we could do that, but...

\textbf{Example (JA)}:\\
A「この(順番:じゅんばん)、(入:い)れ(替:か)えた(方:ほう)が、(効率:こうりつ)よくないですか？」 B「そうできればいいんだけど、この(日:ひ)はまだ(材料:ざいりょう)が(届:とど)いていなくてね」\\
\textbf{Example (EN)}:\\
A: Wouldn't it be more efficient if we switched the order? B: It would be nice if we could do that, but the materials haven't arrived yet by that day.

\textbf{Notes (JA)}:\\
「そうできればいいんだけど」は、何かをしたいけれども、それが難しい事情があることを示す表現。「そうできるといいんだけど」とも(言:い)います。

\textbf{Notes (EN)}:\\
'Sou dekireba ii n da kedo' is a colloquial expression used when you would like to do something, but there are circumstances that make it difficult. It can also be said as 'Sou dekiru to ii n da kedo.'

\newpage

\section*{732. It sounds like you can do it, but...}

\textbf{Date}: 2026.04.12 (Sun)\\
\textbf{Index}: samodekiruyounakotooitterukedo\\
\textbf{Expression (JA)}: さもできるようなことを(言:い)ってるけど、\\
\textbf{Romanization}: Samo dekiru yō na koto o itteru kedo,\\
\textbf{Expression (EN)}: It sounds like you can do it, but...

\textbf{Variants (JA)}: さもできるようなことを言ってるけど、 / さもできそうなことを言ってるけど、\\
\textbf{Variants (EN)}: It sounds like you can do it, but... / It sounds like you could do it, but...

\textbf{Intent (JA)}: 能力の表現, 疑念の共有, 口語的な表現\\
\textbf{Intent (EN)}: expression of ability, sharing of doubt, colloquial expression

\textbf{Tags (JA)}: 即時文法, 口語, 能力表現, 疑念共有表現, 口語的表現\\
\textbf{Tags (EN)}: immediate grammar, spoken, expression of ability, expression of doubt sharing, colloquial expression

\textbf{Adjusted Expression (JA)}:\\
さもできるようなことをおっしゃっているのですが、\\
\textbf{Adjusted Expression (EN)}:\\
(It sounds like you can do it,) but...

\textbf{Example (JA)}:\\
A「この(プロジェクト:ぷろじぇくと)、(私:わたし)が(担当:たんとう)することになったんだけど、さもできるようなことを(言:い)ってるけど、(実:じつ)はちょっと(不安:ふあん)なんだよね」 B「(大丈夫:だいじょうぶ)ですよ！あなたならきっとできます！」\\
\textbf{Example (EN)}:\\
A: I was assigned to this project, and it sounds like I can do it, but I'm actually a bit anxious about it. B: Don't worry! I'm sure you can do it! 

\textbf{Notes (JA)}:\\
「さもできるようなことを言ってるけど、」は、相手が何かをできるように聞こえることを言っているけれども、実際には疑念や不安があることを表す口語的な表現。「さも」は「as if」の意味で、「できるようなことを言ってる」は「it sounds like you can do it」の意味です。

\textbf{Notes (EN)}:\\
'Samo dekiru yō na koto o itteru kedo,' is a colloquial expression used to indicate that someone is saying something that sounds like they can do it, but in reality, there is doubt or anxiety about it. 'Samo' means 'as if', and 'dekiru yō na koto o itteru' means 'it sounds like you can do it'.

\newpage

\section*{731. Ah, this was the way to do it.}

\textbf{Date}: 2026.04.11 (Sat)\\
\textbf{Index}: akoredeiidanaa\\
\textbf{Expression (JA)}: あ、これでいいんだぁ。\\
\textbf{Romanization}: Ah, kore de ii n da a.\\
\textbf{Expression (EN)}: Ah, this was the way to do it.

\textbf{Variants (JA)}: あ、これでいいんだぁ。 / あ、こうすればいいんだぁ。\\
\textbf{Variants (EN)}: Ah, this was the way to do it. / Ah, so this is how I should do it.

\textbf{Intent (JA)}: 発見の表現, 理解の共有, 口語的な表現\\
\textbf{Intent (EN)}: expression of discovery, sharing of understanding, colloquial expression

\textbf{Tags (JA)}: 即時文法, 口語, 発見表現, 理解共有表現, 口語的表現\\
\textbf{Tags (EN)}: immediate grammar, spoken, expression of discovery, expression of understanding sharing, colloquial expression

\textbf{Adjusted Expression (JA)}:\\
あ、これでいいのですね。\\
\textbf{Adjusted Expression (EN)}:\\
Ah, so this is how I should do it.

\textbf{Example (JA)}:\\
(問題:もんだい)を(解決:かいけつ)する方法がわからなくて、いろいろ(試行錯誤:しこうさくご)してたんだけど、あ、これでいいんだぁって(気:き)づいたよ！\\
\textbf{Example (EN)}:\\
I was struggling to find a solution to the problem and trying various things, but then I realized, ah, this was the way to do it!

\textbf{Notes (JA)}:\\
「あ、これでいいんだぁ」は、何かを発見したときや理解したときに使う口語的な表現。「あ」は驚きや発見を表す感嘆詞で、「これでいいんだぁ」は「this was the way to do it」の意味です。いろいろやってみて、やっと答がわかったときや、何かの方法がわかったときに使います。

\textbf{Notes (EN)}:\\
'Ah, kore de ii n da a.' is a colloquial expression used when you discover or understand something. 'Ah' is an interjection that expresses surprise or discovery, and 'kore de ii n da a' means 'this was the way to do it'. It is used when you have been trying various things and finally figure out the answer or understand how to do something.

\newpage

\section*{730. That's what I wanted to say.}

\textbf{Date}: 2026.04.10 (Fri)\\
\textbf{Index}: soregaiitakatta\\
\textbf{Expression (JA)}: それが(言:い)いたかった。\\
\textbf{Romanization}: Sore ga iitakatta.\\
\textbf{Expression (EN)}: That's what I wanted to say.

\textbf{Variants (JA)}: それが言いたかった。 / それを言いたかった。\\
\textbf{Variants (EN)}: That's what I wanted to say. / That's what I meant to say.

\textbf{Intent (JA)}: 代弁への同意, 語化の回収, 共感\\
\textbf{Intent (EN)}: agreement with someone else's wording, recovery of verbalization, empathy

\textbf{Tags (JA)}: 即時文法, 口語, 代弁表現, 同意表現, 語探し表現\\
\textbf{Tags (EN)}: immediate grammar, spoken, proxy-expression, agreement expression, word-search expression

\textbf{Adjusted Expression (JA)}:\\
まさにそれが(言:い)いたかったのです。\\
\textbf{Adjusted Expression (EN)}:\\
That is exactly what I wanted to say.

\textbf{Example (JA)}:\\
A「あの(レストラン:れすとらん)、(味:あじ)はいいし、(値段:ねだん)も(手頃:てごろ)で、すばらしいんだけど、どこかひとつ、ええと、(何:なん)だろ、(何:なん)て(言:い)ったらいいんだろう...」 B「(要:よう)するに、(雑:ざつ)なんだね」 A「うん、それが(言:い)いたかったんだよ！ (雑:ざつ)、あの(レストラン:れすとらん)！」\\
\textbf{Example (EN)}:\\
A: That restaurant has good food and reasonable prices, and it's great, but there is one thing, um, I do not know how to put it... B: So, basically, it feels sloppy, right? A: Yeah, that's what I wanted to say! Sloppy, that restaurant!

\textbf{Notes (JA)}:\\
「それが言いたかった」は、自分ではうまく言えずに困っていたことを、相手が先に適切な表現で言ってくれたときの口語的な表現です。相手のことばに同意しながら、それこそ自分が言いたかったことだ、という気持ちをその場で伝える言い方です。

\textbf{Notes (EN)}:\\
'Sore ga iitakatta' is a colloquial expression used when you were struggling to find the right words and another person says exactly what you were trying to express. It shows immediate agreement with the other person's wording while also indicating that this was precisely what you had wanted to say.

\newpage

\section*{729. What do you call that?}

\textbf{Date}: 2026.04.09 (Thu)\\
\textbf{Index}: sorettenantteiunno\\
\textbf{Expression (JA)}: それって(何:なん)っていうの？\\
\textbf{Romanization}: Sore tte nan tte iu no?\\
\textbf{Expression (EN)}: What do you call that?

\textbf{Variants (JA)}: それって何っていうの？ / それは何と呼ばれているの？\\
\textbf{Variants (EN)}: What do you call that? / What is that called? / What is that referred to as?

\textbf{Intent (JA)}: 名称の質問, 語探し, 会話ストラテジー\\
\textbf{Intent (EN)}: asking for a name, word search, conversation strategy

\textbf{Tags (JA)}: 即時文法, 口語, 語探し表現, 会話ストラテジー\\
\textbf{Tags (EN)}: immediate grammar, spoken, word-search expression, conversation strategy

\textbf{Adjusted Expression (JA)}:\\
それは(何:なん)と(呼:よ)ばれているのですか？\\
\textbf{Adjusted Expression (EN)}:\\
What is that called?

\textbf{Example (JA)}:\\
(昨日:きのう)、(開店前:かいてんまえ)、(店:みせ)の(前:まえ)にみんな(並:なら)んでて、それって(何:なん)て(言:い)うの？(一番乗:いちばんの)りっていうの？すごい(行列:ぎょうれつ)でね、(驚:おどろ)いたよ。\\
\textbf{Example (EN)}:\\
Yesterday, before the store opened, there were people lined up in front of the store, and I was like, what do you call that? Is it called 'first in line'? It was such a long line that I was surprised.

\textbf{Notes (JA)}:\\
「それって何っていうの？」は、普通は、物や概念の名前を尋ねる表現ですが、適切な呼び方や名詞がわからない、あるいは自信がないときに、その場で質問しながら話を進める方法です。これは、知らないから止まるのではなく、わからなさを抱えたまま運用するやり方です。しかも、それは外国語学習者だけの工夫ではなく、日本語母語話者もふつうに使います。「何っていうの？」は「what is it called?」の意味です。「一番乗り」というのは、「first in line」の意味で、何かの先頭にいることを指します。しかし、大体において、一番乗りで買っても特別に安くなるわけではないのに、一番乗りを目指す人が多い。また、何があるのかわからないけれども、大勢並んでいるから、並んでいる人も時々います。

\textbf{Notes (EN)}:\\
'Sore tte nan tte iu no?' is a strategic expression used when asking for the name of something or a concept, especially when the speaker is unsure of the appropriate term or noun to use. This is a way to continue the conversation while asking questions on the spot, rather than stopping because you don't know something. Moreover, this is not just a strategy used by language learners; native Japanese speakers also use it commonly. 'Nan tte iu no?' means 'what is it called?'. The phrase 'ichiban nori' means 'first in line', referring to being at the front of a line. However, in many cases, being first in line does not necessarily mean getting a special discount, yet many people aim to be first in line. Additionally, sometimes people will join a line simply because there are many people already lined up, even if they don't know what is being offered.

\newpage

\section*{728. I would like you to...}

\textbf{Date}: 2026.04.08 (Wed)\\
\textbf{Index}: tehoshiindesukedo\\
\textbf{Expression (JA)}: ...てほしいんですけど、\\
\textbf{Romanization}: ...te hoshii n desu kedo,\\
\textbf{Expression (EN)}: I would like you to...

\textbf{Variants (JA)}: ...てほしいんですけど、 / ...してほしいんですが、\\
\textbf{Variants (EN)}: I would like you to... / I would like you to do...

\textbf{Intent (JA)}: 依頼の表現, 丁寧な表現, 口語的な表現\\
\textbf{Intent (EN)}: expression of request, polite expression, colloquial expression

\textbf{Tags (JA)}: 即時文法, 口語, 依頼表現, 丁寧表現, 口語的表現\\
\textbf{Tags (EN)}: immediate grammar, spoken, expression of request, polite expression, colloquial expression

\textbf{Adjusted Expression (JA)}:\\
...ていただきたいのですが、\\
\textbf{Adjusted Expression (EN)}:\\
...I would like you to do..., but

\textbf{Example (JA)}:\\
A 「これ、(手伝:てつだ)ってほしいんですけど」 B「(今:いま)、ちょっと(忙:いそが)しいんで、(後:あと)でいいですか？」\\
\textbf{Example (EN)}:\\
A: I would like you to help me with this. B: I'm a bit busy right now, can it wait until later?

\textbf{Notes (JA)}:\\
...てほしいんですけど、は、相手に何かをしてほしいときに使う口語的な表現。丁寧な依頼の形で、「...してほしいんですが、」とも言えます。

\textbf{Notes (EN)}:\\
'...te hoshii n desu kedo,' is a colloquial expression used when you want someone to do something for you. It is a polite form of request and can also be said as '...shite hoshii n desu ga,'.

\newpage

\section*{727. What did you just say?}

\textbf{Date}: 2026.04.07 (Tue)\\
\textbf{Index}: imanannteitta\\
\textbf{Expression (JA)}: (今:いま)、(何:なん)て(言:い)った？\\
\textbf{Romanization}: Ima, nan te itta?\\
\textbf{Expression (EN)}: What did you just say?

\textbf{Variants (JA)}: 今、何て言った？ / 今、何て言いました？\\
\textbf{Variants (EN)}: What did you just say? / What did you just say? (polite)

\textbf{Intent (JA)}: 確認の表現, 驚きの共有, 口語的な表現\\
\textbf{Intent (EN)}: expression of confirmation, sharing of surprise, colloquial expression

\textbf{Tags (JA)}: 即時文法, 口語, 確認表現, 驚き共有表現, 口語的表現\\
\textbf{Tags (EN)}: immediate grammar, spoken, expression of confirmation, expression of surprise sharing, colloquial expression

\textbf{Adjusted Expression (JA)}:\\
(今:いま)、(何:なん)とおっしゃいましたか？\\
\textbf{Adjusted Expression (EN)}:\\
What did you just say? (polite)

\textbf{Example (JA)}:\\
A「この(映画:えいが)、(面白:おもしろ)かった！」 B「(今:いま)、(何:なん)て(言:い)った？(本当:ほんとう)にそう(思:おも)うの？」\\
\textbf{Example (EN)}:\\
A: This movie was really good! B: What did you just say? Do you really think so?

\textbf{Notes (JA)}:\\
「今、何て言った？」は、相手の発言に対して驚いたり、確認したりする口語的な表現。「今」は「just now」の意味で、「何て言った？」は「what did you say?」の意味です。

\textbf{Notes (EN)}:\\
'Ima, nan te itta?' is a colloquial expression used to express surprise or seek confirmation about what someone just said. 'Ima' means 'just now', and 'nan te itta?' means 'what did you say?'.

\newpage

\section*{726. That's a relief.}

\textbf{Date}: 2026.04.06 (Mon)\\
\textbf{Index}: sorewananiyori\\
\textbf{Expression (JA)}: それはなにより\\
\textbf{Romanization}: Sore wa nani yori.\\
\textbf{Expression (EN)}: That's a relief.

\textbf{Variants (JA)}: それはなにより / それはよかった\\
\textbf{Variants (EN)}: That's a relief. / That's good to hear.

\textbf{Intent (JA)}: 安心の表現, 喜びの共有, 口語的な表現\\
\textbf{Intent (EN)}: expression of relief, sharing of joy, colloquial expression

\textbf{Tags (JA)}: 即時文法, 口語, 安心表現, 喜び共有表現, 口語的表現\\
\textbf{Tags (EN)}: immediate grammar, spoken, expression of relief, expression of joy sharing, colloquial expression

\textbf{Adjusted Expression (JA)}:\\
それは(無事:ぶじ)、(終:お)わってよかったですね。\\
\textbf{Adjusted Expression (EN)}:\\
That's good to hear that it ended safely.

\textbf{Example (JA)}:\\
(病気:びょうき)が(治:なお)った。それはなにより。\\
\textbf{Example (EN)}:\\
You recovered from your illness. That's a relief.

\textbf{Notes (JA)}:\\
(「それはなにより」は、相手の状況が良くなったが、最終的によい結果で終わって、安心する気持ちの表現。「なにより」は「何より」の意味で、他のどんなことよりも重要であることを強調するために使われます。

\textbf{Notes (EN)}:\\
'Sore wa nani yori.' is a colloquial expression used to convey a feeling of relief when someone's situation has improved or when something has ended well. 'Nani yori' means 'more than anything else', emphasizing that the outcome is more important than anything else.

\newpage

\section*{725. Well, let's do it this way.}

\textbf{Date}: 2026.04.05 (Sun)\\
\textbf{Index}: jaikoushiyou\\
\textbf{Expression (JA)}: じゃ、こうしよう。\\
\textbf{Romanization}: Ja, kou shiyou.\\
\textbf{Expression (EN)}: Well, let's do it this way.

\textbf{Variants (JA)}: じゃ、こうしよう。 / じゃ、これでいこう。\\
\textbf{Variants (EN)}: Well, let's do it this way. / Well, let's go with this.

\textbf{Intent (JA)}: 提案の表現, 決定の共有, 口語的な表現\\
\textbf{Intent (EN)}: expression of suggestion, sharing of decision, colloquial expression

\textbf{Tags (JA)}: 即時文法, 口語, 提案表現, 決定共有表現, 口語的表現\\
\textbf{Tags (EN)}: immediate grammar, spoken, expression of suggestion, expression of decision sharing, colloquial expression

\textbf{Adjusted Expression (JA)}:\\
では、このようにしましょう。\\
\textbf{Adjusted Expression (EN)}:\\
Well, let's do it this way.

\textbf{Example (JA)}:\\
A「この(問題:もんだい)、どう(解決:かいけつ)する？」 B「じゃ、こうしよう。まずは(原因:げんいん)を(調査:ちょうさ)してみて、それから(対策:たいさく)を(考:かんが)えよう」\\
\textbf{Example (EN)}:\\
A: How should we solve this problem? B: Well, let's do it this way. First, let's investigate the cause, and then we can think about the countermeasures.

\textbf{Notes (JA)}:\\
「じゃ、こうしよう」は、何かを提案したり、決定を共有したりする口語的な表現。「じゃ」は「well」の意味で、「こうしよう」は「let's do it this way」の意味です。

\textbf{Notes (EN)}:\\
'Ja, kou shiyou.' is a colloquial expression used to make a suggestion or share a decision. 'Ja' means 'well', and 'kou shiyou' means 'let's do it this way'.

\newpage

\section*{724. Eh, no way...}

\textbf{Date}: 2026.04.04 (Sat)\\
\textbf{Index}: eesonnna\\
\textbf{Expression (JA)}: えぇ、そんなぁ...\\
\textbf{Romanization}: Ee, sonna...\\
\textbf{Expression (EN)}: Eh, no way...

\textbf{Variants (JA)}: えぇ、そんなぁ... / えぇ、そんな...\\
\textbf{Variants (EN)}: Eh, no way... / Eh, that's...

\textbf{Intent (JA)}: 驚きの表現, 否定の共有, 口語的な表現\\
\textbf{Intent (EN)}: expression of surprise, sharing of negation, colloquial expression

\textbf{Tags (JA)}: 即時文法, 口語, 驚き表現, 否定共有表現, 口語的表現\\
\textbf{Tags (EN)}: immediate grammar, spoken, expression of surprise, expression of negation sharing, colloquial expression

\textbf{Adjusted Expression (JA)}:\\
えぇ、それはちょっと...\\
\textbf{Adjusted Expression (EN)}:\\
Eh, that's a bit...

\textbf{Example (JA)}:\\
A「ここの(支払:しはら)い、よろしく！」 B「えぇ、そんなぁ...」\\
\textbf{Example (EN)}:\\
A: You're paying for this one. B: Oh, no way...

\textbf{Notes (JA)}:\\
「えぇ、そんなぁ...」は、何かに対して驚いたり、否定的な感情を共有したりする口語的な表現。「えぇ」は驚きを表す感嘆詞で、「そんなぁ」は「that's...」の意味で、何かが予想外だったり、受け入れがたいと感じたりするときに使われます。

\textbf{Notes (EN)}:\\
'Ee, sonna...' is a colloquial expression used to convey surprise or share a negative emotion about something. 'Ee' is an interjection that expresses surprise, and 'sonna' means 'that's...' and is used when something is unexpected or feels unacceptable.

\newpage

\section*{723. I don't like it.}

\textbf{Date}: 2026.04.03 (Fri)\\
\textbf{Index}: iyadanaa\\
\textbf{Expression (JA)}: いやだなぁ。\\
\textbf{Romanization}: Iya da naa.\\
\textbf{Expression (EN)}: I don't like it.

\textbf{Variants (JA)}: いやだなぁ。 / いやだな。\\
\textbf{Variants (EN)}: I don't like it. / I don't want to.

\textbf{Intent (JA)}: 否定の表現, 感情の共有, 口語的な表現\\
\textbf{Intent (EN)}: expression of negation, sharing of emotion, colloquial expression

\textbf{Tags (JA)}: 即時文法, 口語, 否定表現, 感情共有表現, 口語的表現\\
\textbf{Tags (EN)}: immediate grammar, spoken, expression of negation, expression of emotion sharing, colloquial expression

\textbf{Adjusted Expression (JA)}:\\
いや。それはちょっと\\
\textbf{Adjusted Expression (EN)}:\\
No, that's a bit...

\textbf{Example (JA)}:\\
A「(明日:あした)、(試験:しけん)だよね！」 B「そうなんだよ、いやだなぁ」\\
\textbf{Example (EN)}:\\
A: The exam is tomorrow, right? B: Yeah, I know. I don't like it.

\textbf{Notes (JA)}:\\
「いやだなぁ」は、何かを嫌がる気持ちを伝える口語的な表現。「いや」は「dislike」の意味で、「だなぁ」は感情を強調するための終助詞です。

\textbf{Notes (EN)}:\\
'Iya da naa.' is a colloquial expression used to convey a feeling of dislike or reluctance towards something. 'Iya' means 'dislike', and 'da naa' is a sentence-ending particle that emphasizes the emotion behind the statement.

\newpage

\section*{722. So, you see,}

\textbf{Date}: 2026.04.02 (Thu)\\
\textbf{Index}: sousurudesune\\
\textbf{Expression (JA)}: そうするとですね、\\
\textbf{Romanization}: Sou suru desu ne,\\
\textbf{Expression (EN)}: So, you see,

\textbf{Variants (JA)}: そうするとですね、 / そうすると、\\
\textbf{Variants (EN)}: So, you see, / So,

\textbf{Intent (JA)}: 説明の導入, 話の展開, 口語的な表現\\
\textbf{Intent (EN)}: introduction of explanation, development of story, colloquial expression

\textbf{Tags (JA)}: 即時文法, 口語, 説明導入表現, 話展開表現, 口語的表現\\
\textbf{Tags (EN)}: immediate grammar, spoken, expression of explanation introduction, expression of story development, colloquial expression

\textbf{Adjusted Expression (JA)}:\\
(そう:そう)するとですね、\\
\textbf{Adjusted Expression (EN)}:\\
(So,) you see,

\textbf{Example (JA)}:\\
そうするとですね、この(方法:ほうほう)が(最適:さいてき)だと(言:い)うことですか。\\
\textbf{Example (EN)}:\\
So, you see, are you saying that this method is the best?

\textbf{Notes (JA)}:\\
「そうするとですね、...」は、説明を始めるときや話を展開するときに使う口語的な表現。「そうする」は「to do so」の意味で、「ですね」は話し手が聞き手に同意を求めるための終助詞です。

\textbf{Notes (EN)}:\\
'Sou suru desu ne, ...' is a colloquial expression used when starting an explanation or developing a story. 'Sō suru' means 'to do so', and 'desu ne' is a sentence-ending particle used to seek agreement from the listener.

\newpage

\section*{721. Anyway,}

\textbf{Date}: 2026.04.01 (Wed)\\
\textbf{Index}: nanihatomoare\\
\textbf{Expression (JA)}: なにはともあれ、\\
\textbf{Romanization}: Nani wa tomo are,\\
\textbf{Expression (EN)}: Anyway,

\textbf{Variants (JA)}: なにはともあれ、 / とにかく、\\
\textbf{Variants (EN)}: Anyway, / In any case,

\textbf{Intent (JA)}: 話題の転換, 結論の提示, 口語的な表現\\
\textbf{Intent (EN)}: topic shift, stating a conclusion, colloquial expression

\textbf{Tags (JA)}: 即時文法, 口語, 話題転換表現, 結論提示表現, 口語的表現\\
\textbf{Tags (EN)}: immediate grammar, spoken, expression of topic shift, expression of conclusion stating, colloquial expression

\textbf{Adjusted Expression (JA)}:\\
いづれにいたしましても、\\
\textbf{Adjusted Expression (EN)}:\\
In any case,

\textbf{Example (JA)}:\\
なにはともあれ、(健康:けんこう)が(第一:だいいち)だよね。\\
\textbf{Example (EN)}:\\
Anyway, health is the most important thing, right?

\textbf{Notes (JA)}:\\
「なにはともあれ」は、話題を転換したり、結論を提示したりするときに使う口語的な表現。「なにはともあれ」は「とにかく」と同じ意味で、どんなことがあっても重要なことを強調するために使われます。

\textbf{Notes (EN)}:\\
'Nani wa tomo are' is a colloquial expression used to shift the topic or state a conclusion. It has the same meaning as 'tonikaku', which emphasizes that something is important regardless of what happens.

\newpage

\section*{720. Can I try it?}

\textbf{Date}: 2026.03.31 (Tue)\\
\textbf{Index}: yattetemiteii\\
\textbf{Expression (JA)}: やってみていい？\\
\textbf{Romanization}: Yatte mite ii?\\
\textbf{Expression (EN)}: Can I try it?

\textbf{Variants (JA)}: やってみてもいい？ / やってみてもいいですか？\\
\textbf{Variants (EN)}: Can I try it? / Is it okay if I try it?

\textbf{Intent (JA)}: 許可の表現, 試行の共有, 口語的な表現\\
\textbf{Intent (EN)}: expression of permission, sharing of attempt, colloquial expression

\textbf{Tags (JA)}: 即時文法, 口語, 許可表現, 試行共有表現, 口語的表現\\
\textbf{Tags (EN)}: immediate grammar, spoken, expression of permission, expression of attempt sharing, colloquial expression

\textbf{Adjusted Expression (JA)}:\\
(やっ:やっ)て(み:み)てもいいですか？\\
\textbf{Adjusted Expression (EN)}:\\
Is it okay if I try it?

\textbf{Example (JA)}:\\
A「この(ゲーム:げーむ)、(面白:おもしろ)そうだね。やってみていい？」 B「うん、いいよ！」\\
\textbf{Example (EN)}:\\
A: This game looks fun. Can I try it? B: Yeah, go ahead!

\textbf{Notes (JA)}:\\
「やってみていい？」は、何かを試してみたいときに使う口語的な表現。「やってみる」は「to try doing something」の意味で、「いい？」は「is it okay?」の意味です。

\textbf{Notes (EN)}:\\
'Yatte mite ii?' is a colloquial expression used when you want to try something. 'Yatte miru' means 'to try doing something', and 'ii?' means 'is it okay?'.

\newpage

\section*{719. I knew it...}

\textbf{Date}: 2026.03.30 (Mon)\\
\textbf{Index}: yappari\\
\textbf{Expression (JA)}: やっぱりぃ...\\
\textbf{Romanization}: Yappari...\\
\textbf{Expression (EN)}: I knew it...

\textbf{Variants (JA)}: やっぱり... / やはり...\\
\textbf{Variants (EN)}: I knew it... / As I thought...

\textbf{Intent (JA)}: 予想の表現, 驚きの共有, 口語的な表現\\
\textbf{Intent (EN)}: expression of expectation, sharing of surprise, colloquial expression

\textbf{Tags (JA)}: 即時文法, 口語, 予想表現, 驚き共有表現, 口語的表現\\
\textbf{Tags (EN)}: immediate grammar, spoken, expression of expectation, expression of surprise sharing, colloquial expression

\textbf{Adjusted Expression (JA)}:\\
(やはり:やはり)...\\
\textbf{Adjusted Expression (EN)}:\\
(As I thought)...

\textbf{Example (JA)}:\\
A「(冷蔵庫:れいぞうこ)の(アイスクリーム:あいすくりーむ)、(食:た)べたの、お(父:とう)さん？」 B「うん、そうだけど...」A「やっぱりぃ...」\\
\textbf{Example (EN)}:\\
A: Dad, did you eat the ice cream in the fridge? B: Yeah, I did... A: I knew it...

\textbf{Notes (JA)}:\\
「やっぱりぃ...」は、何かが予想通りだったときや、驚いたときに使う口語的な表現。「やっぱり」は「やはり」の口語的な形で、「予想通り」の意味です。伸ばした「ぃ」は感情を強調するためのものです。

\textbf{Notes (EN)}:\\
'Yappari...' is a colloquial expression used when something turns out as expected or when the speaker is surprised. 'Yappari' is a colloquial form of 'yahari', which means 'as expected'. The elongated 'i' emphasizes the emotion behind the statement.

\newpage

\section*{718. How about going for lunch?}

\textbf{Date}: 2026.03.29 (Sun)\\
\textbf{Index}: ranchikanai\\
\textbf{Expression (JA)}: (ランチ:らんち)、(行:い)かない？\\
\textbf{Romanization}: Ranchi, ikanai?\\
\textbf{Expression (EN)}: How about going for lunch?

\textbf{Variants (JA)}: (ランチ:らんち)、(食:た)べ(行:い)かない？ / (ランチ:らんち)に(行:い)かない？ / (昼食:ちゅうしょく)に(行:い)かない？\\
\textbf{Variants (EN)}: Do you want to go for lunch? / Shall we go for lunch? / Would you like to go for lunch?

\textbf{Intent (JA)}: 誘いの表現, 食事の共有, 口語的な表現\\
\textbf{Intent (EN)}: expression of invitation, sharing of meal, colloquial expression

\textbf{Tags (JA)}: 即時文法, 誘い表現, 食事共有表現, 口語的表現\\
\textbf{Tags (EN)}: immediate grammar, expression of invitation, expression of meal sharing, colloquial expression

\textbf{Adjusted Expression (JA)}:\\
(ランチ:らんち)に(行:い)きませんか？\\
\textbf{Adjusted Expression (EN)}:\\
Would you like to go for lunch?

\textbf{Example (JA)}:\\
A「お(腹:なか)すいたね。ランチ、(:)行かない？」 B「うん、いいね！」\\
\textbf{Example (EN)}:\\
A: I'm hungry. Do you want to go for lunch? B: Yeah, sounds good!

\textbf{Notes (JA)}:\\
「ランチ、行かない？」は、相手をランチに誘う口語的な表現。「ランチ」は「lunch」の意味で、「行かない？」は「do you want to go?」の意味です。

\textbf{Notes (EN)}:\\
'Ranchi, ikanai?' is a colloquial expression used to invite someone to go for lunch. 'Ranchi' means 'lunch', and 'ikanai?' means 'do you want to go?'.

\newpage

\section*{717. How did you do it?}

\textbf{Date}: 2026.03.28 (Sat)\\
\textbf{Index}: doyattano\\
\textbf{Expression (JA)}: どうやったの？\\
\textbf{Romanization}: Dou yatta no?\\
\textbf{Expression (EN)}: How did you do it?

\textbf{Variants (JA)}: どうやってやったの？ / どうやってそれをしたの？\\
\textbf{Variants (EN)}: How did you do it? / How did you manage to do that?

\textbf{Intent (JA)}: 方法の質問, 驚きの表現, 口語的な表現\\
\textbf{Intent (EN)}: question about method, expression of surprise, colloquial expression

\textbf{Tags (JA)}: 即時文法, 口語, 方法質問表現, 驚き表現, 口語的表現\\
\textbf{Tags (EN)}: immediate grammar, spoken, expression of method question, expression of surprise, colloquial expression

\textbf{Adjusted Expression (JA)}:\\
どのようにしてやったのですか？\\
\textbf{Adjusted Expression (EN)}:\\
How did you do it?

\textbf{Example (JA)}:\\
A「この(プロジェクト:ぷろじぇくと)、(成功:せいこう)させたのはすごいね！どうやったの？」 B「ありがとう！まあ、(大:たい)したことしてないんだけどね」\\
\textbf{Example (EN)}:\\
A: It's amazing that you made this project a success! How did you do it? B: Thanks! Well, I didn't do anything special.

\textbf{Notes (JA)}:\\
「どうやったの？」は、相手が何かを成し遂げたことに対して、その方法や手段を尋ねる口語的な表現。「どうやって」は「how」の意味で、「やったの」は「did it」の意味です。

\textbf{Notes (EN)}:\\
'Dou yatta no?' is a colloquial expression used to ask about the method or means by which someone accomplished something. 'Dou yatte' means 'how', and 'yatta no' means 'did it'.

\newpage

\section*{716. I'll think about it.}

\textbf{Date}: 2026.03.27 (Fri)\\
\textbf{Index}: kangaetoku\\
\textbf{Expression (JA)}: (考:かんが)えとく。\\
\textbf{Romanization}: Kangae toku.\\
\textbf{Expression (EN)}: I'll think about it.

\textbf{Variants (JA)}: 考えておきます。 / 考えておくね。\\
\textbf{Variants (EN)}: I'll think about it. / I'll consider it. / I'll keep it in mind.

\textbf{Intent (JA)}: 検討の表現, 保留の共有, 口語的な表現\\
\textbf{Intent (EN)}: expression of consideration, sharing of reservation, colloquial expression

\textbf{Tags (JA)}: 即時文法, 口語, 検討表現, 保留共有表現, 口語的表現\\
\textbf{Tags (EN)}: immediate grammar, spoken, expression of consideration, expression of reservation sharing, colloquial expression

\textbf{Adjusted Expression (JA)}:\\
(考:かんが)えておきます。\\
\textbf{Adjusted Expression (EN)}:\\
I'll think about it.

\textbf{Example (JA)}:\\
A「(来週:らいしゅう)、(木曜:もくよう)、(来:き)てくれる？」 B「うん、(考:かんが)えとくね」\\
\textbf{Example (EN)}:\\
A: Can you come next Thursday? B: Yeah, I'll think about it.

\textbf{Notes (JA)}:\\
「考えとく」は、検討することを伝える口語的な表現。「考える」は「to think about」の意味で、「とく」は「ておく」の口語的な短縮形で、あらかじめ何かをすることを意味します。しかしながら、「考えとく」は必ずしも実際に考えることを意味するわけではなく、単にその場での返答を保留するための表現として使われることもあります。だから、可能性としては低いかもしれませんし、実際には考えないこともあります。

\textbf{Notes (EN)}:\\
'Kangae toku.' is a colloquial expression used to convey that the speaker will consider something. 'Kangaeru' means 'to think about', and 'toku' is a colloquial contraction of 'te oku', which means to do something in advance. However, 'kangae toku' does not necessarily mean that the speaker will actually think about it; it can also be used as an expression to simply defer a response in the moment. Therefore, it may be less likely that the speaker will actually consider it, and they might not think about it at all.

\newpage

\section*{715. Don't worry about it.}

\textbf{Date}: 2026.03.26 (Thu)\\
\textbf{Index}: kinishinaide\\
\textbf{Expression (JA)}: (気:き)にしないで。\\
\textbf{Romanization}: Ki ni shinaide.\\
\textbf{Expression (EN)}: Don't worry about it.

\textbf{Variants (JA)}: 気にしないでください。 / 気にしないでね。\\
\textbf{Variants (EN)}: Don't worry about it. / Please don't worry about it. / Don't worry about it, okay?

\textbf{Intent (JA)}: 安心の表現, 励ましの共有, 口語的な表現\\
\textbf{Intent (EN)}: expression of reassurance, sharing of encouragement, colloquial expression

\textbf{Tags (JA)}: 即時文法, 口語, 安心表現, 励まし共有表現, 口語的表現\\
\textbf{Tags (EN)}: immediate grammar, spoken, expression of reassurance, expression of encouragement sharing, colloquial expression

\textbf{Adjusted Expression (JA)}:\\
どうぞお(気:き)になさらないでください。\\
\textbf{Adjusted Expression (EN)}:\\
Please don't worry about it.

\textbf{Example (JA)}:\\
A「ごめん、(遅:おそ)くなって」 B「(大丈夫:だいじょうぶ)、(気:き)にしないで」\\
\textbf{Example (EN)}:\\
A: Sorry for being late. B: It's okay, don't worry about it.

\textbf{Notes (JA)}:\\
「気にしないで」は、相手に対して心配や不安を感じさせないようにするための口語的な表現。「気にする」は「to worry about」の意味で、「ないで」は「don't」の意味です。

\textbf{Notes (EN)}:\\
'Ki ni shinaide.' is a colloquial expression used to reassure someone and prevent them from feeling worried or anxious. 'Ki ni suru' means 'to worry about', and 'naide' means 'don't'.

\newpage

\section*{714. Oh, I didn't know that.}

\textbf{Date}: 2026.03.25 (Wed)\\
\textbf{Index}: aashiranakatta\\
\textbf{Expression (JA)}: ああ、知らなかった。\\
\textbf{Romanization}: Aa, shiranakatta.\\
\textbf{Expression (EN)}: Oh, I didn't know that.

\textbf{Variants (JA)}: ああ、そうだったんだ。 / ああ、そういうこと？\\
\textbf{Variants (EN)}: Oh, I see. / Oh, is that so?

\textbf{Intent (JA)}: 驚きの表現, 情報共有の表現, 口語的な表現\\
\textbf{Intent (EN)}: expression of surprise, sharing of information, colloquial expression

\textbf{Tags (JA)}: 即時文法, 口語, 驚き表現, 情報共有表現, 口語的表現\\
\textbf{Tags (EN)}: immediate grammar, spoken, expression of surprise, expression of information sharing, colloquial expression

\textbf{Adjusted Expression (JA)}:\\
ああ、それは(知:し)りませんでした。\\
\textbf{Adjusted Expression (EN)}:\\
Oh, I wasn't aware of that.

\textbf{Example (JA)}:\\
ああ、(知:し)らなかった。それは(初耳:はつみみ)。\\
\textbf{Example (EN)}:\\
Oh, I didn't know that. This is news to me.

\textbf{Notes (JA)}:\\
「ああ、知らなかった」は、何か新しい情報を知ったときや驚いたときに使う口語的な表現。「初耳」は「first time hearing it」の意味で、初めて聞いた情報に対する驚きを表します。

\textbf{Notes (EN)}:\\
'Aa, shiranakatta.' is a colloquial expression used when the speaker learns new information or is surprised by something. 'Hatsumimi' means 'first time hearing it', expressing surprise at hearing information for the first time.

\newpage

\section*{713. Honestly, I'm at a loss.}

\textbf{Date}: 2026.03.24 (Tue)\\
\textbf{Index}: hontomaittayo\\
\textbf{Expression (JA)}: ほんと、まいったよ。\\
\textbf{Romanization}: Honto, maitta yo.\\
\textbf{Expression (EN)}: Honestly, I'm at a loss.

\textbf{Variants (JA)}: (本当:ほんとう)、まいったよ。 / (正直:しょうじき)、まいったよ。\\
\textbf{Variants (EN)}: Honestly, I'm at a loss. / To be honest, I'm at a loss.

\textbf{Intent (JA)}: 感情の表現, 困惑の共有, 口語的な表現\\
\textbf{Intent (EN)}: expression of emotion, sharing of confusion, colloquial expression

\textbf{Tags (JA)}: 即時文法, 口語, 感情表現, 困惑共有表現, 口語的表現\\
\textbf{Tags (EN)}: immediate grammar, spoken, expression of emotion, expression of confusion sharing, colloquial expression

\textbf{Adjusted Expression (JA)}:\\
(本当:ほんとう)に、(困:こま)ってしまいました。\\
\textbf{Adjusted Expression (EN)}:\\
To be honest, I'm at a loss.

\textbf{Example (JA)}:\\
ほんと、まいったよ。この(問題:もんだい)、どうやって(解決:かいけつ)すればいいのか、わからないんだ。\\
\textbf{Example (EN)}:\\
Honestly, I'm at a loss. I don't know how to solve this problem.

\textbf{Notes (JA)}:\\
「ほんと、まいったよ。」は、本当に困惑していることを伝える口語的な表現。「ほんと」は「本当」の略で、「まいった」は「困った」の意味です。「よ」は強調のための終助詞です。

\textbf{Notes (EN)}:\\
'Honto, maitta yo.' is a colloquial expression used to convey that the speaker is genuinely confused or at a loss about something. 'Honto' means 'honestly' or 'really', and 'maitta' means 'at a loss' or 'defeated'. The particle 'yo' is used for emphasis in spoken language.

\newpage

\section*{712. To Get Straight to the Point,}

\textbf{Date}: 2026.03.23 (Mon)\\
\textbf{Index}: ketsuronkaraiuto\\
\textbf{Expression (JA)}: (結論:けつろん)から(言:い)うと、...\\
\textbf{Romanization}: Ketsuron kara iu to, ...\\
\textbf{Expression (EN)}: To get straight to the point, ...

\textbf{Variants (JA)}: (要点:ようてん)から(言:い)うと、... / (結論:けつろん)を(先:さき)に(言:い)うと、...\\
\textbf{Variants (EN)}: To get straight to the point, ... / To give you the bottom line, ... / To state the conclusion first, ...

\textbf{Intent (JA)}: 結論の提示, 要点の共有, 説明の表現\\
\textbf{Intent (EN)}: stating a conclusion, sharing the key point, framing an explanation

\textbf{Tags (JA)}: 即時文法, 口語, 結論提示表現, 要点共有表現, 説明表現\\
\textbf{Tags (EN)}: immediate grammar, spoken language, conclusion-marking expression, key-point-sharing expression, explanatory expression

\textbf{Adjusted Expression (JA)}:\\
(結論:けつろん)から(先:さき)に(申:もう)しますと、...\\
\textbf{Adjusted Expression (EN)}:\\
To state the conclusion first, ...

\textbf{Example (JA)}:\\
(結論:けつろん)から(言:い)うと、この(方法:ほうほう)が(最:もっと)も(効果的:こうかてき)だということがわかりました。\\
\textbf{Example (EN)}:\\
To get straight to the point, we found that this method is the most effective.

\textbf{Notes (JA)}:\\
「結論から言うと、...」は、先に要点や結論を示してから説明に入るときに使う口語的な表現です。話を回りくどくせず、聞き手にまず重要点を伝えたいときによく使われます。

\textbf{Notes (EN)}:\\
'Ketsuron kara iu to, ...' is a colloquial expression used when the speaker wants to state the main point or conclusion before giving the details. It is useful when the speaker wants to avoid a roundabout explanation and present the key point first.

\newpage

\section*{711. For now,}

\textbf{Date}: 2026.03.22 (Sun)\\
\textbf{Index}: imanotokoro\\
\textbf{Expression (JA)}: (今:いま)のところ、...\\
\textbf{Romanization}: ima no tokoro, ...\\
\textbf{Expression (EN)}: For now, ...

\textbf{Variants (JA)}: (現時点:げんじてん)では、... / とりあえず、...\\
\textbf{Variants (EN)}: For now, ... / At this point, ... / For the time being, ...

\textbf{Intent (JA)}: 状況の説明, 暫定的な評価の共有, 対比の表現\\
\textbf{Intent (EN)}: explanation of situation, sharing of provisional evaluation, expression of contrast

\textbf{Tags (JA)}: 即時文法, 口語, 状況説明表現, 暫定的評価共有表現, 対比表現\\
\textbf{Tags (EN)}: immediate grammar, spoken, expression of situation explanation, expression of provisional evaluation sharing, expression of contrast

\textbf{Adjusted Expression (JA)}:\\
(現時点:げんじてん)では、...\\
\textbf{Adjusted Expression (EN)}:\\
At this point, ...

\textbf{Example (JA)}:\\
(今:いま)のところ、この(方法:ほうほう)が(最適:さいてき)だと(思:おも)います。\\
\textbf{Example (EN)}:\\
For now, I think this method is the best.

\textbf{Notes (JA)}:\\
「今のところ...」は、現在の状況や評価を説明する際に使う口語的な表現。「今」は「now」の意味で、「ところ」は「point」や「situation」の意味です。

\textbf{Notes (EN)}:\\
'Ima no tokoro...' is a colloquial expression used when explaining the current situation or evaluation. 'Ima' means 'now', and 'tokoro' means 'point' or 'situation'.

\newpage

\section*{710. We can leave now.}

\textbf{Date}: 2026.03.21 (Sat)\\
\textbf{Index}: moushuppatsudekiruyo\\
\textbf{Expression (JA)}: もう(出発:しゅっぱつ)できるよ。\\
\textbf{Romanization}: mou shuppatsu dekiru yo.\\
\textbf{Expression (EN)}: We can leave now.

\textbf{Variants (JA)}: もう行けるよ。 / もう出発できますよ。 / もう出発できると思います。\\
\textbf{Variants (EN)}: We can go now. / We can leave now. / I think we can leave now.

\textbf{Intent (JA)}: 状況の説明, 行動の提案, 確認の表現\\
\textbf{Intent (EN)}: explanation of situation, suggestion of action, expression of confirmation

\textbf{Tags (JA)}: 即時文法, 口語, 状況説明表現, 行動提案表現, 確認表現\\
\textbf{Tags (EN)}: immediate grammar, spoken, expression of situation explanation, expression of action suggestion, expression of confirmation

\textbf{Adjusted Expression (JA)}:\\
もう(出発:しゅっぱつ)できますよ。\\
\textbf{Adjusted Expression (EN)}:\\
We can leave now.

\textbf{Example (JA)}:\\
もう(出発:しゅっぱつ)できるよ。(準備:じゅんび)、できてる？\\
\textbf{Example (EN)}:\\
We can leave now. Are you ready?

\textbf{Notes (JA)}:\\
「もう出発できるよ」は、今すぐに出発できることを伝える口語的な表現。「もう」は「already」や「now」の意味で、「出発」は「departure」の意味です。「できる」は「can do」の意味で、「よ」は強調のための終助詞です。

\textbf{Notes (EN)}:\\
'Mou shuppatsu dekiru yo.' is a colloquial expression used to convey that you can leave right now. 'Mou' means 'already' or 'now', and 'shuppatsu' means 'departure'. 'Dekiru' means 'can do', and 'yo' is a sentence-ending particle used for emphasis.

\newpage

\section*{709. On the way home, I happened to...}

\textbf{Date}: 2026.03.20 (Fri)\\
\textbf{Index}: kaerutochudeneguuzen\\
\textbf{Expression (JA)}: (帰:かえ)る(途中:とちゅう)でね、(偶然:ぐうぜん)、...\\
\textbf{Romanization}: kaeru tochuu de ne, guuzen, ...\\
\textbf{Expression (EN)}: On the way home, I happened to...

\textbf{Variants (JA)}: (帰:かえ)る(途中:とちゅう)で、(偶然:ぐうぜん)、... / (帰:かえ)ろうとしてたらね、(偶然:ぐうぜん)、...\\
\textbf{Variants (EN)}: On the way home, I happened to... / While I was on my way home, I happened to...

\textbf{Intent (JA)}: 経験の共有, 偶然の表現, 状況の説明\\
\textbf{Intent (EN)}: sharing of experience, expression of coincidence, explanation of situation

\textbf{Tags (JA)}: 即時文法, 口語, 経験共有表現, 偶然表現, 状況説明表現\\
\textbf{Tags (EN)}: immediate grammar, spoken, expression of experience sharing, expression of coincidence, expression of situation explanation

\textbf{Adjusted Expression (JA)}:\\
(帰:かえ)る(途中:とちゅう)で(偶然:ぐうぜん)、\\
\textbf{Adjusted Expression (EN)}:\\
On the way home, I happened to...

\textbf{Example (JA)}:\\
(帰:かえ)る(途中:とちゅう)でね、(偶然:ぐうぜん)、この(店:みせ)を見つけたんだ。\\
\textbf{Example (EN)}:\\
On the way home, I happened to find this store.

\textbf{Notes (JA)}:\\
「帰る途中でね、偶然、...」は、帰宅する途中で何か予期しないことが起こったことを伝える表現。「帰る途中で」は「on the way home」の意味で、「偶然」は「happened to」の意味です。

\textbf{Notes (EN)}:\\
'Kaeru tochū de ne, gūzen, ...' is an expression used to convey that something unexpected happened while on the way home. 'Kaeru tochū de' means 'on the way home', and 'gūzen' means 'happened to'.

\newpage

\section*{708. I clean on weekends.}

\textbf{Date}: 2026.03.19 (Thu)\\
\textbf{Index}: shuumatsusoujishiteru\\
\textbf{Expression (JA)}: (週末:しゅうまつ)、(掃除:そうじ)してる。\\
\textbf{Romanization}: shūmatsu, sōji shiteru.\\
\textbf{Expression (EN)}: I clean on weekends.

\textbf{Variants (JA)}: (週末:しゅうまつ)は(掃除:そうじ)してる。 / (週末:しゅうまつ)に(掃除:そうじ)してる。\\
\textbf{Variants (EN)}: I clean on weekends. / On weekends, I clean.

\textbf{Intent (JA)}: 習慣の表現, 時間の共有, 活動の説明\\
\textbf{Intent (EN)}: expression of habit, sharing of time, explanation of activity

\textbf{Tags (JA)}: 即時文法, 口語, 習慣表現, 時間共有表現, 活動説明表現\\
\textbf{Tags (EN)}: immediate grammar, spoken, expression of habit, expression of time sharing, expression of activity explanation

\textbf{Adjusted Expression (JA)}:\\
(週末:しゅうまつ)は(掃除:そうじ)しています。\\
\textbf{Adjusted Expression (EN)}:\\
(On weekends), I clean.

\textbf{Example (JA)}:\\
(週末:しゅうまつ)、(掃除:そうじ)してる。\\
\textbf{Example (EN)}:\\
I clean on weekends. 

\textbf{Notes (JA)}:\\
「週末、掃除してる」は、話し手が週末に掃除をする習慣があることを伝える表現。「してる」は「している」の口語的な短縮形で、現在の習慣や状態を表します。

\textbf{Notes (EN)}:\\
'Shūmatsu, sōji shiteru.' is a colloquial expression used to convey that the speaker has a habit of cleaning on weekends. 'Shiteru' is a colloquial contraction of 'shite iru', which indicates a current habit or state.

\newpage

\section*{707. It was the opposite for me, but...}

\textbf{Date}: 2026.03.18 (Wed)\\
\textbf{Index}: bokuwagyakudeshitakedo\\
\textbf{Expression (JA)}: (僕:ぼく)は(逆:ぎゃく)でしたけど、...\\
\textbf{Romanization}: boku wa gyaku deshita kedo, ...\\
\textbf{Expression (EN)}: It was the opposite for me, but...

\textbf{Variants (JA)}: (私:わたし)は(逆:ぎゃく)でしたけど、... / (自分:じぶん)は(逆:ぎゃく)でしたけど、...\\
\textbf{Variants (EN)}: It was the opposite for me, but... / I had the opposite experience, but...

\textbf{Intent (JA)}: 経験の共有, 対比の表現, 状況の説明\\
\textbf{Intent (EN)}: sharing of experience, expression of contrast, explanation of situation

\textbf{Tags (JA)}: 即時文法, 口語, 経験共有表現, 対比表現, 状況説明表現\\
\textbf{Tags (EN)}: immediate grammar, spoken, expression of experience sharing, expression of contrast, expression of situation explanation

\textbf{Adjusted Expression (JA)}:\\
(私:わたし)は(逆:ぎゃく)でしたが、...\\
\textbf{Adjusted Expression (EN)}:\\
(I had the opposite experience), but...

\textbf{Example (JA)}:\\
(僕:ぼく)は(逆:ぎゃく)でしたけど、この(方法:ほうほう)、(効果的:こうかてき)だと思います。\\
\textbf{Example (EN)}:\\
It was the opposite for me, but I think this method is effective.

\textbf{Notes (JA)}:\\
「僕は逆でしたけど、...」は、話し手が自分の経験や意見が一般的なものとは異なることを伝える際に使う口語的な表現。「僕」は「I」の意味で、「逆」は「opposite」の意味です。

\textbf{Notes (EN)}:\\
'Boku wa gyaku deshita kedo, ...' is a colloquial expression used when the speaker wants to convey that their experience or opinion is different from the general one. 'Boku' means 'I', and 'gyaku' means 'opposite'.

\newpage

\section*{706. In addition,}

\textbf{Date}: 2026.03.17 (Tue)\\
\textbf{Index}: purasushite\\
\textbf{Expression (JA)}: (プラス:ぷらす)して、...\\
\textbf{Romanization}: purasu shite, ...\\
\textbf{Expression (EN)}: In addition, ...

\textbf{Variants (JA)}: (さらに:さらに)加えて、... / (それに:それに)加えて、...\\
\textbf{Variants (EN)}: In addition, ... / Furthermore, ... / Moreover, ...

\textbf{Intent (JA)}: 追加の情報の表現, 強調の共有, 説明の表現\\
\textbf{Intent (EN)}: expression of additional information, sharing of emphasis, expression of explanation

\textbf{Tags (JA)}: 即時文法, 口語, 追加情報表現, 強調共有表現, 説明表現\\
\textbf{Tags (EN)}: immediate grammar, spoken, expression of additional information, expression of emphasis sharing, expression of explanation

\textbf{Adjusted Expression (JA)}:\\
さらに(加:くわ)えると、...\\
\textbf{Adjusted Expression (EN)}:\\
Furthermore, ...

\textbf{Example (JA)}:\\
(プラス:ぷらす)して、この(機能:きのう)は(ユーザー:ゆーざー)からも(好評:こうひょう)です。\\
\textbf{Example (EN)}:\\
In addition, this feature is also popular among users.

\textbf{Notes (JA)}:\\
「プラスして、...」は、ある情報や意見にさらに追加の内容を付け加える際に使う表現。「プラス」は「plus」の意味で、「して」は「and」の意味です。

\textbf{Notes (EN)}:\\
'Purasu shite, ...' is a colloquial expression used when adding additional information or opinions to something already mentioned. 'Purasu' means 'plus', and 'shite' means 'and'.

\newpage

\section*{705. By default,}

\textbf{Date}: 2026.03.16 (Mon)\\
\textbf{Index}: deforutode\\
\textbf{Expression (JA)}: (デフォルト:でふぉると)で、...\\
\textbf{Romanization}: deforuto de, ...\\
\textbf{Expression (EN)}: By default, ...

\textbf{Variants (JA)}: (初期設定:しょきせってい)では、... / (通常:つうじょう)は、...\\
\textbf{Variants (EN)}: By default, ... / In the default setting, ... / Usually, ...

\textbf{Intent (JA)}: 説明の表現, 状況の共有, 一般的な評価\\
\textbf{Intent (EN)}: expression of explanation, sharing of situation, general evaluation

\textbf{Tags (JA)}: 即時文法, 口語, 説明表現, 状況共有表現, 一般的評価表現\\
\textbf{Tags (EN)}: immediate grammar, spoken, expression of explanation, expression of situation sharing, expression of general evaluation

\textbf{Adjusted Expression (JA)}:\\
(初期値:しょきち)として、...が(設定:せってい)されています。\\
\textbf{Adjusted Expression (EN)}:\\
As a default value, ... is set.

\textbf{Example (JA)}:\\
(デフォルト:でふぉると)で、この(設定:せってい)は(有効:ゆうこう)になっています。\\
\textbf{Example (EN)}:\\
By default, this setting is enabled.

\textbf{Notes (JA)}:\\
「デフォルトで、...」は、ある設定や状態が初期状態であることを説明する表現。

\textbf{Notes (EN)}:\\
'Deforuto de, ...' is a colloquial expression used to explain that a certain setting or state is the default or initial state.

\newpage

\section*{704. This is the main thing right now, but...}

\textbf{Date}: 2026.03.15 (Sun)\\
\textbf{Index}: imawakoregameinnandesukedo\\
\textbf{Expression (JA)}: (今:いま)は、これが(メイン:めいん)なんですけど、...\\
\textbf{Romanization}: ima wa, kore ga mein nan desu kedo, ...\\
\textbf{Expression (EN)}: This is the main thing right now, but...

\textbf{Variants (JA)}: 今は、これがメインなんですが、... / 今は、これが主なものなんですけど、...\\
\textbf{Variants (EN)}: This is the main thing right now, but... / This is the primary thing right now, but...

\textbf{Intent (JA)}: 状況の説明, 主題の提示, 対比の表現\\
\textbf{Intent (EN)}: explanation of situation, presentation of topic, expression of contrast

\textbf{Tags (JA)}: 即時文法, 口語, 状況説明表現, 主題提示表現, 対比表現\\
\textbf{Tags (EN)}: immediate grammar, spoken, expression of situation explanation, expression of topic presentation, expression of contrast

\textbf{Adjusted Expression (JA)}:\\
(現在:げんざい)は、これが(主:おも)なものなんですが、...\\
\textbf{Adjusted Expression (EN)}:\\
(Currently), this is the main thing, but...

\textbf{Example (JA)}:\\
(今:いま)は、これが(メイン:めいん)なんですけど、(将来:しょうらい)はもっと(大:おお)きな(プロジェクト:ぷろじぇくと)をやりたいと思っています。\\
\textbf{Example (EN)}:\\
This is the main thing right now, but in the future, I want to work on a bigger project.

\textbf{Notes (JA)}:\\
「今は、これがメインなんですけど、...」は、現在の状況や主題を説明する表現。

\textbf{Notes (EN)}:\\
'Ima wa, kore ga mein nan desu kedo, ...' is an expression used when explaining the current situation or topic.

\newpage

\section*{703. When you think about it in total,}

\textbf{Date}: 2026.03.14 (Sat)\\
\textbf{Index}: totarudekangaeruto\\
\textbf{Expression (JA)}: (トータル:とーたる)で(考:かんが)えると、...\\
\textbf{Romanization}: tōtaru de kangaeru to, ...\\
\textbf{Expression (EN)}: When you think about it in total, ...

\textbf{Variants (JA)}: (全体:ぜんたい)で(考:かんが)えると、... / (総合的:そうごうてき)に(考:かんが)えると、...\\
\textbf{Variants (EN)}: When you think about it as a whole, ... / When you think about it comprehensively, ...

\textbf{Intent (JA)}: 総合的な評価の表現, 全体像の共有, 結論の提示\\
\textbf{Intent (EN)}: expression of comprehensive evaluation, sharing of overall picture, presentation of conclusion

\textbf{Tags (JA)}: 即時文法, 口語, 総合的評価表現, 全体像共有表現, 結論提示表現\\
\textbf{Tags (EN)}: immediate grammar, spoken, expression of comprehensive evaluation, expression of overall picture sharing, expression of conclusion presentation

\textbf{Adjusted Expression (JA)}:\\
(全体的:ぜんたいてき)に(考:かんが)えると、...\\
\textbf{Adjusted Expression (EN)}:\\
When you think about it as a whole, ...

\textbf{Example (JA)}:\\
(トータル:とーたる)で(考:かんが)えると、この(プロジェクト:ぷろじぇくと)、(成功:せいこう)だと思います。\\
\textbf{Example (EN)}:\\
When you think about it in total, I think this project is a success.

\textbf{Notes (JA)}:\\
「トータルで考えると、...」は、話し手がある事柄や状況を全体的に評価する際に使う口語的な表現。「トータル」は「total」の意味で、「考えると」は「when you think about it」の意味です。

\textbf{Notes (EN)}:\\
'Tōtaru de kangaeru to, ...' is a colloquial expression used when the speaker wants to evaluate a certain matter or situation as a whole. 'Tōtaru' means 'total', and 'kangaeru to' means 'when you think about it'.

\newpage

\section*{702. Is that enough?}

\textbf{Date}: 2026.03.13 (Fri)\\
\textbf{Index}: soredetariru\\
\textbf{Expression (JA)}: それで(足:た)りる？\\
\textbf{Romanization}: sore de tariru?\\
\textbf{Expression (EN)}: Is that enough?

\textbf{Variants (JA)}: それで足りますか？ / それで十分ですか？ / それで大丈夫ですか？\\
\textbf{Variants (EN)}: Is that enough? / Is that sufficient? / Is that okay?

\textbf{Intent (JA)}: 確認の表現, 十分性の共有, 質問の表現\\
\textbf{Intent (EN)}: expression of confirmation, sharing of sufficiency, expression of question

\textbf{Tags (JA)}: 即時文法, 口語, 確認表現, 十分性共有表現, 質問表現\\
\textbf{Tags (EN)}: immediate grammar, spoken, expression of confirmation, expression of sufficiency sharing, expression of question

\textbf{Adjusted Expression (JA)}:\\
それで(足:た)りますか？\\
\textbf{Adjusted Expression (EN)}:\\
Is that enough?

\textbf{Example (JA)}:\\
A「(割:わ)り(勘:かん)でいい？」 B「うん、じゃ、これ、(私:わたし)の(分:ぶん)、それで(足:た)りる？」 A「うん、ありがとう」\\
\textbf{Example (EN)}:\\
A: Is it okay to split the bill? B: Yes, then, for my part, is that enough? A: Yes, thank you.

\textbf{Notes (JA)}:\\
「それで足りる？」は、話し手があるものや情報が十分であるかどうかを確認する際に使う口語的な表現。「それで」は「with that」の意味で、「足りる？」は「is it enough?」の意味です。「わりかん」は「split the bill」の意味で、食事の代金を人数で割ることを指します。

\textbf{Notes (EN)}:\\
'Sore de tariru?' is a colloquial expression used when the speaker wants to confirm whether something or information is sufficient. 'Sore de' means 'with that', and 'tariru?' means 'is it enough?'. 'Warikan' means 'split the bill', referring to dividing the cost of a meal among the number of people.

\newpage

\section*{701. Is that what you mean?}

\textbf{Date}: 2026.03.12 (Thu)\\
\textbf{Index}: souiukoto\\
\textbf{Expression (JA)}: そういうこと？\\
\textbf{Romanization}: sou iu koto?\\
\textbf{Expression (EN)}: Is that what you mean?

\textbf{Variants (JA)}: そういうことですか？ / そういう意味ですか？ / そういうことなんですか？\\
\textbf{Variants (EN)}: Is that what you mean? / Do you mean that? / Is that what you're saying?

\textbf{Intent (JA)}: 確認の表現, 理解の共有, 質問の表現\\
\textbf{Intent (EN)}: expression of confirmation, sharing of understanding, expression of question

\textbf{Tags (JA)}: 即時文法, 口語, 確認表現, 理解共有表現, 質問表現\\
\textbf{Tags (EN)}: immediate grammar, spoken, expression of confirmation, expression of understanding sharing, expression of question

\textbf{Adjusted Expression (JA)}:\\
それはそういうことですか？\\
\textbf{Adjusted Expression (EN)}:\\
Is that what you mean?

\textbf{Example (JA)}:\\
A「この(問題:もんだい)、(解決:かいけつ)するには、(新:あたら)しい(方法:ほうほう)が必要だと(思:おも)う」 B「そういうこと？」 A「うん、(今:いま)のやり(方:かた)じゃ、(時間:じかん)がかかりすぎるからね」\\
\textbf{Example (EN)}:\\
A: I think we need a new method to solve this problem. B: Is that what you mean? A: Yes, because the current way takes too much time.

\textbf{Notes (JA)}:\\
「そういうこと？」は、相手の発言や意図を確認する際に使う口語的な表現。「そういう」は「that kind of」の意味で、「こと」は「thing」の意味です。

\textbf{Notes (EN)}:\\
'Sou iu koto?' is a colloquial expression used when confirming someone's statement or intention. 'Sou iu' means 'that kind of', and 'koto' means 'thing'.

\newpage

\section*{700. Do you want more?}

\textbf{Date}: 2026.03.11 (Wed)\\
\textbf{Index}: mottohoshii\\
\textbf{Expression (JA)}: もっとほしい？\\
\textbf{Romanization}: motto hoshii?\\
\textbf{Expression (EN)}: Do you want more?

\textbf{Variants (JA)}: もっとほしいですか？ / もっと欲しいと思いますか？ / もっと欲しいですか？\\
\textbf{Variants (EN)}: Do you want more? / Do you think you want more? / Do you want more, exactly?

\textbf{Intent (JA)}: 質問の表現, 欲求の共有, 情報収集の表現\\
\textbf{Intent (EN)}: expression of question, sharing of desire, expression of information gathering

\textbf{Tags (JA)}: 即時文法, 口語, 質問表現, 欲求共有表現, 情報収集表現\\
\textbf{Tags (EN)}: immediate grammar, spoken, expression of question, expression of desire sharing, expression of information gathering

\textbf{Adjusted Expression (JA)}:\\
もう(少:すこ)しいかがですか？\\
\textbf{Adjusted Expression (EN)}:\\
Would you like a little more?

\textbf{Example (JA)}:\\
A「この(ケーキ:けーき)、(美味:おい)しいね」 B「うん、もっとほしい？まだ、あるよ」 A「うん、もうちょっと(食:た)べたいな」\\
\textbf{Example (EN)}:\\
A: This cake is delicious, isn't it? B: Yes, do you want more? There's still some left. A: Yes, I want to eat a little more.

\textbf{Notes (JA)}:\\
「もっとほしい？」は、さらに欲しいかどうかを尋ねる際に使う表現。「もっと」は「more」の意味で、「ほしい？」は「do you want...?」の意味です。

\textbf{Notes (EN)}:\\
'Motto hoshii?' is a colloquial expression used when asking if someone wants more of something. 'Motto' means 'more', and 'hoshii?' means 'do you want...?'.

\newpage

\section*{699. Let me do it?}

\textbf{Date}: 2026.03.10 (Tue)\\
\textbf{Index}: yarasete\\
\textbf{Expression (JA)}: やらせて？\\
\textbf{Romanization}: yarasete?\\
\textbf{Expression (EN)}: Let me do it?

\textbf{Variants (JA)}: やらせてください。 / やらせてもらえますか？ / やらせてくれませんか？\\
\textbf{Variants (EN)}: Let me do it, please. / Could you let me do it? / Would you mind letting me do it?

\textbf{Intent (JA)}: 依頼の表現, 丁寧な要求の共有, 協力の表現\\
\textbf{Intent (EN)}: expression of request, sharing of polite demand, expression of cooperation

\textbf{Tags (JA)}: 即時文法, 口語, 依頼表現, 丁寧な要求共有表現, 協力表現\\
\textbf{Tags (EN)}: immediate grammar, spoken, expression of request, expression of polite demand sharing, expression of cooperation

\textbf{Adjusted Expression (JA)}:\\
やらせてください。\\
\textbf{Adjusted Expression (EN)}:\\
Please let me do it.

\textbf{Example (JA)}:\\
A「この(仕事:しごと)、(私:わたし)にやらせて？」 B「いいよ、お願い」\\
\textbf{Example (EN)}:\\
A: Can I do this work? B: Sure, please go ahead.

\textbf{Notes (JA)}:\\
「やらせて？」は、何かをすることを許可してほしいときに使う口語的な表現。

\textbf{Notes (EN)}:\\
'Yarasete?' is a colloquial expression used when you want to ask for permission to do something. It is a shortened form of 'やらせてください' (Please let me do it) and is often used in casual conversations.

\newpage

\section*{698. What do you use it for?}

\textbf{Date}: 2026.03.09 (Mon)\\
\textbf{Index}: naninitsukauno\\
\textbf{Expression (JA)}: (何:なに)に(使:つか)うの？\\
\textbf{Romanization}: nani ni tsukau no?\\
\textbf{Expression (EN)}: What do you use it for?

\textbf{Variants (JA)}: (何:なに)に(使:つか)いますか？ / (何:なに)に(使用:しよう)するんですか？\\
\textbf{Variants (EN)}: What do you use it for? / What do you use it for, exactly?

\textbf{Intent (JA)}: 質問の表現, 目的の共有, 情報収集の表現\\
\textbf{Intent (EN)}: expression of question, sharing of purpose, expression of information gathering

\textbf{Tags (JA)}: 即時文法, 口語, 質問表現, 目的共有表現, 情報収集表現\\
\textbf{Tags (EN)}: immediate grammar, spoken, expression of question, expression of purpose sharing, expression of information gathering

\textbf{Adjusted Expression (JA)}:\\
(何:なに)に(使用:しよう)しますか？\\
\textbf{Adjusted Expression (EN)}:\\
What do you use it for, exactly?

\textbf{Example (JA)}:\\
A「これ、(借:か)りてもいい？」 B「え、(何:なに)に(使:つか)うの？」 A「ちょっと、うどんをね」 B「じゃ、こっちの(小:ちい)さい(方:ほう)がいいよ」\\
\textbf{Example (EN)}:\\
A: Can I borrow this? B: Huh, what do you use it for? A: Just for udon. B: Then, this smaller one is better.

\textbf{Notes (JA)}:\\
「何に使うの？」は、道具の目的や用途を尋ねる表現。「何に」は「what for」の意味で、「使うの？」は「do you use it?」の意味です。

\textbf{Notes (EN)}:\\
'Nani ni tsukau no?' is a colloquial expression used to ask about the purpose or use of a tool or item. 'Nani ni' means 'what for', and 'tsukau no?' means 'do you use it?'.

\newpage

\section*{697. Isn't it natural?}

\textbf{Date}: 2026.03.08 (Sun)\\
\textbf{Index}: atarimaejanai\\
\textbf{Expression (JA)}: (当:あ)たり(前:まえ)じゃないですか。\\
\textbf{Romanization}: atarimae janai desu ka.\\
\textbf{Expression (EN)}: Isn't it natural?

\textbf{Variants (JA)}: (当:あ)たり(前:まえ)じゃない？ / (当然:とうぜん)じゃないですか？ / もちろんですよ。\\
\textbf{Variants (EN)}: Isn't it natural? / Isn't it obvious? / Of course.

\textbf{Intent (JA)}: 同意の表現, 当然の評価の共有, 提案の表現\\
\textbf{Intent (EN)}: expression of agreement, sharing of obvious evaluation, expression of suggestion

\textbf{Tags (JA)}: 即時文法, 口語, 同意表現, 当然の評価共有表現, 提案表現\\
\textbf{Tags (EN)}: immediate grammar, spoken, expression of agreement, expression of obvious evaluation sharing, expression of suggestion

\textbf{Adjusted Expression (JA)}:\\
(当:あ)たり(前:まえ)だと(思:おも)います。\\
\textbf{Adjusted Expression (EN)}:\\
I think it's natural.

\textbf{Example (JA)}:\\
(当:あ)たり(前:まえ)じゃないですか。そんなこと(言:い)われなくてもわかりますよね。\\
\textbf{Example (EN)}:\\
Isn't it natural? You don't need to be told that, right?

\textbf{Notes (JA)}:\\
「当たり前」は「natural」の意味で、あることが当然であることを示します。「じゃないですか」は「isn't it?」の意味で、相手に同意を求める表現です。

\textbf{Notes (EN)}:\\
'Atarimae janai desu ka.' means 'Isn't it natural?' in English. 'Atarimae' means 'natural', indicating that something is obvious or should be expected. 'Janai desu ka' is a way to ask for agreement from the listener, similar to 'isn't it?' in English.

\newpage

\section*{696. Rather, it's better not to look.}

\textbf{Date}: 2026.03.07 (Sat)\\
\textbf{Index}: mushirominaihoga\\
\textbf{Expression (JA)}: むしろ、(見:み)ない(方:ほう)がいい。\\
\textbf{Romanization}: mushiro, minai hou ga ii.\\
\textbf{Expression (EN)}: Rather, it's better not to look.

\textbf{Variants (JA)}: むしろ、見ない方がいいですよ。 / むしろ、見ない方がいいと思います。 / むしろ、見ない方がいいかもしれません。\\
\textbf{Variants (EN)}: Rather, it's better not to look. / Rather, I think it's better not to look. / Rather, it might be better not to look.

\textbf{Intent (JA)}: 提案の表現, 否定的な評価の共有, 状況の説明\\
\textbf{Intent (EN)}: expression of suggestion, sharing of negative evaluation, explanation of situation

\textbf{Tags (JA)}: 即時文法, 口語, 提案表現, 否定的評価共有表現, 状況説明表現\\
\textbf{Tags (EN)}: immediate grammar, spoken, expression of suggestion, expression of negative evaluation sharing, expression of situation explanation

\textbf{Adjusted Expression (JA)}:\\
むしろ、(見:み)ない(方:ほう)がいいですね。\\
\textbf{Adjusted Expression (EN)}:\\
Rather, it's better not to look, isn't it?

\textbf{Example (JA)}:\\
A「この(ヒント:ひんと)、(見:み)た(方:ほう)がいい？」 B「いや、むしろ、(見:み)ない(方:ほう)がいい。わかんなくなるから」\\
\textbf{Example (EN)}:\\
A: Should I look at this hint? B: No, rather, it's better not to look. It will just make it more confusing.

\textbf{Notes (JA)}:\\
「むしろ、見ない方がいい。」は、ある行動や選択を避けることを提案する際に使う口語的な表現。「むしろ」は「rather」の意味で、「見ない方がいい」は「it's better not to look」の意味です。

\textbf{Notes (EN)}:\\
'Mushiro, minai hou ga ii.' is a colloquial expression used when the speaker wants to suggest avoiding a certain action or choice. 'Mushiro' means 'rather', and 'minai hou ga ii' means 'it's better not to look'.

\newpage

\section*{695. Isn't there no choice but to do it?}

\textbf{Date}: 2026.03.06 (Fri)\\
\textbf{Index}: surushikanaijanai\\
\textbf{Expression (JA)}: するしかないじゃないですか。\\
\textbf{Romanization}: suru shika nai ja nai desu ka.\\
\textbf{Expression (EN)}: Isn't there no choice but to do it?

\textbf{Variants (JA)}: やるしかないじゃないですか。 / 行くしかないじゃないですか。 / 食べるしかないじゃないですか。\\
\textbf{Variants (EN)}: Isn't there no choice but to do it? / Isn't there no choice but to go? / Isn't there no choice but to eat?

\textbf{Intent (JA)}: 必然性の表現, 選択肢の共有, 状況の説明\\
\textbf{Intent (EN)}: expression of inevitability, sharing of options, explanation of situation

\textbf{Tags (JA)}: 即時文法, 口語, 必然性表現, 選択肢共有表現, 状況説明表現\\
\textbf{Tags (EN)}: immediate grammar, spoken, expression of inevitability, expression of options sharing, expression of situation explanation

\textbf{Adjusted Expression (JA)}:\\
する(以外:いがい)の(選択:せんたく)はありませんよね。\\
\textbf{Adjusted Expression (EN)}:\\
There is no choice but to do it, isn't there?

\textbf{Example (JA)}:\\
まず、(先:さき)にこの(問題:もんだい)、(解決:かいけつ)するしかないじゃないですか。\\
\textbf{Example (EN)}:\\
First, isn't there no choice but to solve this problem?

\textbf{Notes (JA)}:\\
「するしかないんじゃないですか」は、話し手がある行動や選択が避けられないと感じていることを伝える際に使う口語的な表現。

\textbf{Notes (EN)}:\\
'(Suru shika nai) n janai desu ka.' is a colloquial expression used when the speaker wants to convey that they feel a certain action or choice is inevitable.

\newpage

\section*{694. Some people think that way, you know.}

\textbf{Date}: 2026.03.05 (Thu)\\
\textbf{Index}: souiufoonikangaeruhitomoiruyone\\
\textbf{Expression (JA)}: そういうふうに(考:かんが)える(人:ひと)もいるよね。\\
\textbf{Romanization}: sou iu fuu ni kangaeru hito mo iru yo ne.\\
\textbf{Expression (EN)}: Some people think that way, you know.

\textbf{Variants (JA)}: そういうふうに考える人もいるね。 / そういう考え方の人もいるよ。 / そういう見方をする人もいるよね。\\
\textbf{Variants (EN)}: Some people do think that way. / There are people who see it that way.

\textbf{Intent (JA)}: 断定を避ける応答, 一般化による返答, 伝聞内容への控えめな応答\\
\textbf{Intent (EN)}: non-committal response, answer by generalizing, reserved response to reported opinions

\textbf{Tags (JA)}: 即時文法, 口語, 応答表現, 一般化表現\\
\textbf{Tags (EN)}: immediate grammar, spoken, response expression, generalization

\textbf{Adjusted Expression (JA)}:\\
そういうふうに(考:かんが)える(人:ひと)もいますね。\\
\textbf{Adjusted Expression (EN)}:\\
There are people who think that way.

\textbf{Example (JA)}:\\
A「(都会:とかい)より(地方:ちほう)の(方:ほう)が(暮:く)らしやすいと(思:おも)うって、(本当:ほんとう)？」 B「そういうふうに(考:かんが)える(人:ひと)もいるよね」\\
\textbf{Example (EN)}:\\
A: Is it true that some people say it's easier to live in the countryside than in the city? B: Some people think that way, you know.

\textbf{Notes (JA)}:\\
伝聞や一般的に言われていることを確認されたときに、「そういう見方をする人もいる」と一般化して答える表現。内容の真偽を断定せず、ある考え方の存在だけを認める言い方。

\textbf{Notes (EN)}:\\
Used when responding to a question about something reportedly said or believed by others. The speaker does not confirm the truth of the statement but acknowledges that such a view exists.

\newpage

\section*{693. It might take a little longer.}

\textbf{Date}: 2026.03.04 (Wed)\\
\textbf{Index}: mouchottokakarukana\\
\textbf{Expression (JA)}: もうちょっとかかるかな。\\
\textbf{Romanization}: mou chotto kakaru kana.\\
\textbf{Expression (EN)}: It might take a little longer.

\textbf{Variants (JA)}: もうちょっと時間がかかるかな。 / もうちょっと手間がかかるかな。 / もうちょっと時間が必要かな。\\
\textbf{Variants (EN)}: It might take a little more time. / It might take a little more effort. / It might require a little more time.

\textbf{Intent (JA)}: 予測の表現, 不確実性の共有, 状況の説明\\
\textbf{Intent (EN)}: expression of prediction, sharing of uncertainty, explanation of situation

\textbf{Tags (JA)}: 即時文法, 口語, 予測表現, 不確実性共有表現, 状況説明表現\\
\textbf{Tags (EN)}: immediate grammar, spoken, expression of prediction, expression of uncertainty sharing, expression of situation explanation

\textbf{Adjusted Expression (JA)}:\\
もう(少:すこ)し、(時間:じかん)がかかるかもしれませんね。\\
\textbf{Adjusted Expression (EN)}:\\
It might take a little more time, don't you think?

\textbf{Example (JA)}:\\
A「この(プロジェクト:ぷろじぇくと)、(完了:かんりょう)までどれくらい(時間:じかん)が(必要:ひつよう)ですか？」 B「うーん、もうちょっとかかるかな」\\
\textbf{Example (EN)}:\\
A: How much time do we need to complete this project? B: Hmm, it might take a little longer.

\textbf{Notes (JA)}:\\
「もうちょっとかかるかな」は、話し手がある作業の完了までにもう少し時間が必要だと予測していることを伝える際に使う口語的な表現。「もうちょっと」は「a little more」の意味で、ある程度の追加の時間や努力が必要であることを示します。時間が長いか短かいかは述べていません。

\textbf{Notes (EN)}:\\
'Mou chotto kakaru kana.' is a colloquial expression used when the speaker wants to convey that they predict that a certain task will require a little more time to complete. 'Mou chotto' means 'a little more', indicating that some additional time or effort is needed. It does not specify whether the time will be long or short.

\newpage

\section*{692. It's really fun.}

\textbf{Date}: 2026.03.03 (Tue)\\
\textbf{Index}: jitsunitanoshii\\
\textbf{Expression (JA)}: (実:じつ)にたのしい。\\
\textbf{Romanization}: jitsu ni tanoshii.\\
\textbf{Expression (EN)}: It's really fun.

\textbf{Variants (JA)}: 本当に楽しい。 / とても楽しい。 / すごく楽しい。\\
\textbf{Variants (EN)}: It's really fun. / It's very fun. / It's so much fun.

\textbf{Intent (JA)}: 感情の表現, 楽しさの共有, ポジティブな評価\\
\textbf{Intent (EN)}: expression of emotion, sharing of enjoyment, positive evaluation

\textbf{Tags (JA)}: 即時文法, 口語, 感情表現, 楽しさ共有表現, ポジティブ評価表現\\
\textbf{Tags (EN)}: immediate grammar, spoken, expression of emotion, expression of enjoyment sharing, expression of positive evaluation

\textbf{Adjusted Expression (JA)}:\\
(本当:ほんとう)に(楽:たの)しいですね。\\
\textbf{Adjusted Expression (EN)}:\\
(It's) really fun, isn't it?

\textbf{Example (JA)}:\\
A「この(ゲーム:げーむ)、すごいね」 B「うん、(実:じつ)に(楽:たの)しい」\\
\textbf{Example (EN)}:\\
A: This game is amazing, isn't it? B: Yes, it's really fun, isn't it?

\textbf{Notes (JA)}:\\
「実にたのしい。」は、話し手がある事柄や状況が非常に楽しいと感じていることを認める際に使う口語的な表現。「ね」「よ」などの終助詞は使わない。直感的に自分の感想を言う表現です。

\textbf{Notes (EN)}:\\
'Jitsu ni tanoshii.' is a colloquial expression used when the speaker wants to acknowledge that they feel a certain matter or situation is very enjoyable. It does not use sentence-ending particles like 'ne' or 'yo'. It is an expression that intuitively conveys one's feelings.

\newpage

\section*{691. Could you please...?}

\textbf{Date}: 2026.03.02 (Mon)\\
\textbf{Index}: temoraemasuka\\
\textbf{Expression (JA)}: ...てもらえますか？\\
\textbf{Romanization}: ...te moraemasu ka?\\
\textbf{Expression (EN)}: Could you please...?

\textbf{Variants (JA)}: ...ていただけますか？ / ...てくれますか？ / ...てくれると助かります。\\
\textbf{Variants (EN)}: Could you please...? / Would you mind...? / It would be helpful if you could...

\textbf{Intent (JA)}: 依頼の表現, 丁寧な要求の共有, 協力の表現\\
\textbf{Intent (EN)}: expression of request, sharing of polite demand, expression of cooperation

\textbf{Tags (JA)}: 即時文法, 口語, 依頼表現, 丁寧な要求共有表現, 協力表現\\
\textbf{Tags (EN)}: immediate grammar, spoken, expression of request, expression of polite demand sharing, expression of cooperation

\textbf{Adjusted Expression (JA)}:\\
...ていただけますか？\\
\textbf{Adjusted Expression (EN)}:\\
...te itadakemasu ka?

\textbf{Example (JA)}:\\
この(書類:しょるい)を(確認:かくにん)してもらえますか？\\
\textbf{Example (EN)}:\\
Could you please check this document?

\textbf{Notes (JA)}:\\
「...てもらえますか？」は、話し手が聞き手に対して何かをしてほしいときに使う口語的な表現。「もらえますか？」は「could you please...?」の意味で、可能形が使われていることに注意。

\textbf{Notes (EN)}:\\
'...te moraemasu ka?' is a colloquial expression used when the speaker wants the listener to do something. '...te' is the te-form of a verb, used to express a request or demand. 'Moraemasu ka?' means 'could you please...?', and it's important to note that the potential form is used here.

\newpage

\section*{690. It's a bit awkward to ask this suddenly, but...}

\textbf{Date}: 2026.03.01 (Sun)\\
\textbf{Index}: kyuunikonnakotokikuarenandakedo\\
\textbf{Expression (JA)}: (急:きゅう)にこんなこと(聞:き)くのも、あれ、なんだけど、...\\
\textbf{Romanization}: kyū ni konna koto kiku no mo, are, nan da kedo, ...\\
\textbf{Expression (EN)}: It's a bit awkward to ask this suddenly, but...

\textbf{Variants (JA)}: 急にこんなこと聞くのも、あれ、なんだけど、ちょっと聞いてもいい？ / 急にこんなこと聞くのも、あれ、なんだけど、ちょっと質問してもいい？ / 急にこんなこと聞くのも、あれ、なんだけど、ちょっと相談してもいい？\\
\textbf{Variants (EN)}: It's a bit awkward to ask this suddenly, but can I ask you something? / It's a bit awkward to ask this suddenly, but can I ask you a question? / It's a bit awkward to ask this suddenly, but can I consult with you about something?

\textbf{Intent (JA)}: 謝罪の表現, 質問の共有, 相談の表現\\
\textbf{Intent (EN)}: expression of apology, sharing of question, expression of consultation

\textbf{Tags (JA)}: 即時文法, 口語, 謝罪表現, 質問共有表現, 相談表現\\
\textbf{Tags (EN)}: immediate grammar, spoken, expression of apology, expression of question sharing, expression of consultation

\textbf{Adjusted Expression (JA)}:\\
(突然:とつぜん)こんなことを(聞:き)くのは、(申:もう)し(訳:わけ)ないんですが、...\\
\textbf{Adjusted Expression (EN)}:\\
(Suddenly) asking something like this is a bit apologetic, but...

\textbf{Example (JA)}:\\
(急:)にこんなことを(聞:き)くのは、あれなんですが、それって事実？\\
\textbf{Example (EN)}:\\
It's a bit awkward to ask this suddenly, but is that true?

\textbf{Notes (JA)}:\\
「急にこんなこと聞くのも、あれ、なんだけど、...」は、話し手がある質問や相談をする際に、その質問や相談が突然であることを謝罪しつつ、相手に対して基礎的な質問や想定外の質問をするために使う口語的な表現。「急に」は「suddenly」の意味で、「あれ」は「a bit awkward」の意味で、その質問や相談が少し気まずいかもしれないことを示します。そもそも「あれ」って何のことを指すのかわかりませんよね。でも、こういうんですよ。

\textbf{Notes (EN)}:\\
'Kyuu ni konna koto kiku no mo, are, nan da kedo, ...' is a colloquial expression used when the speaker wants to ask a certain question or consult about something, while apologizing for the suddenness of the question or consultation. 'Kyuu ni' means 'suddenly', and 'are' means 'a bit awkward', indicating that the question or consultation might be a little uncomfortable. It's not clear what 'are' refers to, but that's how it's said.

\newpage

\section*{689. You came at just the right time.}

\textbf{Date}: 2026.02.28 (Sat)\\
\textbf{Index}: choudoiitokoronikita\\
\textbf{Expression (JA)}: ちょうど、いいところに(来:き)た。\\
\textbf{Romanization}: choudo, ii tokoro ni kita.\\
\textbf{Expression (EN)}: You came at just the right time.

\textbf{Variants (JA)}: ちょうど、いいところに来てくれたね。 / ちょうど、いいところに来てくれてありがとう。 / ちょうど、いいところに来てくれて助かったよ。\\
\textbf{Variants (EN)}: You came at just the right time, thank you. / You came at just the right time, I appreciate it. / You came at just the right time, it was a great help.

\textbf{Intent (JA)}: 歓迎の表現, 感謝の共有, 状況の説明\\
\textbf{Intent (EN)}: expression of welcome, sharing of gratitude, explanation of situation

\textbf{Tags (JA)}: 即時文法, 口語, 歓迎表現, 感謝共有表現, 状況説明表現\\
\textbf{Tags (EN)}: immediate grammar, spoken, expression of welcome, expression of gratitude sharing, expression of situation explanation

\textbf{Adjusted Expression (JA)}:\\
ちょうど、いいところに(来:き)ましたね。\\
\textbf{Adjusted Expression (EN)}:\\
Thank you for coming at just the right time.

\textbf{Example (JA)}:\\
A「ちょうど、いいところに来たね」 B「そう？(手伝:てつだ)おうか？」 A「うん、ありがとう」\\
\textbf{Example (EN)}:\\
A: You came at just the right time. B: Really? Do you want me to help? A: Yes, thank you.

\textbf{Notes (JA)}:\\
「ちょうど、いいところに来たね。」は、「ちょうど」は「just」の意味で、ある状況やタイミングが非常に適していることを示します。だいたいは話し手にとってはいいかもしれませんが、言われた人にとってはよくないことのほうが多いです。

\textbf{Notes (EN)}:\\
'Choudo, ii tokoro ni kita ne.' is a colloquial expression where 'choudo' means 'just', indicating that a certain situation or timing is very appropriate. It is often used when the situation may be good for the speaker but may not be good for the person being spoken to.

\newpage

\section*{688. One of the reasons I started ...}

\textbf{Date}: 2026.02.27 (Fri)\\
\textbf{Index}: hajimetariyuunohitotsutoshitewa\\
\textbf{Expression (JA)}: (始:はじ)めた(理由:りゆう)のひとつとしては、\\
\textbf{Romanization}: ...hajimeta riyū no hitotsu to shite wa,\\
\textbf{Expression (EN)}: One of the reasons I started...

\textbf{Variants (JA)}: (始:はじ)めた(理由:りゆう)のひとつとしては、(興味:きょうみ)があったからです。 / (始:はじ)めた(理由:りゆう)のひとつとしては、(挑戦:ちょうせん)したかったからです。 / (始:はじ)めた(理由:りゆう)のひとつとしては、(成長:せいちょう)したかったからです。\\
\textbf{Variants (EN)}: One of the reasons I started is because I was interested. / One of the reasons I started is because I wanted to challenge myself. / One of the reasons I started is because I wanted to grow.

\textbf{Intent (JA)}: 動機の表現, 理由の共有, 自己紹介の一部\\
\textbf{Intent (EN)}: expression of motivation, sharing of reason, part of self-introduction

\textbf{Tags (JA)}: 即時文法, 口語, 動機表現, 理由共有表現, 自己紹介表現\\
\textbf{Tags (EN)}: immediate grammar, spoken, expression of motivation, expression of reason sharing, expression of self-introduction

\textbf{Adjusted Expression (JA)}:\\
(始:はじ)めた(理由:りゆう)のひとつは、(興味:きょうみ)があったからです。\\
\textbf{Adjusted Expression (EN)}:\\
One of the reasons I started is because I was interested.

\textbf{Example (JA)}:\\
この(プロジェクト:ぷろじぇくと)を始めた理由のひとつとしては、(新:あたら)しいことに(挑戦:ちょうせん)したかったからです。\\
\textbf{Example (EN)}:\\
One of the reasons I started this project is because I wanted to challenge myself with something new.

\textbf{Notes (JA)}:\\
「(始:はじ)めた(理由:りゆう)のひとつとしては、」は、話し手がある行動やプロジェクトを始めた理由の一部を説明する際に使う口語的な表現。「始めた理由のひとつ」は「one of the reasons I started」の意味で、複数の理由がある中の一つを指します。「としては」は「as for」の意味で、その理由を説明するための表現です。

\textbf{Notes (EN)}:\\
'(Hajimeta riyuu) no hitotsu to shite wa,' is a colloquial expression used when the speaker wants to explain one of the reasons for starting a certain action or project. 'Hajimeta riyuu no hitotsu' means 'one of the reasons I started', referring to one of multiple reasons. 'Toshite wa' means 'as for', serving as an expression to explain that reason.

\newpage

\section*{687. No one see such a thing...}

\textbf{Date}: 2026.02.26 (Thu)\\
\textbf{Index}: nannteminna\\
\textbf{Expression (JA)}: ...なんてみんな(見:み)ていませんから、...\\
\textbf{Romanization}: ...nannte minna mite imasen kara, ...\\
\textbf{Expression (EN)}: No one see such a thing...

\textbf{Variants (JA)}: ...なんてみんな見ていませんから、気にしないでください。 / ...なんてみんな見ていませんから、心配しないでください。 / ...なんてみんな見ていませんから、気にすることはありません。\\
\textbf{Variants (EN)}: No one see such a thing, so please don't worry about it. / No one see such a thing, so please don't be concerned about it. / No one see such a thing, so there's no need to worry about it.

\textbf{Intent (JA)}: 否定の表現, 安心の共有, 状況の説明\\
\textbf{Intent (EN)}: expression of negation, sharing of reassurance, explanation of situation

\textbf{Tags (JA)}: 即時文法, 口語, 否定表現, 安心共有表現, 状況説明表現\\
\textbf{Tags (EN)}: immediate grammar, spoken, expression of negation, expression of reassurance sharing, expression of situation explanation

\textbf{Adjusted Expression (JA)}:\\
...なんてみんな見ていませんから、気にしないでください。\\
\textbf{Adjusted Expression (EN)}:\\
No one see such a thing, so please don't worry about it.

\textbf{Example (JA)}:\\
A「この(ミス:みす)、(大勢:おおぜい)の人に(見:み)られたんだ」 B「いえいえ、(他人:ひと)の(ミス:みす)なんてみんな(見:み)ていませんから、(気:き)にしないでください」\\
\textbf{Example (EN)}:\\
A: This mistake was seen by a lot of people. B: No,  no one care about someone's mistakes, so never mind.

\textbf{Notes (JA)}:\\
「...なんてみんな見ていませんから、...」は、話し手がある出来事や状況が他の人に見られていないことを伝える際に使う口語的な表現。「なんて」は「such a thing」の意味で、あまり重要でないことを示します。「みんな見ていませんから」は「no one see such a thing」の意味で、その出来事や状況が他の人に見られていないことを強調します。

\textbf{Notes (EN)}:\\
'...nannte minna mite imasen kara, ...' is a colloquial expression used when the speaker wants to convey that a certain event or situation has not been seen by others. 'Nannte' means 'such a thing', indicating that it is not something important. 'Minna mite imasen kara' means 'no one see such a thing', emphasizing that the event or situation has not been seen by others.

\newpage

\section*{686. When you put it that way, ...}

\textbf{Date}: 2026.02.25 (Wed)\\
\textbf{Index}: souiufuuniwareruto\\
\textbf{Expression (JA)}: そういうふうに(言:い)われると、...\\
\textbf{Romanization}: sou iu fuu ni iwareru to, ...\\
\textbf{Expression (EN)}: When you put it that way, ...

\textbf{Variants (JA)}: そういうふうに言われると、確かにそうですね。 / そういうふうに言われると、考えさせられますね。 / そういうふうに言われると、納得できますね。\\
\textbf{Variants (EN)}: When you put it that way, that's certainly true. / When you put it that way, it makes me think. / When you put it that way, I can understand.

\textbf{Intent (JA)}: 同意の表現, 考えの共有, 理解の促進\\
\textbf{Intent (EN)}: expression of agreement, sharing of thoughts, encouragement of understanding

\textbf{Tags (JA)}: 即時文法, 口語, 同意表現, 考え共有表現, 理解促進表現\\
\textbf{Tags (EN)}: immediate grammar, spoken, expression of agreement, expression of thought sharing, expression of encouragement of understanding

\textbf{Adjusted Expression (JA)}:\\
そういうふうに(言:い)われると、(確:たし)かにそうですね。\\
\textbf{Adjusted Expression (EN)}:\\
When you put it that way, that's certainly true.

\textbf{Example (JA)}:\\
A「この(問題:もんだい)、(複雑:ふくざつ)に(見:み)えるけど、ここに(線:せん)を(引:ひ)けば、(簡単:かんたん)だよね」 B「そういうふうに(言:い)われると、(確:たし)かにそうですね」\\
\textbf{Example (EN)}:\\
A: This quiz looks complicated, but if you draw a line here, it's simple, right? B: When you put it that way, that's certainly true.

\textbf{Notes (JA)}:\\
「そういうふうに言われると、...」は、話し手がある発言や意見に対して同意や理解を示す際に使う口語的な表現。「そういうふうに」は「when you put it that way」の意味で、前の文や状況を受けて、相手の意見や考え方を尊重するニュアンスがあります。

\textbf{Notes (EN)}:\\
'Sou iu fū ni iwareru to, ...' is a colloquial expression used when the speaker wants to show agreement or understanding towards a certain statement or opinion. 'Sou iu fū ni' means 'when you put it that way', indicating that the speaker is taking into account the previous sentence or situation and has a nuance of respecting the other person's opinion or way of thinking.

\newpage

\section*{685. Well, isn't it good?}

\textbf{Date}: 2026.02.24 (Tue)\\
\textbf{Index}: maainnjanai\\
\textbf{Expression (JA)}: まあ、いいんじゃないですか。\\
\textbf{Romanization}: Maa, ii n janai desu ka.\\
\textbf{Expression (EN)}: Well, isn't it good?

\textbf{Variants (JA)}: まあ、いいんじゃない？ / まあ、いいんじゃないでしょうか？ / まあ、いいんじゃないかな？\\
\textbf{Variants (EN)}: Well, isn't it good? / Well, isn't it good, do you think? / Well, isn't it good, I wonder?

\textbf{Intent (JA)}: 同意の表現, 肯定的な評価の共有, 提案の表現\\
\textbf{Intent (EN)}: expression of agreement, sharing of positive evaluation, expression of suggestion

\textbf{Tags (JA)}: 即時文法, 口語, 同意表現, 肯定的評価共有表現, 提案表現\\
\textbf{Tags (EN)}: immediate grammar, spoken, expression of agreement, expression of positive evaluation sharing, expression of suggestion

\textbf{Adjusted Expression (JA)}:\\
とりあえずは、それでよいのではないかと(思:おも)います。\\
\textbf{Adjusted Expression (EN)}:\\
For now, I think that would be good enough.

\textbf{Example (JA)}:\\
A「この(アイデア:あいであ)、どう(思:おも)う？」 B「まあ、いいんじゃないですか」\\
\textbf{Example (EN)}:\\
A: What do you think about this idea? B: Well, isn't it good?

\textbf{Notes (JA)}:\\
「まあ、いいんじゃないですか。」は、話し手がある事柄に対して同意や肯定的な評価を表現する際に使う口語的な表現。「まあ」は「well」の意味で、話し手がある程度の躊躇や控えめな態度を示すことがあります。「いいんじゃないですか」は「isn't it good?」の意味で、話し手がある事柄が良いと感じていることを伝えますが、同時に聞き手の意見も求めるニュアンスがあります。

\textbf{Notes (EN)}:\\
'Mā, īn janai desu ka.' is a colloquial expression used when the speaker wants to express agreement or a positive evaluation about a certain matter. 'Mā' means 'well', indicating that the speaker may have some hesitation or a modest attitude. 'Iin janai desu ka' means 'isn't it good?', conveying that the speaker feels that a certain matter is good, but also has a nuance of seeking the listener's opinion.

\newpage

\section*{684. That alone would be enough, but...}

\textbf{Date}: 2026.02.23 (Mon)\\
\textbf{Index}: soredakedeiinonini\\
\textbf{Expression (JA)}: それだけでいいのに...\\
\textbf{Romanization}: sore dake de ii no ni.\\
\textbf{Expression (EN)}: That alone would be enough, but...

\textbf{Variants (JA)}: それだけで十分なのに。 / それだけで満足なのに。 / それだけで完璧なのに。\\
\textbf{Variants (EN)}: That alone would be sufficient, but... / That alone would be satisfying, but... / That alone would be perfect, but...

\textbf{Intent (JA)}: 願望の表現, 理想の共有, 状況の説明\\
\textbf{Intent (EN)}: expression of desire, sharing of ideal, explanation of situation

\textbf{Tags (JA)}: 即時文法, 口語, 願望表現, 理想共有表現, 状況説明表現\\
\textbf{Tags (EN)}: immediate grammar, spoken, expression of desire, expression of ideal sharing, expression of situation explanation

\textbf{Adjusted Expression (JA)}:\\
それだけで(十分:じゅうぶん)です。\\
\textbf{Adjusted Expression (EN)}:\\
That alone is sufficient.

\textbf{Example (JA)}:\\
A「これに(新:あたら)しい(機能:きのう)を(追加:ついか)しておいたよ」 B「それだけでいいのに」\\
\textbf{Example (EN)}:\\
A: I added a new feature to this. B: That alone would be enough.

\textbf{Notes (JA)}:\\
「それだけでいいのに。」は、特定の要素や条件があれば十分だと感じていることを伝える際に使う口語的な表現。「それだけ」は「that alone」の意味で、特定の要素や条件を指します。「いいのに」は「would be enough」の意味で、その要素や条件があれば十分だと感じていることを示します。

\textbf{Notes (EN)}:\\
'Sore dake de ii no ni.' is a colloquial expression used when the speaker wants to convey that they feel a certain element or condition alone would be sufficient. 'Sore dake' means 'that alone', referring to a specific element or condition. 'Ii no ni' means 'would be enough', indicating that the speaker feels that if that element or condition is present, it would be sufficient.

\newpage

\section*{683. Perhaps, ...}

\textbf{Date}: 2026.02.22 (Sun)\\
\textbf{Index}: moshikashitara\\
\textbf{Expression (JA)}: もしかしたら、...\\
\textbf{Romanization}: moshikashitara, ...\\
\textbf{Expression (EN)}: Perhaps, ...

\textbf{Variants (JA)}: もしかすると、... / もしかして、...\\
\textbf{Variants (EN)}: Perhaps, ... / Maybe, ...

\textbf{Intent (JA)}: 可能性の表現, 推測の共有, 不確実性の伝達\\
\textbf{Intent (EN)}: expression of possibility, sharing of speculation, conveyance of uncertainty

\textbf{Tags (JA)}: 即時文法, 口語, 可能性表現, 推測共有表現, 不確実性伝達表現\\
\textbf{Tags (EN)}: immediate grammar, spoken, expression of possibility, expression of speculation sharing, expression of uncertainty conveyance

\textbf{Adjusted Expression (JA)}:\\
もしかしたら、それは、(本当:ほんとう)かもしれません。\\
\textbf{Adjusted Expression (EN)}:\\
Perhaps that might be true.

\textbf{Example (JA)}:\\
A「もしかしたら、(今日:きょう)、(雨:あめ)が(降:ふ)るかもしれないね」 B「うん、(傘:かさ)を(持:も)っていったほうがいいかもね」\\
\textbf{Example (EN)}:\\
A: Perhaps it might rain today. B: Yes, maybe you should take an umbrella.

\textbf{Notes (JA)}:\\
「もしかしたら、...」は、話し手がある事柄について可能性や推測を表現したいときに使う口語的な表現。「もしかしたら」は「perhaps」の意味で、不確実な状況や予測を示します。

\textbf{Notes (EN)}:\\
'Moshikashitara, ...' is a colloquial expression used when the speaker wants to express possibility or speculation about a certain matter. 'Moshikashitara' means 'perhaps', indicating an uncertain situation or prediction.

\newpage

\section*{682. It's not that I'm saying that, ...}

\textbf{Date}: 2026.02.21 (Sat)\\
\textbf{Index}: souiukotooitterunjanakute\\
\textbf{Expression (JA)}: そういうことを(言:い)ってるんじゃなくて、...\\
\textbf{Romanization}: sou iu koto o itteru n janakute, ...\\
\textbf{Expression (EN)}: It's not that I'm saying that, ...

\textbf{Variants (JA)}: そういうことを言ってるんじゃなくて、実は... / そういうことを言ってるんじゃなくて、つまり... / そういうことを言ってるんじゃなくて、要するに...\\
\textbf{Variants (EN)}: It's not that I'm saying that, but actually... / It's not that I'm saying that, but in other words... / It's not that I'm saying that, but to sum up...

\textbf{Intent (JA)}: 誤解の訂正の表現, 説明の共有, 話題の転換\\
\textbf{Intent (EN)}: expression of correction of misunderstanding, sharing of explanation, topic shift

\textbf{Tags (JA)}: 即時文法, 口語, 誤解訂正表現, 説明共有表現, 話題転換表現\\
\textbf{Tags (EN)}: immediate grammar, spoken, expression of correction of misunderstanding, expression of explanation sharing, expression of topic shift

\textbf{Adjusted Expression (JA)}:\\
そういうことを(言:い)っているわけではなくて、...\\
\textbf{Adjusted Expression (EN)}:\\
It's not that I'm saying that, ...

\textbf{Example (JA)}:\\
A 「つまり、この仕事、やりたくないんですね」 B 「いやいや、そういうことを(言:い)ってるんじゃなくて、(問題:もんだい)があるって、この(仕事:しごと)には」\\
\textbf{Example (EN)}:\\
A: So, you don't want to do this job, right? B: No, no, it's not that I'm saying that, but there are issues with this job.

\textbf{Notes (JA)}:\\
「そういうことを言ってるんじゃなくて、...」は、話し手がある発言や意見に対して誤解が生じていると感じた際に、その誤解を訂正し、実際に伝えたいことを説明するために使う口語的な表現。

\textbf{Notes (EN)}:\\
'Sou iu koto o itteru n janakute, ...' is a colloquial expression used when the speaker feels that there is a misunderstanding about a certain statement or opinion, and wants to correct that misunderstanding and explain what they actually want to convey.

\newpage

\section*{681. Isn't it...?}

\textbf{Date}: 2026.02.20 (Fri)\\
\textbf{Index}: masuyone\\
\textbf{Expression (JA)}: ますよね。\\
\textbf{Romanization}: ...masu yo ne.\\
\textbf{Expression (EN)}: Isn't it...?

\textbf{Variants (JA)}: ですよね。\\
\textbf{Variants (EN)}: Isn't it...? / Right...?

\textbf{Intent (JA)}: 同意の表現, 確認の共有, 共感の表現\\
\textbf{Intent (EN)}: expression of agreement, sharing of confirmation, expression of empathy

\textbf{Tags (JA)}: 即時文法, 口語, 同意表現, 確認共有表現\\
\textbf{Tags (EN)}: immediate grammar, spoken, expression of agreement, expression of confirmation sharing

\textbf{Adjusted Expression (JA)}:\\
そうしますよね。\\
\textbf{Adjusted Expression (EN)}:\\
That's right, isn't it?

\textbf{Example (JA)}:\\
A「(普通:ふつう)、こういうときは(簡単:かんたん)に(同意:どうい)します」 B「ますよね」\\
\textbf{Example (EN)}:\\
A: Normally, in situations like this, we easily agree. B: That's right, isn't it?

\textbf{Notes (JA)}:\\
「...ますよね」は、聞き手がある発言に対して同意する際に使う口語的な表現。

\textbf{Notes (EN)}:\\
'...masu yo ne.' is a colloquial expression used when the listener wants to agree with a certain statement.

\newpage

\section*{680. If there is someone nearby who can have such a conversation, ...}

\textbf{Date}: 2026.02.19 (Thu)\\
\textbf{Index}: sonnakaiwagadekiruhitoga\\
\textbf{Expression (JA)}: そんな(会話:かいわ)ができる(人:ひと)が、(近:ちか)くにいると、...\\
\textbf{Romanization}: sonna kaiwa ga dekiru hito ga, chikaku ni iru to, ...\\
\textbf{Expression (EN)}: If there is someone nearby who can have such a conversation, ...

\textbf{Variants (JA)}: そんな会話ができる人が、近くにいるといいな。 / そんな会話ができる人が、近くにいると助かるな。 / そんな会話ができる人が、近くにいると嬉しいな。\\
\textbf{Variants (EN)}: I hope there is someone nearby who can have such a conversation. / It would be helpful if there is someone nearby who can have such a conversation. / I would be happy if there is someone nearby who can have such a conversation.

\textbf{Intent (JA)}: 願望の表現, 理想の共有, 状況の説明\\
\textbf{Intent (EN)}: expression of desire, sharing of ideal, explanation of situation

\textbf{Tags (JA)}: 即時文法, 口語, 願望表現, 理想共有表現, 状況説明表現\\
\textbf{Tags (EN)}: immediate grammar, spoken, expression of desire, expression of ideal sharing, expression of situation explanation

\textbf{Adjusted Expression (JA)}:\\
そんな(会話:かいわ)ができる(人:ひと)が、(近:ちか)くにいるといいですね。\\
\textbf{Adjusted Expression (EN)}:\\
It would be nice if there is someone nearby who can have such a conversation.

\textbf{Example (JA)}:\\
A「(最近:さいきん)、(面白:おもしろ)い(話:はなし)を(聞:き)いた？」 B「うん、(友達:ともだち)と(深:ふか)い(会話:かいわ)ができる(人:ひと)が、(近:ちか)くにいるといいなって(思:おも)うよ」\\
\textbf{Example (EN)}:\\
A: Have you heard any interesting stories lately? B: Yes, I think it would be nice if there is someone nearby who can have deep conversations with me.

\textbf{Notes (JA)}:\\
「そんな会話)ができる人が、近くにいると、...」は、話し手が特定の会話ができる人が近くにいることを願望や理想として表現する際に使う口語的な表現。

\textbf{Notes (EN)}:\\
'Sonna kaiwa ga dekiru hito ga, chikaku ni iru to, ...' is a colloquial expression used when the speaker wants to express a desire or ideal for someone who can have a certain type of conversation to be nearby.

\newpage

\section*{679. What...?}

\textbf{Date}: 2026.02.18 (Wed)\\
\textbf{Index}: ttenani\\
\textbf{Expression (JA)}: ...って、(何:なに)？\\
\textbf{Romanization}: ...tte, nani?\\
\textbf{Expression (EN)}: What...?

\textbf{Variants (JA)}: ...って、何のこと？ / ...って、何を意味するの？ / ...って、何を指してるの？\\
\textbf{Variants (EN)}: What does ... mean? / What is ... referring to? / What is ... about?

\textbf{Intent (JA)}: 疑問の表現, 説明の要求, 情報の共有\\
\textbf{Intent (EN)}: expression of doubt, request for explanation, sharing of information

\textbf{Tags (JA)}: 即時文法, 口語, 疑問表現, 説明要求表現, 情報共有表現\\
\textbf{Tags (EN)}: immediate grammar, spoken, expression of doubt, expression of explanation request, expression of information sharing

\textbf{Adjusted Expression (JA)}:\\
...とは、(何:なん)のことですか？\\
\textbf{Adjusted Expression (EN)}:\\
What does ... mean?

\textbf{Example (JA)}:\\
A「この(機能:きのう)、(使:つか)ってみた？」 B「うん、それって、(何:なに)？」\\
\textbf{Example (EN)}:\\
A: Have you tried using this feature? B: Yes, what is that about?

\textbf{Notes (JA)}:\\
「...って、何？」の「...って」は前の文や状況を受けて、「what about ...?」や「what is ...?」の意味を持ちます。「何？」は「what?」の意味です。説明を求める表現です。

\textbf{Notes (EN)}:\\
'...tte, nani?' is a colloquial expression used when the speaker wants to ask for an explanation about a certain matter or situation. '...tte' serves as a conjunction that connects the previous sentence or situation, meaning 'what about ...?' or 'what is ...?'. 'Nani?' means 'what?', indicating a request for explanation.

\newpage

\section*{678. And then, when I realized it, ...}

\textbf{Date}: 2026.02.17 (Tue)\\
\textbf{Index}: dekigatsuitara\\
\textbf{Expression (JA)}: で、(気:き)がついたら、...\\
\textbf{Romanization}: de,  ki ga tsuitara, ...\\
\textbf{Expression (EN)}: And then, when I realized it, ...

\textbf{Variants (JA)}: で、(気:き)づいたときには、... / で、(気:き)づかないうちに、... / で、(気:き)づけば、...\\
\textbf{Variants (EN)}: And then, when I realized it, ... / And then, without noticing, ... / And then, when I noticed, ...

\textbf{Intent (JA)}: 時間の経過の表現, 出来事の説明, 状況の共有\\
\textbf{Intent (EN)}: expression of time passage, explanation of events, sharing of situation

\textbf{Tags (JA)}: 即時文法, 時間経過表現, 出来事説明表現, 状況共有表現\\
\textbf{Tags (EN)}: immediate grammar, expression of time passage, expression of event explanation, expression of situation sharing

\textbf{Adjusted Expression (JA)}:\\
それで、(気:き)がついたときには、もう(数日:すうじつ)が(経:た)っていました。\\
\textbf{Adjusted Expression (EN)}:\\
So, when I realized it, several days had already passed.

\textbf{Example (JA)}:\\
A「で、(気:き)がついたら、(朝:あさ)まで(映画:えいが)を(見:み)てしまってて」 B「で、そのまま(寝:ね)ないで(来:き)たの？」\\
\textbf{Example (EN)}:\\
A: And then, when I realized it, I had ended up watching movies until morning. B: So, you came without sleeping?

\textbf{Notes (JA)}:\\
「で、(気:き)がついたら、...」は、「で」は前の文や状況を受けて、「and then」の意味を持ちます。「気がついたら」は「when I realized it」の意味で、時間の経過や出来事の突然だったことを伝える表現です。

\textbf{Notes (EN)}:\\
'De, ki ga tsuitara, ...' is a colloquial expression where 'de' serves as a conjunction meaning 'and then', connecting the previous sentence or situation. 'Ki ga tsuitara' means 'when I realized it', conveying the passage of time or the suddenness of an event.

\newpage

\section*{677. Calm down.}

\textbf{Date}: 2026.02.16 (Mon)\\
\textbf{Index}: ochitsuite\\
\textbf{Expression (JA)}: (落:おち)ついて。\\
\textbf{Romanization}: ochitsuite.\\
\textbf{Expression (EN)}: Calm down.

\textbf{Variants (JA)}: (落:おち)ついてください。 / (落:おち)つきましょう。 / まあ、(落:おち)ついて。\\
\textbf{Variants (EN)}: Please calm down. / Calm down, please. / Calm down.

\textbf{Intent (JA)}: 安心の表現, 感情のコントロールの促し, 状況の共有\\
\textbf{Intent (EN)}: expression of reassurance, encouragement of emotional control, sharing of situation

\textbf{Tags (JA)}: 即時文法, 口語, 安心表現, 感情コントロール促し表現, 状況共有表現\\
\textbf{Tags (EN)}: immediate grammar, spoken, expression of reassurance, expression of encouragement of emotional control, expression of situation sharing

\textbf{Adjusted Expression (JA)}:\\
(落:おち)ついてください。\\
\textbf{Adjusted Expression (EN)}:\\
Please calm down.

\textbf{Example (JA)}:\\
(落:おち)ついて。(大丈夫:だいじょうぶ)だから。\\
\textbf{Example (EN)}:\\
Calm down. It's going to be okay.

\textbf{Notes (JA)}:\\
話し手が相手に対して安心感を与え、感情をコントロールするよう促す際に使う口語的な表現。「落ちつく」という複合動詞です。

\textbf{Notes (EN)}:\\
'Ochitsuite.' is a colloquial expression used when the speaker wants to provide reassurance to the listener and encourage them to control their emotions. It is a compound verb formed from 'ochi' (to fall) and 'tsuku' (to attach), meaning 'to calm down' or 'to settle down'.

\newpage

\section*{676. I don't need that kind of thing.}

\textbf{Date}: 2026.02.15 (Sun)\\
\textbf{Index}: sonnamoniranaiyo\\
\textbf{Expression (JA)}: そんなもん、いらないよ。\\
\textbf{Romanization}: sonna mon, iranai yo.\\
\textbf{Expression (EN)}: I don't need that kind of thing.

\textbf{Variants (JA)}: そんなもの、いらないよ。 / そんなこと、いらないよ。 / そんなの、いらないよ。\\
\textbf{Variants (EN)}: I don't need that kind of thing. / I don't need that kind of matter. / I don't need that kind of stuff.

\textbf{Intent (JA)}: 否定の表現, 必要性の否定, 感情の共有\\
\textbf{Intent (EN)}: expression of negation, denial of necessity, sharing of feelings

\textbf{Tags (JA)}: 即時文法, 口語, 否定表現, 必要性否定表現, 感情共有表現\\
\textbf{Tags (EN)}: immediate grammar, spoken, expression of negation, expression of denial of necessity, expression of feelings sharing

\textbf{Adjusted Expression (JA)}:\\
そんなものは、(必要:ひつよう)ありません。\\
\textbf{Adjusted Expression (EN)}:\\
I don't need that kind of thing.

\textbf{Example (JA)}:\\
A「(新:あたら)しい(ガジェット:がじぇっと)を(買:か)ったよ」 B「そんなもん、いらないよ」\\
\textbf{Example (EN)}:\\
A: I bought a new gadget. B: I don't need that kind of thing.

\textbf{Notes (JA)}:\\
「そんなもん、いらないよ。」は、話し手がある物事や状況に対して否定的な感情を表現する際に使う口語的な表現。「そんなもん」は「that kind of thing」の意味で、特定の物事や状況を指します。「いらないよ」は「I don't need it」の意味で、必要性の否定を示します。

\textbf{Notes (EN)}:\\
'Sonna mon, iranai yo.' is a colloquial expression used when the speaker wants to express a negative feeling towards a certain matter or situation. 'Sonna mon' means 'that kind of thing', referring to a specific matter or situation. 'Iranai yo' means 'I don't need it', indicating a denial of necessity.

\newpage

\section*{675. I almost...}

\textbf{Date}: 2026.02.14 (Sat)\\
\textbf{Index}: tokorodeshitayo\\
\textbf{Expression (JA)}: ...ところでしたよ。\\
\textbf{Romanization}: tokoro deshita yo.\\
\textbf{Expression (EN)}: I almost...

\textbf{Variants (JA)}: (危:あぶ)ないところでしたよ。 / (助:たす)かったところでしたよ。 / (間一髪:かんいっぱつ)でしたよ。\\
\textbf{Variants (EN)}: I was in a dangerous situation. / I was saved just in time. / It was a close call.

\textbf{Intent (JA)}: 危険の表現, 救済の共有, 緊迫感の伝達\\
\textbf{Intent (EN)}: expression of danger, sharing of rescue, conveyance of tension

\textbf{Tags (JA)}: 即時文法, 口語, 危険表現, 救済共有表現, 緊迫感伝達表現\\
\textbf{Tags (EN)}: immediate grammar, spoken, expression of danger, expression of rescue sharing, expression of tension conveyance

\textbf{Adjusted Expression (JA)}:\\
(危:あぶ)ない(状況:じょうきょう)でした。\\
\textbf{Adjusted Expression (EN)}:\\
It was a dangerous situation.

\textbf{Example (JA)}:\\
A「(車:くるま)に(轢:ひ)かれそうになったんだ」 B「ええ、(危:あぶ)ないところでしたよ」\\
\textbf{Example (EN)}:\\
A: I almost got hit by a car. B: Wow, that was a close call.

\textbf{Notes (JA)}:\\
ところでしたよ。は、話し手がある出来事や状況が非常に危険だったことを伝える際に使う口語的な表現。「ところでしたよ。」は「I almost...」の意味で、特定の行動や出来事が起こりそうになったことを示します。

\textbf{Notes (EN)}:\\
'Tokoro deshita yo.' is a colloquial expression used when the speaker wants to convey that a certain event or situation was very dangerous. 'Tokoro deshita yo.' means 'I almost...', indicating that a specific action or event was about to happen. The speaker shares with the listener that the situation was dangerous.

\newpage

\section*{674. What is this?}

\textbf{Date}: 2026.02.13 (Fri)\\
\textbf{Index}: nandakore\\
\textbf{Expression (JA)}: (何:なん)だこれ？\\
\textbf{Romanization}: nanda kore?\\
\textbf{Expression (EN)}: What is this?

\textbf{Variants (JA)}: これは何？ / これって何？ / これ、何？\\
\textbf{Variants (EN)}: What is this? / What is this thing? / What is this?

\textbf{Intent (JA)}: 驚きの表現, 疑問の共有, 状況の説明要求\\
\textbf{Intent (EN)}: expression of surprise, sharing of doubt, request for explanation of situation

\textbf{Tags (JA)}: 即時文法, 口語, 驚き表現, 疑問共有表現, 状況説明要求表現\\
\textbf{Tags (EN)}: immediate grammar, spoken, expression of surprise, expression of doubt sharing, expression of situation explanation request

\textbf{Adjusted Expression (JA)}:\\
これは(何:なん)ですか？\\
\textbf{Adjusted Expression (EN)}:\\
What is this?

\textbf{Example (JA)}:\\
A「(新:あたら)しい(ガジェット:がじぇっと)を(買:か)ったよ」 B「(何:なん)だこれ？」\\
\textbf{Example (EN)}:\\
A: I bought a new gadget. B: What is this?

\textbf{Notes (JA)}:\\
「なんだこれ？」は、話し手が予期しない物事や状況に対して驚きや疑問を感じた際に使う口語的な表現。「なんだ」は「what is」の意味で、驚きや疑問を示します。「これ」は特定の物事や状況を指します。話し手は、相手に対してその対象についての説明を求めています。変なものを見たときなどにも使われる表現です。

\textbf{Notes (EN)}:\\
'Nanda kore?' is a colloquial expression used when the speaker feels surprise or doubt towards an unexpected matter or situation. 'Nanda' means 'what is', indicating surprise or doubt. 'Kore' refers to a specific matter or situation. The speaker is requesting an explanation about that object from the listener. It is also used when seeing something strange.

\newpage

\section*{673. It's not something you/I should ask about.}

\textbf{Date}: 2026.02.12 (Thu)\\
\textbf{Index}: kikuyounakotojanaishi\\
\textbf{Expression (JA)}: (聞:き)くようなことじゃないし。\\
\textbf{Romanization}: kiku you na koto janai shi.\\
\textbf{Expression (EN)}: It's not something you/I should ask about.

\textbf{Variants (JA)}: (尋:たず)ねるようなことじゃないし。 / (話:はな)すようなことじゃないし。 / (触:さわ)れるようなことじゃないし。\\
\textbf{Variants (EN)}: It's not something you/I should inquire about. / It's not something you/I should talk about. / It's not something you/I should touch on.

\textbf{Intent (JA)}: 話題の制限の表現, プライバシーの保護, 会話の方向性の提示\\
\textbf{Intent (EN)}: expression of topic restriction, protection of privacy, presentation of conversation direction

\textbf{Tags (JA)}: 即時文法, 口語, 話題制限表現, プライバシー保護表現, 会話方向提示表現\\
\textbf{Tags (EN)}: immediate grammar, spoken, expression of topic restriction, expression of privacy protection, expression of conversation direction presentation

\textbf{Adjusted Expression (JA)}:\\
(聞:き)くべきようなことではない。\\
\textbf{Adjusted Expression (EN)}:\\
It's something that shouldn't be asked about.

\textbf{Example (JA)}:\\
A「(彼女:かのじょ)、(子供:こども)いるの？」 B「(知:し)らない。(聞:き)くようなことじゃないし、ちょっと(プライベート:ぷらいべーと)なことだし」\\
\textbf{Example (EN)}:\\
A: Does she have children? B: I don't know. It's not something you/I should ask about, and it's a bit of a private matter.

\textbf{Notes (JA)}:\\
「(聞:き)くようなことじゃないし」は、話し手がある話題や質問が適切でないと感じていることを伝える際に使う口語的な表現。

\textbf{Notes (EN)}:\\
'(Kiku) you na koto janai shi.' is a colloquial expression used when the speaker wants to convey that a certain topic or question is not appropriate.

\newpage

\section*{672. Overflowing with ... love}

\textbf{Date}: 2026.02.11 (Wed)\\
\textbf{Index}: ...ai ga afureru\\
\textbf{Expression (JA)}: ...(愛:あい)があふれる。\\
\textbf{Romanization}: ...ai ga afureru.\\
\textbf{Expression (EN)}: Overflowing with ... love

\textbf{Variants (JA)}: ...(愛情:あいじょう)があふれる。 / ...(思:おも)いやりがあふれる。 / ...(優:やさ)しさがあふれる。\\
\textbf{Variants (EN)}: Overflowing with ... affection. / Overflowing with ... compassion. / Overflowing with ... kindness.

\textbf{Intent (JA)}: 感情の表現, 愛情の共有, 肯定的な評価\\
\textbf{Intent (EN)}: expression of emotion, sharing of affection, positive evaluation

\textbf{Tags (JA)}: 即時文法, 口語, 感情表現, 愛情共有表現, 肯定的評価表現\\
\textbf{Tags (EN)}: immediate grammar, spoken, expression of emotion, expression of affection sharing, expression of positive evaluation

\textbf{Adjusted Expression (JA)}:\\
...には、(深:ふか)い(愛:あい)があります。\\
\textbf{Adjusted Expression (EN)}:\\
There is deep ... love.

\textbf{Example (JA)}:\\
A「この(映画:えいが)、(感動:かんどう)的:てき)だったね」 B「うん、...(愛:あい)があふれる(ストーリー:すとーりー)で、(心:こころ)に(響:ひび)いたよ」\\
\textbf{Example (EN)}:\\
A: This movie was moving, wasn't it? B: Yes, the story overflowing with ... love resonated with my heart.

\textbf{Notes (JA)}:\\
...(愛:あい)があふれる。は、話し手が特定の対象や状況に対して深い愛情や感情が満ち溢れていることを表現する際に使う口語的な表現。「あふれる」は「overflowing」の意味で、感情の豊かさや強さを示します。

\textbf{Notes (EN)}:\\
'... (ai) ga afureru.' is a colloquial expression used when the speaker wants to express that deep affection or emotion is overflowing towards a specific object or situation. 'Afureru' means 'overflowing', indicating the richness and intensity of emotions.

\newpage

\section*{671. It's an incredible app.}

\textbf{Date}: 2026.02.10 (Tue)\\
\textbf{Index}: tondemo nai apuri desu yo\\
\textbf{Expression (JA)}: とんでもない(アプリ:あぷり)ですよ。\\
\textbf{Romanization}: tondemo nai apuri desu yo.\\
\textbf{Expression (EN)}: It's an incredible app.

\textbf{Variants (JA)}: すごい(アプリ:あぷり)ですよ。 / 驚:おどろ)くべき(アプリ:あぷり)ですよ。 / 画期的:かっきてき)な(アプリ:あぷり)ですよ。\\
\textbf{Variants (EN)}: It's an amazing app. / It's a surprising app. / It's a groundbreaking app.

\textbf{Intent (JA)}: 賞賛の表現, 感動の共有, 肯定的な評価\\
\textbf{Intent (EN)}: expression of praise, sharing of emotion, positive evaluation

\textbf{Tags (JA)}: 即時文法, 口語, 賞賛表現, 感動共有表現, 肯定的評価表現\\
\textbf{Tags (EN)}: immediate grammar, spoken, expression of praise, expression of emotion sharing, expression of positive evaluation

\textbf{Adjusted Expression (JA)}:\\
それは、(素晴:すば)らしい(アプリ:あぷり)ですね。\\
\textbf{Adjusted Expression (EN)}:\\
That's a wonderful app.

\textbf{Example (JA)}:\\
A「この(新:あたら)しい(アプリ:あぷり)、(使:つか)ってみた？」 B「うん、とんでもない(アプリ:あぷり)ですよ。すごく(便利:べんり)で(役立:やくだ)つ」\\
\textbf{Example (EN)}:\\
A: Have you tried using this new app? B: Yes, it's an incredible app. It's very convenient and useful.

\textbf{Notes (JA)}:\\
「とんでもない...ですよ」は、話し手が特定の...に対して賞賛や感動を表現する際に使う口語的な表現。「とんでもない」は「incredible」や「amazing」の意味で、肯定的な評価を示します。

\textbf{Notes (EN)}:\\
'Tondemo nai... desu yo' is a colloquial expression used when the speaker wants to express praise or amazement towards a specific ... . 'Tondemo nai' means 'incredible' or 'amazing', indicating a positive evaluation.

\newpage

\section*{670. There are things like ...}

\textbf{Date}: 2026.02.09 (Mon)\\
\textbf{Index}: ...toka arun desu yo\\
\textbf{Expression (JA)}: ...とかあるんですよ。\\
\textbf{Romanization}: ...toka arun desu yo.\\
\textbf{Expression (EN)}: There are things like ...

\textbf{Variants (JA)}: ...などありますよ。 / ...とか(存在:そんざい)するんですよ。 / ...とか(見:み)かけるんですよ。\\
\textbf{Variants (EN)}: There are things such as ... / There exist things like ... / You can find things like ...

\textbf{Intent (JA)}: 例示の表現, 情報の共有, 状況の説明\\
\textbf{Intent (EN)}: expression of exemplification, sharing of information, explanation of situation

\textbf{Tags (JA)}: 即時文法, 口語, 例示表現, 情報共有表現, 状況説明表現\\
\textbf{Tags (EN)}: immediate grammar, spoken, expression of exemplification, expression of information sharing, expression of situation explanation

\textbf{Adjusted Expression (JA)}:\\
...などが(存在:そんざい)します。\\
\textbf{Adjusted Expression (EN)}:\\
There exist things such as ...

\textbf{Example (JA)}:\\
A「(最近:さいきん)、(面白:おもしろ)い(本:ほん)、(読:よ)んだ？」 B「うん、(今迄:いままで)にあまり(見掛:みか)けない(ジャンル:じゃんる)とかあるんですよ」\\
\textbf{Example (EN)}:\\
A: Have you read any interesting books lately? B: Yes, there are genres that I haven't seen much until now.

\textbf{Notes (JA)}:\\
...とかあるんですよ。は、話し手が特定の例や種類を挙げて情報を共有する際に使う口語的な表現。「...とか」は「things like ...」の意味で、具体的な例を示します。

\textbf{Notes (EN)}:\\
'... toka arun desu yo.' is a colloquial expression used when the speaker wants to share information by providing specific examples or types. 

\newpage

\section*{669. Before I knew it, I was ... .}

\textbf{Date}: 2026.02.08 (Sun)\\
\textbf{Index}: kizuitarashitemashita\\
\textbf{Expression (JA)}: (気:き)づいたら、...してました。\\
\textbf{Romanization}: kizuitara ... shitemashita.\\
\textbf{Expression (EN)}: Before I knew it, I was ... .

\textbf{Variants (JA)}: (気:き)づいたときには、...していました。 / (気:き)づかないうちに、...していました。 / (気:き)づいた瞬間に、...していました。\\
\textbf{Variants (EN)}: When I realized it, I was ... . / Without noticing, I was ... . / The moment I noticed, I was ... .

\textbf{Intent (JA)}: 無意識の行動の表現, 時間の経過の共有, 出来事の説明\\
\textbf{Intent (EN)}: expression of unconscious actions, sharing of time passage, explanation of events

\textbf{Tags (JA)}: 即時文法, 口語, 無意識行動表現, 時間経過共有表現, 出来事説明表現\\
\textbf{Tags (EN)}: immediate grammar, spoken, expression of unconscious actions, expression of time passage sharing, expression of event explanation

\textbf{Adjusted Expression (JA)}:\\
(気:き)づいたときには、(仕事:しごと)を(終:お)えていました。\\
\textbf{Adjusted Expression (EN)}:\\
When I realized it, I had finished my work.

\textbf{Example (JA)}:\\
A「(昨日:きのう)、(寝:ね)るつもりが、(気:き)づいたら、(朝:あさ)まで(映画:えいが)を(見:み)てしまってたよ」 B「それは(大変:たいへん)だね」\\
\textbf{Example (EN)}:\\
A: I intended to sleep yesterday, but before I knew it, I ended up watching movies until morning. B: That's tough.

\textbf{Notes (JA)}:\\
(気:き)づいたら、...してました。は、話し手が無意識のうちにある行動をしてしまったことを伝える際に使う口語的な表現。「気づいたら」は「before I knew it」の意味で、時間の経過や出来事の突然の認識を示します。「...してました。」は「I was ...」の意味で、特定の行動を指します。

\textbf{Notes (EN)}:\\
'Kizuitara, ... shitemashita.' is a colloquial expression used when the speaker wants to convey that they performed a certain action unconsciously. 'Kizuitara' means 'before I knew it', indicating the passage of time or sudden recognition of an event. '... shitemashita.' means 'I was ...', referring to a specific action.

\newpage

\section*{668. While using ...}

\textbf{Date}: 2026.02.07 (Sat)\\
\textbf{Index}: ...o tsukai tsutsu,\\
\textbf{Expression (JA)}: ...を(使:つか)いつつ、\\
\textbf{Romanization}: ...o tsukai tsutsu,\\
\textbf{Expression (EN)}: While using ...

\textbf{Variants (JA)}: ...を(利用:りよう)しつつ、 / ...を(活用:かつよう)しつつ、 / ...を(用:もち)いながら、\\
\textbf{Variants (EN)}: While utilizing ... / While making use of ... / While employing ...

\textbf{Intent (JA)}: 同時進行の表現, 行動の説明, 状況の共有\\
\textbf{Intent (EN)}: expression of simultaneous progression, explanation of action, sharing of situation

\textbf{Tags (JA)}: 即時文法, 口語, 同時進行表現, 行動説明表現, 状況共有表現\\
\textbf{Tags (EN)}: immediate grammar, spoken, expression of simultaneous progression, expression of action explanation, expression of situation sharing

\textbf{Adjusted Expression (JA)}:\\
...を(使:つか)いながら、(作業:さぎょう)をしています。\\
\textbf{Adjusted Expression (EN)}:\\
I am working while using ...

\textbf{Example (JA)}:\\
A「(新:あたら)しい(ソフト:そふと)、(使:つか)ってみた？」 B「うん、...(機能:きのう)、(使:つか)いつつ、(操作:そうさ)に(慣:な)れてるところ」\\
\textbf{Example (EN)}:\\
A: Have you tried using the new software? B: Yes, I'm getting used to the operation while using its features.

\textbf{Notes (JA)}:\\
...を使いつつ、は、話し手がある行動や状況が同時に進行していることを説明する際に使う口語的な表現。「...を使いつつ、」は「while using ...」の意味で、特定の行動や状況を指します。

\textbf{Notes (EN)}:\\
'... o tsukai tsutsu,' is a colloquial expression used when the speaker wants to explain that a certain action or situation is progressing simultaneously. '... o tsukai tsutsu,' means 'while using ...', referring to a specific action or situation.

\newpage

\section*{667. Just because I/you/we...,}

\textbf{Date}: 2026.02.06 (Fri)\\
\textbf{Index}: shita bakari ni\\
\textbf{Expression (JA)}: ...したばかりに、\\
\textbf{Romanization}: ... shita bakari ni,\\
\textbf{Expression (EN)}: Just because I/you/we...,

\textbf{Variants (JA)}: ...しただけで、 / ...したせいで、 / ...したために、\\
\textbf{Variants (EN)}: Just by doing..., / Because of doing..., / Due to doing...,

\textbf{Intent (JA)}: 原因の表現, 結果の説明, 状況の共有\\
\textbf{Intent (EN)}: expression of cause, explanation of result, sharing of situation

\textbf{Tags (JA)}: 即時文法, 口語, 原因表現, 結果説明表現, 状況共有表現\\
\textbf{Tags (EN)}: immediate grammar, spoken, expression of cause, expression of result explanation, expression of situation sharing

\textbf{Adjusted Expression (JA)}:\\
...したために、(悪:わる)い(結果:けっか)になってしまいました。\\
\textbf{Adjusted Expression (EN)}:\\
Due to doing..., a bad result occurred.

\textbf{Example (JA)}:\\
A「(昨日:きのう)、(夜遅:よるおそ)くまで(遊:あそ)んだばかりに、(眠:ねむ)くてたまらないよ」 B「それは(大変:たいへん)」\\
\textbf{Example (EN)}:\\
A: Just because I played until late at night yesterday, I'm so sleepy. B: That's tough.

\textbf{Notes (JA)}:\\
「...したばかりに、」は、話し手がある行動や出来事が原因で特定の結果や状況が生じたことを説明する際に使う口語的な表現。

\textbf{Notes (EN)}:\\
'... shita bakari ni,' is a colloquial expression used when the speaker wants to explain that a certain action or event caused a specific result or situation.

\newpage

\section*{666. I'm totally fine.}

\textbf{Date}: 2026.02.05 (Thu)\\
\textbf{Index}: zenzen heiki desu\\
\textbf{Expression (JA)}: ぜんぜん(平気:へいき)です。\\
\textbf{Romanization}: zenzen heiki desu.\\
\textbf{Expression (EN)}: I'm totally fine.

\textbf{Variants (JA)}: まったく(問題:もんだい)ないです。 / 全然(大丈夫:だいじょうぶ)です。 / ぜんぜん(気:き)にしていません。\\
\textbf{Variants (EN)}: There's absolutely no problem. / I'm totally okay. / I don't mind at all.

\textbf{Intent (JA)}: 安心の表現, 自己肯定の伝達, 肯定的な状態の共有\\
\textbf{Intent (EN)}: expression of reassurance, conveyance of self-affirmation, sharing of positive state

\textbf{Tags (JA)}: 即時文法, 口語, 安心表現, 自己肯定伝達表現, 肯定的状態共有表現\\
\textbf{Tags (EN)}: immediate grammar, spoken, expression of reassurance, expression of self-affirmation conveyance, expression of positive state sharing

\textbf{Adjusted Expression (JA)}:\\
まったく(問題:もんだい)ありません。\\
\textbf{Adjusted Expression (EN)}:\\
There is absolutely no problem.

\textbf{Example (JA)}:\\
A「(心配:しんぱい)しないで、(大丈夫:だいじょうぶ)？」 B「うん、ぜんぜん(平気:へいき)です」\\
\textbf{Example (EN)}:\\
A: Don't worry, are you okay? B: Yes, I'm totally fine.

\textbf{Notes (JA)}:\\
「ぜんぜん平気です。」は、話し手が自分の状態が全く問題ないことを伝え、安心感を与える際に使う口語的な表現。「ぜんぜん」は「totally」や「absolutely」の意味で、強調のニュアンスを持ちます。「平気です。」は「I'm fine」の意味で、肯定的な状態を示します。話し手は、相手に対して自分の状態が良好であることを共有しています。

\textbf{Notes (EN)}:\\
'Zenzen heiki desu.' is a colloquial expression used when the speaker wants to convey that their condition is completely fine and provide reassurance. 'Zenzen' means 'totally' or 'absolutely', carrying a nuance of emphasis. 'Heiki desu.' means 'I'm fine', indicating a positive state. The speaker shares with the listener that their condition is good.

\newpage

\section*{665. It's totally fine.}

\textbf{Date}: 2026.02.04 (Wed)\\
\textbf{Index}: zenzen ii desu yo\\
\textbf{Expression (JA)}: ぜんぜんいいですよ。\\
\textbf{Romanization}: zenzen ii desu yo.\\
\textbf{Expression (EN)}: It's totally fine.

\textbf{Variants (JA)}: まったく(問題:もんだい)ないですよ。 / 全然(大丈夫:だいじょうぶ)ですよ。 / ぜんぜん(気:き)にしないでください。\\
\textbf{Variants (EN)}: There's absolutely no problem. / It's totally okay. / Please don't worry about it at all.

\textbf{Intent (JA)}: 安心の表現, 許容の伝達, 肯定的な応答\\
\textbf{Intent (EN)}: expression of reassurance, conveyance of acceptance, positive response

\textbf{Tags (JA)}: 即時文法, 口語, 安心表現, 許容伝達表現, 肯定的応答表現\\
\textbf{Tags (EN)}: immediate grammar, spoken, expression of reassurance, expression of acceptance conveyance, expression of positive response

\textbf{Adjusted Expression (JA)}:\\
まったく(問題:もんだい)ありませんよ。\\
\textbf{Adjusted Expression (EN)}:\\
There is absolutely no problem.

\textbf{Example (JA)}:\\
A「(遅:おく)れてごめんね」 B「ううん、ぜんぜんいいですよ」\\
\textbf{Example (EN)}:\\
A: Sorry for being late. B: No, it's totally fine.

\textbf{Notes (JA)}:\\
「ぜんぜんいいですよ。」は、話し手が相手に対して安心感を与え、許容する意図を伝える際に使う口語的な表現。「ぜんぜん」は「totally」や「absolutely」の意味で、強調のニュアンスを持ちます。「いいですよ。」は「it's fine」の意味で、肯定的な応答を示します。話し手は、相手に対してその行動や状況が問題ないことを共有しています。

\textbf{Notes (EN)}:\\
'Zenzen ii desu yo.' is a colloquial expression used when the speaker wants to provide reassurance to the listener and convey acceptance. 'Zenzen' means 'totally' or 'absolutely', carrying a nuance of emphasis. 'Ii desu yo.' means 'it's fine', indicating a positive response. The speaker shares with the listener that their actions or situation are not a problem.

\newpage

\section*{664. Isn't it amazing?}

\textbf{Date}: 2026.02.03 (Tue)\\
\textbf{Index}: sugoi janai desu ka\\
\textbf{Expression (JA)}: すごいじゃないですか。\\
\textbf{Romanization}: sugoi ja nai desu ka.\\
\textbf{Expression (EN)}: Isn't it amazing?

\textbf{Variants (JA)}: (素晴:すば)らしいじゃないですか。 / (驚:おどろ)くべきじゃないですか。 / (感動:かんどう)的:てき)じゃないですか。\\
\textbf{Variants (EN)}: Isn't it wonderful? / Isn't it surprising? / Isn't it impressive?

\textbf{Intent (JA)}: 賞賛の表現, 感動の共有, 肯定的な評価\\
\textbf{Intent (EN)}: expression of praise, sharing of emotion, positive evaluation

\textbf{Tags (JA)}: 即時文法, 口語, 賞賛表現, 感動共有表現, 肯定的評価表現\\
\textbf{Tags (EN)}: immediate grammar, spoken, expression of praise, expression of emotion sharing, expression of positive evaluation

\textbf{Adjusted Expression (JA)}:\\
それは、(素晴:すば)らしいですね。\\
\textbf{Adjusted Expression (EN)}:\\
That's wonderful.

\textbf{Example (JA)}:\\
A「(新:あたら)しい(車:くるま)、(買:か)ったんだ」 B「へえ、すごいじゃないですか。どんな(車:くるま)？」\\
\textbf{Example (EN)}:\\
A: I bought a new car. B: Wow, isn't it amazing? What kind of car is it?

\textbf{Notes (JA)}:\\
「すごいじゃないですか。」は、話し手が相手の行動や成果に対して賞賛や感動を表現する際に使う口語的な表現。「すごい」は「amazing」や「wonderful」の意味で、肯定的な評価を示します。「じゃないですか。」は強調のニュアンスを持ち、話し手の感情を強く伝えます。話し手は、相手に対してその対象が高く評価されていることを共有しています。

\textbf{Notes (EN)}:\\
'Sugoi ja nai desu ka.' is a colloquial expression used when the speaker wants to express praise or emotion towards the listener's actions or achievements. 'Sugoi' means 'amazing' or 'wonderful', indicating a positive evaluation. 'Ja nai desu ka.' carries a nuance of emphasis, strongly conveying the speaker's feelings. The speaker shares with the listener that the object in question is being highly evaluated.

\newpage

\section*{663. I take walks and do various things,...}

\textbf{Date}: 2026.02.02 (Mon)\\
\textbf{Index}: sanpo shitari shiteru n desu kedo\\
\textbf{Expression (JA)}: (散歩:さんぽ)したりしてるんですけど、\\
\textbf{Romanization}: sanpo shitari shiteru n desu kedo,\\
\textbf{Expression (EN)}: I take walks and do various things,...

\textbf{Variants (JA)}: (散歩:さんぽ)したり、(買:か)い(物:もの)したりしてるんですけど、 / (運動:うんどう)したりしてるんですけど、 / (読書:どくしょ)したりしてるんですけど、\\
\textbf{Variants (EN)}: I take walks and do shopping,... / I exercise and do various things,... / I read books and do various things,...

\textbf{Intent (JA)}: 活動の共有, 日常生活の説明, 会話の導入\\
\textbf{Intent (EN)}: sharing of activities, explanation of daily life, introduction to conversation

\textbf{Tags (JA)}: 即時文法, 口語, 活動共有表現, 日常生活説明表現, 会話導入表現\\
\textbf{Tags (EN)}: immediate grammar, spoken, expression of activity sharing, expression of daily life explanation, expression of conversation introduction

\textbf{Adjusted Expression (JA)}:\\
たとえば、(散歩:さんぽ)、(読書:どくしょ)、(料理:りょうり)、などさまざまなことをしています。\\
\textbf{Adjusted Expression (EN)}:\\
For example, I do various things such as taking walks, reading books, and cooking.

\textbf{Example (JA)}:\\
A「(最近:さいきん)、(暇:ひま)なの？」 B「うん、(散歩:さんぽ)したりしてるんですけど、(何:なに)かいい(趣味:しゅみ)があれば(教:おし)えてほしいな」\\
\textbf{Example (EN)}:\\
A: Are you free lately? B: Yes, I take walks and do various things, but if you have any good hobbies, I'd like you to tell me.

\textbf{Notes (JA)}:\\
(散歩:さんぽ)したりしてるんですけど、は、さまざまなことをしていることを伝える際に使う口語的な表現。「したりしてるんですけど、」は「doing various things,...」の意味で、具体的な活動を列挙する前置きとして使われます。この場合の散歩は、いろいろなことの一例として挙げられています。

\textbf{Notes (EN)}:\\
'(Sanpo) shitari shiteru n desu kedo,' is a colloquial expression used when the speaker wants to convey that they are doing various things. 'Shitari shiteru n desu kedo,' means 'doing various things,...', and is used as a prelude to listing specific activities. In this case, taking walks is mentioned as one example of various activities.

\newpage

\section*{662. Just right.}

\textbf{Date}: 2026.02.01 (Sun)\\
\textbf{Index}: choudoii\\
\textbf{Expression (JA)}: (丁度:ちょうど)いい。\\
\textbf{Romanization}: choudo ii.\\
\textbf{Expression (EN)}: Just right.

\textbf{Variants (JA)}: (ちょうど:丁度)いいね。 / (ぴったり)だね。 / (適切:てきせつ)だね。\\
\textbf{Variants (EN)}: Just right. / It's perfect. / It's appropriate.

\textbf{Intent (JA)}: 適合性の表現, 満足感の共有, 評価の提示\\
\textbf{Intent (EN)}: expression of suitability, sharing of satisfaction, presentation of evaluation

\textbf{Tags (JA)}: 即時文法, 口語, 適合性表現, 満足感共有表現, 評価提示表現\\
\textbf{Tags (EN)}: immediate grammar, spoken, expression of suitability, expression of satisfaction sharing, expression of evaluation presentation

\textbf{Adjusted Expression (JA)}:\\
(ちょうど:丁度)よいです。\\
\textbf{Adjusted Expression (EN)}:\\
It is just right.

\textbf{Example (JA)}:\\
A「この(服:ふく)、(サイズ:さいず)、(丁度:ちょうど)いいね」 B「うん、(似合:にあ)ってるよ」\\
\textbf{Example (EN)}:\\
A: This clothing is just right in size. B: Yes, it suits you.

\textbf{Notes (JA)}:\\
(丁度:ちょうど)いい。は、話し手がある物事や状況が適切であり、満足できる状態であることを伝える際に使う口語的な表現。「(丁度:ちょうど)いい。」は「just right」の意味で、適合性や満足感を示します。話し手は、相手に対してその対象が評価されていることを共有しています。

\textbf{Notes (EN)}:\\
'(Choudo) ii.' is a colloquial expression used when the speaker wants to convey that a certain matter or situation is appropriate and satisfactory. '(Choudo) ii.' means 'just right', indicating suitability and satisfaction. The speaker shares with the listener that the object in question is being evaluated positively.

\newpage

\section*{661. Ah, somehow, I said something weird,}

\textbf{Date}: 2026.01.31 (Sat)\\
\textbf{Index}: a nanka hen na koto icchatte\\
\textbf{Expression (JA)}: ああ、なんか、(変:へん)なこと、(言:い)っちゃって、\\
\textbf{Romanization}: ā, nanka, hen na koto, icchatte,\\
\textbf{Expression (EN)}: Ah, somehow, I said something weird,

\textbf{Variants (JA)}: ああ、なんか、おかしなこと、(言:い)っちゃって、 / ああ、なんか、(変:へん)なこと、(話:はな)しちゃって、 / ああ、なんか、(変:へん)なこと、(言:い)ってしまって、\\
\textbf{Variants (EN)}: Ah, somehow, I said something strange, / Ah, somehow, I talked about something weird, / Ah, somehow, I ended up saying something weird,

\textbf{Intent (JA)}: 謝罪の表現, 自己反省の共有, 会話の和らげ\\
\textbf{Intent (EN)}: expression of apology, sharing of self-reflection, softening of conversation

\textbf{Tags (JA)}: 即時文法, 口語, 謝罪表現, 自己反省共有表現, 会話和らげ表現\\
\textbf{Tags (EN)}: immediate grammar, spoken, expression of apology, expression of self-reflection sharing, expression of conversation softening

\textbf{Adjusted Expression (JA)}:\\
(少々:しょうしょう)、(不適切:ふてきせつ)な(発言:はつげん)しまして、(申:もう)し(訳:わけ)ありませんでした。\\
\textbf{Adjusted Expression (EN)}:\\
I made a slightly inappropriate remark, and I apologize.

\textbf{Example (JA)}:\\
A「(昨日:きのう)、(会:あ)ったとき、(変:へん)なこと、(言:い)っちゃって、ごめんね」 B「ううん、(気:き)にしないで」\\
\textbf{Example (EN)}:\\
A: Sorry for saying something weird when we met yesterday. B: No, don't worry about it.

\textbf{Notes (JA)}:\\
「ああ、なんか、変なこと、言っちゃって、」は、話し手が自分の発言が不適切または奇妙であったことを認識し、そのことについて謝罪や自己反省を表現する際に使う口語的な表現。「ああ」は驚きや気づきを示し、「なんか」は曖昧さを加えます。「変なこと、言っちゃって、」は「I said something weird」の意味で、話し手の発言に対する自己評価を示します。話し手は、相手に対して自分の発言について謝罪し、会話の雰囲気を和らげようとしています。

\textbf{Notes (EN)}:\\
'Ah, nanka, hen na koto, icchatte,' is a colloquial expression used when the speaker recognizes that their statement was inappropriate or strange and wants to express an apology or self-reflection about it. 'Ah' indicates surprise or realization, while 'nanka' adds ambiguity. 'Hen na koto, icchatte,' means 'I said something weird', indicating the speaker's self-assessment of their statement. The speaker apologizes to the listener for their statement and tries to soften the atmosphere of the conversation.

\newpage

\section*{660. I have a feeling...}

\textbf{Date}: 2026.01.30 (Fri)\\
\textbf{Index}: sonnakanjigasuru\\
\textbf{Expression (JA)}: そんな(感:かん)じがする。\\
\textbf{Romanization}: sonna kanji ga suru.\\
\textbf{Expression (EN)}: I have a feeling...

\textbf{Variants (JA)}: なんだか感じがする。 / どうも感じがする。 / 妙な感じがする。\\
\textbf{Variants (EN)}: I have a feeling... / Somehow, I feel... / I have a strange feeling...

\textbf{Intent (JA)}: 直感の表現, 予感の共有, 感覚の伝達\\
\textbf{Intent (EN)}: expression of intuition, sharing of premonition, conveyance of sensation

\textbf{Tags (JA)}: 即時文法, 口語, 直感表現, 予感共有表現, 感覚伝達表現\\
\textbf{Tags (EN)}: immediate grammar, spoken, expression of intuition, expression of premonition sharing, expression of sensation conveyance

\textbf{Adjusted Expression (JA)}:\\
なんとなくそんな(感:かん)じがします。\\
\textbf{Adjusted Expression (EN)}:\\
Somehow, I have that kind of feeling.

\textbf{Example (JA)}:\\
A「(彼女:かのじょ)、(来:こ)ないかなあ？」 B「うん、そんな(感:かん)じがする。(来:き)たとして、ずいぶん(遅:おく)れて(来:く)るんじゃない？」\\
\textbf{Example (EN)}:\\
A: I wonder if she won't come. B: Yes, I have a feeling that even if she comes, she'll be quite late.

\textbf{Notes (JA)}:\\
(そんな感じがする。)は、話し手がある状況や出来事に対して直感的な感覚を持っていることを伝える際に使う口語的な表現。「そんな」は「that kind of」の意味で、特定の状況や感覚を指します。「感じがする。」は「I have a feeling」の意味で、話し手の直感や予感を示します。話し手は、相手に対して自分の感覚を共有し、状況に対する直感的な理解を伝えています。

\textbf{Notes (EN)}:\\
'Sonna kanji ga suru.' is a colloquial expression used when the speaker wants to convey that they have an intuitive feeling about a certain situation or event. 'Sonna' means 'that kind of', referring to a specific situation or sensation. 'Kanji ga suru.' means 'I have a feeling', indicating the speaker's intuition or premonition. The speaker shares their sensation with the listener and conveys their intuitive understanding of the situation.

\newpage

\section*{659. Switch gears there.}

\textbf{Date}: 2026.01.29 (Thu)\\
\textbf{Index}: kirikaete\\
\textbf{Expression (JA)}: そこは(切:き)り(替:か)えて。\\
\textbf{Romanization}: soko wa kirikaete.\\
\textbf{Expression (EN)}: Switch gears there.

\textbf{Variants (JA)}: そこは(変:か)えて。 / そこは(転換:てんかん)して。 / そこは(切替:きりか)しなさい。\\
\textbf{Variants (EN)}: Switch gears there. / Change it there. / Make a switch there.

\textbf{Intent (JA)}: 行動の指示, 状況の変化の促進, 柔軟な対応の提案\\
\textbf{Intent (EN)}: instruction of action, promotion of situation change, suggestion of flexible response

\textbf{Tags (JA)}: 即時文法, 口語, 行動指示表現, 状況変化促進表現, 柔軟対応提案表現\\
\textbf{Tags (EN)}: immediate grammar, spoken, expression of action instruction, expression of situation change promotion, expression of flexible response suggestion

\textbf{Adjusted Expression (JA)}:\\
そういうときは、(考:かんが)え(方:かた)を(柔軟:じゅうなん)にして(対応:たいおう)してください。\\
\textbf{Adjusted Expression (EN)}:\\
In such cases, please think flexibly and respond accordingly.

\textbf{Example (JA)}:\\
A「(仕事:しごと)が(忙:いそが)しくて、(疲:つか)れたよ」 B「そういうときは、そこは(切:き)り(替:か)えて、(気分:きぶん)を(変:か)えたほうがいいよ」\\
\textbf{Example (EN)}:\\
A: I'm busy with work and tired. B: In such cases, switch gears there and change your mood.

\textbf{Notes (JA)}:\\
「そこは切り替えて。」は、話し手が相手に対して特定の状況や感情から別のものに意識や行動を変えるよう促す際に使う口語的な表現。「そこは」は「there」の意味で、特定の状況や感情を指します。「切り替えて。」は「switch gears」の意味で、意識や行動を変えることを示します。話し手は、相手に対して柔軟な対応を提案し、状況の変化を促しています。

\textbf{Notes (EN)}:\\
'Soko wa kirikaete.' is a colloquial expression used when the speaker wants to encourage the listener to change their focus or actions from a specific situation or emotion to another. 'Soko wa' means 'there', referring to a specific situation or emotion. 'Kirikaete.' means 'switch gears', indicating a change in focus or actions. The speaker suggests a flexible response to the listener and promotes a change in the situation.

\newpage

\section*{658. That's wrong.}

\textbf{Date}: 2026.01.28 (Wed)\\
\textbf{Index}: sorechigauyo\\
\textbf{Expression (JA)}: それ、(違:ちが)うよ。\\
\textbf{Romanization}: sore, chigau yo.\\
\textbf{Expression (EN)}: That's wrong.

\textbf{Variants (JA)}: それ、(間違:まちが)ってるよ。 / それ、(違:ちが)いますよ。 / それ、(誤解:ごかい)だよ。\\
\textbf{Variants (EN)}: That's wrong. / That's incorrect. / That's a misunderstanding.

\textbf{Intent (JA)}: 誤りの指摘, 訂正の表現, 理解の促進\\
\textbf{Intent (EN)}: pointing out mistakes, expression of correction, promotion of understanding

\textbf{Tags (JA)}: 即時文法, 口語, 誤り指摘表現, 訂正表現, 理解促進表現\\
\textbf{Tags (EN)}: immediate grammar, spoken, expression of mistake pointing out, expression of correction, expression of understanding promotion

\textbf{Adjusted Expression (JA)}:\\
それは、(間違:まちが)っています。\\
\textbf{Adjusted Expression (EN)}:\\
That is incorrect.

\textbf{Example (JA)}:\\
A「これは、(私:わたし)がもらったもの」 B「ううん、それ、(違:ちが)うよ。みんながもらったものだよ」\\
\textbf{Example (EN)}:\\
A: This is mine. I got this. B: No, that's wrong. Everyone got it.

\textbf{Notes (JA)}:\\
「それ、違うよ。」は、話し手が相手の発言や行動に対して誤りを指摘し、訂正を行う際に使う口語的な表現。「それ」は「that」の意味で、特定の対象を指します。「違うよ。」は「that's wrong」の意味で、相手の認識や行動が正しくないことを伝えています。話し手は、相手に対して正しい情報を提供し、理解を促進しています。

\textbf{Notes (EN)}:\\
'Sore, chigau yo.' is a colloquial expression used when the speaker wants to point out a mistake in the listener's statement or action and make a correction. 'Sore' means 'that', referring to a specific object. 'Chigau yo.' means 'that's wrong', conveying that the listener's recognition or action is incorrect. The speaker provides the correct information to the listener and promotes understanding.

\newpage

\section*{657. Wasn't there...?}

\textbf{Date}: 2026.01.27 (Tue)\\
\textbf{Index}: arimasen deshita kke\\
\textbf{Expression (JA)}: ...ありませんでしたっけ？\\
\textbf{Romanization}: ... arimasen deshita kke?\\
\textbf{Expression (EN)}: Wasn't there...?

\textbf{Variants (JA)}: ...なかったっけ？ / ...じゃなかったっけ？ / ...ではなかったっけ？\\
\textbf{Variants (EN)}: Wasn't there...? / Wasn't it...? / Was it not...?

\textbf{Intent (JA)}: 記憶の確認, 過去の出来事の照会, 情報の再確認\\
\textbf{Intent (EN)}: confirmation of memory, inquiry about past events, reconfirmation of information

\textbf{Tags (JA)}: 即時文法, 口語, 記憶確認表現, 過去出来事照会表現, 情報再確認表現\\
\textbf{Tags (EN)}: immediate grammar, spoken, expression of memory confirmation, expression of past event inquiry, expression of information reconfirmation

\textbf{Adjusted Expression (JA)}:\\
...はありませんでしたか？\\
\textbf{Adjusted Expression (EN)}:\\
Wasn't there...?

\textbf{Example (JA)}:\\
A「(昨日:きのう)、あの(人:ひと)、(会:あ)ったよね」 B「うん、(確:たし)かに...(名前:なまえ)、(思:おも)い(出:だ)せないけど、...(彼:かれ)、(新:あたら)しい(部長:ぶちょう)じゃなかったっけ？」\\
\textbf{Example (EN)}:\\
A: I met that person yesterday, right? B: Yeah, I can't quite remember the name, but wasn't he the new department manager?

\textbf{Notes (JA)}:\\
...ありませんでしたっけ？は、話し手が過去の出来事や情報に関して、自分の記憶を確認し、相手に再確認を求める際に使う口語的な表現。「...ありませんでしたっけ？」は「wasn't there...?」の意味で、過去の事柄について尋ねています。話し手は、相手に対して自分の記憶が正しいかどうかを確認し、情報を再度確認しています。

\textbf{Notes (EN)}:\\
'... arimasen deshita kke?' is a colloquial expression used when the speaker wants to confirm their memory regarding past events or information and seek reconfirmation from the listener. '... arimasen deshita kke?' means 'wasn't there...?', asking about past matters. The speaker is checking with the listener to see if their memory is correct and reconfirming the information.

\newpage

\section*{656. Could you...?}

\textbf{Date}: 2026.01.26 (Mon)\\
\textbf{Index}: tekurenai\\
\textbf{Expression (JA)}: ...てくれない？\\
\textbf{Romanization}: ... te kurenai?\\
\textbf{Expression (EN)}: Could you...?

\textbf{Variants (JA)}: ...てくれませんか？ / ...てもらえますか？ / ...ていただけますか？\\
\textbf{Variants (EN)}: Could you...? / Would you mind...? / Could I ask you to...?

\textbf{Intent (JA)}: 依頼の表現, 丁寧なお願い, 協力の要請\\
\textbf{Intent (EN)}: expression of request, polite request, solicitation of cooperation

\textbf{Tags (JA)}: 即時文法, 口語, 依頼表現, 丁寧お願い表現, 協力要請表現\\
\textbf{Tags (EN)}: immediate grammar, spoken, expression of request, expression of polite request, expression of solicitation of cooperation

\textbf{Adjusted Expression (JA)}:\\
...ていただけますか？\\
\textbf{Adjusted Expression (EN)}:\\
Could I ask you to...?

\textbf{Example (JA)}:\\
A「(手伝:てつだ)ってくれる？」 B「うん、いいよ」\\
\textbf{Example (EN)}:\\
A: Could you help me? B: Sure, no problem.

\textbf{Notes (JA)}:\\
...てくれない？は、話し手が相手に対して何かをしてもらいたいときに使うかなりカジュアルな依頼表現。「...てくれない？」は「could you...?」の意味で、相手に対してお願いしています。話し手は、相手に協力を求める際にこの表現を使い、親しみやすい雰囲気を作り出しています。

\textbf{Notes (EN)}:\\
'... te kurenai?' is a fairly casual request expression used when the speaker wants the listener to do something for them. '... te kurenai?' means 'could you...?', making a request to the listener. The speaker uses this expression when seeking cooperation from the listener, creating a friendly atmosphere.

\newpage

\section*{655. In the end, isn't it...?}

\textbf{Date}: 2026.01.25 (Sun)\\
\textbf{Index}: kekkyokujanai katte\\
\textbf{Expression (JA)}: (結局:けっきょく)、...じゃないかって。\\
\textbf{Romanization}: kekkyoku, ... ja nai ka tte.\\
\textbf{Expression (EN)}: In the end, isn't it...?

\textbf{Variants (JA)}: (最終的:さいしゅうてき)に、...じゃないかって。 / (結局:けっきょく)のところ、...じゃないかって。 / (要:よう)するに、...じゃないかって。\\
\textbf{Variants (EN)}: In the end, isn't it...? / Ultimately, isn't it...? / In short, isn't it...?

\textbf{Intent (JA)}: 結論の表現, 要約の提示, 確認の意図\\
\textbf{Intent (EN)}: expression of conclusion, presentation of summary, intention of confirmation

\textbf{Tags (JA)}: 即時文法, 口語, 結論表現, 要約提示表現, 確認表現\\
\textbf{Tags (EN)}: immediate grammar, spoken, expression of conclusion, expression of summary presentation, expression of confirmation

\textbf{Adjusted Expression (JA)}:\\
(結局:けっきょく)、...ではないかと思います。\\
\textbf{Adjusted Expression (EN)}:\\
In the end, I think it is ...

\textbf{Example (JA)}:\\
A「この(問題:もんだい)、(解決:かいけつ)できる？」 B「うーん、(結局:けっきょく)、(時間:じかん)が(掛:か)かるし、(費用:ひよう)も(高:たか)いから、(難:むずか)しいんじゃないかって」\\
\textbf{Example (EN)}:\\
A: Can we solve this problem? B: Hmm, in the end, I think it's difficult because it takes time and costs a lot.

\textbf{Notes (JA)}:\\
「結局、...じゃないかって。」は、話し手がある議論や状況の最終的な結論を要約し、確認する際に使う口語的な表現。「結局」は「in the end」や「ultimately」の意味で、最終的な結果を示します。「...じゃないかって。」は「isn't it...?」の意味で、相手に対して確認を求めています。

\textbf{Notes (EN)}:\\
'Kekkyoku, ... ja nai ka tte.' is a colloquial expression used when the speaker wants to summarize and confirm the final conclusion of a discussion or situation. 'Kekkyoku' means 'in the end' or 'ultimately', indicating the final result. '... ja nai ka tte.' means 'isn't it...?', seeking confirmation from the listener.

\newpage

\section*{654. To put it very roughly,...}

\textbf{Date}: 2026.01.24 (Sat)\\
\textbf{Index}: sugokuzatsuni iukedo\\
\textbf{Expression (JA)}: すごく、(雑:ざつ)に(言:い)うけど、...\\
\textbf{Romanization}: sugoku, zatsu ni iu kedo, ...\\
\textbf{Expression (EN)}: To put it very roughly,...

\textbf{Variants (JA)}: 大まかに(言:い)うと、... / ざっくり(言:い)うと、... / 簡単に(言:い)うと、...\\
\textbf{Variants (EN)}: To put it very roughly,... / To say it broadly,... / To say it simply,...

\textbf{Intent (JA)}: 概要の提示, 情報の簡略化, 注意喚起\\
\textbf{Intent (EN)}: presentation of overview, simplification of information, attention call

\textbf{Tags (JA)}: 即時文法, 口語, 概要提示表現, 情報簡略化表現, 注意喚起表現\\
\textbf{Tags (EN)}: immediate grammar, spoken, expression of overview presentation, expression of information simplification, expression of attention call

\textbf{Adjusted Expression (JA)}:\\
(大:おお)まかに(言:い)うと、...\\
\textbf{Adjusted Expression (EN)}:\\
To put it broadly,...

\textbf{Example (JA)}:\\
A「この(計画:けいかく)、(詳:くわ)しく(説明:せつめい)して」 B「まずね、すごく(雑:ざつ)に(言:い)うけど、(要:よう)は(利益:りえき)を(最大化:さいだいか)することが(目的:もくてき)だよ」\\
\textbf{Example (EN)}:\\
A: Please explain this plan in detail. B: First, to put it very roughly, the main point is to maximize profits.

\textbf{Notes (JA)}:\\
「すごく、雑に言うけど、...」は、話し手がある複雑な事柄や情報を簡略化して伝える際に使う口語的な表現。「すごく」は「very」の意味で、強調のニュアンスを持ちます。「雑に言うけど、...」は「to put it roughly,...」の意味で、詳細を省略して概要を伝えることを示します。

\textbf{Notes (EN)}:\\
'To put it very roughly,...' is a colloquial expression used when the speaker wants to convey a complex matter or information in a simplified manner. 'Sugoku' means 'very', adding an emphasis nuance. 'Zatsu ni iu kedo,...' means 'to put it roughly,...', indicating that the speaker is omitting details and providing an overview. In this case, the speaker is giving a broad explanation of the plan's main objective without going into specifics.

\newpage

\section*{653. Strangely,...}

\textbf{Date}: 2026.01.23 (Fri)\\
\textbf{Index}: myouni\\
\textbf{Expression (JA)}: (妙:みょう)に、...\\
\textbf{Romanization}: myou ni, ...\\
\textbf{Expression (EN)}: Strangely,...

\textbf{Variants (JA)}: 不思議:ふしぎ)に、... / 変:へん)に、... / 奇妙:きみょう)に、...\\
\textbf{Variants (EN)}: Strangely,... / Wonderfully,... / Oddly,...

\textbf{Intent (JA)}: 異常性の表現, 感覚の強調, 状況の説明\\
\textbf{Intent (EN)}: expression of abnormality, emphasis of sensation, explanation of situation

\textbf{Tags (JA)}: 即時文法, 口語, 異常性表現, 感覚強調表現, 状況説明表現\\
\textbf{Tags (EN)}: immediate grammar, spoken, expression of abnormality, expression of sensation emphasis, expression of situation explanation

\textbf{Adjusted Expression (JA)}:\\
(妙:みょう)に(感じ:かんじ)られるのですが、...\\
\textbf{Adjusted Expression (EN)}:\\
It feels strange, but ...

\textbf{Example (JA)}:\\
A「(最近:さいきん)、(疲:つか)れやすいなあ」 B「うん、(妙:みょう)に(体:からだ)が(重:おも)いし、(眠:ねむ)れないこともあるよね」\\
\textbf{Example (EN)}:\\
A: I've been getting tired easily lately. B: Yeah, my body feels strangely heavy, and sometimes I can't sleep.

\textbf{Notes (JA)}:\\
「妙に、...」は、話し手がある感覚や状況が通常とは異なることを強調して表現する際に使う口語的な表現。「妙に」は「strangely」や「oddly」の意味で、異常性や違和感を示します。

\textbf{Notes (EN)}:\\
'Myou ni, ...' is a colloquial expression used when the speaker wants to emphasize that a certain sensation or situation is different from the usual. 'Myou ni' means 'strangely' or 'oddly', indicating abnormality or discomfort.

\newpage

\section*{652. Speaking of...}

\textbf{Date}: 2026.01.22 (Thu)\\
\textbf{Index}: toieba\\
\textbf{Expression (JA)}: ...といえば、...\\
\textbf{Romanization}: ... to ieba, ...\\
\textbf{Expression (EN)}: Speaking of...

\textbf{Variants (JA)}: ...と言えば、... / ...って言えば、... / ...といえば、...ね。\\
\textbf{Variants (EN)}: Speaking of... / If you say... / Speaking of..., right?

\textbf{Intent (JA)}: 話題の転換, 関連情報の提供, 思い出の喚起\\
\textbf{Intent (EN)}: topic shift, provision of related information, evocation of memories

\textbf{Tags (JA)}: 即時文法, 口語, 話題転換表現, 関連情報提供表現, 思い出喚起表現\\
\textbf{Tags (EN)}: immediate grammar, spoken, expression of topic shift, expression of related information provision, expression of memory evocation

\textbf{Adjusted Expression (JA)}:\\
(関連:かんれん)するかもしれませんが、...と(言:い)えば、...\\
\textbf{Adjusted Expression (EN)}:\\
It may be related, but speaking of ...

\textbf{Example (JA)}:\\
A「(昨日:きのう)、(映画:えいが)を(見:み)たよ」 B「へえ、(映画:えいが)といえば、(最近:さいきん)、あの(監督:かんとく)の(新作:しんさく)が(公開:こうかい)されたね」\\
\textbf{Example (EN)}:\\
A: I watched a movie yesterday. B: Oh, speaking of movies, that director's new work was released recently.

\textbf{Notes (JA)}:\\
...といえば、...」は、話し手がある話題から関連する別の話題に転換し、関連情報を提供したり、思い出を喚起したりする際に使う口語的な表現。「...といえば」は「speaking of ...」の意味で、特定の対象を指します。話し手は、相手に対して新しい話題を提示し、会話を広げています。

\textbf{Notes (EN)}:\\
'... to ieba, ...' is a colloquial expression used when the speaker wants to shift from one topic to a related one, provide related information, or evoke memories. '... to ieba' means 'speaking of ...', referring to a specific object. The speaker introduces a new topic to the listener and expands the conversation.

\newpage

\section*{651. ... is necessary.}

\textbf{Date}: 2026.01.21 (Wed)\\
\textbf{Index}: hitsuyaodesu\\
\textbf{Expression (JA)}: ...は(必要:ひつよう)。\\
\textbf{Romanization}: ... wa hitsuyou.\\
\textbf{Expression (EN)}: ... is necessary.

\textbf{Variants (JA)}: ...が(必要:ひつよう)。 / ...は(必須:ひっす)。 / ...が(必須:ひっす)。\\
\textbf{Variants (EN)}: ... is necessary. / ... is essential. / ... is required.

\textbf{Intent (JA)}: 必要性の表現, 条件の提示, 重要性の強調\\
\textbf{Intent (EN)}: expression of necessity, presentation of conditions, emphasis of importance

\textbf{Tags (JA)}: 即時文法, 口語, 必要性表現, 条件提示表現, 重要性強調表現\\
\textbf{Tags (EN)}: immediate grammar, spoken, expression of necessity, expression of condition presentation, expression of importance emphasis

\textbf{Adjusted Expression (JA)}:\\
...は(必要:ひつよう)です。\\
\textbf{Adjusted Expression (EN)}:\\
... is necessary.

\textbf{Example (JA)}:\\
A「この(書類:しょるい)、(提出:ていしゅつ)しなきゃ？」 B「うん、これは(必要:ひつよう)」\\
\textbf{Example (EN)}:\\
A: Do I have to submit this document? B: Yes, this is necessary.

\textbf{Notes (JA)}:\\
...は必要。は、話し手がある事柄や物事が欠かせないものであることを伝える際に使う口語的な表現。もう少し詳しくこのことばの構成を説明すると、「...は」は「... is」の意味で、特定の対象を指します。「必要」は「necessary」や「essential」の意味です。話し手は、相手に対してその対象が重要であることを強調しています。会話ですぐに返答するには「...は必要」で終わる形の方が自然です。

\textbf{Notes (EN)}:\\
'... wa hitsuyou.' is a colloquial expression used when the speaker wants to convey that a certain matter or thing is indispensable. To explain the structure of this phrase in more detail, '... wa' means '... is', referring to a specific object. 'Hitsuyou' means 'necessary' or 'essential'. The speaker emphasizes to the listener that the object in question is important. In conversation, it is more natural to end with the form '... wa hitsuyou' for an immediate response.

\newpage

\section*{650. As for this, ...}

\textbf{Date}: 2026.01.20 (Tue)\\
\textbf{Index}: korebakariwa\\
\textbf{Expression (JA)}: こればかりは、...\\
\textbf{Romanization}: kore bakari wa, ...\\
\textbf{Expression (EN)}: As for this, ...

\textbf{Variants (JA)}: これだけは、... / どうすることも、... / これはちょっと、...\\
\textbf{Variants (EN)}: As for this, ... / No way to do anything about it, ... / This is a bit, ...

\textbf{Intent (JA)}: 限定の表現, 例外の強調, 特別な扱いの示唆, 不可能\\
\textbf{Intent (EN)}: expression of limitation, emphasis of exception, suggestion of special treatment, impossibility

\textbf{Tags (JA)}: 即時文法, 口語, 限定表現, 例外強調表現, 特別扱い示唆表現, 不可能表現\\
\textbf{Tags (EN)}: immediate grammar, spoken, expression of limitation, expression of exception emphasis, expression of special treatment suggestion, expression of impossibility

\textbf{Adjusted Expression (JA)}:\\
これに関しては、いかなることもいたしかねます。\\
\textbf{Adjusted Expression (EN)}:\\
As for this, there is nothing I can do.

\textbf{Example (JA)}:\\
A「この(問題:もんだい)、(解決:かいけつ)できる？」 B「ごめんね、こればかりは、どうすることも...」\\
\textbf{Example (EN)}:\\
A: Can you solve this problem? B: Sorry, as for this, there is nothing I can do.

\textbf{Notes (JA)}:\\
「こればかりは、...」は、話し手が特定の事柄に対して例外的な扱いを示し、その事柄に関しては自分の力ではどうすることもできないことを伝える際に使う口語的な表現。「こればかりは」は「as for this only」の意味で、特定の対象を強調します。それだけで後を言わなくても、困った表情をすれば、不可能であることを示しています。

\textbf{Notes (EN)}:\\
'Kore bakari wa, ...' is a colloquial expression used when the speaker wants to indicate special treatment for a specific matter and convey that there is nothing they can do about it. 'Kore bakari wa' means 'as for this only', emphasizing a specific object. Even without saying anything further, a troubled expression can indicate that it is impossible.

\newpage

\section*{649. What do you mean by that?}

\textbf{Date}: 2026.01.19 (Mon)\\
\textbf{Index}: soredouiu kotodesuka\\
\textbf{Expression (JA)}: それ、どういうことですか？\\
\textbf{Romanization}: sore, dou iu koto desu ka?\\
\textbf{Expression (EN)}: What do you mean by that?

\textbf{Variants (JA)}: それ、どういう意味ですか？ / それ、どういう風に？ / それ、どういうこと？\\
\textbf{Variants (EN)}: What do you mean by that? / What does that mean? / What do you mean?

\textbf{Intent (JA)}: 意味の確認, 説明の要求, 理解の促進\\
\textbf{Intent (EN)}: confirmation of meaning, request for explanation, promotion of understanding

\textbf{Tags (JA)}: 即時文法, 口語, 意味確認表現, 説明要求表現, 理解促進表現\\
\textbf{Tags (EN)}: immediate grammar, spoken, expression of meaning confirmation, expression of explanation request, expression of understanding promotion

\textbf{Adjusted Expression (JA)}:\\
それは、どういう(意味:いみ)でしょうか？\\
\textbf{Adjusted Expression (EN)}:\\
What does that mean?

\textbf{Example (JA)}:\\
A「(今日:きょう)は、(無料:むりょう)です」 B「それ、どういうことですか？」\\
\textbf{Example (EN)}:\\
A: Today is free of charge. B: What do you mean by that?

\textbf{Notes (JA)}:\\
「それ、どういうことですか？」は、話し手が相手の発言や行動の意味を確認し、説明を求める際に使う口語的な表現。「それ」は「that」の意味で、特定の対象を指します。「どういうことですか？」は「what do you mean?」の意味で、相手の意図や意味を尋ねています。話し手は、相手の発言や行動の背後にある意図や意味を理解しようとしています。

\textbf{Notes (EN)}:\\
'Sore, dou iu koto desu ka?' is a colloquial expression used when the speaker wants to confirm the meaning of the listener's statement or action and request an explanation. 'Sore' means 'that', referring to a specific object. 'Dou iu koto desu ka?' means 'what do you mean?', asking about the listener's intent or meaning. The speaker is trying to understand the intent or meaning behind the listener's statement or action.

\newpage

\section*{648. What on earth}

\textbf{Date}: 2026.01.18 (Sun)\\
\textbf{Index}: ittaizentai\\
\textbf{Expression (JA)}: いったい(全体:ぜんたい)\\
\textbf{Romanization}: ittai zentai\\
\textbf{Expression (EN)}: What on earth

\textbf{Variants (JA)}: いったいいつ？ / いったい誰が？ / いったいなんで？\\
\textbf{Variants (EN)}: When on earth / Who on earth / Why on earth

\textbf{Intent (JA)}: 強調の表現, 疑問の強調, 驚きの表現\\
\textbf{Intent (EN)}: expression of emphasis, emphasis of question, expression of surprise

\textbf{Tags (JA)}: 即時文法, 口語, 強調表現, 疑問強調表現, 驚き表現\\
\textbf{Tags (EN)}: immediate grammar, spoken, expression of emphasis, expression of question emphasis, expression of surprise

\textbf{Adjusted Expression (JA)}:\\
どういうわけで、このようなことになったのでしょうか。\\
\textbf{Adjusted Expression (EN)}:\\
How on earth did this happen?

\textbf{Example (JA)}:\\
A「(彼:かれ)、(行:い)っちゃったよ」 B「いったい(全体:ぜんたい)、(何:なに)が(起:お)こったの？」\\
\textbf{Example (EN)}:\\
A: He left. B: What on earth happened?

\textbf{Notes (JA)}:\\
「いったい全体」は、話し手が驚きや困惑を表現し、ある出来事や状況に対して強い疑問を抱いていることを示す際に使う口語的な表現。「いったい」は「on earth」や「in the world」の意味で、強調のニュアンスを持ちます。「全体」は「whole」や「entirety」の意味ですが、この場合は強調の一部として機能します。

\textbf{Notes (EN)}:\\
'Ittai zentai' is a colloquial expression used when the speaker wants to express surprise or confusion and indicate that they have strong questions about a certain event or situation. 'Ittai' means 'on earth' or 'in the world', carrying a nuance of emphasis. 'Zentai' means 'whole' or 'entirety', but in this case, it functions as part of the emphasis.

\newpage

\section*{647. It might take a little longer.}

\textbf{Date}: 2026.01.17 (Sat)\\
\textbf{Index}: mouchottokakaru\\
\textbf{Expression (JA)}: もうちょっと(掛:か)かる。\\
\textbf{Romanization}: mou chotto kakaru.\\
\textbf{Expression (EN)}: It might take a little longer.

\textbf{Variants (JA)}: もう少し(掛:か)かる。 / あと少し(掛:か)かる。 / もうちょっと(時間:じかん)が(必要:ひつよう)だ。\\
\textbf{Variants (EN)}: It might take a little longer. / It will take a little more time. / It needs a bit more time.

\textbf{Intent (JA)}: 時間の見積もり, 進捗の報告, 期待値の調整\\
\textbf{Intent (EN)}: estimation of time, reporting of progress, adjustment of expectations

\textbf{Tags (JA)}: 即時文法, 口語, 時間見積もり表現, 進捗報告表現, 期待値調整表現\\
\textbf{Tags (EN)}: immediate grammar, spoken, expression of time estimation, expression of progress reporting, expression of expectation adjustment

\textbf{Adjusted Expression (JA)}:\\
もう(少々:しょうしょう)、(時間:じかん)が(掛:か)かります。\\
\textbf{Adjusted Expression (EN)}:\\
It will take a little more time.

\textbf{Example (JA)}:\\
A「(仕事:しごと)、(終:お)わった？」 B「もうちょっと(掛:か)かる」\\
\textbf{Example (EN)}:\\
A: Is the work finished? B: It might take a little longer.

\textbf{Notes (JA)}:\\
「もうちょっと掛かる」は、話し手がある作業やプロセスに対して、完了までにもう少し時間が必要であることを伝える際に使う口語的な表現。「もうちょっと」は「a little more」や「a bit more」の意味です。「掛かる」は「to take time」の意味です。話し手は、相手に対して進捗状況を報告し、期待値を調整しています。

\textbf{Notes (EN)}:\\
'Mou chotto kakaru.' is a colloquial expression used when the speaker wants to convey that a certain task or process requires a little more time to complete. 'Mou chotto' means 'a little more' or 'a bit more'. 'Kakaru' means 'to take time'. The speaker is reporting the progress to the listener and adjusting their expectations.

\newpage

\section*{646. Because I'm an amateur}

\textbf{Date}: 2026.01.16 (Fri)\\
\textbf{Index}: doshiroutonande\\
\textbf{Expression (JA)}: ど(素人:しろうと)なんで、...。\\
\textbf{Romanization}: doshirouto nande, ...\\
\textbf{Expression (EN)}: Because I'm an amateur, ...

\textbf{Variants (JA)}: ど(素人:しろうと)なんですけど、... / ど(素人:しろうと)だから、... / ど(素人:しろうと)なので、...\\
\textbf{Variants (EN)}: Because I'm an amateur, ... / Since I'm an amateur, ... / As I'm an amateur, ...

\textbf{Intent (JA)}: 自己紹介の表現, 能力の制限の伝達, 謙遜の意図\\
\textbf{Intent (EN)}: expression of self-introduction, communication of limitation of ability, intention of humility

\textbf{Tags (JA)}: 即時文法, 口語, 自己紹介表現, 能力制限伝達, 謙遜表現\\
\textbf{Tags (EN)}: immediate grammar, spoken, expression of self-introduction, communication of ability limitation, expression of humility

\textbf{Adjusted Expression (JA)}:\\
まったくの(素人:しろうと)なので、(期待:きたい)しないでください。\\
\textbf{Adjusted Expression (EN)}:\\
Please don't expect much, as I'm a complete amateur.

\textbf{Example (JA)}:\\
A「この(仕事:しごと)、(手伝:てつだ)ってくれる？」 B「ごめんね、ど(素人:しろうと)なんで、あまり(役:やく)に(立:た)てないかも」\\
\textbf{Example (EN)}:\\
A: Can you help me with this work? B: Sorry, because I'm an amateur, I might not be of much help.

\textbf{Notes (JA)}:\\
「ど素人なんで、...」は、話し手が自分の能力や経験に対して謙遜の意を表し、相手に対して過度な期待を持たせないようにする際に使う口語的な表現。「ど素人」は「complete amateur」で、全くの初心者であることを示します。「なんで、...」は「because...」の意味で、理由を述べています。話し手は、相手に対して自分の限界を認識させ、謙虚な態度を示しています。

\textbf{Notes (EN)}:\\
'Doshirouto nande, ...' is a colloquial expression used when the speaker wants to express humility about their abilities or experience and to avoid raising excessive expectations in the listener. 'Doshirouto' means 'complete amateur', indicating that the speaker is a complete beginner. 'Nande, ...' means 'because...', providing a reason. The speaker acknowledges their limitations to the listener and demonstrates a humble attitude.

\newpage

\section*{645. That's a different thing.}

\textbf{Date}: 2026.01.15 (Thu)\\
\textbf{Index}: sorebetsumono\\
\textbf{Expression (JA)}: それ、(別物:べつもの)。\\
\textbf{Romanization}: sore, betsu-mono.\\
\textbf{Expression (EN)}: That's a different thing.

\textbf{Variants (JA)}: それ、(別:べつ)の(物:もの)。 / それ、(違:ちが)う(物:もの)。 / それ、(別物:べつもの)です。\\
\textbf{Variants (EN)}: That's a different thing. / That's another thing. / That's a different item.

\textbf{Intent (JA)}: 区別の表現, 誤解の訂正, 情報の明確化\\
\textbf{Intent (EN)}: expression of distinction, correction of misunderstanding, clarification of information

\textbf{Tags (JA)}: 即時文法, 口語, 区別表現, 誤解訂正表現\\
\textbf{Tags (EN)}: immediate grammar, spoken, expression of distinction, expression of misunderstanding correction

\textbf{Adjusted Expression (JA)}:\\
それとは(別:べつ)の(物:もの)です。\\
\textbf{Adjusted Expression (EN)}:\\
That is a different thing.

\textbf{Example (JA)}:\\
A「これが、(以前:いぜん)、(話:はな)してたものなの？」 B「ううん、それ、(別物:べつもの)」\\
\textbf{Example (EN)}:\\
A: Is this what you were talking about before? B: No, that's a different thing.

\textbf{Notes (JA)}:\\
「それ、別物」は、話し手がある事柄や物体が他のものと異なることを強調し、誤解を訂正する際に使う口語的な表現。もう少し詳しくこのことばの構成を説明すると、「それ」は「that」の意味で、特定の対象を指します。「別物」は「different thing」や「another item」の意味です。話し手は、相手に対してその対象が他のものと区別されるべきであることを伝えています。

\textbf{Notes (EN)}:\\
'Sore, betsu-mono.' is a colloquial expression used when the speaker wants to emphasize that a certain matter or object is different from others and to correct a misunderstanding. To explain the structure of this phrase in more detail, 'sore' means 'that', referring to a specific object. 'Betsu-mono' means 'different thing' or 'another item'. The speaker is conveying to the listener that the object in question should be distinguished from others.

\newpage

\section*{644. It doesn't budge at all.}

\textbf{Date}: 2026.01.14 (Wed)\\
\textbf{Index}: bikutomoshinai\\
\textbf{Expression (JA)}: (ビク:びく)ともしない。\\
\textbf{Romanization}: biku to mo shinai.\\
\textbf{Expression (EN)}: It doesn't budge at all.

\textbf{Variants (JA)}: 全く(動:うご)かない。 / まったく(動:うご)かない。 / 全然(動:うご)かない。\\
\textbf{Variants (EN)}: It doesn't budge at all. / It doesn't move at all. / It won't move at all.

\textbf{Intent (JA)}: 不動の表現, 強調の意図, 状況の説明\\
\textbf{Intent (EN)}: expression of immobility, intention of emphasis, explanation of situation

\textbf{Tags (JA)}: 即時文法, 口語, 不動表現, 強調表現\\
\textbf{Tags (EN)}: immediate grammar, spoken, expression of immobility, expression of emphasis

\textbf{Adjusted Expression (JA)}:\\
まったく(動:うご)きません。\\
\textbf{Adjusted Expression (EN)}:\\
It doesn't move at all.

\textbf{Example (JA)}:\\
A「この(ドア:どあ)、(開:あ)けてみて」 B「うーん、(ビク:びく)ともしない」\\
\textbf{Example (EN)}:\\
A: Try opening this door. B: Hmm, it doesn't budge at all.

\textbf{Notes (JA)}:\\
「ビクともしない。」は、話し手がある物体や状況が全く動かないことを強調して表現する際に使う口語的な表現。もう少し詳しくこのことばの構成を説明すると、「ビクともしない」は「not budging at all」の意味で、完全な不動状態を示します。話し手は、物体や状況の動かない性質を強調しています。

\textbf{Notes (EN)}:\\
'Biku to mo shinai.' is a colloquial expression used when the speaker wants to emphasize that a certain object or situation does not move at all. To explain the structure of this phrase in more detail, 'biku to mo shinai' means 'not budging at all', indicating a state of complete immobility. The speaker emphasizes the unmovable nature of the object or situation.

\newpage

\section*{643. let me do this much}

\textbf{Date}: 2026.01.13 (Tue)\\
\textbf{Index}: koreguraiwasasete\\
\textbf{Expression (JA)}: これぐらいはさせて\\
\textbf{Romanization}: kore gurai wa sasete,\\
\textbf{Expression (EN)}: let me do this much

\textbf{Variants (JA)}: これくらいはさせて / これだけはさせて / これぐらいはさせてください。\\
\textbf{Variants (EN)}: let me do this much / please let me do this much / let me at least do this much, please.

\textbf{Intent (JA)}: 謙遜の表現, 気持ちの伝達, 依頼の意図\\
\textbf{Intent (EN)}: humble expression, communication of feelings, intention of request

\textbf{Tags (JA)}: 即時文法, 口語, 謙遜表現, 依頼表現\\
\textbf{Tags (EN)}: immediate grammar, spoken, humble expression, request

\textbf{Adjusted Expression (JA)}:\\
これぐらいは(私:わたし)に(任:まか)せてください。\\
\textbf{Adjusted Expression (EN)}:\\
Please let me handle at least this much.

\textbf{Example (JA)}:\\
A「これ、(手伝:てつだ)ってくれてありがとう」 B「これぐらいはさせて」\\
\textbf{Example (EN)}:\\
A: Thank you for helping me with this. B: Let me do this much.

\textbf{Notes (JA)}:\\
「これぐらいはさせて」は、いつもお世話になっている人からの頼まれごとに対して、謙遜しつつも自分の力でできる範囲で手伝いたいという気持ちの表現。「これぐらいは」は「at least this much」の意味。「させて」は「let me do it」の意味。

\textbf{Notes (EN)}:\\
'Kore gurai wa sasete' is an expression used when responding to a request from someone who is always taking care of you, expressing a humble yet willing attitude to help within your own capabilities. 'Kore gurai wa' means 'at least this much', and 'sasete' means 'let me do it'.

\newpage

\section*{642. sorry that this is all I can do}

\textbf{Date}: 2026.01.12 (Mon)\\
\textbf{Index}: koreguraishikadekinakute\\
\textbf{Expression (JA)}: これぐらいしかできなくて、\\
\textbf{Romanization}: kore gurai shika dekinakute,\\
\textbf{Expression (EN)}: sorry that this is all I can do

\textbf{Variants (JA)}: これくらいしかできなくて、 / これだけしかできなくて、 / これぐらいしかできなくて、すみません。\\
\textbf{Variants (EN)}: sorry that this is all I can do / sorry that I can only do this much / sorry that this is all I can do, sorry.

\textbf{Intent (JA)}: 謙遜の表現, 能力の制限の伝達, 謝罪の意図\\
\textbf{Intent (EN)}: expression of humility, communication of limitation of ability, intention of apology

\textbf{Tags (JA)}: 即時文法, 口語, 謙遜表現, 能力制限伝達, 謝罪表現\\
\textbf{Tags (EN)}: immediate grammar, spoken, expression of humility, communication of ability limitation, expression of apology

\textbf{Adjusted Expression (JA)}:\\
これぐらいしかできなくて、(申:もう)し(訳:わけ)ありません。\\
\textbf{Adjusted Expression (EN)}:\\
I'm sorry that this is all I can do.

\textbf{Example (JA)}:\\
A「これ、(手伝:てつだ)ってくれてありがとう」 B「ごめんね、これぐらいしかできなくて」\\
\textbf{Example (EN)}:\\
A: Thank you for helping me with this. B: Sorry that this is all I can do.

\textbf{Notes (JA)}:\\
「これぐらいしかできなくて、」は、話し手が自分の能力や行動に対して謙遜の意を表し、謝罪の意図を伝える際に使う口語的な表現。もう少し詳しくこのことばの構成を説明すると、「これぐらいしか」は「only this much」の意味で、自分の限界を示します。「できなくて、」は「I can't do it」の意味で、能力の制限を伝えています。話し手は、相手に対して自分の限界を認識させ、謝罪の気持ちを表現しています。

\textbf{Notes (EN)}:\\
'Kore gurai shika dekinakute,' is a colloquial expression used when the speaker wants to express humility about their abilities or actions and convey an intention of apology. To explain the structure of this phrase in more detail, 'kore gurai shika' means 'only this much', indicating the speaker's limitations. 'Dekinakute,' means 'I can't do it', communicating the limitation of ability. The speaker acknowledges their limitations to the listener and expresses a feeling of apology.

\newpage

\section*{641. I'm glad I did it.}

\textbf{Date}: 2026.01.11 (Sun)\\
\textbf{Index}: shitoiteyokatta\\
\textbf{Expression (JA)}: しといてよかった。\\
\textbf{Romanization}: shitoite yokatta.\\
\textbf{Expression (EN)}: I'm glad I did it.

\textbf{Variants (JA)}: やっておいてよかった。 / しておいてよかった。 / やっといてよかった。\\
\textbf{Variants (EN)}: I'm glad I did it. / I'm glad I took care of it. / I'm glad I got it done.

\textbf{Intent (JA)}: 行動の肯定, 過去の選択の評価, 安心感の表現\\
\textbf{Intent (EN)}: affirmation of action, evaluation of past choice, expression of relief

\textbf{Tags (JA)}: 即時文法, 口語, 行動肯定表現, 過去選択評価表現, 安心感表現\\
\textbf{Tags (EN)}: immediate grammar, spoken, expression of action affirmation, expression of past choice evaluation, expression of relief

\textbf{Adjusted Expression (JA)}:\\
しておいて(正解:せいかい)でした。\\
\textbf{Adjusted Expression (EN)}:\\
It was the right what to do.

\textbf{Example (JA)}:\\
A「これ、(最後:さいご)の(一:ひと)つ？」 B「まだあるよ、たくさん(買:か)っといてよかったね」\\
\textbf{Example (EN)}:\\
A: Is this the last one? B: There are still more, I'm glad I bought a lot.

\textbf{Notes (JA)}:\\
「しといてよかった。」は、話し手が過去に行った行動や選択を肯定的に評価し、その結果に対して安心感を表現する際に使う口語的な表現。もう少し詳しくこのことばの構成を説明すると、「しといて」は「しておく」の口語形で、事前に行動を完了させることを意味します。「よかった」は「I'm glad」や「It was good」の意味です。話し手は、過去の選択が正しかったことを認識し、その結果に満足しています。

\textbf{Notes (EN)}:\\
'Shitoite yokatta.' is a colloquial expression used when the speaker wants to affirmatively evaluate an action or choice made in the past and express a sense of relief about the outcome. To explain the structure of this phrase in more detail, 'shitoite' is the colloquial form of 'shite oku', meaning to complete an action in advance. 'Yokatta' means 'I'm glad' or 'It was good'. The speaker recognizes that their past choice was correct and is satisfied with the result.

\newpage

\section*{640. Somehow, it's weird.}

\textbf{Date}: 2026.01.10 (Sat)\\
\textbf{Index}: aead\\
\textbf{Expression (JA)}: なんか、(変:へん)。\\
\textbf{Romanization}: nanka, hen.\\
\textbf{Expression (EN)}: Somehow, it's weird.

\textbf{Variants (JA)}: なんとなく、(変:へん)。 / なんだか、(変:へん)。 / なんとなく、おかしい。\\
\textbf{Variants (EN)}: Somehow, it's weird. / Somehow, it's strange. / Somehow, it's odd.

\textbf{Intent (JA)}: 違和感の表現, 感覚の共有, 評価の提示\\
\textbf{Intent (EN)}: expression of discomfort, sharing of feelings, presentation of evaluation

\textbf{Tags (JA)}: 即時文法, 口語, 違和感表現, 感覚共有表現, 評価提示表現\\
\textbf{Tags (EN)}: immediate grammar, spoken, expression of discomfort, expression of feeling sharing, expression of evaluation presentation

\textbf{Adjusted Expression (JA)}:\\
なんとなく、(変:へん)です。\\
\textbf{Adjusted Expression (EN)}:\\
Somehow, it's strange.

\textbf{Example (JA)}:\\
A「この(デザイン:でざいん)、どう？」 B「うーん、なんか、(変:へん)」\\
\textbf{Example (EN)}:\\
A: How about this design? B: Hmm, somehow, it's weird.

\textbf{Notes (JA)}:\\
「なんか、変。」は、話し手がある事柄に対して違和感を抱いていることを表現する際に使う口語的な表現。もう少し詳しくこのことばの構成を説明すると、「なんか」は「somehow」や「somewhat」の意味で、漠然とした感覚を示します。「変」は「weird」や「strange」の意味です。話し手は、違和感を共有し、その事柄が自分にとって適切でないことを伝えています。

\textbf{Notes (EN)}:\\
'Nanka, hen.' is a colloquial expression used when the speaker wants to express that they have a sense of discomfort about something. To explain the structure of this phrase in more detail, 'nanka' means 'somehow' or 'somewhat', indicating a vague feeling. 'Hen' means 'weird' or 'strange'. The speaker shares their discomfort and conveys that the matter is not suitable for them.

\newpage

\section*{639. How have you been?}

\textbf{Date}: 2026.01.09 (Fri)\\
\textbf{Index}: genkishiteru\\
\textbf{Expression (JA)}: (元気:げんき)してる？\\
\textbf{Romanization}: genki shiteru?\\
\textbf{Expression (EN)}: How have you been?

\textbf{Variants (JA)}: (元気:げんき)ですか？ / (調子:ちょうし)はどう？ / (最近:さいきん)、(元気:げんき)？\\
\textbf{Variants (EN)}: How have you been? / How are you? / How's it going?

\textbf{Intent (JA)}: 健康状態の確認, 近況の共有, 会話の開始\\
\textbf{Intent (EN)}: confirmation of health status, sharing of recent news, initiation of conversation

\textbf{Tags (JA)}: 即時文法, 口語, 健康確認表現, 近況共有表現, 会話開始表現\\
\textbf{Tags (EN)}: immediate grammar, spoken, expression of health confirmation, expression of recent news sharing, expression of conversation initiation

\textbf{Adjusted Expression (JA)}:\\
(お元気:おげんき)ですか？\\
\textbf{Adjusted Expression (EN)}:\\
How are you?

\textbf{Example (JA)}:\\
A「(久:ひさ)しぶり！(元気:げんき)してる？」 B「うん、(元気:げんき)だよ」\\
\textbf{Example (EN)}:\\
A: Long time no see! How have you been? B: Yeah, I'm fine.

\textbf{Notes (JA)}:\\
「元気してる？」は、話し手が相手の健康状態や近況を確認し、会話を始める際に使う口語的な表現。もう少し詳しくこのことばの構成を説明すると、「元気」は「health」や「well-being」の意味で、「してる？」は「are you doing?」の意味です。話し手は、相手の状態を尋ね、会話のきっかけを作っています。

\textbf{Notes (EN)}:\\
'Genki shiteru?' is a colloquial expression used when the speaker wants to check the listener's health status and recent news, and to initiate a conversation. To explain the structure of this phrase in more detail, 'genki' means 'health' or 'well-being', and 'shiteru?' means 'are you doing?'. The speaker is inquiring about the listener's condition and creating an opportunity for conversation.

\newpage

\section*{638. It doesn't click.}

\textbf{Date}: 2026.01.08 (Thu)\\
\textbf{Index}: pintokonai\\
\textbf{Expression (JA)}: (ピン:ぴん)と(来:こ)ない。\\
\textbf{Romanization}: pin to konai.\\
\textbf{Expression (EN)}: It doesn't click.

\textbf{Variants (JA)}: (ピン:ぴん)と(来:こ)ません。 / しっくり(来:こ)ない。 / なんとなく(違:ちが)う(感:かん)じがする。\\
\textbf{Variants (EN)}: It doesn't click. / It doesn't quite make sense. / I don't really get it.

\textbf{Intent (JA)}: 合わない感覚の表現, 違和感の共有, 直感的に違うと感じることの伝達\\
\textbf{Intent (EN)}: sharing a sense of mismatch, sharing intuitive unease, expressing lack of intuitive alignment

\textbf{Tags (JA)}: 即時文法, 口語, 感覚表現, 違和感表現, 直感表現\\
\textbf{Tags (EN)}: immediate grammar, spoken, expression of feeling, expression of discomfort, expression of intuition

\textbf{Adjusted Expression (JA)}:\\
しっくり(来:こ)ない。\\
\textbf{Adjusted Expression (EN)}:\\
It doesn't fit well.

\textbf{Example (JA)}:\\
A「この(デザイン:でざいん)、どう(思:おも)う？」 B「うーん、(ピン:ぴん)と(来:こ)ない」\\
\textbf{Example (EN)}:\\
A: What do you think of this design? B: Hmm, it doesn't click.

\textbf{Notes (JA)}:\\
「ピンと来ない」は、話し手がある事柄に対して直感的に合わない感覚を抱いていることを表現する際に使う口語的な表現。話し手は、違和感を共有し、その事柄が自分にとって適切でないことを伝えています。

\textbf{Notes (EN)}:\\
'Pin to konai.' is a colloquial expression used when the speaker wants to express that they have an intuitive sense of mismatch about something. The speaker shares their discomfort and conveys that the matter is not suitable for them.

\newpage

\section*{637. I made fish.}

\textbf{Date}: 2026.01.07 (Wed)\\
\textbf{Index}: sakana tsukuttayo\\
\textbf{Expression (JA)}: (魚:さかな)、(作:つく)ったよ。\\
\textbf{Romanization}: sakana, tsukutta yo.\\
\textbf{Expression (EN)}: I made fish.

\textbf{Variants (JA)}: (魚:さかな)、(作:つく)りましたよ。 / (魚:さかな)、(調理:ちょうり)したよ。 / (魚:さかな)、(料理:りょうり)したよ。\\
\textbf{Variants (EN)}: I made fish. / I cooked fish. / I prepared fish.

\textbf{Intent (JA)}: 料理の報告, 成果の共有, 食事の提案\\
\textbf{Intent (EN)}: reporting of cooking, sharing of achievement, suggestion of meal

\textbf{Tags (JA)}: 即時文法, 口語, 料理報告表現, 成果共有表現, 食事提案表現\\
\textbf{Tags (EN)}: immediate grammar, spoken, expression of cooking report, expression of achievement sharing, expression of meal suggestion

\textbf{Adjusted Expression (JA)}:\\
(魚:さかな)を(調理:ちょうり)しました。\\
\textbf{Adjusted Expression (EN)}:\\
I cooked fish.

\textbf{Example (JA)}:\\
A「(夕食:ゆうしょく)、何にする？」 B「(魚:さかな)、(作:つく)ったよ」\\
\textbf{Example (EN)}:\\
A: What should we have for dinner? B: I made fish.

\textbf{Notes (JA)}:\\
「魚、作ったよ。」は、話し手が料理した内容を報告し、食事の提案をする際に使う口語的な表現。もう少し詳しくこのことばの構成を説明すると、「魚」は「fish」の意味で、「作ったよ」は「I made」や「I cooked」の意味です。話し手は、相手に料理の内容を伝え、食事の選択肢として提案しています。

\textbf{Notes (EN)}:\\
'Sakana, tsukutta yo.' is a colloquial expression used when the speaker wants to report what they have cooked and suggest it as a meal. To explain the structure of this phrase in more detail, 'sakana' means 'fish', and 'tsukutta yo' means 'I made' or 'I cooked'. The speaker is conveying the content of the cooking to the listener and proposing it as a meal option.

\newpage

\section*{636. What I just saw was,}

\textbf{Date}: 2026.01.06 (Tue)\\
\textbf{Index}: imamitanowa\\
\textbf{Expression (JA)}: (今:いま)、(見:み)たのは、\\
\textbf{Romanization}: ima, mita no wa,\\
\textbf{Expression (EN)}: What I just saw was,

\textbf{Variants (JA)}: さっき、(見:み)たのは、 / たった(今:いま)、(見:み)たのは、 / (最近:さいきん)、(見:み)たのは、\\
\textbf{Variants (EN)}: What I just saw was, / What I recently saw was, / What I saw a moment ago was,

\textbf{Intent (JA)}: 情報の共有, 経験の報告, 注意喚起, 話題の導入, 名詞句の強調\\
\textbf{Intent (EN)}: sharing of information, reporting of experience, drawing attention, introduction of topic, emphasis on noun phrase

\textbf{Tags (JA)}: 即時文法, 口語, 情報共有表現, 経験報告表現, 注意喚起表現\\
\textbf{Tags (EN)}: immediate grammar, spoken, expression of information sharing, expression of experience reporting, expression of drawing attention

\textbf{Adjusted Expression (JA)}:\\
(先程:さきほど)、(見:み)たのは、\\
\textbf{Adjusted Expression (EN)}:\\
What I just saw was,

\textbf{Example (JA)}:\\
A「(今:いま)、(見:み)たのは、何だったの」 B「あれは(流:なが)れ(星:ぼし)だよ」\\
\textbf{Example (EN)}:\\
A: What was that I just saw? B: That was a shooting star.

\textbf{Notes (JA)}:\\
「今、見たのは、」は、話し手が最近経験した出来事や情報を共有し、注意を引くために使う口語的な表現。もう少し詳しくこのことばの構成を説明すると、「今」は「just now」や「recently」の意味で、時間的な近さを示します。「見たのは、」は「what I saw was」の意味で、経験したことを報告しています。話し手は、相手の注意を引きつけ、その後に続く情報に関心を持たせることを意図しています。

\textbf{Notes (EN)}:\\
'Ima, mita no wa,' is a colloquial expression used when the speaker wants to share a recent experience or information and draw attention. To explain the structure of this phrase in more detail, 'ima' means 'just now' or 'recently', indicating temporal proximity. 'Mita no wa,' means 'what I saw was', reporting on an experienced event. The speaker intends to capture the listener's attention and generate interest in the information that follows.

\newpage

\section*{635. Still mulling it over}

\textbf{Date}: 2026.01.05 (Mon)\\
\textbf{Index}: yayanayanderu\\
\textbf{Expression (JA)}: やや、(悩:なや)んでる、\\
\textbf{Romanization}: yaya, nayanderu,\\
\textbf{Expression (EN)}: Still mulling it over,

\textbf{Variants (JA)}: ちょっと、(悩:なや)んでる、 / 少し、(悩:なや)んでる、 / まあ、(悩:なや)んでる、\\
\textbf{Variants (EN)}: Still thinking about it, / Still mulling it over, / Still pondering it, / I'm going to sleep on it.

\textbf{Intent (JA)}: 思考過程の表現, 決定の保留, 内省の共有\\
\textbf{Intent (EN)}: expression of thought process, postponement of decision, sharing of introspection

\textbf{Tags (JA)}: 即時文法, 口語, 思考表現, 決定保留\\
\textbf{Tags (EN)}: immediate grammar, spoken, expression of thought, postponement of decision

\textbf{Adjusted Expression (JA)}:\\
まだ検討しているところです。\\
\textbf{Adjusted Expression (EN)}:\\
I'm still looking into it.

\textbf{Example (JA)}:\\
A「この(提案:ていあん)、どうする？」 B「やや、(悩:なや)んでる、」\\
\textbf{Example (EN)}:\\
A: What should we do about this proposal? B: Still mulling it over, ...

\textbf{Notes (JA)}:\\
「やや、悩んでる、」は、話し手が感情的に苦しんでいる状態を表すものではなく、ある事柄について判断を下す前の処理過程にあることを、やや距離を置いて共有する表現である。「やや」は思考の程度を示すだけでなく、深刻な悩みではないことを相手に示す調整の役割を持つ。この表現は、話し手が自分の内部処理を対象化し、聞き手に過度な心配を与えずに決定を保留していることを伝える点に特徴がある。

\textbf{Notes (EN)}:\\
'Yaya, nayanderu,' is not used to express emotional distress or worry. Rather, it is a colloquial expression that reports the speaker's current processing state before making a decision, with a certain distance. The adverb 'yaya' functions not only to indicate degree, but also to signal that the situation is not serious. By using this expression, the speaker objectifies their internal thought process and shares it without inviting unnecessary concern from the listener.

\newpage

\section*{634. Do you have time?}

\textbf{Date}: 2026.01.04 (Sun)\\
\textbf{Index}: jikanaru\\
\textbf{Expression (JA)}: (時間:じかん)、ある？\\
\textbf{Romanization}: jikan aru?\\
\textbf{Expression (EN)}: Do you have time?

\textbf{Variants (JA)}: (時間:じかん)、ありますか？ / (暇:ひま)、ありますか？ / (時間:じかん)、(大丈夫:だいじょうぶ)ですか？\\
\textbf{Variants (EN)}: Do you have time? / Are you free? / Is your time okay?

\textbf{Intent (JA)}: 時間の確認, 予定の調整, 会話の開始\\
\textbf{Intent (EN)}: confirmation of time, adjustment of schedule, initiation of conversation

\textbf{Tags (JA)}: 即時文法, 口語, 時間確認, 予定調整\\
\textbf{Tags (EN)}: immediate grammar, spoken, time confirmation, schedule adjustment

\textbf{Adjusted Expression (JA)}:\\
(今:いま)、お(時間:じかん)、ございますか？\\
\textbf{Adjusted Expression (EN)}:\\
Do you have time now?

\textbf{Example (JA)}:\\
A「(今:いま)、(時間:じかん)、ある？」 B「うん、いいよ」\\
\textbf{Example (EN)}:\\
A: Do you have time now? B: Yeah, it's fine.

\textbf{Notes (JA)}:\\
「時間ある？」は、話し手が相手の時間的な余裕を確認し、予定の調整や会話の開始を促す際に使う口語的な表現。もう少し詳しくこのことばの構成を説明すると、「時間」は「time」の意味で、「ある？」は「do you have?」の意味です。話し手は、相手の都合を尋ね、会話や活動の開始を意図しています。

\textbf{Notes (EN)}:\\
'Jikan aru?' is a colloquial expression used when the speaker wants to confirm the listener's availability of time and to prompt the adjustment of schedule or initiation of conversation. To explain the structure of this phrase in more detail, 'jikan' means 'time', and 'aru?' means 'do you have?'. The speaker is inquiring about the listener's convenience and intends to start a conversation or activity.

\newpage

\section*{633. So, do you have something else?}

\textbf{Date}: 2026.01.03 (Sat)\\
\textbf{Index}: janankaarimasu\\
\textbf{Expression (JA)}: じゃあ、なんかあります？\\
\textbf{Romanization}: jaa, nanka arimasu?\\
\textbf{Expression (EN)}: So, do you have something else?

\textbf{Variants (JA)}: じゃあ、他にありますか？ / じゃあ、何かありますか？ / じゃあ、ほかに何かありますか？\\
\textbf{Variants (EN)}: So, do you have something else? / Well then, is there anything else? / So, do you have any other options?

\textbf{Intent (JA)}: 追加の提案の促し, 選択肢の確認, 会話の継続\\
\textbf{Intent (EN)}: prompting additional suggestions, confirming options, continuation of conversation

\textbf{Tags (JA)}: 即時文法, 口語, 追加提案, 選択肢確認\\
\textbf{Tags (EN)}: immediate grammar, spoken, additional suggestion, option confirmation

\textbf{Adjusted Expression (JA)}:\\
では、他に何かありますか？\\
\textbf{Adjusted Expression (EN)}:\\
Well then, is there anything else?

\textbf{Example (JA)}:\\
A「これはあまりよくないよ」 B「じゃあ、なんかあります？」\\
\textbf{Example (EN)}:\\
A: This one isn't very good. B: Hmm, so, do you have something else?

\textbf{Notes (JA)}:\\
「じゃあ、なんかあります？」は、話し手が相手に対して追加の提案や選択肢を求める際に使う口語的な表現。もう少し詳しくこのことばの構成を説明すると、「じゃあ」は「well then」や「so」の意味で、会話の流れを示します。「なんかあります？」は「do you have something else?」や「is there anything else?」の意味です。話し手は、会話を継続し、相手からの追加の提案を促しています。

\textbf{Notes (EN)}:\\
'Jaa, nanka arimasu?' is a colloquial expression used when the speaker wants to ask the listener for additional suggestions or options. To explain the structure of this phrase in more detail, 'jaa' means 'well then' or 'so', indicating the flow of conversation. 'Nanka arimasu?' means 'do you have something else?' or 'is there anything else?'. The speaker is continuing the conversation and prompting the listener for additional suggestions.

\newpage

\section*{632. You know, there's no one knowing the answer.}

\textbf{Date}: 2026.01.02 (Fri)\\
\textbf{Index}: wakaruhitonantedaremoinai\\
\textbf{Expression (JA)}: わかる(人:ひと)なんて(誰:だれ)もいない。\\
\textbf{Romanization}: wakaruhito nante dare mo inai.\\
\textbf{Expression (EN)}: You know, there's no one knowing the answer.

\textbf{Variants (JA)}: わかる(人:ひと)なんて(誰:だれ)もいません。 / (答:こた)えが(分:わ)かる(人:ひと)なんて(誰:だれ)もいない。 / (理解:りかい)できる(人:ひと)なんて(誰:だれ)もいない。\\
\textbf{Variants (EN)}: You know, there's no one knowing the answer. / There's no one who understands the answer. / No one can comprehend it.

\textbf{Intent (JA)}: 知識の欠如の表現, 共通認識の強調, 不確実性の示唆\\
\textbf{Intent (EN)}: expression of lack of knowledge, emphasis on common understanding, suggestion of uncertainty

\textbf{Tags (JA)}: 即時文法, 口語, 知識欠如表現, 共通認識強調\\
\textbf{Tags (EN)}: immediate grammar, spoken, expression of lack of knowledge, emphasis on common understanding

\textbf{Adjusted Expression (JA)}:\\
(答:こた)えが(分:わ)かる(人:ひと)なんて(誰:だれ)もいません。\\
\textbf{Adjusted Expression (EN)}:\\
There's no one who knows the answer.

\textbf{Example (JA)}:\\
A「この(問題:もんだい)、できる？」 B「これは、(難:むずか)しいよ、わかる(人:ひと)なんて(誰:だれ)もいないよ」\\
\textbf{Example (EN)}:\\
A: Can you solve this quiz? B: This is difficult, you know, there's no one knowing the answer.

\textbf{Notes (JA)}:\\
「わかる人なんて誰もいない。」は、話し手がある事柄に対して知識や理解が欠如していることを表現し、共通認識を強調する際に使う口語的な表現。もう少し詳しくこのことばの構成を説明すると、「わかる人なんて」は「no one who knows」の意味で、「誰もいない」は「there is no one」の意味です。

\textbf{Notes (EN)}:\\
'Wakaruhito nante dare mo inai.' is a colloquial expression used when the speaker wants to express a lack of knowledge or understanding about a certain matter and to emphasize common understanding. To explain the structure of this phrase in more detail, 'wakaruhito nante' means 'no one who knows', and 'dare mo inai' means 'there is no one'.

\newpage

\section*{631. I wonder about that?}

\textbf{Date}: 2026.01.01 (Thu)\\
\textbf{Index}: doudaro\\
\textbf{Expression (JA)}: それ、どうだろ？\\
\textbf{Romanization}: sore, dou daro?\\
\textbf{Expression (EN)}: I wonder about that?

\textbf{Variants (JA)}: それ、どうかな？ / それ、どう思う？ / それ、どうなんだろう？\\
\textbf{Variants (EN)}: I wonder about that? / What do you think about that? / I wonder how that is?

\textbf{Intent (JA)}: 疑問の表現, 意見の求め, 不確実性の示唆\\
\textbf{Intent (EN)}: expression of doubt, seeking opinion, suggestion of uncertainty

\textbf{Tags (JA)}: 即時文法, 口語, 疑問表現, 意見求め\\
\textbf{Tags (EN)}: immediate grammar, spoken, expression of doubt, seeking opinion

\textbf{Adjusted Expression (JA)}:\\
それは、どうでしょうか？\\
\textbf{Adjusted Expression (EN)}:\\
I wonder about that?

\textbf{Example (JA)}:\\
A「この(計画:けいかく)、(成功:せいこう)するよ」 B「それ、どうだろ？」\\
\textbf{Example (EN)}:\\
A: This plan will succeed. B: I wonder about that?

\textbf{Notes (JA)}:\\
「それ、どうだろ？」は、話し手がある事柄に対して疑問を抱き、相手の意見を求める際に使う口語的な表現。もう少し詳しくこのことばの構成を説明すると、「それ」は特定の事柄を指し、「どうだろ？」は「I wonder about that?」や「What do you think about that?」の意味です。話し手は、不確実性を示唆し、相手の考えを聞きたいと考えています。

\textbf{Notes (EN)}:\\
'Sore, dou daro?' is a colloquial expression used when the speaker has doubts about a certain matter and wants to seek the opinion of the listener. To explain the structure of this phrase in more detail, 'sore' refers to a specific matter, and 'dou daro?' means 'I wonder about that?' or 'What do you think about that?'. The speaker suggests uncertainty and is interested in hearing the listener's thoughts.

\newpage

\section*{630. There's no other choice.}

\textbf{Date}: 2025.12.31 (Wed)\\
\textbf{Index}: soreshikanai\\
\textbf{Expression (JA)}: それしかない。\\
\textbf{Romanization}: sore shika nai.\\
\textbf{Expression (EN)}: There's no other choice.

\textbf{Variants (JA)}: 他に選択肢がない。 / これ以外に方法がない。 / これ以外の選択肢はない。\\
\textbf{Variants (EN)}: There's no other choice. / There are no other options. / There's no alternative.

\textbf{Intent (JA)}: 選択肢の欠如の表現, 決定の強調, 状況の受容\\
\textbf{Intent (EN)}: expression of lack of options, emphasis on decision, acceptance of situation

\textbf{Tags (JA)}: 即時文法, 口語, 選択肢欠如表現, 決定強調\\
\textbf{Tags (EN)}: immediate grammar, spoken, expression of lack of options, emphasis on decision

\textbf{Adjusted Expression (JA)}:\\
(他:ほか)に(選択肢:せんたくし)がありません。\\
\textbf{Adjusted Expression (EN)}:\\
There are no other options.

\textbf{Example (JA)}:\\
A「この(問題:もんだい)、どうする？やめる？」 B「うーん、それしかない」\\
\textbf{Example (EN)}:\\
A: What should we do about this problem? Quit? B: Hmm, there's no other choice.

\textbf{Notes (JA)}:\\
「それしかない」は、話し手が選択肢が限られていることを強調し、状況を受け入れる際に使う口語的な表現。

\textbf{Notes (EN)}:\\
'Sore shika nai.' is a colloquial expression used when the speaker wants to emphasize that options are limited and to accept the situation. 

\newpage

\section*{629. Either way}

\textbf{Date}: 2025.12.30 (Tue)\\
\textbf{Index}: dochirademo\\
\textbf{Expression (JA)}: どちらでも、\\
\textbf{Romanization}: dochira demo,\\
\textbf{Expression (EN)}: Either way,

\textbf{Variants (JA)}: どちらにしても、 / どっちでも、 / いずれでも、\\
\textbf{Variants (EN)}: Either way, / Any one is fine, / Whichever you prefer, ...

\textbf{Intent (JA)}: 選択の自由の提示, 中立的立場の表明, 柔軟な対応の示唆\\
\textbf{Intent (EN)}: presentation of freedom of choice, expression of neutral stance, suggestion of flexible response

\textbf{Tags (JA)}: 即時文法, 口語, 選択表現, 中立表現, 柔軟表現\\
\textbf{Tags (EN)}: immediate grammar, spoken, expression of choice, expression of neutrality, expression of flexibility

\textbf{Adjusted Expression (JA)}:\\
どちらでも(結構:けっこう)です\\
\textbf{Adjusted Expression (EN)}:\\
Either one is fine.

\textbf{Example (JA)}:\\
A「お(昼:ひる)は(和食:わしょく)？(洋食:ようしょく)？」 B「どちらでも、」\\
\textbf{Example (EN)}:\\
A: For lunch, Japanese food or Western food? B: Either way, ...

\textbf{Notes (JA)}:\\
「どちらでも、」は、話し手が選択肢に対して中立的な立場を示し、相手に選択の自由を与える際に使う口語的な表現。もう少し詳しくこのことばの構成を説明すると、「どちらでも」は「either way」や「whichever」の意味で、選択肢に対して特定の偏りがないことを示しています。話し手は、柔軟な対応を示唆し、相手に選択を委ねています。

\textbf{Notes (EN)}:\\
'Dochira demo,' is a colloquial expression used when the speaker wants to indicate a neutral stance towards options and give the listener freedom of choice. To explain the structure of this phrase in more detail, 'dochira demo' means 'either way' or 'whichever', indicating that there is no specific bias towards any option. The speaker suggests a flexible response and entrusts the choice to the listener.

\newpage

\section*{628. Somehow, it might be good.}

\textbf{Date}: 2025.12.29 (Mon)\\
\textbf{Index}: nankaikimo\\
\textbf{Expression (JA)}: なんか、(良:い)いかも、\\
\textbf{Romanization}: nanka, ii kamo.\\
\textbf{Expression (EN)}: Somehow, it might be good.

\textbf{Variants (JA)}: なんとなく、(良:い)いかもしれない。 / なんだか、(良:い)いかもね。 / なんとなく、(良:い)い気がする。\\
\textbf{Variants (EN)}: Somehow, it might be good. / Somehow, it can be good. / Somehow, I feel it might be good.

\textbf{Intent (JA)}: 肯定的な可能性の表現, 漠然とした感覚の共有, 評価生成途中の提示\\
\textbf{Intent (EN)}: expression of positive possibility, sharing of vague feelings, presentation of evaluation in progress

\textbf{Tags (JA)}: 即時文法, 口語, 肯定表現, 可能性表現, 漠然表現, 言いさし表現\\
\textbf{Tags (EN)}: immediate grammar, spoken, expression of affirmation, expression of possibility, vague expression, trailing-off expression

\textbf{Adjusted Expression (JA)}:\\
なんとなく、(良:い)いかもしれません。\\
\textbf{Adjusted Expression (EN)}:\\
Somehow, it could be good.

\textbf{Example (JA)}:\\
A「この(ヘッドホン:へっどほん)、どう？」 B「うん、なんか、(良:い)いかも」\\
\textbf{Example (EN)}:\\
A: How about these headphones? B: Yeah, somehow, they might be good.

\textbf{Notes (JA)}:\\
「なんか、いいかも」は、話し手がある物事に対して肯定的な可能性を感じていることを表現する際に使う口語的な表現。もう少し詳しくこのことばの構成を説明すると、「なんか」は「somehow」や「somewhat」の意味で、漠然とした感覚を示します。「いいかも」は「it might be good」の意味で、肯定的な可能性を示しています。これは言いさしの表現であり、インタラクションの継続を意図的に開いたままにして、自分の評価がまだ完全に形成されていないことを示しています。そうすることによって、相手からの反応や意見を促す効果もあります。

\textbf{Notes (EN)}:\\
'Nanka, ii kamo.' is a colloquial expression used when the speaker wants to express that they feel a positive possibility about something. To explain the structure of this phrase in more detail, 'nanka' functions like 'somehow', indicating an evaluation whose reason is not yet articulated. 'Ii kamo' means 'it might be good', indicating a positive possibility. This is a trailing-off expression, intentionally leaving the interaction open to indicate that the speaker's evaluation is not yet fully formed. By doing so, it also has the effect of prompting a response or opinion from the listener.

\newpage

\section*{627. and,.. yah,...}

\textbf{Date}: 2025.12.28 (Sun)\\
\textbf{Index}: te..hai\\
\textbf{Expression (JA)}: ...て、はい、...\\
\textbf{Romanization}: ...te, hai,...\\
\textbf{Expression (EN)}: ...and, yah,...

\textbf{Variants (JA)}: ...して、はい、... / ...やって、はい、... / ...行って、はい、...\\
\textbf{Variants (EN)}: ...and, yah,... / ...doing it, yah,... / ...going, yah,...

\textbf{Intent (JA)}: 話の継続, 同意の表現, 会話の流れ維持\\
\textbf{Intent (EN)}: continuation of conversation, expression of agreement, maintenance of conversational flow

\textbf{Tags (JA)}: 即時文法, 口語, 会話継続, 言いさし表現, 中途半端表現\\
\textbf{Tags (EN)}: immediate grammar, spoken, conversation continuation, trailing-off expression, incomplete expression

\textbf{Adjusted Expression (JA)}:\\
...して、そうですね、...\\
\textbf{Adjusted Expression (EN)}:\\
...and, well,...

\textbf{Example (JA)}:\\
A「(昨日:きのう)、(パーティ:ぱーてぃ)、どうだった」 B「まあ、(食:た)べて(飲:の)んで、はい、」 A「いつもと(同:おな)じかぁ」\\
\textbf{Example (EN)}:\\
A: How was the party yesterday? B: Well, we ate and drank, and, yah,... A: Same as always, huh?

\textbf{Notes (JA)}:\\
「...て、はい、...」は、話し手が会話を続ける際に使う言いさしの表現であり、相手に同意や理解を促すための口語的な表現。もう少し詳しくこのことばの構成を説明すると、「て」は動詞をつなぐのに使い、やったことの例を示しましょう。形容詞で感想を述べるときには「たのしくて」のように「くて」の形を使いましょう。「はい」は相手の反応を促すための間投詞です。話し手は、会話の流れを維持しつつ、相手に共感や同意を求めています。

\textbf{Notes (EN)}:\\
'...te, hai,...' is a trailing-off expression used by the speaker to continue the conversation and to prompt agreement or understanding from the listener in a colloquial manner. To explain the structure of this phrase in more detail, 'te' is used to connect verbs and to provide examples of actions taken. When expressing feelings with adjectives, use the 'kute' form, such as 'tanoshikute'. 'Hai' is an interjection used to elicit a response from the listener. The speaker maintains the flow of conversation while seeking empathy or agreement from the listener.

\newpage

\section*{626. You got a good one, didn't you?}

\textbf{Date}: 2025.12.27 (Sat)\\
\textbf{Index}: iinogettoshitane\\
\textbf{Expression (JA)}: いいの(ゲット:げっと)したね。\\
\textbf{Romanization}: ii no getto shita ne.\\
\textbf{Expression (EN)}: You got a good one, didn't you?

\textbf{Variants (JA)}: いいもの(手:て)に入:い)れたね / いい(品:しな)を(買:か)ったね / いい(物:もの)を(手:て)に(入:い)れたね\\
\textbf{Variants (EN)}: You got a good one, didn't you? / You got a good item, didn't you? / You bought a good product, didn't you?

\textbf{Intent (JA)}: 称賛の表現, 肯定的評価, 所有物の認識\\
\textbf{Intent (EN)}: expression of praise, positive evaluation, recognition of possession

\textbf{Tags (JA)}: 即時文法, 口語, 称賛表現, 肯定評価\\
\textbf{Tags (EN)}: immediate grammar, spoken, expression of praise, positive evaluation

\textbf{Adjusted Expression (JA)}:\\
いいものを(手:て)に(入:い)れましたね。\\
\textbf{Adjusted Expression (EN)}:\\
You obtained a good item, didn't you?

\textbf{Example (JA)}:\\
A「これ、(新:あたら)しい(キーボード:きーぼーど)だよ」 B「いいの(ゲット:げっと)したね」\\
\textbf{Example (EN)}:\\
A: This is a new keyboard. B: You got a good one, didn't you?

\textbf{Notes (JA)}:\\
「いいのゲットしたね」は、話し手が相手の所有物や取得物に対して称賛や肯定的な評価を表現する際に使う口語的な表現。もう少し詳しくこのことばの構成を説明すると、「いいの」は「a good one」の意味で、「ゲットした」は「got」や「obtained」の意味です。話し手は、相手が良いものを手に入れたことを認識し、それを称賛しています。

\textbf{Notes (EN)}:\\
'Ii no getto shita ne.' is a colloquial expression used when the speaker wants to express praise or positive evaluation of the listener's possession or acquisition. To explain the structure of this phrase in more detail, 'ii no' means 'a good one', and 'getto shita' means 'got' or 'obtained'. The speaker recognizes that the listener has obtained something good and praises them for it.

\newpage

\section*{625. Huh, isn't that my favorite one?}

\textbf{Date}: 2025.12.26 (Fri)\\
\textbf{Index}: esorebokunosukinayatsujanaidesuka\\
\textbf{Expression (JA)}: え、それ、(僕:ぼく)の(好:す)きなやつじゃないですか。\\
\textbf{Romanization}: e, sore, boku no suki na yatsu janai desu ka.\\
\textbf{Expression (EN)}: Huh, isn't that my favorite one?

\textbf{Variants (JA)}: あれ、それ、(私:わたし)の(好:す)きなやつじゃないですか。 / おお、それ、(俺:おれ)の(好:す)きなやつじゃないですか。 / うわ、それ、(あたし)の(好:す)きなやつじゃないですか。\\
\textbf{Variants (EN)}: Huh, isn't that my favorite one? / Oh, isn't that my favorite one? / Wow, isn't that my favorite one?

\textbf{Intent (JA)}: 驚きの表現, 喜びの表現, 所有物の認識\\
\textbf{Intent (EN)}: expression of surprise, expression of joy, recognition of possession

\textbf{Tags (JA)}: 即時文法, 口語, 驚き表現, 喜び表現\\
\textbf{Tags (EN)}: immediate grammar, spoken, expression of surprise, expression of joy

\textbf{Adjusted Expression (JA)}:\\
あれは、(私:わたし)の(好:す)きなものではありませんか。\\
\textbf{Adjusted Expression (EN)}:\\
Isn't that my favorite one?

\textbf{Example (JA)}:\\
A「これ、(新:あたら)しい(ゲーム:げーむ)だよ」 B「え、それ、(僕:ぼく)の(好:す)きなやつじゃないですか」\\
\textbf{Example (EN)}:\\
A: This is a new game. B: Huh, isn't that my favorite one?

\textbf{Notes (JA)}:\\
「え、それ、僕の好きなやつじゃないですか。」は、話し手が驚きや喜びを表現し、自分の所有物や好みを認識する際に使う口語的な表現。もう少し詳しくこのことばの構成を説明すると、「え」は驚きを表す感嘆詞で、「それ」は指示代名詞、「僕の好きなやつ」は「my favorite one」の意味です。少々、長い表現ですが、自分の心境を自然に表わしています。

\textbf{Notes (EN)}:\\
'E, sore, boku no suki na yatsu janai desu ka.' is a colloquial expression used when the speaker wants to express surprise or joy and recognize their possession or preference. To explain the structure of this phrase in more detail, 'e' is an interjection expressing surprise, 'sore' is a demonstrative pronoun, and 'boku no suki na yatsu' means 'my favorite one'. Although it is a somewhat long expression, it naturally conveys the speaker's feelings.

\newpage

\section*{624. In society's view,}

\textbf{Date}: 2025.12.25 (Thu)\\
\textbf{Index}: yononakakiniha\\
\textbf{Expression (JA)}: (世:よ)の(中:なか)(的:てき)には、\\
\textbf{Romanization}: yo no naka teki ni wa,\\
\textbf{Expression (EN)}: In society's view,

\textbf{Variants (JA)}: 社会的には、 / 世間的には、 / 一般的には、\\
\textbf{Variants (EN)}: In society's view, / From a societal perspective, / Generally speaking,

\textbf{Intent (JA)}: 社会的視点の提示, 一般的認識の表明, 価値観の共有\\
\textbf{Intent (EN)}: presentation of societal viewpoint, expression of general perception, sharing of values

\textbf{Tags (JA)}: 即時文法, 口語, 社会的視点, 一般認識\\
\textbf{Tags (EN)}: immediate grammar, spoken, societal viewpoint, general perception

\textbf{Adjusted Expression (JA)}:\\
(社会:しゃかい)の(見解:けんかい)では、\\
\textbf{Adjusted Expression (EN)}:\\
In society's view,

\textbf{Example (JA)}:\\
A「この(仕事:しごと)、(大変:たいへん)だよね」 B「うん、(世:よ)の(中:なか)(的:てき)には、(当:あ)たり前なんだろうけど」\\
\textbf{Example (EN)}:\\
A: This job is tough, isn't it? B: Yeah, from society's view, it's probably considered normal.

\textbf{Notes (JA)}:\\
「世の中的には、」は、社会全体の一般的な認識や価値観を述べる際に使う口語的な表現。もう少し詳しくこのことばの構成を説明すると、「世の中」は「society」や「the world」の意味で、「的には」は「from the perspective of」や「in terms of」の意味です。話し手は、社会的な視点から物事を捉えていることを示しています。

\textbf{Notes (EN)}:\\
'Yo no naka teki ni wa,' is a colloquial expression used when stating the general perception or values of society as a whole. To explain the structure of this phrase in more detail, 'yo no naka' means 'society' or 'the world', and 'teki ni wa' means 'from the perspective of' or 'in terms of'. The speaker indicates that they are viewing matters from a societal perspective.

\newpage

\section*{623. In my opinion,}

\textbf{Date}: 2025.12.24 (Wed)\\
\textbf{Index}: bokutekiniha\\
\textbf{Expression (JA)}: ぼく的には、\\
\textbf{Romanization}: boku teki ni wa,\\
\textbf{Expression (EN)}: In my opinion,

\textbf{Variants (JA)}: わたし的には、 / おれ的には、 / あたし的には、\\
\textbf{Variants (EN)}: In my opinion, / From my perspective, / As for me,

\textbf{Intent (JA)}: 個人の意見表明, 主観的視点の提示, 自己表現\\
\textbf{Intent (EN)}: expression of personal opinion, presentation of subjective viewpoint, self-expression

\textbf{Tags (JA)}: 即時文法, 口語, 意見表明, 主観視点\\
\textbf{Tags (EN)}: immediate grammar, spoken, expression of opinion, subjective viewpoint

\textbf{Adjusted Expression (JA)}:\\
(私:わたし)の(意見:いけん)では、\\
\textbf{Adjusted Expression (EN)}:\\
In my opinion,

\textbf{Example (JA)}:\\
A「この(映画:えいが)、(面白:おもしろ)かった？」 B「ぼく(的:てき)には、あまり(良:よ)くなかったかな」\\
\textbf{Example (EN)}:\\
A: Was this movie interesting? B: In my opinion, it wasn't very good.

\textbf{Notes (JA)}:\\
「ぼく的には、」は、話し手が自分の主観的な意見や視点を述べる際に使う口語的な表現。もう少し詳しくこのことばの構成を説明すると、「ぼく」は男性が自分を指す一人称で、「的には」は「from the perspective of」や「in terms of」の意味です。話し手は、自分の意見が他の人と異なる可能性があることを示しています。

\textbf{Notes (EN)}:\\
'Boku teki ni wa,' is a colloquial expression used when the speaker is stating their own subjective opinion or viewpoint. To explain the structure of this phrase in more detail, 'boku' is a first-person pronoun used by males, and 'teki ni wa' means 'from the perspective of' or 'in terms of'. The speaker indicates that their opinion may differ from that of others.

\newpage

\section*{622. cheaply and easily}

\textbf{Date}: 2025.12.23 (Tue)\\
\textbf{Index}: yasuyasuto\\
\textbf{Expression (JA)}: (安安:やすやす)と\\
\textbf{Romanization}: yasu yasu to\\
\textbf{Expression (EN)}: cheaply and easily

\textbf{Variants (JA)}: 手軽に / 簡単に / 巧妙に\\
\textbf{Variants (EN)}: easily / simply / cleverly

\textbf{Intent (JA)}: 容易, 騙す, 手口\\
\textbf{Intent (EN)}: ease, deceive, tactics

\textbf{Tags (JA)}: 即時文法, 口語, 騙す, 容易\\
\textbf{Tags (EN)}: immediate grammar, spoken, deception, easiness

\textbf{Adjusted Expression (JA)}:\\
(簡単:かんたん)に\\
\textbf{Adjusted Expression (EN)}:\\
easily

\textbf{Example (JA)}:\\
A「これはなかなか(壊:こわ)れませんよ」 B「そう(安安:やすやす)と(騙:だま)されるもんか」\\
\textbf{Example (EN)}:\\
A: This is quite durable. B: I won't be easily deceived.

\textbf{Notes (JA)}:\\
「安安と」は、物事が非常に簡単に行われることであって、価格が非常に安いことを意味する口語ではありません。この表現は、物事が容易に達成されることを強調するために使われます。もう少し詳しくこのことばの構成を説明すると、「安安」は「easily」や「without difficulty」の意味で、「と」は副詞を形成する助詞です。だいたいあまり良いことには使われません。

\textbf{Notes (EN)}:\\
'Yasu yasu to' is not a colloquial expression meaning that something is very cheap in price, but rather that something is done very easily. This expression is used to emphasize that something is achieved with ease. To explain the structure of this phrase in more detail, 'yasu yasu' means 'easily' or 'without difficulty', and 'to' is a particle that forms an adverb. It is generally not used for positive contexts.

\newpage

\section*{621. No, that's...}

\textbf{Date}: 2025.12.22 (Mon)\\
\textbf{Index}: iyasorewa\\
\textbf{Expression (JA)}: いや、それは...\\
\textbf{Romanization}: iya, sore wa...\\
\textbf{Expression (EN)}: No, that's...

\textbf{Variants (JA)}: いや、それはちょっと... / いや、それはなんていえばいいか... / いや、それはねぇ...\\
\textbf{Variants (EN)}: No, that's a bit... / No, how should I say that... / No, that's well...

\textbf{Intent (JA)}: 即答を避ける, 申し訳なさの表明, 相手との関係を保つ\\
\textbf{Intent (EN)}: avoiding a direct response, expressing apologetic stance, maintaining interpersonal harmony

\textbf{Tags (JA)}: 即時文法, 口語, 否定表現, 躊躇表現\\
\textbf{Tags (EN)}: immediate grammar, spoken, expression of negation, expression of hesitation

\textbf{Adjusted Expression (JA)}:\\
いいえ、それはどういえばよいのでしょうか、どうも...\\
\textbf{Adjusted Expression (EN)}:\\
No, how should I say that, well...

\textbf{Example (JA)}:\\
A「これ、先日のお礼です」 B「いや、それは...」\\
\textbf{Example (EN)}:\\
A: This is a thank you for the other day. B: No, that's...

\textbf{Notes (JA)}:\\
「いや、それは...」は、発話内容を完成させる前に、相手の行為に対する申し訳なさや距離調整を先に示す言いさしである。明確な否定や理由説明を行わなくても、この時点でインタラクションは成立しており、続きを言わない選択も含めて機能する。

\textbf{Notes (EN)}:\\
'Iya, sore wa...' is a phrase used to express an apologetic stance or to adjust interpersonal distance before completing the utterance. Even without providing a clear negation or explanation, the interaction is established at this point, and the choice to not continue can also function effectively.

\newpage

\section*{620. It's not just about being usable or working,}

\textbf{Date}: 2025.12.21 (Sun)\\
\textbf{Index}: tsukaerebaiiugokebaiitteiunjanakute\\
\textbf{Expression (JA)}: (使:つか)えれば(良:い)い、(動:うご)けば(良:い)いっていうんじゃなくて、\\
\textbf{Romanization}: tsukaereba ii, ugokeba ii tte iu n janakute,\\
\textbf{Expression (EN)}: It's not just about being usable or working,

\textbf{Variants (JA)}: (使:つか)えるだけでなく、(動:うご)くだけでなく、 / (使:つか)えれば(良:よ)いとか、(動:うご)けば(良:よ)いとかじゃなくて、 / (使:つか)えることや(動:うご)くことだけが重要ではなくて、\\
\textbf{Variants (EN)}: It's not just about being usable or working, / It's not just about being able to use it or having it work, / It's not only important that it can be used or that it works,

\textbf{Intent (JA)}: 価値観の表現, 品質の重視, 機能性の超越\\
\textbf{Intent (EN)}: expression of values, emphasis on quality, transcendence of functionality

\textbf{Tags (JA)}: 即時文法, 口語, 価値観表現, 品質重視\\
\textbf{Tags (EN)}: immediate grammar, spoken, expression of values, emphasis on quality

\textbf{Adjusted Expression (JA)}:\\
(使:つか)えることや(動:うご)くことだけが(重要:じゅうよう)ではありません、\\
\textbf{Adjusted Expression (EN)}:\\
It's not only important that it can be used or that it works,

\textbf{Example (JA)}:\\
A「この(ソフト:そふと)、(使:つか)えれば(良:い)い、(動:うご)けば(良:い)っていうんじゃなくて、(使:つか)いやすさも(重視:じゅうし)してほしい」 B「そうだね、(見:み)た目も(大事:だいじ)だよね」\\
\textbf{Example (EN)}:\\
A: It's not just about being usable or working, I also want ease of use to be emphasized. B: That's right, appearance is also important.

\textbf{Notes (JA)}:\\
「使えればいい、動けばいいっていうんじゃなくて、」は、物事の価値や重要性が単に機能的な側面にとどまらないことを強調する際に使う口語的な表現。もう少し詳しくこのことばの構成を説明すると、「使えればいい」は「it's okay if it can be used」の意味で、「動けばいい」は「it's okay if it works」の意味です。話し手は、機能性だけでなく、品質や使いやすさなど他の要素も重視していることを示しています。

\textbf{Notes (EN)}:\\
'Tsukaereba ii, ugokeba ii tte iu n janakute,' is a colloquial expression used to emphasize that the value or importance of something goes beyond just its functional aspects. To explain the structure of this phrase in more detail, 'tsukaereba ii' means 'it's okay if it can be used', and 'ugokeba ii' means 'it's okay if it works'. The speaker indicates that they value not only functionality but also other elements such as quality and ease of use.

\newpage

\section*{619. That's not true.}

\textbf{Date}: 2025.12.20 (Sat)\\
\textbf{Index}: sonnakotonaiyo\\
\textbf{Expression (JA)}: そんなことないよ。\\
\textbf{Romanization}: sonna koto nai yo.\\
\textbf{Expression (EN)}: That's not true.

\textbf{Variants (JA)}: そうじゃないよ。 / そんなことはないよ。 / そういうわけではないよ。\\
\textbf{Variants (EN)}: That's not the case. / That's not so. / That's not the reason.

\textbf{Intent (JA)}: 否定の表現, 反論の提示, 意見の違い\\
\textbf{Intent (EN)}: expression of negation, presentation of counterargument, difference of opinion

\textbf{Tags (JA)}: 即時文法, 口語, 否定表現, 反論提示\\
\textbf{Tags (EN)}: immediate grammar, spoken, expression of negation, presentation of counterargument

\textbf{Adjusted Expression (JA)}:\\
そうではありません。\\
\textbf{Adjusted Expression (EN)}:\\
That's not the case.

\textbf{Example (JA)}:\\
A「(彼:かれ)は(本当:ほんとう)に(忙:いそが)しいんだよ」 B「そんなことないよ、(昨日:きのう)も(遊:あそ)んでたじゃない」\\
\textbf{Example (EN)}:\\
A: He is really busy. B: That's not true, he was hanging out yesterday too.

\textbf{Notes (JA)}:\\
「そんなことないよ。」は、ある意見や主張に対して否定的な立場を示す際に使う口語的な表現。主観的な意見や一般的な認識に対して反論を提示する場合に用いられます。

\textbf{Notes (EN)}:\\
'Sonna koto nai yo.' is a colloquial expression used to indicate a negative stance towards a certain opinion or assertion. It is employed when presenting a counterargument against subjective opinions or general perceptions.

\newpage

\section*{618. Even if you say that.}

\textbf{Date}: 2025.12.19 (Fri)\\
\textbf{Index}: sonnakotoittatte\\
\textbf{Expression (JA)}: そんなこと(言:い)ったって。\\
\textbf{Romanization}: sonna koto itta tte.\\
\textbf{Expression (EN)}: Even if you say that.

\textbf{Variants (JA)}: そう言われても。 / そんなことを言われても。 / そう言ったところで。\\
\textbf{Variants (EN)}: Even if you say so. / Even if you are told that. / Even if you say that.

\textbf{Intent (JA)}: 譲歩の表現, 反論の提示, 意見の違い\\
\textbf{Intent (EN)}: expression of concession, presentation of counterargument, difference of opinion

\textbf{Tags (JA)}: 即時文法, 口語, 譲歩表現, 反論提示\\
\textbf{Tags (EN)}: immediate grammar, spoken, expression of concession, presentation of counterargument

\textbf{Adjusted Expression (JA)}:\\
そう言われましても、\\
\textbf{Adjusted Expression (EN)}:\\
Even if you say so,

\textbf{Example (JA)}:\\
A「(彼:かれ)は(本当:ほんとう)に(忙:いそが)しいんだよ」 B「そんなこと(言:い)ったって、(約束:やくそく)は(守:まも)ってほしい」\\
\textbf{Example (EN)}:\\
A: He is really busy. B: Even if you say that, I want him to keep his promise.

\textbf{Notes (JA)}:\\
「そんなこと言ったって。」は、ある意見や主張に対して譲歩的な立場を示す際に使う口語的な表現。主観的な意見や一般的な認識に対して反論を提示する場合に用いられます。

\textbf{Notes (EN)}:\\
'Sonna koto itta tte.' is a colloquial expression used to indicate a concessive stance towards a certain opinion or assertion. It is employed when presenting a counterargument against subjective opinions or general perceptions.

\newpage

\section*{617. That's apples and oranges.}

\textbf{Date}: 2025.12.18 (Thu)\\
\textbf{Index}: sorehasorekorehakore\\
\textbf{Expression (JA)}: それはそれ、これはこれ\\
\textbf{Romanization}: sore wa sore, kore wa kore.\\
\textbf{Expression (EN)}: That's that, and this is this.

\textbf{Variants (JA)}: それぞれ別々のものだ。 / それとこれは別物だ。 / それはそれ、これは別だ。\\
\textbf{Variants (EN)}: That's apples and oranges. / That's a different matter. / That and this are not the same thing.

\textbf{Intent (JA)}: 区別の表現, 独立性の強調, 対比の提示\\
\textbf{Intent (EN)}: expression of distinction, emphasis on independence, presentation of contrast

\textbf{Tags (JA)}: 即時文法, 口語, 区別表現, 対比提示\\
\textbf{Tags (EN)}: immediate grammar, spoken, expression of distinction, presentation of contrast

\textbf{Adjusted Expression (JA)}:\\
それらは(別:べつ)の(事柄:ことがら)として(扱:あつか)うべきです。\\
\textbf{Adjusted Expression (EN)}:\\
They should be treated as separate matters.

\textbf{Example (JA)}:\\
A「さっきも(僕:ぼく)が(払:はら)ったんだから、いっしょに(払:はら)うよ」 B「ちがう、それはそれ、これはこれ」\\
\textbf{Example (EN)}:\\
A: I already paid earlier, so I'll pay together. B: No, that's apples and oranges.

\textbf{Notes (JA)}:\\
「それはそれ、これはこれ。」は、二つの事柄や状況が異なるものであり、互いに影響を与えないことを強調する際に使う口語的な表現。もう少し詳しくこのことばの構成を説明すると、「それはそれ」は「that is that」の意味で、「これはこれ」は「this is this」の意味です。英語では 'That's apples and oranges.' や 'That's a different matter.' がよく対応します。

\textbf{Notes (EN)}:\\
'Sore wa sore, kore wa kore.' is a colloquial expression used to emphasize that two matters or situations are different and do not affect each other. To explain the structure of this phrase in more detail, 'sore wa sore' means 'that is that', and 'kore wa kore' means 'this is this'. In English, it often corresponds to 'That's apples and oranges.' or 'That's a different matter.'

\newpage

\section*{616. No, that's not necessarily the case.}

\textbf{Date}: 2025.12.17 (Wed)\\
\textbf{Index}: iya,soutohakagiranaiyo\\
\textbf{Expression (JA)}: いや、そうとは(限:かぎ)らないよ。\\
\textbf{Romanization}: iya, sō to wa kagiranai yo.\\
\textbf{Expression (EN)}: No, that's not necessarily the case.

\textbf{Variants (JA)}: いや、(必:かなら)ずしもそうではないよ。 / いや、(一概:いちがい)にそうとは言えないよ。 / いや、(全:まった)くそうとは(限:かぎ)らないよ。\\
\textbf{Variants (EN)}: No, that's not always the case. / No, you can't say that for sure. / No, that's not entirely true.

\textbf{Intent (JA)}: 否定の表現, 反論の提示, 意見の違い\\
\textbf{Intent (EN)}: expression of negation, presentation of counterargument, difference of opinion

\textbf{Tags (JA)}: 即時文法, 口語, 否定表現, 反論提示\\
\textbf{Tags (EN)}: immediate grammar, spoken, expression of negation, presentation of counterargument

\textbf{Adjusted Expression (JA)}:\\
いいえ、そうとは(限:かぎ)りません。\\
\textbf{Adjusted Expression (EN)}:\\
No, that's not necessarily the case.

\textbf{Example (JA)}:\\
A「だいたいの(人:ひと)がこれでいいと思っている」 B「いや、そうとは(限:かぎ)らないよ。」\\
\textbf{Example (EN)}:\\
A: Most people think this is fine. B: No, that's not necessarily the case.

\textbf{Notes (JA)}:\\
「いや、そうとは限らないよ。」は、ある意見や主張に対して否定的な立場を示す際に使う口語的な表現。主観的な意見や一般的な認識に対して反論を提示する場合に用いられます。

\textbf{Notes (EN)}:\\
'Iya, sō to wa kagiranai yo.' is a colloquial expression used to indicate a negative stance towards a certain opinion or assertion. It is employed when presenting a counterargument against subjective opinions or general perceptions.

\newpage

\section*{615. The reason is...}

\textbf{Date}: 2025.12.16 (Tue)\\
\textbf{Index}: nandekatouiuto\\
\textbf{Expression (JA)}: なんで...かと(言:い)うと、\\
\textbf{Romanization}: nande... ka to iu to,\\
\textbf{Expression (EN)}: The reason is...

\textbf{Variants (JA)}: なぜ...かと(言:い)うと、 / ...の(理由:りゆう)は、 / ...の(原因:げんいん)は、\\
\textbf{Variants (EN)}: Why... is because... / The reason for... is... / The cause of... is...

\textbf{Intent (JA)}: 説明の導入, 理由の提示, 背景の説明\\
\textbf{Intent (EN)}: introduction of explanation, presentation of reason, explanation of background

\textbf{Tags (JA)}: 即時文法, 口語, 説明導入, 理由提示\\
\textbf{Tags (EN)}: immediate grammar, spoken, introduction of explanation, presentation of reason

\textbf{Adjusted Expression (JA)}:\\
...の(理由:りゆう)は、\\
\textbf{Adjusted Expression (EN)}:\\
The reason for... is...

\textbf{Example (JA)}:\\
A「(最近:さいきん)、(元気:げんき)がないね」 B「なんでかと言うと、(仕事:しごと)が(忙:いそが)しくて..」\\
\textbf{Example (EN)}:\\
A: You haven't been energetic lately. B: The reason is that I've been busy with work...

\textbf{Notes (JA)}:\\
「なんで...かと言うと、」は、ある事柄や状況の理由や背景を説明する際に使う口語的な表現。もう少し詳しくこのことばの構成を説明すると、「なんで」は「why」や「the reason is」の意味で、「かと言うと」は「is because」や「the reason for... is」の意味です。

\textbf{Notes (EN)}:\\
'Nande... ka to iu to,' is a colloquial expression used when explaining the reason or background of a certain matter or situation. To explain the structure of this phrase in more detail, 'nande' means 'why' or 'the reason is', and 'ka to iu to' means 'is because' or 'the reason for... is'.

\newpage

\section*{614. isn't .. really energetic?}

\textbf{Date}: 2025.12.15 (Mon)\\
\textbf{Index}: sugoigenkijanaidesuka\\
\textbf{Expression (JA)}: (今:いま)、..がすごい(元気:げんき)じゃないですか\\
\textbf{Romanization}: ima, .. ga sugoi genki janai desu ka\\
\textbf{Expression (EN)}: isn't .. really energetic?

\textbf{Variants (JA)}: (最近:さいきん)、..がすごく(元気:げんき)ですよね / (最近:さいきん)、..がとても(元気:げんき)じゃないですか / (今:いま)、..がかなり(元気:げんき)じゃないですか\\
\textbf{Variants (EN)}: lately, .. is really energetic, isn't it? / lately, .. is very energetic, isn't it? / now, .. is quite energetic, isn't it?

\textbf{Intent (JA)}: 確認の表現, 感嘆の表現, 状況の報告\\
\textbf{Intent (EN)}: expression of confirmation, expression of admiration, reporting of situation

\textbf{Tags (JA)}: 即時文法, 口語, 確認表現, 感嘆表現, 一気に言うフレーズ\\
\textbf{Tags (EN)}: immediate grammar, spoken, expression of confirmation, expression of admiration, phrases said in one breath

\textbf{Adjusted Expression (JA)}:\\
(最近:さいきん)、..がとても(活発:かっぱつ)です\\
\textbf{Adjusted Expression (EN)}:\\
lately, .. is very active

\textbf{Example (JA)}:\\
A「(今:いま)、DTMが、(すごい)(元気:げんき)じゃないですか」 B「うん、(素人:しろうと)でも(簡単:かんたん)に(作:つく)れるからね」\\
\textbf{Example (EN)}:\\
A: Isn't DTM really energetic now? B: Yeah, because even amateurs can easily create music with it.

\textbf{Notes (JA)}:\\
「今、..がすごい元気じゃないですか」は、ある人や物事が非常に活発でエネルギッシュであることを確認し、感嘆する際に使う口語的な表現。もう少し詳しくこのことばの構成を説明すると、「今」は「now」や「lately」の意味で、「すごい元気じゃないですか」は「isn't really energetic?」や「is very active, isn't it?」の意味です。DTM（デスクトップミュージック）やプログラム言語など、技術的な分野での進歩や普及に関連して使われることが多い。

\textbf{Notes (EN)}:\\
'Ima, .. ga sugoi genki janai desu ka' is a colloquial expression used to confirm and express admiration for a person or thing that is very active and energetic. To explain the structure of this phrase in more detail, 'ima' means 'now' or 'lately', and 'sugoi genki janai desu ka' means 'isn't really energetic?' or 'is very active, isn't it?'. It is often used in relation to advancements or popularization in technical fields such as DTM (desktop music) or programming languages.

\newpage

\section*{613. even if...}

\textbf{Date}: 2025.12.14 (Sun)\\
\textbf{Index}: tokorode,ta\\
\textbf{Expression (JA)}: 〜たところで\\
\textbf{Romanization}: 〜ta tokoro de\\
\textbf{Expression (EN)}: even if...

\textbf{Variants (JA)}: 〜ても / 〜たって / 〜たとしても\\
\textbf{Variants (EN)}: even if \\textasciitilde{} / even though \\textasciitilde{} / even assuming \\textasciitilde{}

\textbf{Intent (JA)}: 譲歩の表現, 条件の提示, 否定的な結論\\
\textbf{Intent (EN)}: expression of concession, presentation of condition, negative conclusion

\textbf{Tags (JA)}: 即時文法, 口語, 譲歩表現, 条件提示\\
\textbf{Tags (EN)}: immediate grammar, spoken, expression of concession, presentation of condition

\textbf{Adjusted Expression (JA)}:\\
〜ても\\
\textbf{Adjusted Expression (EN)}:\\
even if \\textasciitilde{}

\textbf{Example (JA)}:\\
A「(言:い)ったところで、(意味:いみ)がない」 B「そうかもしれないけど、(伝:つた)えないと..」\\
\textbf{Example (EN)}:\\
A: Even if I say it, it has no meaning. B: That may be true, but if you don't convey it...

\textbf{Notes (JA)}:\\
「〜たところで」は、ある行動や状況が起こったとしても、その結果が期待されるものではない、または意味がないことを表現する際に使う口語的な表現。もう少し詳しくこのことばの構成を説明すると、「たところで」は「even if \\textasciitilde{}」や「even though \\textasciitilde{}」の意味で、話し手が述べた内容に対して譲歩的な立場を示します。

\textbf{Notes (EN)}:\\
'\\textasciitilde{}Ta tokoro de' is a colloquial expression used to indicate that even if a certain action or situation occurs, the result is not what is expected or is meaningless. To explain the structure of this phrase in more detail, 'ta tokoro de' means 'even if \\textasciitilde{}' or 'even though \\textasciitilde{}', indicating a concessive stance in relation to what the speaker has stated.

\newpage

\section*{612. not that..but...}

\textbf{Date}: 2025.12.13 (Sat)\\
\textbf{Index}: janakute\\
\textbf{Expression (JA)}: 〜じゃなくて、\\
\textbf{Romanization}: 〜janakute,\\
\textbf{Expression (EN)}: not that... but...

\textbf{Variants (JA)}: 〜ではなくて、 / 〜じゃないけど、 / 〜ではないけれども、\\
\textbf{Variants (EN)}: not that..., but... / it's not that..., but... / it's not that..., however...

\textbf{Intent (JA)}: 否定の表現, 対比の強調, 説明の導入\\
\textbf{Intent (EN)}: expression of negation, emphasis on contrast, introduction of explanation

\textbf{Tags (JA)}: 即時文法, 口語, 否定表現, 対比強調\\
\textbf{Tags (EN)}: immediate grammar, spoken, expression of negation, emphasis on contrast

\textbf{Adjusted Expression (JA)}:\\
〜ではなく、\\
\textbf{Adjusted Expression (EN)}:\\
not that..., but...

\textbf{Example (JA)}:\\
A「いいとか、(悪:わる)いとかじゃなくて..」 B「でも、(結局:けっきょく)、したくないんでしょ」\\
\textbf{Example (EN)}:\\
A: It's not that it's good or bad... B: But in the end, you don't want to do it, do you?

\textbf{Notes (JA)}:\\
「〜じゃなくて、」は、ある事柄や状況を否定しつつ、別の事柄や状況を強調する際に使う口語的な表現。もう少し詳しくこのことばの構成を説明すると、「じゃなくて」は「not that... but...」の意味で、話し手が述べた内容に対して否定的な立場を示します。

\textbf{Notes (EN)}:\\
'\\textasciitilde{}Janakute,' is a colloquial expression used to negate a certain matter or situation while emphasizing another matter or situation. To explain the structure of this phrase in more detail, 'janakute' means 'not that... but...', indicating a negative stance in relation to what the speaker has stated.

\newpage

\section*{611. only this much / just this much}

\textbf{Date}: 2025.12.12 (Fri)\\
\textbf{Index}: tattakoredake\\
\textbf{Expression (JA)}: たったこれだけ\\
\textbf{Romanization}: tatta kore dake\\
\textbf{Expression (EN)}: only this much / just this much

\textbf{Variants (JA)}: これだけしかない / これだけしか持っていない / これだけしか残っていない\\
\textbf{Variants (EN)}: there's only this much / I only have this much / only this much is left

\textbf{Intent (JA)}: 限定の表現, 量の強調, 驚きの表現\\
\textbf{Intent (EN)}: expression of limitation, emphasis on quantity, expression of surprise

\textbf{Tags (JA)}: 即時文法, 口語, 限定表現, 量強調\\
\textbf{Tags (EN)}: immediate grammar, spoken, expression of limitation, emphasis on quantity

\textbf{Adjusted Expression (JA)}:\\
これだけしかありません\\
\textbf{Adjusted Expression (EN)}:\\
there's only this much

\textbf{Example (JA)}:\\
A「はい、お(小遣:こづか)い」 B「え、たったこれだけ？」\\
\textbf{Example (EN)}:\\
A: Here is your allowance. B: Huh, only this much?

\textbf{Notes (JA)}:\\
「たったこれだけ」は、ある物事や数量が非常に少ないことを強調する際に使う口語的な表現。もう少し詳しくこのことばの構成を説明すると、「たった」は「only」や「just」の意味で、「これだけ」は「this much」や「this amount」の意味です。あまりにも少なくて、驚きや失望を感じる場合に使われることが多い。

\textbf{Notes (EN)}:\\
'Tatta kore dake' is a colloquial expression used to emphasize that a certain thing or quantity is very small. To explain the structure of this phrase in more detail, 'tatta' means 'only' or 'just', and 'kore dake' means 'this much' or 'this amount'. It is often used when the amount is so small that it evokes surprise or disappointment.

\newpage

\section*{610. It could be that way / I guess there are such things}

\textbf{Date}: 2025.12.11 (Thu)\\
\textbf{Index}: souiunomoarukana\\
\textbf{Expression (JA)}: そういうのもあるかな\\
\textbf{Romanization}: sou iu no mo aru kana\\
\textbf{Expression (EN)}: It could be that way / I guess there are such things

\textbf{Variants (JA)}: そういうこともあるかも / そういう場合もあるかな / そういう状況もあるかもしれない\\
\textbf{Variants (EN)}: That could be the case / There might be such situations / Such circumstances may exist

\textbf{Intent (JA)}: 可能性の表現, 状況の認識, 柔軟な思考\\
\textbf{Intent (EN)}: expression of possibility, recognition of situations, flexible thinking

\textbf{Tags (JA)}: 即時文法, 口語, 可能性表現, 状況認識\\
\textbf{Tags (EN)}: immediate grammar, spoken, expression of possibility, recognition of situations

\textbf{Adjusted Expression (JA)}:\\
そのようなことも(考:かんが)えられますね\\
\textbf{Adjusted Expression (EN)}:\\
That could be the case

\textbf{Example (JA)}:\\
A「この(花:はな)、(店先:みせさき)に(置:お)くよりも、(中:なか)の(方:ほう)がよくない？」 B「そうねぇ、そういうのもあるかな」\\
\textbf{Example (EN)}:\\
A: Wouldn't it be better to place these flowers inside rather than at the storefront? B: Yeah, I guess there are such things.

\textbf{Notes (JA)}:\\
「そういうのもあるかな」は、ある状況や可能性を認識し、それに対して柔軟な思考を示す際に使う口語的な表現。もう少し詳しくこのことばの構成を説明すると、「そういうの」は「such things」や「that kind of thing」の意味で、「..もあるかな」は「could be」や「might exist」の意味です。

\textbf{Notes (EN)}:\\
'Sou iu no mo aru kana' is a colloquial expression used to recognize a certain situation or possibility and to indicate flexible thinking in response. To explain the structure of this phrase in more detail, 'sou iu no' means 'such things' or 'that kind of thing', and '..mo aru kana' means 'could be' or 'might exist'. 

\newpage

\section*{609. to lend me / to lend us}

\textbf{Date}: 2025.12.10 (Wed)\\
\textbf{Index}: kashitekureru\\
\textbf{Expression (JA)}: (貸:か)してくれる\\
\textbf{Romanization}: kashite kureru\\
\textbf{Expression (EN)}: to lend me / to lend us

\textbf{Variants (JA)}: (貸:か)してもらえる / (貸:か)してくださる / (貸:か)してあげる\\
\textbf{Variants (EN)}: can you lend me/us / would you kindly lend me/us / I will lend you/us

\textbf{Intent (JA)}: 依頼の表現, 助けを求める, 親しみの表現\\
\textbf{Intent (EN)}: expression of request, seeking help, expression of familiarity

\textbf{Tags (JA)}: 即時文法, 口語, 依頼表現, 助け求める\\
\textbf{Tags (EN)}: immediate grammar, spoken, expression of request, seeking help

\textbf{Adjusted Expression (JA)}:\\
(貸:か)してください\\
\textbf{Adjusted Expression (EN)}:\\
please lend me/us

\textbf{Example (JA)}:\\
A「これ、(貸:か)してくれる？」 B「うん、でも、(必:かなら)ず、(返:かえ)してよ」\\
\textbf{Example (EN)}:\\
A: Can you lend me this? B: Yeah, but make sure to return it.

\textbf{Notes (JA)}:\\
「貸してくれる」は、物やサービスを一時的に借りることを依頼する際に使う口語的な表現。もう少し詳しくこのことばの構成を説明すると、「貸して」は「to lend」の意味で、「くれる」は「to give me/us」の意味です。「借りる」という動詞がありますが、これは「to borrow」の意味であり、話し手が物を借りる側であることを示します。

\textbf{Notes (EN)}:\\
'Kashite kureru' is a colloquial expression used when requesting to borrow something temporarily. To explain the structure of this phrase in more detail, 'kashite' means 'to lend', and 'kureru' means 'to give me/us'. There is also the verb 'kariru', which means 'to borrow', indicating that the speaker is on the borrowing side.

\newpage

\section*{608. doesn't fit / doesn't suit}

\textbf{Date}: 2025.12.09 (Tue)\\
\textbf{Index}: awanai\\
\textbf{Expression (JA)}: (合:あ)わない\\
\textbf{Romanization}: awanai\\
\textbf{Expression (EN)}: doesn't fit / doesn't suit

\textbf{Variants (JA)}: 似合わない / 適切でない / 調子が合わない\\
\textbf{Variants (EN)}: doesn't match / is not appropriate / is out of sync

\textbf{Intent (JA)}: 不適合の表現, 不一致の表現, 否定的な評価\\
\textbf{Intent (EN)}: expression of incompatibility, expression of mismatch, negative evaluation

\textbf{Tags (JA)}: 即時文法, 口語, 不適合表現, 否定的評価\\
\textbf{Tags (EN)}: immediate grammar, spoken, expression of incompatibility, negative evaluation

\textbf{Adjusted Expression (JA)}:\\
(合:あ)いません\\
\textbf{Adjusted Expression (EN)}:\\
does not fit

\textbf{Example (JA)}:\\
A「この(服:ふく)、(合:あ)わないよね」 B「うん、(色:いろ)も(柄:がら)も(好:す)きじゃないし」\\
\textbf{Example (EN)}:\\
A: These clothes don't fit, do they? B: Yeah, I don't like the color or the pattern either.

\textbf{Notes (JA)}:\\
「合わない」は、物事や状況が適切でない、または一致しないことを表現する際に使う口語的な表現。もう少し詳しくこのことばの構成を説明すると、「合わない」は「doesn't fit」や「doesn't suit」の意味で、話し手が述べた内容に対して不適合や不一致を示します。ものがその場に適していない場合や、期待に沿わない場合、そして、対人関係においても使われ、「あの人とは合わない」というように使われることもあります。

\textbf{Notes (EN)}:\\
'Awanai' is a colloquial expression used to indicate that something is not appropriate or does not match. To explain the structure of this phrase in more detail, 'awanai' means 'doesn't fit' or 'doesn't suit', indicating incompatibility or mismatch in relation to what the speaker has stated. It can be used when something is not suitable for a particular situation, does not meet expectations, and can also be applied in interpersonal relationships, as in 'I don't get along with that person.'

\newpage

\section*{607. to make a long story short}

\textbf{Date}: 2025.12.08 (Mon)\\
\textbf{Index}: hayaihanashi\\
\textbf{Expression (JA)}: (早:はや)い(話:はなし)\\
\textbf{Romanization}: hayai hanashi\\
\textbf{Expression (EN)}: to make a long story short

\textbf{Variants (JA)}: (手短:てみじか)に(言:い)うと / (要:よう)するに / (簡単:かんたん)に(言:い)うと\\
\textbf{Variants (EN)}: to put it briefly / in short / to say it simply

\textbf{Intent (JA)}: 要約の表現, 簡潔な説明, 話のまとめ\\
\textbf{Intent (EN)}: expression of summary, concise explanation, conclusion of story

\textbf{Tags (JA)}: 即時文法, 口語, 要約表現, 簡潔説明\\
\textbf{Tags (EN)}: immediate grammar, spoken, expression of summary, concise explanation

\textbf{Adjusted Expression (JA)}:\\
(手短:てみじか)に(言:い)いますと\\
\textbf{Adjusted Expression (EN)}:\\
to put it briefly

\textbf{Example (JA)}:\\
A「(昨日:きのう)の(会議:かいぎ)、(長:なが)かったね」 B「うん、(早:はや)い(話:はなし)、(誰:だれ)かが(社長:しゃちょう)のお(弁当:べんとう)を(買:か)ってくるということ？」\\
\textbf{Example (EN)}:\\
A: Yesterday's meeting was long, wasn't it? B: Yeah, to make a long story short, someone ended up being the one to run the errand for the president's bento, right?

\textbf{Notes (JA)}:\\
「早い話」は、長い説明や複雑な経緯を省いて結論だけを述べる際に使う口語的な表現。単に要点をまとめるだけでなく、長い過程に比べて結末がくだらない場合などに、軽い皮肉やユーモアを含むこともある。「早い」は「quick」や「fast」、「話」は「story」や「talk」の意味で、英語では「to make a long story short」や「to put it briefly」がよく対応する。

\textbf{Notes (EN)}:\\
'Hayai hanashi' is a colloquial expression used when the speaker wants to skip long explanations and present only the conclusion. It not only summarizes the main point but may also carry a light sense of irony or humor, especially when the final outcome is trivial compared to the long process leading up to it. Literally, 'hayai' means 'quick' or 'fast', and 'hanashi' means 'story' or 'talk'. Natural English equivalents include 'to make a long story short' and 'to put it briefly'.

\newpage

\section*{606. I ended up ... / I accidentally ...}

\textbf{Date}: 2025.12.07 (Sun)\\
\textbf{Index}: socchiniicchatta\\
\textbf{Expression (JA)}: そっちに(行:い)っちゃった\\
\textbf{Romanization}: socchi ni icchatta\\
\textbf{Expression (EN)}: I ended up going that way

\textbf{Variants (JA)}: そっちに(行:い)っちゃったんだよね / そっちに(行:い)っちゃったよ / そっちに(行:い)っちゃったんだ\\
\textbf{Variants (EN)}: I accidentally went that way / I ended up going in that direction / I went that way

\textbf{Intent (JA)}: 結果の表現, 予期せぬ行動の説明, 状況の報告\\
\textbf{Intent (EN)}: expression of result, explanation of unexpected action, reporting of situation

\textbf{Tags (JA)}: 即時文法, 口語, 結果表現, 行動説明\\
\textbf{Tags (EN)}: immediate grammar, spoken, expression of result, explanation of action

\textbf{Adjusted Expression (JA)}:\\
(話題:わだい)が(逸:そ)れています\\
\textbf{Adjusted Expression (EN)}:\\
The topic has strayed

\textbf{Example (JA)}:\\
A「(猫:ねこ)が(好:す)きなんですねぇ」 B「(何:なん)の(話:はなし)だったっけ？また、(そっち:そっち)に(行:い)っちゃった」\\
\textbf{Example (EN)}:\\
A: You like cats, don't you? B: What were we talking about? Again, I ended up going in that direction.

\textbf{Notes (JA)}:\\
「そっちに行っちゃった」は、話題や注意が別の方向に逸れたことを表現する際に使う口語的な表現。もう少し詳しくこのことばの構成を説明すると、「そっちに」は「that way」や「in that direction」の意味で、「行っちゃった」は「ended up going」や「accidentally went」の意味です。つまり、「そっちに行っちゃった」は文脈によって 'I ended up ...' とも 'I accidentally ...' とも対応し、英語では結果としてそちらの方向へ行ってしまったことを表します。

\textbf{Notes (EN)}:\\
'Socchi ni icchatta' is a colloquial expression used to indicate that the topic or attention has strayed in another direction. To explain the structure of this phrase in more detail, 'socchi ni' means 'that way' or 'in that direction', and 'icchatta' means 'ended up going' or 'accidentally went'. Therefore, depending on the context, 'socchi ni icchatta' can correspond to both 'I ended up ...' and 'I accidentally ...', expressing that one has gone in that direction as a result.

\newpage

\section*{605. I told you so...}

\textbf{Date}: 2025.12.06 (Sat)\\
\textbf{Index}: tteittanoni\\
\textbf{Expression (JA)}: って(言:い)ったのに\\
\textbf{Romanization}: tte itta no ni\\
\textbf{Expression (EN)}: I told you so...

\textbf{Variants (JA)}: って(言:い)ったじゃん / って(言:い)ったでしょ / って(言:い)ったのにさ\\
\textbf{Variants (EN)}: I told you, didn't I? / I told you, right? / I told you so, you know

\textbf{Intent (JA)}: 注意喚起, 不満の表現, 確認の強調\\
\textbf{Intent (EN)}: attention-getting, warning tone, calling attention to something

\textbf{Tags (JA)}: 即時文法, 口語, 注意喚起, 不満表現\\
\textbf{Tags (EN)}: immediate grammar, spoken, attention calling, expression of dissatisfaction

\textbf{Adjusted Expression (JA)}:\\
そのように(申:もう)し(上:あ)げました\\
\textbf{Adjusted Expression (EN)}:\\
I did mention that earlier

\textbf{Example (JA)}:\\
A「(忘:わす)れ(物:もの)しちゃった」 B「だから、(忘:わす)れるよって(言:い)ったのに」\\
\textbf{Example (EN)}:\\
A: I forgot it. B: I told you you'd forget it.

\textbf{Notes (JA)}:\\
「って言ったのに」は、以前に伝えた内容が守られなかったり、期待通りに行かなかった場合に、そのことを指摘する際に使う口語的な表現。もう少し詳しくこのことばの構成を説明すると、「って」は引用を示す助詞で、「言ったのに」は「I told you so」の意味です。つまり、「って言ったのに」は「I told you so...」や「I did mention that earlier」と同じ意味になります。だいたいは相手を非難することばなので、使うときは注意してくださいというよりも、使わない方がいいかもしれませんね。冗談ならいいですが。

\textbf{Notes (EN)}:\\
'Tte itta no ni' is a colloquial expression used to point out when previously conveyed information was not followed or when things did not go as expected. To explain the structure of this phrase in more detail, 'tte' is a particle indicating quotation, and 'itta no ni' means 'I told you so'. Therefore, 'tte itta no ni' has the same meaning as 'I told you so...' or 'I did mention that earlier.' It is generally used to criticize the other person, so be cautious when using it; in fact, it might be better not to use it at all, unless it's in a joking context.

\newpage

\section*{604. regardless of...}

\textbf{Date}: 2025.12.05 (Fri)\\
\textbf{Index}: tokakankeinaku\\
\textbf{Expression (JA)}: 〜とか(関係:かんけい)なく、\\
\textbf{Romanization}: toka kankei naku\\
\textbf{Expression (EN)}: regardless of...

\textbf{Variants (JA)}: 〜とか(気:き)に(し:し)なく / 〜とか(問題:もんだい)なく / 〜とか(影響:えいきょう)なく\\
\textbf{Variants (EN)}: regardless of \\textasciitilde{} / without worrying about \\textasciitilde{} / without concern for \\textasciitilde{}

\textbf{Intent (JA)}: 無関心の表現, 影響の否定, 条件の無視\\
\textbf{Intent (EN)}: expression of indifference, denial of influence, ignoring conditions

\textbf{Tags (JA)}: 即時文法, 口語, 無関心表現, 条件無視\\
\textbf{Tags (EN)}: immediate grammar, spoken, expression of indifference, ignoring conditions

\textbf{Adjusted Expression (JA)}:\\
〜などを(問:と)わず、\\
\textbf{Adjusted Expression (EN)}:\\
regardless of ...

\textbf{Example (JA)}:\\
A「(ラグビー:らぐびー)って(天気:てんき)、(天気:てんき)とか(関係:かんけい)なく、(ゲーム:げーむ)するの？」 B「ああ、(雨:あめ)でもやるよ。うん、(関係:かんけい)なくね」\\
\textbf{Example (EN)}:\\
A: In rugby, do they play games regardless of the weather? B: Yeah, they play even in the rain. Yeah, regardless of it.

\textbf{Notes (JA)}:\\
「〜とか関係なく？」は、ある事柄や状況に対して無関心であることや、その影響を受けないことを即時的に表現する口語的な表現。もう少し詳しくこのことばの構成を説明すると、「とか」は「regardless of」や「without worrying about」の意味で、「関係なく」は「without concern for」や「ignoring conditions」の意味です。つまり、「〜とか関係なく？」は「regardless of ...?」や「without worrying about ...?」と同じ意味になります。

\textbf{Notes (EN)}:\\
'\\textasciitilde{}Toka kankei naku?' is a colloquial expression used to immediately express indifference to a certain matter or situation, or that one is not influenced by it. To explain the structure of this phrase in more detail, 'toka' means 'regardless of' or 'without worrying about', and 'kankei naku' means 'without concern for' or 'ignoring conditions'. Therefore, '\\textasciitilde{}toka kankei naku?' has the same meaning as 'regardless of ...?' or 'without worrying about ...?'.

\newpage

\section*{603. hardly / scarcely / almost not}

\textbf{Date}: 2025.12.04 (Thu)\\
\textbf{Index}: hotonndo..nai\\
\textbf{Expression (JA)}: (ほとんど)...ない\\
\textbf{Romanization}: hotonndo...nai\\
\textbf{Expression (EN)}: hardly / scarcely / almost not

\textbf{Variants (JA)}: (ほぼ)...ない / (めったに)...ない / (ほとんど)...しない\\
\textbf{Variants (EN)}: (almost)...not / (rarely)...not / (hardly)...do not

\textbf{Intent (JA)}: 否定の強調, 頻度の表現, 程度の表現\\
\textbf{Intent (EN)}: emphasis on negation, expression of frequency, expression of degree

\textbf{Tags (JA)}: 即時文法, 口語, 否定強調, 頻度表現\\
\textbf{Tags (EN)}: immediate grammar, spoken, emphasis on negation, expression of frequency

\textbf{Adjusted Expression (JA)}:\\
(ほとんど)...ません\\
\textbf{Adjusted Expression (EN)}:\\
hardly ...

\textbf{Example (JA)}:\\
A「(週末:しゅうまつ)、(遊:あそ)びに(行:い)く？」 B「うーん、(ほとんど)(行:い)かない」\\
\textbf{Example (EN)}:\\
A: Are you going out on the weekend? B: Hmm, I hardly go.

\textbf{Notes (JA)}:\\
「ほとんど...ない」は、ある事柄や行動が非常に少ない頻度で行われることや、ほとんど存在しないことを強調する際に使う口語的な表現。もう少し詳しくこのことばの構成を説明すると、「ほとんど」は「hardly」や「scarcely」、「almost not」の意味で、「ない」は否定を示します。つまり、「ほとんど...ない」は「hardly ...」や「scarcely ...」と同じ意味になります。

\textbf{Notes (EN)}:\\
'Hotonndo...nai' is a colloquial expression used to emphasize that a certain matter or action occurs very infrequently or is almost non-existent. To explain the structure of this phrase in more detail, 'hotonndo' means 'hardly', 'scarcely', or 'almost not', and 'nai' indicates negation. Therefore, 'hotonndo...nai' has the same meaning as 'hardly ...' or 'scarcely ...'.

\newpage

\section*{602. not very / not much / not so}

\textbf{Date}: 2025.12.03 (Wed)\\
\textbf{Index}: annmari..nai\\
\textbf{Expression (JA)}: あんまり...ない\\
\textbf{Romanization}: anmari...nai\\
\textbf{Expression (EN)}: not very / not much / not so

\textbf{Variants (JA)}: あまり...ない / そんなに...ない / それほど...ない\\
\textbf{Variants (EN)}: not really ... / not that much ... / not to that extent ...

\textbf{Intent (JA)}: 否定の強調, 程度の表現, 評価の抑制\\
\textbf{Intent (EN)}: emphasis on negation, expression of degree, moderation of evaluation

\textbf{Tags (JA)}: 即時文法, 口語, 否定強調, 程度表現\\
\textbf{Tags (EN)}: immediate grammar, spoken, emphasis on negation, expression of degree

\textbf{Adjusted Expression (JA)}:\\
あまり...ません\\
\textbf{Adjusted Expression (EN)}:\\
not very ...

\textbf{Example (JA)}:\\
A「(料理:りょうり)、する？」 B「うーん、あんまり」\\
\textbf{Example (EN)}:\\
A: Do you cook? B: Hmm, not really.

\textbf{Notes (JA)}:\\
「あんまり...ない」は、ある事柄や行動があまり頻繁に行われないことや、程度が低いことを強調する際に使う口語的な表現。もう少し詳しくこのことばの構成を説明すると、「あんまり」は「not very」や「not much」の意味で、「ない」は否定を示します。つまり、「あんまり...ない」は「not very ...」や「not much ...」と同じ意味になります。

\textbf{Notes (EN)}:\\
'Anmari...nai' is a colloquial expression used to emphasize that a certain matter or action is not done very often or is of low degree. To explain the structure of this phrase in more detail, 'anmari' means 'not very' or 'not much', and 'nai' indicates negation. Therefore, 'anmari...nai' has the same meaning as 'not very ...' or 'not much ...'.

\newpage

\section*{601. worthless / trivial / silly}

\textbf{Date}: 2025.12.02 (Tue)\\
\textbf{Index}: kudaranai\\
\textbf{Expression (JA)}: (下:くだ)らない\\
\textbf{Romanization}: kudaranai\\
\textbf{Expression (EN)}: worthless / trivial / silly

\textbf{Variants (JA)}: (無価値:むかち)な / (取:と)るに(足:た)らない / (くだら)なくて(意味:いみ)がない\\
\textbf{Variants (EN)}: worthless / not worth considering / silly and meaningless

\textbf{Intent (JA)}: 否定的な評価, 価値の否定, 軽蔑の表現\\
\textbf{Intent (EN)}: negative evaluation, denial of value, expression of contempt

\textbf{Tags (JA)}: 即時文法, 口語, 否定的評価, 価値否定\\
\textbf{Tags (EN)}: immediate grammar, spoken, negative evaluation, denial of value

\textbf{Adjusted Expression (JA)}:\\
(無価値:むかち)な\\
\textbf{Adjusted Expression (EN)}:\\
worthless

\textbf{Example (JA)}:\\
A「この(映画:えいが)、(下:くだ)らないよね」 B「うん、でも、その(下:くだ)らないところがいいんじゃない」\\
\textbf{Example (EN)}:\\
A: This movie is silly, isn't it? B: Yeah, but isn't that silliness what makes it good?

\textbf{Notes (JA)}:\\
「くだらない」は、物事や出来事が価値がない、または重要でないことを表現する際に使う口語的な表現。もう少し詳しくこのことばの構成を説明すると、「くだらない」は「worthless」や「trivial」、「silly」の意味で、話し手が述べた内容に対して否定的な評価を示します。つまり、「くだらない」は「worthless」や「trivial」や「silly」と同じ意味になります。

\textbf{Notes (EN)}:\\
'Kudaranai' is a colloquial expression used to indicate that something is worthless or unimportant. To explain the structure of this phrase in more detail, 'kudaranai' means 'worthless', 'trivial', or 'silly', indicating a negative evaluation in relation to what the speaker has stated. Therefore, 'kudaranai' has the same meaning as 'worthless', 'trivial', or 'silly'.

\newpage

\section*{600. quite interesting / pretty interesting}

\textbf{Date}: 2025.12.01 (Mon)\\
\textbf{Index}: kekkoomoshiroi\\
\textbf{Expression (JA)}: (結構:けっこう)(面白:おもしろ)い\\
\textbf{Romanization}: kekko omoshiroi\\
\textbf{Expression (EN)}: quite interesting / pretty interesting

\textbf{Variants (JA)}: (かなり)(面白:おもしろ)い / (相当:そうとう)(面白:おもしろ)い / (なかなか)(面白:おもしろ)い\\
\textbf{Variants (EN)}: fairly interesting / considerably interesting / rather interesting

\textbf{Intent (JA)}: 評価の強調, 興味の表現, 肯定的な意見\\
\textbf{Intent (EN)}: emphasis on evaluation, expression of interest, positive opinion

\textbf{Tags (JA)}: 即時文法, 口語, 評価強調, 興味表現\\
\textbf{Tags (EN)}: immediate grammar, spoken, emphasis on evaluation, expression of interest

\textbf{Adjusted Expression (JA)}:\\
(非常:ひじょう)に(面白:おもしろ)い\\
\textbf{Adjusted Expression (EN)}:\\
very interesting

\textbf{Example (JA)}:\\
A「この(本:ほん)、(結構:けっこう)(面白:おもしろ)いよ」 B「へえ、(読:よ)んでみようかな」\\
\textbf{Example (EN)}:\\
A: This book is quite interesting. B: Oh, I might give it a read.

\textbf{Notes (JA)}:\\
「結構面白い」は、物事や出来事が予想以上に興味深いことを表現する際に使う口語的な表現。もう少し詳しくこのことばの構成を説明すると、「結構」は「quite」や「pretty」の意味で、「面白い」は「interesting」の意味です。つまり、「結構面白い」は「quite interesting」や「pretty interesting」と同じ意味になります。

\textbf{Notes (EN)}:\\
'Kekko omoshiroi' is a colloquial expression used to indicate that something is more interesting than expected. To explain the structure of this phrase in more detail, 'kekko' means 'quite' or 'pretty', and 'omoshiroi' means 'interesting'. Therefore, 'kekko omoshiroi' has the same meaning as 'quite interesting' or 'pretty interesting'.

\newpage

\section*{599. long-lasting / durable}

\textbf{Date}: 2025.11.30 (Sun)\\
\textbf{Index}: nagamochi\\
\textbf{Expression (JA)}: (長持:ながも)ち\\
\textbf{Romanization}: nagamochi\\
\textbf{Expression (EN)}: long-lasting / durable

\textbf{Variants (JA)}: (耐久:たいきゅう)性(高:たか)い / (丈夫:じょうぶ)な / (持続:じぞく)力(強:つよ)い\\
\textbf{Variants (EN)}: high durability / sturdy / strong / strong endurance

\textbf{Intent (JA)}: 耐久性の表現, 品質の強調, 持続力の表現\\
\textbf{Intent (EN)}: expression of durability, emphasis on quality, expression of endurance

\textbf{Tags (JA)}: 即時文法, 口語, 耐久性表現, 品質強調\\
\textbf{Tags (EN)}: immediate grammar, spoken, expression of durability, emphasis on quality

\textbf{Adjusted Expression (JA)}:\\
(耐久:たいきゅう)性が(高:たか)い\\
\textbf{Adjusted Expression (EN)}:\\
high durability

\textbf{Example (JA)}:\\
A「この(靴:くつ)、(長持:ながも)ちすんの？」 B「うん、3(年:ねん)は(持:も)つよ」\\
\textbf{Example (EN)}:\\
A: Are these shoes durable? B: Yeah, they'll last for three years.

\textbf{Notes (JA)}:\\
「長持ち」は、物や製品が長期間使用できることを表現する際に使う口語的な表現。もう少し詳しくこのことばの構成を説明すると、「長持ち」は「long-lasting」や「durable」の意味で、話し手が述べた内容に対して耐久性や持続力を示します。つまり、「長持ち」は「long-lasting」や「durable」と同じ意味になります。「長く持ちますか」という表現もありますが、そちらは少しフォーマルな響きがあります。動詞「持つ」はhaveという意味の方が一般的ですが、ここでは「持ちますか？」と聞くと「持つ」という動詞の意味が変わってしまうので注意してください。

\textbf{Notes (EN)}:\\
'Nagamochi' is a colloquial expression used to indicate that an object or product can be used for a long period of time. To explain the structure of this phrase in more detail, 'nagamochi' means 'long-lasting' or 'durable', indicating durability or endurance in relation to what the speaker has stated. Therefore, 'nagamochi' has the same meaning as 'long-lasting' or 'durable'. There is also the expression 'naga ku mochimasu ka', but that has a slightly more formal tone. While the verb 'motsu' generally means 'to have', in this context, asking 'mochimasu ka?' changes the meaning of the verb, so be cautious.

\newpage

\section*{598. suffer quietly}

\textbf{Date}: 2025.11.29 (Sat)\\
\textbf{Index}: nakineiri\\
\textbf{Expression (JA)}: (泣:な)き(寝入:ねい)り\\
\textbf{Romanization}: naki-neiri\\
\textbf{Expression (EN)}: suffer quietly

\textbf{Variants (JA)}: (我慢:がまん)の(限界:げんかい) / (耐:た)え(忍:しの)ぶ / (静:しず)かに(耐:た)える\\
\textbf{Variants (EN)}: the limit of endurance / to endure patiently / to endure quietly

\textbf{Intent (JA)}: 忍耐, 我慢, 内面的な苦しみの表現\\
\textbf{Intent (EN)}: patience, endurance, expression of internal suffering

\textbf{Tags (JA)}: 即時文法, 口語, 忍耐, 我慢\\
\textbf{Tags (EN)}: immediate grammar, spoken, patience, endurance

\textbf{Adjusted Expression (JA)}:\\
(理不尽:りふじん)にも(為:な)す(術:すべ)がない\\
\textbf{Adjusted Expression (EN)}:\\
to have no way to deal with such unfairness

\textbf{Example (JA)}:\\
A「(問:と)い(合:あ)わせても(結局:けっきょく)『(仕様:しよう)です』で(終:お)わっちゃった」 B「それって(泣:な)き(寝入:ねい)りだよね」\\
\textbf{Example (EN)}:\\
A: Even after I inquired, it just ended with 'That's how it is.' B: Yeah, that's suffering quietly, isn't it?

\textbf{Notes (JA)}:\\
「泣き寝入り」は、不当な扱いや不公平な状況に対して、声を上げずに我慢することを表現する口語的な表現。もう少し詳しくこのことばの構成を説明すると、「泣き」は「crying」の意味で、「寝入り」は「going to sleep」の意味です。つまり、「泣き寝入り」は「suffer quietly」や「endure silently」と同じ意味になります。

\textbf{Notes (EN)}:\\
'Naki-neiri' is a colloquial expression used to describe enduring unfair treatment or situations without raising one's voice. To explain the structure of this phrase in more detail, 'naki' means 'crying', and 'neiri' means 'going to sleep'. Therefore, 'naki-neiri' has the same meaning as 'suffer quietly' or 'endure silently'.

\newpage

\section*{597. to be available}

\textbf{Date}: 2025.11.28 (Fri)\\
\textbf{Index}: aiteru\\
\textbf{Expression (JA)}: (空:あ)いてる\\
\textbf{Romanization}: aiteru\\
\textbf{Expression (EN)}: to be available

\textbf{Variants (JA)}: (空:あ)いていますか / (席:せき)が(空:あ)いてる\\
\textbf{Variants (EN)}: is available? / is there an open seat?

\textbf{Intent (JA)}: 状態の表現, 質問, 確認\\
\textbf{Intent (EN)}: expression of state, question, confirmation

\textbf{Tags (JA)}: 即時文法, 口語, 状態表現, 質問\\
\textbf{Tags (EN)}: immediate grammar, spoken, expression of state, question

\textbf{Adjusted Expression (JA)}:\\
(空:あ)いていますか\\
\textbf{Adjusted Expression (EN)}:\\
is it available?

\textbf{Example (JA)}:\\
A「その(日:ひ)、(空:あ)いてる？」 B「うん、(何:なに)？」 A「(映画:えいが)、(見:み)に(行:い)こ」\\
\textbf{Example (EN)}:\\
A: Are you free that day? B: Yeah, what is it? A: Let's go see a movie.

\textbf{Notes (JA)}:\\
「空いてる」は、場所や時間が利用可能であることを即時的に表現する口語的な表現。もう少し詳しくこのことばの構成を説明すると、「空いてる」は「to be open」や「to be available」の意味で、話し手が述べた内容に対して利用可能な状態を示します。つまり、「空いてる」は「is open」や「is available?」と同じ意味になります。

\textbf{Notes (EN)}:\\
'Aiteru' is a colloquial expression used to immediately express that a place or time is available. To explain the structure of this phrase in more detail, 'aiteru' means 'to be open' or 'to be available', indicating the state of availability in relation to what the speaker has stated. Therefore, 'aiteru' has the same meaning as 'is open' or 'is available?'.

\newpage

\section*{596. to go somewhere to see sth}

\textbf{Date}: 2025.11.27 (Thu)\\
\textbf{Index}: shiniiku\\
\textbf{Expression (JA)}: しに(行:い)く\\
\textbf{Romanization}: mi ni iku\\
\textbf{Expression (EN)}: to go somewhere to do sth

\textbf{Variants (JA)}: (買:か)いに(行:い)く / (食:た)べに(行:い)く / (遊:あそ)びに(行:い)く\\
\textbf{Variants (EN)}: to go somewhere to buy sth / to go somewhere to eat sth / to go somewhere to play/hang out

\textbf{Intent (JA)}: 目的の表現, 行動の表現, 移動の表現\\
\textbf{Intent (EN)}: expression of purpose, expression of action, expression of movement

\textbf{Tags (JA)}: 即時文法, 口語, 目的表現, 行動表現\\
\textbf{Tags (EN)}: immediate grammar, spoken, expression of purpose, expression of action

\textbf{Adjusted Expression (JA)}:\\
しに(行:い)きます\\
\textbf{Adjusted Expression (EN)}:\\
to go to see sth

\textbf{Example (JA)}:\\
A「(新:あたら)しい(映画:えいが)、(見:み)に(行:い)く？」 B「うん、(行:い)こう」\\
\textbf{Example (EN)}:\\
A: Shall we go see the new movie? B: Yeah, let's go.

\textbf{Notes (JA)}:\\
「〜に行く」は、特定の目的を持ってどこかへ行くことを表現する際に使う口語的な表現。もう少し詳しくこのことばの構成を説明すると、「に」は目的地や目的を示す助詞で、「行く」は移動することを意味します。つまり、「〜に行く」は「to go somewhere to do sth」や「to go somewhere for a purpose」と同じ意味になります。これを速く言ってみてください。たぶん「みにーく」になりますよね。「買いにく」、「食べにく」、「遊びにく」も同様です。そちらの方が自然な発音になります。

\textbf{Notes (EN)}:\\
'\\textasciitilde{}Ni iku' is a colloquial expression used to indicate going somewhere with a specific purpose. To explain the structure of this phrase in more detail, 'ni' is a particle that indicates the destination or purpose, and 'iku' means to move or go. Therefore, '\\textasciitilde{}ni iku' has the same meaning as 'to go somewhere to do sth' or 'to go somewhere for a purpose.' Try saying it quickly. It probably sounds like 'mini-ku.' The same applies to 'kainiku,' 'tabeniku,' and 'asobiniku,' which also sound more natural when pronounced that way.

\newpage

\section*{595. soft framed evaluation}

\textbf{Date}: 2025.11.26 (Wed)\\
\textbf{Index}: arikanashika\\
\textbf{Expression (JA)}: (アリ:あり)か(ナシ:なし)か（でいえば）\\
\textbf{Romanization}: ari ka nashi ka (de ieba)\\
\textbf{Expression (EN)}: if I had to say yes or no

\textbf{Variants (JA)}: アリかナシかでいえば / アリ寄り / まあアリかな\\
\textbf{Variants (EN)}: if pressed, I'd say it's acceptable / leaning yes / I'd say it's fine

\textbf{Intent (JA)}: 即時的な評価, ソフトな肯定の提示, 理由の省略を可能にする判断, 選択肢フレーミングによる肯定誘導\\
\textbf{Intent (EN)}: immediate evaluation, soft-positive judgment, judgment that allows omission of reasons, framing that guides toward acceptance

\textbf{Tags (JA)}: 即時文法, 評価, ソフト肯定, 語用論\\
\textbf{Tags (EN)}: immediate grammar, evaluation, soft positive, pragmatics

\textbf{Adjusted Expression (JA)}:\\
(賛成:さんせい)か(反対:はんたい)かをはっきり(述:の)べなければならないならば、\\
\textbf{Adjusted Expression (EN)}:\\
to state explicitly whether one agrees or disagrees

\textbf{Example (JA)}:\\
A「この(案:あん)、(アリ:あり)か(ナシ:なし)かでいえば？」 B「うーん、(アリ:あり)か(ナシ:なし)かっていえば、(アリ:あり)かな。」\\
\textbf{Example (EN)}:\\
A: So, for this idea, if you had to say yes or no? B: Hmm... if I had to say, I'd say it's a yes.

\textbf{Notes (JA)}:\\
「アリかナシか（でいえば）」は、二択を提示しているように見えて、実際には「アリ」方向へ評価を誘導するソフト肯定の構文である。否定（ナシ）は理由負荷が高く、語用論的に言いづらいため、自然会話ではこの構文の後件が「ナシ」になることはほとんどない。話者が強い賛成や積極的な理由を持っていなくても、『弱い支持／弱い応援』を伝えることができる点に特徴がある。心理学のframing効果とも対応し、即時文法が持つ判断の即時性と理由省略の機能が顕著に現れる表現である。

\textbf{Notes (EN)}:\\
The phrase 'ari ka nashi ka (de ieba)' appears to offer a binary choice, but pragmatically it functions as a soft-positive construction. In natural Japanese conversation, giving a negative answer ('nashi') would require explanation, so the framing itself guides the listener toward a soft yes. The expression allows the speaker to show weak support or gentle approval without stating explicit reasons. This makes it an excellent example of Immediate Grammar: immediate evaluation, omission of justification, and a framing effect that biases the judgment toward acceptance.

\newpage

\section*{594. only / just / best}

\textbf{Date}: 2025.11.25 (Tue)\\
\textbf{Index}: kagiru\\
\textbf{Expression (JA)}: 〜に(限:かぎ)る\\
\textbf{Romanization}: 〜ni kagiru\\
\textbf{Expression (EN)}: only / just / best

\textbf{Variants (JA)}: 〜だけだ / 〜に決まっている / 〜が一番だ\\
\textbf{Variants (EN)}: it's only \\textasciitilde{} / it's definitely \\textasciitilde{} / \\textasciitilde{} is the best

\textbf{Intent (JA)}: 限定, 強調, 評価\\
\textbf{Intent (EN)}: limitation, emphasis, evaluation

\textbf{Tags (JA)}: 即時文法, 口語, 限定, 強調\\
\textbf{Tags (EN)}: immediate grammar, spoken, limitation, emphasis

\textbf{Adjusted Expression (JA)}:\\
〜が一番よい。\\
\textbf{Adjusted Expression (EN)}:\\
it's best to \\textasciitilde{}.

\textbf{Example (JA)}:\\
A「(夏:なつ)の(暑:あつ)いとき、(何:なに)が(一番:いちばん)？」 B「やっぱり、(冷:つめ)たい(飲:の)み物に限るよね」\\
\textbf{Example (EN)}:\\
A: What's the best thing during hot summers? B: After all, it's best to have a cold drink.

\textbf{Notes (JA)}:\\
「〜に限る」は、ある選択肢や行動が最も適切であることを強調する際に使う口語的な表現。もう少し詳しくこのことばの構成を説明すると、「に限る」は「only」や「just」、「best」の意味で、話し手が述べた内容に対して最も適した選択肢や行動を示します。つまり、「〜に限る」は「only \\textasciitilde{}」や「just \\textasciitilde{}」と同じ意味になります。

\textbf{Notes (EN)}:\\
'\\textasciitilde{}Ni kagiru' is a colloquial expression used to emphasize that a certain choice or action is the most appropriate. To explain the structure of this phrase in more detail, 'ni kagiru' means 'only', 'just', or 'best', indicating the most suitable choice or action in relation to what the speaker has stated. Therefore, '\\textasciitilde{}ni kagiru' has the same meaning as 'only \\textasciitilde{}' or 'just \\textasciitilde{}'.

\newpage

\section*{593. ...though...}

\textbf{Date}: 2025.11.24 (Mon)\\
\textbf{Index}: keredomo\\
\textbf{Expression (JA)}: 〜けれども...\\
\textbf{Romanization}: 〜keredomo.\\
\textbf{Expression (EN)}: ...though...

\textbf{Variants (JA)}: 〜けど... / 〜だけれども... / 〜ないけど...\\
\textbf{Variants (EN)}: ...but... / ...however... / ...though not...

\textbf{Intent (JA)}: 条件の追加, 対比, 譲歩\\
\textbf{Intent (EN)}: addition of conditions, contrast, concession

\textbf{Tags (JA)}: 即時文法, 口語, 条件追加, 対比\\
\textbf{Tags (EN)}: immediate grammar, spoken, addition of conditions, contrast

\textbf{Adjusted Expression (JA)}:\\
〜ことになりますけれども。\\
\textbf{Adjusted Expression (EN)}:\\
--though...

\textbf{Example (JA)}:\\
A「(新:あたら)しい(レストラン:れすとらん)、(美味:おい)しいよ」 B「でも、(高:たか)いんだよね、おいしいけれども...」\\
\textbf{Example (EN)}:\\
A: The new restaurant is delicious. B: But it's expensive, it's delicious though...

\textbf{Notes (JA)}:\\
「〜けれども...」は、ある事柄に対して条件を追加したり、対比を示したりする際に使う口語的な表現。もう少し詳しくこのことばの構成を説明すると、「けれども」は「but」や「however」の意味で、話し手が述べた内容に対して何かを付け加えたり、対照的な情報を提供したりする役割を果たします。つまり、「〜けれども...」は「...though...」や「...but...」と同じ意味になります。

\textbf{Notes (EN)}:\\
'\\textasciitilde{}Keredomo...' is a colloquial expression used to add conditions or indicate contrast regarding a certain matter. To explain the structure of this phrase in more detail, 'keredomo' means 'but' or 'however', serving the role of adding something to what the speaker has stated or providing contrasting information. Therefore, '\\textasciitilde{}keredomo...' has the same meaning as '...though...' or '...but...'.

\newpage

\section*{592. want to try to do}

\textbf{Date}: 2025.11.23 (Sun)\\
\textbf{Index}: mitainaa\\
\textbf{Expression (JA)}: 〜みたいなあ\\
\textbf{Romanization}: 〜mitai naa\\
\textbf{Expression (EN)}: want to try to do

\textbf{Variants (JA)}: 〜したいなあ / 〜やってみたいなあ / 〜経験してみたいなあ\\
\textbf{Variants (EN)}: I want to do \\textasciitilde{} / I want to try doing \\textasciitilde{} / I want to experience \\textasciitilde{}

\textbf{Intent (JA)}: 願望の表現, 希望の表現, 感情表現\\
\textbf{Intent (EN)}: expression of desire, expression of hope, emotional expression

\textbf{Tags (JA)}: 即時文法, 口語, 願望表現, 希望表現\\
\textbf{Tags (EN)}: immediate grammar, spoken, expression of desire, expression of hope

\textbf{Adjusted Expression (JA)}:\\
〜したいと思います。\\
\textbf{Adjusted Expression (EN)}:\\
I want to do \\textasciitilde{}.

\textbf{Example (JA)}:\\
A「(スカイダイビング:すかいだいびんぐ)って(楽:たの)しいらしいよ」 B「うん、(一度:いちど)、やってみたいなあ」\\
\textbf{Example (EN)}:\\
A: I heard skydiving is fun. B: Yeah, I want to try it once.

\textbf{Notes (JA)}:\\
「みたいなあ」は、何かをしたいという願望や希望を即時的に表現する口語的な表現。もう少し詳しくこのことばの構成を説明すると、「みたいなあ」は「I want to try to do \\textasciitilde{}」や「I want to experience ...」の意味で、話し手の願望や希望を強調します。

\textbf{Notes (EN)}:\\
'Mitai naa' is a colloquial expression used to immediately express a desire or hope to do something. To explain the structure of this phrase in more detail, '\\textasciitilde{}mitai naa' means 'I want to try to do ...' or 'I want to experience ...', emphasizing the speaker's desire or hope.

\newpage

\section*{591. that's all / only that much}

\textbf{Date}: 2025.11.22 (Sat)\\
\textbf{Index}: soredake\\
\textbf{Expression (JA)}: それだけ\\
\textbf{Romanization}: sore dake\\
\textbf{Expression (EN)}: that's all / only that much

\textbf{Variants (JA)}: それだけです / それだけしかない / それだけで十分\\
\textbf{Variants (EN)}: that's all there is / only that much / that's enough

\textbf{Intent (JA)}: 限定, 強調, 完結\\
\textbf{Intent (EN)}: limitation, emphasis, conclusion

\textbf{Tags (JA)}: 即時文法, 口語, 限定, 強調\\
\textbf{Tags (EN)}: immediate grammar, spoken, limitation, emphasis

\textbf{Adjusted Expression (JA)}:\\
それだけです。\\
\textbf{Adjusted Expression (EN)}:\\
That's all.

\textbf{Example (JA)}:\\
A「(今日:きょう)の(会議:かいぎ)、どうだった？」 B「うん、(重要:じゅうよう)な(話:はなし)はそれだけ」\\
\textbf{Example (EN)}:\\
A: How was today's meeting? B: Yeah, that's all the important stuff.

\textbf{Notes (JA)}:\\
「それだけ」は、ある事柄や情報が限定されていることを強調する際に使う口語的な表現。もう少し詳しくこのことばの構成を説明すると、「それ」は「that」の意味で、「だけ」は「only」や「just」の意味です。つまり、「それだけ」は「that's all」や「only that much」と同じ意味になります。

\textbf{Notes (EN)}:\\
'Sore dake' is a colloquial expression used to emphasize that a certain matter or information is limited. To explain the structure of this phrase in more detail, 'sore' means 'that', and 'dake' means 'only' or 'just'. Therefore, 'sore dake' has the same meaning as 'that's all' or 'only that much'.

\newpage

\section*{590. that's what I keep saying to him/her / I've been telling him/her that}

\textbf{Date}: 2025.11.21 (Fri)\\
\textbf{Index}: souitterundesukedone\\
\textbf{Expression (JA)}: そう(言:い)ってるんですけどね\\
\textbf{Romanization}: sou itterun desu kedo ne\\
\textbf{Expression (EN)}: that's what I keep saying to him/her / I've been telling him/her that

\textbf{Variants (JA)}: そう話しているんですけどね / そう伝えているんですけどね / そう説明しているんですけどね\\
\textbf{Variants (EN)}: that's what I've been telling him/her / that's what I've been conveying to him/her / that's what I've been explaining to him/her

\textbf{Intent (JA)}: 説明, 強調, 情報提供\\
\textbf{Intent (EN)}: explanation, emphasis, information provision

\textbf{Tags (JA)}: 即時文法, 口語, 説明, 強調\\
\textbf{Tags (EN)}: immediate grammar, spoken, explanation, emphasis

\textbf{Adjusted Expression (JA)}:\\
そう(言:い)っていますが、\\
\textbf{Adjusted Expression (EN)}:\\
I've been saying that, but

\textbf{Example (JA)}:\\
A「(彼:かれ)、また(遅刻:ちこく)？」 B「うん、(遅刻:)しちゃだめだよって、いつも、そう(言:い)ってるんですけどね」\\
\textbf{Example (EN)}:\\
A: Is he late again? B: Yeah, I always tell him, 'You shouldn't be late,' that's what I keep saying to him.

\textbf{Notes (JA)}:\\
「そう言ってるんですけどね」は、相手に対して何かを伝えたり説明したりしていることを強調する際に使う口語的な表現。もう少し詳しくこのことばの構成を説明すると、「そう言ってる」は「that's what I keep saying」や「I've been telling」の意味で、「んですけどね」は「but you know」や「though」の意味です。つまり、「そう言ってるんですけどね」は「that's what I keep saying to him/her」や「I've been telling him/her that」と同じ意味になります。

\textbf{Notes (EN)}:\\
'Sou itterun desu kedo ne' is a colloquial expression used to emphasize that you have been conveying or explaining something to someone. To explain the structure of this phrase in more detail, 'sou itteru' means 'that's what I keep saying' or 'I've been telling', and 'n desu kedo ne' means 'but you know' or 'though'. Therefore, 'sou itterun desu kedo ne' has the same meaning as 'that's what I keep saying to him/her' or 'I've been telling him/her that'.

\newpage

\section*{589. this is really something / this is serious}

\textbf{Date}: 2025.11.20 (Thu)\\
\textbf{Index}: korehontonohanashi\\
\textbf{Expression (JA)}: これ、(本当:ほんとう)の(話:はなし)\\
\textbf{Romanization}: kore hontou no hanashi\\
\textbf{Expression (EN)}: this is really something / this is serious / I'm telling you the truth

\textbf{Variants (JA)}: これは本物の話だ / これは真実の話だ / これは本格的な話だ\\
\textbf{Variants (EN)}: this is a genuine story / this is a true story / this is a serious matter

\textbf{Intent (JA)}: 強調, 驚きの表現, 情報提供\\
\textbf{Intent (EN)}: emphasis, expression of surprise, information provision

\textbf{Tags (JA)}: 即時文法, 口語, 強調, 驚きの表現\\
\textbf{Tags (EN)}: immediate grammar, spoken, emphasis, expression of surprise

\textbf{Adjusted Expression (JA)}:\\
これは本当の話です。\\
\textbf{Adjusted Expression (EN)}:\\
This is a true story.

\textbf{Example (JA)}:\\
A「たったの30(分:ぷん)で(全部:ぜんぶ)、(覚:おぼ)えたんだって。これ、(本当:ほんとう)の(話:はなし)」 B「えっ、(ホント:ほんと)？」\\
\textbf{Example (EN)}:\\
A: I heard he memorized everything in just 30 minutes. I'm telling you, this is really something. B: Really?

\textbf{Notes (JA)}:\\
「これ本当の話」は、ある情報や出来事が真実であることを強調する際に使う口語的な表現。もう少し詳しくこのことばの構成を説明すると、「これ」は「this」の意味で、「本当の話」は「true story」や「real story」の意味です。つまり、「これ本当の話」は「this is really something」や「this is serious」と同じ意味になります。カタカナで「ホント」と表記されることも多いです。

\textbf{Notes (EN)}:\\
'Kore hontou no hanashi' is a colloquial expression used to emphasize that certain information or events are true. To explain the structure of this phrase in more detail, 'kore' means 'this', and 'hontou no hanashi' means 'true story' or 'real story'. Therefore, 'kore hontou no hanashi' has the same meaning as 'this is really something' or 'this is serious'. It is also often written in katakana as 'honto'.

\newpage

\section*{588. I was wondering if...}

\textbf{Date}: 2025.11.19 (Wed)\\
\textbf{Index}: kanatoomotte\\
\textbf{Expression (JA)}: 〜かなと(思:おも)って。\\
\textbf{Romanization}: 〜kana to omotte.\\
\textbf{Expression (EN)}: I was wondering if...

\textbf{Variants (JA)}: 〜かしらと思って / 〜かなと思いまして / 〜かもと思って\\
\textbf{Variants (EN)}: I wonder if... / I was thinking that... / I thought maybe...

\textbf{Intent (JA)}: 疑問の表現, 丁寧な表現, 提案の前置き\\
\textbf{Intent (EN)}: expression of doubt, polite expression, prelude to suggestion

\textbf{Tags (JA)}: 即時文法, 口語, 疑問表現, 丁寧な表現\\
\textbf{Tags (EN)}: immediate grammar, spoken, expression of doubt, polite expression

\textbf{Adjusted Expression (JA)}:\\
〜かもしれないと思います。\\
\textbf{Adjusted Expression (EN)}:\\
I think it might be...

\textbf{Example (JA)}:\\
A「(明日:あした)、(雨:あめ)が(降:ふ)るかなと(思:おも)って...」 B「そうだね、たぶん」\\
\textbf{Example (EN)}:\\
A: I was wondering if it might rain tomorrow... B: Yeah, probably.

\textbf{Notes (JA)}:\\
「〜かなと思って。」は、何かについて疑問を持ったり、考えたりしていることを丁寧に表現する口語的な表現。もう少し詳しくこのことばの構成を説明すると、「〜かな」は「I wonder if...」や「I was thinking that...」の意味で、「と思って」は「I thought...」や「I was thinking...」の意味です。つまり、「〜かなと思って」は「I was wondering if...」や「I was thinking that...」と同じ意味になります。

\textbf{Notes (EN)}:\\
'\\textasciitilde{}Kana to omotte.' is a colloquial expression used to politely express doubt or contemplation about something. To explain the structure of this phrase in more detail, '\\textasciitilde{}kana' means 'I wonder if...' or 'I was thinking that...', and 'to omotte' means 'I thought...' or 'I was thinking...'. Therefore, '\\textasciitilde{}kana to omotte' has the same meaning as 'I was wondering if...' or 'I was thinking that...'.

\newpage

\section*{587. it's not like that / that's not the case}

\textbf{Date}: 2025.11.18 (Tue)\\
\textbf{Index}: souiukotojanai\\
\textbf{Expression (JA)}: そういうことじゃない\\
\textbf{Romanization}: sou iu koto janai\\
\textbf{Expression (EN)}: it's not like that / that's not the case

\textbf{Variants (JA)}: そんなことはない / そうではない / 違うよ\\
\textbf{Variants (EN)}: that's not true / that's not the case / that's not what I mean

\textbf{Intent (JA)}: 否定, 訂正, 説明\\
\textbf{Intent (EN)}: negation, correction, explanation

\textbf{Tags (JA)}: 即時文法, 口語, 否定, 訂正\\
\textbf{Tags (EN)}: immediate grammar, spoken, negation, correction

\textbf{Adjusted Expression (JA)}:\\
そうではありません\\
\textbf{Adjusted Expression (EN)}:\\
that's not the case

\textbf{Example (JA)}:\\
A「(私:わたし)が(悪:わる)かったんだ」 B「いや、そういうことじゃないよ。(君:きみ)のせいじゃないよ」\\
\textbf{Example (EN)}:\\
A: It was my fault. B: No, it's not like that. It's not your fault.

\textbf{Notes (JA)}:\\
「そういうことじゃない」は、相手の捉え方や解釈が誤っているときに、それをやんわり否定して軌道修正する口語的な表現です。「そういうこと」は「そのような事情・その解釈」の意味で、「じゃない」は「〜ではない」という否定を表します。したがって、「そういうことじゃない」は直訳すると「that's not the case」や「it's not like that」に相当し、「あなたの受け取り方は違うよ」「話の本質はそこではないよ」といったニュアンスを含みます。相手が自分を責めたり、誤った結論に飛びついたときに、その誤解をやさしく止めるための即時文法的な表現です。

\textbf{Notes (EN)}:\\
'Sou iu koto janai' is a colloquial expression used to gently negate and correct someone's interpretation when it does not match the actual situation. The phrase consists of 'sou iu koto' meaning 'that kind of thing' or 'that interpretation,' and 'janai' meaning 'is not.' Taken together, it corresponds to expressions such as 'it's not like that' or 'that's not the case.' It carries the nuance of softly correcting a misunderstanding; indicating that the person's interpretation is off, that the essence of the matter lies elsewhere, and that the issue is not simply a matter of 'it's your fault.' It is an immediate, spoken-style response used to stop someone from jumping to the wrong conclusion.

\newpage

\section*{586. seems fun / looks enjoyable}

\textbf{Date}: 2025.11.17 (Mon)\\
\textbf{Index}: nannkatanoshisou\\
\textbf{Expression (JA)}: (何:なん)か(楽:たの)しそう\\
\textbf{Romanization}: nanka tanoshisou\\
\textbf{Expression (EN)}: seems fun / looks enjoyable

\textbf{Variants (JA)}: 面白そう / 嬉しそう / ワクワクしそう\\
\textbf{Variants (EN)}: looks interesting / seems happy / looks exciting

\textbf{Intent (JA)}: 感情表現, 評価, 予測\\
\textbf{Intent (EN)}: emotional expression, evaluation, prediction

\textbf{Tags (JA)}: 即時文法, 口語, 感情表現, 評価\\
\textbf{Tags (EN)}: immediate grammar, spoken, emotional expression, evaluation

\textbf{Adjusted Expression (JA)}:\\
(楽:たの)しそうだ\\
\textbf{Adjusted Expression (EN)}:\\
seems fun

\textbf{Example (JA)}:\\
A「(新:あたら)しい(ゲーム:げーむ)、(面白:おもしろ)そうだね」 B「うん、なんか(楽:たの)しそう」\\
\textbf{Example (EN)}:\\
A: The new game looks interesting, doesn't it? B: Yeah, it seems fun.

\textbf{Notes (JA)}:\\
「なんか楽しそう」は、物事が楽しそうに見えることや、楽しさを感じさせる様子を即時的に表現する口語的な表現。もう少し詳しくこのことばの構成を説明すると、「なんか」は「somehow」や「something like that」の意味で、「楽しそう」は「seems fun」や「looks enjoyable」の意味です。つまり、「なんか楽しそう」は「seems fun」や「looks enjoyable」と同じ意味になります。

\textbf{Notes (EN)}:\\
'Nanka tanoshisou' is a colloquial expression used to immediately express that something looks fun or gives a sense of enjoyment. To explain the structure of this phrase in more detail, 'nanka' means 'somehow' or 'something like that', and 'tanoshisou' means 'seems fun' or 'looks enjoyable'. Therefore, 'nanka tanoshisou' has the same meaning as 'seems fun' or 'looks enjoyable'.

\newpage

\section*{585. it is not worth doing / it's not necessary to go to such lengths}

\textbf{Date}: 2025.11.16 (Sun)\\
\textbf{Index}: hodonokotowanai\\
\textbf{Expression (JA)}: (無理:むり)にするほどのことじゃない\\
\textbf{Romanization}: muri ni suru hodo no koto janai\\
\textbf{Expression (EN)}: it is not worth doing / it's not necessary to go to such lengths

\textbf{Variants (JA)}: そこまで(頑張:がんば)ることじゃない / そんなに(無理:むり)することじゃない / そこまで(努力:どりょく)することじゃない\\
\textbf{Variants (EN)}: it's not worth trying that hard / it's not necessary to push yourself that much / it's not worth making that much effort

\textbf{Intent (JA)}: 評価, 助言, 感情表現\\
\textbf{Intent (EN)}: evaluation, advice, emotional expression

\textbf{Tags (JA)}: 即時文法, 口語, 評価, 助言\\
\textbf{Tags (EN)}: immediate grammar, spoken, evaluation, advice

\textbf{Adjusted Expression (JA)}:\\
(無理:むり)をする(必要:ひつよう)はない\\
\textbf{Adjusted Expression (EN)}:\\
there is no need to overdo it

\textbf{Example (JA)}:\\
A「(試験:しけん)のために(徹夜:てつや)で(勉強:べんきょう)するよ」 B「いいや、(無理:むり)にするほどのことじゃないよ。ちゃんと(休:やす)んだほうが(良:い)いよ」\\
\textbf{Example (EN)}:\\
A: I'm going to study all night for the exam. B: No, it's not worth doing that. It's better to get some rest.

\textbf{Notes (JA)}:\\
「無理にするほどのことじゃない」は、ある行動や努力が過剰であることを評価し、その必要性を否定する際に使う口語的な表現。もう少し詳しくこのことばの構成を説明すると、「無理にする」は「to do something excessively」の意味で、「ほどのことじゃない」は「it's not worth doing」の意味です。つまり、「無理にするほどのことじゃない」は「it is not worth doing」や「it's not necessary to go to such lengths」と同じ意味になります。

\textbf{Notes (EN)}:\\
'Muri ni suru hodo no koto janai' is a colloquial expression used to evaluate that a certain action or effort is excessive and to deny its necessity. To explain the structure of this phrase in more detail, 'muri ni suru' means 'to do something excessively', and 'hodo no koto janai' means 'it's not worth doing'. Therefore, 'muri ni suru hodo no koto janai' has the same meaning as 'it is not worth doing' or 'it's not necessary to go to such lengths'.

\newpage

\section*{584. to have something to talk about / to have a matter to discuss}

\textbf{Date}: 2025.11.15 (Sat)\\
\textbf{Index}: hanashigaaru\\
\textbf{Expression (JA)}: (話:はなし)がある\\
\textbf{Romanization}: hanashi ga aru\\
\textbf{Expression (EN)}: to have something to talk about / to have a matter to discuss

\textbf{Variants (JA)}: 相談がある / 用事がある / 伝えたいことがある\\
\textbf{Variants (EN)}: to have a consultation / to have business / to have something to convey

\textbf{Intent (JA)}: 意思表示, コミュニケーション, 情報提供\\
\textbf{Intent (EN)}: expression of intention, communication, information provision

\textbf{Tags (JA)}: 即時文法, 口語, 意思表示, コミュニケーション\\
\textbf{Tags (EN)}: immediate grammar, spoken, expression of intention, communication

\textbf{Adjusted Expression (JA)}:\\
(相談:そうだん)がある\\
\textbf{Adjusted Expression (EN)}:\\
to have a consultation

\textbf{Example (JA)}:\\
A「ちょっと(話:はなし)があるんだけど、(今:いま)いい？」 B「うん、どうしたの？」\\
\textbf{Example (EN)}:\\
A: I have something to talk about, is now a good time? B: Yeah, what's up?

\textbf{Notes (JA)}:\\
「話がある」は、何か伝えたいことや相談したいことがある場合に使う口語的な表現。もう少し詳しくこのことばの構成を説明すると、「話」は「talk」や「matter to discuss」の意味で、「ある」は「to have」の意味です。つまり、「話がある」は「to have something to talk about」や「to have a matter to discuss」と同じ意味になります。

\textbf{Notes (EN)}:\\
'Hanashi ga aru' is a colloquial expression used when there is something to convey or discuss. To explain the structure of this phrase in more detail, 'hanashi' means 'talk' or 'matter to discuss', and 'aru' means 'to have'. Therefore, 'hanashi ga aru' has the same meaning as 'to have something to talk about' or 'to have a matter to discuss'.

\newpage

\section*{583. not good / not quite right}

\textbf{Date}: 2025.11.14 (Fri)\\
\textbf{Index}: anmarida\\
\textbf{Expression (JA)}: あんまりだ\\
\textbf{Romanization}: anmari da\\
\textbf{Expression (EN)}: not good / not quite right

\textbf{Variants (JA)}: ひどい / 不公平な / 不当な\\
\textbf{Variants (EN)}: terrible / unfair / unjust

\textbf{Intent (JA)}: 評価, 否定, 感情表現\\
\textbf{Intent (EN)}: evaluation, negation, emotional expression

\textbf{Tags (JA)}: 即時文法, 口語, 評価, 否定\\
\textbf{Tags (EN)}: immediate grammar, spoken, evaluation, negation

\textbf{Adjusted Expression (JA)}:\\
(酷:ひど)い(状況:じょうきょう)\\
\textbf{Adjusted Expression (EN)}:\\
terrible situation

\textbf{Example (JA)}:\\
A「(先生:せんせい)、(明日:あす)、(急:きゅう)に(試験:しけん)するって」 B「うん、そりゃ、あんまりだよ」\\
\textbf{Example (EN)}:\\
A: The teacher said there will be a sudden exam tomorrow. B: Oh, no way, that's not fair.

\textbf{Notes (JA)}:\\
「あんまりだ」は、物事がひどい状態であることや、不公平であることを即時的に表現する口語的な表現。もう少し詳しくこのことばの構成を説明すると、「あんまり」は「not good」や「not quite right」の意味で、物事が適切でない様子を指します。

\textbf{Notes (EN)}:\\
'Anmari da' is a colloquial expression used to immediately express that something is in a terrible state or is unfair. To explain the structure of this phrase in more detail, 'anmari' means 'not good' or 'not quite right', referring to a state where something is inappropriate.

\newpage

\section*{582. number / turn / order}

\textbf{Date}: 2025.11.13 (Thu)\\
\textbf{Index}: ban\\
\textbf{Expression (JA)}: (番:ばん)\\
\textbf{Romanization}: ban\\
\textbf{Expression (EN)}: number / turn / order

\textbf{Variants (JA)}: 順番 / 回数 / 順位\\
\textbf{Variants (EN)}: order / number of times / ranking

\textbf{Intent (JA)}: 順序の表現, 回数の表現, 順位の表現\\
\textbf{Intent (EN)}: expression of order, expression of number of times, expression of ranking

\textbf{Tags (JA)}: 即時文法, 口語, 順序表現, 回数表現\\
\textbf{Tags (EN)}: immediate grammar, spoken, order expression, number of times expression

\textbf{Adjusted Expression (JA)}:\\
(順番:じゅんばん)\\
\textbf{Adjusted Expression (EN)}:\\
order

\textbf{Example (JA)}:\\
A「(次:つぎ)は(誰:だれ)？」 B「(私:わたし)の(番:ばん)」\\
\textbf{Example (EN)}:\\
A: Who's next? B: It's my turn.

\textbf{Notes (JA)}:\\
「番」は、順序や回数、順位を表す際に使う口語的な表現。もう少し詳しくこのことばの構成を説明すると、「番」は「number」や「turn」、「order」の意味で、物事の順序や回数、順位を指します。

\textbf{Notes (EN)}:\\
'Ban' is a colloquial expression used to indicate order, number of times, or ranking. To explain the structure of this phrase in more detail, 'ban' means 'number', 'turn', or 'order', referring to the sequence, frequency, or ranking of things.

\newpage

\section*{581. I'll take this / I'll have this}

\textbf{Date}: 2025.11.12 (Wed)\\
\textbf{Index}: korenisuru\\
\textbf{Expression (JA)}: これにする\\
\textbf{Romanization}: kore ni suru\\
\textbf{Expression (EN)}: I'll take this / I'll have this

\textbf{Variants (JA)}: これを選ぶ / これを頼む / これに決める\\
\textbf{Variants (EN)}: I'll choose this / I'll order this / I'll decide on this

\textbf{Intent (JA)}: 選択, 意思表示, 決定\\
\textbf{Intent (EN)}: selection, expression of intention, decision

\textbf{Tags (JA)}: 即時文法, 口語, 選択, 意思表示\\
\textbf{Tags (EN)}: immediate grammar, spoken, selection, expression of intention

\textbf{Adjusted Expression (JA)}:\\
これを(選:えら)ぶ\\
\textbf{Adjusted Expression (EN)}:\\
I'll choose this

\textbf{Example (JA)}:\\
A「(何:なに)にする？」 B「うん、これにする」\\
\textbf{Example (EN)}:\\
A: What will you have? B: I'll take this.

\textbf{Notes (JA)}:\\
「これにする」は、選択や意思表示を即時的に表現する口語的な表現。もう少し詳しくこのことばの構成を説明すると、「これ」は「this」の意味で、「にする」は「to choose」や「to decide on」の意味です。つまり、「これにする」は「I'll take this」や「I'll have this」と同じ意味になります。

\textbf{Notes (EN)}:\\
'Kore ni suru' is a colloquial expression used to immediately express selection or intention. To explain the structure of this phrase in more detail, 'kore' means 'this', and 'ni suru' means 'to choose' or 'to decide on'. Therefore, 'kore ni suru' has the same meaning as 'I'll take this' or 'I'll have this'.

\newpage

\section*{580. in the first place / to begin with}

\textbf{Date}: 2025.11.11 (Tue)\\
\textbf{Index}: somosomo\\
\textbf{Expression (JA)}: そもそも\\
\textbf{Romanization}: somosomo\\
\textbf{Expression (EN)}: in the first place / to begin with

\textbf{Variants (JA)}: 最初から / 元々 / 根本的に\\
\textbf{Variants (EN)}: from the beginning / originally / fundamentally

\textbf{Intent (JA)}: 強調, 原因の提示, 議論の開始\\
\textbf{Intent (EN)}: emphasis, presentation of cause, start of discussion

\textbf{Tags (JA)}: 即時文法, 口語, 強調, 原因提示\\
\textbf{Tags (EN)}: immediate grammar, spoken, emphasis, cause presentation

\textbf{Adjusted Expression (JA)}:\\
(本来:ほんらい)は、\\
\textbf{Adjusted Expression (EN)}:\\
originally,

\textbf{Example (JA)}:\\
A「(彼:かれ)は(仕事:しごと)が(遅:おそ)いよね」 B「うん、そもそも(時間:じかん)に(ルーズ:るーず)なんだよ」\\
\textbf{Example (EN)}:\\
A: He's slow at work, isn't he? B: Yeah, to begin with, he's not punctual.

\textbf{Notes (JA)}:\\
「そもそも」は、物事の根本的な原因や起点を示す際に使う口語的な表現。もう少し詳しくこのことばの構成を説明すると、「そもそも」は「in the first place」や「to begin with」の意味で、話の出発点や基本的な前提を強調する際に用いられます。ただし、この表現には、根本から見直すべきだというニュアンスが含まれ、あまりよい意味で使われないこともあります。「そもそもスタートから間違っている」などの使い方が典型的です。

\textbf{Notes (EN)}:\\
'Somosomo' is a colloquial expression used to indicate the fundamental cause or starting point of a matter. To explain the structure of this phrase in more detail, 'somosomo' means 'in the first place' or 'to begin with', and it is used to emphasize the starting point of a discussion or basic premise. However, this expression often carries a nuance that suggests something should be reconsidered from the ground up, and it is not always used in a positive context. For example, it is typically used in phrases like 'the whole thing was wrong from the start.'

\newpage

\section*{579. rather / as long as}

\textbf{Date}: 2025.11.10 (Mon)\\
\textbf{Index}: baii,tsukaerebaii\\
\textbf{Expression (JA)}: (使:つか)えればいい\\
\textbf{Romanization}: tsukaereba ii\\
\textbf{Expression (EN)}: rather / as long as

\textbf{Variants (JA)}: ～ならばいい / ～であればいい / ～さえすればいい\\
\textbf{Variants (EN)}: as long as \\textasciitilde{} / if \\textasciitilde{} is fine / provided that \\textasciitilde{}

\textbf{Intent (JA)}: 条件設定, 妥協, 限定的な要求\\
\textbf{Intent (EN)}: condition setting, compromise, limited requirement

\textbf{Tags (JA)}: 即時文法, 口語, 条件設定, 妥協\\
\textbf{Tags (EN)}: immediate grammar, spoken, condition setting, compromise

\textbf{Adjusted Expression (JA)}:\\
(使用:しよう)できれば(良:よ)い\\
\textbf{Adjusted Expression (EN)}:\\
as long as it can be used

\textbf{Example (JA)}:\\
A「この(ソフト:そふと)、どう？」 B「あまり(良:よ)くない。(使:つか)えればいいとか、(動:うご)けばいいっていうんじゃなくて、(使:つか)いやすいかどうかなんだよ」\\
\textbf{Example (EN)}:\\
A: How's this software? B: It's not very good. It's not just about whether it can be used or works, but whether it's easy to use.

\textbf{Notes (JA)}:\\
「使えればいい」は、ある条件が満たされればそれで十分であることを即時的に表現する口語的な表現。もう少し詳しくこのことばの構成を説明すると、「使えれば」は「if it can be used」の意味で、「いい」は「is fine」や「is good」の意味です。つまり、「使えればいい」は「as long as it can be used」や「if it can be used, it's fine」と同じ意味になります。

\textbf{Notes (EN)}:\\
'Tsukaereba ii' is a colloquial expression used to immediately express that a certain condition being met is sufficient. To explain the structure of this phrase in more detail, 'tsukaereba' means 'if it can be used', and 'ii' means 'is fine' or 'is good'. Therefore, 'tsukaereba ii' has the same meaning as 'as long as it can be used' or 'if it can be used, it's fine'.

\newpage

\section*{578. too much / excessive}

\textbf{Date}: 2025.11.09 (Sun)\\
\textbf{Index}: sugi\\
\textbf{Expression (JA)}: すぎ\\
\textbf{Romanization}: sugi\\
\textbf{Expression (EN)}: too much / excessive

\textbf{Variants (JA)}: 過剰な / 過度な / やりすぎの\\
\textbf{Variants (EN)}: excessive / over the top / too much

\textbf{Intent (JA)}: 評価, 注意喚起, 感情表現\\
\textbf{Intent (EN)}: evaluation, warning, emotional expression

\textbf{Tags (JA)}: 即時文法, 口語, 評価, 注意喚起\\
\textbf{Tags (EN)}: immediate grammar, spoken, evaluation, warning

\textbf{Adjusted Expression (JA)}:\\
(過剰:かじょう)な(状態:じょうたい)\\
\textbf{Adjusted Expression (EN)}:\\
excessive state

\textbf{Example (JA)}:\\
A「(彼:かれ)、(仕事:しごと)し(過:す)ぎだよ」 B「うん、やりすぎ」\\
\textbf{Example (EN)}:\\
A: He's working too much. B: Yeah, it's excessive.

\textbf{Notes (JA)}:\\
「すぎ」は、物事が過剰であることや、限度を超えていることを即時的に表現する口語的な表現。もう少し詳しくこのことばの構成を説明すると、「すぎ」は「too much」や「excessive」の意味で、物事が適切な範囲を超えている様子を指します。

\textbf{Notes (EN)}:\\
'Sugi' is a colloquial expression used to immediately express that something is excessive or beyond limits. To explain the structure of this phrase in more detail, 'sugi' means 'too much' or 'excessive', referring to a state where something exceeds an appropriate range.

\newpage

\section*{577. messy / chaotic}

\textbf{Date}: 2025.11.08 (Sat)\\
\textbf{Index}: mechakucha\\
\textbf{Expression (JA)}: めちゃくちゃ\\
\textbf{Romanization}: mechakucha\\
\textbf{Expression (EN)}: messy / chaotic

\textbf{Variants (JA)}: 混乱した / 無秩序な / めちゃめちゃな\\
\textbf{Variants (EN)}: confused / disorderly / chaotic

\textbf{Intent (JA)}: 状態の描写, 評価, 感情表現\\
\textbf{Intent (EN)}: state description, evaluation, emotional expression

\textbf{Tags (JA)}: 即時文法, 口語, 状態描写, 評価\\
\textbf{Tags (EN)}: immediate grammar, spoken, state description, evaluation

\textbf{Adjusted Expression (JA)}:\\
(混乱:こんらん)している\\
\textbf{Adjusted Expression (EN)}:\\
being confused

\textbf{Example (JA)}:\\
A「(部屋:へや)、めちゃくちゃだね」 B「うん、(片付:かたづ)けないと」\\
\textbf{Example (EN)}:\\
A: Your room is messy, isn't it? B: Yeah, I need to tidy up.

\textbf{Notes (JA)}:\\
「めちゃくちゃ」は、物が散らかっている状態や、物事が整理されていない様子を即時的に表現する口語的な表現。もう少し詳しくこのことばの構成を説明すると、「めちゃくちゃ」は「messy」や「chaotic」の意味で、物が乱雑に置かれている様子や、状況が混乱している様子を指します。

\textbf{Notes (EN)}:\\
'Mechakucha' is a colloquial expression used to immediately describe a state where things are scattered or situations are disorganized. To explain the structure of this phrase in more detail, 'mechakucha' means 'messy' or 'chaotic', referring to a state where things are placed haphazardly or situations are chaotic.

\newpage

\section*{576. somehow / by some means}

\textbf{Date}: 2025.11.07 (Fri)\\
\textbf{Index}: tokanantoka\\
\textbf{Expression (JA)}: とかなんとか\\
\textbf{Romanization}: toka nantoka\\
\textbf{Expression (EN)}: somehow / by some means

\textbf{Variants (JA)}: どうにかして / 何とかして / なんとかなる\\
\textbf{Variants (EN)}: by some means / somehow / it will work out

\textbf{Intent (JA)}: 説明, 曖昧な表現, 不確定な状況\\
\textbf{Intent (EN)}: explanation, ambiguous expression, uncertain situation

\textbf{Tags (JA)}: 即時文法, 口語, 説明, 曖昧な表現\\
\textbf{Tags (EN)}: immediate grammar, spoken, explanation, ambiguous expression

\textbf{Adjusted Expression (JA)}:\\
などということでした。\\
\textbf{Adjusted Expression (EN)}:\\
It was something like that.

\textbf{Example (JA)}:\\
A「(昨日:きのう)の(会議:かいぎ)、どうだった？」 B「うん、(消費税:しょうひぜい)とかなんとかっていう(話:はなし)だったよ」\\
\textbf{Example (EN)}:\\
A: How was the meeting yesterday? B: Yeah, it was something about consumption tax or something.

\textbf{Notes (JA)}:\\
「とかなんとか」は、具体的な内容を曖昧に表現したり、不確定な状況を示したりする際に使う口語的な表現。もう少し詳しくこのことばの構成を説明すると、「とか」は「or something like that」の意味、「なんとか」は「somehow」や「by some means」の意味です。つまり、「とかなんとか」は「somehow」や「by some means」と同じ意味になります。

\textbf{Notes (EN)}:\\
'Toka nantoka' is a colloquial expression used to vaguely express specific content or indicate uncertain situations. To explain the structure of this phrase in more detail, 'toka' means 'or something like that', and 'nantoka' means 'somehow' or 'by some means'. Therefore, 'toka nantoka' has the same meaning as 'somehow' or 'by some means'.

\newpage

\section*{575. to move around / to circulate}

\textbf{Date}: 2025.11.06 (Thu)\\
\textbf{Index}: ugokimawatte\\
\textbf{Expression (JA)}: (動:うご)き(回:まわ)って\\
\textbf{Romanization}: ugoki mawatte\\
\textbf{Expression (EN)}: to move around / to circulate

\textbf{Variants (JA)}: 巡回する / 移動する / 循環する\\
\textbf{Variants (EN)}: to patrol / to move / to circulate

\textbf{Intent (JA)}: 動作の描写, 状態の説明, 行動の指示\\
\textbf{Intent (EN)}: action description, state explanation, action instruction

\textbf{Tags (JA)}: 即時文法, 口語, 動作描写, 状態説明\\
\textbf{Tags (EN)}: immediate grammar, spoken, action description, state explanation

\textbf{Adjusted Expression (JA)}:\\
活発に動いている\\
\textbf{Adjusted Expression (EN)}:\\
being actively in motion

\textbf{Example (JA)}:\\
A「ねえ、(パンダ:ぱんだ)、どうだった？」 B「うん、(動:うご)き(回:まわ)ってて(可愛:かわい)かったよ」\\
\textbf{Example (EN)}:\\
A: Hey, how was the panda? B: Yeah, it was moving around and cute.

\textbf{Notes (JA)}:\\
「動き回って」は、物や人が活発に移動したり、巡回したりする様子を即時的に表現する口語的な表現。もう少し詳しくこのことばの構成を説明すると、「動き」は「to move」の意味、「回って」は「to go around」や「to circulate」の意味です。つまり、「動き回って」は「to move around」や「to circulate」と同じ意味になります。調整後の表現として、「活発に動いている」を紹介しましたが、普通に「動き回っている」と言うこともできます。

\textbf{Notes (EN)}:\\
'Ugoki mawatte' is a colloquial expression used to immediately describe the state of objects or people moving actively or patrolling. To explain the structure of this phrase in more detail, 'ugoki' means 'to move', and 'mawatte' means 'to go around' or 'to circulate'. Therefore, 'ugoki mawatte' has the same meaning as 'to move around' or 'to circulate'. As an adjusted expression, 'being actively in motion' was introduced, but it is also common to simply say 'moving around'.

\newpage

\section*{574. to be irritated / to be edgy}

\textbf{Date}: 2025.11.05 (Wed)\\
\textbf{Index}: piripirishiteru\\
\textbf{Expression (JA)}: (ピリピリ:ぴりぴり)してる\\
\textbf{Romanization}: piripiri shiteru\\
\textbf{Expression (EN)}: to be irritated / to be edgy

\textbf{Variants (JA)}: 緊張している / イライラしている / 神経質になっている\\
\textbf{Variants (EN)}: to be tense / to be irritated / to be nervous

\textbf{Intent (JA)}: 感情表現, 緊張, 不安\\
\textbf{Intent (EN)}: emotional expression, tension, anxiety

\textbf{Tags (JA)}: 即時文法, 口語, 感情表現, 緊張\\
\textbf{Tags (EN)}: immediate grammar, spoken, emotional expression, tension

\textbf{Adjusted Expression (JA)}:\\
(緊張:きんちょう)している\\
\textbf{Adjusted Expression (EN)}:\\
to be tense

\textbf{Example (JA)}:\\
A「(試験:しけん)の(前:まえ)で(ピリピリ:ぴりぴり)してるね」 B「うん、(緊張:きんちょう)して(落:お)ち(着:つ)かないよ」\\
\textbf{Example (EN)}:\\
A: You're edgy before the exam, aren't you? B: Yeah, I'm nervous and can't settle down.

\textbf{Notes (JA)}:\\
「ピリピリしてる」は、緊張や不安などの感情から落ち着かない状態を即時的に表現する口語的な表現。もう少し詳しくこのことばの構成を説明すると、「ピリピリ」は「irritated」や「edgy」の意味で、心が落ち着かない様子を指します。

\textbf{Notes (EN)}:\\
'Piripiri shiteru' is a colloquial expression used to immediately express a state of restlessness due to emotions such as nervousness or anxiety. To explain the structure of this phrase in more detail, 'piripiri' means 'irritated' or 'edgy', referring to a state of unease.

\newpage

\section*{573. to be irritated / to be annoyed}

\textbf{Date}: 2025.11.04 (Tue)\\
\textbf{Index}: irairasuru\\
\textbf{Expression (JA)}: (イライラ:いらいら)する\\
\textbf{Romanization}: iraira suru\\
\textbf{Expression (EN)}: to be irritated / to be annoyed

\textbf{Variants (JA)}: 苛立つ / 腹立たしい / 不快に感じる\\
\textbf{Variants (EN)}: to be frustrated / to be angry / to feel uncomfortable

\textbf{Intent (JA)}: 感情表現, 不満, 怒り\\
\textbf{Intent (EN)}: emotional expression, dissatisfaction, anger

\textbf{Tags (JA)}: 即時文法, 口語, 感情表現, 不満\\
\textbf{Tags (EN)}: immediate grammar, spoken, emotional expression, dissatisfaction

\textbf{Adjusted Expression (JA)}:\\
(苛立:いらだ)つ\\
\textbf{Adjusted Expression (EN)}:\\
to be frustrated

\textbf{Example (JA)}:\\
A「なかなか(連絡:れんらく)が(来:こ)ないね」 B「うん、(イライラ:いらいら)する」\\
\textbf{Example (EN)}:\\
A: They still haven't contacted us. B: Yeah, I'm getting irritated.

\textbf{Notes (JA)}:\\
「イライラする」は、何かに対して不満や怒りを感じることを即時的に表現する口語的な表現。もう少し詳しくこのことばの構成を説明すると、「イライラ」は「irritated」や「annoyed」の意味で、心が落ち着かない様子を指します。

\textbf{Notes (EN)}:\\
'Iraira suru' is a colloquial expression used to immediately express feelings of dissatisfaction or anger towards something. To explain the structure of this phrase in more detail, 'iraira' means 'irritated' or 'annoyed', referring to a state of unease.

\newpage

\section*{572. waste / futility}

\textbf{Date}: 2025.11.03 (Mon)\\
\textbf{Index}: muda\\
\textbf{Expression (JA)}: (無駄:むだ)\\
\textbf{Romanization}: muda\\
\textbf{Expression (EN)}: waste / futility

\textbf{Variants (JA)}: 無意味 / 無益 / 徒労\\
\textbf{Variants (EN)}: meaningless / pointless / futile

\textbf{Intent (JA)}: 評価, 否定, 感情表現\\
\textbf{Intent (EN)}: evaluation, negation, emotional expression

\textbf{Tags (JA)}: 即時文法, 口語, 評価, 否定\\
\textbf{Tags (EN)}: immediate grammar, spoken, evaluation, negation

\textbf{Adjusted Expression (JA)}:\\
(無意味:むいみ)\\
\textbf{Adjusted Expression (EN)}:\\
meaningless

\textbf{Example (JA)}:\\
A「この(仕事:しごと)、(無駄:むだ)じゃない？」 B「うん、(時間:じかん)の(無駄:むだ)だよね」\\
\textbf{Example (EN)}:\\
A: Isn't this work a waste? B: Yeah, it's a waste of time.

\textbf{Notes (JA)}:\\
「無駄」は、労力や時間を費やしても成果が得られないことを即時的に表現する口語的な表現。もう少し詳しくこのことばの構成を説明すると、「無駄」は「waste」や「futility」の意味で、価値のないことや無意味なことを指します。

\textbf{Notes (EN)}:\\
'Muda' is a colloquial expression used to immediately express that effort or time spent does not yield results. To explain the structure of this phrase in more detail, 'muda' means 'waste' or 'futility', referring to something that is worthless or meaningless.

\newpage

\section*{571. to be frustrated / to be impatient}

\textbf{Date}: 2025.11.02 (Sun)\\
\textbf{Index}: yakimokisuru\\
\textbf{Expression (JA)}: やきもきする\\
\textbf{Romanization}: yakimoki suru\\
\textbf{Expression (EN)}: to be frustrated / to be impatient

\textbf{Variants (JA)}: いらいらする / もどかしい / じれったい\\
\textbf{Variants (EN)}: to be irritated / to be vexed / to be anxious

\textbf{Intent (JA)}: 感情表現, 不安, 焦り\\
\textbf{Intent (EN)}: emotional expression, anxiety, impatience

\textbf{Tags (JA)}: 即時文法, 口語, 感情表現, 不安\\
\textbf{Tags (EN)}: immediate grammar, spoken, emotional expression, anxiety

\textbf{Adjusted Expression (JA)}:\\
(苛立:いらだ)つ\\
\textbf{Adjusted Expression (EN)}:\\
to be irritated

\textbf{Example (JA)}:\\
A「(連絡:れんらく)まだ(来:こ)ないね」 B「うん、やきもきする」\\
\textbf{Example (EN)}:\\
A: They still haven't contacted us. B: Yeah, I'm getting impatient.

\textbf{Notes (JA)}:\\
「やきもきする」は、待っている間に不安や焦りを感じることを即時的に表現する口語的な表現。もう少し詳しくこのことばの構成を説明すると、「やきもき」は「frustrated」や「impatient」の意味で、心が落ち着かない様子を指します。

\textbf{Notes (EN)}:\\
'Yakimoki suru' is a colloquial expression used to immediately express feelings of anxiety or impatience while waiting. To explain the structure of this phrase in more detail, 'yakimoki' means 'frustrated' or 'impatient', referring to a state of unease.

\newpage

\section*{570. messy / disorganized}

\textbf{Date}: 2025.11.01 (Sat)\\
\textbf{Index}: gochagocha\\
\textbf{Expression (JA)}: ごちゃごちゃ\\
\textbf{Romanization}: gocha gocha\\
\textbf{Expression (EN)}: messy / disorganized

\textbf{Variants (JA)}: 散らかっている / 乱雑な / 混乱した\\
\textbf{Variants (EN)}: cluttered / disorderly / chaotic

\textbf{Intent (JA)}: 状態の描写, 評価, 感情表現\\
\textbf{Intent (EN)}: state description, evaluation, emotional expression

\textbf{Tags (JA)}: 即時文法, 口語, 状態描写, 評価\\
\textbf{Tags (EN)}: immediate grammar, spoken, state description, evaluation

\textbf{Adjusted Expression (JA)}:\\
(散:ちら)かっている\\
\textbf{Adjusted Expression (EN)}:\\
cluttered

\textbf{Example (JA)}:\\
A「(部屋:へや)、ごちゃごちゃしてるね」 B「うん、(片付:かたづ)けないと」\\
\textbf{Example (EN)}:\\
A: Your room is messy, isn't it? B: Yeah, I need to tidy up.

\textbf{Notes (JA)}:\\
「ごちゃごちゃ」は、物が散らかっている状態や、物事が整理されていない様子を即時的に表現する口語的な表現。もう少し詳しくこのことばの構成を説明すると、「ごちゃごちゃ」は「messy」や「disorganized」の意味で、物が乱雑に置かれている様子や、状況が混乱している様子を指します。

\textbf{Notes (EN)}:\\
'Gocha gocha' is a colloquial expression used to immediately describe a state where things are scattered or situations are disorganized. To explain the structure of this phrase in more detail, 'gocha gocha' means 'messy' or 'disorganized', referring to a state where things are placed haphazardly or situations are chaotic.

\newpage

\section*{569. detailed things / small details}

\textbf{Date}: 2025.10.31 (Fri)\\
\textbf{Index}: komakaikoto\\
\textbf{Expression (JA)}: (細:こま)かいこと\\
\textbf{Romanization}: komakai koto\\
\textbf{Expression (EN)}: detailed things / small details

\textbf{Variants (JA)}: 細部のこと / 重箱の隅 / 小さなこと\\
\textbf{Variants (EN)}: detailed matters / specific details / minor things

\textbf{Intent (JA)}: 説明, 情報提供, 注意喚起\\
\textbf{Intent (EN)}: explanation, information provision, attention

\textbf{Tags (JA)}: 即時文法, 口語, 説明, 情報提供\\
\textbf{Tags (EN)}: immediate grammar, spoken, explanation, information provision

\textbf{Adjusted Expression (JA)}:\\
(些細:ささい)なこと\\
\textbf{Adjusted Expression (EN)}:\\
minor things

\textbf{Example (JA)}:\\
A「(髪:かみ)、ちゃんと(結:むす)んで(行:い)きなさい」 B「そんな(細:こま)かいこと」\\
\textbf{Example (EN)}:\\
A: Make sure to tie your hair properly. B: Such a detailed thing.

\textbf{Notes (JA)}:\\
「細かいこと」は、詳細な情報や小さな事柄を指す口語的な表現。もう少し詳しくこのことばの構成を説明すると、「細かい」は「detailed」や「small」の意味、「こと」は「thing」や「matter」の意味です。つまり、「細かいこと」は「detailed things」や「small details」と同じ意味になりますが、だいたいは「どうでもいいこと」という意味で使われます。

\textbf{Notes (EN)}:\\
'Komakai koto' is a colloquial expression referring to detailed information or small matters. To explain the structure of this phrase in more detail, 'komakai' means 'detailed' or 'small', and 'koto' means 'thing' or 'matter'. Therefore, 'komakai koto' has the same meaning as 'detailed things' or 'small details', but it is generally used to mean 'trivial matters'.

\newpage

\section*{568. to be restless / to be fidgety}

\textbf{Date}: 2025.10.30 (Thu)\\
\textbf{Index}: sowasowasuru\\
\textbf{Expression (JA)}: そわそわする\\
\textbf{Romanization}: sowasowa suru\\
\textbf{Expression (EN)}: to be restless / to be fidgety

\textbf{Variants (JA)}: 落ち着かない / そわそわしている / 不安定な気持ちになる\\
\textbf{Variants (EN)}: to be uneasy / to be fidgety / to feel unsettled

\textbf{Intent (JA)}: 感情表現, 不安, 緊張\\
\textbf{Intent (EN)}: emotional expression, anxiety, nervousness

\textbf{Tags (JA)}: 即時文法, 口語, 感情表現, 不安\\
\textbf{Tags (EN)}: immediate grammar, spoken, emotional expression, anxiety

\textbf{Adjusted Expression (JA)}:\\
(落:お)ち(着:つ)かない\\
\textbf{Adjusted Expression (EN)}:\\
to be uneasy

\textbf{Example (JA)}:\\
A「(試験:しけん)の(前:まえ)で(そわそわ)してるね」 B「うん、(緊張:きんちょう)して(落:お)ち(着:つ)かないよ」\\
\textbf{Example (EN)}:\\
A: You're fidgety before the exam, aren't you? B: Yeah, I'm nervous and can't settle down.

\textbf{Notes (JA)}:\\
「そわそわする」は、緊張や不安などの感情から落ち着かない状態を即時的に表現する口語的な表現。もう少し詳しくこのことばの構成を説明すると、「そわそわ」は「restless」や「fidgety」の意味で、心が落ち着かない様子を指します。

\textbf{Notes (EN)}:\\
'Sowasowa suru' is a colloquial expression used to immediately express a state of restlessness due to emotions such as nervousness or anxiety. To explain the structure of this phrase in more detail, 'sowasowa' means 'restless' or 'fidgety', referring to a state of unease.

\newpage

\section*{567. have to / must / need to}

\textbf{Date}: 2025.10.29 (Wed)\\
\textbf{Index}: nakuccha\\
\textbf{Expression (JA)}: しなくっちゃ\\
\textbf{Romanization}: nakuccha\\
\textbf{Expression (EN)}: have to / must / need to

\textbf{Variants (JA)}: しなければならない / しないといけない / しなくてはいけない\\
\textbf{Variants (EN)}: have to / must / need to

\textbf{Intent (JA)}: 義務, 必要性, 強制\\
\textbf{Intent (EN)}: obligation, necessity, compulsion

\textbf{Tags (JA)}: 即時文法, 口語, 義務, 必要性\\
\textbf{Tags (EN)}: immediate grammar, spoken, obligation, necessity

\textbf{Adjusted Expression (JA)}:\\
しなければならない\\
\textbf{Adjusted Expression (EN)}:\\
have to

\textbf{Example (JA)}:\\
A「(宿題:しゅくだい)は？」 B「うん、(今:いま)からしなくっちゃ」\\
\textbf{Example (EN)}:\\
A: What about your homework? B: Yeah, I have to do it now.

\textbf{Notes (JA)}:\\
「しなくっちゃ」は、「しなければならない」や「しないといけない」と同じ意味で、義務や必要性を即時的に表現する口語的な表現。もう少し詳しくこのことばの構成を説明すると、「しなくっちゃ」は「しなければ」の口語縮約形で、「to have to do」や「must do」の意味、「ちゃ」は「ては」の口語縮約形で、その後ろに「ならない」が省略されている形です。つまり、「しなくっちゃ」は「have to」や「must」と同じ意味になります。

\textbf{Notes (EN)}:\\
'Shinakuccha' has the same meaning as 'shinakereba naranai' or 'shinai to ikenai', used to immediately express obligation or necessity in a colloquial manner. To explain the structure of this phrase in more detail, 'shinakuccha' is a colloquial contraction of 'shinakereba', meaning 'to have to do' or 'must do', and 'cha' is a colloquial contraction of 'te wa', with 'naranai' omitted afterward. Therefore, 'shinakuccha' has the same meaning as 'have to' or 'must'.

\newpage

\section*{566. It's going well / It's working out}

\textbf{Date}: 2025.10.28 (Tue)\\
\textbf{Index}: umakuitteru\\
\textbf{Expression (JA)}: うまくいってる\\
\textbf{Romanization}: umaku itteru\\
\textbf{Expression (EN)}: It's going well / It's working out

\textbf{Variants (JA)}: 順調に進んでいる / うまく進んでいる / 問題なく進んでいる\\
\textbf{Variants (EN)}: It's progressing smoothly / It's going smoothly / It's progressing without problems

\textbf{Intent (JA)}: 状況報告, 評価, 安心感の表現\\
\textbf{Intent (EN)}: status report, evaluation, expression of reassurance

\textbf{Tags (JA)}: 即時文法, 口語, 状況報告, 評価\\
\textbf{Tags (EN)}: immediate grammar, spoken, status report, evaluation

\textbf{Adjusted Expression (JA)}:\\
(順調:じゅんちょう)に(進:すす)んでいる\\
\textbf{Adjusted Expression (EN)}:\\
It's progressing smoothly

\textbf{Example (JA)}:\\
A「(新企画:しんきかく)はどう？」 B「うまくいってるよ、(予定:よてい)通りに(進:すす)んでる」\\
\textbf{Example (EN)}:\\
A: How's the new project going? B: It's going well, progressing as planned.

\textbf{Notes (JA)}:\\
「うまくいってる」は、物事が順調に進んでいることを即時的に表現する口語的な表現。もう少し詳しくこのことばの構成を説明すると、「うまく」は「well」や「smoothly」の意味、「いってる」は「行っている」の口語縮約形で、「to go」や「to progress」の意味です。つまり、「うまくいってる」は「It's going well」や「It's working out」と同じ意味になります。

\textbf{Notes (EN)}:\\
'Umaku itteru' is a colloquial expression used to immediately express that things are progressing smoothly. To explain the structure of this phrase in more detail, 'umaku' means 'well' or 'smoothly', and 'itteru' is a colloquial contraction of 'itte iru', meaning 'to go' or 'to progress'. Therefore, 'umaku itteru' has the same meaning as 'It's going well' or 'It's working out'.

\newpage

\section*{565. Where should we start?}

\textbf{Date}: 2025.10.27 (Mon)\\
\textbf{Index}: dokokaraatarimasuka\\
\textbf{Expression (JA)}: どこから(当:あ)たりますか？\\
\textbf{Romanization}: doko kara atarimasu ka?\\
\textbf{Expression (EN)}: Where should we start?

\textbf{Variants (JA)}: どこから始めますか？ / どこから取り掛かりますか？ / どこから手を付けますか？\\
\textbf{Variants (EN)}: Where do we begin? / Where do we get started? / Where do we tackle first?

\textbf{Intent (JA)}: 確認, 計画立案, 指示\\
\textbf{Intent (EN)}: confirmation, planning, instruction

\textbf{Tags (JA)}: 即時文法, 口語, 確認, 計画立案\\
\textbf{Tags (EN)}: immediate grammar, spoken, confirmation, planning

\textbf{Adjusted Expression (JA)}:\\
どこから(始:はじ)めますか？\\
\textbf{Adjusted Expression (EN)}:\\
Where do we begin?

\textbf{Example (JA)}:\\
A「(仕事:しごと)を(始:はじ)めるにあたって、どこから(当:あ)たりますか？」 B「まずは(資料:しりょう)の(確認:かくにん)から」\\
\textbf{Example (EN)}:\\
A: Where should we start when beginning the work? B: Let's start by reviewing the materials.

\textbf{Notes (JA)}:\\
「どこから当たりますか？」は、仕事やプロジェクトを始める際に、最初に取り組むべき場所やタスクを確認するために使う口語的な表現。もう少し詳しくこのことばの構成を説明すると、「どこから」は「where from」の意味、「当たりますか？」は「to hit」や「to tackle」の意味です。つまり、「どこから当たりますか？」は「Where should we start?」と同じ意味になります。

\textbf{Notes (EN)}:\\
'Doko kara atarimasu ka?' is a colloquial expression used to confirm where to start or which task to tackle first when beginning work or a project. To explain the structure of this phrase in more detail, 'doko kara' means 'where from', and 'atarimasu ka?' means 'to hit' or 'to tackle'. Therefore, 'doko kara atarimasu ka?' has the same meaning as 'Where should we start?'.

\newpage

\section*{564. I mean, ... / You know, ...}

\textbf{Date}: 2025.10.26 (Sun)\\
\textbf{Index}: tteiunomo,...\\
\textbf{Expression (JA)}: っていうのも、...\\
\textbf{Romanization}: tte iunomo,...\\
\textbf{Expression (EN)}: I mean, ... / You know, ...

\textbf{Variants (JA)}: つまり、... / 要するに、... / なんていうか、...\\
\textbf{Variants (EN)}: In other words, ... / To sum up, ... / How should I put it, ...

\textbf{Intent (JA)}: 説明, 補足, 話題転換\\
\textbf{Intent (EN)}: explanation, supplement, topic shift

\textbf{Tags (JA)}: 即時文法, 口語, 説明, 補足\\
\textbf{Tags (EN)}: immediate grammar, spoken, explanation, supplement

\textbf{Adjusted Expression (JA)}:\\
つまり、...\\
\textbf{Adjusted Expression (EN)}:\\
In other words, ...

\textbf{Example (JA)}:\\
A「(今日:きょう)はもう(帰:かえ)るよ、っていうのも、(彼女:かのじょ)の(誕生日:たんじょうび)なんだ」\\
\textbf{Example (EN)}:\\
A: I'm going home now, you know, it's my girlfriend's birthday.

\textbf{Notes (JA)}:\\
「っていうのも、...」は、補足説明や理由を提供する際に使う口語的な表現。もっと本音を漏らすときに使われることにも使われる。

\textbf{Notes (EN)}:\\
'Tte iunomo, ...' is a colloquial expression used to provide additional explanation or reasons. It is also used when revealing more honest feelings.

\newpage

\section*{563. Not sure why, ...}

\textbf{Date}: 2025.10.25 (Sat)\\
\textbf{Index}: doushitahazumika\\
\textbf{Expression (JA)}: どうしたはずみか、...\\
\textbf{Romanization}: doushita hazumi ka, \\
\textbf{Expression (EN)}: Not sure why, 

\textbf{Variants (JA)}: なぜか知らないけど、... / 理由はわからないけど、... / なんとなくだけど、...\\
\textbf{Variants (EN)}: I don't know why, ... / I don't know the reason, ... / For some reason, ...

\textbf{Intent (JA)}: 説明, 理由付け, 感情表現\\
\textbf{Intent (EN)}: explanation, justification, emotional expression

\textbf{Tags (JA)}: 即時文法, 口語, 説明, 理由付け\\
\textbf{Tags (EN)}: immediate grammar, spoken, explanation, justification

\textbf{Adjusted Expression (JA)}:\\
(明確:めいかく)な(理由:りゆう)は(不明:ふめい)だが、...\\
\textbf{Adjusted Expression (EN)}:\\
Although the exact reason is unclear, ...

\textbf{Example (JA)}:\\
どうしたはずみか、(休:やす)みなのに(大学:だいがく)に(来:き)てしまった。\\
\textbf{Example (EN)}:\\
Not sure why, but I ended up coming to the university on a holiday.

\textbf{Notes (JA)}:\\
「どうしたはずみか、...」は、明確な理由がわからないが、ある出来事や行動が起こったことを即時的に表現する口語的な表現。もう少し詳しくこのことばの構成を説明すると、「どうした」は「what happened」の意味、「はずみ」は「impetus」や「trigger」の意味、「か」は疑問を示す助詞です。

\textbf{Notes (EN)}:\\
'Doushita hazumi ka, ...' is a colloquial expression used to immediately express that an event or action occurred without a clear reason. To explain the structure of this phrase in more detail, 'doushita' means 'what happened', 'hazumi' means 'impetus' or 'trigger', and 'ka' is a particle indicating a question.

\newpage

\section*{562. I feel sick / I feel nauseous}

\textbf{Date}: 2025.10.24 (Fri)\\
\textbf{Index}: unzarisuru\\
\textbf{Expression (JA)}: うんざりする\\
\textbf{Romanization}: unzari suru\\
\textbf{Expression (EN)}: I feel sick / I feel nauseous

\textbf{Variants (JA)}: 飽き飽きする / 嫌気がさす / うんざりだ\\
\textbf{Variants (EN)}: I'm fed up / I'm sick of it / I'm tired of it

\textbf{Intent (JA)}: 感情表現, 不満, 嫌悪\\
\textbf{Intent (EN)}: emotional expression, dissatisfaction, disgust

\textbf{Tags (JA)}: 即時文法, 口語, 感情表現, 不満\\
\textbf{Tags (EN)}: immediate grammar, spoken, emotional expression, dissatisfaction

\textbf{Adjusted Expression (JA)}:\\
(嫌気:いやけ)がさす\\
\textbf{Adjusted Expression (EN)}:\\
I'm sick of it

\textbf{Example (JA)}:\\
A「(彼:かれ)の(話:はなし)、またあの(話:はなし)？」 B「うん、(本当:ほんとう)にうんざりするよ」\\
\textbf{Example (EN)}:\\
A: His story, that story again? B: Yeah, I'm really sick of it.

\textbf{Notes (JA)}:\\
「うんざりする」は、同じことが繰り返されたり、嫌なことが続いたりして、強い不満や嫌悪感を感じることを即時的に表現する口語的な表現。もう少し詳しくこのことばの構成を説明すると、「うんざり」は「飽き飽きする」や「嫌気がさす」と同じ意味で、物事に対して強い嫌悪感や不満を感じる状態を指します。

\textbf{Notes (EN)}:\\
'Unzari suru' is a colloquial expression used to immediately express strong dissatisfaction or disgust when the same thing is repeated or unpleasant things continue. To explain the structure of this phrase in more detail, 'unzari' has the same meaning as 'akiaki suru' or 'iyake ga sasu', referring to a state of feeling strong disgust or dissatisfaction towards something.

\newpage

\section*{561. What am I supposed to do?}

\textbf{Date}: 2025.10.23 (Thu)\\
\textbf{Index}: naniyannakuchanaranaindakke\\
\textbf{Expression (JA)}: (何:なに)やんなくちゃならないんだっけ\\
\textbf{Romanization}: nani yannakucha naranain dakke\\
\textbf{Expression (EN)}: What am I supposed to do?

\textbf{Variants (JA)}: 何をしなくちゃならないんだっけ / 何をやらなくちゃならないんだっけ / 何をしなければならないんだっけ\\
\textbf{Variants (EN)}: What do I have to do? / What must I do? / What am I supposed to do?

\textbf{Intent (JA)}: 確認, 思い出し, 自己管理\\
\textbf{Intent (EN)}: confirmation, recall, self-management

\textbf{Tags (JA)}: 即時文法, 口語, 確認, 思い出し\\
\textbf{Tags (EN)}: immediate grammar, spoken, confirmation, recall

\textbf{Adjusted Expression (JA)}:\\
(何:なに)をすべきだったのか、(忘:わす)れてしまった\\
\textbf{Adjusted Expression (EN)}:\\
I forgot what I was supposed to do

\textbf{Example (JA)}:\\
A「(明日:あした)の(会議:かいぎ)の(準備:じゅんび)、(何:なに)やんなくちゃならないんだっけ？」 B「(資料:しりょう)の(印刷:いんさつ)とだけじゃない？」\\
\textbf{Example (EN)}:\\
A: What am I supposed to do to prepare for tomorrow's meeting? B: Isn't it just printing the materials?

\textbf{Notes (JA)}:\\
「何やんなくちゃならないんだっけ」は、「何をしなければならないのか忘れてしまった」という意味で、話し手が自分の義務やタスクを思い出そうとしていることを示す口語的な表現。もう少し詳しくこのことばの構成を説明すると、「何(なに)」は「what」、「やんなくちゃ」は「やらなくては」の口語縮約形で、「しなければならない」の意味、「ならないんだっけ」は「ならないのか忘れてしまった」という意味です。つまり、「何やんなくちゃならないんだっけ」は「何をしなければならないのか忘れてしまった」という意味になります。独り言としてよく使われます。

\textbf{Notes (EN)}:\\
'Nani yannakucha naranain dakke' means 'I forgot what I was supposed to do', indicating that the speaker is trying to recall their duties or tasks. To explain the structure of this phrase in more detail, 'nani' means 'what', 'yannakucha' is a colloquial contraction of 'yaranakute wa', meaning 'must do', and 'naranain dakke' means 'I forgot what I was supposed to do'. Therefore, 'nani yannakucha naranain dakke' means 'I forgot what I was supposed to do'. It is often used as a self-talk.

\newpage

\section*{560. Just a little / Just a bit}

\textbf{Date}: 2025.10.22 (Wed)\\
\textbf{Index}: honnnochotto\\
\textbf{Expression (JA)}: ほんのちょっと\\
\textbf{Romanization}: honn no chotto\\
\textbf{Expression (EN)}: Just a little / Just a bit

\textbf{Variants (JA)}: 少しだけ / わずかに / ちょっぴり\\
\textbf{Variants (EN)}: A little bit / Slightly / Just a tiny bit

\textbf{Intent (JA)}: 数量の表現, 控えめな表現, 強調\\
\textbf{Intent (EN)}: quantity expression, modest expression, emphasis

\textbf{Tags (JA)}: 即時文法, 口語, 数量表現, 控えめな表現\\
\textbf{Tags (EN)}: immediate grammar, spoken, quantity expression, modest expression

\textbf{Adjusted Expression (JA)}:\\
(少:すこ)しだけ\\
\textbf{Adjusted Expression (EN)}:\\
A little bit

\textbf{Example (JA)}:\\
A「(砂糖:さとう)、(使:つか)う？」 B「ほんのちょっと」\\
\textbf{Example (EN)}:\\
A: Do you want some sugar? B: Just a little.

\textbf{Notes (JA)}:\\
「ほんのちょっと」は、数量を控えめに表現する際に使う口語的な表現。質量にも使えるし、時間にも使えるし、距離にも使える。非常に少ないことを強調するニュアンスがある。また、心理的に違和感があるときにも使われることがある。例えば、「ほんのちょっと違う」と言うと、少し違うだけでなく、心理的に受け入れがたいニュアンスが含まれることがある。

\textbf{Notes (EN)}:\\
'Honn no chotto' is a colloquial expression used to modestly express quantity. It can be used for mass, time, and distance. It carries a nuance of emphasizing that something is very small. Additionally, it can also be used when there is a psychological discomfort. For example, saying 'honn no chotto chigau' (just a little different) may imply not only a slight difference but also a nuance of being psychologically unacceptable.

\newpage

\section*{559. Not as I think / Not as expected}

\textbf{Date}: 2025.10.21 (Tue)\\
\textbf{Index}: omouyounihane\\
\textbf{Expression (JA)}: (思:おも)うように(は)ね\\
\textbf{Romanization}: omou you ni hane\\
\textbf{Expression (EN)}: Not as I think / Not as expected

\textbf{Variants (JA)}: 予想通りにはね / 期待通りにはね / 考え通りにはね\\
\textbf{Variants (EN)}: Not as I expected / Not as I anticipated / Not as I considered

\textbf{Intent (JA)}: 比較, 評価, 感想\\
\textbf{Intent (EN)}: comparison, evaluation, impression

\textbf{Tags (JA)}: 即時文法, 口語, 比較, 評価\\
\textbf{Tags (EN)}: immediate grammar, spoken, comparison, evaluation

\textbf{Adjusted Expression (JA)}:\\
(予想:よそう)通りにはね\\
\textbf{Adjusted Expression (EN)}:\\
Not as I expected

\textbf{Example (JA)}:\\
A「(順調:じゅんちょう)に(進:すす)んでる？」 B「うーん、(思:おも)うように(は)ね」\\
\textbf{Example (EN)}:\\
A: Is it going smoothly? B: Hmm, not as I thought.

\textbf{Notes (JA)}:\\
「思うようにはね」は、予想や期待と比較して実際の結果や状況が異なるという意味では、「思うようにいかない」と同じ。日本語の助詞で終わる文は、話し手の感情やニュアンスを強調することが多い。英語では"It's not going as I thought"や"Nothing is going right"などと表現されることが多い。また「ね」「よ」「さ」の前も助詞である場合は、その文が完結していない場合であるので、何かの余韻や続きがあることを示唆している。このように、助詞で終わる文は、文であることには間違いないが、不完全な文であり、話し手の感情やニュアンスを強調する役割を果たしている。一方、完全な文は、述語で終わる文であり、簡単に言えば、句点「。」で終わる文である。こういう文は記録や公式な言及に使われることが多く、私的な感情を話すように表現する形式ではない。

\textbf{Notes (EN)}:\\
'Omou you ni hane' has the same meaning as 'omou you ni ikanai', indicating that the actual result or situation differs from one's expectations or anticipations. In Japanese, sentences ending with particles often emphasize the speaker's emotions or nuances. In English, it is often expressed as "It's not going as I thought" or "Nothing is going right." Additionally, when a sentence ends with a particle such as 'ne', 'yo', or 'sa', it suggests that the sentence is incomplete and implies some lingering thought or continuation. Thus, sentences ending with particles are indeed sentences but are incomplete and serve to emphasize the speaker's emotions or nuances. On the other hand, complete sentences end with predicates and can be simply defined as sentences that end with a period (。). Such sentences are often used in records or official references and are not a form of expression for conveying personal emotions.

\newpage

\section*{558. It doesn't go as I think / It doesn't go as expected}

\textbf{Date}: 2025.10.20 (Mon)\\
\textbf{Index}: omouyouniikanai\\
\textbf{Expression (JA)}: (思:おも)うように(行:い)かない\\
\textbf{Romanization}: omou you ni ikanai\\
\textbf{Expression (EN)}: It doesn't go as I think / It doesn't go as expected

\textbf{Variants (JA)}: 予想通りに行かない / 期待通りに行かない / 考え通りに行かない\\
\textbf{Variants (EN)}: It doesn't go as I expected / It doesn't go as I anticipated / It doesn't go as I considered

\textbf{Intent (JA)}: 比較, 評価, 感想\\
\textbf{Intent (EN)}: comparison, evaluation, impression

\textbf{Tags (JA)}: 即時文法, 口語, 比較, 評価\\
\textbf{Tags (EN)}: immediate grammar, spoken, comparison, evaluation

\textbf{Adjusted Expression (JA)}:\\
(予想:よそう)通りに(行:い)かない\\
\textbf{Adjusted Expression (EN)}:\\
It doesn't go as I expected

\textbf{Example (JA)}:\\
A「(順調:じゅんちょう)に(進:すす)んでる？」 B「うーん、(思:おも)うように(行:い)かないね」\\
\textbf{Example (EN)}:\\
A: Is it going smoothly? B: Hmm, it doesn't go as I thought.

\textbf{Notes (JA)}:\\
「思うように行かない」は、予想や期待と比較して実際の結果や状況が異なるという意味。英語では"It's not going my way"や"Nothing is going right"などと表現されることが多い。少し苛立ちや失望のニュアンスが含まれることがある。

\textbf{Notes (EN)}:\\
'Omou you ni ikanai' means that the actual result or situation differs from one's expectations or anticipations. In English, it is often expressed as "It's not going my way" or "Nothing is going right." It may carry a nuance of frustration or disappointment.

\newpage

\section*{557. than I thought / than expected}

\textbf{Date}: 2025.10.19 (Sun)\\
\textbf{Index}: omottayorine\\
\textbf{Expression (JA)}: (思:おも)ったよりね\\
\textbf{Romanization}: omotta yorine\\
\textbf{Expression (EN)}: than I thought / than expected

\textbf{Variants (JA)}: 予想よりね / 期待よりね / 考えよりね\\
\textbf{Variants (EN)}: than I expected / than I anticipated / than I considered

\textbf{Intent (JA)}: 比較, 評価, 感想\\
\textbf{Intent (EN)}: comparison, evaluation, impression

\textbf{Tags (JA)}: 即時文法, 口語, 比較, 評価\\
\textbf{Tags (EN)}: immediate grammar, spoken, comparison, evaluation

\textbf{Adjusted Expression (JA)}:\\
(予想:よそう)よりも\\
\textbf{Adjusted Expression (EN)}:\\
than I expected

\textbf{Example (JA)}:\\
A「(早:はや)く(着:つ)いたね」 B「うん、(思:おも)ったよりね」\\
\textbf{Example (EN)}:\\
A: We arrived early, didn't we? B: Yeah, than I thought.

\textbf{Notes (JA)}:\\
「思ったよりね」は、予想や期待と比較して実際の結果や状況が異なることを即時的に表現する。

\textbf{Notes (EN)}:\\
'Omotta yorine' is used to immediately express that the actual result or situation differs from one's expectations or anticipations.

\newpage

\section*{556. about one hour / around one hour}

\textbf{Date}: 2025.10.18 (Sat)\\
\textbf{Index}: koichijikanhodo\\
\textbf{Expression (JA)}: (小1時間:こいちじかん)ほど\\
\textbf{Romanization}: koichijikan hodo\\
\textbf{Expression (EN)}: about one hour / around one hour

\textbf{Variants (JA)}: 約1時間 / およそ1時間 / だいたい1時間\\
\textbf{Variants (EN)}: approximately one hour / around one hour / roughly one hour

\textbf{Intent (JA)}: 時間の見積もり, 説明, 情報提供\\
\textbf{Intent (EN)}: time estimation, explanation, information provision

\textbf{Tags (JA)}: 即時文法, 書き言葉, 時間表現, 説明\\
\textbf{Tags (EN)}: immediate grammar, written language, time expression, explanation

\textbf{Adjusted Expression (JA)}:\\
(約1時間:やくいちじかん)\\
\textbf{Adjusted Expression (EN)}:\\
approximately one hour

\textbf{Example (JA)}:\\
A「(会議:かいぎ)は(小1時間:こいちじかん)ほどで」 B「では、またその(頃:ころ)」\\
\textbf{Example (EN)}:\\
A: The meeting will be about one hour. B: See you around that time.

\textbf{Notes (JA)}:\\
「小1時間ほど」は「約1時間」や「およそ1時間」と同じ意味で、時間の見積もりを表現する際に使う。だいたい1時間より少し短い時間を指すことが多く、相手にさほど長い時間ではないことを示す。それでは「小30分ほど」と言うかというと、あまり言わない。なぜかはわからないが、言いにくいことぐらいはわかる。

\textbf{Notes (EN)}:\\
'Koichijikan hodo' has the same meaning as 'yaku ichijikan' or 'oyoso ichijikan', used to express an estimate of time. It often refers to a duration slightly less than one hour, indicating to the listener that it is not a particularly long time. However, people generally do not say 'ko sanjuppun hodo' (about half an hour). The reason for this is unclear, but it is understandable that it may sound awkward.

\newpage

\section*{555. if you say ...}

\textbf{Date}: 2025.10.17 (Fri)\\
\textbf{Index}: tteiutone,\\
\textbf{Expression (JA)}: っていうとね、\\
\textbf{Romanization}: tte iutone,\\
\textbf{Expression (EN)}: if you say ...

\textbf{Variants (JA)}: そんなこといったら、 / そういうと、 / そう言ってしまったら、\\
\textbf{Variants (EN)}: If you say that, ... / If you put it that way, ... / If you say so, ...

\textbf{Intent (JA)}: 説明, 補足, 過剰の指摘\\
\textbf{Intent (EN)}: explanation, supplement, excessive pointing out

\textbf{Tags (JA)}: 即時文法, 口語, 説明, 補足\\
\textbf{Tags (EN)}: immediate grammar, spoken, explanation, supplement

\textbf{Adjusted Expression (JA)}:\\
と言ってしまいますと、\\
\textbf{Adjusted Expression (EN)}:\\
If you say so, ...

\textbf{Example (JA)}:\\
A「だいたい5(時間:じかん)ぐらいかかる、っていうとね、すごく(難:)しいような(感:かん)じがするけど、(実:じつ)は、(簡単:かんたん)」\\
\textbf{Example (EN)}:\\
A: It takes about five hours. If you say that, it sounds really difficult, but actually, it's easy.

\textbf{Notes (JA)}:\\
「っていうとね、...」は、相手もしくは自分の発言に対して補足説明や別の視点を提供する際に使う口語的な表現。もう少し詳しくこのことばの構成を説明すると、「って」「いう」「と」「ね」はそれぞれ独立した言葉で、「って」は「という」の口語縮約形、「いう」は「to say」、「と」は引用を示す助詞、「ね」は話し手の感情や意図を強調する終助詞です。つまり、「っていうとね、...」は「と言ってしまいますと、...」と同じ意味になります。

\textbf{Notes (EN)}:\\
'Tte iutone, ...' is a colloquial expression used to provide additional explanation or a different perspective regarding someone's or one's own statement. To explain the structure of this phrase in more detail, 'tte' is a colloquial contraction of 'to iu' (meaning 'to say'), 'to' is a particle indicating quotation, and 'ne' is a sentence-ending particle that emphasizes the speaker's feelings or intentions. Therefore, 'tte iutone, ...' has the same meaning as 'to itte shimaimasu to, ...' (If you say so, ...).

\newpage

\section*{554. This is the thing / This is the point / This is why ...}

\textbf{Date}: 2025.10.16 (Thu)\\
\textbf{Index}: nazekattuuto\\
\textbf{Expression (JA)}: なぜかっつーと\\
\textbf{Romanization}: naze kattsuu to\\
\textbf{Expression (EN)}: This is the thing / This is the point / This is why ...

\textbf{Variants (JA)}: 理由はね、... / 原因はね、... / ポイントはね、...\\
\textbf{Variants (EN)}: The reason is, ... / The cause is, ... / The point is, ...

\textbf{Intent (JA)}: 説明, 理由付け, 強調\\
\textbf{Intent (EN)}: explanation, justification, emphasis

\textbf{Tags (JA)}: 即時文法, 口語, 説明, 理由付け\\
\textbf{Tags (EN)}: immediate grammar, spoken, explanation, justification

\textbf{Adjusted Expression (JA)}:\\
なぜなら、\\
\textbf{Adjusted Expression (EN)}:\\
The reason is, ...

\textbf{Example (JA)}:\\
A「(今日:きょう)はね、すぐ(家:いえ)に(帰:かえ)る、なぜかっつーと、(彼女:かのじょ)の(誕生日:たんじょうび)なんだ」\\
\textbf{Example (EN)}:\\
A: I'm going home right away today. This is the thing, it's my girlfriend's birthday.

\textbf{Notes (JA)}:\\
「なぜかっつーと、...」は、相手に対して説明や理由付けを即時的に行う際に使う口語的な表現。もう少し詳しくこのことばの構成を説明すると、「なぜ」「か」「っつーと」はそれぞれ独立した言葉で、「なぜ」は「why」、「か」は疑問を示す助詞、「っつーと」は「というと」の口語縮約形です。つまり、「なぜかっつーと、...」は「なぜなら、...」と同じ意味になります。でも、これよく使うんですよ。

\textbf{Notes (EN)}:\\
'Naze kattsuu to, ...' is a colloquial expression used to provide immediate explanation or justification to someone. To explain the structure of this phrase in more detail, 'naze' means 'why', 'ka' is a particle indicating a question, and 'ttsu to' is a colloquial contraction of 'to iu to' (meaning 'if you say'). Therefore, 'naze kattsuu to, ...' has the same meaning as 'nazenara, ...' (The reason is, ...). However, this expression is frequently used in casual conversation.

\newpage

\section*{553. rather than that, ... / I mean, ... / You know, ...}

\textbf{Date}: 2025.10.15 (Wed)\\
\textbf{Index}: toiukane,...\\
\textbf{Expression (JA)}: というかね、...\\
\textbf{Romanization}: toiukane,...\\
\textbf{Expression (EN)}: rather than that, ... / I mean, ... / You know, ...

\textbf{Variants (JA)}: つまりね、... / 要するにね、... / なんていうかね、...\\
\textbf{Variants (EN)}: In other words, ... / To sum up, ... / How should I put it, ...

\textbf{Intent (JA)}: 説明, 補足, 話題転換\\
\textbf{Intent (EN)}: explanation, supplement, topic shift

\textbf{Tags (JA)}: 即時文法, 口語, 説明, 補足\\
\textbf{Tags (EN)}: immediate grammar, spoken, explanation, supplement

\textbf{Adjusted Expression (JA)}:\\
つまり、\\
\textbf{Adjusted Expression (EN)}:\\
In other words, ...

\textbf{Example (JA)}:\\
A「(彼:かれ)って、(本当:ほんとう)に(優:やさ)しいよね」 B「うん、というかね、(優:やさ)しすぎるんだよ」\\
\textbf{Example (EN)}:\\
A: He's really kind, right? B: Yeah, I mean, he's too kind.

\textbf{Notes (JA)}:\\
「というかね、...」は、相手の発言に対して補足説明や別の視点を提供する際に使う口語的な表現。もっと本音を漏らすときに使われることが多い。そして、だいたいあまり良い意味ではないことが多い。

\textbf{Notes (EN)}:\\
'Toiukane,...' is a colloquial expression used to provide additional explanation or a different perspective in response to someone's statement. It is often used when revealing more honest feelings, and it usually carries a somewhat negative connotation.

\newpage

\section*{552. I'm lucky / I'm fortunate}

\textbf{Date}: 2025.10.14 (Tue)\\
\textbf{Index}: tsuiteru\\
\textbf{Expression (JA)}: ついてる\\
\textbf{Romanization}: tsuiteru\\
\textbf{Expression (EN)}: I'm lucky / I'm fortunate

\textbf{Variants (JA)}: 運がついてる / ツイてる / 幸運がついてる\\
\textbf{Variants (EN)}: I'm lucky / I'm fortunate / Luck is on my side

\textbf{Intent (JA)}: 感謝, 喜び, 自己肯定\\
\textbf{Intent (EN)}: gratitude, joy, self-affirmation

\textbf{Tags (JA)}: 即時文法, 口語, 感謝, 喜び\\
\textbf{Tags (EN)}: immediate grammar, spoken, gratitude, joy

\textbf{Adjusted Expression (JA)}:\\
運がいいです\\
\textbf{Adjusted Expression (EN)}:\\
I'm lucky

\textbf{Example (JA)}:\\
A「(宝:たから)くじ、(当:あ)たったんだって？」 B「うん、(本当:ほんとう)についてるよ」\\
\textbf{Example (EN)}:\\
A: I heard you won the lottery? B: Yeah, I'm really lucky.

\textbf{Notes (JA)}:\\
「ついてる」は、運が良いことや幸運に恵まれていることを即時的に表現する口語的な表現。親しい間柄でよく使われます。▼逆に運が悪い場合は「ついてない」を使います。

\textbf{Notes (EN)}:\\
'Tsuiteru' is a colloquial expression used to immediately express that one is lucky or fortunate. It is often used in casual conversation among friends. Conversely, when someone is unlucky, they might use 'tsuitenai' to express that they are not fortunate.

\newpage

\section*{551. Lately ... / These days ...}

\textbf{Date}: 2025.10.13 (Mon)\\
\textbf{Index}: kokosaikin\\
\textbf{Expression (JA)}: ここ(最近:さいきん)\\
\textbf{Romanization}: koko saikin\\
\textbf{Expression (EN)}: Lately ... / These days ...

\textbf{Variants (JA)}: ここ(最近:さいきん)(忙:いそが)しい / ここ(最近:さいきん)(寒:さむ)い / ここ(最近:さいきん)(疲:つか)れる\\
\textbf{Variants (EN)}: Lately I've been busy / These days it's cold / Lately I'm getting tired

\textbf{Intent (JA)}: 近況報告, 話題提供, 感情表現\\
\textbf{Intent (EN)}: recent updates, topic introduction, emotional expression

\textbf{Tags (JA)}: 即時文法, 口語, 近況報告, 話題提供\\
\textbf{Tags (EN)}: immediate grammar, spoken, recent updates, topic introduction

\textbf{Adjusted Expression (JA)}:\\
最近\\
\textbf{Adjusted Expression (EN)}:\\
Lately ... / These days ...

\textbf{Example (JA)}:\\
A「(最近:さいきん)、(忙:いそが)しい？」 B「うん、ここ(最近:さいきん)は(本当:ほんとう)に(忙:いそが)しいよ」\\
\textbf{Example (EN)}:\\
A: Have you been busy lately? B: Yeah, these days I've really been busy.

\textbf{Notes (JA)}:\\
「ここ最近」は「最近」の前に「ここ」を付け加えたもので、話し手が現在の近況や状況を強調して伝えたいときに使う。正直に言うが、なぜ「ここ」を付けるのかというかという明瞭な理由はわからない。「そこそこちょっと前」などという表現はないので、こういうものはルール化できない。

\textbf{Notes (EN)}:\\
'Koko saikin' is formed by adding 'koko' (here) before 'saikin' (recently), used when the speaker wants to emphasize their current situation or recent updates. Honestly, there isn't a clear reason why 'koko' is added; there are no expressions like 'soko soko chotto mae' (there, a little while ago), so such expressions cannot be strictly rule-based.

\newpage

\section*{550. I get dragged into ... / I end up ...}

\textbf{Date}: 2025.10.12 (Sun)\\
\textbf{Index}: hikizurarechau\\
\textbf{Expression (JA)}: 引きずられちゃう\\
\textbf{Romanization}: hikizurarechau\\
\textbf{Expression (EN)}: I get dragged into ... / I end up ...

\textbf{Variants (JA)}: 感情に引きずられちゃう / 過去に引きずられちゃう / 状況に引きずられちゃう\\
\textbf{Variants (EN)}: I get dragged into emotions / I end up dwelling on the past / I get dragged into the situation

\textbf{Intent (JA)}: 感情表現, 自己反省, 状況説明\\
\textbf{Intent (EN)}: emotional expression, self-reflection, situation explanation

\textbf{Tags (JA)}: 即時文法, 口語, 感情表現, 自己反省\\
\textbf{Tags (EN)}: immediate grammar, spoken, emotional expression, self-reflection

\textbf{Adjusted Expression (JA)}:\\
とらわれてしまいます\\
\textbf{Adjusted Expression (EN)}:\\
I end up being caught up in ...

\textbf{Example (JA)}:\\
A「(彼:かれ)の(言:い)い(方:かた)に(引:ひ)きずられちゃうよね」 B「うん、(本当:ほんとう)に。(冷静:れいせい)に(考:かんが)えると(全然:ぜんぜん)(賛成:さんせい)できないんだけどね」\\
\textbf{Example (EN)}:\\
A: You get dragged into his way of speaking, right? B: Yeah, really. When I think about it calmly, I can't agree with it at all.

\textbf{Notes (JA)}:\\
「引きずられちゃう」は「引きずられてしまう」の口語縮約形で、「ひきづられる」とは、望んでいないのに、他者や場の勢いなどに影響され、巻き込まれてしまう様子を意味します。物理的に引きずられる「車にひきずられる」は物理的に力によって引っ張られてしまう状態、「精神的に引きずられる」は過去の嫌な出来事の記憶に囚われてしまう状態、「場の雰囲気に引きずられる(流される)」はその場の勢いや、周囲の人々の意見に押されて行動してしまう状態を示します。

\textbf{Notes (EN)}:\\
'Hikizurarechau' is a colloquial contraction of 'hikizurarete shimau', meaning to be influenced or caught up in something against one's will, such as others or the momentum of a situation. 'Hikizurareru' can refer to being physically dragged (e.g., being dragged by a car) or mentally caught up in past unpleasant memories. It can also refer to being swayed by the atmosphere of a situation or the opinions of those around you.

\newpage

\section*{549. I can't keep up / I can't follow}

\textbf{Date}: 2025.10.11 (Sat)\\
\textbf{Index}: tsuitekenai\\
\textbf{Expression (JA)}: ついてけない\\
\textbf{Romanization}: tsuitekenai\\
\textbf{Expression (EN)}: I can't keep up / I can't follow

\textbf{Variants (JA)}: 理解ついてけない / 話についてけない / 変化についてけない\\
\textbf{Variants (EN)}: I can't keep up with understanding / I can't follow the conversation / I can't keep up with the changes

\textbf{Intent (JA)}: 困惑, 理解不足, 感情表現\\
\textbf{Intent (EN)}: confusion, lack of understanding, emotional expression

\textbf{Tags (JA)}: 即時文法, 口語, 感情表現, 困惑\\
\textbf{Tags (EN)}: immediate grammar, spoken, emotional expression, confusion

\textbf{Adjusted Expression (JA)}:\\
ついていけません\\
\textbf{Adjusted Expression (EN)}:\\
I can't keep up

\textbf{Example (JA)}:\\
A「(最近:さいきん)、(新:あたら)しい(技術:ぎじゅつ)が(次々:つぎつぎ)に(出:で)てくるよね」 B「うん、(本当:ほんとう)についてけないよ」\\
\textbf{Example (EN)}:\\
A: New technologies are coming out one after another these days, right? B: Yeah, I really can't keep up.

\textbf{Notes (JA)}:\\
「ついてけない」は「ついていけない」の口語縮約形で、相手の話や状況に対して困惑や理解不足を即時的に表現する。感情を強調し、親しい間柄で使われることが多い。フォーマルな場では「ついていけません」などに置き換えられる。

\textbf{Notes (EN)}:\\
'Tsuitekenai' is a colloquial contraction of 'tsuite ikenai', used to immediately express confusion or lack of understanding regarding someone's story or situation. It emphasizes emotion and is often used among close acquaintances. In formal contexts, it can be replaced with 'tsuite ikemasen' (I can't keep up).

\newpage

\section*{548. No way! / It can't be!}

\textbf{Date}: 2025.10.10 (Fri)\\
\textbf{Index}: sonnnaa\\
\textbf{Expression (JA)}: そんなあ！\\
\textbf{Romanization}: sonn naa!\\
\textbf{Expression (EN)}: No way! / It can't be!

\textbf{Variants (JA)}: 嘘でしょ！ / まさか！ / 信じられない！\\
\textbf{Variants (EN)}: You're kidding! / No way! / I can't believe it!

\textbf{Intent (JA)}: 驚き, 否定, 感情表現\\
\textbf{Intent (EN)}: surprise, denial, emotional expression

\textbf{Tags (JA)}: 即時文法, 口語, 感情表現, 驚き\\
\textbf{Tags (EN)}: immediate grammar, spoken, emotional expression, surprise

\textbf{Adjusted Expression (JA)}:\\
それはおかしいですね。\\
\textbf{Adjusted Expression (EN)}:\\
Is that really possible?

\textbf{Example (JA)}:\\
A「(今日:きょう)から、また(元:もと)の(値段:ねだん)に(戻:もど)るんだって」 B「そんなあ！」\\
\textbf{Example (EN)}:\\
A: I heard the price is going back to the original from today. B: No way!

\textbf{Notes (JA)}:\\
「そんなあ！」は、驚きや否定を即時的に表現する口語的な表現。親しい間柄でよく使われ、フォーマルな場では基本的にすぐに不平不満は言わない方が良いので、これに相当する表現はあまり使われないが、どうしても使う場合は「それはおかしいですね」と言うとよい。

\textbf{Notes (EN)}:\\
'Sonn naa!' is a colloquial expression used to immediately convey surprise or denial. It is commonly used among close acquaintances. In formal contexts, it is generally better not to express complaints immediately, so there isn't a direct equivalent expression. However, if you must use one, saying 'sore wa okashii desu ne' (Is that really possible?) would be appropriate.

\newpage

\section*{547. You mean ...? / Are you saying ...?}

\textbf{Date}: 2025.10.09 (Thu)\\
\textbf{Index}: ttedounano?\\
\textbf{Expression (JA)}: ってどうなの？\\
\textbf{Romanization}: tte dou nano?\\
\textbf{Expression (EN)}: You mean ...? / Are you saying ...?

\textbf{Variants (JA)}: どう？ / いいの？ / だいじょうぶ？\\
\textbf{Variants (EN)}: How about it? / Is it okay? / Is it alright?

\textbf{Intent (JA)}: 確認, 理解, 驚き\\
\textbf{Intent (EN)}: confirmation, understanding, surprise

\textbf{Tags (JA)}: 即時文法, 口語, 確認, 理解\\
\textbf{Tags (EN)}: immediate grammar, spoken, confirmation, understanding

\textbf{Adjusted Expression (JA)}:\\
ということは、どういうことですか？\\
\textbf{Adjusted Expression (EN)}:\\
So, what does that mean?

\textbf{Example (JA)}:\\
A「(明日:あした)、(雨:あめ)だって」 B「ってどうなの？(中止:ちゅうし)？」\\
\textbf{Example (EN)}:\\
A: They say it's going to rain tomorrow. B: You mean, is it canceled?

\textbf{Notes (JA)}:\\
「ってどうなの？」は、相手の発言に基づいて確認や理解を求める即時的な表現。親しい間柄でよく使われ、フォーマルな場では「ということは、どういうことですか？」などに置き換えられる。

\textbf{Notes (EN)}:\\
'Tte dou nano?' is an immediate expression used to seek confirmation or understanding based on someone's statement. It is commonly used among close acquaintances. In formal contexts, it can be replaced with 'to iu koto wa, dou iu koto desu ka?' (So, what does that mean?).

\newpage

\section*{546. You mean ...? / Are you saying ...?}

\textbf{Date}: 2025.10.08 (Wed)\\
\textbf{Index}: tteiukotoha?\\
\textbf{Expression (JA)}: っていうことは？\\
\textbf{Romanization}: tte iu koto wa?\\
\textbf{Expression (EN)}: You mean ...? / Are you saying ...?

\textbf{Variants (JA)}: つまり？ / 要するに？ / ということは？\\
\textbf{Variants (EN)}: In other words? / To sum up? / So, are you saying...?

\textbf{Intent (JA)}: 確認, 理解, 驚き\\
\textbf{Intent (EN)}: confirmation, understanding, surprise

\textbf{Tags (JA)}: 即時文法, 口語, 確認, 理解\\
\textbf{Tags (EN)}: immediate grammar, spoken, confirmation, understanding

\textbf{Adjusted Expression (JA)}:\\
すなわち、\\
\textbf{Adjusted Expression (EN)}:\\
In other words, ...

\textbf{Example (JA)}:\\
A「(明日:あした)、(雨:あめ)だって」 B「っていうことは、(中止:ちゅうし)？」\\
\textbf{Example (EN)}:\\
A: They say it's going to rain tomorrow. B: You mean it's canceled?

\textbf{Notes (JA)}:\\
「っていうことは？」は、相手の発言に基づいて確認や理解を求める即時的な表現。親しい間柄でよく使われ、フォーマルな場では「すなわち、」などに置き換えられる。

\textbf{Notes (EN)}:\\
'Tte iu koto wa?' is an immediate expression used to seek confirmation or understanding based on someone's statement. It is commonly used among close acquaintances. In formal contexts, it can be replaced with 'sunawachi,' (In other words, ...).

\newpage

\section*{545. Is it over already? / Is that it already?}

\textbf{Date}: 2025.10.07 (Tue)\\
\textbf{Index}: aremouowari?\\
\textbf{Expression (JA)}: あれ！もう(終:お)わり？\\
\textbf{Romanization}: are mou owari?\\
\textbf{Expression (EN)}: Is it over already? / Is that it already?

\textbf{Variants (JA)}: それ！もう(終:お)わり？ / これ！もう(終:お)わり？ / あれ！もう(終:お)わったの？\\
\textbf{Variants (EN)}: Is that it already? / Is this over already? / Is it already finished?

\textbf{Intent (JA)}: 驚き, 確認, 感情表現\\
\textbf{Intent (EN)}: surprise, confirmation, emotional expression

\textbf{Tags (JA)}: 即時文法, 口語, 感情表現, 確認\\
\textbf{Tags (EN)}: immediate grammar, spoken, emotional expression, confirmation

\textbf{Adjusted Expression (JA)}:\\
もう(終:お)わりですか？\\
\textbf{Adjusted Expression (EN)}:\\
Is it over already?

\textbf{Example (JA)}:\\
A「これで、おしまい、(面白:おもしろ)かったね」 B「あれ！もう(終:お)わり？」\\
\textbf{Example (EN)}:\\
A: That's it, it was fun, right? B: Is it over already?

\textbf{Notes (JA)}:\\
「あれもう終わり？」は、相手の発言や状況に対して驚きや確認を即時的に表現する口語的な表現。親しい間柄でよく使われ、フォーマルな場では「もう終わりですか？」などに置き換えられる。

\textbf{Notes (EN)}:\\
'Are mou owari?' is a colloquial expression used to immediately convey surprise or seek confirmation regarding someone's statement or situation. It is commonly used among close acquaintances. In formal contexts, it can be replaced with 'mou owari desu ka?' (Is it over already?).

\newpage

\section*{544. I definitely want to ... / I really want to ...}

\textbf{Date}: 2025.10.06 (Mon)\\
\textbf{Index}: zehishitaidesuyo\\
\textbf{Expression (JA)}: ぜひしたいですよ\\
\textbf{Romanization}: zehi shitai desu yo\\
\textbf{Expression (EN)}: I definitely want to ... / I really want to ...

\textbf{Variants (JA)}: ぜひ行きたいですよ / ぜひ会いたいですよ / ぜひ参加したいですよ\\
\textbf{Variants (EN)}: I definitely want to go / I really want to meet / I definitely want to participate

\textbf{Intent (JA)}: 強い願望, 熱意, 積極性\\
\textbf{Intent (EN)}: strong desire, enthusiasm, proactiveness

\textbf{Tags (JA)}: 即時文法, 口語, 願望, 強調\\
\textbf{Tags (EN)}: immediate grammar, spoken, desire, emphasis

\textbf{Adjusted Expression (JA)}:\\
ぜひ(〜たい)と(思:おも)います\\
\textbf{Adjusted Expression (EN)}:\\
I definitely want to ...

\textbf{Example (JA)}:\\
A「(次回:じかい)は、もっと(時間:じかん)をかけて、(議論:ぎろん)しませんか？」 B「ええ、ぜひしたいですよ！」\\
\textbf{Example (EN)}:\\
A: Next time, shall we take more time for discussion? B: Yes, I definitely want to!

\textbf{Notes (JA)}:\\
「ぜひしたいですよ」は、強い願望や熱意を即時的に表現する口語的な表現。親しい間柄でよく使われ、フォーマルな場では「ぜひ〜たいと思います」などに置き換えられる。

\textbf{Notes (EN)}:\\
'Zehi shitai desu yo' is a colloquial expression used to convey strong desire or enthusiasm immediately. It is commonly used among close acquaintances. In formal contexts, it can be replaced with 'zehi ...tai to omoimasu' (I definitely want to ...).

\newpage

\section*{543. I agree. / That's true. / I think so too.}

\textbf{Date}: 2025.10.05 (Sun)\\
\textbf{Index}: daisansei\\
\textbf{Expression (JA)}: (大賛成:だいさんせい)\\
\textbf{Romanization}: dai sansei\\
\textbf{Expression (EN)}: I agree. / That's true. / I think so too.

\textbf{Variants (JA)}: 大反対 / 大歓迎 / 大歓迎\\
\textbf{Variants (EN)}: I strongly disagree / I warmly welcome / I strongly support

\textbf{Intent (JA)}: 同意, 賛成, 強調\\
\textbf{Intent (EN)}: agreement, approval, emphasis

\textbf{Tags (JA)}: 即時文法, 書き言葉, 同意, 強調\\
\textbf{Tags (EN)}: immediate grammar, written language, agreement, emphasis

\textbf{Adjusted Expression (JA)}:\\
(大:おお)いに(賛成:さんせい)です\\
\textbf{Adjusted Expression (EN)}:\\
I strongly agree

\textbf{Example (JA)}:\\
A「じゃ、(ビール:びーる)、(飲:の)み(行:い)こう！」 B「うん、(大賛成:だいさんせい)！」\\
\textbf{Example (EN)}:\\
A: So, let's go for a beer! B: Yeah, I totally agree!

\textbf{Notes (JA)}:\\
「大賛成」は、相手の意見や提案に対して強く同意・賛成する際に使う表現。書き言葉でよく使われ、フォーマルな場でも適している。会話ではやや堅苦しい印象を与えることがある。フォーマルには「大いに賛成です」などに置き換えられる。

\textbf{Notes (EN)}:\\
'Dai sansei' is used to strongly agree or approve of someone's opinion or proposal. It is commonly found in written language and is suitable for formal contexts. In conversation, it may sound somewhat stiff. In formal settings, it can be replaced with 'ooi ni sansei desu' (I strongly agree).

\newpage

\section*{542. That may be so, but ... / Be that as it may, ...}

\textbf{Date}: 2025.10.04 (Sat)\\
\textbf{Index}: sorehasoretoshite\\
\textbf{Expression (JA)}: それはそれとして\\
\textbf{Romanization}: sore wa sore to shite\\
\textbf{Expression (EN)}: That may be so, but ... / Be that as it may, ...

\textbf{Variants (JA)}: 問題は問題として / 事実は事実として / 意見は意見として\\
\textbf{Variants (EN)}: The problem is the problem, but ... / The fact is the fact, but ... / The opinion is the opinion, but ...

\textbf{Intent (JA)}: 対比, 話題転換, 補足\\
\textbf{Intent (EN)}: contrast, topic shift, supplement

\textbf{Tags (JA)}: 即時文法, 口語, 対比, 話題転換\\
\textbf{Tags (EN)}: immediate grammar, spoken, contrast, topic shift

\textbf{Adjusted Expression (JA)}:\\
それは(別:べつ)の(問題:もんだい)だとしてですね...\\
\textbf{Adjusted Expression (EN)}:\\
That may be so, but ...

\textbf{Example (JA)}:\\
A「(彼:かれ)、(優:やさ)しいよね」 B「うん、それはそれとして、(仕事:しごと)は(厳:きび)しいよね」\\
\textbf{Example (EN)}:\\
A: He's kind, right? B: Yeah, that may be so, but his work is tough.

\textbf{Notes (JA)}:\\
「それはそれとして」は、相手の意見や事実を認めつつも、別の視点や話題に転換したいときに使う表現。親しい間柄でよく使われ、フォーマルな場では「それは別の問題だとしてですね」などに置き換えられる。

\textbf{Notes (EN)}:\\
'Sore wa sore to shite' is used to acknowledge someone's opinion or fact while shifting to a different perspective or topic. It is commonly used among close acquaintances. In formal contexts, it can be replaced with 'sore wa betsu no mondai da to shite desu ne' (That may be so, but ...).

\newpage

\section*{541. I see. / Is that so?}

\textbf{Date}: 2025.10.03 (Fri)\\
\textbf{Index}: sounanda\\
\textbf{Expression (JA)}: そうなんだ\\
\textbf{Romanization}: sou nan da\\
\textbf{Expression (EN)}: I see. / Is that so?

\textbf{Variants (JA)}: 本当そうなんだ / へえ、そうなんだ / なるほど、そうなんだ\\
\textbf{Variants (EN)}: Is that really so? / Oh, is that so? / I see, is that so?

\textbf{Intent (JA)}: 理解, 驚き, 共感\\
\textbf{Intent (EN)}: understanding, surprise, empathy

\textbf{Tags (JA)}: 即時文法, 口語, 理解, 驚き\\
\textbf{Tags (EN)}: immediate grammar, spoken, understanding, surprise

\textbf{Adjusted Expression (JA)}:\\
そうなんですか\\
\textbf{Adjusted Expression (EN)}:\\
I see. / Is that so?

\textbf{Example (JA)}:\\
A「(彼:かれ)、(結婚:けっこん)するんだって」 B「へえ、そうなんだ」\\
\textbf{Example (EN)}:\\
A: I heard he's getting married. B: Oh, is that so?

\textbf{Notes (JA)}:\\
「そうなんだ」は、相手の話を聞いて理解や驚きを即時的に表現する口語的な表現。親しい間柄でよく使われ、フォーマルな場では「そうですか」などに置き換えられる。

\textbf{Notes (EN)}:\\
'Sou nan da' is a colloquial expression used to immediately convey understanding or surprise upon hearing someone's story. It is commonly used among close acquaintances. In formal contexts, it can be replaced with 'sou desu ka' (I see. / Is that so?).

\newpage

\section*{540. That's true, but ... / I agree, but ...}

\textbf{Date}: 2025.10.02 (Thu)\\
\textbf{Index}: sounandakedone\\
\textbf{Expression (JA)}: そうなんだけどね\\
\textbf{Romanization}: sou nan da kedo ne\\
\textbf{Expression (EN)}: That's true, but ... / I agree, but ...

\textbf{Variants (JA)}: いいんだけどね / わかるんだけどね / 助かるんだけどね\\
\textbf{Variants (EN)}: That's fine, but ... / I understand, but ... / That helps, but ...

\textbf{Intent (JA)}: 同意, 反論, 補足\\
\textbf{Intent (EN)}: agreement, objection, supplement

\textbf{Tags (JA)}: 即時文法, 口語, 同意, 反論\\
\textbf{Tags (EN)}: immediate grammar, spoken, agreement, objection

\textbf{Adjusted Expression (JA)}:\\
そうなんですけどね、\\
\textbf{Adjusted Expression (EN)}:\\
That's true, but ...

\textbf{Example (JA)}:\\
A「(彼:かれ)、(優:やさ)しいよね」 B「うん、そうなんだけどね、(優:やさ)しすぎるんだよ」\\
\textbf{Example (EN)}:\\
A: He's kind, right? B: Yeah, that's true, but he's too kind.

\textbf{Notes (JA)}:\\
「そうなんだけどね」は、相手の意見に同意しつつも、何らかの反論や補足を加える際に使う口語的な表現。親しい間柄でよく使われ、フォーマルな場では「そうですが」などに置き換えられる。

\textbf{Notes (EN)}:\\
'Sou nan da kedo ne' is a colloquial expression used to agree with someone's opinion while also adding an objection or additional information. It is commonly used among close acquaintances. In formal contexts, it can be replaced with 'sou desu ga' (That's true, but ...).

\newpage

\section*{539. Even if ... / Assuming that ...}

\textbf{Date}: 2025.10.01 (Wed)\\
\textbf{Index}: oitoitatoshite\\
\textbf{Expression (JA)}: おいといたとして\\
\textbf{Romanization}: oitoita toshite\\
\textbf{Expression (EN)}: even if ... / assuming that ...

\textbf{Variants (JA)}: 忘れといたとして / 無視しといたとして / 考えといたとして\\
\textbf{Variants (EN)}: even if you forget / even if you ignore / even if you consider

\textbf{Intent (JA)}: 仮定, 条件, 思考実験\\
\textbf{Intent (EN)}: hypothesis, condition, thought experiment

\textbf{Tags (JA)}: 即時文法, 口語, 仮定, 条件\\
\textbf{Tags (EN)}: immediate grammar, spoken, hypothesis, condition

\textbf{Adjusted Expression (JA)}:\\
おいておいたとしても\\
\textbf{Adjusted Expression (EN)}:\\
even if you put aside ...

\textbf{Example (JA)}:\\
A「(実現性:じつげんせい)は(置:お)いといたとして、(面白:おもしろ)い(話:はなし)だよね」 B「うん、(確:たし)かに」\\
\textbf{Example (EN)}:\\
A: Even if we put aside the feasibility, it's an interesting story, right? B: Yeah, certainly.

\textbf{Notes (JA)}:\\
「おいといたとして」は「おく」の口語的な連結形「おいとく」に過去形「た」をつけ、さらに仮定の「として」を加えた表現。ある事柄を一時的に脇に置き、仮定や条件を設定して思考実験を行う際に使われる。親しい間柄で使われ、フォーマルな場では「おいておいたとしても」などに置き換えられる。

\textbf{Notes (EN)}:\\
'Oitoita toshite' combines the colloquial connective form 'oitoiku' of 'oku' (to put aside) with the past tense 'ta', followed by the hypothetical 'toshite'. It is used to temporarily set aside a matter and establish a hypothesis or condition for thought experiments. It is often used among close acquaintances. In formal contexts, it can be replaced with 'oite oita toshitemo' (even if you put aside ...).

\newpage

\section*{538. Something like ... / Kind of like ...}

\textbf{Date}: 2025.09.30 (Tue)\\
\textbf{Index}: younamono\\
\textbf{Expression (JA)}: ようなもの\\
\textbf{Romanization}: you na mono\\
\textbf{Expression (EN)}: something like ... / kind of like ...

\textbf{Variants (JA)}: 夢のようなもの / 映画のようなもの / おとぎ話のようなもの\\
\textbf{Variants (EN)}: something like a dream / something like a movie / something like a fairy tale

\textbf{Intent (JA)}: 比喩, 例示, 曖昧表現\\
\textbf{Intent (EN)}: metaphor, example, vague expression

\textbf{Tags (JA)}: 即時文法, 口語, 比喩, 曖昧さ\\
\textbf{Tags (EN)}: immediate grammar, spoken, metaphor, vagueness

\textbf{Adjusted Expression (JA)}:\\
のようなもの\\
\textbf{Adjusted Expression (EN)}:\\
something like ...

\textbf{Example (JA)}:\\
A「(彼:かれ)の(話:はなし)、(夢:ゆめ)のようなもんだよね」 B「うん、(信:しん)じられないよね」\\
\textbf{Example (EN)}:\\
A: His story is something like a dream, right? B: Yeah, it's unbelievable.

\textbf{Notes (JA)}:\\
「ようなもの/ん」は、比喩的に何かを説明したり、例示したりする即時的な表現。具体的な言葉を避け、曖昧さを持たせることで、話し手の感情やニュアンスを伝える。親しい間柄で使われることが多く、フォーマルな場では「のようなもの/ん」などに置き換えられる。

\textbf{Notes (EN)}:\\
'You na mono/n' is an immediate expression used metaphorically to describe or exemplify something. By avoiding specific terms and introducing vagueness, it conveys the speaker's emotions or nuances. It is often used among close acquaintances. In formal contexts, it can be replaced with 'no you na mono' (something like ...).

\newpage

\section*{537. I want to ... / I feel like ...}

\textbf{Date}: 2025.09.29 (Mon)\\
\textbf{Index}: takute\\
\textbf{Expression (JA)}: たくて\\
\textbf{Romanization}: takute\\
\textbf{Expression (EN)}: I want to ... / I feel like ...

\textbf{Variants (JA)}: 食べたくて / 行きたくて / 話したくて\\
\textbf{Variants (EN)}: I want to eat / I want to go / I want to talk

\textbf{Intent (JA)}: 願望, 動機, 感情表現\\
\textbf{Intent (EN)}: desire, motivation, emotional expression

\textbf{Tags (JA)}: 即時文法, 口語, 願望, 動機\\
\textbf{Tags (EN)}: immediate grammar, spoken, desire, motivation

\textbf{Adjusted Expression (JA)}:\\
〜たいので\\
\textbf{Adjusted Expression (EN)}:\\
because I want to ...

\textbf{Example (JA)}:\\
A「(野菜:やさい)だけは、いい(材料:ざいりょう)を(使:つか)いたくて」」 B「へえ、そうなんだ」\\
\textbf{Example (EN)}:\\
A: I just want to use good ingredients for the vegetables. B: Oh, I see.

\textbf{Notes (JA)}:\\
「たくて」は動詞のたい形に接続助詞「て」がついたもので、願望や動機を即時的に表現する。感情を強調し、親しい間柄で使われることが多い。フォーマルな場では「〜たいので」などに置き換えられる。

\textbf{Notes (EN)}:\\
'Takute' combines the 'tai' form of a verb with the connective particle 'te', expressing desire or motivation in an immediate way. It emphasizes emotion and is often used among close acquaintances. In formal contexts, it can be replaced with '...tai node' (because I want to ...).

\newpage

\section*{536. so ... huh? / it's ... really ... huh?}

\textbf{Date}: 2025.09.28 (Sun)\\
\textbf{Index}: nkai\\
\textbf{Expression (JA)}: おいしいんかい\\
\textbf{Romanization}: oishii n kai\\
\textbf{Expression (EN)}: so it's tasty, huh?

\textbf{Variants (JA)}: 高いんかい / 安いんかい / 難しいんかい\\
\textbf{Variants (EN)}: so it's expensive, huh? / so it's cheap, huh? / so it's difficult, huh?

\textbf{Intent (JA)}: 驚き, 確認, 感情表現\\
\textbf{Intent (EN)}: surprise, confirmation, emotional expression

\textbf{Tags (JA)}: 即時文法, 口語, 感情表現, 確認\\
\textbf{Tags (EN)}: immediate grammar, spoken, emotional expression, confirmation

\textbf{Adjusted Expression (JA)}:\\
おいしいのですね\\
\textbf{Adjusted Expression (EN)}:\\
So it's tasty, huh?

\textbf{Example (JA)}:\\
A「この(料理:りょうり)、(美味:おい)しいよ」 B「おいしいんかい！」\\
\textbf{Example (EN)}:\\
A: This dish is tasty. B: So it's tasty, huh!

\textbf{Notes (JA)}:\\
「んかい」は「の＋かい」の口語縮約形で、相手の発言に対して驚きや確認を即時的に表現する。感情を強調し、親しい間柄で使われることが多い。これは関西弁で相手の言ったことに驚いたり、確認したりする際に使われる表現です。フォーマルな場では「のですね」などに置き換えられる。

\textbf{Notes (EN)}:\\
'N kai' is a colloquial contraction of 'no + kai', used to express surprise or seek confirmation in response to someone's statement. It emphasizes emotion and is often used among close acquaintances. This expression is typical in Kansai dialect when reacting to what someone has said. In formal contexts, it can be replaced with 'no desu ne' (So it's ... huh?).

\newpage

\section*{535. Even if you don't say it, ... / Without saying it, ...}

\textbf{Date}: 2025.09.27 (Sat)\\
\textbf{Index}: tohaiwanaimademo,\\
\textbf{Expression (JA)}: とは(言:い)わないまでも、\\
\textbf{Romanization}: to wa iwanai mademo,\\
\textbf{Expression (EN)}: even if you don't say it, ... / without saying it, ...

\textbf{Variants (JA)}: 好きとは言わないまでも / 嫌いとは言わないまでも / 得意とは言わないまでも\\
\textbf{Variants (EN)}: Even if you don't say you like it / Even if you don't say you hate it / Even if you don't say you're good at it

\textbf{Intent (JA)}: 控えめな表現, 婉曲な同意, 前置き\\
\textbf{Intent (EN)}: modest expression, indirect agreement, preface

\textbf{Tags (JA)}: 即時文法, 書き言葉, 婉曲表現, 前置き\\
\textbf{Tags (EN)}: immediate grammar, written language, euphemistic expression, preface

\textbf{Adjusted Expression (JA)}:\\
とは(言:い)いませんが\\
\textbf{Adjusted Expression (EN)}:\\
I wouldn't say that, but ...

\textbf{Example (JA)}:\\
A「(彼:かれ)、(天才:てんさい)だよね」 B「うーん、(天才:てんさい)とは(言:い)わないまでも、(優:すぐ)れてるよね」\\
\textbf{Example (EN)}:\\
A: He's a genius, right? B: Hmm, even if you don't say he's a genius, he's certainly exceptional.

\textbf{Notes (JA)}:\\
「とは言わないまでも、」は、相手の意見に対して完全には同意しないが、ある程度認めるときに使う婉曲的な表現。書き言葉でよく使われ、フォーマルな場面でも適している。会話ではやや堅苦しい印象を与えることがある。フォーマルには「とは言いませんが」などに置き換えられる。

\textbf{Notes (EN)}:\\
'To wa iwanai mademo,' is a euphemistic expression used to partially agree with someone's opinion without fully endorsing it. It is commonly found in written language and is suitable for formal contexts. In conversation, it may sound somewhat stiff. In formal settings, it can be replaced with 'to wa iimasen ga' (I wouldn't say that, but ...).

\newpage

\section*{534. Well, that's that. / So be it.}

\textbf{Date}: 2025.09.26 (Fri)\\
\textbf{Index}: jasorede.\\
\textbf{Expression (JA)}: じゃ、それで。\\
\textbf{Romanization}: ja, sore de.\\
\textbf{Expression (EN)}: Well, that's that. / So be it.

\textbf{Variants (JA)}: じゃ、そうするよ / じゃ、それでいいよ / じゃ、決まりだね\\
\textbf{Variants (EN)}: Well, I'll do that / Well, that's fine / Well, it's decided then

\textbf{Intent (JA)}: 決定, 結論, 受け入れ\\
\textbf{Intent (EN)}: decision, conclusion, acceptance

\textbf{Tags (JA)}: 即時文法, 口語, 結論, 受け入れ\\
\textbf{Tags (EN)}: immediate grammar, spoken, conclusion, acceptance

\textbf{Adjusted Expression (JA)}:\\
それでよろしいです\\
\textbf{Adjusted Expression (EN)}:\\
That will be fine

\textbf{Example (JA)}:\\
A「出発、(明日:あした)に(延期:えんき)しようか」 B「うん、じゃ、それで」\\
\textbf{Example (EN)}:\\
A: So, shall we postpone the departure to tomorrow? B: Yeah, well, that's that.

\textbf{Notes (JA)}:\\
「じゃ、それで」は会話の中で即時的に結論を出し、受け入れるときに使う表現。親しい間柄でよく使われ、決定事項を軽く確認するニュアンスがある。フォーマルな場では「それでよろしいです」などに言い換えられる。

\textbf{Notes (EN)}:\\
'Ja, sore de.' is used to quickly conclude and accept a decision in conversation. It is common among close acquaintances and carries a light tone of confirmation. In formal contexts, it can be replaced with 'sore de yoroshii desu' (That will be fine).

\newpage

\section*{533. Even if ... / For ...}

\textbf{Date}: 2025.09.25 (Thu)\\
\textbf{Index}: nishichaa\\
\textbf{Expression (JA)}: にしちゃあ\\
\textbf{Romanization}: ni shichaa\\
\textbf{Expression (EN)}: even if ... / for ...

\textbf{Variants (JA)}: 若いにしちゃあ / 安いにしちゃあ / 上手いにしちゃあ\\
\textbf{Variants (EN)}: for being young / for being cheap / for being good

\textbf{Intent (JA)}: 評価, 比較, 驚き\\
\textbf{Intent (EN)}: evaluation, comparison, surprise

\textbf{Tags (JA)}: 即時文法, 口語, 比較, 驚き\\
\textbf{Tags (EN)}: immediate grammar, spoken, comparison, surprise

\textbf{Adjusted Expression (JA)}:\\
〜にしては\\
\textbf{Adjusted Expression (EN)}:\\
for ... / considering ...

\textbf{Example (JA)}:\\
A「(彼:かれ)、(若:わか)いよね」 B「うん、(若:わか)いにしちゃあ、(落:お)ち(着:つ)いてるよね」\\
\textbf{Example (EN)}:\\
A: He's young, right? B: Yeah, for being young, he's pretty composed.

\textbf{Notes (JA)}:\\
「にしちゃあ」は「〜にしては」のくだけた口語表現で、ある基準や期待と比較して評価や驚きを即時的に伝える。親しい間柄で使われ、フォーマルな場では「〜にしては」と言い換えられる。

\textbf{Notes (EN)}:\\
'Ni shichaa' is a casual spoken form of '...ni shite wa', used to express evaluation or surprise by comparing against a certain standard or expectation. It is commonly used among close acquaintances. In formal contexts, it can be replaced with '...ni shite wa'.

\newpage

\section*{532. Come on, do it.}

\textbf{Date}: 2025.09.24 (Wed)\\
\textbf{Index}: teyo\\
\textbf{Expression (JA)}: みてきてよ\\
\textbf{Romanization}: mite kite yo\\
\textbf{Expression (EN)}: Go check it out, will you?

\textbf{Variants (JA)}: たべてよ / みてよ / なんとかしてよ / やってよ\\
\textbf{Variants (EN)}: Eat it! / look at it! / Do something about it! / Come on, do it!

\textbf{Intent (JA)}: 依頼, 感情, 圧力, 親しみ\\
\textbf{Intent (EN)}: request, emotion, pressure, familiarity

\textbf{Tags (JA)}: 即時文法, 口語, 終助詞, 感情\\
\textbf{Tags (EN)}: immediate grammar, spoken, sentence-final particle, emotion

\textbf{Adjusted Expression (JA)}:\\
(見:み)てきてください\\
\textbf{Adjusted Expression (EN)}:\\
Could you go and have a look, please?

\textbf{Example (JA)}:\\
A「(外:そと)、(見:み)てきて」 B「えー、(見:み)てきてよ〜」\\
\textbf{Example (EN)}:\\
A: Go check outside. B: Ugh, you go check it out!

\textbf{Notes (JA)}:\\
「..てよ」は、て形に終助詞「よ」がつくことで、親しい相手に対して感情をこめて即時的に依頼・要求する口語表現となる。やわらかくも、どこか圧力のある言い方として使われ、特に子どもや友人同士の会話で自然に使われる。フォーマルには「...てください」と言い換えられる。

\textbf{Notes (EN)}:\\
"...te yo" attaches the sentence-final particle 'yo' to the te-form of a verb, creating a casual yet emotionally charged request. While often friendly, it can imply pressure or insistence, especially among close friends or children. In formal settings, it corresponds to polite requests like '...te kudasai'.

\newpage

\section*{531. Nothing.}

\textbf{Date}: 2025.09.23 (Tue)\\
\textbf{Index}: nanimo\\
\textbf{Expression (JA)}: (何:なに)も\\
\textbf{Romanization}: nani mo\\
\textbf{Expression (EN)}: Nothing.

\textbf{Variants (JA)}: 何もない / 何もなかった / 何もいらない\\
\textbf{Variants (EN)}: There is nothing / There was nothing / I don't need anything

\textbf{Intent (JA)}: 応答, 否定, 曖昧\\
\textbf{Intent (EN)}: response, negation, ambiguity

\textbf{Tags (JA)}: 即時文法, 応答, 省略, 曖昧さ\\
\textbf{Tags (EN)}: immediate grammar, response, ellipsis, ambiguity

\textbf{Adjusted Expression (JA)}:\\
(特:とく)にありません\\
\textbf{Adjusted Expression (EN)}:\\
Nothing in particular

\textbf{Example (JA)}:\\
A「どうしたの？」 B「(何:なに)も。」\\
\textbf{Example (EN)}:\\
A: What's wrong? B: Nothing.

\textbf{Notes (JA)}:\\
「何も」は問いかけに対して短く応答する即時文法の表現で、暗黙的に「何もない」「特にない」という否定を含意する。あえて文を省略することで曖昧さを残し、会話の余地を与える。フォーマルな場面では「特にありません」と言い換えられる。

\textbf{Notes (EN)}:\\
'Nani mo' is an immediate grammar expression used as a short response, implicitly carrying a negation such as 'there is nothing' or 'nothing in particular'. By omitting the full sentence, it leaves room for ambiguity and keeps the conversation open. In formal contexts, it can be expressed as 'tokuni arimasen' (nothing in particular).

\newpage

\section*{530. Ah, so that's your move.}

\textbf{Date}: 2025.09.22 (Mon)\\
\textbf{Index}: soukitaka\\
\textbf{Expression (JA)}: そう、きたか\\
\textbf{Romanization}: sou, kita ka\\
\textbf{Expression (EN)}: so that's how you're playing it

\textbf{Variants (JA)}: そうきたか / おお、そうきたか / そう来るとは\\
\textbf{Variants (EN)}: so that’s how it came / didn't expect that / ah, you went that way

\textbf{Intent (JA)}: 驚き, 受けとめ, 即時評価\\
\textbf{Intent (EN)}: surprise, uptake, instant evaluation

\textbf{Tags (JA)}: 即時文法, 口語, 反応, 感嘆\\
\textbf{Tags (EN)}: immediate grammar, spoken, reaction, exclamation

\textbf{Adjusted Expression (JA)}:\\
そう(:)来ましたか\\
\textbf{Adjusted Expression (EN)}:\\
I see, so that was your move

\textbf{Example (JA)}:\\
A「で、そこでまさかの(:)謝罪文ですよ」 B「...そう、きたか。」\\
\textbf{Example (EN)}:\\
A: So then he suddenly sent an apology letter. B: ...So that's how you're playing it.

\textbf{Notes (JA)}:\\
「そう、きたか」は、相手の意外な行動や展開に対して即座に反応し、受け止めるときに使う表現。感心・納得・軽い驚きを含むことが多く、単なる感嘆ではなく状況の受容を示す。文頭に「ああ」や「へえ」が付くこともある。フォーマルな場では「そう来ましたか」などに言い換えられる。

\textbf{Notes (EN)}:\\
'Sou, kitaka' is used when reacting to an unexpected turn of events or clever move. It expresses not just surprise, but also a kind of impressed acceptance. Often used in casual, reflective tones. In formal settings, it might be phrased as 'Sou kimashita ka' (I see, so that was your move).

\newpage

\section*{529. Hmm, that's kind of...}

\textbf{Date}: 2025.09.21 (Sun)\\
\textbf{Index}: nattotoka\\
\textbf{Expression (JA)}: ...なっと\\
\textbf{Romanization}: …natto\\
\textbf{Expression (EN)}: hmm, that's kind of...

\textbf{Variants (JA)}: むずかしいなっと / こわいなっと / 知らないなっと\\
\textbf{Variants (EN)}: Hmm, that’s kinda hard / That’s kind of scary / I don’t really know, hmm

\textbf{Intent (JA)}: 評価, つぶやき, 即時判断\\
\textbf{Intent (EN)}: evaluation, muttering, immediate judgment

\textbf{Tags (JA)}: 即時文法, 口語, 評価, 曖昧\\
\textbf{Tags (EN)}: immediate grammar, spoken, evaluation, vagueness

\textbf{Adjusted Expression (JA)}:\\
〜だなと(思:おも)って\\
\textbf{Adjusted Expression (EN)}:\\
I thought it was ...

\textbf{Example (JA)}:\\
A「やってみる？」 B「うーん、これは、ちょっと、むずかしいなっと」\\
\textbf{Example (EN)}:\\
A: Wanna give it a try? B: Hmm, this seems kinda hard...

\textbf{Notes (JA)}:\\
「なっと」は「だな＋っと」の口語的な連結・訛りで、評価やつぶやきを即時的に表現する表現。言い切らず、判断を含ませたまま終えるニュアンスがあり、相手に含みを伝える。フォーマルな場では「〜だと思いました」などに置き換えられる。

\textbf{Notes (EN)}:\\
'Natto' in this context is a colloquial blend of 'da na' and 'tto', expressing immediate, often soft or hesitant judgments. It conveys a private thought, slightly left unsaid, and appears in informal spoken Japanese. In formal language, it would be replaced with 'I thought it was...'.

\newpage

\section*{528. Lately / These days}

\textbf{Date}: 2025.09.20 (Sat)\\
\textbf{Index}: kokontoko\\
\textbf{Expression (JA)}: ここんとこ\\
\textbf{Romanization}: kokon toko\\
\textbf{Expression (EN)}: lately / these days

\textbf{Variants (JA)}: ここんとこ忙しくて / ここんとこ会ってないね / ここんとこずっと眠い\\
\textbf{Variants (EN)}: I've been busy lately / We haven't met lately / I've been feeling sleepy these days

\textbf{Intent (JA)}: 近況, 話題導入, 状況共有\\
\textbf{Intent (EN)}: recent status, topic opener, context sharing

\textbf{Tags (JA)}: 即時文法, 口語, 省略, 時間表現\\
\textbf{Tags (EN)}: immediate grammar, spoken, ellipsis, temporal expression

\textbf{Adjusted Expression (JA)}:\\
このところ\\
\textbf{Adjusted Expression (EN)}:\\
recently / these days

\textbf{Example (JA)}:\\
A「(元気:げんき)？」 B「いや、ここんとこちょっと(疲:つか)れてて」\\
\textbf{Example (EN)}:\\
A: How are you? B: Hmm, I've been kind of tired lately.

\textbf{Notes (JA)}:\\
「ここんとこ」は「このところ」「ここ最近」のくだけた口語表現で、時間の幅を曖昧にしたまま近況を伝えるときに使われる。省略形の響きが親しみやすさを与え、話題の導入にも適している。フォーマルな場では「このところ」などに置き換えられる。

\textbf{Notes (EN)}:\\
'Kokontoko' is a colloquial contraction of 'kono tokoro' or 'koko saikin', used to refer to a vaguely recent span of time. It adds a friendly, casual tone and is often used to introduce a current situation. In formal contexts, it can be replaced with 'kono tokoro' or 'saikin'.

\newpage

\section*{527. Do you know?}

\textbf{Date}: 2025.09.19 (Fri)\\
\textbf{Index}: shitteru\\
\textbf{Expression (JA)}: (知:し)ってる？\\
\textbf{Romanization}: shitteru?\\
\textbf{Expression (EN)}: do you know?

\textbf{Variants (JA)}: これ知ってる？ / あの人知ってる？ / その話知ってる？\\
\textbf{Variants (EN)}: Do you know this? / Do you know that person? / Do you know that story?

\textbf{Intent (JA)}: 確認, 話題導入, 共有意識\\
\textbf{Intent (EN)}: confirmation, topic introduction, shared awareness

\textbf{Tags (JA)}: 即時文法, 口語, 確認, 共有\\
\textbf{Tags (EN)}: immediate grammar, spoken, confirmation, shared knowledge

\textbf{Adjusted Expression (JA)}:\\
ご(存:ぞん)じですか\\
\textbf{Adjusted Expression (EN)}:\\
do you happen to know

\textbf{Example (JA)}:\\
A「ねえ、これ(知:し)ってる？」 B「うん、(前:まえ)に(見:み)たよ」\\
\textbf{Example (EN)}:\\
A: Hey, do you know this? B: Yeah, I saw it before.

\textbf{Notes (JA)}:\\
「知ってる？」は相手の知識を即時に確認する口語的な表現で、驚きや話題の導入に使われる。親しい間柄では軽く言えるが、フォーマルな場では「ご存知ですか」に置き換えられる。

\textbf{Notes (EN)}:\\
'Shitteru?' is a colloquial way to check whether the other person knows something. It is often used to introduce a topic or express surprise. In casual contexts, it is light and friendly, while in formal situations it can be replaced with 'gozonji desu ka' (do you happen to know?).

\newpage

\section*{526. What else can you do?}

\textbf{Date}: 2025.09.18 (Thu)\\
\textbf{Index}: hokanandekiru\\
\textbf{Expression (JA)}: (他:ほか)、(何:なに)、できる？\\
\textbf{Romanization}: hoka, nani, dekiru?\\
\textbf{Expression (EN)}: what else can you do?

\textbf{Variants (JA)}: 他、できることある？ / 他、何かできたりする？ / 他には？\\
\textbf{Variants (EN)}: Is there anything else you can do? / Can you do anything else? / Anything else?

\textbf{Intent (JA)}: 確認, 話題継続, 探り\\
\textbf{Intent (EN)}: checking, follow-up, exploratory

\textbf{Tags (JA)}: 即時文法, 口語, 省略, 順序の自由\\
\textbf{Tags (EN)}: immediate grammar, spoken, ellipsis, flexible word order

\textbf{Adjusted Expression (JA)}:\\
(他:ほか)に(何:なに)ができますか？\\
\textbf{Adjusted Expression (EN)}:\\
What else can you do?

\textbf{Example (JA)}:\\
A「(英語:えいご)できるんだ？」 B「うん」 A「(他:ほか)、(何:なに)、できる？」\\
\textbf{Example (EN)}:\\
A: You speak English? B: Yeah. A: What else can you do?

\textbf{Notes (JA)}:\\
「他、何、できる？」は、助詞や語順を省略しつつ、相手の能力や選択肢を探る即時的な表現。質問というより話題継続や共有関心を示す働きを持つ。文法的には未完文だが、口語会話では自然で柔らかい印象を与える。フォーマルには「他に何ができますか？」とする。

\textbf{Notes (EN)}:\\
'Hoka, nani, dekiru?' is an immediate and casual way of asking what else someone can do. It omits particles and rearranges word order freely, often as part of a continuing conversation. While technically incomplete, it sounds natural and gentle in spoken Japanese. In formal contexts, it becomes 'Hoka ni nani ga dekimasu ka?'

\newpage

\section*{525. You're good at it! / Not bad!}

\textbf{Date}: 2025.09.17 (Wed)\\
\textbf{Index}: dekirenee\\
\textbf{Expression (JA)}: できるねぇ\\
\textbf{Romanization}: dekiru nee\\
\textbf{Expression (EN)}: you're good at it! / not bad!

\textbf{Variants (JA)}: よくできるねぇ / そんなこともできるねぇ / 一人でできるねぇ\\
\textbf{Variants (EN)}: You can really do it! / You can even do that! / You can do it all by yourself!

\textbf{Intent (JA)}: 感心, 評価, 驚き\\
\textbf{Intent (EN)}: admiration, praise, surprise

\textbf{Tags (JA)}: 即時文法, 口語, 評価, 感嘆\\
\textbf{Tags (EN)}: immediate grammar, spoken, evaluation, exclamation

\textbf{Adjusted Expression (JA)}:\\
(:)上手ですね / よくできますね\\
\textbf{Adjusted Expression (EN)}:\\
you're skilled / you do it well

\textbf{Example (JA)}:\\
A「(:)見て、この(:)絵、どう！」 B「できるねぇ！」\\
\textbf{Example (EN)}:\\
A: Look at this painting, how do you think? B: Not bad!

\textbf{Notes (JA)}:\\
「できるねぇ」は相手の能力や行動を目の前で評価・感心する即時的な表現。砕けた口語で親しい間柄で使われることが多く、驚きや共感を含む。フォーマルな場では「上手ですね」「よくできますね」と言い換えられる。

\textbf{Notes (EN)}:\\
'Dekiru nee' is an immediate expression of praise or admiration for someone's ability or action. It is casual and used in close relationships, often with a tone of surprise or shared feeling. In formal contexts, it can be replaced with 'jouzu desu ne' (you're skilled) or 'yoku dekimasu ne' (you do it well).

\newpage

\section*{524. Well, you see...}

\textbf{Date}: 2025.09.15 (Mon)\\
\textbf{Index}: sorewane\\
\textbf{Expression (JA)}: うーん、それはね\\
\textbf{Romanization}: uun, sore wa ne\\
\textbf{Expression (EN)}: well, you see...

\textbf{Variants (JA)}: うーん、それはね、ちょっと違うよ / うーん、それはね、説明が難しいんだ / うーん、それはね、わからない\\
\textbf{Variants (EN)}: Well, you see, that's a bit different / Well, you see, it's hard to explain / Well, you see, I don't know

\textbf{Intent (JA)}: 思考中, 前置き, 柔らかい説明\\
\textbf{Intent (EN)}: thinking, preface, soft explanation

\textbf{Tags (JA)}: 即時文法, 口語, 前置き, ためらい\\
\textbf{Tags (EN)}: immediate grammar, spoken, preface, hesitation

\textbf{Adjusted Expression (JA)}:\\
それについてはですね\\
\textbf{Adjusted Expression (EN)}:\\
regarding that, you see

\textbf{Example (JA)}:\\
A「なんで(行:い)かなかったの？」 B「うーん、それはね、ちょっと...」\\
\textbf{Example (EN)}:\\
A: Why didn't you go? B: Well, you see, it's a bit...

\textbf{Notes (JA)}:\\
「うーん、それはね」は考えながら答えるときの前置きで、相手にやわらかく説明の余地を残す即時表現。ためらいや思考中のニュアンスを含み、会話に自然な間を作る。フォーマルな場では「それについてはですね」などに置き換えられる。

\textbf{Notes (EN)}:\\
'Uun, sore wa ne' is a prefatory phrase used while thinking, signaling that an explanation will follow. It conveys hesitation or consideration, and creates a natural pause in conversation. In formal contexts, it can be replaced with 'sore ni tsuite wa desu ne' (regarding that...).

\newpage

\section*{523. Oops, I ... / I ended up ...}

\textbf{Date}: 2025.09.14 (Sun)\\
\textbf{Index}: chatta\\
\textbf{Expression (JA)}: ちゃった\\
\textbf{Romanization}: chatta\\
\textbf{Expression (EN)}: oops, I ... / I ended up ...

\textbf{Variants (JA)}: 寝ちゃった / 忘れちゃった / 食べちゃった\\
\textbf{Variants (EN)}: I fell asleep (oops) / I forgot (oops) / I ate it (oops)

\textbf{Intent (JA)}: 失敗, 予期せぬ行動, 軽い後悔\\
\textbf{Intent (EN)}: mistake, unexpected action, mild regret

\textbf{Tags (JA)}: 即時文法, 口語, 縮約, 感情\\
\textbf{Tags (EN)}: immediate grammar, spoken, contraction, emotion

\textbf{Adjusted Expression (JA)}:\\
〜してしまった\\
\textbf{Adjusted Expression (EN)}:\\
ended up doing ...

\textbf{Example (JA)}:\\
A「(宿題:しゅくだい)どうした？」 B「(忘:わす)れちゃった！」\\
\textbf{Example (EN)}:\\
A: What about the homework? B: I forgot (oops)!

\textbf{Notes (JA)}:\\
「ちゃった」は「〜してしまった」の口語縮約形で、失敗や予期せぬ行動を軽く即時的に伝える表現。深刻な後悔ではなく、照れや軽い謝罪のニュアンスが多い。フォーマルな場では「〜してしまいました」と言い換えられる。

\textbf{Notes (EN)}:\\
'Chatta' is the colloquial contraction of '...shite shimatta.' It expresses mistakes or unintended actions with a light, immediate tone, often implying embarrassment or mild apology rather than deep regret. In formal contexts, it is expressed as '...shite shimaimashita.'

\newpage

\section*{522. You got me? / Busted!}

\textbf{Date}: 2025.09.13 (Sat)\\
\textbf{Index}: bareta\\
\textbf{Expression (JA)}: (バレ:ばれ)た？\\
\textbf{Romanization}: bareta?\\
\textbf{Expression (EN)}: You got me? / Busted!

\textbf{Variants (JA)}: 隠してたけどバレた？ / 嘘ついたのバレた？ / 気持ちバレた？\\
\textbf{Variants (EN)}: I tried to hide it, but you noticed? / Did you catch my lie? / Did you notice my feelings?

\textbf{Intent (JA)}: 本音露呈, 秘密発覚, からかい\\
\textbf{Intent (EN)}: true feelings revealed, secret exposed, playful admission

\textbf{Tags (JA)}: 即時文法, 口語, からかい, 感情\\
\textbf{Tags (EN)}: immediate grammar, spoken, playful, emotion

\textbf{Adjusted Expression (JA)}:\\
(気:き)づかれましたか？\\
\textbf{Adjusted Expression (EN)}:\\
Did you notice?

\textbf{Example (JA)}:\\
A「(本当:ほんとう)は(行:い)きたくなかったんでしょ？」 B「(バレ:ばれ)た？」\\
\textbf{Example (EN)}:\\
A: You actually didn't want to go, right? B: Busted!

\textbf{Notes (JA)}:\\
「バレた？」は自分の本音や秘密が相手に見抜かれたときに使う即時的な表現。多くの場合、深刻さよりも軽い照れやからかいを含んで使われる。フォーマルな場では「気づかれましたか？」などに置き換えられる。

\textbf{Notes (EN)}:\\
'Bareta?' is used when your true feelings or a secret have been noticed by someone else. It usually carries a playful or embarrassed tone rather than seriousness. In formal contexts, it can be expressed as 'Did you notice?'

\newpage

\section*{521. Right? / Isn't it?}

\textbf{Date}: 2025.09.12 (Fri)\\
\textbf{Index}: dayone\\
\textbf{Expression (JA)}: だよね\\
\textbf{Romanization}: da yo ne\\
\textbf{Expression (EN)}: right? / isn't it?

\textbf{Variants (JA)}: いい天気だよね / 難しかったよね / 楽しかったよね\\
\textbf{Variants (EN)}: It's nice weather, right? / That was difficult, wasn't it? / It was fun, wasn't it?

\textbf{Intent (JA)}: 同意確認, 共感, 会話のつなぎ\\
\textbf{Intent (EN)}: confirmation, empathy, conversational link

\textbf{Tags (JA)}: 即時文法, 口語, 共感, 確認\\
\textbf{Tags (EN)}: immediate grammar, spoken, empathy, confirmation

\textbf{Adjusted Expression (JA)}:\\
そうですよね\\
\textbf{Adjusted Expression (EN)}:\\
that's true

\textbf{Example (JA)}:\\
A「ふつう、こんなことしないよ！」 B「だよね！」\\
\textbf{Example (EN)}:\\
A: Normally, nobody would do such a thing! B: Exactly!

\textbf{Notes (JA)}:\\
「だよね」は相手の発話に即時的に同意・共感を示す表現で、親しい間柄でよく使われる。単独で強い共感を表すこともできる。フォーマルな場面では「そうですね」「〜ではありませんか」に置き換えられる。

\textbf{Notes (EN)}:\\
'Da yo ne' is used to show immediate agreement or empathy with what the other person said, common in casual speech. It can stand alone to strongly express empathy. In formal contexts, it can be replaced with 'sou desu ne' or '...de wa arimasen ka'.

\newpage

\section*{520. Normally / Usually}

\textbf{Date}: 2025.09.11 (Thu)\\
\textbf{Index}: itsumonara\\
\textbf{Expression (JA)}: いつもなら\\
\textbf{Romanization}: itsumo nara\\
\textbf{Expression (EN)}: normally / usually

\textbf{Variants (JA)}: いつもなら寝ている時間 / いつもなら怒るところ / いつもなら行かない店\\
\textbf{Variants (EN)}: Normally I'd be asleep at this time / Normally I'd get angry / Normally I wouldn't go to this shop

\textbf{Intent (JA)}: 基準提示, 比較, 状況の違い\\
\textbf{Intent (EN)}: baseline, comparison, situation difference

\textbf{Tags (JA)}: 即時文法, 口語, 習慣, 比較\\
\textbf{Tags (EN)}: immediate grammar, spoken, habit, comparison

\textbf{Adjusted Expression (JA)}:\\
(通常:つうじょう)であれば\\
\textbf{Adjusted Expression (EN)}:\\
under normal circumstances

\textbf{Example (JA)}:\\
A「(今日:きょう)は(遅:おそ)いね」 B「いつもならもう(帰:かえ)ってるんだけど」\\
\textbf{Example (EN)}:\\
A: You're late today. B: Normally I'd already be home.

\textbf{Notes (JA)}:\\
「いつもなら」は普段の習慣や基準を持ち出して、今回の状況との差を示す表現。即時的に比較を持ち出すため、驚きや違和感、状況説明を自然に伝えることができる。フォーマルでは「通常であれば」と言い換えられる。

\textbf{Notes (EN)}:\\
'Itsumo nara' brings up one's usual habit or baseline to contrast with the current situation. It conveys surprise, discomfort, or explanation. In formal contexts, it can be replaced with 'tsuujou de areba' (under normal circumstances).

\newpage

\section*{519. By all means / No matter what}

\textbf{Date}: 2025.09.10 (Wed)\\
\textbf{Index}: nantoshitemo\\
\textbf{Expression (JA)}: なんとしても\\
\textbf{Romanization}: nantoshitemo\\
\textbf{Expression (EN)}: by all means / no matter what

\textbf{Variants (JA)}: なんとしても成功させたい / なんとしても行かなければならない / なんとしても守り抜く\\
\textbf{Variants (EN)}: I want to succeed by all means / I must go no matter what / I will protect it at all costs

\textbf{Intent (JA)}: 強い意志, 決意, 必然\\
\textbf{Intent (EN)}: strong will, determination, inevitability

\textbf{Tags (JA)}: 即時文法, 口語, 強調, 意志\\
\textbf{Tags (EN)}: immediate grammar, spoken, emphasis, will

\textbf{Adjusted Expression (JA)}:\\
どのような(手段:しゅだん)を(取:と)っても\\
\textbf{Adjusted Expression (EN)}:\\
by any means / whatever the means

\textbf{Example (JA)}:\\
A「(無理:むり)じゃない？」 B「なんとしてもやり(遂:と)げる！」\\
\textbf{Example (EN)}:\\
A: Isn't it impossible? B: I'll get it done no matter what!

\textbf{Notes (JA)}:\\
「なんとしても」は強い決意や必然性を即時に表す表現。通常はポジティブな意志に使われるが、場合によっては義務感や切迫感を含むこともある。フォーマルでは「どうしても」「必ず」などに置き換えられる。

\textbf{Notes (EN)}:\\
'Nantoshitemo' expresses strong determination or inevitability. It is usually used for positive resolve but can also convey obligation or urgency. In formal contexts, 'doushitemo' or 'kanarazu' can be used instead.

\newpage

\section*{518. If you like / If anything}

\textbf{Date}: 2025.09.09 (Tue)\\
\textbf{Index}: nannara\\
\textbf{Expression (JA)}: なんなら\\
\textbf{Romanization}: nannara\\
\textbf{Expression (EN)}: if you like / if anything

\textbf{Variants (JA)}: なんなら手伝いますよ / なんなら帰ってもいい / なんならやめてもいいよ\\
\textbf{Variants (EN)}: If you like, I can help / If anything, you can go home / If necessary, you can quit

\textbf{Intent (JA)}: 提案, 強調, 柔らかい逆接\\
\textbf{Intent (EN)}: suggestion, emphasis, soft contrast

\textbf{Tags (JA)}: 即時文法, 口語, 提案, 含み\\
\textbf{Tags (EN)}: immediate grammar, spoken, suggestion, implication

\textbf{Adjusted Expression (JA)}:\\
もしよければ / (必要:ひつよう)なら\\
\textbf{Adjusted Expression (EN)}:\\
if you like / if necessary

\textbf{Example (JA)}:\\
A「(明日:あす)、(荷物:)(多:)いんだ」 B「なんなら(車:)で(送:)るよ」\\
\textbf{Example (EN)}:\\
A: I'll have a lot of luggage tomorrow. B: If you like, I can give you a ride.

\textbf{Notes (JA)}:\\
「なんなら」は直訳しにくい即時文法的表現で、相手に柔らかく提案や選択肢を示す。場合によっては逆接的に「むしろ」「かえって」の意味も持つ。フォーマルな場では「もしよければ」「必要なら」と言い換えられる。

\textbf{Notes (EN)}:\\
'Nannara' is an immediate grammar expression that gently offers a suggestion or option. Depending on context, it can also carry the sense of 'rather' or 'if anything.' In formal contexts, it can be expressed as 'moshi yokereba' (if you like) or 'hitsuyou nara' (if necessary).

\newpage

\section*{517. Little by little / It's about time}

\textbf{Date}: 2025.09.08 (Mon)\\
\textbf{Index}: bochibochi\\
\textbf{Expression (JA)}: ぼちぼち\\
\textbf{Romanization}: bochibochi\\
\textbf{Expression (EN)}: little by little / it's about time

\textbf{Variants (JA)}: ぼちぼち行こか / 仕事はぼちぼちやってます / 景気はぼちぼちやな\\
\textbf{Variants (EN)}: Shall we get going? / I'm working little by little / The economy is so-so

\textbf{Intent (JA)}: 緩やかな進行, そろそろの合図, 控えめな評価\\
\textbf{Intent (EN)}: gradual progress, signal for it's time, modest evaluation

\textbf{Tags (JA)}: 即時文法, 口語, 緩和, 進行\\
\textbf{Tags (EN)}: immediate grammar, spoken, softening, progress

\textbf{Adjusted Expression (JA)}:\\
(少:すこ)しずつ / そろそろ\\
\textbf{Adjusted Expression (EN)}:\\
gradually / it's about time

\textbf{Example (JA)}:\\
A「(最近:さいきん)どう？」 B「まあ、ぼちぼちかな」\\
\textbf{Example (EN)}:\\
A: How are things lately? B: Well, little by little, I guess.

\textbf{Notes (JA)}:\\
「ぼちぼち」は関西弁由来の口語表現で、全国的に使われるようになった。『そろそろ』の意味で行動を始める合図にも、『まあまあ』『少しずつ』のような控えめな評価にも使える。即時的に状況を和らげて伝える即時文法的表現。フォーマルな場では「少しずつ」「そろそろ」などに置き換えられる。

\textbf{Notes (EN)}:\\
'Bochi bochi' originated in Kansai dialect but is now widely understood. It can signal 'it's about time' to start something, or express modest evaluation like 'so-so' or 'little by little.' It softens statements and conveys progress in an immediate, conversational way. In formal contexts, it is replaced with 'sukoshi zutsu' (gradually) or 'sorosoro' (it's about time).

\newpage

\section*{516. That's it}

\textbf{Date}: 2025.09.07 (Sun)\\
\textbf{Index}: soregasore\\
\textbf{Expression (JA)}: それがそれ\\
\textbf{Romanization}: sore ga sore\\
\textbf{Expression (EN)}: that's it

\textbf{Variants (JA)}: それがそれなんだよ / それがそれだから困る / それがそれってこと\\
\textbf{Variants (EN)}: That's it, you know / That's it, that's the problem / That's what it is

\textbf{Intent (JA)}: 納得, 共感, 問題の指摘\\
\textbf{Intent (EN)}: acceptance, empathy, pointing out an issue

\textbf{Tags (JA)}: 即時文法, 口語, 納得, 含み\\
\textbf{Tags (EN)}: immediate grammar, spoken, acceptance, implication

\textbf{Adjusted Expression (JA)}:\\
それがつまり...ということです\\
\textbf{Adjusted Expression (EN)}:\\
in other words, that is the point

\textbf{Example (JA)}:\\
A「それって、(東急:とうきゅう)じゃないよね...」 B「それがそれなんだよ。」\\
\textbf{Example (EN)}:\\
A: That's not Tokyu, right... B: That's it.

\textbf{Notes (JA)}:\\
「それがそれ」は、一見あいまいに見えるが、相手の発話をそのまま受け止め『まさにそれが答え／事情だ』と即時に納得を示す表現。場面によっては「ご名答」に近いニュアンスになる。説明用に「噂のお店って...」と前置きすることもあるが、実際の自然なやりとりでは「それって、東急じゃないよね...」「それがそれなんだよ」のようにシンプルに発話される。フォーマルな場面では「つまり、それが問題です」などに言い換えられる。

\textbf{Notes (EN)}:\\
'Sore ga sore' looks vague, but functions as an immediate grammar expression to confirm what the other person has just said as the key point. It is close to saying 'Exactly' or 'That's the answer,' often carrying a nuance of 'spot on.' While explanatory contexts may add background like 'That trendy shop...', in natural conversation it appears simply as in 'That's not Tokyu, right...' / 'That's it.' In formal contexts, it can be expressed as 'In other words, that is the issue.'

\newpage

\section*{515. Fairly / relatively}

\textbf{Date}: 2025.09.06 (Sat)\\
\textbf{Index}: warito\\
\textbf{Expression (JA)}: わりと\\
\textbf{Romanization}: warito\\
\textbf{Expression (EN)}: fairly / relatively

\textbf{Variants (JA)}: わりと安い / わりと静か / わりとよくできている\\
\textbf{Variants (EN)}: fairly cheap / relatively quiet / rather well done

\textbf{Intent (JA)}: 評価, 意外性, 口語的和らげ\\
\textbf{Intent (EN)}: evaluation, unexpectedness, colloquial softening

\textbf{Tags (JA)}: 即時文法, 口語, 評価, 意外性\\
\textbf{Tags (EN)}: immediate grammar, spoken, evaluation, unexpectedness

\textbf{Adjusted Expression (JA)}:\\
(比較的:ひかくてき)\\
\textbf{Adjusted Expression (EN)}:\\
comparatively / relatively

\textbf{Example (JA)}:\\
A「この(店:みせ)、(高:たか)い？」 B「ううん、わりと(安:やす)いよ。」\\
\textbf{Example (EN)}:\\
A: Is this place expensive? B: No, it's fairly cheap.

\textbf{Notes (JA)}:\\
「わりと」は思ったより程度が軽いことを即時的に評価する表現。日常会話で多用され、肯定的にも否定的にも使える。フォーマルな場面では「比較的」と言い換えられる。

\textbf{Notes (EN)}:\\
'Warito' is an immediate grammar expression used to evaluate something as less extreme than expected. It is common in casual conversation and can be either positive or negative. In formal contexts, it is replaced with 'hikakuteki' (comparatively).

\newpage

\section*{514. Taking the trouble to...}

\textbf{Date}: 2025.09.05 (Fri)\\
\textbf{Index}: wazawaza\\
\textbf{Expression (JA)}: わざわざ\\
\textbf{Romanization}: wazawaza\\
\textbf{Expression (EN)}: taking the trouble to / going out of one's way

\textbf{Variants (JA)}: わざわざ来てくれてありがとう / わざわざ言わなくてもいい / わざわざ持ってきてくれた\\
\textbf{Variants (EN)}: Thank you for taking the trouble to come / You didn’t have to say it / You went out of your way to bring it

\textbf{Intent (JA)}: 感謝, 非難, 強調\\
\textbf{Intent (EN)}: gratitude, criticism, emphasis

\textbf{Tags (JA)}: 即時文法, 口語, 感謝, 強調\\
\textbf{Tags (EN)}: immediate grammar, spoken, gratitude, emphasis

\textbf{Adjusted Expression (JA)}:\\
(特別:とくべつ)に\\
\textbf{Adjusted Expression (EN)}:\\
especially / particularly

\textbf{Example (JA)}:\\
A「わざわざ(届:とど)けてくれて(助:たす)かったよ。」 B「いいえ、ついでですから。」\\
\textbf{Example (EN)}:\\
A: Thanks for going out of your way to deliver it. B: No, it was on the way.

\textbf{Notes (JA)}:\\
「わざわざ」は相手の行為が手間や労力をかけていることを即座に認識して表す即時文法的表現。文脈によって感謝にも非難にもなる。フォーマルな文脈では「特別に」「あえて」などに言い換えられる。

\textbf{Notes (EN)}:\\
'Wazawaza' is an immediate grammar expression used to highlight that an action involves special effort or trouble. Depending on context, it can express gratitude or criticism. In formal contexts, it can be replaced with 'tokubetsu ni' (especially) or 'aete' (intentionally).

\newpage

\section*{513. Oh, so that’s how it was}

\textbf{Date}: 2025.09.04 (Thu)\\
\textbf{Index}: soudattanda\\
\textbf{Expression (JA)}: そうだったんだ\\
\textbf{Romanization}: sou datta n da\\
\textbf{Expression (EN)}: Oh, so that’s how it was

\textbf{Variants (JA)}: そうだったんですね / そうだったのか / そうだったんだね\\
\textbf{Variants (EN)}: Oh, I see, that’s how it was / So that’s what happened / So that’s how it was, huh

\textbf{Intent (JA)}: 納得, 理解, 気づき\\
\textbf{Intent (EN)}: acceptance, understanding, realization

\textbf{Tags (JA)}: 即時文法, 口語, 納得, 発見\\
\textbf{Tags (EN)}: immediate grammar, spoken, acceptance, realization

\textbf{Adjusted Expression (JA)}:\\
そうだったのですね\\
\textbf{Adjusted Expression (EN)}:\\
I see, that was the case

\textbf{Example (JA)}:\\
A「(彼:かれ)、(転校:てんこう)してたんだよ。」 B「そうだったんだ！」\\
\textbf{Example (EN)}:\\
A: He had transferred schools. B: Oh, so that’s how it was!

\textbf{Notes (JA)}:\\
「そうだったんだ」は、相手の説明を聞いて初めて理解や納得が生じたときに使う即時文法的表現。過去形を用いているが、今この瞬間の気づきを表しているのが特徴。形式的には過去の事実を確認しているようでありながら、実際は『なるほど』『そういうことか』と同じく感情の定型句として使われる。フォーマルな場面では「そうだったのですね」と言い換えられる。

\textbf{Notes (EN)}:\\
'Sou datta n da' is an immediate grammar expression used when new information brings about understanding or acceptance. Although it employs the past form 'datta,' it conveys the present realization of 'I see' or 'So that’s what happened.' It functions as a fixed expression of realization rather than a literal past tense. In formal contexts, it becomes 'sou datta no desu ne.'

\newpage

\section*{512. If it had..., then...}

\textbf{Date}: 2025.09.03 (Wed)\\
\textbf{Index}: tanara\\
\textbf{Expression (JA)}: ...たなら、...\\
\textbf{Romanization}: ta nara, ...\\
\textbf{Expression (EN)}: If it had..., then...

\textbf{Variants (JA)}: 行けたなら、会えたのに / 知っていたなら、準備した / 一緒だったなら、楽しかった\\
\textbf{Variants (EN)}: If I had been able to go, I could have seen you / If I had known, I would have prepared / If we had been together, it would have been fun

\textbf{Intent (JA)}: 仮定, 後悔, 感情の余韻\\
\textbf{Intent (EN)}: hypothetical, regret, emotional lingering

\textbf{Tags (JA)}: 即時文法, 口語, 仮定, 感情\\
\textbf{Tags (EN)}: immediate grammar, spoken, hypothetical, emotion

\textbf{Adjusted Expression (JA)}:\\
もし...たならば、...したでしょう\\
\textbf{Adjusted Expression (EN)}:\\
If it had..., then it would have...

\textbf{Example (JA)}:\\
(勇気:ゆうき)があったなら、(言:い)えたのに...\\
\textbf{Example (EN)}:\\
If I had had the courage, I could have said it...

\textbf{Notes (JA)}:\\
「...たなら、...」は、過去の事実とは異なる仮定を提示し、後悔や残念さを表す即時文法的表現。多くの場合、文末を省略して感情を余韻として残す。フォーマルな文脈では「もし...たならば、...したでしょう」と言い換えられる。

\textbf{Notes (EN)}:\\
'...ta nara, ...' is an immediate grammar expression presenting a counterfactual condition, often conveying regret or disappointment. It commonly trails off to leave the feeling implicit. In formal contexts, it is expressed as 'moshi ...ta naraba, ...shita deshou'.

\newpage

\section*{511. I'm telling you...}

\textbf{Date}: 2025.09.02 (Tue)\\
\textbf{Index}: datteba\\
\textbf{Expression (JA)}: ...だって(ば)\\
\textbf{Romanization}: ...datte(ba)\\
\textbf{Expression (EN)}: I'm telling you...

\textbf{Variants (JA)}: 本当だって / 行ったんだってば / 違うんだって\\
\textbf{Variants (EN)}: It really is, I'm telling you / I went, I'm telling you / It's not true, I said

\textbf{Intent (JA)}: 強調, 念押し, 不満\\
\textbf{Intent (EN)}: emphasis, insistence, frustration

\textbf{Tags (JA)}: 即時文法, 口語, 強調, 感情\\
\textbf{Tags (EN)}: immediate grammar, spoken, emphasis, emotion

\textbf{Adjusted Expression (JA)}:\\
(確:たし)かにそうなんです\\
\textbf{Adjusted Expression (EN)}:\\
It is certainly so

\textbf{Example (JA)}:\\
A「ほんとにできるの？」 B「できるんだってば！」\\
\textbf{Example (EN)}:\\
A: Can you really do it? B: I'm telling you, I can!

\textbf{Notes (JA)}:\\
「...だって(ば)」は、伝聞の「だって」と異なり、話し手がすでに述べたことを繰り返し強調し、相手に理解を迫る即時文法的表現。しばしば苛立ちや感情を伴う。フォーマルな場面では「確かにそうです」「本当です」などに置き換えられる。

\textbf{Notes (EN)}:\\
Unlike the quotative 'datte' (as in 'he said that...'), '...datte(ba)' is used when the speaker insists on something they already said, pressing the listener to accept it. It often carries frustration or emotion. In formal contexts, it can be replaced with 'It is certainly so' or 'Indeed, it is so'.

\newpage

\section*{510. I managed to...}

\textbf{Date}: 2025.09.01 (Mon)\\
\textbf{Index}: nantoka\\
\textbf{Expression (JA)}: なんとか\\
\textbf{Romanization}: nantoka\\
\textbf{Expression (EN)}: somehow / manage to

\textbf{Variants (JA)}: なんとかのりこえた / なんとか卒業できた / なんとか生き延びた\\
\textbf{Variants (EN)}: I managed to get through it somehow / I somehow managed to graduate / I somehow survived

\textbf{Intent (JA)}: 困難克服, 安堵, ぎりぎりの成功\\
\textbf{Intent (EN)}: overcoming hardship, relief, barely succeeding

\textbf{Tags (JA)}: 即時文法, 口語, 克服, 安堵\\
\textbf{Tags (EN)}: immediate grammar, spoken, overcoming, relief

\textbf{Adjusted Expression (JA)}:\\
かろうじて\\
\textbf{Adjusted Expression (EN)}:\\
barely / just managed to

\textbf{Example (JA)}:\\
A「(大変:たいへん)だったね。」 B「うん、なんとか、のりこえたよ。」\\
\textbf{Example (EN)}:\\
A: That must have been tough. B: Yeah, I managed to get through it somehow.

\textbf{Notes (JA)}:\\
「なんとか」は困難をかろうじて克服したときに用いられる即時文法的表現。明確な手段を示さずに『どうにかできた』という安堵の気持ちを伝える。フォーマルな文脈では「かろうじて」「どうにか」などに置き換えられる。

\textbf{Notes (EN)}:\\
'Nantoka' is used as an immediate grammar expression when overcoming difficulties with relief. It avoids specifying the means and simply conveys that the speaker 'somehow managed it'. In formal contexts, it can be replaced with 'karoujite' (barely) or 'dounika' (somehow).

\newpage

\section*{509. How about coming?}

\textbf{Date}: 2025.08.31 (Sun)\\
\textbf{Index}: konai\\
\textbf{Expression (JA)}: (来:こ)ない？\\
\textbf{Romanization}: konai?\\
\textbf{Expression (EN)}: How about coming?

\textbf{Variants (JA)}: 行かない？ / 食べない？ / 飲まない？\\
\textbf{Variants (EN)}: How about going? / How about eating? / How about having a drink?

\textbf{Intent (JA)}: 誘い, 提案, 親しさ\\
\textbf{Intent (EN)}: invitation, suggestion, closeness

\textbf{Tags (JA)}: 即時文法, 口語, 誘い, 否定疑問\\
\textbf{Tags (EN)}: immediate grammar, spoken, invitation, negative question

\textbf{Adjusted Expression (JA)}:\\
(来:き)ませんか？\\
\textbf{Adjusted Expression (EN)}:\\
Would you like to come?

\textbf{Example (JA)}:\\
A「お(昼:ひる)どうする？」 B「(ラーメン:らーめん)(食:た)べに(行:い)かない？」\\
\textbf{Example (EN)}:\\
A: What should we do for lunch? B: How about going for ramen?

\textbf{Notes (JA)}:\\
「V-ない？」は形の上では否定疑問だが、実際には相手への軽い誘いを表す即時文法的な表現。親しい会話で多用され、強制ではなく共感や仲間意識を含む。フォーマルな場面では「V-ませんか？」に言い換えられる。

【日本語と英語の違い】
- 日本語の「行かない？」は、ほぼ自動的に「一緒にしようよ」という勧誘になる。
- 英語の "Not going?" "Not coming?" は、本来「行かないの？」「来ないの？」という確認や驚きの意味で、勧誘にはならない。
- 文脈によって軽い促しに聞こえることはあるが、日本語ほど便利で普遍的な勧誘表現ではない。
- 英語で勧誘するには "Shall we go?" "How about going?" "Do you want to come?" などを使う。

\textbf{Notes (EN)}:\\
Although 'V-nai?' is grammatically a negative question, it functions as an immediate grammar expression to casually invite someone. It is widely used in friendly speech and implies closeness rather than obligation. In formal contexts, it is replaced by 'V-masen ka?'.

[Difference between Japanese and English]
- Japanese 'ikanai?' almost automatically means 'Shall we...?', functioning as a universal invitation.
- English 'Not going?' or 'Not coming?' literally means 'Aren't you going/coming?', mainly expressing confirmation or surprise, not an invitation.
- In some contexts, it can sound like a nudge to go, but it is not a standard or reliable way to invite.
- Natural ways to invite in English are 'Shall we go?', 'How about going?', 'Do you want to come?', etc.

\newpage

\section*{508. If it's alright with you...}

\textbf{Date}: 2025.08.30 (Sat)\\
\textbf{Index}: moshiyokattara\\
\textbf{Expression (JA)}: もし、よかったら、\\
\textbf{Romanization}: moshi, yokattara,\\
\textbf{Expression (EN)}: If it's alright with you...

\textbf{Variants (JA)}: もし、よかったら、一緒に行きませんか？ / もし、よかったら、使ってください。 / もし、よかったら、見てみて。\\
\textbf{Variants (EN)}: If it's alright with you, would you come with me? / If it's alright with you, please use it. / If it's alright with you, have a look.

\textbf{Intent (JA)}: 控えめな提案, 配慮, 誘い\\
\textbf{Intent (EN)}: tentative suggestion, consideration, invitation

\textbf{Tags (JA)}: 即時文法, 口語, 丁寧, 提案\\
\textbf{Tags (EN)}: immediate grammar, spoken, polite, suggestion

\textbf{Adjusted Expression (JA)}:\\
もしよろしければ、\\
\textbf{Adjusted Expression (EN)}:\\
If you don't mind, ...

\textbf{Example (JA)}:\\
A「(宿題:しゅくだい)、むずかしかった？」 B「うん。もし、よかったら、あとで(一緒:いっしょ)にやらない？」\\
\textbf{Example (EN)}:\\
A: Was the homework difficult? B: Yeah. If it's alright with you, want to do it together later?

\textbf{Notes (JA)}:\\
「もし、よかったら、」は、相手の都合や気持ちを尊重しつつ提案するときの即時文法的表現。文末を切ることで柔らかさや遠慮がにじみ出る。フォーマルな場面では「もしよろしければ、」に言い換えられる。

\textbf{Notes (EN)}:\\
The phrase 'moshi, yokattara,' is an immediate grammar expression used when making a suggestion while showing consideration for the listener. The trailing comma adds hesitation and softness. In formal contexts, it is expressed as 'moshi yoroshikereba,'.

\newpage

\section*{507. I never thought that...}

\textbf{Date}: 2025.08.29 (Fri)\\
\textbf{Index}: towaomowanakatta\\
\textbf{Expression (JA)}: ..とは(思:おも)わなかった\\
\textbf{Romanization}: ...to wa omowanakatta\\
\textbf{Expression (EN)}: I never thought that...

\textbf{Variants (JA)}: 勝てるとは思わなかった / 君に会えるとは思わなかった / こんなに時間がかかるとは思わなかった\\
\textbf{Variants (EN)}: I never thought we could win / I never thought I would see you / I never thought it would take this long

\textbf{Intent (JA)}: 意外性, 驚き, 回想\\
\textbf{Intent (EN)}: unexpectedness, surprise, retrospection

\textbf{Tags (JA)}: 即時文法, 口語, 意外性, 感情\\
\textbf{Tags (EN)}: immediate grammar, spoken, unexpectedness, emotion

\textbf{Adjusted Expression (JA)}:\\
...とは(思:)いませんでした\\
\textbf{Adjusted Expression (EN)}:\\
I did not expect that...

\textbf{Example (JA)}:\\
A「(優勝:ゆうしょう)できたね！」 B「うん、まさか(勝:か)てるとは(思:おも)わなかった」\\
\textbf{Example (EN)}:\\
A: We won the championship! B: Yeah, I never thought we could win.

\textbf{Notes (JA)}:\\
「...とは思わなかった」は、予想外の出来事や意外な出会いを表す即時文法的表現。話し手の驚きや感慨を直接的に伝える。フォーマルな場面では「...とは思いませんでした」と言い換えられる。

\textbf{Notes (EN)}:\\
The phrase '...to wa omowanakatta' is an immediate grammar expression showing surprise at an unexpected outcome or encounter. It conveys the speaker's astonishment or reflection. In formal contexts, it can be expressed as '...to wa omoimasen deshita'.

\newpage

\section*{506. Well, but you know...}

\textbf{Date}: 2025.08.28 (Thu)\\
\textbf{Index}: nandakedone\\
\textbf{Expression (JA)}: ...なんだけどね\\
\textbf{Romanization}: ...nan da kedo ne\\
\textbf{Expression (EN)}: Well, but you know...

\textbf{Variants (JA)}: 行きたかったんだけどね / 頑張ったんだけどね / いい試合だったんだけどね\\
\textbf{Variants (EN)}: I wanted to go, but you know... / I tried hard, but you know... / It was a good game, but you know...

\textbf{Intent (JA)}: 事情説明, 言いさし, 共感呼びかけ\\
\textbf{Intent (EN)}: explanation, trailing off, seeking empathy

\textbf{Tags (JA)}: 即時文法, 口語, 含み, 共感\\
\textbf{Tags (EN)}: immediate grammar, spoken, implication, empathy

\textbf{Adjusted Expression (JA)}:\\
...なのですが\\
\textbf{Adjusted Expression (EN)}:\\
However, ...

\textbf{Example (JA)}:\\
A「(試合:しあい)、(負:ま)けちゃったね？」 B「いい(プレイ:ぷれい)はあったんだけどね...」\\
\textbf{Example (EN)}:\\
A: We lost the game, huh? B: There were some good plays, but you know...

\textbf{Notes (JA)}:\\
「...なんだけどね」は、話を途中で切ることで含みを残し、聞き手に理解や共感を促す表現。言い訳や弁解、惜しい気持ちなど、話者の感情がにじむ。フォーマルな場面では「…なのですが」と言い換えられる。

\textbf{Notes (EN)}:\\
The phrase '...nan da kedo ne' leaves the sentence trailing off, inviting the listener to infer the unsaid part. It often conveys excuses, disappointment, or an appeal for empathy. In formal contexts, it can be expressed as '...nano desu ga'.

\newpage

\section*{505. How does it end up like that?}

\textbf{Date}: 2025.08.27 (Wed)\\
\textbf{Index}: nandesounaruno\\
\textbf{Expression (JA)}: なんでそうなるの？\\
\textbf{Romanization}: nande sou naru no?\\
\textbf{Expression (EN)}: Why does it end up like that?

\textbf{Variants (JA)}: なんでそうなる？ / どうしてそうなるの？ / なんでそんなことになるの？\\
\textbf{Variants (EN)}: Why does it end up like that? / How come it turns out that way? / Why does it go that way?

\textbf{Intent (JA)}: 驚き, 不満, 納得できなさ\\
\textbf{Intent (EN)}: surprise, discontent, incredulity

\textbf{Tags (JA)}: 即時文法, 口語, 疑問, 感情\\
\textbf{Tags (EN)}: immediate grammar, spoken, question, emotion

\textbf{Adjusted Expression (JA)}:\\
どうしてそのようなことになるのですか？\\
\textbf{Adjusted Expression (EN)}:\\
Why does it end in that way?

\textbf{Example (JA)}:\\
A「(パソコン:ぱそこん)(直:なお)そうとしたら、もっと(壊:こわ)れた」 B「えっ、なんでそうなるの？」\\
\textbf{Example (EN)}:\\
A: I tried to fix my computer, but it broke even more. B: What? How does it end up like that?

\textbf{Notes (JA)}:\\
「なんでそうなるの？」は、予想外の結果や納得できない状況に対して即時的に発せられる感情的な疑問。純粋な質問ではなく、驚きや不満を含み、相手に説明を求める圧力も持つ。フォーマルな場面では「どうしてそのようなことになるのですか？」のように言い換えられる。

\textbf{Notes (EN)}:\\
"Nande sou naru no?" is an immediate grammar expression used when confronted with an unexpected or unreasonable result. It is less a request for information and more an incredulous reaction that pressures the other person to explain. In formal contexts, it can be rephrased as "Doushite sonoyouna koto ni naru no desu ka?"

\newpage

\section*{504. You should've known if you saw it}

\textbf{Date}: 2025.08.26 (Tue)\\
\textbf{Index}: tetara-hazude\\
\textbf{Expression (JA)}: ...てたら、...はずで\\
\textbf{Romanization}: ...tetara, ...hazude\\
\textbf{Expression (EN)}: If you'd been watching, you should have known

\textbf{Variants (JA)}: 見てたら、わかるはずで / 聞いてたら、知ってるはずで / やってたら、気づくはずで\\
\textbf{Variants (EN)}: If you had seen it, you should have known / If you had heard it, you'd know / If you had tried it, you'd have noticed

\textbf{Intent (JA)}: 根拠提示, 軽い主張, 含み\\
\textbf{Intent (EN)}: grounded assertion, soft criticism, implication

\textbf{Tags (JA)}: 即時文法, 口語, 感情, 根拠\\
\textbf{Tags (EN)}: immediate grammar, spoken, emotion, evidentiality

\textbf{Adjusted Expression (JA)}:\\
...ていたなら、..はずだったのだが\\
\textbf{Adjusted Expression (EN)}:\\
If you had been ..., then you should have ...

\textbf{Example (JA)}:\\
A「え？そんなことあったの？」 B「それは(見:み)てたら、わかるはずで...」\\
\textbf{Example (EN)}:\\
A: Huh? That happened? B: If you were watching, you should've noticed...

\textbf{Notes (JA)}:\\
「...てたら、...はずで」は、話し手が何らかの観察や経験を根拠に「〜は当然である」と判断する表現。文末の「で」で言い切らずに余韻を残すことで、相手に「気づくべきだった」という含意をやわらかく伝える。即時文法として、場面に即応する主観的な評価や不満を含むことが多い。

\textbf{Notes (EN)}:\\
The expression '...tetara, ...hazude' implies that something should have been obvious if the listener had seen or experienced it. Ending with 'de' softens the claim and leaves it open-ended, often suggesting a mild reproach or implicit expectation. It is an example of immediate grammar, expressing the speaker’s real-time judgment.

\newpage

\section*{503. I had no choice but to...}

\textbf{Date}: 2025.08.25 (Mon)\\
\textbf{Index}: shikanakute\\
\textbf{Expression (JA)}: しかなくて\\
\textbf{Romanization}: shika nakute\\
\textbf{Expression (EN)}: I had no choice but to...

\textbf{Variants (JA)}: 行くしかなくて / 待つしかなくて / 謝るしかなくて\\
\textbf{Variants (EN)}: I had no choice but to go / I had no choice but to wait / I had no choice but to apologize

\textbf{Intent (JA)}: 限定, 言い訳, 必然\\
\textbf{Intent (EN)}: limitation, excuse, inevitability

\textbf{Tags (JA)}: 即時文法, 口語, 制約\\
\textbf{Tags (EN)}: immediate grammar, spoken, constraint

\textbf{Adjusted Expression (JA)}:\\
〜する(以外:いがい)に(方法:ほうほう)はなくて\\
\textbf{Adjusted Expression (EN)}:\\
there was no other option but to ...

\textbf{Example (JA)}:\\
A「どうして(遅:おく)れたの？」 B「(電車:でんしゃ)が(止:と)まって、(歩:ある)くしかなくて...」\\
\textbf{Example (EN)}:\\
A: Why were you late? B: The train stopped, so I had no choice but to walk...

\textbf{Notes (JA)}:\\
「しかなくて」は『〜する以外に方法がない』ことを即時的に説明するときに使われる表現です。多くの場合、言い訳や事情説明に使われ、文を途中で切って余韻を残すこともできます。フォーマルな場面では「〜するほかなくて」と表現されることもあります。

\textbf{Notes (EN)}:\\
"Shika nakute" is an immediate grammar expression used to say there was no other option. It often appears in excuses or explanations, and can also trail off to leave the situation implicit. In formal contexts, it can be expressed as "...suru hoka nakute."

\newpage

\section*{502. As soon as}

\textbf{Date}: 2025.08.24 (Sun)\\
\textbf{Index}: tosugu\\
\textbf{Expression (JA)}: ...とすぐ、\\
\textbf{Romanization}: ... to sugu,\\
\textbf{Expression (EN)}: as soon as ...

\textbf{Variants (JA)}: …とたん / ..てすぐに / ..と同時に\\
\textbf{Variants (EN)}: the moment ... / immediately after ...ing / at the same time as ...

\textbf{Intent (JA)}: 即時性, 連鎖, 予期せぬ出来事\\
\textbf{Intent (EN)}: immediacy, sequence, unexpected event

\textbf{Tags (JA)}: 即時文法, 接続, 時間\\
\textbf{Tags (EN)}: immediate grammar, conjunction, time

\textbf{Adjusted Expression (JA)}:\\
...と(間髪:かんぱつ)(入:い)れず\\
\textbf{Adjusted Expression (EN)}:\\
... without a moment's delay

\textbf{Example (JA)}:\\
(玄関:げんかん)を(開:あ)けるとすぐ、(猫:ねこ)が(飛:と)びかかってきた。\\
\textbf{Example (EN)}:\\
As soon as I opened the door, the cat jumped at me.

\textbf{Notes (JA)}:\\
「…とすぐ、」は前件が起こった直後に、ほとんど間をおかず後件が起こることを表す即時文法の接続表現です。単なる条件ではなく、前件と後件が連続して必然的に結びつくことを示します。英語の "then" に近く、会話でも文章でも自然に使われます。

\textbf{Notes (EN)}:\\
"... to sugu," is an immediate grammar conjunction indicating that the second event occurs almost immediately after the first. It signifies a close, inevitable connection between the two events, similar to "then" in English. It is natural in both spoken and written Japanese.

\newpage

\section*{501. In that case}

\textbf{Date}: 2025.08.23 (Sat)\\
\textbf{Index}: sorenara\\
\textbf{Expression (JA)}: それなら\\
\textbf{Romanization}: sore nara\\
\textbf{Expression (EN)}: In that case

\textbf{Variants (JA)}: それなら行こう / それなら大丈夫 / それならやめよう\\
\textbf{Variants (EN)}: In that case, let's go / Then it's fine / In that case, let's stop

\textbf{Intent (JA)}: 条件転換, 提案, 同意\\
\textbf{Intent (EN)}: conditional shift, suggestion, agreement

\textbf{Tags (JA)}: 即時文法, 口語, 条件\\
\textbf{Tags (EN)}: immediate grammar, spoken, condition

\textbf{Adjusted Expression (JA)}:\\
そうなら\\
\textbf{Adjusted Expression (EN)}:\\
If that is so

\textbf{Example (JA)}:\\
A「(今日:きょう)は(雨:あめ)みたいだよ」 B「それなら、(家:いえ)で(映画:えいが)を(見:み)よう」\\
\textbf{Example (EN)}:\\
A: Looks like it's raining today. B: In that case, let's watch a movie at home.

\textbf{Notes (JA)}:\\
「それなら」は相手の発言や状況を受けて、その条件に合わせて行動や判断を即時的に切り替える表現です。日常会話では提案や同意の導入によく用いられ、柔らかく相手に合わせるニュアンスがあります。学習者は「名詞句＋なら」の条件文として難しく考えがちですが、実際には英語の "then" に近い感覚で、単純に会話をつなぐ即時反応のことばとして理解すれば十分です。

\textbf{Notes (EN)}:\\
"Sore nara" is an immediate grammar expression used to shift actions or decisions according to the given condition. In daily conversation, it often introduces suggestions or agreement, carrying a nuance of adapting to the other person's statement or the situation. Learners sometimes overanalyze it as a nominal conditional, but in practice it works very much like the English "then", a simple, immediate connector.

\newpage

\section*{500. Shall I...?}

\textbf{Date}: 2025.08.22 (Fri)\\
\textbf{Index}: shiyouka\\offer\\
\textbf{Expression (JA)}: しようか\\
\textbf{Romanization}: shiyou ka\\
\textbf{Expression (EN)}: Shall I...?

\textbf{Variants (JA)}: 手伝おうか / 持ってこようか / 開けようか\\
\textbf{Variants (EN)}: Shall I help? / Shall I bring it? / Shall I open it?

\textbf{Intent (JA)}: 申し出, 配慮, 気遣い\\
\textbf{Intent (EN)}: offer, consideration, thoughtfulness

\textbf{Tags (JA)}: 即時文法, 口語, 申し出\\
\textbf{Tags (EN)}: immediate grammar, spoken, offer

\textbf{Adjusted Expression (JA)}:\\
いたしましょうか\\
\textbf{Adjusted Expression (EN)}:\\
Would you like me to...?

\textbf{Example (JA)}:\\
A「(荷物:にもつ)、(重:おも)そうだね、(持:も)とうか？」\\
\textbf{Example (EN)}:\\
A: Your bag looks heavy. You want me to carry it?

\textbf{Notes (JA)}:\\
「しようか」はここでは相手のために自分が行動を申し出る即時文法の表現です。「〜しませんか」のような共同行動の誘いではなく、相手への気遣いを示す用法です。フォーマルな場では「いたしましょうか」と表現されます。

\textbf{Notes (EN)}:\\
Here, "Shiyou ka" is an immediate grammar expression used to offer to do something for the listener. Unlike the invitational use ("Shall we...?"), it expresses consideration or willingness to help. In formal settings, it appears as "Itashimashou ka".

\newpage

\section*{499. Unexpectedly}

\textbf{Date}: 2025.08.21 (Thu)\\
\textbf{Index}: omoigakezu\\
\textbf{Expression (JA)}: (思:おも)いがけず\\
\textbf{Romanization}: omoigakezu\\
\textbf{Expression (EN)}: unexpectedly

\textbf{Variants (JA)}: 思いがけず会った / 思いがけず助かった / 思いがけず褒められた\\
\textbf{Variants (EN)}: I met (someone) unexpectedly / I was saved unexpectedly / I was praised unexpectedly

\textbf{Intent (JA)}: 予想外, 偶然, 驚き\\
\textbf{Intent (EN)}: unexpected, coincidence, surprise

\textbf{Tags (JA)}: 即時文法, 副詞, 予期せぬ出来事\\
\textbf{Tags (EN)}: immediate grammar, adverb, unexpected event

\textbf{Adjusted Expression (JA)}:\\
(予想外:よそうがい)に\\
\textbf{Adjusted Expression (EN)}:\\
unpredictably

\textbf{Example (JA)}:\\
思いがけず、(駅:えき)で(先生:せんせい)に(会:あ)ったよ。\\
\textbf{Example (EN)}:\\
Unexpectedly, I met my teacher.

\textbf{Notes (JA)}:\\
「思いがけず」は予想していなかった出来事が起こったときに即時的に使われる副詞です。ポジティブにもネガティブにも使え、偶然性や意外性を表す柔らかい表現です。文語にも日常会話にも自然に使われます。

\textbf{Notes (EN)}:\\
"Omoigakezu" is an adverb used when something happens unexpectedly. It can describe both positive and negative events, expressing coincidence or surprise. It is natural in both spoken and written Japanese.

\newpage

\section*{498. The moment}

\textbf{Date}: 2025.08.20 (Wed)\\
\textbf{Index}: totan\\
\textbf{Expression (JA)}: とたん\\
\textbf{Romanization}: totan\\
\textbf{Expression (EN)}: the moment

\textbf{Variants (JA)}: 立ったとたん / 言ったとたん / 出かけたとたん\\
\textbf{Variants (EN)}: the moment (I) stood up / the moment (I) said it / the moment (I) left

\textbf{Intent (JA)}: 即時性, 転換, 予期せぬ出来事\\
\textbf{Intent (EN)}: immediacy, change, unexpected occurrence

\textbf{Tags (JA)}: 即時文法, 接続, 時間\\
\textbf{Tags (EN)}: immediate grammar, conjunction, time

\textbf{Adjusted Expression (JA)}:\\
〜した(瞬間:しゅんかん)\\
\textbf{Adjusted Expression (EN)}:\\
the instant ... happened

\textbf{Example (JA)}:\\
A「(外:そと)に(出:で)たとたん、(雨:あめ)が(降:ふ)り(出:だ)したよ」 B「うわ、(タイミング:たいみんぐ)(悪:わる)いね」\\
\textbf{Example (EN)}:\\
A: The moment I stepped outside, it started raining. B: Wow, bad timing.

\textbf{Notes (JA)}:\\
「とたん」は出来事が起こった直後に、別のことが起こる即時的な接続表現です。予期しない変化や突然の出来事を述べるのに多く用いられます。形式としては「動詞た形＋とたん（に）」で用いられます。日常会話でも書き言葉でも広く使われます。

\textbf{Notes (EN)}:\\
"Totan" is an immediate grammar conjunction indicating that one event occurs immediately after another. It is often used for unexpected or sudden changes. The typical form is "verb in ta-form + totan (ni)". It is common in both spoken and written Japanese.

\newpage

\section*{497. With all my might}

\textbf{Date}: 2025.08.19 (Tue)\\
\textbf{Index}: omoikkiri\\
\textbf{Expression (JA)}: (思:おも)いっきり\\
\textbf{Romanization}: omoikkiri\\
\textbf{Expression (EN)}: with all my might

\textbf{Variants (JA)}: 思いっきり走る / 思いっきり笑う / 思いっきり怒る\\
\textbf{Variants (EN)}: run with all my might / laugh out loud / get really angry

\textbf{Intent (JA)}: 強調, 感情, 全力\\
\textbf{Intent (EN)}: emphasis, emotion, full force

\textbf{Tags (JA)}: 即時文法, 口語, 強調\\
\textbf{Tags (EN)}: immediate grammar, spoken, emphasis

\textbf{Adjusted Expression (JA)}:\\
(力:ちから)いっぱい\\
\textbf{Adjusted Expression (EN)}:\\
to the fullest

\textbf{Example (JA)}:\\
A「どうやって(投:な)げるの？」 B「(思:おも)いっきり(投:な)げて！」\\
\textbf{Example (EN)}:\\
A: How should I throw it? B: Throw it with all your might!

\textbf{Notes (JA)}:\\
「思いっきり」は気持ちや力を出し切ることを即時的に表す副詞です。行動に感情を込めたいときに使われ、強さ・激しさ・熱中を伴います。文脈によってポジティブ（思いっきり楽しむ）にもネガティブ（思いっきり叱る）にも使われます。

\textbf{Notes (EN)}:\\
"Omoikkiri" is an adverb expressing doing something with full force or emotion. It conveys intensity, passion, or wholeheartedness. Depending on context, it can be positive (enjoy to the fullest) or negative (scold severely).

\newpage

\section*{496. Likewise}

\textbf{Date}: 2025.08.18 (Mon)\\
\textbf{Index}: kochirakoso\\
\textbf{Expression (JA)}: こちらこそ\\
\textbf{Romanization}: kochirakoso\\
\textbf{Expression (EN)}: Likewise

\textbf{Variants (JA)}: こちらこそ、ありがとう / こちらこそ、すみません / こちらこそ、よろしく\\
\textbf{Variants (EN)}: Likewise, thank you / No, I'm the one who should apologize / Likewise, nice to meet you

\textbf{Intent (JA)}: 応答, 謙遜, 相互性\\
\textbf{Intent (EN)}: response, modesty, reciprocity

\textbf{Tags (JA)}: 即時文法, 口語, 応答\\
\textbf{Tags (EN)}: immediate grammar, spoken, response

\textbf{Adjusted Expression (JA)}:\\
(私:わたし)の(方:ほう)こそ\\
\textbf{Adjusted Expression (EN)}:\\
It is I who should say so

\textbf{Example (JA)}:\\
A「ありがとう！」 B「こちらこそ、ありがとう」\\
\textbf{Example (EN)}:\\
A: Thank you! B: Likewise, thank you.

\textbf{Notes (JA)}:\\
「こちらこそ」は相手の感謝・謝罪・挨拶に対して、謙遜や相互性を示す即時文法の表現です。文脈に応じて「ありがとう」「すみません」「よろしく」に続けて使うことが多く、相手との関係を和やかにする働きを持ちます。

\textbf{Notes (EN)}:\\
"Kochira koso" is an immediate grammar expression used in response to gratitude, apology, or greetings. It conveys modesty and reciprocity, often followed by phrases like "arigatou" (thank you), "sumimasen" (sorry), or "yoroshiku" (pleased to meet you). It helps smooth and warm interpersonal interactions.

\newpage

\section*{495. See?}

\textbf{Date}: 2025.08.17 (Sun)\\
\textbf{Index}: horane\\
\textbf{Expression (JA)}: ほらね\\
\textbf{Romanization}: hora ne\\
\textbf{Expression (EN)}: See?

\textbf{Variants (JA)}: ほらね！ / ほらね、言ったでしょ / ほらね、だから言ったんだ\\
\textbf{Variants (EN)}: See! / See, I told you / See, that's why I said it

\textbf{Intent (JA)}: 確認, 共感, 得意\\
\textbf{Intent (EN)}: confirmation, shared feeling, self-assertion

\textbf{Tags (JA)}: 即時文法, 口語, 確認\\
\textbf{Tags (EN)}: immediate grammar, spoken, confirmation

\textbf{Adjusted Expression (JA)}:\\
(言:い)った(通:とお)りでしょう\\
\textbf{Adjusted Expression (EN)}:\\
As I told you

\textbf{Example (JA)}:\\
A「やっぱり(雨:あめ)が(降:ふ)ってきたね」 B「ほらね」\\
\textbf{Example (EN)}:\\
A: It started raining after all. B: See?

\textbf{Notes (JA)}:\\
「ほらね」は相手に自分の予想や発言が当たったことを示す即時文法表現です。共感的に使えば仲間意識を強め、得意げに使えば自己主張が強くなります。声の調子や文脈によってニュアンスが大きく変わります。

\textbf{Notes (EN)}:\\
"Hora ne" is an immediate grammar expression used to point out that one's prediction or statement has come true. Used warmly, it reinforces a sense of shared understanding, while in a proud tone it emphasizes self-assertion. Its nuance depends heavily on intonation and context.

\newpage

\section*{494. Since we're at it...}

\textbf{Date}: 2025.08.16 (Sat)\\
\textbf{Index}: dousenara\\
\textbf{Expression (JA)}: どうせなら\\
\textbf{Romanization}: douse nara\\
\textbf{Expression (EN)}: Since we're at it...

\textbf{Variants (JA)}: どうせなら行こう / どうせならやってみよう / どうせなら言えばいい\\
\textbf{Variants (EN)}: Since we're at it, let's go / Since we're at it, let's try / Since we're at it, you might as well say it

\textbf{Intent (JA)}: 提案, 気分, 即時判断\\
\textbf{Intent (EN)}: suggestion, feeling, immediate judgment

\textbf{Tags (JA)}: 即時文法, 口語, 提案\\
\textbf{Tags (EN)}: immediate grammar, spoken, suggestion

\textbf{Adjusted Expression (JA)}:\\
せっかくなら\\
\textbf{Adjusted Expression (EN)}:\\
Might as well

\textbf{Example (JA)}:\\
A「もう(遅:おそ)いよ」 B「どうせなら、(泊:と)まっていこうよ」\\
\textbf{Example (EN)}:\\
A: It's already late. B: Since we're at it, we might as well stay over.

\textbf{Notes (JA)}:\\
「どうせなら」は、その場の状況を前提にして行動や提案をする即時文法の表現です。「せっかくなら」「いっそ」と近いニュアンスがあり、日常会話では気分や軽い判断を反映する柔らかい響きを持ちます。選択肢の中でより良い方を選ぶときや、消極的な状況を前向きに転換するときに使われます。

\textbf{Notes (EN)}:\\
"Douse nara" is an immediate grammar expression used to make a suggestion or decision based on the current situation. Similar in nuance to "sekkaku nara" (since we have the chance) or "isso" (might as well), it carries a light and flexible tone in conversation. Often used to pick the better option among available ones or to turn a passive situation into a positive action.

\newpage

\section*{493. Got a moment now?}

\textbf{Date}: 2025.08.15 (Fri)\\
\textbf{Index}: imachottoii\\
\textbf{Expression (JA)}: (今:いま)、ちょっと、いい\\
\textbf{Romanization}: ima, chotto, ii\\
\textbf{Expression (EN)}: Got a moment now?

\textbf{Variants (JA)}: 今、ちょっといいですか / 今、いい？ / 今、少し、いい？\\
\textbf{Variants (EN)}: Do you have a moment now? / Is now a good time? / Got a sec?

\textbf{Intent (JA)}: 呼びかけ, 配慮, 会話開始\\
\textbf{Intent (EN)}: attention call, consideration, conversation start

\textbf{Tags (JA)}: 即時文法, 口語, 配慮\\
\textbf{Tags (EN)}: immediate grammar, spoken, consideration

\textbf{Adjusted Expression (JA)}:\\
(今:いま)、(少:すこ)し、(よろ)しいですか。\\
\textbf{Adjusted Expression (EN)}:\\
May I have a moment now?

\textbf{Example (JA)}:\\
A「(今:いま)、ちょっと、いい？」 B「うん、どうしたの？」\\
\textbf{Example (EN)}:\\
A: Got a moment now? B: Yeah, what's up?

\textbf{Notes (JA)}:\\
「今、ちょっと、いい」は相手の都合をうかがいながら会話を始める即時文法の呼びかけ表現です。文末の「いい」は省略的で、疑問形の助詞や丁寧語を加えることで柔らかさや丁寧さを調整できます。日常会話では相手が答えやすい間を作るために、文中で小さな間を置くことが多いです。英語でいう"I have a quick question."のように簡単な質問をするときにも使われます。

\textbf{Notes (EN)}:\\
"Ima, chotto, ii" is an immediate grammar expression used to start a conversation while checking if the other person is available. The final "ii" is abbreviated, and adding a question particle or polite form can adjust the tone. In casual speech, brief pauses between words create space for the listener to respond. It's similar to saying "I have a quick question" in English, used for simple inquiries.

\newpage

\section*{492. It's nothing.}

\textbf{Date}: 2025.08.14 (Thu)\\
\textbf{Index}: nandemonai\\
\textbf{Expression (JA)}: なんでもない\\
\textbf{Romanization}: nandemo nai\\
\textbf{Expression (EN)}: It's nothing.

\textbf{Variants (JA)}: 全然なんでもない / 別になんでもない / なんでもないよ\\
\textbf{Variants (EN)}: It's totally nothing / It's really nothing / It's nothing, really

\textbf{Intent (JA)}: 否定, 安心, ごまかし\\
\textbf{Intent (EN)}: denial, reassurance, evasion

\textbf{Tags (JA)}: 即時文法, 口語, 感情\\
\textbf{Tags (EN)}: immediate grammar, spoken, emotion

\textbf{Adjusted Expression (JA)}:\\
(大:たい)したことはない\\
\textbf{Adjusted Expression (EN)}:\\
It's not a big deal

\textbf{Example (JA)}:\\
A「どうしたの？」 B「なんでもない」\\
\textbf{Example (EN)}:\\
A: What's wrong? B: It's nothing.

\textbf{Notes (JA)}:\\
「なんでもない」は特別な理由や問題がないことを即時的に表す口語表現です。安心させるためにも、ごまかすためにも使われ、文脈によってニュアンスが変わります。語気を柔らかくすると親しみが増し、強く言うと突き放す印象になります。また、発言を軽く取り消すときにも使われます。

\textbf{Notes (EN)}:\\
"Nandemo nai" is a colloquial expression used to state immediately that there is no particular reason or problem. It can be used to reassure or to deflect, with nuances changing depending on context. A softer tone adds warmth, while a stronger tone can feel dismissive. It can also be used to lightly retract a statement such as "Never mind."

\newpage

\section*{491. I'm glad I did it.}

\textbf{Date}: 2025.08.13 (Wed)\\
\textbf{Index}: yatteteyokatta\\
\textbf{Expression (JA)}: やっててよかった\\
\textbf{Romanization}: yattete yokatta\\
\textbf{Expression (EN)}: I'm glad I did it.

\textbf{Variants (JA)}: 行っててよかった / 買っててよかった / 聞いててよかった\\
\textbf{Variants (EN)}: I'm glad I went / I'm glad I bought it / I'm glad I heard it

\textbf{Intent (JA)}: 満足, 安心, 即時評価\\
\textbf{Intent (EN)}: satisfaction, relief, immediate evaluation

\textbf{Tags (JA)}: 即時文法, 口語, 満足\\
\textbf{Tags (EN)}: immediate grammar, spoken, satisfaction

\textbf{Adjusted Expression (JA)}:\\
やっていてよかった\\
\textbf{Adjusted Expression (EN)}:\\
I'm glad I have done it

\textbf{Example (JA)}:\\
A「(事前:じぜん)に(予約:よやく)してた？」 B「うん、やっててよかったよ」\\
\textbf{Example (EN)}:\\
A: Did you book it in advance? B: Yeah, I'm glad I did.

\textbf{Notes (JA)}:\\
「やっててよかった」は「やっていてよかった」の口語縮約形で、事前に行った行動が結果的に良かったと感じるときに使います。動詞のて形＋いる（継続）＋て＋よかった（良かった）から構成されます。日常会話では「…ていて」が「…てて」と縮まり、感情のこもった即時評価として機能します。

\textbf{Notes (EN)}:\\
"Yattete yokatta" is a colloquial contraction of "yatte ite yokatta", used when an action taken beforehand turns out to have been beneficial. It is formed from the te-form of a verb + iru (continuative) + te + yokatta (it was good). In casual speech, "...te ite" shortens to "...tete", giving an immediate and emotional evaluation.

\newpage

\section*{490. I should've done it.}

\textbf{Date}: 2025.08.12 (Tue)\\
\textbf{Index}: yattokya\\
\textbf{Expression (JA)}: やっときゃよかった\\
\textbf{Romanization}: yattokya yokatta\\
\textbf{Expression (EN)}: I should've done it.

\textbf{Variants (JA)}: 行っときゃよかった / 買っときゃよかった / 言っときゃよかった\\
\textbf{Variants (EN)}: I should've gone / I should've bought it / I should've said it

\textbf{Intent (JA)}: 後悔, 未実行, 即時評価\\
\textbf{Intent (EN)}: regret, not done, immediate evaluation

\textbf{Tags (JA)}: 即時文法, 口語, 後悔\\
\textbf{Tags (EN)}: immediate grammar, spoken, regret

\textbf{Adjusted Expression (JA)}:\\
やっておけばよかった\\
\textbf{Adjusted Expression (EN)}:\\
I should have done it

\textbf{Example (JA)}:\\
A 「あと2(点:てん)(足:た)りませんでした」B「もうちょっと(勉強:べんきょう)、やっときゃよかった」\\
\textbf{Example (EN)}:\\
A: I didn't pass by 2 points. B: I should have studied a bit more.

\textbf{Notes (JA)}:\\
「やっときゃよかった」は「やっておけばよかった」の口語縮約形で、やらなかったことへの後悔を即時的に表す表現です。動詞のて形＋おく（事前に行う）＋ば（仮定）＋よかった（良かったのに）から構成されます。日常会話では「おく」が縮まって「とく」になり、「〜ときゃ」という形で使われます。

\textbf{Notes (EN)}:\\
"Yattokya yokatta" is a colloquial contraction of "yatte okeba yokatta", expressing immediate regret for not having done something. It is formed from the te-form of a verb + oku (to do in advance) + ba (conditional) + yokatta (it would have been good). In casual speech, "oku" shortens to "toku", and the expression appears as "-tokya".

\newpage

\section*{489. I'll do it in advance.}

\textbf{Date}: 2025.08.11 (Mon)\\
\textbf{Index}: toku\\
\textbf{Expression (JA)}: ...とく\\
\textbf{Romanization}: ... toku\\
\textbf{Expression (EN)}: (do something) in advance

\textbf{Variants (JA)}: これ、買っとく / 先に言っとく / 用意しとく\\
\textbf{Variants (EN)}: I'll buy this in advance / I'll tell you first / I'll prepare it beforehand

\textbf{Intent (JA)}: 準備, 配慮, 先回り\\
\textbf{Intent (EN)}: preparation, consideration, anticipation

\textbf{Tags (JA)}: 即時文法, 口語, 話しことば\\
\textbf{Tags (EN)}: immediate grammar, spoken, colloquial

\textbf{Adjusted Expression (JA)}:\\
...ておく\\
\textbf{Adjusted Expression (EN)}:\\
do (something) beforehand

\textbf{Example (JA)}:\\
A「(飲:の)み(物:もの)、どうする？」 B「あとで、(買:か)っとくよ」\\
\textbf{Example (EN)}:\\
A: What about drinks? B: I'll buy them later.

\textbf{Notes (JA)}:\\
「...とく」は「...ておく」の省略形で、事前の準備や先回りの行動を即時的に表す口語表現です。日常会話では動詞のて形に直接付けて使われます。フォーマルな場では省略せず「...ておきましょう」を用います。

\textbf{Notes (EN)}:\\
"... toku" is a colloquial contraction of "...te oku", indicating doing something in advance or in preparation. Common in casual speech, it attaches directly to the -te form of verbs. In formal contexts, the full form "...te okimashou" is used.

\newpage

\section*{488. I'm waiting for you.}

\textbf{Date}: 2025.08.10 (Sun)\\
\textbf{Index}: matteruyo\\
\textbf{Expression (JA)}: (待:ま)ってるよ\\
\textbf{Romanization}: matteru yo\\
\textbf{Expression (EN)}: I'm waiting for you

\textbf{Variants (JA)}: 待ってるよね / 待ってるよー / 待ってるよん\\
\textbf{Variants (EN)}: I'll be waiting / Waiting for you / You know I'll wait

\textbf{Intent (JA)}: 安心, 期待, やさしさ\\
\textbf{Intent (EN)}: reassurance, expectation, kindness

\textbf{Tags (JA)}: 即時文法, 感情, 話しことば\\
\textbf{Tags (EN)}: immediate grammar, emotion, spoken language

\textbf{Adjusted Expression (JA)}:\\
お(待:ま)ちしております\\
\textbf{Adjusted Expression (EN)}:\\
I will be waiting for you

\textbf{Example (JA)}:\\
A「(先:さき)に(行:い)ってて」 B「うん、(待:ま)ってるよ」\\
\textbf{Example (EN)}:\\
A: Go ahead without me. B: Sure, I'll be waiting for you.

\textbf{Notes (JA)}:\\
「待ってるよ」は相手に安心感を与える即時文法の表現です。単独で使用でき、日常会話では語尾を伸ばしたり、ね・よねを付けて親しみを加えることもあります。フォーマルな場では「お待ちしております」と言い換えます。

\textbf{Notes (EN)}:\\
"Matteru yo" reassures the listener that you will wait for them. It stands alone as a complete utterance, and in casual conversation, it can be softened with elongated endings or particles like 'ne' or 'yone'. In formal settings, it is replaced with 'O-machi shite orimasu'.

\newpage

\section*{487. More}

\textbf{Date}: 2025.08.09 (Sat)\\
\textbf{Index}: motto\\
\textbf{Expression (JA)}: もっと\\
\textbf{Romanization}: motto\\
\textbf{Expression (EN)}: more

\textbf{Variants (JA)}: もっと食べたい / もっと見せて / もっと早く\\
\textbf{Variants (EN)}: I want more / Show me more / Faster

\textbf{Intent (JA)}: 追加, 強調, 要求\\
\textbf{Intent (EN)}: addition, emphasis, request

\textbf{Tags (JA)}: 即時文法, 感情, 話しことば\\
\textbf{Tags (EN)}: immediate grammar, emotion, spoken language

\textbf{Adjusted Expression (JA)}:\\
さらに\\
\textbf{Adjusted Expression (EN)}:\\
more

\textbf{Example (JA)}:\\
A「もっと(食:た)べる？」 B「うん！」\\
\textbf{Example (EN)}:\\
A: Do you want more to eat? B: Yeah!

\textbf{Notes (JA)}:\\
「もっと」は数量や程度を強めたいときに即時的に使える副詞です。単独でも応答として成立し、感情や欲求をそのまま表せます。文脈によって名詞・動詞・形容詞などにかかります。

\textbf{Notes (EN)}:\\
"Motto" is an adverb used to indicate wanting more of something, either in quantity or degree. It can stand alone as an immediate response, directly expressing emotion or desire, and can modify nouns, verbs, or adjectives depending on context.

\newpage

\section*{486. I hear that a lot.}

\textbf{Date}: 2025.08.08 (Fri)\\
\textbf{Index}: yokukiku\\
\textbf{Expression (JA)}: よく(聞:き)く\\
\textbf{Romanization}: yoku kiku\\
\textbf{Expression (EN)}: I often hear that

\textbf{Variants (JA)}: よく聞くね / よく聞くやつ / よく聞く話\\
\textbf{Variants (EN)}: You hear that a lot / Common phrase / I’ve heard that before

\textbf{Intent (JA)}: 共感, 観察, 自然な反応\\
\textbf{Intent (EN)}: shared feeling, observation, natural reaction

\textbf{Tags (JA)}: 即時文法, 感情, 話しことば\\
\textbf{Tags (EN)}: immediate grammar, emotion, spoken language

\textbf{Adjusted Expression (JA)}:\\
(頻繁:ひんぱん)に(耳:みみ)にします。\\
\textbf{Adjusted Expression (EN)}:\\
I hear it frequently.

\textbf{Example (JA)}:\\
A「こんな(話:はなし)、よく(聞:き)くよね」 B「うん、たしかに」\\
\textbf{Example (EN)}:\\
A: We hear this kind of story a lot, right? B: Yeah, definitely.

\textbf{Notes (JA)}:\\
「よく聞く」は、よく耳にすることばや話題に対して反応する即時文法の表現です。「よく聞くね」のように「ね」をつけることで、共感や共有のニュアンスが加わります。会話をつなげる導入としても自然です。

\textbf{Notes (EN)}:\\
"Yoku kiku" is used to react to things you've often heard. It expresses recognition and connection with familiar topics. Adding "ne" (as in "Yoku kiku ne") adds warmth and shared feeling, making it useful in casual conversation.

\newpage

\section*{485. I see that a lot.}

\textbf{Date}: 2025.08.07 (Thu)\\
\textbf{Index}: yokumiru\\
\textbf{Expression (JA)}: よく(見:み)る\\
\textbf{Romanization}: yoku miru\\
\textbf{Expression (EN)}: I often see that

\textbf{Variants (JA)}: よく見るね / よく見るやつ / よく見る光景\\
\textbf{Variants (EN)}: You see that a lot / Common sight / I keep seeing that

\textbf{Intent (JA)}: 観察, 親しみ, 共感\\
\textbf{Intent (EN)}: observation, familiarity, shared feeling

\textbf{Tags (JA)}: 即時文法, 口語, 共感\\
\textbf{Tags (EN)}: immediate grammar, spoken, shared impression

\textbf{Adjusted Expression (JA)}:\\
(頻繁:ひんぱん)に(見:み)かけます。\\
\textbf{Adjusted Expression (EN)}:\\
I see it frequently.

\textbf{Example (JA)}:\\
A「こんな(写真:しゃしん)、よく(見:み)るよね」 B「うん、たしかに」\\
\textbf{Example (EN)}:\\
A: You see this kind of photo a lot, right? B: Yeah, definitely.

\textbf{Notes (JA)}:\\
「よく見る」は、日常的に目にするものや経験を共有するときに使われる即時的な観察表現です。「よく見るね」には共感や親しみがこもることが多く、特に誰かとの話のきっかけやつながりに役立ちます。

\textbf{Notes (EN)}:\\
"Yoku miru" means to frequently see something. It's often used to comment on common or shared experiences, and when used in conversation, it creates a bond over mutual recognition. Adding "ne" (as in "yoku miru ne") can soften the tone and enhance shared feeling.

\newpage

\section*{484. Feeling reluctant / Heavy-hearted}

\textbf{Date}: 2025.08.06 (Wed)\\
\textbf{Index}: kigaomoi\\
\textbf{Expression (JA)}: (気:き)が(重:おも)い\\
\textbf{Romanization}: Ki ga omoi\\
\textbf{Expression (EN)}: Feeling reluctant / Heavy-hearted

\textbf{Variants (JA)}: なんか気が重い / 明日が気が重い / 正直、気が重い\\
\textbf{Variants (EN)}: I feel reluctant / I'm dreading tomorrow / Honestly, I'm not looking forward to it

\textbf{Intent (JA)}: 憂鬱, 負担感, ためらい\\
\textbf{Intent (EN)}: gloominess, sense of burden, reluctance

\textbf{Tags (JA)}: 即時文法, 感情, 比喩\\
\textbf{Tags (EN)}: immediate grammar, emotion, metaphor

\textbf{Adjusted Expression (JA)}:\\
(気持:きも)ちが(沈:しず)んでいます。\\
\textbf{Adjusted Expression (EN)}:\\
I’m feeling down.

\textbf{Example (JA)}:\\
A「(明日:あす)、(会議:かいぎ)だよね」 B「そうねぇ、(正直:しょうじき)、(気:き)が(重:おも)い」\\
\textbf{Example (EN)}:\\
A: We have a meeting tomorrow, right? B: Yeah... honestly, I'm not looking forward to it.

\textbf{Notes (JA)}:\\
「気が重い」は、やらなければならないことに対して気分が沈んでいる様子を表します。身体的な重さではなく、心理的な負担を比喩的に「重い」と言い換えた表現で、日常会話ではため息や間を伴って使われることが多いです。即時文法としては、状況に即して感情をそのまま口に出す形です。

\textbf{Notes (EN)}:\\
"Ki ga omoi" literally means "the spirit feels heavy" and figuratively expresses reluctance or emotional burden toward an upcoming task or situation. It's a metaphorical phrase often paired with sighs or pauses, and in immediate grammar, it conveys an unfiltered emotional state tied to the moment.

\newpage

\section*{483. At least that's better}

\textbf{Date}: 2025.08.05 (Tue)\\
\textbf{Index}: zuttomashi\\
\textbf{Expression (JA)}: ずっとまし\\
\textbf{Romanization}: Zutto mashi\\
\textbf{Expression (EN)}: Way better / Much better (comparatively)

\textbf{Variants (JA)}: まだまし / こっちの方がまし / 〜よりずっとまし\\
\textbf{Variants (EN)}: still better / this one’s better at least / much better than \\textasciitilde{}

\textbf{Intent (JA)}: 比較, 妥協, 否定の中の肯定\\
\textbf{Intent (EN)}: comparison, consolation, reluctant approval

\textbf{Tags (JA)}: 即時文法, 話しことば, 感情\\
\textbf{Tags (EN)}: immediate grammar, spoken language, emotion

\textbf{Adjusted Expression (JA)}:\\
こちらの(方:ほう)がはるかに(良:よ)いです。\\
\textbf{Adjusted Expression (EN)}:\\
This one is far better.

\textbf{Example (JA)}:\\
A「こっちも(失敗:しっぱい)したけど」 B「でも、(前:まえ)よりはずっとましじゃん」\\
\textbf{Example (EN)}:\\
A: This one failed too, though... B: But still, it's way better than before.

\textbf{Notes (JA)}:\\
「ずっとまし」は、完全な満足ではないが、比較の中で相対的に「良い」と言いたいときに使われる表現です。否定的な状況の中に小さな肯定を見出す、話しことばらしいやわらかい響きを持ちます。丁寧な文体ではあまり使われず、自然な会話で親しみや共感を込めるために用いられます。

\textbf{Notes (EN)}:\\
"Zutto mashi" is a casual way to express that something is 'way better' in comparison, especially in less-than-ideal situations. It reflects a soft tone of consolation or reluctant approval. It's rarely used in formal writing but very common in spontaneous, empathetic speech.

\newpage

\section*{482. Step by step}

\textbf{Date}: 2025.08.04 (Mon)\\
\textbf{Index}: ippoippo\\
\textbf{Expression (JA)}: (一歩:いっぽ)(一歩:いっぽ)\\
\textbf{Romanization}: Ippo ippo\\
\textbf{Expression (EN)}: Step by step

\textbf{Variants (JA)}: 一歩ずつ / 少しずつ / ゆっくりと\\
\textbf{Variants (EN)}: One step at a time / Bit by bit / Slowly but surely

\textbf{Intent (JA)}: 進行, 努力, 忍耐\\
\textbf{Intent (EN)}: progress, effort, perseverance

\textbf{Tags (JA)}: 即時文法, 反復, 感情\\
\textbf{Tags (EN)}: immediate grammar, repetition, emotion

\textbf{Adjusted Expression (JA)}:\\
(段階:だんかい)を(追:お)って(着実:ちゃくじつ)に(進:すす)んでいます。\\
\textbf{Adjusted Expression (EN)}:\\
I'm making steady progress, step by step.

\textbf{Example (JA)}:\\
A「すぐにはうまくならないけど...」 B「うん、(一歩:いっぽ)(一歩:いっぽ)、やっていこう」\\
\textbf{Example (EN)}:\\
A: I won't get good at it right away... B: Yeah, let's take it step by step.

\textbf{Notes (JA)}:\\
「一歩一歩」は、あせらずに着実に進もうとする気持ちを表します。繰り返しの表現によって、励ましや自分への言い聞かせのニュアンスが生まれます。声に出して言うことで、今この瞬間の決意や落ち着きを伝える即時的な表現となります。

\textbf{Notes (EN)}:\\
"Ippo ippo" conveys the idea of making calm and steady progress. The repetition emphasizes patience and encouragement, often spoken aloud as a way of reinforcing determination or calming oneself. It's an immediate expression of commitment to keep going, little by little.

\newpage

\section*{481. I wonder if this is okay...}

\textbf{Date}: 2025.08.03 (Sun)\\
\textbf{Index}: koredeiinokana\\
\textbf{Expression (JA)}: これでいいのかな\\
\textbf{Romanization}: Kore de ii no kana\\
\textbf{Expression (EN)}: I wonder if this is okay

\textbf{Variants (JA)}: これでよかったかな / こうしてよかったのかな / これで大丈夫かな\\
\textbf{Variants (EN)}: Was this the right choice? / Was it okay to do it this way? / Is this alright, I wonder?

\textbf{Intent (JA)}: 不安, 確認, つぶやき\\
\textbf{Intent (EN)}: uncertainty, seeking reassurance, murmur

\textbf{Tags (JA)}: 即時文法, 感情, 内語\\
\textbf{Tags (EN)}: immediate grammar, emotion, inner speech

\textbf{Adjusted Expression (JA)}:\\
これで(問題:もんだい)ないかどうか、(確認:かくにん)しています。\\
\textbf{Adjusted Expression (EN)}:\\
I'm checking to see if this is okay.

\textbf{Example (JA)}:\\
A「(クレジット:くれじっと)で(払:はら)っておいたけど、これでいいのかな」\\
\textbf{Example (EN)}:\\
A: I've paid with a credit card, but I wonder if this is okay.

\textbf{Notes (JA)}:\\
「これでいいのかな」は、自分の選択や判断に対する一時的な不安や確認をやわらかく表現することばです。「のかな」という言い回しによって、疑問をやさしく伝えたり、自分に問いかけたりする響きを持ちます。相手に聞いているようでもあり、自分に向けてのつぶやきでもあります。

\textbf{Notes (EN)}:\\
"Kore de ii no kana" gently expresses uncertainty or self-doubt about a decision. The phrase softens the question with "no kana," which adds a touch of hesitation or introspection. It can be directed to others for reassurance, or used as a form of self-talk to confirm one's own action or choice.

\newpage

\section*{480. Hmm, I can't decide.}

\textbf{Date}: 2025.08.02 (Sat)\\
\textbf{Index}: mayouna\\
\textbf{Expression (JA)}: うーん、(迷:まよ)うなぁ\\
\textbf{Romanization}: Uun, mayou naa\\
\textbf{Expression (EN)}: Hmm, I can't decide.

\textbf{Variants (JA)}: 迷うなあ / どうしようかな / どっちにしようかな\\
\textbf{Variants (EN)}: I’m torn... / What should I do? / Which one should I pick...

\textbf{Intent (JA)}: 迷い, 独り言, 感情の共有\\
\textbf{Intent (EN)}: indecision, self-talk, emotional sharing

\textbf{Tags (JA)}: 即時文法, 感情, 思考\\
\textbf{Tags (EN)}: immediate grammar, emotion, thinking

\textbf{Adjusted Expression (JA)}:\\
どちらにするか(決:き)めかねています。\\
\textbf{Adjusted Expression (EN)}:\\
I'm having trouble deciding between the two.

\textbf{Example (JA)}:\\
A「どれにする？」 B「うーん、(迷:まよ)うなぁ。どっちも(良:よ)さそうだし...」\\
\textbf{Example (EN)}:\\
A: Which one will you pick? B: Hmm, I can't decide. They both seem nice...

\textbf{Notes (JA)}:\\
「うーん、迷うなぁ」は、日常会話でよく聞かれる独り言風の表現で、決断がつかない気持ちを素直に表します。「うーん」で悩んでいる様子を、「なぁ」で感情をにじませます。相手に判断を委ねたり、共感を求めたりする場面でも使われます。

\textbf{Notes (EN)}:\\
"Uun, mayou naa" is a common spoken phrase used to express indecision. The drawn-out "Uun" signals deep thinking, and "naa" adds a soft emotional tone. It's often used aloud even when speaking to oneself, and it invites empathy or shared decision-making in conversation.

\newpage

\section*{479. Hmm... What should I do?}

\textbf{Date}: 2025.08.01 (Fri)\\
\textbf{Index}: doushiyou kana\\
\textbf{Expression (JA)}: どうしようかな\\
\textbf{Romanization}: Dou shiyou kana\\
\textbf{Expression (EN)}: Hmm... What should I do?

\textbf{Variants (JA)}: どうしよっかな / どうしたもんかな / どうしようか\\
\textbf{Variants (EN)}: What to do... / I wonder what I should do / Shall I...?

\textbf{Intent (JA)}: 迷い, 独り言, 思考の始まり\\
\textbf{Intent (EN)}: indecision, murmuring, thought initiation

\textbf{Tags (JA)}: 即時文法, 感情, 思考\\
\textbf{Tags (EN)}: immediate grammar, emotion, thinking

\textbf{Adjusted Expression (JA)}:\\
どうすればよいか(考:かんが)えています。\\
\textbf{Adjusted Expression (EN)}:\\
I'm considering what to do.

\textbf{Example (JA)}:\\
A「(映画:えいが)、(見:み)に(行:い)く？」 B「うーん、どうしようかな...」\\
\textbf{Example (EN)}:\\
A: Wanna go see a movie? B: Hmm... What should I do?

\textbf{Notes (JA)}:\\
「どうしようかな」は、思考や迷いの瞬間に自然に出てくる表現です。「かな」によって、独り言のようなやわらかさが生まれ、相手との共有よりも、自分の内面との対話を表します。小声でつぶやくように使われることも多く、非常に人間らしい発話です。このようななくてもよいような発話が実は人間関係にはとても重要な役割を果たします。

\textbf{Notes (EN)}:\\
"Dou shiyou kana" is often used when someone is thinking aloud or feeling indecisive. The "kana" softens the tone, making it feel like a private murmur rather than a request for advice. It reflects a moment of internal dialogue, a hallmark of natural, human speech. Such seemingly unnecessary utterances play a crucial role in human relationships, fostering connection and understanding.

\newpage

\section*{478. Not impossible / It's not that I wouldn't...}

\textbf{Date}: 2025.07.31 (Thu)\\
\textbf{Index}: nakuha nai\\
\textbf{Expression (JA)}: なくはない\\
\textbf{Romanization}: naku wa nai\\
\textbf{Expression (EN)}: not impossible / not that I wouldn’t

\textbf{Variants (JA)}: やってみなくはない / 行かなくはない / 食べなくはない\\
\textbf{Variants (EN)}: I might try it / I could go / I don't exactly dislike it

\textbf{Intent (JA)}: あいまい, 可能性, 含み\\
\textbf{Intent (EN)}: ambiguity, possibility, implication

\textbf{Tags (JA)}: 即時文法, 二重否定, 感情のぼかし\\
\textbf{Tags (EN)}: immediate grammar, double negative, emotional nuance

\textbf{Adjusted Expression (JA)}:\\
〜することも(考:かんが)えられる\\
\textbf{Adjusted Expression (EN)}:\\
It’s possible that I might...

\textbf{Example (JA)}:\\
A「いっしょに(行:い)く？」 B「うーん、(行:い)かなくはないけど...」\\
\textbf{Example (EN)}:\\
A: Wanna come with me? B: Hmm... I could go, I guess...

\textbf{Notes (JA)}:\\
「なくはない」は否定の中に肯定を含む二重否定表現で、あいまいな肯定や控えめな同意を示します。はっきりとは言えないが、完全には否定しない、というニュアンスがあり、感情や状況に揺れがあるときによく使われます。

\textbf{Notes (EN)}:\\
"Naku wa nai" is a double negative that implies a mild or hesitant affirmation. It doesn’t strongly agree, but it avoids complete denial. Often used when the speaker feels uncertain or is softening their response, conveying emotional hesitation or consideration.

\newpage

\section*{477. Quite / Pretty much / Fairly}

\textbf{Date}: 2025.07.30 (Wed)\\
\textbf{Index}: kanari\\
\textbf{Expression (JA)}: かなり\\
\textbf{Romanization}: Kanari\\
\textbf{Expression (EN)}: quite / pretty much / fairly

\textbf{Variants (JA)}: かなりいい / かなり(疲:つか)れた / かなり(難:むずか)しかった\\
\textbf{Variants (EN)}: pretty good / got quite tired / was fairly difficult

\textbf{Intent (JA)}: 程度, 感覚, 強調\\
\textbf{Intent (EN)}: degree, perception, emphasis

\textbf{Tags (JA)}: 即時文法, 副詞, 評価\\
\textbf{Tags (EN)}: immediate grammar, adverb, evaluation

\textbf{Adjusted Expression (JA)}:\\
相当〜\\
\textbf{Adjusted Expression (EN)}:\\
considerably / substantially

\textbf{Example (JA)}:\\
A「この(テスト:てすと)、どうだった？」 B「うん、かなり(難:むずか)しかったよ」\\
\textbf{Example (EN)}:\\
A: How was the test? B: Yeah, it was quite difficult.

\textbf{Notes (JA)}:\\
「かなり」は数量や程度を強調する副詞ですが、話し手がその場で感じた強さや印象を、あいまいさを保ったまま伝える特徴があります。「すごく」と違って客観的なトーンもあり、評価や感想を即時的に伝える際によく使われます。

\textbf{Notes (EN)}:\\
"Kanari" is an adverb used to emphasize degree or intensity. Unlike more subjective terms like "sugoku" (super), it has a slightly more objective tone and is often used to express one's impression or judgment on the spot, maintaining a flexible, nuanced sense of scale.

\newpage

\section*{476. They say / It seems...}

\textbf{Date}: 2025.07.29 (Tue)\\
\textbf{Index}: rashii\\
\textbf{Expression (JA)}: ...らしい\\
\textbf{Romanization}: ...rashii\\
\textbf{Expression (EN)}: They say / It seems...

\textbf{Variants (JA)}: (雨:あめ)らしい / テストは(明日:あした)らしい / もう(終:お)わったらしいよ\\
\textbf{Variants (EN)}: Seems like rain / Apparently the test is tomorrow / Looks like it's already over

\textbf{Intent (JA)}: 伝聞, 推測, 控えめ\\
\textbf{Intent (EN)}: hearsay, inference, tentative

\textbf{Tags (JA)}: 即時文法, 判断, 情報共有\\
\textbf{Tags (EN)}: immediate grammar, judgment, shared info

\textbf{Adjusted Expression (JA)}:\\
...とのことです\\
\textbf{Adjusted Expression (EN)}:\\
According to what I heard...

\textbf{Example (JA)}:\\
A「(駅:えき)、すごく(混:こ)んでたよ」 B「なんか、(事故:じこ)があったらしい」\\
\textbf{Example (EN)}:\\
A: The station was super crowded. B: Seems like there was an accident or something.

\textbf{Notes (JA)}:\\
「...らしい」は、自分が直接経験していないことを話すときに使います。文末に置かれることで、断定を避けたやわらかい伝え方になります。ニュース、噂、気配など、さまざまな情報源が背景にあるときに便利な表現です。「...らしいよ」「...らしかった」などでも使われます。

\textbf{Notes (EN)}:\\
"...rashii" expresses information the speaker didn't directly experience. It softens the tone, making the statement feel less assertive and more like a shared observation or rumor. It's often used with a sense of uncertainty or indirect knowledge, such as in "rashii yo" or "rashikatta".

\newpage

\section*{475. It looks like...}

\textbf{Date}: 2025.07.28 (Mon)\\
\textbf{Index}: mitai\\
\textbf{Expression (JA)}: ...みたい\\
\textbf{Romanization}: ...mitai\\
\textbf{Expression (EN)}: It looks like...

\textbf{Variants (JA)}: (雨:あめ)みたい / (大丈夫:だいじょうぶ)みたい / もう(終:お)わったみたい\\
\textbf{Variants (EN)}: Looks like rain / Seems okay / Looks like it's already over

\textbf{Intent (JA)}: 推測, 観察, 親しみ\\
\textbf{Intent (EN)}: guess, observation, affection

\textbf{Tags (JA)}: 即時文法, 判断, 自然な会話\\
\textbf{Tags (EN)}: immediate grammar, judgment, natural speech

\textbf{Adjusted Expression (JA)}:\\
...のようです\\
\textbf{Adjusted Expression (EN)}:\\
It appears that...

\textbf{Example (JA)}:\\
A「(人:ひと)、(集:あつ)まってるね」 B「なんか、(事故:じこ)みたい」\\
\textbf{Example (EN)}:\\
A: People are gathering there. B: Looks like there was an accident or something.

\textbf{Notes (JA)}:\\
「...みたい」は、確実ではないけれど、目にしたこと・感じたことから判断して話すときに使われます。話し手の主観や直感をそのまま言葉にしたような形で、即時文法の代表例といえます。「...みたいなんだけど」「...みたいね」など、語尾との組み合わせでもニュアンスが変わります。

\textbf{Notes (EN)}:\\
"...mitai" is used when expressing something based on appearance or intuition, without certainty. It reflects the speaker’s immediate impression, making it a hallmark of immediate grammar. Variants like "...mitai na n da kedo" or "...mitai ne" further soften the tone and add emotional nuance.

\newpage

\section*{474. Since a moment ago...}

\textbf{Date}: 2025.07.27 (Sun)\\
\textbf{Index}: sakkikara\\
\textbf{Expression (JA)}: さっきから\\
\textbf{Romanization}: Sakki kara\\
\textbf{Expression (EN)}: Since a moment ago

\textbf{Variants (JA)}: さっきからずっと / さっきから気になってた / さっきから見てたよ\\
\textbf{Variants (EN)}: This whole time / I’ve been wondering since earlier / I’ve been watching you since just now

\textbf{Intent (JA)}: 継続, 気づき, 感情の蓄積\\
\textbf{Intent (EN)}: continuity, awareness, emotional buildup

\textbf{Tags (JA)}: 即時文法, 感情, 自然な会話\\
\textbf{Tags (EN)}: immediate grammar, emotion, natural speech

\textbf{Adjusted Expression (JA)}:\\
先ほどからずっと...\\
\textbf{Adjusted Expression (EN)}:\\
I've been... since earlier

\textbf{Example (JA)}:\\
A「さっきから、(言:い)おうと(思:おも)ってたんだけど...」 B「うん、なに？」\\
\textbf{Example (EN)}:\\
A: I've been meaning to say this for a while now. B: Yeah? What is it?

\textbf{Notes (JA)}:\\
「さっきから」は、話し手が心の中に抱えていた感情や気づきをやわらかく切り出すための表現です。「ずっと」や「なんか」といった言葉とよく共起し、軽い焦り、気づき、または思いやりを含むことがあります。即時文法としては、過去から今に続く気持ちをそのまま言葉にする形式です。

\textbf{Notes (EN)}:\\
"Sakki kara" introduces a thought or feeling that has been lingering since a short while ago. It often co-occurs with words like "zutto" (all along) or "nanka" (somehow), and conveys mild persistence, emotional buildup, or gentle concern. As an instance of immediate grammar, it verbalizes an ongoing internal state in real time.

\newpage

\section*{473. About a moment ago...}

\textbf{Date}: 2025.07.26 (Sat)\\
\textbf{Index}: sakki\\
\textbf{Expression (JA)}: さっき\\
\textbf{Romanization}: Sakki\\
\textbf{Expression (EN)}: A moment ago

\textbf{Variants (JA)}: さっきね / さっきのことだけど / さっき言ってた\\
\textbf{Variants (EN)}: Just earlier / About what we said a moment ago / You mentioned earlier

\textbf{Intent (JA)}: つなぎ, 気づき, やわらかさ\\
\textbf{Intent (EN)}: connection, awareness, softness

\textbf{Tags (JA)}: 即時文法, 感情, 自然な会話\\
\textbf{Tags (EN)}: immediate grammar, emotion, natural speech

\textbf{Adjusted Expression (JA)}:\\
さきほど...\\
\textbf{Adjusted Expression (EN)}:\\
Earlier, ...

\textbf{Example (JA)}:\\
A「さっき、(言:い)いすぎたかも」 B「ううん、(気:き)にしてないよ」\\
\textbf{Example (EN)}:\\
A: I might've said too much earlier. B: No, don’t worry about it.

\textbf{Notes (JA)}:\\
「さっき」は直前の時間をやわらかく指す表現で、会話の流れを自然に保ちます。「さっきね」のように語尾に「ね」をつけることで、相手との共有感や距離の近さが強まります。あやまり、照れ、確認など、感情とともに使われることが多いです。

\textbf{Notes (EN)}:\\
"Sakki" refers gently to something that just happened. It helps the conversation flow naturally, especially when used with "ne" (as in "Sakki ne"). Often used when softening apologies, expressing affection, or continuing a thought from earlier.

\newpage

\section*{472. You remembered.}

\textbf{Date}: 2025.07.25 (Fri)\\
\textbf{Index}: oboetetekureta\\
\textbf{Expression (JA)}: (覚:おぼ)えててくれたんだ\\
\textbf{Romanization}: Oboetete kureta nda\\
\textbf{Expression (EN)}: You remembered for me.

\textbf{Variants (JA)}: 覚えててくれたね / 覚えててくれたの？ / 覚えててくれてたんだ\\
\textbf{Variants (EN)}: You remembered, huh? / Did you remember? / You had remembered for me.

\textbf{Intent (JA)}: 感謝, 共有, 驚き\\
\textbf{Intent (EN)}: gratitude, shared feeling, surprise

\textbf{Tags (JA)}: 即時文法, 感情, 親しみ\\
\textbf{Tags (EN)}: immediate grammar, emotion, closeness

\textbf{Adjusted Expression (JA)}:\\
(覚:おぼ)えていてくださって、ありがとうございます。\\
\textbf{Adjusted Expression (EN)}:\\
Thank you for remembering.

\textbf{Example (JA)}:\\
A「(誕生日:たんじょうび)、(覚:おぼ)えててくれたんだ」 B「もちろんだよ」\\
\textbf{Example (EN)}:\\
A: You remembered my birthday! B: Of course I did!

\textbf{Notes (JA)}:\\
「覚えててくれたんだ」は、相手の記憶や心づかいに気づいたときに思わず出る表現で、強い感謝の気持ちを自然に伝えます。「ありがとう」を言わなくても、その感情は十分に伝わる、即時文法的なやさしさがあります。

\textbf{Notes (EN)}:\\
"Oboetete kureta nda" is an expression that naturally conveys heartfelt gratitude when you realize someone has remembered something for you. Even without the word “thank you,” the feeling comes through clearly, making it a gentle and effective phrase of immediate grammar.

\newpage

\section*{471. You waited for me.}

\textbf{Date}: 2025.07.24 (Thu)\\
\textbf{Index}: matettekureta\\
\textbf{Expression (JA)}: (待:ま)っててくれたんだ\\
\textbf{Romanization}: Mattete kureta nda\\
\textbf{Expression (EN)}: You waited for me.

\textbf{Variants (JA)}: 待っててくれたね / 待っててくれたの？ / 待ってくれてたんだ\\
\textbf{Variants (EN)}: You stayed and waited, huh? / Were you waiting for me? / You’d been waiting for me.

\textbf{Intent (JA)}: 感謝, 共有, 驚き\\
\textbf{Intent (EN)}: gratitude, shared feeling, surprise

\textbf{Tags (JA)}: 即時文法, 感情, 親しみ\\
\textbf{Tags (EN)}: immediate grammar, emotion, closeness

\textbf{Adjusted Expression (JA)}:\\
お(待:ま)ちくださり、どうもありがとうございます。\\
\textbf{Adjusted Expression (EN)}:\\
Thank you for waiting.

\textbf{Example (JA)}:\\
A「(試験:しけん)、どうだった？」 B「ああ、(待:ま)っててくれたんだ」\\
\textbf{Example (EN)}:\\
A: How was the exam? B: Oh, you have been waiting for me.

\textbf{Notes (JA)}:\\
「待っててくれたんだ」は、相手のやさしさに気づいて、感謝やうれしさを込めて返すことばです。「んだ」には気づきや驚き、感情の高まりが含まれています。「待っててくれてたんだ」とも言えますが、「て」の音が多く言いにくいかもしれませんね。「待っててくれたんだ」も正しい言い方です。

\textbf{Notes (EN)}:\\
"Mattete kureta nda" is a phrase that expresses gratitude and warmth when you realize someone has been waiting for you. The "nda" ending adds a sense of realization or surprise, conveying emotional depth. You can also say "Mattete kurete ta nda," but it might be a bit harder to pronounce due to the extra syllables. Both forms are correct.

\newpage

\section*{470. Thanks for doing that.}

\textbf{Date}: 2025.07.23 (Wed)\\
\textbf{Index}: tekuretane\\
\textbf{Expression (JA)}: ...てくれたね\\
\textbf{Romanization}: ...te kureta ne\\
\textbf{Expression (EN)}: You did that for me, didn't you?

\textbf{Variants (JA)}: 手伝ってくれたね / 待っててくれたね / 覚えててくれたね\\
\textbf{Variants (EN)}: You helped me, right? / You waited for me, didn’t you? / You remembered, didn’t you?

\textbf{Intent (JA)}: 感謝, 共有, 親密さ\\
\textbf{Intent (EN)}: gratitude, shared moment, closeness

\textbf{Tags (JA)}: 即時文法, 感情, 親しみ\\
\textbf{Tags (EN)}: immediate grammar, emotion, affection

\textbf{Adjusted Expression (JA)}:\\
〜してくれて、ありがとう。\\
\textbf{Adjusted Expression (EN)}:\\
Thank you for doing that.

\textbf{Example (JA)}:\\
A「いつも、(一緒:いっしょ)にいてくれたね」 B「ぼくが(一緒:いっしょ)にいたかっただけだよ」\\
\textbf{Example (EN)}:\\
A: Thanks so much for always being with me. B: I just wanted to be with you.

\textbf{Notes (JA)}:\\
「...てくれたね」は、相手の行動に対する感謝やあたたかい気持ちを表す表現です。文末の「ね」によって、相手との共有感や親しさが自然に生まれます。感謝を直接言うよりも、思い出をふりかえるような、やさしいトーンになります。

\textbf{Notes (EN)}:\\
"...te kureta ne" conveys warm gratitude for something someone did for you. The sentence-final particle "ne" adds a sense of shared memory or closeness. It's softer than directly saying thank you, it reflects appreciation through emotional connection.

\newpage

\section*{469. Even if I have to...}

\textbf{Date}: 2025.07.22 (Tue)\\
\textbf{Index}: tedemo\\
\textbf{Expression (JA)}: ...てでも\\
\textbf{Romanization}: ...te demo\\
\textbf{Expression (EN)}: Even if I have to...

\textbf{Variants (JA)}: 行ってでも / 奪ってでも / 泣いてでも\\
\textbf{Variants (EN)}: Even if I go / Even if I have to take it / Even if I cry

\textbf{Intent (JA)}: 強調, 決意, 目的\\
\textbf{Intent (EN)}: emphasis, determination, goal-oriented

\textbf{Tags (JA)}: 即時文法, 意志, 過程強調\\
\textbf{Tags (EN)}: immediate grammar, willpower, process emphasis

\textbf{Adjusted Expression (JA)}:\\
たとえ〜してでも、やりとげたいです。\\
\textbf{Adjusted Expression (EN)}:\\
I want to do it, even if I have to go through 〜.

\textbf{Example (JA)}:\\
A「その(本:ほん)、もう(手:て)に(入:はい)らないよ」 B「いかなる(大金:たいきん)を支払ってでも(手:て)に(入:い)れたい！」\\
\textbf{Example (EN)}:\\
A: That book is no longer available. B: I want to get it, even if I have to pay any amount of money!

\textbf{Notes (JA)}:\\
「...てでも」は、困難や犠牲を伴ってでも目的を達成しようとする強い意志を表す表現です。文法的には動詞のて形に「でも」をつけて使われます。やや演劇的・ドラマ的な響きもありますが、日常会話でも強調したいときに使われます。

\textbf{Notes (EN)}:\\
“...te demo” expresses strong determination to achieve something, even at a cost. It attaches "demo" to the te-form of a verb. It has a dramatic tone but is also used casually to express persistence or resolve.

\newpage

\section*{468. Not *that* far...}

\textbf{Date}: 2025.07.21 (Mon)\\
\textbf{Index}: madewane\\
\textbf{Expression (JA)}: ...まではね\\
\textbf{Romanization}: ...made wa ne\\
\textbf{Expression (EN)}: Not that far...

\textbf{Variants (JA)}: そこまではね / やるけど...そこまではね / 行くけど...泊まりまではね\\
\textbf{Variants (EN)}: I can do it, but not that far / I'm okay up to a point / That's too much for me

\textbf{Intent (JA)}: ためらい, 限定, 控えめ\\
\textbf{Intent (EN)}: hesitation, limitation, modesty

\textbf{Tags (JA)}: 即時文法, ためらい, 条件づけ\\
\textbf{Tags (EN)}: immediate grammar, hesitation, conditional

\textbf{Adjusted Expression (JA)}:\\
それ以上はちょっと難しいです。\\
\textbf{Adjusted Expression (EN)}:\\
Going beyond that might be difficult.

\textbf{Example (JA)}:\\
A「じゃ、いっしょに(泊:と)まって(行:い)こうよ」 B「うーん、泊まりまではね...」\\
\textbf{Example (EN)}:\\
A: Then let's stay over together. B: Hmm... not that far.

\textbf{Notes (JA)}:\\
「...まではね」は、ある行動を完全には否定しないが、そこまでには踏み込めないという気持ちを表す即時表現です。相手に配慮しつつ、自分の限界をさりげなく示すために使われます。

\textbf{Notes (EN)}:\\
“...made wa ne” implies hesitation or reluctance to go beyond a certain point. It's a way to set a soft boundary without outright refusal, often used in casual, emotionally aware speech.

\newpage

\section*{467. What is that!?}

\textbf{Date}: 2025.07.20 (Sun)\\
\textbf{Index}: nani-sore\\
\textbf{Expression (JA)}: 何、それ！\\
\textbf{Romanization}: Nani, sore!\\
\textbf{Expression (EN)}: What is that!?

\textbf{Variants (JA)}: なにそれ / なに、それ / 何それ！ / 何、それ！\\
\textbf{Variants (EN)}: What the heck is that? / Seriously? / What is that!? / No way!

\textbf{Intent (JA)}: 驚き, 共感, 呆れ\\
\textbf{Intent (EN)}: surprise, empathy, exasperation

\textbf{Tags (JA)}: 即時文法, 感情, 感嘆\\
\textbf{Tags (EN)}: immediate grammar, emotion, exclamation

\textbf{Adjusted Expression (JA)}:\\
それは何ですか？\\
\textbf{Adjusted Expression (EN)}:\\
What is that?

\textbf{Example (JA)}:\\
A「あんなに親切にしたのに何も言わないんだよ」 B「(何:なに)、それ！」\\
\textbf{Example (EN)}:\\
A: I was so nice to him, but he doesn't say anything. B: What is that!?

\textbf{Notes (JA)}:\\
「何、それ！」は即時的な驚きや共感、呆れを表す感嘆的な応答です。「それは何ですか？」のような文法的な質問とは異なり、感情が先行し、文末表現も省略されます。イントネーションや間の取り方によって意味の幅も広がります。

\textbf{Notes (EN)}:\\
"Nani, sore!" is an exclamatory, immediate response that expresses surprise, empathy, or mild disbelief. Unlike the more neutral "What is that?", this expression focuses on emotional reaction. It omits formal grammar and relies heavily on tone and timing.

\newpage

\section*{466. Might be nice}

\textbf{Date}: 2025.07.19 (Sat)\\
\textbf{Index}: iikamo\\
\textbf{Expression (JA)}: いいかも\\
\textbf{Romanization}: Ii kamo\\
\textbf{Expression (EN)}: Might be nice

\textbf{Variants (JA)}: それ、いいかも / いいかもしれない / けっこういいかも\\
\textbf{Variants (EN)}: That might be good / Could be nice / Sounds like a good idea

\textbf{Intent (JA)}: 判断, 柔らかい肯定, 共感\\
\textbf{Intent (EN)}: judgment, soft agreement, shared feeling

\textbf{Tags (JA)}: 即時文法, 評価, 控えめ\\
\textbf{Tags (EN)}: immediate grammar, assessment, tentative

\textbf{Adjusted Expression (JA)}:\\
それは(良:よ)いかもしれませんね。\\
\textbf{Adjusted Expression (EN)}:\\
That might be a good idea.

\textbf{Example (JA)}:\\
A「ちょっと(遠出:とおで)してみる？」 B「それ、いいかも」\\
\textbf{Example (EN)}:\\
A: Wanna go for a little trip? B: That might be nice.

\textbf{Notes (JA)}:\\
「いいかも」は、「かもしれない」の略式で、提案や考えに対して肯定的な可能性を示します。即時に判断を返しながらも、柔らかく控えめな印象を与えるため、日常会話でよく使われます。

\textbf{Notes (EN)}:\\
"Ii kamo" is a casual form of "ii kamo shirenai" (might be good), used to express a soft, tentative agreement. It offers immediate feedback while remaining non-assertive, making it common in everyday speech.

\newpage

\section*{465. Can I...?}

\textbf{Date}: 2025.07.18 (Fri)\\
\textbf{Index}: temoii\\
\textbf{Expression (JA)}: 〜てもいい？\\
\textbf{Romanization}: …temo ii?\\
\textbf{Expression (EN)}: Can I...?

\textbf{Variants (JA)}: 見てもいい？ / 食べてもいい？ / 入ってもいい？\\
\textbf{Variants (EN)}: Can I look? / Can I eat it? / Can I come in?

\textbf{Intent (JA)}: 許可, 確認, 丁寧な依頼\\
\textbf{Intent (EN)}: permission, checking, polite request

\textbf{Tags (JA)}: 即時文法, 依頼, 丁寧, 疑問\\
\textbf{Tags (EN)}: immediate grammar, request, politeness, question

\textbf{Adjusted Expression (JA)}:\\
〜てもよろしいですか？\\
\textbf{Adjusted Expression (EN)}:\\
Would it be all right if I...?

\textbf{Example (JA)}:\\
A「それ、(見:み)てもいい？」 B「うん、いいよ。」\\
\textbf{Example (EN)}:\\
A: Can I look at that? B: Sure.

\textbf{Notes (JA)}:\\
「〜てもいい？」は、相手に何かをしてよいかを尋ねるカジュアルな表現です。親しい間柄ではよく使われますが、丁寧に言う場合は「〜てもよろしいですか？」に言い換えます。子どもが親に何かを頼むときにもよく使われる、即時的で自然な聞き方です。

\textbf{Notes (EN)}:\\
The phrase "...temo ii?" is a casual way to ask for permission, often used among friends or family. It literally means "Is it okay if I...?" In more formal contexts, it’s replaced with "...temo yoroshii desu ka?" This expression is commonly used in everyday conversation, especially by children when asking adults for permission.

\newpage

\section*{464. No matter what}

\textbf{Date}: 2025.07.17 (Thu)\\
\textbf{Index}: doushitemo\\
\textbf{Expression (JA)}: どうしても\\
\textbf{Romanization}: Doushitemo.\\
\textbf{Expression (EN)}: No matter what.

\textbf{Variants (JA)}: どうしても無理 / どうしても必要 / どうしてもやりたい\\
\textbf{Variants (EN)}: I just can’t. / I really have to. / I absolutely must.

\textbf{Intent (JA)}: 強調, 決意, 不可避\\
\textbf{Intent (EN)}: emphasis, determination, inevitability

\textbf{Tags (JA)}: 即時文法, 強調, 断固\\
\textbf{Tags (EN)}: immediate grammar, emphasis, firmness

\textbf{Adjusted Expression (JA)}:\\
どうしても、(必要:ひつよう)です。\\
\textbf{Adjusted Expression (EN)}:\\
I really need to.

\textbf{Example (JA)}:\\
A「(来:こ)れないの？」 B「うん、どうしても」\\
\textbf{Example (EN)}:\\
A: Can't you come? B: No, I just can't.

\textbf{Notes (JA)}:\\
「どうしても。」は強い意思や避けられない事情を表す短い応答です。肯定でも否定でも使われ、文脈で意味が変わります。単独で使うと「譲れない」「無理だ」という強い気持ちが含まれます。

\textbf{Notes (EN)}:\\
"Doushitemo." shows strong will or inevitability. It can be used in affirmative or negative contexts. Alone, it often conveys that something is non-negotiable or impossible.

\newpage

\section*{463. That again.}

\textbf{Date}: 2025.07.16 (Wed)\\
\textbf{Index}: sorebakkari\\
\textbf{Expression (JA)}: それ、ばっかり。\\
\textbf{Romanization}: Sore, bakkari.\\
\textbf{Expression (EN)}: That again.

\textbf{Variants (JA)}: そればっかりじゃん / またそれ / そればかり\\
\textbf{Variants (EN)}: Always that. / Just that. / You’re always doing that.

\textbf{Intent (JA)}: 批判, 呆れ, 指摘\\
\textbf{Intent (EN)}: criticism, exasperation, pointing out

\textbf{Tags (JA)}: 即時文法, 指摘, カジュアル\\
\textbf{Tags (EN)}: immediate grammar, remark, casual

\textbf{Adjusted Expression (JA)}:\\
そればかりですね。\\
\textbf{Adjusted Expression (EN)}:\\
You only do that.

\textbf{Example (JA)}:\\
A「ねぇ、ラーメン、(食:た)ベ(行:い)こ？」 B「それ、ばっかり」\\
\textbf{Example (EN)}:\\
A: Hey, let's go eat ramen? B: That again.

\textbf{Notes (JA)}:\\
「それ、ばっかり。」は相手の行動や話題が偏っていることを軽く批判・指摘する表現です。やや呆れを込めたニュアンスがありますが、口調によって強さは変わります。

\textbf{Notes (EN)}:\\
"Sore, bakkari." is used to point out that someone is always doing or talking about the same thing. It conveys a light criticism or exasperation, but tone can soften or strengthen it.

\newpage

\section*{462. So what?}

\textbf{Date}: 2025.07.15 (Tue)\\
\textbf{Index}: dattara\\
\textbf{Expression (JA)}: だったら？\\
\textbf{Romanization}: Dattara?\\
\textbf{Expression (EN)}: So what?

\textbf{Variants (JA)}: それで？ / だから？ / もしそうなら？\\
\textbf{Variants (EN)}: And? / Then what? / If so, then what?

\textbf{Intent (JA)}: 反論, 問い返し, 挑発\\
\textbf{Intent (EN)}: rebuttal, questioning, provocation

\textbf{Tags (JA)}: 即時文法, 対抗, 短答\\
\textbf{Tags (EN)}: immediate grammar, challenge, short reply

\textbf{Adjusted Expression (JA)}:\\
もしそうなら、どうするつもりですか？\\
\textbf{Adjusted Expression (EN)}:\\
If that’s the case, what are you going to do?

\textbf{Example (JA)}:\\
A「それ、(絶対:ぜったい)、(間違:まちが)ってるよ」 B「だったら？」\\
\textbf{Example (EN)}:\\
A: That's definitely wrong. B: So what?

\textbf{Notes (JA)}:\\
「だったら？」は相手の主張を受けたうえで、次の説明や行動を問い返す表現です。やや挑戦的・突き放す印象があります。でも、短かいので、相手の発言に対してすぐに反応することができます。

\textbf{Notes (EN)}:\\
"Dattara?" is used to question or challenge the other person’s statement, implying "If that’s the case, so what?" It can sound provocative or dismissive. It's a short response that allows for immediate reaction to the other person's claim.

\newpage

\section*{461. Right.}

\textbf{Date}: 2025.07.14 (Mon)\\
\textbf{Index}: desune\\
\textbf{Expression (JA)}: ですね\\
\textbf{Romanization}: Desu ne.\\
\textbf{Expression (EN)}: Right.

\textbf{Variants (JA)}: そうですね / ですねー / ええ、ですね\\
\textbf{Variants (EN)}: Yeah. / It is. / Sure is.

\textbf{Intent (JA)}: 同意, 確認, 共感\\
\textbf{Intent (EN)}: agreement, acknowledgment, empathy

\textbf{Tags (JA)}: 即時文法, あいづち, 共感\\
\textbf{Tags (EN)}: immediate grammar, backchannel, empathy

\textbf{Adjusted Expression (JA)}:\\
そうですね。\\
\textbf{Adjusted Expression (EN)}:\\
That’s true.

\textbf{Example (JA)}:\\
A「(暑:あつ)いですね」 B「ですね」\\
\textbf{Example (EN)}:\\
A: It's hot. B: Right.

\textbf{Notes (JA)}:\\
「ですね。」は相手の発言に軽く同意や共感を示す短い応答です。自分の意見を付け加えず、会話を穏やかにつなぎます。この発話は相手の発話に続けてすぐに言わなければなりません。

\textbf{Notes (EN)}:\\
"Desu ne." is a short, soft way to show agreement or empathy without adding further opinion. It keeps the conversation smooth and polite. This phrase should be used immediately after the other person's statement.

\newpage

\section*{460. Piece of cake}

\textbf{Date}: 2025.07.13 (Sun)\\
\textbf{Index}: asameshimae\\
\textbf{Expression (JA)}: (朝飯前:あさめしまえ)\\
\textbf{Romanization}: Asameshimae\\
\textbf{Expression (EN)}: Piece of cake

\textbf{Variants (JA)}: 朝飯前だよ / そんなの朝飯前 / 朝飯前です\\
\textbf{Variants (EN)}: It’s easy / No sweat / It’s nothing

\textbf{Intent (JA)}: 容易さ, 余裕, 即答\\
\textbf{Intent (EN)}: ease, confidence, immediate response

\textbf{Tags (JA)}: 即時文法, 評価, カジュアル\\
\textbf{Tags (EN)}: immediate grammar, assessment, casual

\textbf{Adjusted Expression (JA)}:\\
(簡単:かんたん)にできます。\\
\textbf{Adjusted Expression (EN)}:\\
I can do it easily.

\textbf{Example (JA)}:\\
A「この(問題:もんだい)できる？」 B「(朝飯前:あさめしまえ)だよ。」\\
\textbf{Example (EN)}:\\
A: Can you solve this problem? B: Piece of cake.

\textbf{Notes (JA)}:\\
「朝飯前」は「とても簡単」という意味のカジュアルな表現です。即答で余裕を伝えるときに使います。「飯」は「ごはん」の意味ですが、「めし、おねがいします」というと乱暴な言い方になります。しかし、「飯」が乱暴な表現だからと言って「朝飯前」の表現を「朝ごはん前」と言ってはいけません。「飯」はカジュアルで男性的な響きを持つため、丁寧な場ではあまり使いませんが、「朝飯前」は慣用句なので失礼にはなりません。

\textbf{Notes (EN)}:\\
"Asameshimae" literally means "before breakfast" and is a casual way to say something is very easy, like "a piece of cake." It conveys confidence and ease in response. The word "meshi" (飯) is casual and can sound rough if used in other contexts, but in this phrase, it's perfectly acceptable. Avoid saying "asagohan mae" as it doesn't carry the same meaning. "meshi" has a casual, masculine tone, so it's not typically used in formal situations, but as a set phrase, "asameshimae" is not considered rude.

\newpage

\section*{459. Yeah, I've had that experience}

\textbf{Date}: 2025.07.12 (Sat)\\
\textbf{Index}: aruaru\\
\textbf{Expression (JA)}: あるある\\
\textbf{Romanization}: Aru aru\\
\textbf{Expression (EN)}: Yeah, I've had that experience

\textbf{Variants (JA)}: あるよ / うん、あるある\\
\textbf{Variants (EN)}: Yeah, I have / I’ve had that

\textbf{Intent (JA)}: 経験確認, 共感, 即答\\
\textbf{Intent (EN)}: confirmation of experience, empathy, immediate response

\textbf{Tags (JA)}: 即時文法, 共感, 経験\\
\textbf{Tags (EN)}: immediate grammar, empathy, experience

\textbf{Adjusted Expression (JA)}:\\
そういう(経験:けいけん)はあります。\\
\textbf{Adjusted Expression (EN)}:\\
I’ve had that kind of experience.

\textbf{Example (JA)}:\\
A「そういう(経験:けいけん)、ある？」 B「あるある。」\\
\textbf{Example (EN)}:\\
A: Have you had that kind of experience? B: Yeah, I have.

\textbf{Notes (JA)}:\\
「あるある」は経験を明示的に尋ねられたときに共感を込めて即答する言い方です。「そんなことある？」のように疑念を含む質問には使いません。

\textbf{Notes (EN)}:\\
"Aru aru" here is used as a quick, empathetic response when someone asks if you’ve had a similar experience. It’s not typically used when the question expresses disbelief, like "Does that really happen?"

\newpage

\section*{458. In that case}

\textbf{Date}: 2025.07.11 (Fri)\\
\textbf{Index}: tonaruto\\
\textbf{Expression (JA)}: となると\\
\textbf{Romanization}: To naru to\\
\textbf{Expression (EN)}: In that case

\textbf{Variants (JA)}: そうなると / じゃあ、となると / それとなると\\
\textbf{Variants (EN)}: If that’s the case / Then / That means

\textbf{Intent (JA)}: 推論, 整理, 条件\\
\textbf{Intent (EN)}: inference, reasoning, conditional

\textbf{Tags (JA)}: 即時文法, 推論, 状況整理\\
\textbf{Tags (EN)}: immediate grammar, inference, situation

\textbf{Adjusted Expression (JA)}:\\
そういうことになりますと、\\
\textbf{Adjusted Expression (EN)}:\\
If that is the situation,

\textbf{Example (JA)}:\\
A「(田中:たなか)さん、(明日:あす)(来:こ)ないって。」 B「となると、(会議:かいぎ)は(延期:えんき)かな。」\\
\textbf{Example (EN)}:\\
A: Tanaka-san says he won’t come tomorrow. B: In that case, the meeting might be postponed.

\textbf{Notes (JA)}:\\
「となると」は前の情報を受けて、次の推測や判断を示す表現です。話の切り替えや整理に使われます。

\textbf{Notes (EN)}:\\
"To naru to" is used to indicate reasoning or inference based on what was just said. It connects the previous statement to a new conclusion or consideration.

\newpage

\section*{457. Apparently}

\textbf{Date}: 2025.07.10 (Thu)\\
\textbf{Index}: rashii\\
\textbf{Expression (JA)}: らしい\\
\textbf{Romanization}: Rashii\\
\textbf{Expression (EN)}: Apparently

\textbf{Variants (JA)}: 〜らしいよ / 〜らしいね / 〜らしいです\\
\textbf{Variants (EN)}: It seems / I heard / Apparently

\textbf{Intent (JA)}: 伝聞, 推測, 共有\\
\textbf{Intent (EN)}: hearsay, supposition, sharing

\textbf{Tags (JA)}: 即時文法, 伝聞, 推測\\
\textbf{Tags (EN)}: immediate grammar, hearsay, supposition

\textbf{Adjusted Expression (JA)}:\\
そうだそうです。\\
\textbf{Adjusted Expression (EN)}:\\
That's what I heard.

\textbf{Example (JA)}:\\
A「(明日:あした)、(雨:あめ)、(降:ふ)るかな？」 B「ええ、(降:ふ)るらしい。」\\
\textbf{Example (EN)}:\\
A: Is it going to rain tomorrow? B: Yeah, I heard it's going to rain.

\textbf{Notes (JA)}:\\
「らしい」は他から得た情報や伝聞を表す言葉です。話し手が完全に確信していないことを柔らかく伝えます。

\textbf{Notes (EN)}:\\
"Rashii" expresses that you heard or suppose something without being entirely sure. It’s used to soften statements.

\newpage

\section*{456. Looks like it’s over}

\textbf{Date}: 2025.07.09 (Wed)\\
\textbf{Index}: owattappoi\\
\textbf{Expression (JA)}: (終:お)わったっぽい\\
\textbf{Romanization}: Owatta ppoi\\
\textbf{Expression (EN)}: Looks like it’s over

\textbf{Variants (JA)}: 終わったみたい / もう終わったっぽい / 終わった感じ\\
\textbf{Variants (EN)}: Seems like it’s over / I think it’s over / Looks like it’s finished

\textbf{Intent (JA)}: 推量, 確認, 軽い言い方\\
\textbf{Intent (EN)}: supposition, confirmation, casual

\textbf{Tags (JA)}: 即時文法, 推量, カジュアル\\
\textbf{Tags (EN)}: immediate grammar, supposition, casual

\textbf{Adjusted Expression (JA)}:\\
(終:お)わったようです。\\
\textbf{Adjusted Expression (EN)}:\\
It seems to have ended.

\textbf{Example (JA)}:\\
A「まだやってる？」 B「(終:お)わったっぽいよ。」\\
\textbf{Example (EN)}:\\
A: Is it still going? B: Looks like it’s over.

\textbf{Notes (JA)}:\\
「終わったっぽい」は「〜っぽい」で推測や不確かな感じを表すカジュアルな言い方です。軽い確認や共有のときによく使います。

\textbf{Notes (EN)}:\\
"Owatta ppoi" uses "ppoi" to express a guess or impression in a casual tone. It’s often used when you’re not completely sure but want to share what you think.

\newpage

\section*{455. Oh, I see}

\textbf{Date}: 2025.07.08 (Tue)\\
\textbf{Index}: sokka\\
\textbf{Expression (JA)}: そっか\\
\textbf{Romanization}: Sokka\\
\textbf{Expression (EN)}: Oh, I see

\textbf{Variants (JA)}: そっかそっか / そっかー / ああ、そっか\\
\textbf{Variants (EN)}: I get it / I see / Ah, okay

\textbf{Intent (JA)}: 納得, 気づき, 共感\\
\textbf{Intent (EN)}: understanding, realization, empathy

\textbf{Tags (JA)}: 即時文法, 気づき, 相づち\\
\textbf{Tags (EN)}: immediate grammar, realization, backchannel

\textbf{Adjusted Expression (JA)}:\\
そうですか。\\
\textbf{Adjusted Expression (EN)}:\\
I see.

\textbf{Example (JA)}:\\
A「もう(帰:かえ)ったよ」 B「そっか」\\
\textbf{Example (EN)}:\\
A: They already left. B: Oh, I see.

\textbf{Notes (JA)}:\\
「そっか」は理解や納得を軽く表すカジュアルな応答です。相手の話に共感しながら反応する時に使います。

\textbf{Notes (EN)}:\\
"Sokka" is a casual way to show you understand or realize something. It conveys light agreement and acceptance.

\newpage

\section*{454. Didn't they come?}

\textbf{Date}: 2025.07.07 (Mon)\\
\textbf{Index}: kitanjanai\\
\textbf{Expression (JA)}: (来:き)たんじゃない？\\
\textbf{Romanization}: Kita n ja nai?\\
\textbf{Expression (EN)}: Didn't they come?

\textbf{Variants (JA)}: 来たんじゃないかな？ / 来たと思うよ / 来たんじゃん？\\
\textbf{Variants (EN)}: I think they came. / Did they come? / Maybe they came.

\textbf{Intent (JA)}: 推量, 確認, 共有\\
\textbf{Intent (EN)}: supposition, confirmation, shared knowledge

\textbf{Tags (JA)}: 即時文法, 推量, 問いかけ\\
\textbf{Tags (EN)}: immediate grammar, supposition, question

\textbf{Adjusted Expression (JA)}:\\
(来:き)たのではないでしょうか。\\
\textbf{Adjusted Expression (EN)}:\\
Perhaps they came.

\textbf{Example (JA)}:\\
A「もうみんな(集:あつ)まった？」 B「(来:き)たんじゃない？」\\
\textbf{Example (EN)}:\\
A: Is everyone here already? B: Didn't they come?

\textbf{Notes (JA)}:\\
「来たんじゃない？」は推量と確認を軽い調子で示す表現です。相手に同意や反応を促す柔らかい問いかけです。

\textbf{Notes (EN)}:\\
"Kita n janai?" is a casual way to suggest or confirm that something has happened, inviting agreement or response from the listener.

\newpage

\section*{453. I don't know}

\textbf{Date}: 2025.07.06 (Sun)\\
\textbf{Index}: wakannai\\
\textbf{Expression (JA)}: わかんない\\
\textbf{Romanization}: Wakannai\\
\textbf{Expression (EN)}: I don't know

\textbf{Variants (JA)}: わからない / うーん、わかんない / さあ、わかんない\\
\textbf{Variants (EN)}: I’m not sure / I have no idea / I don’t really know

\textbf{Intent (JA)}: 否定, ためらい, 即時応答\\
\textbf{Intent (EN)}: negation, hesitation, immediate response

\textbf{Tags (JA)}: 即時文法, 否定, 軽い\\
\textbf{Tags (EN)}: immediate grammar, negation, casual

\textbf{Adjusted Expression (JA)}:\\
よく(分:わ)かりません。\\
\textbf{Adjusted Expression (EN)}:\\
I'm not sure.

\textbf{Example (JA)}:\\
A「これどうやって(使:つか)うの？)」 B「わかんない。」\\
\textbf{Example (EN)}:\\
A: How do you use this? B: I don't know.

\textbf{Notes (JA)}:\\
「わかんない」は「分からない」の口語的・縮約形で、カジュアルで即座に答える否定表現です。場面によってはややそっけない印象を与えることもあります。

\textbf{Notes (EN)}:\\
"Wakannai" is the casual contracted form of "wakaranai" (I don't know). It's used in quick, informal replies and can sound a bit blunt depending on tone.

\newpage

\section*{452. Sorry, sorry}

\textbf{Date}: 2025.07.05 (Sat)\\
\textbf{Index}: gomengomen\\
\textbf{Expression (JA)}: ごめんごめん\\
\textbf{Romanization}: Gomen gomen\\
\textbf{Expression (EN)}: Sorry, sorry

\textbf{Variants (JA)}: ごめんね / あ、ごめんごめん / ごめんごめんね\\
\textbf{Variants (EN)}: Oh, sorry / Sorry about that / Oops, sorry

\textbf{Intent (JA)}: 謝罪, 軽い気まずさ, 親しみ\\
\textbf{Intent (EN)}: apology, slight awkwardness, friendly

\textbf{Tags (JA)}: 即時文法, 謝罪, 軽い\\
\textbf{Tags (EN)}: immediate grammar, apology, light

\textbf{Adjusted Expression (JA)}:\\
(申:もう)し(訳:わけ)ない、(気:き)づかなかったよ。\\
\textbf{Adjusted Expression (EN)}:\\
I’m sorry, I didn’t notice.

\textbf{Example (JA)}:\\
A「(あれ？(ドア:どあ)(閉:し)めた？)」 B「ごめんごめん、(閉:し)めてくるね。」\\
\textbf{Example (EN)}:\\
A: Did you close the door? B: Sorry, sorry, I’ll go close it.

\textbf{Notes (JA)}:\\
「ごめんごめん」はカジュアルで軽い謝罪を表す言い方です。繰り返すことで「気にしないでね」という柔らかい気持ちを添えます。

\textbf{Notes (EN)}:\\
"Gomen gomen" is a casual, friendly way to apologize. Repeating it softens the tone and shows you don't want the other person to feel bad.

\newpage

\section*{451. I wish that were true...}

\textbf{Date}: 2024.07.04 (Fri)\\
\textbf{Index}: dattaraiikedo\\
\textbf{Expression (JA)}: だったらいいけど...\\
\textbf{Romanization}: Dattara ii kedo...\\
\textbf{Expression (EN)}: I wish that were true...

\textbf{Variants (JA)}: そうだったらいいけど / そうならいいけど / そうだといいけど\\
\textbf{Variants (EN)}: I hope so... / That would be nice... / It'd be great if that were true...

\textbf{Intent (JA)}: 期待, 不安, ためらい\\
\textbf{Intent (EN)}: hope, uncertainty, hesitation

\textbf{Tags (JA)}: 即時文法, 感情, 控えめ\\
\textbf{Tags (EN)}: immediate grammar, emotion, tentative

\textbf{Adjusted Expression (JA)}:\\
そうであれば(嬉:うれ)しいのですが、どうでしょうか。\\
\textbf{Adjusted Expression (EN)}:\\
It would make me happy if that were the case, but I’m not sure.

\textbf{Example (JA)}:\\
A「(明日:あした)は(晴:は)れるかな。」 B「だったらいいけど…」\\
\textbf{Example (EN)}:\\
A: Do you think it’ll be sunny tomorrow? B: I wish that were true...

\textbf{Notes (JA)}:\\
「だったらいいけど…」は、希望を述べながらも不安や現実への慎重さを含む表現です。言い切らないことで控えめな気持ちを伝えます。

\textbf{Notes (EN)}:\\
"Dattara ii kedo..." expresses a wish mixed with uncertainty. The lack of certainty softens the statement.

\newpage

\section*{450. That's not true / No way}

\textbf{Date}: 2024.07.03 (Thu)\\
\textbf{Index}: sonnakotonaitte\\
\textbf{Expression (JA)}: そんなことないって\\
\textbf{Romanization}: Sonna koto nai tte\\
\textbf{Expression (EN)}: That's not true

\textbf{Variants (JA)}: そんなことないよ / そんなことないから / そんなことないってば\\
\textbf{Variants (EN)}: That's not true / No way / It's not like that

\textbf{Intent (JA)}: 否定, 励まし, 共感\\
\textbf{Intent (EN)}: negation, encouragement, empathy

\textbf{Tags (JA)}: 即時文法, 感情, 反応\\
\textbf{Tags (EN)}: immediate grammar, emotion, reaction

\textbf{Adjusted Expression (JA)}:\\
そんなことはないと(思:おも)いますよ。\\
\textbf{Adjusted Expression (EN)}:\\
I don't think that's the case.

\textbf{Example (JA)}:\\
A「(私:わたし)、(何:なに)やっても(ダメ:だめ)なんだよね。」 B「そんなことないって！」\\
\textbf{Example (EN)}:\\
A: I feel like I fail at everything. B: That's not true!

\textbf{Notes (JA)}:\\
「そんなことないって」は、相手の否定的な発言をやさしく否定し、励ます表現です。親しい会話でよく使います。

\textbf{Notes (EN)}:\\
"Sonna koto nai tte" gently denies someone's negative statement and offers encouragement. It's common in close conversations.

\newpage

\section*{449. Examplifying things}

\textbf{Date}: 2024.07.02 (Wed)\\
\textbf{Index}: toka\\
\textbf{Expression (JA)}: (本:ほん)とか(雑誌:ざっし)とか、...\\
\textbf{Romanization}: Hon toka zasshi toka,...\\
\textbf{Expression (EN)}: Books or magazine or kinda stuff, ...

\textbf{Variants (JA)}: ...とか... / ...など... / ...なり...\\
\textbf{Variants (EN)}: ...and...and... / ...or...or... / ...such as...

\textbf{Intent (JA)}: 列挙, 例示, 曖昧\\
\textbf{Intent (EN)}: listing, example, vagueness

\textbf{Tags (JA)}: 即時文法, 列挙, 感情\\
\textbf{Tags (EN)}: immediate grammar, listing, emotion

\textbf{Adjusted Expression (JA)}:\\
(本:ほん)や(雑誌:ざっし)などを(整理:せいり)しました。\\
\textbf{Adjusted Expression (EN)}:\\
I sorted out books, magazines, and such.

\textbf{Example (JA)}:\\
A「(全部:ぜんぶ)、(捨:す)てちゃったの？」 B「ああ、(本:ほん)とか(雑誌:ざっし)とか...いろいろ。」\\
\textbf{Example (EN)}:\\
A: Did you throw everything away? B: Yeah, books, magazines, and all sorts of stuff...

\textbf{Notes (JA)}:\\
「...とか...とか」は、複数のものを挙げて曖昧に例示する表現です。言い切らないことで柔らかい印象を与えます。

\textbf{Notes (EN)}:\\
"...or...or..." is used to list multiple things vaguely. Not stating precisely makes the expression softer.

\newpage

\section*{448. I ate too much}

\textbf{Date}: 2025.07.01 (Tue)\\
\textbf{Index}: tabesugita\\
\textbf{Expression (JA)}: (食:た)べすぎた\\
\textbf{Romanization}: Tabesugita\\
\textbf{Expression (EN)}: I ate too much

\textbf{Variants (JA)}: たべすぎた / 食べすぎちゃった / 食べすぎてしまった\\
\textbf{Variants (EN)}: I ate too much / I overate / I stuffed myself

\textbf{Intent (JA)}: 後悔, 反省, 気づき\\
\textbf{Intent (EN)}: regret, reflection, realization

\textbf{Tags (JA)}: 即時文法, 感情, 身体感覚\\
\textbf{Tags (EN)}: immediate grammar, emotion, physical sensation

\textbf{Adjusted Expression (JA)}:\\
かなり多量に食べてしまいました。\\
\textbf{Adjusted Expression (EN)}:\\
I ended up eating quite a lot.

\textbf{Example (JA)}:\\
A「どうしたの？」B「(食:た)べすぎた。(苦:くる)しい」\\
\textbf{Example (EN)}:\\
A: What's wrong? B: I ate too much... My stomach hurts.

\textbf{Notes (JA)}:\\
「(食:た)べすぎた」は、その瞬間に体の違和感や後悔を感じて出る表現です。言い訳や説明より先に、感情が発話される表現です。

\textbf{Notes (EN)}:\\
"Tabesugita" is an immediate expression of realization and regret about overeating. It often comes out before any explanation.

\newpage

\section*{447. If possible / If you can}

\textbf{Date}: 2025.06.30 (Mon)\\
\textbf{Index}: dekireba\\
\textbf{Expression (JA)}: できれば\\
\textbf{Romanization}: Dekireba\\
\textbf{Expression (EN)}: If possible

\textbf{Variants (JA)}: できれば / もしできれば / できたら\\
\textbf{Variants (EN)}: If possible / If you can / If it's okay

\textbf{Intent (JA)}: 希望, 依頼, 遠慮\\
\textbf{Intent (EN)}: hope, request, politeness

\textbf{Tags (JA)}: 即時文法, 依頼, 丁寧\\
\textbf{Tags (EN)}: immediate grammar, request, politeness

\textbf{Adjusted Expression (JA)}:\\
もし(可能:かのう)であればお(願:ねが)いします。\\
\textbf{Adjusted Expression (EN)}:\\
If it's possible, please.

\textbf{Example (JA)}:\\
A「できれば、(今日中:きょうじゅう)」\\
\textbf{Example (EN)}:\\
A: If possible, by today.

\textbf{Notes (JA)}:\\
「できれば」は相手に配慮しつつ自分の希望を伝える表現です。言い切りを避けることで、柔らかい印象になります。

\textbf{Notes (EN)}:\\
"Dekireba" expresses a polite hope or request, softening the tone by avoiding a direct command.

\newpage

\section*{446. Oops / Damn / Oh no}

\textbf{Date}: 2025.06.29 (Sun)\\
\textbf{Index}: shimatta\\
\textbf{Expression (JA)}: しまった\\
\textbf{Romanization}: Shimatta\\
\textbf{Expression (EN)}: Oops

\textbf{Variants (JA)}: しまった / あ、しまった / うわ、しまった\\
\textbf{Variants (EN)}: Oops / Oh no / Damn

\textbf{Intent (JA)}: 後悔, 気づき, 失敗\\
\textbf{Intent (EN)}: regret, realization, mistake

\textbf{Tags (JA)}: 即時文法, 感情, 反応\\
\textbf{Tags (EN)}: immediate grammar, emotion, reaction

\textbf{Adjusted Expression (JA)}:\\
(失敗:しっぱい)してしまった。\\
\textbf{Adjusted Expression (EN)}:\\
I made a mistake.

\textbf{Example (JA)}:\\
A「あ、しまった」 B「どうしたの」A「(携帯:けいたい)、(忘:わす)れた」\\
\textbf{Example (EN)}:\\
A: Oops. B: What happened? A: I forgot my phone.

\textbf{Notes (JA)}:\\
「しまった」は、その瞬間に失敗や後悔に気づいたときに出る短い感情表現です。状況を説明するより先に感情が口をついて出ます。

\textbf{Notes (EN)}:\\
"Shimatta" is a short exclamation expressing sudden realization of a mistake or regret. It's used immediately upon noticing an error, often before explaining the situation.

\newpage

\section*{445. feeling uneasy}

\textbf{Date}: 2025.06.28 (Sat)\\
\textbf{Index}: hiyahiya\\
\textbf{Expression (JA)}: ひやひやする\\
\textbf{Romanization}: Hiyahiya suru\\
\textbf{Expression (EN)}: I feel uneasy

\textbf{Variants (JA)}: ひやひやする / 胸がひやひやする / 冷や汗が出る\\
\textbf{Variants (EN)}: I feel uneasy / I feel nervous / It makes me break into a cold sweat

\textbf{Intent (JA)}: 不安, 恐れ, 緊張\\
\textbf{Intent (EN)}: anxiety, fear, nervousness

\textbf{Tags (JA)}: 即時文法, 感情, 身体感覚\\
\textbf{Tags (EN)}: immediate grammar, emotion, physical sensation

\textbf{Adjusted Expression (JA)}:\\
(危:あぶ)なくて(冷:ひ)や(汗:あせ)が(出:で)る。\\
\textbf{Adjusted Expression (EN)}:\\
It feels dangerous, and I break into a cold sweat.

\textbf{Example (JA)}:\\
A「(道:みち)が(凍:こお)っていて、(歩:ある)くのがひやひやするね」\\
\textbf{Example (EN)}:\\
A: The road is frozen, so I feel nervous walking on it.

\textbf{Notes (JA)}:\\
「ひやひやする」は、危険や不安を感じて、冷や汗が出るような気持ちを表すカジュアルな表現です。ちょっと怖い状況や緊張感を伝えるときに使います。

\textbf{Notes (EN)}:\\
"Hiyahiya suru" is a casual expression that describes feeling anxious or nervous, as if breaking into a cold sweat. It conveys a sense of fear or unease in potentially dangerous situations.

\newpage

\section*{444. being so excited}

\textbf{Date}: 2025.06.27 (Fri)\\
\textbf{Index}: wakuwaku\\
\textbf{Expression (JA)}: わくわくする\\
\textbf{Romanization}: Wakuwaku suru\\
\textbf{Expression (EN)}: I am so excited

\textbf{Variants (JA)}: わくわくする / 胸がわくわくする / 楽しみでわくわくする\\
\textbf{Variants (EN)}: I am so excited / I feel thrilled / I am looking forward to it

\textbf{Intent (JA)}: 期待, 楽しみ, 高揚感\\
\textbf{Intent (EN)}: anticipation, excitement, elevation

\textbf{Tags (JA)}: 即時文法, 感情, 身体感覚\\
\textbf{Tags (EN)}: immediate grammar, emotion, physical sensation

\textbf{Adjusted Expression (JA)}:\\
(気持:きも)ちが(盛:も)り(上:あ)がっている。\\
\textbf{Adjusted Expression (EN)}:\\
I am feeling uplifted.

\textbf{Example (JA)}:\\
A「(息子:むすこ)は、(明日:あす)、(球場:きゅうじょう)に(行:い)けるので、わくわくしているよ」\\
\textbf{Example (EN)}:\\
A: My son is so excited because he can go to the stadium tomorrow.

\textbf{Notes (JA)}:\\
「わくわくしている」は、何かを楽しみにして心が高まる気持ちを表すカジュアルな表現です。前向きな期待感を伝えるときに使います。自分以外の感情を表現するときには「わくわくしている」と「わくわくしている？」のように「している」を使いますが、自分の感情を表現するときには「わくわくする」と「わくわくしている」のように両方使えます。

\textbf{Notes (EN)}:\\
"Wakuwaku shita" is a casual expression describing the uplifting feeling of excited anticipation. It is used to convey a positive sense of looking forward to something. When expressing someone else's emotions, you can use "wakuwaku shite iru" or "wakuwaku shite iru?" but when expressing your own feelings, both "wakuwaku suru" and "wakuwaku shite iru" can be used.

\newpage

\section*{443. My heart was pounding}

\textbf{Date}: 2025.06.26 (Thu)\\
\textbf{Index}: dokidoki\\
\textbf{Expression (JA)}: どきどきした\\
\textbf{Romanization}: Dokidoki shita\\
\textbf{Expression (EN)}: My heart was pounding

\textbf{Variants (JA)}: どきどきした / 胸がどきどきした / 心臓がばくばくした\\
\textbf{Variants (EN)}: My heart was pounding / My heart was racing / My heart thumped

\textbf{Intent (JA)}: 緊張, 期待, 不安\\
\textbf{Intent (EN)}: nervousness, anticipation, anxiety

\textbf{Tags (JA)}: 即時文法, 感情, 身体感覚\\
\textbf{Tags (EN)}: immediate grammar, emotion, physical sensation

\textbf{Adjusted Expression (JA)}:\\
(緊張:きんちょう)して(心臓:しんぞう)がばくばくしてしまった。\\
\textbf{Adjusted Expression (EN)}:\\
I was so nervous that my heart started racing.

\textbf{Example (JA)}:\\
A「(初:はじ)めての(デート:でーと)で、どきどきした」\\
\textbf{Example (EN)}:\\
A: I was so nervous on my first date.

\textbf{Notes (JA)}:\\
「どきどきした」は、緊張や期待、不安などの感情を表すカジュアルな表現です。特に、心臓の鼓動が速くなる様子を表現します。カジュアルな会話でよく使われ、相手に感情を伝える役割があります。

\textbf{Notes (EN)}:\\
"Dokidoki shita" is a casual expression that conveys feelings of nervousness, anticipation, or anxiety. It particularly describes the sensation of a racing heartbeat. This phrase is commonly used in informal conversations to express emotions to the listener.

\newpage

\section*{442. That's what I wanted to say}

\textbf{Date}: 2025.06.25 (Wed)\\
\textbf{Index}: iitaisorega\\
\textbf{Expression (JA)}: (言:い)いたい、それが\\
\textbf{Romanization}: Iitai, Sore ga\\
\textbf{Expression (EN)}: That's what I wanted to say

\textbf{Variants (JA)}: 言いたいんだよ、それが / 言いたいの、それが / まさに言いたかった、それが\\
\textbf{Variants (EN)}: That's what I wanted to say / Exactly that! / That's the point I was trying to make

\textbf{Intent (JA)}: 強調, 納得, 共感\\
\textbf{Intent (EN)}: emphasis, agreement, shared feeling

\textbf{Tags (JA)}: 即時文法, 倒置, 感情, 倒置-述部先行\\
\textbf{Tags (EN)}: immediate grammar, inversion, emotion, structure-pattern

\textbf{Adjusted Expression (JA)}:\\
私はそれを言いたかったんです。\\
\textbf{Adjusted Expression (EN)}:\\
That's what I've been trying to say all along.

\textbf{Example (JA)}:\\
A「(言:い)いたいんだよ、それが」B「うん、(分:わか)るよ」\\
\textbf{Example (EN)}:\\
A: That's what I wanted to say. B: Yeah, I get it.

\textbf{Notes (JA)}:\\
「言いたい、それが」は、即時に表現するときに通常の語順を変えることで、強調や感情を伝えるカジュアルな表現です。このような表現を倒置と呼び、重要な情報を通常の語順とは関係なく、前に置くことができます。しかし、後ろに移動した後の関係がよく分かるように、助詞をつけると文がより明確になります。「これを食べる」のように「を」の場合は、だいたいにおいて「食べる、これ」のように「を」を省略しますが、他の関係、たとえば「が」「は」「に」「で」「を(場所)」の場合は、後ろに移動した後も助詞ををつけることを自然な感じがします。

\textbf{Notes (EN)}:\\
"Iitai, sore ga" is a casual expression that emphasizes or conveys emotion by changing the usual word order. This type of expression is called inversion, where important information is placed at the front regardless of the normal word order. However, adding particles after moving the information to the back makes the sentence clearer. For example, in "kore o taberu" (to eat this), the object particle "o" can often be omitted when inverted to "taberu, kore." However, for other particles like "ga," "wa," "ni," "de," and "o (location)," it feels more natural to keep them even after inversion.

\newpage

\section*{441. completely / entirely}

\textbf{Date}: 2025.06.24 (Tue)\\
\textbf{Index}: sukkari\\
\textbf{Expression (JA)}: すっかり\\
\textbf{Romanization}: Sukkari\\
\textbf{Expression (EN)}: completely / entirely

\textbf{Variants (JA)}: すっかり / 完全に / まるっきり / すっかりと\\
\textbf{Variants (EN)}: completely / entirely / fully / totally

\textbf{Intent (JA)}: 完了, 変化, 驚き\\
\textbf{Intent (EN)}: completion, change, surprise

\textbf{Tags (JA)}: 即時文法, 副詞, 状態変化\\
\textbf{Tags (EN)}: immediate grammar, adverb, change of state

\textbf{Adjusted Expression (JA)}:\\
完全に日が沈み、あたりは薄闇に包まれていた。\\
\textbf{Adjusted Expression (EN)}:\\
The sun had completely set, and everything was shrouded in dusk.

\textbf{Example (JA)}:\\
A「(宿題:しゅくだい)、(終:お)わった？」B「ああ、(宿題:しゅくだい)ね、すっかり、(忘:わす)れてた」\\
\textbf{Example (EN)}:\\
A: Did you finish your homework? B: Oh, the homework? I completely forgot about it.

\textbf{Notes (JA)}:\\
「すっかり」は、何かが完全に終わったり、変わったりしたことを表すカジュアルな表現です。「completely」や「entirely」と訳されます。特に、何かが完全に変化したり、忘れ去られたりしたときに使われます。

\textbf{Notes (EN)}:\\
"Sukkari" is a casual expression that indicates something has completely ended or changed. It can be translated as "completely" or "entirely." This phrase is often used when something has undergone a complete transformation or has been forgotten.

\newpage

\section*{440. Before I knew it / When I realized it}

\textbf{Date}: 2025.06.23 (Mon)\\
\textbf{Index}: kigatsuitara\\
\textbf{Expression (JA)}: (気:き)がついたら\\
\textbf{Romanization}: Ki ga tsuitara\\
\textbf{Expression (EN)}: Before I knew it / When I realized it

\textbf{Variants (JA)}: 気がついたら / いつの間にか / 知らないうちに / ふと気づくと\\
\textbf{Variants (EN)}: Before I knew it / When I realized it / Before I noticed / Suddenly I noticed

\textbf{Intent (JA)}: 無意識, 時間経過, 驚き, 気づき\\
\textbf{Intent (EN)}: unconscious, passage of time, surprise, awareness

\textbf{Tags (JA)}: 即時文法, 時制, 気づき\\
\textbf{Tags (EN)}: immediate grammar, tense, realization

\textbf{Adjusted Expression (JA)}:\\
知らないうちに時間が過ぎていて、卒業の年を迎えていた。\\
\textbf{Adjusted Expression (EN)}:\\
Time had passed without me realizing it, and I found myself in my final year of school.

\textbf{Example (JA)}:\\
A「(気:き)がついたら、もう(夏休:なつやす)みだね」B「うん、(早:はや)いね」\\
\textbf{Example (EN)}:\\
A: Before I knew it, it's already summer vacation. B: Yeah, time flies.

\textbf{Notes (JA)}:\\
「気がついたら」は、何かに気づくまでの時間の経過や無意識のうちに起こったことを表すカジュアルな表現です。「Before I knew it」や「When I realized it」と訳されます。特に、何かがいつの間にか変わっていたり、時間が経っていたりすることを伝えるときに使われます。

\textbf{Notes (EN)}:\\
"Ki ga tsuitara" is a casual expression that indicates the passage of time or something that happened unconsciously until one becomes aware of it. It can be translated as "Before I knew it" or "When I realized it." This phrase is often used to convey that something has changed or that time has passed without the speaker noticing.

\newpage

\section*{439. awkward / uncomfortable}

\textbf{Date}: 2025.06.22 (Sun)\\
\textbf{Index}: kimazui\\
\textbf{Expression (JA)}: (気:き)まずい\\
\textbf{Romanization}: Kimazui\\
\textbf{Expression (EN)}: awkward / uncomfortable

\textbf{Variants (JA)}: きまずい / ちょっときまずいね / なんか空気が重い / 気まずい / 気まずくなった / ちょっと気まずい雰囲気になった\\
\textbf{Variants (EN)}: awkward / awkward silence / felt uncomfortable / a bit uncomfortable / kind of tense

\textbf{Intent (JA)}: 感情, 対人関係, 気まずさ, 遠慮, 場の空気, 沈黙, 誤解, 不快感\\
\textbf{Intent (EN)}: emotion, interpersonal, awkwardness, reserve, social tension, silence, misunderstanding, discomfort

\textbf{Tags (JA)}: 即時文法, 対人, 感情, 状況描写\\
\textbf{Tags (EN)}: immediate grammar, interpersonal, emotion, situation description

\textbf{Adjusted Expression (JA)}:\\
いたたまれない(雰囲気:ふんいき)になってしまった\\
\textbf{Adjusted Expression (EN)}:\\
It became an unbearable atmosphere.

\textbf{Example (JA)}:\\
A「(昨日:きのう)の(話:はなし)、(気:き)まずいね」B「うん、ちょっとね」\\
\textbf{Example (EN)}:\\
A: Yesterday's conversation was awkward, wasn't it? B: Yeah, a bit.

\textbf{Notes (JA)}:\\
「きまずい」は、対人関係において気まずさや不快感を表すカジュアルな表現です。「awkward」や「uncomfortable」と訳されます。特に、会話が途切れたり、誤解が生じたりしたときに使われることが多いです。相手との関係がぎくしゃくしているときに使うことが一般的です。

\textbf{Notes (EN)}:\\
"Kimazui" is a casual expression that indicates awkwardness or discomfort in interpersonal relationships. It can be translated as "awkward" or "uncomfortable." This phrase is often used when conversations break down or misunderstandings occur. It is commonly used when the relationship with the other person feels strained.

\newpage

\section*{438. No thanks / I don't need it}

\textbf{Date}: 2025.06.21 (Sat)\\
\textbf{Index}: iranai\\
\textbf{Expression (JA)}: いらない\\
\textbf{Romanization}: Iranai\\
\textbf{Expression (EN)}: No thanks / I don't need it

\textbf{Variants (JA)}: いらない / もういらない / ううん、いらない / それは、いらないかな\\
\textbf{Variants (EN)}: No thanks / I don't need it / I'm good / Not necessary

\textbf{Intent (JA)}: 拒否, 遠慮, 完了, 冷却\\
\textbf{Intent (EN)}: refusal, restraint, completion, detachment

\textbf{Tags (JA)}: 即時文法, 会話, 反応, 対人関係\\
\textbf{Tags (EN)}: immediate grammar, conversation, reaction, interpersonal

\textbf{Adjusted Expression (JA)}:\\
お(気持:きも)ちだけいただいておきます。(今:いま)は(大丈夫:だいじょうぶ)です。\\
\textbf{Adjusted Expression (EN)}:\\
I appreciate the gesture, but I'm fine for now.

\textbf{Example (JA)}:\\
A「これ、いる？」B「ううん、いらない」\\
\textbf{Example (EN)}:\\
A: Do you need this? B: No thanks, I don't need it.

\textbf{Notes (JA)}:\\
「いらない」は、何かを必要としないことを表すカジュアルな表現です。「No thanks」や「I don't need it」と訳されます。特に、相手からの提案や贈り物に対して使われることが多いです。相手に対して遠慮や拒否の意志を示す役割があります。

\textbf{Notes (EN)}:\\
"Iranai" is a casual expression that indicates that something is not needed. It can be translated as "No thanks" or "I don't need it." This phrase is often used in response to offers or gifts from others. It serves to show restraint or refusal towards the other person.

\newpage

\section*{437. All of a sudden / Out of nowhere}

\textbf{Date}: 2025.06.20 (Fri)\\
\textbf{Index}: ikinari\\
\textbf{Expression (JA)}: いきなり\\
\textbf{Romanization}: Ikinari\\
\textbf{Expression (EN)}: All of a sudden / Out of nowhere

\textbf{Variants (JA)}: いきなり / いきなり話しかけられた / いきなり走り出した / いきなり泣き出した\\
\textbf{Variants (EN)}: All of a sudden / Suddenly, someone talked to me / Suddenly started running / Suddenly burst into tears

\textbf{Intent (JA)}: 驚き, 即時反応, 状況変化\\
\textbf{Intent (EN)}: surprise, immediate reaction, sudden change

\textbf{Tags (JA)}: 即時文法, 副詞, 感情\\
\textbf{Tags (EN)}: immediate grammar, adverb, emotion

\textbf{Adjusted Expression (JA)}:\\
(突然:とつぜん)の(出来事:できごと)に(驚:おどろ)きましたが、すぐに(対応:たいおう)しました。\\
\textbf{Adjusted Expression (EN)}:\\
I was surprised by the sudden event, but I responded immediately.

\textbf{Example (JA)}:\\
いきなり(質問:しつもん)されてもね。(何:なん)て(言:い)えばいいかね？」\\
\textbf{Example (EN)}:\\
Even if you ask me a question all of a sudden, I don't know what to say.

\textbf{Notes (JA)}:\\
「いきなり」は、予期せぬ出来事や行動を表すカジュアルな表現です。「All of a sudden」や「Out of nowhere」と訳されます。カジュアルな会話でよく使われ、相手に対して驚きを伝える役割があります。特に、何かが突然起こったときに使われることが多いです。

\textbf{Notes (EN)}:\\
"Ikinari" is a casual expression that indicates an unexpected event or action. It can be translated as "All of a sudden" or "Out of nowhere." This phrase is commonly used in informal conversations to convey surprise to the listener. It is often used when something happens suddenly.

\newpage

\section*{436. That, I Understand / I Know That}

\textbf{Date}: 2025.06.19 (Thu)\\
\textbf{Index}: sorewakaru\\
\textbf{Expression (JA)}: それ、わかる\\
\textbf{Romanization}: Sore, wakaru\\
\textbf{Expression (EN)}: That, I know / I understand that

\textbf{Variants (JA)}: それ、わかる / 本当にそれ、わかる / まさにそれだよね / 分かる分かる\\
\textbf{Variants (EN)}: I know that / That really resonates / Exactly, I understand / I totally get it

\textbf{Intent (JA)}: 共感, 理解, 同調, 安心感の提供\\
\textbf{Intent (EN)}: empathy, understanding, agreement, providing reassurance

\textbf{Tags (JA)}: 即時文法, 共感, 会話, リアクション\\
\textbf{Tags (EN)}: immediate grammar, empathy, conversation, reaction

\textbf{Adjusted Expression (JA)}:\\
本当にその通りで、私も同じ気持ちです。\\
\textbf{Adjusted Expression (EN)}:\\
Exactly, I feel the same way.

\textbf{Example (JA)}:\\
A「この(状況:じょうきょう)、ちょっと(難:むずか)しいね」B「うん、でも、それ、わかる」\\
\textbf{Example (EN)}:\\
A: This situation is a bit tough. B: Yeah, I totally get that.

\textbf{Notes (JA)}:\\
「それ、わかる」は、相手の意見や感情に共感するカジュアルな表現です。「I know that」や「I understand that」と訳されます。カジュアルな会話でよく使われ、相手に対して理解を示す役割があります。特に、相手の気持ちを受け入れることで、安心感を提供することができます。

\textbf{Notes (EN)}:\\
"Sore, wakaru" is a casual expression used to empathize with the other person's opinion or feelings. It can be translated as "I know that" or "I understand that." This phrase is commonly used in informal conversations to show understanding towards the other person. By accepting the other person's feelings, it can provide a sense of reassurance.

\newpage

\section*{435. Just a little thing / Nothing big}

\textbf{Date}: 2025.06.18 (Wed)\\
\textbf{Index}: chottoshita.koto\\
\textbf{Expression (JA)}: ちょっとしたこと\\
\textbf{Romanization}: Chotto shita koto\\
\textbf{Expression (EN)}: Just a little thing / Nothing big

\textbf{Variants (JA)}: ちょっとしたこと / ちょっとしたことで / そんなちょっとしたことで / ちょっとしたことで気が滅入る\\
\textbf{Variants (EN)}: Just a little thing / Something minor / Over something trivial / A small thing can really get to you

\textbf{Intent (JA)}: 軽視, 謙遜, 共感, 心の揺れ, 感情のきっかけ\\
\textbf{Intent (EN)}: downplaying, humility, empathy, emotional trigger, small cause

\textbf{Tags (JA)}: 即時文法, 感情, 評価, 状況描写\\
\textbf{Tags (EN)}: immediate grammar, emotion, evaluation, situation description

\textbf{Adjusted Expression (JA)}:\\
ほんの些細なことが、心に残ってしまいました。\\
\textbf{Adjusted Expression (EN)}:\\
Something trivial ended up lingering in my mind.

\textbf{Example (JA)}:\\
ちょっとしたことだが、(簡単:かんたん)じゃない。\\
\textbf{Example (EN)}:\\
Because of a trivial matter, it's not easy.

\textbf{Notes (JA)}:\\
「ちょっとしたこと」は、何かが小さくて重要でないことを示すカジュアルな表現です。「Just a little thing」や「Nothing big」と訳されます。カジュアルな会話でよく使われ、相手に対して軽い気持ちで話しかける役割があります。しかし、時には小さなことが大きな影響を与えることもあるので、注意が必要だという意味にもなります。

\textbf{Notes (EN)}:\\
"Chotto shita koto" is a casual expression that indicates something is small and not significant. It can be translated as "Just a little thing" or "Nothing big." This phrase is commonly used in informal conversations to approach the listener with a light-hearted tone. However, it can also imply that even small matters can have a significant impact, so caution is advised.

\newpage

\section*{434. Just a little longer / Please wait a bit}

\textbf{Date}: 2025.06.17 (Tue)\\
\textbf{Index}: mousukoshi.matte\\
\textbf{Expression (JA)}: もう(少:すこ)し、(待:ま)って\\
\textbf{Romanization}: Mō sukoshi, matte\\
\textbf{Expression (EN)}: Just a little longer / Please wait a bit

\textbf{Variants (JA)}: もう少し、待って / あとちょっと、待って / ちょっとだけ、待って / 少しだけ、待ってて\\
\textbf{Variants (EN)}: Just a little longer / Wait a bit more / Please wait a little longer / Hold on a second

\textbf{Intent (JA)}: 依頼, 猶予, 調整\\
\textbf{Intent (EN)}: request, delay, adjustment

\textbf{Tags (JA)}: 即時文法, 会話, 時間\\
\textbf{Tags (EN)}: immediate grammar, conversation, time

\textbf{Adjusted Expression (JA)}:\\
もう少し時間をいただければ、対応できます。\\
\textbf{Adjusted Expression (EN)}:\\
If I could have a little more time, I’ll be able to take care of it.

\textbf{Example (JA)}:\\
A「もうすぐ(終:お)わるから、もう(少:すこ)し、(待:ま)って」B「うん、(大丈夫:だいじょうぶ)だよ」\\
\textbf{Example (EN)}:\\
A: It'll be done soon, so please wait a bit. B: Okay, no problem.

\textbf{Notes (JA)}:\\
「もう少し、待って」は、相手に時間をかけてもらうように頼むカジュアルな表現です。特に、何かを終わらせるためにもう少し時間が必要なときに使われます。

\textbf{Notes (EN)}:\\
"Mou sukoshi, matte" is a casual expression used to ask someone to give you a little more time. It is particularly used when you need a bit more time to finish something.

\newpage

\section*{433. Just now / A moment ago}

\textbf{Date}: 2025.06.16 (Mon)\\
\textbf{Index}: sakki\\
\textbf{Expression (JA)}: さっき\\
\textbf{Romanization}: Sakki\\
\textbf{Expression (EN)}: Just now / A moment ago

\textbf{Variants (JA)}: さっき / ついさっき / さっきまで / さっき言ってたやつ\\
\textbf{Variants (EN)}: Just now / A moment ago / Until a moment ago / What I mentioned earlier

\textbf{Intent (JA)}: 直前, 切り出し, 再開, 回想, 時間感覚\\
\textbf{Intent (EN)}: just before, starting point, resumption, recall, sense of time

\textbf{Tags (JA)}: 即時文法, 時間, 会話の流れ\\
\textbf{Tags (EN)}: immediate grammar, time, conversational flow

\textbf{Adjusted Expression (JA)}:\\
さっき駅で見かけた人、どこかで見たことある気がする。\\
\textbf{Adjusted Expression (EN)}:\\
The person I saw at the station just now looked familiar somehow.

\textbf{Example (JA)}:\\
さっきまでは雨だったのに、今は晴れてる。\\
\textbf{Example (EN)}:\\
It was raining just a moment ago, but now it's sunny.

\textbf{Notes (JA)}:\\
「さっき」は、直前の出来事や時間を指すカジュアルな表現で、会話の流れをスムーズにする役割があります。特に、何かを思い出したり、再開したりするときに使われます。

\textbf{Notes (EN)}:\\
"Sakki" is a casual expression that refers to a recent event or time, serving to smooth the flow of conversation. It is particularly useful when recalling something or resuming a discussion.

\newpage

\section*{432. Together, always}

\textbf{Date}: 2025.06.15 (Sun)\\
\textbf{Index}: zutto.issho\\
\textbf{Expression (JA)}: ずっと、いっしょ\\
\textbf{Romanization}: Zutto, issho\\
\textbf{Expression (EN)}: Together, always

\textbf{Variants (JA)}: ずっと、いっしょ / いつまでも、いっしょ / これからも、ずっといっしょ\\
\textbf{Variants (EN)}: Together, always / Forever together / We’ll always be together

\textbf{Intent (JA)}: 愛情, 約束, 安心\\
\textbf{Intent (EN)}: affection, commitment, reassurance

\textbf{Tags (JA)}: 即時文法, 感情, 親密さ\\
\textbf{Tags (EN)}: immediate grammar, emotion, closeness

\textbf{Adjusted Expression (JA)}:\\
これから先も、あなたと一緒にいたいと思っています。\\
\textbf{Adjusted Expression (EN)}:\\
I want to stay with you from now on, always.

\textbf{Example (JA)}:\\
A「これからも？」B「うん、ずっと、いっしょ。」\\
\textbf{Example (EN)}:\\
A: From now on too? B: Yeah, together, always.

\textbf{Notes (JA)}:\\
「ずっと、いっしょ」は、愛情や約束を強調する表現で、相手との永続的な関係を示します。カジュアルな会話でよく使われ、親密さや安心感を伝える役割があります。「いっしょ」には「一緒」という意味があり、相手と共にいるという意味もありますが、「同じ」という意味もあり、「シャツの色、いっしょ」のように使えます。

\textbf{Notes (EN)}:\\
"Zutto, issho" is an expression that emphasizes affection and commitment, indicating a lasting relationship with the other person. It is commonly used in casual conversations to convey closeness and reassurance. The term "issho" means "together," but it can also mean "the same," as in "Shatsu no iro, issho./the color of the shirt is the same."

\newpage

\section*{431. Anytime / Whenever}

\textbf{Date}: 2025.06.14 (Sat)\\
\textbf{Index}: itsudemo\\
\textbf{Expression (JA)}: いつでも\\
\textbf{Romanization}: Itsudemo\\
\textbf{Expression (EN)}: Anytime / Whenever

\textbf{Variants (JA)}: いつでも / いつでもいいよ / いつでもどうぞ / いつでも言ってね\\
\textbf{Variants (EN)}: Anytime / Whenever's fine / Feel free anytime / Just let me know whenever

\textbf{Intent (JA)}: 柔軟, 受容, 歓迎, 許可\\
\textbf{Intent (EN)}: flexibility, openness, invitation, permission

\textbf{Tags (JA)}: 即時文法, 応答, 寛容, 会話\\
\textbf{Tags (EN)}: immediate grammar, response, tolerance, conversation

\textbf{Adjusted Expression (JA)}:\\
ご都合の良いときに、いつでもご連絡ください。\\
\textbf{Adjusted Expression (EN)}:\\
Please feel free to contact me at your convenience.

\textbf{Example (JA)}:\\
A「じゃ、こんどいい？」B「うん、いつでも」\\
\textbf{Example (EN)}:\\
A: So, how about next time? B: Yeah, anytime.

\textbf{Notes (JA)}:\\
「いつでも」は、相手に選択の自由を与えることで、柔らかな応答を実現する表現です。具体的な時間を明示しなくても会話が成立する点で、即時文法の特徴が見られます。時間を指定しないまま、会話のテンポや関係性に委ねる姿勢が込められています。

\textbf{Notes (EN)}:\\
"Itsudemo" is a flexible and accepting expression that allows the other person to choose the time. It exemplifies immediate grammar in that it permits smooth conversation without specifying details like time. It conveys a stance of openness and willingness.

\newpage

\section*{430. For now / Just in case / Tentatively}

\textbf{Date}: 2025.06.13 (Fri)\\
\textbf{Index}: ichiou\\
\textbf{Expression (JA)}: (一応:いちおう)\\
\textbf{Romanization}: Ichiou\\
\textbf{Expression (EN)}: For now / Just in case / Tentatively

\textbf{Variants (JA)}: 一応 / 一応、確認だけしとくね / 一応、終わったことにしておこう / 一応はね\\
\textbf{Variants (EN)}: For now / Just in case / Tentatively / Kind of

\textbf{Intent (JA)}: 保留, 控えめ, 不確か, 前置き\\
\textbf{Intent (EN)}: tentative, hedging, uncertainty, preliminary

\textbf{Tags (JA)}: 即時文法, あいまい, 会話, 態度\\
\textbf{Tags (EN)}: immediate grammar, ambiguity, conversation, stance

\textbf{Adjusted Expression (JA)}:\\
概ね、計画し、準備したものは完了したのですが、まだ調整が若干必要となる見込みです。\\
\textbf{Adjusted Expression (EN)}:\\
The planning and preparations are largely complete, but some minor adjustments may still be necessary.

\textbf{Example (JA)}:\\
A「(準備:じゅんび)できた？」B「うん、(一応:いちおう)ね」\\
\textbf{Example (EN)}:\\
A: Are you ready? B: Yeah, for now.

\textbf{Notes (JA)}:\\
「一応」は、何かが形式上、または暫定的に済んでいることを伝えるときによく使われます。自信がなかったり、完全ではないと感じているときにも使われ、「たぶん大丈夫だけど、まだ不安」な気持ちが込められています。つまり、「一応」の裏の意味としては、「まだ不十分」だということですね。「一応」と言われたら、「もう一度、確認だけしとくね」と言って、ダブルチェックしたほうがよいかもしれません。

\textbf{Notes (EN)}:\\
The expression "ichiou" implies that something is done in a formal or tentative manner. It is often used when the speaker is not entirely confident or feels that something is not yet complete. The underlying meaning of "ichiou" suggests that there may still be some uncertainty or incompleteness. When someone says "ichiou," it might be wise to respond with, "Let me double-check just in case."

\newpage

\section*{429. While I'm at it / Incidentally}

\textbf{Date}: 2025.06.12 (Thu)\\
\textbf{Index}: tsuideni\\
\textbf{Expression (JA)}: ついでに\\
\textbf{Romanization}: Tsuideni\\
\textbf{Expression (EN)}: While I'm at it / Incidentally

\textbf{Variants (JA)}: ついでに / そのついでに / ついでと言ってはなんだけど / 話のついでに\\
\textbf{Variants (EN)}: While I'm at it / Along the way / As an aside / By the way

\textbf{Intent (JA)}: 追加, 軽さ, 策略, 含み\\
\textbf{Intent (EN)}: addition, casualness, tact, implication

\textbf{Tags (JA)}: 即時文法, 会話, 接続, 行動\\
\textbf{Tags (EN)}: immediate grammar, conversation, connection, action

\textbf{Adjusted Expression (JA)}:\\
この(機会:きかい)を(利用:りよう)して、もうひとつご(提案:ていあん)するのはいかがでしょうか。\\
\textbf{Adjusted Expression (EN)}:\\
May I take this opportunity to make one additional suggestion?

\textbf{Example (JA)}:\\
A「(コンビニ:こんびに)、(行:い)くの？」B「うん、(ついで:)に(何:なに)か(買:か)って(来:こ)ようか」\\
\textbf{Example (EN)}:\\
A: You're going to the convenience store? B: Yeah, do you want me to pick something up while I'm at it?

\textbf{Notes (JA)}:\\
「ついでに」は本来の目的に付け足す行動や発言に使われます。この例のように、何かをするついでに誰かのために、別の仕事をしてあげようと提案するときに使いますよ。

\textbf{Notes (EN)}:\\
"Tsuideni" is used to indicate an action or statement that is added to the main purpose. In this example, it is used to suggest doing an additional task for someone while going somewhere or doing something else.

\newpage

\section*{428. It depends}

\textbf{Date}: 2025.06.11 (Wed)\\
\textbf{Index}: baai.ni.yorikeri\\
\textbf{Expression (JA)}: (場合:ばあい)によりけり\\
\textbf{Romanization}: Baai ni yorikeri\\
\textbf{Expression (EN)}: It depends on the case

\textbf{Variants (JA)}: 場合によりけり / 場合によっては / 状況しだい / ケースバイケース\\
\textbf{Variants (EN)}: It depends on the case / Depending on the situation / Depends on the circumstances / Case by case

\textbf{Intent (JA)}: 柔軟, 保留, 思慮, 配慮\\
\textbf{Intent (EN)}: flexibility, non-commitment, consideration, carefulness

\textbf{Tags (JA)}: 即時文法, 判断, 返答, 含み\\
\textbf{Tags (EN)}: immediate grammar, judgment, response, implication

\textbf{Adjusted Expression (JA)}:\\
この(件:けん)は、(場合:ばあい)によりけりだと(思:おも)います。\\
\textbf{Adjusted Expression (EN)}:\\
I think it depends on the case.

\textbf{Example (JA)}:\\
A「(この:この)問題、(どう:どう)する？」B「(場合:ばあい)によりけりだね」\\
\textbf{Example (EN)}:\\
A: What should we do about this issue? B: It depends on the case.

\textbf{Notes (JA)}:\\
「場合によりけり」は、状況や条件によって異なることを示すカジュアルな表現です。「It depends on the case」や「Depending on the situation」と訳されます。即時的に簡単に判断できないことは多いと思います。とりあえず「場合によりけり」と言っておけば、その場をしのげ、その時間稼ぎで、慎重に自分の考えを調整することができるかもしれませんね。

\textbf{Notes (EN)}:\\
"Baai ni yorikeri" is a casual expression that indicates that something varies depending on the situation or conditions. It can be translated as "It depends on the case" or "Depending on the situation." It is often used when an immediate judgment cannot be made easily. By saying "baai ni yorikeri," one can buy time to carefully adjust their thoughts and responses.

\newpage

\section*{427. For example}

\textbf{Date}: 2025.06.10 (Tue)\\
\textbf{Index}: tatoeba\\
\textbf{Expression (JA)}: たとえば\\
\textbf{Romanization}: Tatoeba\\
\textbf{Expression (EN)}: For example

\textbf{Variants (JA)}: たとえば / たとえば、そうだなあ / たとえば...って言われても / たとえば、ね？\\
\textbf{Variants (EN)}: For example / Well, for instance... / Say, for example... / I mean, like...

\textbf{Intent (JA)}: 導入, 例示, あいまい, 話題転換\\
\textbf{Intent (EN)}: introduction, example, ambiguity, topic shift

\textbf{Tags (JA)}: 即時文法, あいまい, 会話, 展開\\
\textbf{Tags (EN)}: immediate grammar, ambiguity, conversation, development

\textbf{Adjusted Expression (JA)}:\\
たとえばの話ですが、こういうふうにも考えられるかもしれません。\\
\textbf{Adjusted Expression (EN)}:\\
For example, one might consider it this way.

\textbf{Example (JA)}:\\
A「(何:なに)か(例:れい)をあげてみて」B「うーん、たとえば、(映画:えいが)とか？」\\
\textbf{Example (EN)}:\\
A: Can you give me an example? B: Hmm, for example, movies?

\textbf{Notes (JA)}:\\
「たとえば」は、何かを説明するときに具体的な例を挙げるためのカジュアルな表現です。「For example」と訳されます。カジュアルな会話でよく使われ、相手に対して具体的なイメージを提供する役割があります。しかし、厳密な表現が頭に思い浮かばないときに使われることもあります。たとえば、何かを説明しようとしているが、具体的な例が思い浮かばないのだから、抽象的な話になっているわけで、そこでなんでもかんでも「たとえば...」と繰り返し質問すると、あなたは嫌われるかもしれませんね。

\textbf{Notes (EN)}:\\
"Tatoeba" is a casual expression used to provide specific examples when explaining something. It can be translated as "For example." This phrase is commonly used in informal conversations to give the listener a concrete image of what is being discussed. However, it is also used when a precise expression does not come to mind. For instance, when trying to explain something but struggling to think of a specific example, the conversation may become abstract. Repeatedly asking "For example..." in such situations might lead to you being disliked.

\newpage

\section*{426. Some days are just like that}

\textbf{Date}: 2025.06.09 (Mon)\\
\textbf{Index}: sonnahi\\
\textbf{Expression (JA)}: そんな(日:ひ)もある\\
\textbf{Romanization}: Sonna hi mo aru\\
\textbf{Expression (EN)}: Some days are just like that

\textbf{Variants (JA)}: そんな日もある / そういうこともある / まあ、あるよね、そういう日 / たまにはそういう日もあるさ\\
\textbf{Variants (EN)}: Some days are like that / These things happen / It happens / Well, that’s just how it goes

\textbf{Intent (JA)}: 受容, 慰め, 共感, ゆるし\\
\textbf{Intent (EN)}: acceptance, comfort, empathy, forgiveness

\textbf{Tags (JA)}: 即時文法, 感情, 日常\\
\textbf{Tags (EN)}: immediate grammar, emotion, daily life

\textbf{Adjusted Expression (JA)}:\\
今日はうまくいかないことが多かったけれど、そういう日もあるものです。\\
\textbf{Adjusted Expression (EN)}:\\
Today didn't go well in many ways, but there are days like that.

\textbf{Example (JA)}:\\
「あー、なんか(全然:ぜんぜん)だめだった...」「うん、そんな(日:ひ)もあるよ」\\
\textbf{Example (EN)}:\\
Ugh, nothing went right today... Yeah, some days are like that.

\textbf{Notes (JA)}:\\
「そんな日もある」は、思いどおりにならない状況を受け入れる姿勢を表すカジュアルな表現である。「そんな日ってどんな日？」と聞かれると実際には困る。言語を使うと何でも具体的に言えるわけではない。具体的に言えない場合には、「そんな」「こんな」「こういう」という実際に見ている、体験していることを差してコミュニケーションをすすめることも必要である。

\textbf{Notes (EN)}:\\
"Sonna hi mo aru" is a casual expression that conveys an attitude of acceptance towards situations that do not go as planned. When asked, "What kind of day is that?" it can be difficult to provide a specific answer. Language does not always allow for concrete descriptions, and in cases where specifics cannot be provided, using terms like "sonna" (that kind of) or "konna" (this kind of) based on what one has seen or experienced can be necessary for effective communication.

\newpage

\section*{424. That being said / Putting that aside}

\textbf{Date}: 2025.06.08 (Sun)\\
\textbf{Index}: sorehasoretoshite\\
\textbf{Expression (JA)}: それはそれとして\\
\textbf{Romanization}: Sore wa sore to shite\\
\textbf{Expression (EN)}: That being said / Putting that aside

\textbf{Variants (JA)}: それはそれとして / ともかくとして / それはそれで / それは置いといて\\
\textbf{Variants (EN)}: That being said / Be that as it may / Anyway / Putting that aside

\textbf{Intent (JA)}: 話題転換, 譲歩, 切り分け\\
\textbf{Intent (EN)}: topic shift, concession, separation

\textbf{Tags (JA)}: 即時文法, 構文省略, スタンス\\
\textbf{Tags (EN)}: immediate grammar, ellipsis, stance

\textbf{Adjusted Expression (JA)}:\\
その(件:けん)はさておき、(次:つぎ)の(問題:もんだい)に(入:はい)りましょう。\\
\textbf{Adjusted Expression (EN)}:\\
Setting that matter aside, let us move on to the next issue.

\textbf{Example (JA)}:\\
「まあ、それはそれとして、(今日:きょう)の(予定:よてい)どうする？」\\
\textbf{Example (EN)}:\\
Anyway, putting that aside, what's the plan for today?

\textbf{Notes (JA)}:\\
「それはそれとして」は、ある話題を一旦脇に置いて別の話題に移るときに使われるカジュアルな表現です。「That being said」や「Putting that aside」と訳されます。カジュアルな会話でよく使われ、相手に対して話題を切り替える役割があります。しかし、自分の都合が悪くなるとこの表現を使いたがる人もいますので、注意が必要です。特に、自分の不都合な話題を摩り替えたいのでしょう。

\textbf{Notes (EN)}:\\
"Sore wa sore to shite" is a casual expression used to set aside a topic and shift to another one. It can be translated as "That being said" or "Putting that aside." This phrase is commonly used in informal conversations to indicate a change of topic. However, some people tend to use this expression when they want to avoid discussing something inconvenient for them, so caution is advised. It is particularly used when they want to divert attention from an unfavorable topic.

\newpage

\section*{423. Not everything / Not all of it}

\textbf{Date}: 2025.06.07 (Sat)\\
\textbf{Index}: zenbugazenbu\\
\textbf{Expression (JA)}: ぜんぶがぜんぶ\\
\textbf{Romanization}: Zenbu ga zenbu\\
\textbf{Expression (EN)}: Not everything / Not all of it

\textbf{Variants (JA)}: ぜんぶがぜんぶ / すべてがすべて / みんながみんな / 何もかもがそうってわけじゃない\\
\textbf{Variants (EN)}: Not everything / Not all of it / It's not always the case / Not everyone is like that

\textbf{Intent (JA)}: 部分否定, 慎重, バランス, 説明の前置き\\
\textbf{Intent (EN)}: partial negation, caution, balance, hedging

\textbf{Tags (JA)}: 即時文法, スタンス, 慎重さ\\
\textbf{Tags (EN)}: immediate grammar, stance, cautiousness

\textbf{Adjusted Expression (JA)}:\\
すべてが必ずしも悪いとは言い切れません。\\
\textbf{Adjusted Expression (EN)}:\\
Not everything can necessarily be said to be bad.

\textbf{Example (JA)}:\\
A「(人:ひと)は(みんな:みんな)そうだよね？」B「いや、ぜんぶがぜんぶそうってわけじゃないよ」\\
\textbf{Example (EN)}:\\
A: People are all like that, right? B: No, not everyone is like that.

\textbf{Notes (JA)}:\\
「ぜんぶがぜんぶ」は、すべての事例や人が同じではないことを強調するカジュアルな表現です。「Not everything」や「Not all of it」と訳されます。カジュアルな会話でよく使われ、相手に対して慎重な態度を示す役割があります。特に、一般化や偏見を避けるために使われることが多いです。

\textbf{Notes (EN)}:\\
"Zenbu ga zenbu" is a casual expression used to emphasize that not all cases or people are the same. It can be translated as "Not everything" or "Not all of it." This phrase is commonly used in informal conversations to indicate a cautious attitude towards generalizations or stereotypes, often serving to avoid bias.

\newpage

\section*{422. It’ll work out somehow / It'll be fine}

\textbf{Date}: 2025.06.06 (Fri)\\
\textbf{Index}: dounika.naru\\
\textbf{Expression (JA)}: どうにかなる\\
\textbf{Romanization}: Dounika naru\\
\textbf{Expression (EN)}: It'll work out somehow / It'll be fine

\textbf{Variants (JA)}: どうにかなる / なんとかなる / まあ何とかなるさ / きっと大丈夫\\
\textbf{Variants (EN)}: It’ll work out somehow / It’ll be fine / Things will turn out okay / It’ll be all right

\textbf{Intent (JA)}: 楽観, 安心, 慰め, 励まし\\
\textbf{Intent (EN)}: optimism, reassurance, comfort, encouragement

\textbf{Tags (JA)}: 即時文法, 感情, 状況描写, 慰め\\
\textbf{Tags (EN)}: immediate grammar, emotion, situation description, comfort

\textbf{Adjusted Expression (JA)}:\\
状況は厳しいかもしれませんが、最終的には何とか解決できると思います。\\
\textbf{Adjusted Expression (EN)}:\\
Although the situation may be difficult, I believe that things will eventually be resolved.

\textbf{Example (JA)}:\\
A「(締切:しめきり)に(間:ま)に(合:あ)うかな…」B「(大丈夫:だいじょうぶ)、どうにかなるよ。」\\
\textbf{Example (EN)}:\\
A: Do you think we'll make the deadline...? B: Don't worry, it'll be fine.

\textbf{Notes (JA)}:\\
「どうにかなる」は、状況が厳しいときでも、最終的には何とかなるという楽観的な気持ちを表すカジュアルな表現です。「It'll work out somehow」や「It'll be fine」と訳されます。カジュアルな会話でよく使われ、相手に対して安心感を与える役割があるかもしれませんが、根拠もないのに、「どうにかなる」というと無責任に思われるかもしれません。私の友人に口癖でこういう人がいますが、だいたい「どうにもならない」ことのようです。

\textbf{Notes (EN)}:\\
"Dounika naru" is a casual expression that conveys an optimistic feeling that things will work out in the end, even in difficult situations. It can be translated as "It'll work out somehow" or "It'll be fine." This phrase is commonly used in informal conversations to provide reassurance to the listener. However, saying "Dounika naru" without any basis may come across as irresponsible. I have a friend who often uses this phrase, but it usually turns out that things do not work out as expected.

\newpage

\section*{421. whatever ... is / no matter what ... is}

\textbf{Date}: 2025.06.05 (Thu)\\
\textbf{Index}: nan.de.ari\\
\textbf{Expression (JA)}: ...は(何:なん)であれ\\
\textbf{Romanization}: ... wa nan de are\\
\textbf{Expression (EN)}: whatever ... is / no matter what ... is

\textbf{Variants (JA)}: 何であれ / どんなものであれ / 何であっても / いずれにせよ\\
\textbf{Variants (EN)}: whatever it is / no matter what / in any case / regardless of what it is

\textbf{Intent (JA)}: 条件不問, 一般化, 包括, 話題転換\\
\textbf{Intent (EN)}: regardless of condition, generalization, inclusiveness, topic shift

\textbf{Tags (JA)}: 即時文法, 条件, 一般化, 文脈転換\\
\textbf{Tags (EN)}: immediate grammar, condition, generalization, context shift

\textbf{Adjusted Expression (JA)}:\\
いかなる(理由:りゆう)があったとして、...。\\
\textbf{Adjusted Expression (EN)}:\\
Regardless of the reason, ...

\textbf{Example (JA)}:\\
A「(結果:けっか)が(何:なん)であれ、もう(終:お)わったことだ」\\
\textbf{Example (EN)}:\\
A: Whatever the result is, it's already over.

\textbf{Notes (JA)}:\\
「...は何であれ」は、条件や状況に関係なく、何かを包括的に述べる表現です。「whatever ... is」や「no matter what ... is」と訳されます。「あれ」は、めずらしい形で、こういう定型表現のときに使われますが、単独ではあまり使われません。

\textbf{Notes (EN)}:\\
"... wa nan de are" is an expression that describes something inclusively, regardless of conditions or situations. It can be translated as "whatever ... is" or "no matter what ... is." The term "are" is used in this fixed expression but is not commonly used on its own.

\newpage

\section*{420. feel relieved}

\textbf{Date}: 2025.06.04 (Wed)\\
\textbf{Index}: hotto.suru\\
\textbf{Expression (JA)}: ほっとする\\
\textbf{Romanization}: Hotto suru\\
\textbf{Expression (EN)}: feel relieved

\textbf{Variants (JA)}: ほっとした / ほっとする / ほっとして / ほっとしている\\
\textbf{Variants (EN)}: felt relieved / feel relieved / what a relief / that's a relief

\textbf{Intent (JA)}: 安堵, 緊張の解放, 安心, 気持ちの変化\\
\textbf{Intent (EN)}: relief, release of tension, reassurance, change of feeling

\textbf{Tags (JA)}: 即時文法, 感情, 心情, 状況描写\\
\textbf{Tags (EN)}: immediate grammar, emotion, feeling, situation description

\textbf{Adjusted Expression (JA)}:\\
やっと安心できました。\\
\textbf{Adjusted Expression (EN)}:\\
I can finally relax.

\textbf{Example (JA)}:\\
A「(試験:しけん)が(終:お)わった。ほっとした」B「うん、(結果:けっか)が(良:よ)かったらいいね」\\
\textbf{Example (EN)}:\\
A: The exam is over. I feel relieved. B: Yeah, I hope the results are good.

\textbf{Notes (JA)}:\\
「ほっとする」は、緊張や不安が解消されて安心することを表すカジュアルな表現です。「feel relieved」や「what a relief」と訳されます。カジュアルな会話でよく使われ、相手に対して安心感を伝える役割があります。

\textbf{Notes (EN)}:\\
"Hotto suru" is a casual expression used to convey the feeling of relief when tension or anxiety is released. It can be translated as "feel relieved" or "what a relief." 

\newpage

\section*{419. Let's get started}

\textbf{Date}: 2025.06.03 (Tue)\\
\textbf{Index}: hajimeyou\\
\textbf{Expression (JA)}: さあ、(始:はじ)めよう\\
\textbf{Romanization}: Saa, hajimeyou\\
\textbf{Expression (EN)}: Alright, let’s get started

\textbf{Variants (JA)}: さあ、始めよう / さあ、始めましょう / さあ、いこう\\
\textbf{Variants (EN)}: Alright, let's begin / Okay, let's get started / Let's go

\textbf{Intent (JA)}: 開始, 決意, 勢い\\
\textbf{Intent (EN)}: start, determination, momentum

\textbf{Tags (JA)}: 即時文法, 行動, 開始\\
\textbf{Tags (EN)}: immediate grammar, action, start

\textbf{Adjusted Expression (JA)}:\\
さて、始めましょうか。\\
\textbf{Adjusted Expression (EN)}:\\
Well then, shall we begin?

\textbf{Example (JA)}:\\
A「さあ、(始:はじ)めよう」B「うん、でも、(何:なに)を？」\\
\textbf{Example (EN)}:\\
A: Alright, let's get started. B: Sure, but what should we do?

\textbf{Notes (JA)}:\\
「さあ、始めよう」は、何かを始めるときの決意や勢いを表すカジュアルな表現です。「Alright, let's get started」や「Okay, let's begin」と訳されます。カジュアルな会話でよく使われ、相手に対して行動を促す役割があります。

\textbf{Notes (EN)}:\\
"Saa, hajimeyou" is a casual expression used to convey determination or momentum when starting something. It can be translated as "Alright, let's get started" or "Okay, let's begin."

\newpage

\section*{418. repeatedly / over and over}

\textbf{Date}: 2025.06.02 (Mon)\\
\textbf{Index}: kurikaeshikurikaeshi\\
\textbf{Expression (JA)}: (繰:く)り(返:かえ)し(繰:く)り(返:かえ)し\\
\textbf{Romanization}: Kurikaeshi kurikaeshi\\
\textbf{Expression (EN)}: repeatedly / over and over

\textbf{Variants (JA)}: 繰り返し繰り返し / 繰り返し / 何度も何度も / 何度も\\
\textbf{Variants (EN)}: repeatedly / over and over / time and again / again and again

\textbf{Intent (JA)}: 反復, 持続, 強調\\
\textbf{Intent (EN)}: repetition, continuity, emphasis

\textbf{Tags (JA)}: 即時文法, 反復, 状況描写, 強調\\
\textbf{Tags (EN)}: immediate grammar, repetition, situation description, emphasis

\textbf{Adjusted Expression (JA)}:\\
何度も繰り返して、やっと理解できた。\\
\textbf{Adjusted Expression (EN)}:\\
After repeating it over and over, I finally understood.

\textbf{Example (JA)}:\\
A「(同:おな)じことを(繰:く)り(返:かえ)し(言:い)ってるね」B「うん、(何度:なんど)も(繰:く)り(返:かえ)し(言:い)うのが(大事:だいじ)なんだよ」\\
\textbf{Example (EN)}:\\
A: You keep saying the same thing. B: Yeah, it's important to repeat it over and over.

\textbf{Notes (JA)}:\\
「繰り返し繰り返し」は、何かを何度も行うことを強調するカジュアルな表現です。「repeatedly」や「over and over」と訳されます。カジュアルな会話でよく使われ、相手に対して反復的な行動や状況を伝える役割があります。

\textbf{Notes (EN)}:\\
"Kurikaeshi kurikaeshi" is a casual expression used to emphasize doing something repeatedly. It can be translated as "repeatedly" or "over and over." This phrase is commonly used in informal conversations to convey repetitive actions or situations to the listener. 

\newpage

\section*{417. nonstop / without pause}

\textbf{Date}: 2025.06.01 (Sun)\\
\textbf{Index}: hikkirinashini\\
\textbf{Expression (JA)}: ひっきりなしに\\
\textbf{Romanization}: Hikkiri nashi ni\\
\textbf{Expression (EN)}: nonstop / without pause

\textbf{Variants (JA)}: ひっきりなしに / ひっきりなしで / ひっきりなしだ\\
\textbf{Variants (EN)}: nonstop / without pause / constantly / incessantly

\textbf{Intent (JA)}: 連続, 途切れない, 多忙\\
\textbf{Intent (EN)}: continuity, uninterrupted, busyness

\textbf{Tags (JA)}: 即時文法, 状況描写, 頻度\\
\textbf{Tags (EN)}: immediate grammar, situation description, frequency

\textbf{Adjusted Expression (JA)}:\\
休む間もなく次の仕事が始まった。\\
\textbf{Adjusted Expression (EN)}:\\
Without pause, the next task began.

\textbf{Example (JA)}:\\
ひっきりなしに(来客:らいきゃく)があって(休:やす)む(暇:ひま)がない。\\
\textbf{Example (EN)}:\\
Visitors keep coming nonstop, leaving no time to rest.

\textbf{Notes (JA)}:\\
「ひっきりなしに」は、何かが途切れることなく続く様子を表すカジュアルな表現です。「nonstop」や「without pause」と訳されます。カジュアルな会話でよく使われ、相手に対して連続的な状況を伝える役割があります。「ひっきり」ということばは「引き切る」から来たことばといわれていますが、「ひっきり」だけで使われることはありません。もちろん、「ひっきりありに」ということばもありませんよ。

\textbf{Notes (EN)}:\\
"Hikkiri nashi ni" is a casual expression used to describe something that continues without interruption. It can be translated as "nonstop" or "without pause." This phrase is commonly used in informal conversations to convey a continuous situation to the listener. The term "hikkiri" is believed to originate from "hiki kiru," but it is not used on its own, and there is no phrase like "hikkiri ari ni."

\newpage

\section*{416. That's fine as it is, isn't it?}

\textbf{Date}: 2025.05.31 (Sat)\\
\textbf{Index}: sorehasorede\\
\textbf{Expression (JA)}: それはそれで...いいんじゃない？\\
\textbf{Romanization}: Sore wa sore de... iin janai?\\
\textbf{Expression (EN)}: That's fine as it is, isn't it?

\textbf{Variants (JA)}: それはそれでいいんじゃない？ / それはそれでアリじゃない？ / まあ、それはそれでいいか\\
\textbf{Variants (EN)}: That's fine as it is, isn't it? / That's acceptable, right? / Well, that’s okay, I guess

\textbf{Intent (JA)}: 受容, 妥協, 緩和\\
\textbf{Intent (EN)}: acceptance, compromise, softening

\textbf{Tags (JA)}: 即時文法, 受容, 妥協, ゆるし\\
\textbf{Tags (EN)}: immediate grammar, acceptance, compromise, forgiveness

\textbf{Adjusted Expression (JA)}:\\
それはそれで、まあいいんじゃないでしょうか。\\
\textbf{Adjusted Expression (EN)}:\\
That's fine as it is, I suppose.

\textbf{Example (JA)}:\\
A「(失敗作:しっぱいさく)の(ケーキ:けーき)、(見:み)た(目:め)はぐちゃぐちゃだよ」 B「でも、それはそれで、(手作:てづく)り(感:かん)あっていいんじゃない？」\\
\textbf{Example (EN)}:\\
A: 'This failed cake looks like a total mess...' B: 'But, you know, that’s fine as it is ,  it has a homemade charm!'

\textbf{Notes (JA)}:\\
「それはそれで...いいんじゃない？」は、相手の意見や状況を受け入れつつ、少しの妥協や緩和を示すカジュアルな表現です。「That's fine as it is, isn't it?」と訳されます。カジュアルな会話でよく使われ、相手に対して受容的な態度を示す役割があります。通常の基準ではマイナス評価される要素も、視点を変えることで肯定される場面で使われます。

\textbf{Notes (EN)}:\\
"Sore wa sore de... iin janai?" is a casual expression used to accept the other person's opinion or situation while showing a bit of compromise or softening. It can be translated as "That's fine as it is, isn't it?" This phrase is commonly used in informal conversations to indicate an accepting attitude towards the listener. It often appears in situations where something normally seen as negative is reframed in a positive light.

\newpage

\section*{415. Not at all / Not in the least}

\textbf{Date}: 2025.05.30 (Fri)\\
\textbf{Index}: sappari\\
\textbf{Expression (JA)}: さっぱり\\
\textbf{Romanization}: Sappari\\
\textbf{Expression (EN)}: Not at all / Not in the least

\textbf{Variants (JA)}: さっぱり / さっぱりわからない / さっぱりしない\\
\textbf{Variants (EN)}: Not at all / Not in the least / I don't understand at all

\textbf{Intent (JA)}: 否定, 不満, 軽い驚き\\
\textbf{Intent (EN)}: negation, dissatisfaction, mild surprise

\textbf{Tags (JA)}: 即時文法, 否定, 状況依存型, 軽い驚き\\
\textbf{Tags (EN)}: immediate grammar, negation, context-dependent, mild surprise

\textbf{Adjusted Expression (JA)}:\\
まったくうまくいきません。\\
\textbf{Adjusted Expression (EN)}:\\
It doesn't go well at all.

\textbf{Example (JA)}:\\
A「(今日:きょう)の(成果:せいか)は？」B「さっぱり」\\
\textbf{Example (EN)}:\\
A: How was today's result? B: Not at all.

\textbf{Notes (JA)}:\\
「さっぱり」は、何かが全くうまくいかない、または理解できないことを表すカジュアルな表現です。「Not at all」や「Not in the least」と訳されます。カジュアルな会話でよく使われ、相手に対して否定的な感情を伝える役割があります。

\textbf{Notes (EN)}:\\
"Sappari" is a casual expression used to indicate that something is not going well at all or that one does not understand something. It can be translated as "Not at all" or "Not in the least." .

\newpage

\section*{414. this and that / such and such}

\textbf{Date}: 2025.05.29 (Thu)\\
\textbf{Index}: dounokouno\\
\textbf{Expression (JA)}: どうのこうの\\
\textbf{Romanization}: Dō no kō no\\
\textbf{Expression (EN)}: this and that / such and such

\textbf{Variants (JA)}: どうのこうの / どうこう / ああだこうだ / あれこれ\\
\textbf{Variants (EN)}: this and that / such and such / blah blah / this or that

\textbf{Intent (JA)}: 文句, 細かいこと, 批判, 不明な事情の濁し\\
\textbf{Intent (EN)}: complaint, nitpicking, criticism, glossing over unclear details

\textbf{Tags (JA)}: 即時文法, 状況依存, 繰り返し, 軽い文句, 事情の省略\\
\textbf{Tags (EN)}: immediate grammar, context-dependent, repetition, light complaint, skipping details

\textbf{Adjusted Expression (JA)}:\\
詳細は存じませんが、税金に関する何らかの話題だったかと記憶しております。\\
\textbf{Adjusted Expression (EN)}:\\
I'm not certain of the details, but I believe it was some matter regarding taxes.

\textbf{Example (JA)}:\\
よくわからないが、(税金:ぜいきん)がどうのこうのという(話:はなし)だったかな？\\
\textbf{Example (EN)}:\\
I don't really know, but I think it was something about taxes and such.

\textbf{Notes (JA)}:\\
「どうのこうの」は、話の細部や事情がはっきりわからないとき、または細かい説明を省略したいときに用いる表現である。例えば「税金がどうのこうのという話だった」という場合、具体的な内容は把握していないが、何かそういう話題があったことを伝えるニュアンスになる。

\textbf{Notes (EN)}:\\
"Dou no kou no" is used when the speaker does not know the details or prefers to skip over detailed explanations. For example, saying "It was something about taxes and such" conveys that the speaker is aware there was some discussion but does not provide or know the specifics.

\newpage

\section*{413. After all / In the end}

\textbf{Date}: 2025.05.28 (Tue)\\
\textbf{Index}: nandakanda de\\
\textbf{Expression (JA)}: なんだかんだで\\
\textbf{Romanization}: Nandakanda de\\
\textbf{Expression (EN)}: After all / In the end

\textbf{Variants (JA)}: なんだかんだで / なんだかんだ言って / なんだかんだの末に\\
\textbf{Variants (EN)}: After all / In the end / When all is said and done

\textbf{Intent (JA)}: 結論, まとめ, 軽いニュアンス\\
\textbf{Intent (EN)}: conclusion, summary, light nuance

\textbf{Tags (JA)}: 即時文法, 話し言葉, カジュアル, 結論\\
\textbf{Tags (EN)}: immediate grammar, spoken language, casual, conclusion

\textbf{Adjusted Expression (JA)}:\\
結局のところ、最終的には\\
\textbf{Adjusted Expression (EN)}:\\
Ultimately, in the end

\textbf{Example (JA)}:\\
なんだかんだで(彼:かれ)が(勝:か)つんだよね。\\
\textbf{Example (EN)}:\\
After all, he ends up winning, you know.

\textbf{Notes (JA)}:\\
「なんだかんだで」は軽い口語表現で、細かいことがあったけれど結局〜だという結論やまとめを表すときに使われる。友人同士の会話などカジュアルな場面でよく耳にする。

\textbf{Notes (EN)}:\\
"Nandakanda de" is a casual spoken expression meaning "after all" or "in the end," used to summarize that despite various details, something ultimately happens. It’s often heard in casual conversations among friends.

\newpage

\section*{412. And then / So then}

\textbf{Date}: 2025.05.27 (Tue)\\
\textbf{Index}: demotte\\
\textbf{Expression (JA)}: でもって\\
\textbf{Romanization}: Demotte\\
\textbf{Expression (EN)}: And then / So then

\textbf{Variants (JA)}: でもって\\
\textbf{Variants (EN)}: And then / So then / Plus / Also

\textbf{Intent (JA)}: 接続, 話のつなぎ, 順接, 補足\\
\textbf{Intent (EN)}: connection, linking, addition, sequence

\textbf{Tags (JA)}: 即時文法, 接続, 話し言葉, カジュアル\\
\textbf{Tags (EN)}: immediate grammar, connection, spoken language, casual

\textbf{Adjusted Expression (JA)}:\\
それで、さらに、\\
\textbf{Adjusted Expression (EN)}:\\
Moreover, additionally, furthermore

\textbf{Example (JA)}:\\
(昨日:きのう)は(雨:あめ)だった、でもって(風:かぜ)も(強:つよ)かった。\\
\textbf{Example (EN)}:\\
Yesterday it rained, and then the wind was strong too.

\textbf{Notes (JA)}:\\
「でもって」は、話の流れをつなぐために使われるカジュアルな表現です。「And then」や「So then」と訳されます。カジュアルな会話でよく使われ、相手に対して話の流れを示す役割があります。この表現は、あまり堅苦しくなく、親しい友人同士の会話でよく使われます。

\textbf{Notes (EN)}:\\
"Demotte" is a casual expression used to connect the flow of conversation. It can be translated as "And then" or "So then." This phrase is commonly used in informal conversations to indicate the continuation of a story or topic. It is not too formal and is often used among close friends.

\newpage

\section*{411. I caught a cold}

\textbf{Date}: 2025.05.26 (Mon)\\
\textbf{Index}: kazehiita\\
\textbf{Expression (JA)}: (風邪:かぜ)ひいた\\
\textbf{Romanization}: Kaze hiita\\
\textbf{Expression (EN)}: I caught a cold

\textbf{Variants (JA)}: 風邪ひいた / 風邪ひいちゃった / 風邪ひいちゃいました / 風邪ひいてしまった\\
\textbf{Variants (EN)}: I caught a cold / I got a cold / I ended up catching a cold / I happened to catch a cold

\textbf{Intent (JA)}: 病気, 軽い後悔, 注意不足\\
\textbf{Intent (EN)}: illness, light regret, lack of attention

\textbf{Tags (JA)}: 即時文法, 省略, 状況依存型, 病気, 軽い後悔, 注意不足\\
\textbf{Tags (EN)}: immediate grammar, ellipsis, context-dependent, illness, light regret, lack of attention

\textbf{Adjusted Expression (JA)}:\\
風邪を引いてしまいました。\\
\textbf{Adjusted Expression (EN)}:\\
I have caught a cold.

\textbf{Example (JA)}:\\
A「(風邪:かぜ)ひいた」B「大丈夫？(薬:くすり)飲んだ？」\\
\textbf{Example (EN)}:\\
A: I caught a cold. B: Are you okay? Did you take medicine?

\textbf{Notes (JA)}:\\
「風邪ひいた」は、病気や体調不良をカジュアルに伝える表現です。「I caught a cold」や「I got a cold」と訳されます。カジュアルな会話でよく使われ、相手に対して体調の悪さを伝える役割があります。この表現は、あまり堅苦しくなく、親しい友人同士の会話でよく使われます。

\textbf{Notes (EN)}:\\
"Kaze hiita" is a casual expression used to convey illness or poor health. It can be translated as "I caught a cold" or "I got a cold." This phrase is commonly used in informal conversations to inform the listener about feeling unwell. It is not too formal and is often used among close friends.

\newpage

\section*{410. What are you doing?}

\textbf{Date}: 2025.05.25 (Sun)\\
\textbf{Index}: naniyattenno\\
\textbf{Expression (JA)}: (何:なに)やってんの？\\
\textbf{Romanization}: Nani yatten no?\\
\textbf{Expression (EN)}: What are you doing?

\textbf{Variants (JA)}: 何やってんの？ / 何やってるの？ / 何してんの？ / 何してるの？\\
\textbf{Variants (EN)}: What are you doing? / What are you up to? / What are you working on? / What are you engaged in?

\textbf{Intent (JA)}: 疑問, 関心, 親密さ\\
\textbf{Intent (EN)}: question, interest, intimacy

\textbf{Tags (JA)}: 即時文法, 省略, 状況依存型, 疑問, 関心, 親密さ\\
\textbf{Tags (EN)}: immediate grammar, ellipsis, context-dependent, question, interest, intimacy

\textbf{Adjusted Expression (JA)}:\\
今、何をしているのですか？\\
\textbf{Adjusted Expression (EN)}:\\
What are you doing right now?

\textbf{Example (JA)}:\\
A「(何:なに)やってんの？」B「(勉強:べんきょう)してるよ」\\
\textbf{Example (EN)}:\\
A: What are you doing? B: I'm studying.

\textbf{Notes (JA)}:\\
「何やってんの？」は、相手の行動や状況に対する疑問をカジュアルに尋ねる表現です。「What are you doing?」や「What are you up to?」と訳されます。カジュアルな会話でよく使われ、相手に対して関心を示す役割があります。

\textbf{Notes (EN)}:\\
"Nani yatten no?" is a casual expression used to ask about someone's actions or situation. It can be translated as "What are you doing?" or "What are you up to?" This phrase is commonly used in informal conversations to show interest in the listener's activities.

\newpage

\section*{409. Delicious}

\textbf{Date}: 2025.05.24 (Sat)\\
\textbf{Index}: uma\\
\textbf{Expression (JA)}: うま！\\
\textbf{Romanization}: Uma!\\
\textbf{Expression (EN)}: Delicious

\textbf{Variants (JA)}: うま！ / うまい！ / うっま！ / うまいよ！ / めっちゃうま！\\
\textbf{Variants (EN)}: Delicious! / Tasty! / Yummy! / So good!

\textbf{Intent (JA)}: 感情の表出, 喜び, 驚き, 強調, 若者言葉\\
\textbf{Intent (EN)}: expression of emotion, joy, surprise, emphasis, youth slang

\textbf{Tags (JA)}: 即時文法, 感情表出, 状況依存型\\
\textbf{Tags (EN)}: immediate grammar, emotion expression, context-dependent

\textbf{Adjusted Expression (JA)}:\\
おいしい！\\
\textbf{Adjusted Expression (EN)}:\\
It's delicious!

\textbf{Example (JA)}:\\
A「(食:た)べてみて」B「うま！」\\
\textbf{Example (EN)}:\\
A: Try it! B: Delicious!

\textbf{Notes (JA)}:\\
「うま！」は、食べ物が美味しいときに使うカジュアルな表現です。「Delicious!」や「Tasty!」と訳されます。カジュアルな会話でよく使われ、相手に対して感情を表現する役割があります。この表現は、あまり堅苦しくなく、親しい友人同士の会話でよく使われます。反射的に感情が口に出てしまう場面には、イ形容詞の「うまい」の「い」を発音せずに「うま！」と短縮して使うことが多いです。

\textbf{Notes (EN)}:\\
"Uma!" is a casual expression used to indicate that food is delicious. It can be translated as "Delicious!" or "Tasty!" This phrase is commonly used in informal conversations to express emotions to the listener. It is not too formal and is often used among close friends. In spontaneous emotional reactions, the adjective 'umai' is often shortened to 'uma!' without pronouncing the 'i'.

\newpage

\section*{408. accidentally}

\textbf{Date}: 2025.05.23 (Fri)\\
\textbf{Index}: tsui.*teshimatta\\
\textbf{Expression (JA)}: つい、...してしまった\\
\textbf{Romanization}: Tsui ... shite shimatta\\
\textbf{Expression (EN)}: accidentally ...

\textbf{Variants (JA)}: つい、...してしまった / つい、...しちゃった / つい、...しちゃいました / つい、...してしまいました\\
\textbf{Variants (EN)}: accidentally ... / I accidentally ... / I ended up ... / I happened to ...

\textbf{Intent (JA)}: 失敗, 軽い後悔, 注意不足\\
\textbf{Intent (EN)}: mistake, light regret, lack of attention

\textbf{Tags (JA)}: 即時文法, 省略, 状況依存型, 失敗, 軽い後悔, 注意不足\\
\textbf{Tags (EN)}: immediate grammar, ellipsis, context-dependent, mistake, light regret, lack of attention

\textbf{Adjusted Expression (JA)}:\\
(意図:いと)せずに、...してしまいました。\\
\textbf{Adjusted Expression (EN)}:\\
I have done ... unintentionally.

\textbf{Example (JA)}:\\
A「おいしいんで、つい、いっぱい(食:た)べちゃった」\\
\textbf{Example (EN)}:\\
A: It was so delicious, I ended up eating a lot.

\textbf{Notes (JA)}:\\
「つい、...してしまった」は、注意不足や軽い後悔を表すカジュアルな表現です。「accidentally ...」や「I accidentally ...」と訳されます。カジュアルな会話でよく使われ、相手に対して失敗を伝える役割があります。

\textbf{Notes (EN)}:\\
"Tsui ... shite shimatta" is a casual expression used to convey a lack of attention or light regret. It can be translated as "accidentally ..." or "I ended up ..."

\newpage

\section*{407. Oops}

\textbf{Date}: 2025.05.22 (Thu)\\
\textbf{Index}: ukkari\\
\textbf{Expression (JA)}: うっかり\\
\textbf{Romanization}: Ukkari\\
\textbf{Expression (EN)}: Oops

\textbf{Variants (JA)}: うっかり / うっかりして / うっかりしてた / うっかりしちゃった\\
\textbf{Variants (EN)}: Oops / I accidentally / I carelessly / I forgot

\textbf{Intent (JA)}: 失敗, 軽い後悔, 注意不足\\
\textbf{Intent (EN)}: mistake, light regret, lack of attention

\textbf{Tags (JA)}: 即時文法, 省略, 状況依存型, 失敗, 軽い後悔, 注意不足\\
\textbf{Tags (EN)}: immediate grammar, ellipsis, context-dependent, mistake, light regret, lack of attention

\textbf{Adjusted Expression (JA)}:\\
注意を怠ってしまいました。\\
\textbf{Adjusted Expression (EN)}:\\
I neglected to pay attention.

\textbf{Example (JA)}:\\
A「(うっかり)して、(傘:かさ)、(忘:わす)れちゃった」B「それ、(安物:やすもの)でしょ？」\\
\textbf{Example (EN)}:\\
A: I accidentally forgot my umbrella. B: That was a cheap one, right?

\textbf{Notes (JA)}:\\
「うっかり」は、注意不足や軽い後悔を表すカジュアルな表現です。「Oops」や「I accidentally」と訳されます。カジュアルな会話でよく使われ、相手に対して失敗を伝える役割があります。この表現は、あまり堅苦しくなく、親しい友人同士の会話でよく使われます。

\textbf{Notes (EN)}:\\
"Ukkari" is a casual expression used to convey a lack of attention or light regret. It can be translated as "Oops" or "I accidentally." This phrase is commonly used in informal conversations to inform the listener about a mistake. It is not too formal and is often used among close friends.

\newpage

\section*{406. What should I do?}

\textbf{Date}: 2025.05.21 (Wed)\\
\textbf{Index}: dousuru?\\
\textbf{Expression (JA)}: どうする？\\
\textbf{Romanization}: Dousuru?\\
\textbf{Expression (EN)}: What should I do?

\textbf{Variants (JA)}: どうする？ / どうしよう？ / どうするの？ / どうしようかな？\\
\textbf{Variants (EN)}: What should I do? / What will I do? / What should I do about it? / What should I think?

\textbf{Intent (JA)}: 選択, 不安, 迷い\\
\textbf{Intent (EN)}: choice, anxiety, hesitation

\textbf{Tags (JA)}: 即時文法, 省略, 状況依存型, 選択, 不安, 迷い\\
\textbf{Tags (EN)}: immediate grammar, ellipsis, context-dependent, choice, anxiety, hesitation

\textbf{Adjusted Expression (JA)}:\\
どうしようか考えています。\\
\textbf{Adjusted Expression (EN)}:\\
I'm thinking about what to do.

\textbf{Example (JA)}:\\
A「(明日:あした)、(雨:あめ)だって。どうする？」B「(雨:あめ)でも、行かなくちゃ」\\
\textbf{Example (EN)}:\\
A: Tomorrow, they say it's going to rain. What should we do? B: Even if it rains, we have to go.

\textbf{Notes (JA)}:\\
「どうする？」は、選択や決断を求めるカジュアルな表現です。「What should I do?」や「What will I do?」と訳されます。カジュアルな会話でよく使われ、相手に対して選択肢を提示する役割があります。この表現は、あまり堅苦しくなく、親しい友人同士の会話でよく使われます。

\textbf{Notes (EN)}:\\
"Dousuru?" is a casual expression used to ask for choices or decisions. It can be translated as "What should I do?" or "What will I do?" This phrase is commonly used in informal conversations to present options to the listener. It is not too formal and is often used among close friends.

\newpage

\section*{405. It's delicious, but...}

\textbf{Date}: 2025.05.20 (Tue)\\
\textbf{Index}: oishiinoni\\
\textbf{Expression (JA)}: おいしいのに...\\
\textbf{Romanization}: Oishii noni...\\
\textbf{Expression (EN)}: It's delicious, but...

\textbf{Variants (JA)}: おいしいのに / おいしいのになあ / おいしいのにね / おいしかったのに\\
\textbf{Variants (EN)}: It's delicious, but... / It was so good, but... / Even though it's tasty... / Even though I liked it...

\textbf{Intent (JA)}: 感情の矛盾, 切なさ, 共有, 同情\\
\textbf{Intent (EN)}: emotional conflict, bittersweet, shared feeling, sympathy

\textbf{Tags (JA)}: 即時文法, 省略, 感情表出, 状況依存型\\
\textbf{Tags (EN)}: immediate grammar, ellipsis, emotion, context-dependent

\textbf{Adjusted Expression (JA)}:\\
おいしいのに、誰も食べてくれません。\\
\textbf{Adjusted Expression (EN)}:\\
Although it's delicious, nobody eats it.

\textbf{Example (JA)}:\\
A「(納豆:なっとう)、きらい」B「え、おいしいのに...」\\
\textbf{Example (EN)}:\\
A: Natto(fermented soybeans) is gross. B: Huh, but it's delicious...

\textbf{Notes (JA)}:\\
「おいしいのに」は、感情の矛盾や切なさを表現するカジュアルな表現です。「It's delicious, but...」や「It was so good, but...」と訳されます。カジュアルな会話でよく使われ、相手に対して感情を共有する役割があります。「けれども」と違うのは、状況的に相手の意見に対して「そうではない」と否定的な意味合いを持つことです。

\textbf{Notes (EN)}:\\
"Oishii noni" is a casual expression used to convey emotional conflict or bittersweet feelings. It can be translated as "It's delicious, but..." or "It was so good, but..." This phrase is commonly used in informal conversations to share emotions with the listener. Unlike the word "keredomo," it carries a more negative connotation, suggesting that the speaker disagrees with the listener's opinion.

\newpage

\section*{404. I might be able to}

\textbf{Date}: 2025.05.19 (Mon)\\
\textbf{Index}: dekirukamo\\
\textbf{Expression (JA)}: できるかも\\
\textbf{Romanization}: Dekiru kamo\\
\textbf{Expression (EN)}: I might be able to

\textbf{Variants (JA)}: できるかも / できるかもしれない / できるかもね / できるかもよ\\
\textbf{Variants (EN)}: I might be able to / Maybe I can / I guess I could / I might manage

\textbf{Intent (JA)}: 可能性, 控えめ, 柔らかさ\\
\textbf{Intent (EN)}: possibility, tentativeness, softening

\textbf{Tags (JA)}: 即時文法, 推量, 婉曲, 控えめ, 柔らかさ\\
\textbf{Tags (EN)}: immediate grammar, guessing, euphemism, tentativeness, softening

\textbf{Adjusted Expression (JA)}:\\
もしかしたら、できるかもしれません。\\
\textbf{Adjusted Expression (EN)}:\\
Perhaps I might be able to do it.

\textbf{Example (JA)}:\\
A「この課題、終わりそう？」 B「うん、できるかも。」\\
\textbf{Example (EN)}:\\
A: “Do you think you can finish the assignment?” B: “Yeah, I might be able to.”

\textbf{Notes (JA)}:\\
「できるかも」は、即時的な思いつきや予測に基づいて、自信はないが可能性はあるという態度を示す表現。丁寧に言うと「できるかもしれません」になるが、即時文法では判断の途中で止まっている状態をそのまま発話する。

\textbf{Notes (EN)}:\\
'Dekiru kamo' expresses a tentative possibility without strong confidence. It reflects the speaker’s real-time hesitation or ongoing evaluation, unlike the more formal ‘dekiru kamo shiremasen,’ which is often used in adjusted speech.

\newpage

\section*{403. As expected}

\textbf{Date}: 2025.05.18 (Sun)\\
\textbf{Index}: yappari\\
\textbf{Expression (JA)}: やっぱり\\
\textbf{Romanization}: Yappari\\
\textbf{Expression (EN)}: As expected

\textbf{Variants (JA)}: やっぱり / やっぱ / やっぱね / やっぱさ\\
\textbf{Variants (EN)}: As expected / Just as I thought

\textbf{Intent (JA)}: 確信, 同意, 親密\\
\textbf{Intent (EN)}: certainty, agreement, intimacy

\textbf{Tags (JA)}: 即時文法, 確信, 同意, 親密\\
\textbf{Tags (EN)}: immediate grammar, certainty, agreement, intimacy

\textbf{Adjusted Expression (JA)}:\\
やっぱりそうだと思います。\\
\textbf{Adjusted Expression (EN)}:\\
I think that's true.

\textbf{Example (JA)}:\\
A「(明日:あした)は(晴:は)れるかな？」B「やっぱり(晴:は)れないよ」\\
\textbf{Example (EN)}:\\
A: Will it be sunny tomorrow? B: As expected, it won't be.

\textbf{Notes (JA)}:\\
「やっぱり」は、相手の言葉に対して軽く同意するカジュアルな表現です。「As expected」や「Just as I thought」と訳されます。カジュアルな会話でよく使われ、相手に対して軽い同意を示す役割があります。この表現は、あまり本気ではないが、相手の言葉を受け入れるときに使われます。

\textbf{Notes (EN)}:\\
"Yappari" is a casual expression used to lightly agree with someone. It can be translated as "As expected" or "Just as I thought." This phrase is commonly used in informal conversations to show a light agreement with the speaker. It is often used when you don't take the statement too seriously but still want to acknowledge it.

\newpage

\section*{402. I have to eat}

\textbf{Date}: 2025.05.17 (Sat)\\
\textbf{Index}: tabenakucha\\
\textbf{Expression (JA)}: (食:た)べなくちゃ\\
\textbf{Romanization}: Tabenakucha\\
\textbf{Expression (EN)}: I have to eat

\textbf{Variants (JA)}: 食べなくちゃ / 食べなきゃ / 食べなきゃいけない / 食べなきゃならない\\
\textbf{Variants (EN)}: I have to eat / I must eat / I need to eat / I gotta eat

\textbf{Intent (JA)}: 義務, 必要, 意志\\
\textbf{Intent (EN)}: obligation, necessity, intention

\textbf{Tags (JA)}: 即時文法, 省略, 状況依存型, 義務, 必要, 意志\\
\textbf{Tags (EN)}: immediate grammar, ellipsis, context-dependent, obligation, necessity, intention

\textbf{Adjusted Expression (JA)}:\\
食べなければなりません。\\
\textbf{Adjusted Expression (EN)}:\\
I have to eat.

\textbf{Example (JA)}:\\
A「(食:た)べなくちゃ」B「(食:た)べなきゃいけないよ」\\
\textbf{Example (EN)}:\\
A: I have to eat. B: You need to eat.

\textbf{Notes (JA)}:\\
「...なくちゃ」は、義務や必要性を表すカジュアルな表現です。「...なきゃ」や「...なければならない」と訳されます。カジュアルな会話でよく使われ、相手に対して義務や必要性を伝える役割があります。この表現は、あまり堅苦しくなく、親しい友人同士の会話でよく使われます。

\textbf{Notes (EN)}:\\
"...nakucha" is a casual expression used to convey obligation or necessity. It can be translated as "I have to..." or "I need to..." This phrase is commonly used in informal conversations to express the speaker's obligation or necessity. It is not too formal and is often used among close friends.

\newpage

\section*{401. No doubt about it}

\textbf{Date}: 2025.05.16 (Fri)\\
\textbf{Index}: machigainai\\
\textbf{Expression (JA)}: (間違:まちが)いない\\
\textbf{Romanization}: Machigai nai\\
\textbf{Expression (EN)}: No doubt about it

\textbf{Variants (JA)}: 間違いない / 間違いないって / 間違いないよ / 間違いないから\\
\textbf{Variants (EN)}: No doubt about it / I'm telling you / I'm sure of it / Definitely

\textbf{Intent (JA)}: 確信, 強調, 同意の要求\\
\textbf{Intent (EN)}: certainty, emphasis, seeking agreement

\textbf{Tags (JA)}: 即時文法, 確信, 強調, 断言, 親密\\
\textbf{Tags (EN)}: immediate grammar, certainty, emphasis, assertion, intimacy

\textbf{Adjusted Expression (JA)}:\\
それは確実です／間違いありません\\
\textbf{Adjusted Expression (EN)}:\\
That is certain / There is no mistake about it

\textbf{Example (JA)}:\\
A「(明日:あした)は(晴:は)れるかな？」B「(間違:まちが)いないよ」\\
\textbf{Example (EN)}:\\
A: Will it be sunny tomorrow? B: No doubt about it.

\textbf{Notes (JA)}:\\
「間違いない」は、直感的な確信や経験に基づく判断を即時に表現する形式で、友人同士の親密な会話で多用される。末尾に「って」「よ」「から」などを加えることで、柔らかさや共感の度合いを調整する。

\textbf{Notes (EN)}:\\
'Machigai nai' expresses immediate certainty based on intuition or past experience. Common in casual or intimate conversations, it often ends with particles like 'tte', 'yo', or 'kara' to adjust tone and seek agreement.

\newpage

\section*{400. So, about that...}

\textbf{Date}: 2025.05.15 (Thu)\\
\textbf{Index}: sorehamataatode,\\
\textbf{Expression (JA)}: それはまたあとで、\\
\textbf{Romanization}: Sore wa mata ato de\\
\textbf{Expression (EN)}: So, about that later...

\textbf{Variants (JA)}: それはまたあとで / それはまた後でね\\
\textbf{Variants (EN)}: So, about that later / We'll talk about that later

\textbf{Intent (JA)}: 話の導入, 親密, 共有\\
\textbf{Intent (EN)}: introduction, intimacy, sharing

\textbf{Tags (JA)}: 即時文法, 省略, 状況依存型, 話の導入, 親密, 共有\\
\textbf{Tags (EN)}: immediate grammar, ellipsis, context-dependent, introduction, intimacy, sharing

\textbf{Adjusted Expression (JA)}:\\
それはまたのちほどお話しします。\\
\textbf{Adjusted Expression (EN)}:\\
I'll talk about that later.

\textbf{Example (JA)}:\\
A「どんな(問題:もんだい)が(試験:しけん)に出るんですか」B「それはまたあとで」\\
\textbf{Example (EN)}:\\
A: What kind of questions will be on the exam? B: We'll talk about that later.

\textbf{Notes (JA)}:\\
「それはまたあとで」は、話の導入や親密さを示すカジュアルな表現です。「So, about that later...」や「We'll talk about that later」と訳されます。カジュアルな会話でよく使われ、相手に対して話の導入を示す役割があります。

\textbf{Notes (EN)}:\\
"Sore wa mata ato de" is a casual expression used to introduce a topic or to show intimacy. It can be translated as "So, about that later..." or "We'll talk about that later." This phrase is commonly used in informal conversations to indicate the beginning of a story or to create a sense of closeness with the listener.

\newpage

\section*{399. So, you know...}

\textbf{Date}: 2025.05.14 (Wed)\\
\textbf{Index}: soregasa\\
\textbf{Expression (JA)}: それがさ、\\
\textbf{Romanization}: Sore ga sa\\
\textbf{Expression (EN)}: So, you know...

\textbf{Variants (JA)}: それがさ、 / それがさあ、 / それがね、 / それがねえ、\\
\textbf{Variants (EN)}: So, you know... / So, like... / So, you see... / So, you know what?

\textbf{Intent (JA)}: 話の導入, 親密, 共有\\
\textbf{Intent (EN)}: introduction, intimacy, sharing

\textbf{Tags (JA)}: 即時文法, 省略, 状況依存型, 話の導入, 親密, 共有\\
\textbf{Tags (EN)}: immediate grammar, ellipsis, context-dependent, introduction, intimacy, sharing

\textbf{Adjusted Expression (JA)}:\\
それで、お(話:はなし)というのは、...\\
\textbf{Adjusted Expression (EN)}:\\
So, the story is ...

\textbf{Example (JA)}:\\
A「(無理:むり)だって(思:おも)ってたんだけど、それがさ、やってみた(案外:あんがい)らうまくいったんだよね」B「え、ほんとに？」\\
\textbf{Example (EN)}:\\
A: I thought it was impossible, but you know what? When I tried it, it actually went surprisingly well. B: Oh, really?

\textbf{Notes (JA)}:\\
「それがさ、」は、直前に語られた内容とは異なる展開や、予想外の事実を切り出すときに使われるカジュアルな話し言葉です。前提となる内容（たとえば「無理だよと言われた」など）に対して、それを覆すような展開や、意外性を含んだ事実を語る導入句として機能します。英語では "So, you know what?" や "“But then..." のような語り出しに近く、相手の注意を引きつける語用効果があります。単なる話題転換ではなく、「聞いて！意外なことが起きたの」という語り手の感情や態度が表現されるのが特徴です。

\textbf{Notes (EN)}:\\
"Sore ga sa" is a casual expression used to introduce an unexpected or contrasting development following the previous statement. It typically functions as a narrative pivot, drawing the listener's attention to a surprising or emotionally significant event that contradicts or complicates what came before. In English, it is similar to phrases like "So, you know what?" or "But then..." Unlike neutral topic shifts, this expression carries the speaker's emotional engagement and signals a moment worth listening to.

\newpage

\section*{398. So, you know...}

\textbf{Date}: 2025.05.13 (Tue)\\
\textbf{Index}: dene\\
\textbf{Expression (JA)}: でね、\\
\textbf{Romanization}: De ne\\
\textbf{Expression (EN)}: So, you know...

\textbf{Variants (JA)}: でね、 / でさ、 / でさあ、 / でねえ、\\
\textbf{Variants (EN)}: So, you know... / So, like... / So, you see... / So, you know what?

\textbf{Intent (JA)}: 話の導入, 親密, 共有\\
\textbf{Intent (EN)}: introduction, intimacy, sharing

\textbf{Tags (JA)}: 即時文法, 省略, 状況依存型, 話の導入, 親密, 共有\\
\textbf{Tags (EN)}: immediate grammar, ellipsis, context-dependent, introduction, intimacy, sharing

\textbf{Adjusted Expression (JA)}:\\
そこで、お(話:はなし)というのは、...\\
\textbf{Adjusted Expression (EN)}:\\
So, the story is ...

\textbf{Example (JA)}:\\
A「きのう、(買:か)い(物:もの)に(行:い)ったの、でね、そこで、(誰:だれ)に(会:あ)ったと(思:おも)う？」\\
\textbf{Example (EN)}:\\
A: Yesterday, I went shopping, and you know what? I met someone there.

\textbf{Notes (JA)}:\\
「でね」は、話の導入や親密さを示すカジュアルな表現です。「So, you know...」や「So, like...」と訳されます。カジュアルな会話でよく使われ、相手に対して話の導入を示す役割があります。

\textbf{Notes (EN)}:\\
"De ne" is a casual expression used to introduce a story or to show intimacy. It can be translated as "So, you know..." or "So, like..." This phrase is commonly used in informal conversations to indicate the beginning of a story or to create a sense of closeness with the listener.

\newpage

\section*{397. Yeah, yeah}

\textbf{Date}: 2025.05.12 (Mon)\\
\textbf{Index}: haihai\\
\textbf{Expression (JA)}: はいはい\\
\textbf{Romanization}: Hai hai\\
\textbf{Expression (EN)}: Yeah, yeah

\textbf{Variants (JA)}: はいはい / はいはい、わかりました / はいはい、わかりましたよ / はいはい、わかりましたよ、もう\\
\textbf{Variants (EN)}: Yeah, yeah / Okay, okay, I got it / Yeah, yeah, I understand / Yeah, yeah, I got it already

\textbf{Intent (JA)}: 同意, あきらめ, 妥協\\
\textbf{Intent (EN)}: agreement, resignation, compromise

\textbf{Tags (JA)}: 即時文法, 省略, 状況依存型, 同意, あきらめ, 妥協, 繰り返し\\
\textbf{Tags (EN)}: immediate grammar, ellipsis, context-dependent, agreement, resignation, compromise, repetition

\textbf{Adjusted Expression (JA)}:\\
はい、わかりました。\\
\textbf{Adjusted Expression (EN)}:\\
Yes, I understand.

\textbf{Example (JA)}:\\
A「(明日:あした)こそ、やってよ」B「はいはい、わかりました」\\
\textbf{Example (EN)}:\\
A: Tomorrow, you have to do it. B: Yeah, yeah, I got it.

\textbf{Notes (JA)}:\\
「はいはい」は、相手の言葉に対して軽く同意するカジュアルな表現です。「Yeah, yeah」や「Okay, okay」と訳されます。カジュアルな会話でよく使われ、相手に対して軽い同意を示す役割があります。この表現は、あまり本気ではないが、相手の言葉を受け入れるときに使われます。

\textbf{Notes (EN)}:\\
"Hai hai" is a casual expression used to lightly agree with someone. It can be translated as "Yeah, yeah" or "Okay, okay." This phrase is commonly used in informal conversations to show a light agreement with the speaker. It is often used when you don't take the statement too seriously but still want to acknowledge it.

\newpage

\section*{396. If I had to say}

\textbf{Date}: 2025.05.11 (Sun)\\
\textbf{Index}: aeteiunara\\
\textbf{Expression (JA)}: あえて(言:い)うなら\\
\textbf{Romanization}: Aete iu nara\\
\textbf{Expression (EN)}: If I had to say

\textbf{Variants (JA)}: あえて言うなら / あえて言えば / あえて言うと / あえて言うならば\\
\textbf{Variants (EN)}: If I had to say / If I had to put it that way / If I had to say it / If I had to express it

\textbf{Intent (JA)}: 自己開示, 率直, 親密\\
\textbf{Intent (EN)}: self-disclosure, frankness, intimacy

\textbf{Tags (JA)}: 即時文法, 省略, 状況依存型, 自己開示, 率直, 親密\\
\textbf{Tags (EN)}: immediate grammar, ellipsis, context-dependent, self-disclosure, frankness, intimacy

\textbf{Adjusted Expression (JA)}:\\
あえて申し上げれるのなら、\\
\textbf{Adjusted Expression (EN)}:\\
If I were to express it more formally, ...

\textbf{Example (JA)}:\\
A「どっちが(好:す)き？」B「(あえて言:い)うなら、どちらも(好:す)きじゃない」\\
\textbf{Example (EN)}:\\
A: Which do you like? B: If I had to say, I don't like either.

\textbf{Notes (JA)}:\\
「あえて言うなら」は、あえて自分の意見や気持ちを伝えるための表現です。「If I had to say」や「If I had to put it that way」と訳されます。カジュアルな会話でよく使われ、相手に対して自分の気持ちを率直に伝える役割があります。この表現は言いにくいこと、否定的なことをどうしても言わなければならないときに使えばよいでしょう。一方で、どちらも良くて選びにくいときに、無理に一方を選ぶ場合には、「どちらかと言えば、」と言うとやわらかい印象になります。

\textbf{Notes (EN)}:\\
"Aete iu nara" is an expression used to convey one's opinion or feelings in a frank manner. It can be translated as "If I had to say" or "If I had to put it that way." This phrase is commonly used in casual conversations to express one's feelings openly to the listener. It is particularly useful when you need to say something difficult or negative. Conversely, if you are in a situation where both options are good and you find it hard to choose, you can use "dochira ka to ieba" to soften the impression.

\newpage

\section*{395. To be honest}

\textbf{Date}: 2025.05.10 (Sat)\\
\textbf{Index}: shoujiki\\
\textbf{Expression (JA)}: (正直:しょうじき)、\\
\textbf{Romanization}: Shoujiki\\
\textbf{Expression (EN)}: To be honest

\textbf{Variants (JA)}: 正直に言うと / 正直に言えば / 正直なところ / 正直な話\\
\textbf{Variants (EN)}: To be honest / Honestly speaking / To be frank / In all honesty

\textbf{Intent (JA)}: 率直, 自己開示, 親密\\
\textbf{Intent (EN)}: frankness, self-disclosure, intimacy

\textbf{Tags (JA)}: 即時文法, 省略, 状況依存型, 率直, 自己開示, 親密\\
\textbf{Tags (EN)}: immediate grammar, ellipsis, context-dependent, frankness, self-disclosure, intimacy

\textbf{Adjusted Expression (JA)}:\\
正直に言いますと、そう思います。\\
\textbf{Adjusted Expression (EN)}:\\
frankly speaking, I think so.

\textbf{Example (JA)}:\\
A「(正直:しょうじき)、あまり(好:す)きじゃない」B「え、そうなの？」\\
\textbf{Example (EN)}:\\
A: To be honest, I don't really like it. B: Oh, really?

\textbf{Notes (JA)}:\\
「正直」は、率直に自分の気持ちや意見を伝えるための表現です。「To be honest」や「Honestly speaking」と訳されます。カジュアルな会話でよく使われ、相手に対して自分の気持ちを率直に伝える役割があります。

\textbf{Notes (EN)}:\\
"Shoujiki" is an expression used to convey one's feelings or opinions frankly. It can be translated as "to be honest" or "honestly speaking." This phrase is commonly used in casual conversations to express one's feelings openly to the listener. 

\newpage

\section*{394. To put it simply}

\textbf{Date}: 2025.05.09 (Fri)\\
\textbf{Index}: ketsuronkaraieba\\
\textbf{Expression (JA)}: (結論:けつろん)から(言:い)えば\\
\textbf{Romanization}: Ketsuron kara ieba\\
\textbf{Expression (EN)}: To put it simply

\textbf{Variants (JA)}: 結論から言えば / 結論から言うと / 結論だけ言えば / 一言で言えば\\
\textbf{Variants (EN)}: To put it simply / To sum it up / If I had to say just the conclusion / To make a long story short

\textbf{Intent (JA)}: 要点提示, 話の切り出し, 前置き省略\\
\textbf{Intent (EN)}: main point first, discourse entry, omit preface

\textbf{Tags (JA)}: 即時文法, 視点提示, 談話標識, 要点優先, 省略, 口語\\
\textbf{Tags (EN)}: immediate grammar, viewpoint marking, discourse marker, point-first, ellipsis, spoken

\textbf{Adjusted Expression (JA)}:\\
結論から言うと、そういうことです。\\
\textbf{Adjusted Expression (EN)}:\\
To put it simply, that's how it is.

\textbf{Example (JA)}:\\
A「(結論:けつろん)から(言:い)えば、それはできないということ」B「じゃ、これまでは何だったの？」\\
\textbf{Example (EN)}:\\
A: To put it simply, that means it can't be done. B: Then what was all this about?

\textbf{Notes (JA)}:\\
「結論から言えば」は、話の要点を先に伝えるための表現です。カジュアルな会話でよく使われ、相手に対して要点を明確に伝える役割があります。「結論から言うと」や「一言で言えば」といった意味合いも含まれますが、あきらかに「今、すぐに」という意味ではありません。だいたい、「結論から言えば、...」と聞けば、すべてがすべてではありませんが、あまり良いことが待っているようには思えませんね。

\textbf{Notes (EN)}:\\
"Ketsuron kara ieba" is an expression used to convey the main point of a conversation or discussion first. It is commonly used in casual conversations to clarify the main point for the listener. This phrase can also imply meanings like "to sum it up" or "to make a long story short," but it does not necessarily mean "right now." Generally, when you hear "ketsuron kara ieba, ...", it doesn't always mean something good is coming.

\newpage

\section*{393. I don't care}

\textbf{Date}: 2025.05.08 (Thu)\\
\textbf{Index}: doudemoii\\
\textbf{Expression (JA)}: もう、どうでもいい\\
\textbf{Romanization}: Mou, dou demo ii\\
\textbf{Expression (EN)}: I don't care anymore

\textbf{Variants (JA)}: もう、どうでもいい / もう、どうでもよくなった / もう、どうでもいいや / もう、どうでもいいよ\\
\textbf{Variants (EN)}: I don't care anymore / I don't care about it anymore / I don't care at all / I don't care anymore, you know?

\textbf{Intent (JA)}: あきらめ, 自己納得, 妥協\\
\textbf{Intent (EN)}: resignation, self-acceptance, compromise

\textbf{Tags (JA)}: 即時文法, 省略, 状況依存型, あきらめ, 自己納得, 妥協\\
\textbf{Tags (EN)}: immediate grammar, ellipsis, context-dependent, resignation, self-acceptance, compromise

\textbf{Adjusted Expression (JA)}:\\
もう、どうでもいいと思います。\\
\textbf{Adjusted Expression (EN)}:\\
I think I don't care anymore.

\textbf{Example (JA)}:\\
A「お(好:す)きなものを(買:か)ってきましょうね」B「もう、どうでもいい」\\
\textbf{Example (EN)}:\\
A: Let's buy what you like. B: I don't care anymore.

\textbf{Notes (JA)}:\\
「もう、どうでもいい」は、状況を軽く受け入れることを伝えるカジュアルな表現です。「I don't care anymore」や「I don't care about it anymore」と訳されます。意にそぐわないことを妥協して受け入れる状況でよく使えます。

\textbf{Notes (EN)}:\\
"Mou, dou demo ii" is a casual expression used to convey a sense of resignation or acceptance of a situation, often in a light-hearted or humorous way. It can be translated as "I don't care anymore" or "I don't care about it anymore." This phrase is commonly used in informal conversations when someone is faced with an unexpected or undesirable situation but chooses to accept it anyway. 

\newpage

\section*{392. After that}

\textbf{Date}: 2025.05.07 (Wed)\\
\textbf{Index}: arekara\\
\textbf{Expression (JA)}: あれから\\
\textbf{Romanization}: Are kara\\
\textbf{Expression (EN)}: After that

\textbf{Variants (JA)}: あれから / それから / その後\\
\textbf{Variants (EN)}: after that / since then / afterwards

\textbf{Intent (JA)}: 回想, 導入, 共有, 親密\\
\textbf{Intent (EN)}: recollection, introduction, sharing, intimacy

\textbf{Tags (JA)}: 即時文法, 話題導入, 状況依存型, 親密, 回想\\
\textbf{Tags (EN)}: immediate grammar, topic-introduction, context-dependent, intimacy, recollection

\textbf{Adjusted Expression (JA)}:\\
あれから、みんなで食事をして帰りました。\\
\textbf{Adjusted Expression (EN)}:\\
After that, we all had a meal and went home.

\textbf{Example (JA)}:\\
A「きのう、あれから、どうした？」B「うん、(財布:さいふ)も、(カード:かーど)も、(現金:げんきん)も、(全部:ぜんぶ)、いっしょに(見:み)つかったよ」\\
\textbf{Example (EN)}:\\
A: So, what happened after that yesterday? B: Yeah, I found my wallet, card, and cash all together.

\textbf{Notes (JA)}:\\
「あれから」は、過去の出来事を思い出して話し始めるときに使う表現です。この表現は、会話の中で相手に対して親密さや共感を示す役割も果たします。「あれ」というのは会話をしている両者が共通して知っていることを指します。たとえば、AさんとBさんが「昨日、あれからどうした？」と話している場合、「あれ」は「昨日の出来事」を指しています。つまり、AさんとBさんは「昨日の出来事」を共通の認識として持っているということです。

\textbf{Notes (EN)}:\\
"Are kara" is used to begin recalling a past event in conversation. It also serves to express a sense of closeness or shared memory with the listener. The term "are" refers to something that both speakers are aware of, such as a shared experience or event. For example, if A and B are talking about "what happened after that yesterday," "are" refers to "yesterday's event." This means that A and B share a common understanding of "yesterday's event."

\newpage

\section*{391. From now on}

\textbf{Date}: 2025.05.06 (Tue)\\
\textbf{Index}: korekara\\
\textbf{Expression (JA)}: これから\\
\textbf{Romanization}: kore kara\\
\textbf{Expression (EN)}: From now on

\textbf{Variants (JA)}: これから / 今から / ここから\\
\textbf{Variants (EN)}: from now on / starting now / from here

\textbf{Intent (JA)}: 開始, 意志, 未来, 区切り\\
\textbf{Intent (EN)}: beginning, intention, future, transition

\textbf{Tags (JA)}: 即時文法, 時の区切り, 意志表明, 状況転換, モーダル表現\\
\textbf{Tags (EN)}: immediate grammar, temporal boundary, statement of intent, context shift, modal expression

\textbf{Adjusted Expression (JA)}:\\
これから始めます。\\
\textbf{Adjusted Expression (EN)}:\\
I will start from now.

\textbf{Example (JA)}:\\
A「ねえ、いつやるの？」B「これから」\\
\textbf{Example (EN)}:\\
A: Hey, when are you doing it? B: From now on.

\textbf{Notes (JA)}:\\
「これから」は、これからの行動や出来事を示す表現です。未来の意志や計画を伝える際に使われます。カジュアルな会話でよく使われ、相手に対して行動の開始を伝える役割があります。「すぐに」や「今から」といった意味合いも含まれますが、あきらかに「今、すぐに」という意味ではありませんね。

\textbf{Notes (EN)}:\\
"Kore kara" is an expression that indicates actions or events that will happen from now on. It is used to convey future intentions or plans. This phrase is commonly used in casual conversations to inform the listener about the start of an action. It can also imply meanings like "immediately" or "from now," but it does not necessarily mean "right now."

\newpage

\section*{390. So, about that time...}

\textbf{Date}: 2025.05.05 (Mon)\\
\textbf{Index}: anotokine\\
\textbf{Expression (JA)}: あのときね\\
\textbf{Romanization}: Ano toki ne\\
\textbf{Expression (EN)}: So back then, you know...

\textbf{Variants (JA)}: あのときね / あのときさ / あのときはね\\
\textbf{Variants (EN)}: So back then, you know... / At that time, remember? / Back in that moment...

\textbf{Intent (JA)}: 回想, 導入, 共有, 親密\\
\textbf{Intent (EN)}: recollection, introduction, sharing, intimacy

\textbf{Tags (JA)}: 即時文法, 話題導入, 状況依存型, 親密, 回想\\
\textbf{Tags (EN)}: immediate grammar, topic-introduction, context-dependent, intimacy, recollection

\textbf{Adjusted Expression (JA)}:\\
あのときのことを思い出しました。\\
\textbf{Adjusted Expression (EN)}:\\
I remembered what happened back then.

\textbf{Example (JA)}:\\
A「けがしたの？」B「うん、(駅:えき)で」A「ああ、あのときね」\\
\textbf{Example (EN)}:\\
A: Did you get hurt? B: Yeah, at the station. A: Oh, that time, right?

\textbf{Notes (JA)}:\\
「あのときね」は、過去の出来事を思い出して話し始めるときに使う表現です。この表現は、会話の中で相手に対して親密さや共感を示す役割も果たします。

\textbf{Notes (EN)}:\\
"Ano toki ne" is used to begin recalling a past event in conversation. It also serves to express a sense of closeness or shared memory with the listener.

\newpage

\section*{389. いくらなんでも...}

\textbf{Date}: 2025.05.04 (Sun)\\
\textbf{Index}: ikuranandemo\\
\textbf{Expression (JA)}: いくらなんでも...\\
\textbf{Romanization}: Ikura nan demo...\\
\textbf{Expression (EN)}: No matter how you look at it... / Come on, seriously...

\textbf{Variants (JA)}: いくらなんでも / いくらなんでもさ / いくらなんでもそれは / いくらなんでもひどい\\
\textbf{Variants (EN)}: No matter how you look at it / Come on, seriously / That’s just too much / This is going too far

\textbf{Intent (JA)}: 驚き, 非難, 拒否, 感情の爆発, 同意の拒絶\\
\textbf{Intent (EN)}: surprise, criticism, rejection, emotional outburst, refusal to agree

\textbf{Tags (JA)}: 即時文法, 感情表現, 状況依存型, 非難, 確認, 曖昧な終止\\
\textbf{Tags (EN)}: immediate grammar, emotional expression, context-dependent, criticism, confirmation, unfinished ending

\textbf{Adjusted Expression (JA)}:\\
いくらなんでも、それはひどすぎます。\\
\textbf{Adjusted Expression (EN)}:\\
No matter how you look at it, that's just too much.

\textbf{Example (JA)}:\\
いくらなんでも、それはないんじゃない？\\
\textbf{Example (EN)}:\\
Come on, that’s just too much, don’t you think?

\textbf{Notes (JA)}:\\
「いくらなんでも」は、文法的には副詞的に前置されるが、即時文法では後続の文を曖昧に省略したり、語尾をぼかしたりして、感情の高ぶりや言いにくさを表現する。しばしば相手に対する批判・非難の前置きとなり、共感や同意を求める機能もある。

\textbf{Notes (EN)}:\\
Although 'ikura nan demo' functions as a concessive adverb, in immediate grammar it often appears at the beginning or middle of an utterance, trailing off to imply criticism or disbelief. It softens or signals an emotional explosion while inviting empathy or support from the listener.

\newpage

\section*{388. Seriously}

\textbf{Date}: 2025.05.03 (Sat)\\
\textbf{Index}: maji\\
\textbf{Expression (JA)}: (マジ:まじ)\\
\textbf{Romanization}: Maji\\
\textbf{Expression (EN)}: Seriously

\textbf{Variants (JA)}: まじ / マジ / マジで / 本気で / 本当に\\
\textbf{Variants (EN)}: Seriously / For real / Really / Honestly / No kidding

\textbf{Intent (JA)}: 驚き, 確認, 興味, 親密\\
\textbf{Intent (EN)}: surprise, confirmation, interest, intimacy

\textbf{Tags (JA)}: 即時文法, 省略, 状況依存型, 驚き, 確認, 親密, 若者語, リアルタイム性\\
\textbf{Tags (EN)}: immediate grammar, ellipsis, context-dependent, surprise, confirmation, intimacy, youth language, strong real-time nature

\textbf{Adjusted Expression (JA)}:\\
(本気:ほんき)ですか。\\
\textbf{Adjusted Expression (EN)}:\\
Are you serious?

\textbf{Example (JA)}:\\
A「(明日:あした)、(試験:しけん)、あるよ」B「(マジ:まじ)？」\\
\textbf{Example (EN)}:\\
A: There's an exam tomorrow. B: Seriously?

\textbf{Notes (JA)}:\\
「マジ」は、何かに対する驚きや真剣さを伝えるカジュアルな表現です。「seriously」や「for real」と訳され、カジュアルな会話で信じられないことを表現したり、発言の真実性を強調したりするためによく使われます。このフレーズは友人や仲間の間でよく使われ、会話の中で親密さや近さを示すこともあります。よりフォーマルな場面では「本当に」や「本当ですか」と言うことが一般的です。

\textbf{Notes (EN)}:\\
"Maji" is a casual expression used to convey surprise or seriousness about something. It can be translated as "seriously" or "for real," and is often used in informal conversations to express disbelief or to emphasize the truth of a statement. This phrase is commonly used among friends or peers, and can also imply a sense of intimacy or closeness in the conversation. In more formal situations, one might use phrases like "hontou ni" or "hontou desu ka" instead, which are more polite ways to express a similar sentiment.

\newpage

\section*{387. Lend me ...}

\textbf{Date}: 2025.05.02 (Fri)\\
\textbf{Index}: kashite\\
\textbf{Expression (JA)}: (貸:か)して、...\\
\textbf{Romanization}: Kashite, ...\\
\textbf{Expression (EN)}: Lend me ...

\textbf{Variants (JA)}: 貸して / ちょっと貸して / 貸してくれる？ / 貸してよ / 貸してください\\
\textbf{Variants (EN)}: Lend me / Let me borrow / Can I borrow? / Let me have it / Please lend me

\textbf{Intent (JA)}: 要求, 確認, 関心, 親密, ラフさ（casualness）\\
\textbf{Intent (EN)}: request, confirmation, interest, intimacy, casualness

\textbf{Tags (JA)}: 即時文法, 省略, 命令, 依頼, 口語\\
\textbf{Tags (EN)}: Immediate Grammar, ellipsis, imperative, request, spoken

\textbf{Adjusted Expression (JA)}:\\
それを(貸:か)していただけませんか。\\
\textbf{Adjusted Expression (EN)}:\\
Please lend me that.

\textbf{Example (JA)}:\\
A「これ、(貸:か)して」B「(貸:か)してあげるよ」\\
\textbf{Example (EN)}:\\
A: Can you lend me this? B: Sure, I'll lend it to you.

\textbf{Notes (JA)}:\\
「貸して」は、誰かに何かを貸したり借りたりするように頼むカジュアルな表現です。「lend me」や「let me borrow」と訳され、カジュアルな会話で一時的に何かを使いたいという気持ちを表すためによく使われます。このフレーズは、友人からアイテムを借りるときなど、日常的な状況でよく使われます。

\textbf{Notes (EN)}:\\
"Kashite" is a casual expression used to request someone to lend or borrow something. It can be translated as "lend me" or "let me borrow," and is often used in informal conversations to express a desire to use something temporarily. This phrase is commonly used in everyday situations, such as when asking for a favor or borrowing an item from a friend.

\newpage

\section*{386. Oh well}

\textbf{Date}: 2025.05.01 (Thu)\\
\textbf{Index}: maikka\\
\textbf{Expression (JA)}: ま、いっか\\
\textbf{Romanization}: Ma, ikka\\
\textbf{Expression (EN)}: Oh well. / Never mind.

\textbf{Variants (JA)}: ま、いっか / ま、いっか / ま、いいか\\
\textbf{Variants (EN)}: Oh well. / Never mind. / It's okay. / It's all right. / It's fine.

\textbf{Intent (JA)}: 妥協, あきらめ, 自己納得\\
\textbf{Intent (EN)}: compromise, resignation, self-acceptance

\textbf{Tags (JA)}: 即時文法, 省略, 状況依存型, 妥協, あきらめ, 自己納得, リアルタイム性\\
\textbf{Tags (EN)}: immediate grammar, ellipsis, context-dependent, compromise, resignation, self-acceptance, strong real-time nature

\textbf{Adjusted Expression (JA)}:\\
とりあえず、そうしましょう。\\
\textbf{Adjusted Expression (EN)}:\\
For now, let's do it that way.

\textbf{Example (JA)}:\\
A「(今日:きょう)は(カレー:かれー)よ」B「え、(今日:きょう)も？ま、いっか」\\
\textbf{Example (EN)}:\\
A: It's curry today. B: Oh, again? Oh well.

\textbf{Notes (JA)}:\\
「ま、いっか」は、状況を軽く受け入れることを伝えるカジュアルな表現です。「Oh well」や「Never mind」と訳されます。意にそぐわないことを妥協して受け入れる状況でよく使えます。このフレーズは、予期しない状況や望ましくない状況に直面したときに使われますが、あえて受け入れることを選ぶときによく使われます。妥協や自己納得の感覚を示すこともあります。よりフォーマルな場面では「それでいいです」や「それで構わない」と言うことが一般的です。

\textbf{Notes (EN)}:\\
"Maikka" is a casual expression used to convey a sense of resignation or acceptance of a situation, often in a light-hearted or humorous way. It can be translated as "Oh well" or "Never mind." This phrase is commonly used in informal conversations when someone is faced with an unexpected or undesirable situation but chooses to accept it anyway. It can also imply a sense of compromise or self-acceptance, indicating that the speaker is willing to let go of their initial expectations. In more formal situations, one might use phrases like "sore de ii desu" or "sore de kamawanai" instead, which are more polite ways to express a similar sentiment.

\newpage

\section*{385. Almost all}

\textbf{Date}: 2025.04.30 (Wed)\\
\textbf{Index}: daitai\\
\textbf{Expression (JA)}: だいたい\\
\textbf{Romanization}: Daitai\\
\textbf{Expression (EN)}: Almost all

\textbf{Variants (JA)}: だいたい / ほとんど / おおよそ / 大体 / ほぼ\\
\textbf{Variants (EN)}: Almost all / Mostly / Roughly / About / Nearly

\textbf{Intent (JA)}: 概算, おおよそ, 大体, ほぼ\\
\textbf{Intent (EN)}: approximation, roughly, about, almost

\textbf{Tags (JA)}: 即時文法, 省略, 状況依存型, 概算, おおよそ, リアルタイム性\\
\textbf{Tags (EN)}: immediate grammar, ellipsis, context-dependent, approximation, roughly, strong immediacy

\textbf{Adjusted Expression (JA)}:\\
ほぼすべて、(完了:かんりょう)しました。\\
\textbf{Adjusted Expression (EN)}:\\
Almost all of it is done.

\textbf{Example (JA)}:\\
A「(宿題:しゅくだい)は、(終:お)わった？」B「うん、だいたい」\\
\textbf{Example (EN)}:\\
A: Did you finish your homework? B: Yeah, almost all of it.

\textbf{Notes (JA)}:\\
「だいたい」は、何かの概算やおおよその見積もりを示すカジュアルな表現です。「almost all」や「mostly」、「roughly」と訳され、カジュアルな会話であまり具体的に言わずに一般的な考えを伝えるためによく使われます。このフレーズは、時間、数量、程度について話す日常的な状況でよく使われます。「だいたい」ということは別の見方をすれば、「まだ終わっていない」ということですね。

\textbf{Notes (EN)}:\\
"Daitai" is a casual expression used to indicate an approximation or a rough estimate of something. It can be translated as "almost all," "mostly," or "roughly," and is often used in informal conversations to convey a general idea without being too specific. This phrase is common in everyday situations, such as when discussing time, quantity, or degree. In another way, saying "daitai" can also imply that something is "not finished yet.

\newpage

\section*{384. どういうこと？}

\textbf{Date}: 2025.04.29 (Tue)\\
\textbf{Index}: douiukoto\\
\textbf{Expression (JA)}: どういうこと？\\
\textbf{Romanization}: Dou iu koto?\\
\textbf{Expression (EN)}: What do you mean?

\textbf{Variants (JA)}: どういうこと？ / どういう意味？ / どういうことなの？ / どういうことだろう？ / どういうことなんだろう？\\
\textbf{Variants (EN)}: What do you mean? / What does that mean? / What does it mean? / What is that supposed to mean? / What are you talking about?

\textbf{Intent (JA)}: 確認, 疑問, 説明要求, 理解, 興味\\
\textbf{Intent (EN)}: confirmation, question, request for explanation, understanding, interest

\textbf{Tags (JA)}: 即時文法, 省略, 状況依存型, 疑問, 説明要求, リアルタイム性\\
\textbf{Tags (EN)}: immediate grammar, ellipsis, context-dependent, question, request for explanation, strong real-time nature

\textbf{Adjusted Expression (JA)}:\\
どういうことですか？\\
\textbf{Adjusted Expression (EN)}:\\
What do you mean?

\textbf{Example (JA)}:\\
A「(予定:よてい)は、中止」B「え、どういうこと？」\\
\textbf{Example (EN)}:\\
A: The plan is canceled. B: Huh? What do you mean?

\textbf{Notes (JA)}:\\
「どういうこと？」は、何かについての説明や明確化を求めるカジュアルな表現です。「What do you mean?」や「What does that mean?」と訳され、発言や状況についての詳細を知りたいときによく使われます。よりフォーマルな場面では「どういう意味ですか？」と言います。

\textbf{Notes (EN)}:\\
"Dou iu koto?" is a casual expression used to ask for clarification or explanation about something. It can be translated as "What do you mean?" or "What does that mean?" This phrase is often used in informal conversations when someone is confused or needs more information about a statement or situation. In more formal situations, one might use "dou iu imi desu ka?" instead, which is a more polite way to ask the same question.

\newpage

\section*{383. そのつもりだった}

\textbf{Date}: 2025.04.28 (Mon)\\
\textbf{Index}: tsumoridatta\\
\textbf{Expression (JA)}: そのつもりだった\\
\textbf{Romanization}: Tsumoridatta\\
\textbf{Expression (EN)}: I was going to

\textbf{Variants (JA)}: そのつもりだった / そうするつもりだった / やるつもりだった / 行くつもりだった\\
\textbf{Variants (EN)}: I meant to / I was going to / I intended to / I was about to (but...)

\textbf{Intent (JA)}: 過去の意図, 過去の約束, 確認, 同意, 承諾\\
\textbf{Intent (EN)}: past intention, past promise, confirmation, agreement, acceptance

\textbf{Tags (JA)}: 即時文法, 感情表現, 状況依存型, 省略, 意図表明, 過去の意図\\
\textbf{Tags (EN)}: immediate grammar, emotional expression, context-dependent, ellipsis, expression of intention, past intention

\textbf{Adjusted Expression (JA)}:\\
["〜する予定でした","〜する意向でした","〜しようと考えていました"]\\
\textbf{Adjusted Expression (EN)}:\\
["I had planned to do ～.","I intended to do ～.","I had the intention to do ～."]

\textbf{Example (JA)}:\\
A「(行:い)くつもりだったの？」B「(行:い)くつもりだったけど、(行:い)けなかった」\\
\textbf{Example (EN)}:\\
A: Were you going to go? B: I meant to go, but I couldn't.

\textbf{Notes (JA)}:\\
「つもりだった」は、実行されなかった過去の意図や計画を示すカジュアルな表現です。「つもり」は「intention」や「plan」を意味し、「だった」は過去形を表します。このフレーズは、以前の意図を果たせなかったことに対する後悔や失望を表現するためによく使われます。重要なことは、この「つもり」は第1人称の意図しか表せないことです。たとえば、「行くつもりだった」と言った場合、相手は「行くつもりだったの？」と聞き返すことができます。これは、相手の意図を確認するための質問です。

\textbf{Notes (EN)}:\\
"Tsumoridatta" is a casual expression used to indicate a past intention or plan that was not fulfilled. It combines the word "tsumori" (meaning "intention" or "plan") with the past tense form "datta." This phrase is often used in informal conversations to express regret or disappointment about being unable to carry out a previous intention. The most important thing to note is that this "tsumori" can only express the first-person intention. For example, if you say "iku tsumori datta," the other person can ask back, "iku tsumori datta no?" This is a question to confirm the other person's intention.

\newpage

\section*{382. そのつもり}

\textbf{Date}: 2025.04.27 (Sun)\\
\textbf{Index}: tsumori\\
\textbf{Expression (JA)}: そのつもり\\
\textbf{Romanization}: Sono tsumori\\
\textbf{Expression (EN)}: That's the intention

\textbf{Variants (JA)}: そのつもり / ええ、そのつもり / もちろん、そのつもり / うん、そのつもり / ちゃんとそのつもり / ばっちりそのつもり / はい、そのつもりです\\
\textbf{Variants (EN)}: I will / Yeah, I will / Of course I will / Sure, I will / Absolutely, I will / Definitely, I will / Yes, I will

\textbf{Intent (JA)}: 意図, 約束, 確認, 同意, 承諾\\
\textbf{Intent (EN)}: intention, promise, confirmation, agreement, consent

\textbf{Tags (JA)}: 即時文法, 省略, 状況依存型, 意図表現, 確認応答\\
\textbf{Tags (EN)}: immediate grammar, omission, context-dependent, expression of intention, confirmation response

\textbf{Adjusted Expression (JA)}:\\
そうするつもりでいます。\\
\textbf{Adjusted Expression (EN)}:\\
I intend to do so.

\textbf{Example (JA)}:\\
A「(明日:あす)も(練習:れんしゅう)する？」 B「ええ、そのつもり。」\\
\textbf{Example (EN)}:\\
A: Will you practice tomorrow too? B: Yeah, I will.

\textbf{Notes (JA)}:\\
「そのつもり」は、行動の意図を即時的に肯定するカジュアルな表現です。「その」は「that」を、「つもり」は「意図」や「計画」を意味します。

\textbf{Notes (EN)}:\\
"Sono tsumori" is a casual, real-time expression to affirm one's intent to act. "Sono" means "that," and "tsumori" means "intention" or "plan."

\newpage

\section*{381. (君:きみ)の(番:ばん)だよ}

\textbf{Date}: 2025.04.26 (Sat)\\
\textbf{Index}: kiminobandayo\\
\textbf{Expression (JA)}: (君:きみ)の(番:ばん)だよ\\
\textbf{Romanization}: Kimi no ban da yo\\
\textbf{Expression (EN)}: (It's) your turn

\textbf{Variants (JA)}: 次は君だよ / 君の出番だよ / あなたの番ですよ / 次はあなたですよ\\
\textbf{Variants (EN)}: It's your turn. / You're up. / Now it's your turn. / Your move. / You're next.

\textbf{Intent (JA)}: 順番を伝える, 行動を促す, 交代を知らせる, 親しみを込めて呼びかける, タイミングを示す\\
\textbf{Intent (EN)}: indicating the turn, encouraging action, notifying a handover, calling warmly, indicating timing

\textbf{Tags (JA)}: 即時文法, 省略, 状況依存型, 促し, 交代, リアルタイム性\\
\textbf{Tags (EN)}: immediate grammar, ellipsis, context-dependent, encouragement, handover, strong real-time nature

\textbf{Adjusted Expression (JA)}:\\
(次:つぎ)の(方:かた)？\\
\textbf{Adjusted Expression (EN)}:\\
Next, please?

\textbf{Example (JA)}:\\
A「(君:きみ)の(番:ばん)だよ」B「(分:わ)かってるよ」\\
\textbf{Example (EN)}:\\
A: It's your turn. B: I know.

\textbf{Notes (JA)}:\\
「君の番だよ」は、誰かが何かをする番であることを示すカジュアルな表現です。「君」は「you」を意味し、「番」は「turn」や「order」を意味します。

\textbf{Notes (EN)}:\\
"Kimi no ban da yo" is a casual expression used to indicate that it is someone's turn to do something. It combines the word "kimi" (meaning "you") with "ban" (meaning "turn" or "order"). 

\newpage

\section*{380. You'll still have another chance}

\textbf{Date}: 2025.04.25 (Fri)\\
\textbf{Index}: tsugi ga aru\\
\textbf{Expression (JA)}: (次:つぎ)がある\\
\textbf{Romanization}: Tsugi ga aru\\
\textbf{Expression (EN)}: You'll still have another chance

\textbf{Variants (JA)}: 次があるさ / またチャンスあるよ / また今度がある / 今度こそ\\
\textbf{Variants (EN)}: There's always next time / You can try again / You'll get another chance / Better luck next time / Next time for sure

\textbf{Intent (JA)}: 励まし, 慰め, 再挑戦\\
\textbf{Intent (EN)}: encouragement, comfort, retry

\textbf{Tags (JA)}: 即時文法, 省略, 励まし, 慰め, 口語\\
\textbf{Tags (EN)}: Immediate Grammar, ellipsis, encouragement, comfort, spoken

\textbf{Adjusted Expression (JA)}:\\
今回は残念でしたが、またの機会があります\\
\textbf{Adjusted Expression (EN)}:\\
Although this time was unfortunate, there will be another opportunity.

\textbf{Example (JA)}:\\
A「(失敗:しっぱい)しちゃった」B「(次:つぎ)があるよ」\\
\textbf{Example (EN)}:\\
A: I failed. B: You'll still have another chance.

\textbf{Notes (JA)}:\\
「次がある」は、将来に別の機会やチャンスがあることを伝えるカジュアルな表現です。「次があるさ」や「またチャンスあるよ」と訳され、カジュアルな会話で失敗や落胆の後に励ましや慰めを提供するためによく使われます。

\textbf{Notes (EN)}:\\
"Tsugi ga aru" is a casual expression used to convey the idea that there will be another opportunity or chance in the future. It can be translated as "there's always next time" or "you can try again," and is often used in informal conversations to provide encouragement or comfort after a setback or disappointment. 

\newpage

\section*{379. Show me ...}

\textbf{Date}: 2025.04.24 (Thu)\\
\textbf{Index}: misete\\
\textbf{Expression (JA)}: (見:み)せて、...\\
\textbf{Romanization}: Misete, ...\\
\textbf{Expression (EN)}: Show me ...

\textbf{Variants (JA)}: 見せて / ちょっと見せて / 見せてくれる？ / 見せてよ / 見せなさい\\
\textbf{Variants (EN)}: Show me / Let me see / Can I see? / Let me have a look / Give me a peek

\textbf{Intent (JA)}: 要求, 確認, 関心\\
\textbf{Intent (EN)}: request, confirmation, interest

\textbf{Tags (JA)}: 即時文法, 省略, 命令, 依頼, 口語\\
\textbf{Tags (EN)}: Immediate Grammar, ellipsis, imperative, request, spoken

\textbf{Adjusted Expression (JA)}:\\
それを(見:み)せてください。\\
\textbf{Adjusted Expression (EN)}:\\
Please show me that.

\textbf{Example (JA)}:\\
A「これ、(新:あたら)しい(本:ほん)？ちょっと(見:み)せて？」\\
\textbf{Example (EN)}:\\
A: Is this a new book? Let me see it?

\textbf{Notes (JA)}:\\
「見せて」は、誰かに何かを見せたり、表示したりするように頼むカジュアルな表現です。「show me」や「let me see」と訳され、カジュアルな会話で何かに対する好奇心や関心を表すためによく使われます。

\textbf{Notes (EN)}:\\
"Misete" is a casual expression used to request someone to show or display something. It can be translated as "show me" or "let me see," and is often used in informal conversations to express curiosity or interest in something. 

\newpage

\section*{378. All right, even if we got some troubles}

\textbf{Date}: 2025.04.23 (Wed)\\
\textbf{Index}: kekka, o-rai\\
\textbf{Expression (JA)}: (結果:けっか)、(オーライ:おーらい)\\
\textbf{Romanization}: Kekka, o-rai\\
\textbf{Expression (EN)}: All right, even if we got some troubles

\textbf{Variants (JA)}: まあ、なんとかなった / うまくいった\\
\textbf{Variants (EN)}: All good / We're good / It worked out / All's well that ends well / No problem in the end / Made it

\textbf{Intent (JA)}: 結果, 解決, 安堵\\
\textbf{Intent (EN)}: result, solution, relief

\textbf{Tags (JA)}: 即時文法, 省略, 状況依存型, 結果, 安堵, 解決, 終結\\
\textbf{Tags (EN)}: immediate grammar, omission, context-dependent, result, solution, light tone, relief, closure

\textbf{Adjusted Expression (JA)}:\\
(終:お)わりよければすべてよし\\
\textbf{Adjusted Expression (EN)}:\\
All's well that ends well.

\textbf{Example (JA)}:\\
A「いろいろ(大変:たいへん)だったね」B「(結果:けっか)、(オーライ:おーらい)ですよ」\\
\textbf{Example (EN)}:\\
A: It was tough, wasn't it? B: All right, even if we got some troubles.

\textbf{Notes (JA)}:\\
「結果、オーライ」は、困難や課題があったにもかかわらず、最終的にすべてがうまくいったことを示すカジュアルな表現です。「結果」は「result」や「outcome」を意味し、「オーライ」は「all right」や「okay」を意味します。このフレーズは、障害を克服した後の安堵感や満足感を表現するためによく使われます。途中での困難も何もない人が時に使っているのを見ますが、こういう人のことを楽天家というんですよ。

\textbf{Notes (EN)}:\\
"Kekka, o-rai" is a casual expression used to indicate that despite difficulties or challenges, everything turned out well in the end. It combines the word "kekka" (meaning "result" or "outcome") with "o-rai" (meaning "all right" or "okay"). This phrase is often used in informal conversations to express relief or satisfaction after overcoming obstacles. It can also imply a sense of acceptance or resignation to the situation. Sometimes, I saw people using this phrase even when they didn't face any difficulties, but those people are called optimists.

\newpage

\section*{377. But,...}

\textbf{Date}: 2025.04.22 (Tue)\\
\textbf{Index}: demo,...\\
\textbf{Expression (JA)}: でも、...\\
\textbf{Romanization}: Demo, ...\\
\textbf{Expression (EN)}: But ...

\textbf{Variants (JA)}: でも、... / しかし、... / だけど、... / けど、... / とはいえ、... / それでも、...\\
\textbf{Variants (EN)}: But ... / However ... / Though ... / Still ... / Nevertheless ...

\textbf{Intent (JA)}: 否定, 反論, 対照, 対比, 逆接\\
\textbf{Intent (EN)}: negation, refutation, contrast, comparison, adversative

\textbf{Tags (JA)}: 即時文法, 省略, 状況依存型, 否定, 反論, リアルタイム性\\
\textbf{Tags (EN)}: immediate grammar, omission, context-dependent, negation, refutation, strong real-time nature

\textbf{Adjusted Expression (JA)}:\\
しかしながら、...\\
\textbf{Adjusted Expression (EN)}:\\
However,...

\textbf{Example (JA)}:\\
A「おいしいね」B「でも、あまり、(見:み)た(目:め)がきれいじゃない」\\
\textbf{Example (EN)}:\\
A: It's delicious. B: But it doesn't look very good.

\textbf{Notes (JA)}:\\
「でも」は、対照的な考えを導入したり、反論を表現したりするためのカジュアルな表現です。「but」や「however」と訳され、カジュアルな会話で方向転換や反対意見を示すためによく使われます。このフレーズは、意見を述べたり提案をしたりする日常的な状況でよく使われます。よりフォーマルな場面では「しかし」と言うことが一般的です。2音節ですから、「でも」の方が言いやすいですね。

\textbf{Notes (EN)}:\\
"Demo" is a casual expression used to introduce a contrasting idea or to express disagreement. It can be translated as "but" or "however," and it is often used in informal conversations to indicate a change in direction or to present an opposing viewpoint. The phrase is commonly used in everyday situations, such as when discussing opinions or making suggestions. In more formal situations, one might use "shikashi" instead, which is a more polite way to express the same idea. Additionally, "demo" is often used in a light-hearted or humorous context, making it a versatile expression in casual conversations. It is typically two syllables long, making it easier to say than the more formal alternatives.

\newpage

\section*{376. I can do it}

\textbf{Date}: 2025.04.21 (Mon)\\
\textbf{Index}: ikeru, ikeru\\
\textbf{Expression (JA)}: (行:い)ける、(行:い)ける\\
\textbf{Romanization}: Ikeru, ikeru\\
\textbf{Expression (EN)}: I can do it

\textbf{Variants (JA)}: できる / いける / やれる / 可能だ\\
\textbf{Variants (EN)}: Go for it / I'm good / I can do it / I can go / I can make it / It's possible

\textbf{Intent (JA)}: 可能, 行動\\
\textbf{Intent (EN)}: possibility, action

\textbf{Tags (JA)}: 即時文法, 省略, 状況依存型, 可能性, 行動, リアルタイム性\\
\textbf{Tags (EN)}: immediate grammar, omission, context-dependent, possibility, action, strong real-time nature

\textbf{Adjusted Expression (JA)}:\\
(行:い)けると(思:おも)います。\\
\textbf{Adjusted Expression (EN)}:\\
I think I can do it.

\textbf{Example (JA)}:\\
A「(行:い)ける？」B「(行:い)ける、(行:い)ける」\\
\textbf{Example (EN)}:\\
A: Can you do it? B: I can do it, I can do it.

\textbf{Notes (JA)}:\\
「行ける、行ける」は、何かをすることができるという考えを伝えたり、励ましの気持ちを表現したりするカジュアルな表現です。「行ける」は「I can do it」や「Go for it」と訳され、カジュアルな会話で他人の能力に自信を持っていることを表すためによく使われます。このフレーズは、友人や同僚に挑戦するように励ます日常的な状況でよく使われます。よりフォーマルな場面では「できます」と言うことが一般的ですが、そういうときには、具体的に何ができるかを述べた方がいいですね。他にも、「行ける」と言うときは、「大丈夫」とか「おいしくていくつでも食べられる」などの意味でも使われます。

\textbf{Notes (EN)}:\\
"Ikeru, ikeru" is a casual expression used to convey the idea of being able to do something or to express encouragement. It can be translated as "I can do it" or "Go for it," and is often used in informal conversations to show confidence in someone's ability. This phrase is common in everyday situations, such as when encouraging a friend or colleague to take on a challenge. In more formal situations, one might use "dekimasu" instead, which is a more polite way to express the same idea, typically with a clearer subject or object. Additionally, depending on context, "ikeru" can also mean "It's okay" or even "It's so good I could eat a lot of it."

\newpage

\section*{375. Without thinking ...}

\textbf{Date}: 2025.04.20 (Sun)\\
\textbf{Index}: omowazu,..\\
\textbf{Expression (JA)}: (思:おも)わず、...\\
\textbf{Romanization}: Omowazu, ...\\
\textbf{Expression (EN)}: Without thinking ...

\textbf{Variants (JA)}: 無意識に / つい / はからずも / 反射的に / 直感的に / 無意識のうちに / 反射的に / 意に反して / 考えずに / 咄嗟に\\
\textbf{Variants (EN)}: instinctively / unconsciously / reflexively / involuntarily

\textbf{Intent (JA)}: 行動, 反応, 意識, 心理, 感情\\
\textbf{Intent (EN)}: action, reaction, awareness, consciousness, emotion, psychology

\textbf{Tags (JA)}: 即時文法, 省略, 状況依存型, 行動, 反応, 意識, リアルタイム性\\
\textbf{Tags (EN)}: immediate grammar, omission, context-dependent, action, reaction, awareness, strong real-time nature

\textbf{Adjusted Expression (JA)}:\\
(配慮:はいりょ)も(十分:じゅうぶん)にせずに、(発言:はつげん)してしまいました。\\
\textbf{Adjusted Expression (EN)}:\\
I spoke without thinking about the consideration.

\textbf{Example (JA)}:\\
A「(思:おも)わず、(本当:ほんとう)のこと、(言:い)っちゃった」\\
\textbf{Example (EN)}:\\
A: I said the truth without thinking.

\textbf{Notes (JA)}:\\
「思わず」は、意図せずに、咄嗟に反応してしまうことを表す表現です。「考えるより先に行動が出る」ような瞬間に使われ、英語では「without thinking」や「instinctively」と訳されます。

たとえば、「思わず本当のことを言ってしまった」という場合、本人にとっては後から「言わなければよかった」と思うような場面かもしれません。周囲の人が口にしない本音を代弁する形になったのでしょう。皆も内心では同じことを思っていたかもしれません。

ですが、「思わず言ってしまった」ことは、気にしても仕方ありません。こうした即時的な反応は、誰にでもあることです。

\textbf{Notes (EN)}:\\
"Omowazu" refers to an action or reaction that occurs instinctively or unintentionally, often before any conscious thought. It conveys a sense of doing something before realizing it, and is typically translated as "without thinking" or "instinctively."

For example, in the phrase "I said the truth without thinking," the speaker might later regret having spoken. Perhaps they voiced what others were thinking but dared not say aloud. In that case, while it might have been better to stay silent, it’s already done.

Since such instinctive reactions are common and natural, there's no point in worrying too much. It happens to everyone.

\newpage

\section*{374. Not yet ...}

\textbf{Date}: 2025.04.19 (Sat)\\
\textbf{Index}: ieiemadamada,..\\
\textbf{Expression (JA)}: いえいえ、まだまだ、...\\
\textbf{Romanization}: Ie ie, mada mada, ...\\
\textbf{Expression (EN)}: Not yet ...

\textbf{Variants (JA)}: いえいえ、そんな、... / とんでもないです、... / まだまだですよ、... / いやいや、全然です、... / いえ、私なんてまだまだ、... / 恐縮です、... / とてもそんな、... / まだまだ未熟で、... / とても及びません、...\\
\textbf{Variants (EN)}: Oh no, not at all... / I'm still learning, really... / Not even close... / I’m flattered, but no... / Far from it... / I’m not there yet... / I have a long way to go... / You're too kind... / I'm still just a beginner...

\textbf{Intent (JA)}: 謙遜, 否定\\
\textbf{Intent (EN)}: humility, negation

\textbf{Tags (JA)}: 即時文法, 省略, 状況依存型, 謙遜, 否定, 繰り返し, リアルタイム性\\
\textbf{Tags (EN)}: immediate grammar, omission, context-dependent, humility, negation, repetition, strong real-time nature

\textbf{Adjusted Expression (JA)}:\\
いえいえ、そんなことはないですよ。\\
\textbf{Adjusted Expression (EN)}:\\
Oh no, that's not true at all.

\textbf{Example (JA)}:\\
A「(ピアノ:ぴあの)、お(上手:じょうず)ですね」B「いえいえ、まだまだ...」\\
\textbf{Example (EN)}:\\
A: You're good at piano. B: Oh no, not at all...

\textbf{Notes (JA)}:\\
「いえいえ、まだまだ」は、褒められたときに謙遜や控えめな気持ちを表現するカジュアルな表現です。「いえ」は「no」を意味し、「まだ」は「not yet」や「still」を意味します。このフレーズは、カジュアルな会話で自分の能力や業績を控えめに表現するためによく使われます。文脈によって冗談として使われることもあります。よりフォーマルな場面では「いえ、そんなことはないです」と言うことが一般的です。だいたい、本当に上手であったとしても「いえいえ、まだまだ」と謙遜するのが日本人の美徳です。しかし、図々しい人はどこにでもいるもので、「いえいえ、まだまだ」と言いながら、そのあとすぐ、すごいテクニックを見せ付ける人がいます。

\textbf{Notes (EN)}:\\
"Ie ie, mada mada" is a casual expression used to express humility or modesty in response to a compliment. It combines the word "ie" (meaning "no") with "mada" (meaning "not yet" or "still"). This phrase is often used in informal conversations to downplay one's abilities or achievements. The phrase can also be used in a joking manner, depending on the context. In more formal situations, one might use "ie, sonna koto wa nai desu" instead, which is a more polite way to express the same idea. In Japan, it is considered a virtue to be humble, and even if someone is truly skilled, they might say "ie ie, mada mada" to express modesty. However, there are always some people who, while saying "ie ie, mada mada," immediately show off their amazing skills.

\newpage

\section*{373. Well, ...}

\textbf{Date}: 2025.04.18 (Fri)\\
\textbf{Index}: sate,..\\
\textbf{Expression (JA)}: さて、...\\
\textbf{Romanization}: Sate, ...\\
\textbf{Expression (EN)}: Well, ...

\textbf{Variants (JA)}: さて、... / それでは、... / じゃあ、... / では、...\\
\textbf{Variants (EN)}: Well, ... / Then, ... / So, ... / Now, ...

\textbf{Intent (JA)}: 提案, 疑問, 談話標識\\
\textbf{Intent (EN)}: suggestion, question, discourse marker

\textbf{Tags (JA)}: 即時文法, 省略, 状況依存型, 新規発話, リアルタイム性, 話題転換\\
\textbf{Tags (EN)}: immediate grammar, omission, context-dependent, new utterance, strong real-time nature, topic change, transition

\textbf{Adjusted Expression (JA)}:\\
それでは、つぎです。\\
\textbf{Adjusted Expression (EN)}:\\
Well, let's move on to the next one.

\textbf{Example (JA)}:\\
A「さて、まずはじめに、(問題:もんだい)です」\\
\textbf{Example (EN)}:\\
A: Well, first of all, here's a question.

\textbf{Notes (JA)}:\\
「さて」は、新しい話題に移るときや、何かを紹介するときに使われるカジュアルな表現です。「well」や「now」と訳され、カジュアルな会話で焦点を変えたり、次に来ることに備えてリスナーを準備させたりするためによく使われます。このフレーズは、プレゼンテーションを始めたり、新しいテーマを紹介したりする日常的な状況でよく使われます。フォーマルな場面では「それでは」と言うことが多いかもしれません。

\textbf{Notes (EN)}:\\
"Sate" is a casual expression used to transition to a new topic or to introduce something. It can be translated as "well" or "now," and it is often used in informal conversations to signal a change in focus or to prepare the listener for what is coming next. The phrase is commonly used in everyday situations, such as when starting a presentation or introducing a new subject. In more formal situations, one might use "soredewa" instead, which is a more polite way to express the same idea.

\newpage

\section*{372. That was thoughtful!}

\textbf{Date}: 2025.04.17 (Thu)\\
\textbf{Index}: kigakiteru\\
\textbf{Expression (JA)}: (気:き)が(利:き)いてる！\\
\textbf{Romanization}: Ki ga kiiteru!\\
\textbf{Expression (EN)}: That was thoughtful!

\textbf{Variants (JA)}: よく(気:き)がつく / よく(気:き)が(利:き)く / (気:き)がまわる / (気配:くば)りがある\\
\textbf{Variants (EN)}: How thoughtful! / So considerate! / You're so attentive. / That was really nice of you.

\textbf{Intent (JA)}: 称賛, 感謝, 評価\\
\textbf{Intent (EN)}: praise, appreciation, evaluation

\textbf{Tags (JA)}: 即時文法, 評価, 会話, 省略, 文末\\
\textbf{Tags (EN)}: immediate grammar, evaluation, conversation, ellipsis, sentence-final

\textbf{Adjusted Expression (JA)}:\\
(気:き)が(利:き)いていますね。\\
\textbf{Adjusted Expression (EN)}:\\
That's very thoughtful of you.

\textbf{Example (JA)}:\\
A「あ、(コーヒー:こーひー)に(チョコレート:ちょこれーと)がついてる！」B「ちょっと、(気:き)が(利:き)いてるね」\\
\textbf{Example (EN)}:\\
A: A piece of chocolate is on the side of the coffee! B: Oh, that was thoughtful!

\textbf{Notes (JA)}:\\
「気が利いてる」は、誰かの思いやりや配慮を称賛するためのカジュアルな表現です。「気」は「spirit」や「mind」を意味し、「利いてる」は「to be attentive」や「to be considerate」を意味します。このフレーズは、他人を思いやる行動や振る舞いだけでなく、ちょっとした小さなサービスに感心したときに、思わず発する表現です。でも、この表現は「気が利いている」人に直接言わず、聞こえるように言った方が効果的ですよ。人間関係ってそんなもんじゃないですか？

\textbf{Notes (EN)}:\\
"Kigakiteru" is a casual expression used to praise someone for their thoughtfulness or consideration. It combines the word "ki" (meaning "spirit" or "mind") with "kiteru" (meaning "to be attentive" or "to be considerate"). This phrase is often used in informal settings to express appreciation for someone's actions or behavior that show care and attention to others. It can also be used when someone does something small but thoughtful. But it's more effective to say this expression in a way that the person can hear it rather than directly to her. After all, isn't that how human relationships work?

\newpage

\section*{371. In the blink of an eye}

\textbf{Date}: 2025.04.16 (Wed)\\
\textbf{Index}: attoiumani\\
\textbf{Expression (JA)}: (あっ:あっ)という(間:ま)に\\
\textbf{Romanization}: A tto iu mani\\
\textbf{Expression (EN)}: In the blink of an eye

\textbf{Variants (JA)}: たちまち / (一瞬:いっしゅん)のうちに / (瞬:またた)く(間:ま)に / (即座:そくざ)に\\
\textbf{Variants (EN)}: In an instant / In a moment / In a flash / Immediately

\textbf{Intent (JA)}: 時間, 即時性\\
\textbf{Intent (EN)}: time, immediacy

\textbf{Tags (JA)}: 即時文法, 省略, 状況依存型, 時間, 即時性\\
\textbf{Tags (EN)}: immediate grammar, omission, context-dependent, time, immediacy

\textbf{Adjusted Expression (JA)}:\\
あっという(間:ま)に、(終:お)わってしまいました。\\
\textbf{Adjusted Expression (EN)}:\\
It ended in the blink of an eye.

\textbf{Example (JA)}:\\
A「おいしかった？」B「(あっ:あっ)という(間:ま)に、(食:た)べちゃった」\\
\textbf{Example (EN)}:\\
A: Was it delicious? B: In the blink of an eye, I ate it.

\textbf{Notes (JA)}:\\
「あっという間に」は、非常に早く、または予期せぬ出来事を表現するカジュアルな表現です。「あっ」は人が言語を1つ発する音、すなわち、とても短い時間を意味し、「いう間」は「a moment to say ...」を意味します。このフレーズは、日常的な状況で、突然起こった出来事や経験について話すときに使われます。文脈によって冗談として使われることもあります。

\textbf{Notes (EN)}:\\
"Atto iu mani" is a casual expression used to describe something that happens very quickly or unexpectedly. The word "atto" means the sound of a person saying something, which implies a very short time, and "iu mani" means "a moment to say ..." This phrase is often used in everyday situations to talk about events or experiences that occurred suddenly. Depending on the context, it can also be used jokingly.

\newpage

\section*{370. It's up to you}

\textbf{Date}: 2025.04.15 (Tue)\\
\textbf{Index}: jibunshidai\\
\textbf{Expression (JA)}: (自分:じぶん)(次第:しだい)\\
\textbf{Romanization}: Jibun shidai\\
\textbf{Expression (EN)}: It's up to you

\textbf{Variants (JA)}: (自分:じぶん)(次第:しだい) / あなた(次第:しだい) / 君(次第:しだい)\\
\textbf{Variants (EN)}: It's up to you / It's your choice / It's your decision / You decide

\textbf{Intent (JA)}: 選択, 決定, 責任\\
\textbf{Intent (EN)}: choice, decision, responsibility

\textbf{Tags (JA)}: 即時文法, 省略, 状況依存型, 選択, 責任, リアルタイム性\\
\textbf{Tags (EN)}: immediate grammar, omission, context-dependent, choice, responsibility, strong real-time nature

\textbf{Adjusted Expression (JA)}:\\
それは、(本人:ほんにん)の(気持:きも)ち(次第:しだい)ですね。\\
\textbf{Adjusted Expression (EN)}:\\
It's up to the person involved.

\textbf{Example (JA)}:\\
それをやるかどうかは、(自分:じぶん)(次第:しだい)。\\
\textbf{Example (EN)}:\\
Whether to do it or not is up to you.

\textbf{Notes (JA)}:\\
「自分次第」は、決定や選択が個人に委ねられていることを示すカジュアルな表現です。「自分」は「oneself」を、「次第」は「depending on」や「based on」を意味します。このフレーズは、カジュアルな会話で、話し手が他の人に決定を委ねていることを表現するためによく使われます。また、選択に対する責任感や所有感を暗示することもあります。また、結果や成功が自分自身の努力や意欲にかかっていることを暗に示す使い方もされます。

\textbf{Notes (EN)}:\\
"Jibun shidai" is a casual expression used to indicate that the decision or choice is up to the individual. It combines the word "jibun" (meaning "oneself") with "shidai" (meaning "depending on" or "based on"). This phrase is often used in informal conversations to express that the speaker is leaving the decision to someone else. It can also imply a sense of responsibility or ownership over the choice being made. It is sometimes used to imply that the outcome or success also depends on one's own efforts or motivation.

\newpage

\section*{369. I mean ...}

\textbf{Date}: 2025.04.14 (Mon)\\
\textbf{Index}: teiuka\\
\textbf{Expression (JA)}: て(言:い)うか、(チャンス:ちゃんす)がないっていうか、...\\
\textbf{Romanization}: Te iu ka, chansu ga nai tte iu ka, ...\\
\textbf{Expression (EN)}: I mean, ...

\textbf{Variants (JA)}: むしろ、... / じゃなくて、... / というよりも、...\\
\textbf{Variants (EN)}: Rather, ... / Not that, ... / More than that, ...

\textbf{Intent (JA)}: 説明, 強調, 言い換え\\
\textbf{Intent (EN)}: explanation, emphasis, rephrasing

\textbf{Tags (JA)}: 即時文法, 省略, 状況依存型, 説明, リアルタイム性\\
\textbf{Tags (EN)}: immediate grammar, omission, context-dependent, explanation, strong real-time quality

\textbf{Adjusted Expression (JA)}:\\
と言いますか、チャンスがないと申しますか、...\\
\textbf{Adjusted Expression (EN)}:\\
Or rather, I should say, there just aren't many opportunities...

\textbf{Example (JA)}:\\
A「いい(仕事:しごと)がないの？」B「て(言:い)うか、そもそも(チャンス:ちゃんす)(自体:じたい)がないっていうか、...」\\
\textbf{Example (EN)}:\\
A: Is there no good job? B: I mean, there are no chances at all, ...

\textbf{Notes (JA)}:\\
「て言うか」は、相手の表現が核心から外れていて、それを修正し、自分の言いたいことを表現するフレーズです。「言う」は「to say」を意味し、「か」は疑問や説明のニュアンスを加える助詞です。このフレーズは、友人や家族などのカジュアルな場面で使われることが多く、追加の文脈を提供したり、ポイントを強調したりするために使われます。よりフォーマルな場面では「と言いますか、...」と言うことが一般的です。

関連語句として、「そもそも」は話の前提に立ち返る際に用いられる語で、現在の話題の前提そのものを問い直す働きがあります。「自体」は主語や主題の中の焦点を強調する語で、特に否定文において「◯◯ですらない」ような語感を伴います。これらはいずれも即時的な焦点調整に関わる語句であり、会話の中で素早く文脈を調整するのに役立ちます。

\textbf{Notes (EN)}:\\
"Te iu ka" is a casual expression to guide the conversation or to correct the previous statement. It can be translated as "I mean" or "rather," and it is often used in informal conversations to clarify or emphasize a point. The phrase is commonly used in everyday situations, such as when discussing plans or making guesses. In more formal situations, one might use "to iimasu ka" instead, which is a more polite way to express the same idea.

Related expressions include "somosomo" (meaning "in the first place") which allows the speaker to question or reset the premise of the current conversation. Another is "jitai" (meaning "itself"), which is used to highlight the very essence of a subject, often in negative contexts like "There’s no chance itself." Both phrases function as immediate, real-time focus operators that support impromptu corrections or emphasis in conversation.

\newpage

\section*{368. Surely}

\textbf{Date}: 2025.04.13 (Sun)\\
\textbf{Index}: kitto\\
\textbf{Expression (JA)}: きっと、(正解:せいかい)\\
\textbf{Romanization}: Kitto, seikai\\
\textbf{Expression (EN)}: Surely, it's the answer.

\textbf{Variants (JA)}: きっと / たぶん / おそらく / 確かに\\
\textbf{Variants (EN)}: Surely / Probably / Perhaps / Certainly

\textbf{Intent (JA)}: 確信, 予測\\
\textbf{Intent (EN)}: certainty, prediction

\textbf{Tags (JA)}: 即時文法, 省略, 状況依存型, 確信, リアルタイム性\\
\textbf{Tags (EN)}: immediate grammar, omission, context-dependent, certainty, strong real-time nature

\textbf{Adjusted Expression (JA)}:\\
きっと、正解だと思います。\\
\textbf{Adjusted Expression (EN)}:\\
I think it's surely the answer.

\textbf{Example (JA)}:\\
A「(正解:せいかい)は、きっと、これだ」B「きっと、そうだね」\\
\textbf{Example (EN)}:\\
A: Surely, the answer is this. B: Surely, that's right.

\textbf{Notes (JA)}:\\
「きっと」は、何かに対する確信や強い信念を表現するカジュアルな表現です。「きっと」は「surely」や「probably」と訳され、カジュアルな会話で予測や仮定に自信を持っていることを表すためによく使われます。このフレーズは、計画を話し合ったり、推測したりする日常的な状況でよく使われます。「正解」は「answer」を意味しますが、問題の答という意味以外にも、「迷っていたけれどもそれを選んでよかった」と思うときに使えます。

\textbf{Notes (EN)}:\\
"Kitto" is a casual expression used to convey certainty or strong belief about something. It can be translated as "surely" or "probably," and it is often used in informal conversations to express confidence in a prediction or assumption. The phrase is commonly used in everyday situations, such as when discussing plans or making guesses. "Seikai" means "answer," but it can also be used to express the feeling of having made the right choice after being uncertain.

\newpage

\section*{367. Don't be silly}

\textbf{Date}: 2025.04.12 (Sat)\\
\textbf{Index}: nawakenaidesho\\
\textbf{Expression (JA)}: な(訳:わけ)ないでしょ\\
\textbf{Romanization}: Na wake nai desho\\
\textbf{Expression (EN)}: Don't be silly

\textbf{Variants (JA)}: まさか / ちょっと待ってよ / ありえない / どうして、そうなる？\\
\textbf{Variants (EN)}: Come on / No way / That's impossible / How could that happen?

\textbf{Intent (JA)}: 否定, 驚き\\
\textbf{Intent (EN)}: negation, surprise

\textbf{Tags (JA)}: 即時文法, 省略, 状況依存型, 否定, リアルタイム性\\
\textbf{Tags (EN)}: immediate grammar, omission, context-dependent, negation, strong real-time nature

\textbf{Adjusted Expression (JA)}:\\
そんなことはないと思います。\\
\textbf{Adjusted Expression (EN)}:\\
I don't think that's the case.

\textbf{Example (JA)}:\\
A「あの(人:ひと)、(実:じつ)は(宇宙人:うちゅうじん)なんじゃない？」B「な(訳:わけ)ないでしょ！(漫画:まんが)の(読:よ)みすぎ！」\\
\textbf{Example (EN)}:\\
A: That person is actually an alien, right? B: Don't be silly! You're reading too many comics!

\textbf{Notes (JA)}:\\
「な訳ないでしょ」は、何かを否定したり、信じられないことを表現するカジュアルな表現です。「な訳」は「that kind of thing」を意味し、「ない」は否定形、「でしょ」は強調や感情を加える助詞です。このフレーズは、友人や家族などのカジュアルな場面で使われることが多く、状況に対する驚きや疑念を表現します。文脈によっては冗談として使われることもあります。よりフォーマルな場面では「そんなことはないでしょう」と言うことが一般的です。

\textbf{Notes (EN)}:\\
"Nawakenaidesho" is a casual expression used to deny something or express disbelief. It combines the word "nawa" (meaning "that kind of thing") with the negative form "kenai" (meaning "not possible") and the particle "desho," which adds emphasis or emotion to the statement. This phrase is often used in informal settings, such as among friends or family, to express surprise or disbelief at a situation. The phrase can also be used in a joking manner, depending on the context. In more formal situations, one might use "sonna koto wa nai deshou" instead, which is a more polite way to express disbelief.

\newpage

\section*{366. ruine everything}

\textbf{Date}: 2025.04.11 (Fri)\\
\textbf{Index}: motomokomonai\\
\textbf{Expression (JA)}: (元:もと)も(子:こ)もない\\
\textbf{Romanization}: Motomokomonai\\
\textbf{Expression (EN)}: ruin everything

\textbf{Variants (JA)}: (台:だい)なし / すべてが(水:みず)の(泡:あわ)\\
\textbf{Variants (EN)}: ruin everything / nothing left / everything is gone

\textbf{Intent (JA)}: 否定\\
\textbf{Intent (EN)}: negation

\textbf{Tags (JA)}: 即時文法, 省略, 状況依存型, 否定, リアルタイム性\\
\textbf{Tags (EN)}: immediate grammar, omission, context-dependent, negation, strong real-time nature

\textbf{Adjusted Expression (JA)}:\\
(台:だい)なしにしないで\\
\textbf{Adjusted Expression (EN)}:\\
Don't ruin everything.

\textbf{Example (JA)}:\\
A「(無理:むり)しちゃいけない。(体:からだ)、(壊:こわ)したら(元:もと)も(子:こ)もない」\\
\textbf{Example (EN)}:\\
A: Don't overdo it. If you break your body, there will be nothing left.

\textbf{Notes (JA)}:\\
「元も子もない」は、何かがうまくいかないときに、元の状態や大切なものが失われてしまうことを表現するフレーズです。「元」は「origin」や「source」を、「子」は「child」や「offspring」を意味します。このフレーズは、カジュアルな会話で誰かに無理をしないように警告したり、自分自身を大切にするように促すためによく使われます。文脈によっては冗談として使われることもあります。

\textbf{Notes (EN)}:\\
"Motomokomonai" is a phrase used to express the idea that if something goes wrong, there will be nothing left or everything will be ruined. "Moto" means "origin" or "source," and "ko" means "child" or "offspring." This phrase is often used in casual conversations to warn someone not to overdo it or to take care of themselves. It can also be used in a joking manner, depending on the context. 

\newpage

\section*{365. What now?}

\textbf{Date}: 2025.04.10 (Thu)\\
\textbf{Index}: dousurebaii\\
\textbf{Expression (JA)}: どうすればいい？\\
\textbf{Romanization}: Dousureba ii?\\
\textbf{Expression (EN)}: What should I do?

\textbf{Variants (JA)}: どうすればいい / どうしよう / どうしたらいい / どしたらいい / どうする\\
\textbf{Variants (EN)}: What should I do? / What can I do? / What do I do? / What am I supposed to do? / What now?

\textbf{Intent (JA)}: 提案, 疑問\\
\textbf{Intent (EN)}: suggestion, question

\textbf{Tags (JA)}: 即時文法, 省略, 状況依存型, 疑問文, リアルタイム性\\
\textbf{Tags (EN)}: immediate grammar, omission, context-dependent, question, strong real-time nature

\textbf{Adjusted Expression (JA)}:\\
どうすればいいのか、(考:かんが)えています。\\
\textbf{Adjusted Expression (EN)}:\\
I'm thinking about what to do.

\textbf{Example (JA)}:\\
A「ああ、もう(時間:じかん)がない、どうすればいい？」B「ちょっと、(落:お)ちついて」\\
\textbf{Example (EN)}:\\
A: Oh no, there's no time left. What should I do? B: Calm down a bit.

\textbf{Notes (JA)}:\\
「どうすればいい」は、困ったときにアドバイスや提案を求めるカジュアルな表現です。「する」の条件形「すれば」と、感情を強調するための口語的な助詞「いい」を組み合わせたものです。このフレーズは、友人や家族などのカジュアルな場面でよく使われ、緊急性や切迫感を伝えることができます。文脈によっては冗談として使われることもあります。よりフォーマルな場面では「どうしようか」と言うことが一般的です。「落ちついて」も、カジュアルな表現で、相手に冷静になるよう促す言葉です。「落ち着いて」は「calm down」を意味し、友人や家族との会話でよく使われます。もちろん、これらのことばは自分自身に対して、独り言として使うことができます。

\textbf{Notes (EN)}:\\
"Dousurebaii" is a casual expression used to ask for advice or suggestions when in a difficult situation. It combines the verb "suru" (to do) with the conditional form "ba" and the informal particle "ii," which adds emphasis or emotion to the request. This phrase is often used in informal settings, such as among friends or family, and can convey a sense of urgency or desperation. The phrase can also be used in a joking manner, depending on the context. In more formal situations, one might use "dou shiyou ka" instead, which is a more polite way to ask for advice. "Ochitsuite" is also a casual expression used to encourage someone to calm down. It means "calm down" and is often used in conversations with friends or family. Of course, these phrases can also be used as self-talk.

\newpage

\section*{364. Cut it out}

\textbf{Date}: 2025.04.09 (Wed)\\
\textbf{Index}: yameteyo\\
\textbf{Expression (JA)}: ちょっと、やめてよ\\
\textbf{Romanization}: Chotto, yamete yo\\
\textbf{Expression (EN)}: Oh, no, cut it out.

\textbf{Variants (JA)}: やめてよ / やめて / やめてください\\
\textbf{Variants (EN)}: Cut it out / Stop it / Please stop

\textbf{Intent (JA)}: 拒否, 感情, 照れ, 冗談\\
\textbf{Intent (EN)}: refusal, emotion, embarrassment, joking

\textbf{Tags (JA)}: 即時文法, 感情表現, 状況依存型, リアルタイム性, 呼応\\
\textbf{Tags (EN)}: immediate grammar, emotional expression, context-dependent, real-time nature, responsive

\textbf{Adjusted Expression (JA)}:\\
そういうことはやめてください。\\
\textbf{Adjusted Expression (EN)}:\\
Please don't say such things.

\textbf{Example (JA)}:\\
A「(今日:きょう)は(髪型:かみがた)、すてきだね」B「ちょっと、やめてよ」\\
\textbf{Example (EN)}:\\
A: Your hairstyle looks great today. B: Oh, no, cut it out.

\textbf{Notes (JA)}:\\
「やめてよ」は、誰かに何かをやめてほしいと頼むカジュアルな表現です。「やめる」の連用形「やめて」と、感情を強調するための口語的な助詞「よ」を組み合わせたものです。このフレーズは、友人や家族などのカジュアルな場面でよく使われ、照れくささや遊び心を伝えることができます。文脈によっては冗談として使われることもあります。よりフォーマルな場面では「やめてください」と言うことが一般的で、本当に止めてほしいときに使います。

\textbf{Notes (EN)}:\\
"Yamete yo" is a casual expression used to ask someone to stop doing something. It combines the verb "yameru" (to stop) in the te-form "yamete" with the informal particle "yo," which adds emphasis or emotion to the request. This phrase is often used in informal settings, such as among friends or family, and can convey a sense of embarrassment or playfulness. The phrase can also be used in a joking manner, depending on the context. In more formal situations, one might use "yamete kudasai" instead, which is a more polite way to make a sincere request for someone to stop.

\newpage

\section*{363. Thank you for doing ...}

\textbf{Date}: 2025.04.08 (Tue)\\
\textbf{Index}: kurete\\
\textbf{Expression (JA)}: (来:き)てくれてありがとう\\
\textbf{Romanization}: Kite kurete arigatou\\
\textbf{Expression (EN)}: Thank you for coming

\textbf{Variants (JA)}: \\
\textbf{Variants (EN)}: 

\textbf{Intent (JA)}: 感謝, 招待\\
\textbf{Intent (EN)}: gratitude, invitation

\textbf{Tags (JA)}: 即時文法, 省略, 状況依存型, 感謝, リアルタイム性\\
\textbf{Tags (EN)}: immediate grammar, omission, context-dependent, gratitude, strong real-time nature

\textbf{Adjusted Expression (JA)}:\\
(来:き)てくださり、ありがとうございます。\\
\textbf{Adjusted Expression (EN)}:\\
Thank you for coming.

\textbf{Example (JA)}:\\
A「(来:き)てくれてありがとう」B「いえいえ、呼んでくれてありがとう」\\
\textbf{Example (EN)}:\\
A: Thank you for coming. B: No, thank you for inviting me.

\textbf{Notes (JA)}:\\
「来てくれてありがとう」は、誰かの存在や参加に感謝するためのカジュアルな表現です。「来る」の連用形「来て」と「くれて」を組み合わせたもので、話し手の利益のために行われた行動を示します。このフレーズは、友人や家族などのカジュアルな場面で、イベントや集まりに来てくれたことへの感謝を表すためによく使われます。「～てくれてありがとう」あるいはその丁寧「～てくださり、ありがとうございます」は、相手の行動が話し手にとって嬉しいことを伝える文型です。友達、家族以外には「～てくださり、ありがとうございます」と言うのが一般的です。

\textbf{Notes (EN)}:\\
"Kite kurete arigatou" is a casual expression used to express gratitude for someone's presence or participation. It combines the verb "kuru" (to come) in the te-form "kite" with "kurete," which indicates that the action was done for the speaker's benefit. This phrase is commonly used in casual settings, such as among friends or family, to express appreciation for someone's attendance at an event or gathering. The phrase "\\textasciitilde{}te kurete arigatou" or its more polite form "\\textasciitilde{}te kudasari, arigatou gozaimasu" conveys that the speaker is pleased with the action taken by the other person. It is generally used in formal situations with people other than friends or family.

\newpage

\section*{362. Basically}

\textbf{Date}: 2025.04.07 (Mon)\\
\textbf{Index}: kihontekiniha\\
\textbf{Expression (JA)}: (基本的:きほんてき)には\\
\textbf{Romanization}: Kihonteki ni wa\\
\textbf{Expression (EN)}: Basically

\textbf{Variants (JA)}: \\
\textbf{Variants (EN)}: 

\textbf{Intent (JA)}: 基本, 基礎, 評価, 枠づけ\\
\textbf{Intent (EN)}: basic, foundation, evaluation, framing

\textbf{Tags (JA)}: 即時文法, 省略, 状況依存型, 構造外前置き句, リアルタイム性, 評価的発話\\
\textbf{Tags (EN)}: immediate grammar, omission, context-dependent, extrasyntactic premodifier, strong real-time nature, evaluative utterance

\textbf{Adjusted Expression (JA)}:\\
(基本的:きほんてき)には、(賛成:さんせい)です。\\
\textbf{Adjusted Expression (EN)}:\\
Basically, I agree.

\textbf{Example (JA)}:\\
A「(基本的:きほんてき)には、(賛成:さんせい)です」B「ということは、何か問題でも？」\\
\textbf{Example (EN)}:\\
A: Basically, I agree. B: So, is there a problem?

\textbf{Notes (JA)}:\\
[\\{"title":"「基本的には」","note":"「基本的には」は、基本的な視点や立場を表現するためのフレーズです。「基本的」は「basic」や「fundamental」を意味し、「には」は話し手の立場や視点を示す文法構造です。このフレーズは、評価や意見を枠付けるためによく使われ、後続の文に対する前提を提供します。構造外前置き句として機能し、主文の枠組みを整える役割を果たします。「基本的には」と発話することから、何か問題があることを示唆しています。もしかしたら、反対していないというニュアンスを持つかもしれません。このような表現は、即時文法において非常に重要な役割を果たします。「構造外前置き句」という名称は広く定着した専門用語ではありませんが、即時文法の特徴を記述するうえで極めて有効な概念です。"\\},\\{"title":"構造外前置き句について","note":"「構造外前置き句」とは、文の冒頭に置かれ、話し手の立場や話の枠組みを示すものの、後続の構文には影響を与えない句や語句を指す記述的な名称です。これらは談話の流れを整えたり、話し手の態度を示したりする役割を持ちますが、文の主構造（主語・述語・修飾関係）には含まれません。"\\},\\{"title":"features","note":"文の冒頭に置かれる; 構文には組み込まれない（構造外）; 意味的にも主文とは独立して発話される; 談話の枠組み、話し手の立場、心構えを示す"\\},\\{"title":"examples","note":"一般的には、そういう傾向がある。『一般的には』が構造外前置き句として機能し、全体の前提を緩やかに提示している。"\\},\\{"title":"正直に言うと、それはちょっと違うと思う。","note":"『正直に言うと』が話し手の態度を表す挿入句。後続の文の構文には影響していない。"\\},\\{"title":"まあ、それはそれとして、次の話題に移ろう。","note":"『まあ、それはそれとして』は話題転換のマーカーであり、即時的な談話のコントロールを行っている。"\\},\\{"title":"基本的には、賛成です。","note":"『基本的には』が主文の判断に対する枠付けを行っている。"\\},\\{"title":"『構造外前置き句』","note":"『構造外前置き句』という名称は広く定着した専門用語ではありませんが、即時文法の特徴を記述するうえで極めて有効な概念であり、言語横断的にも観察される現象です。"\\}]

\textbf{Notes (EN)}:\\
"Kihonteki ni wa" is a phrase used to express a basic or foundational perspective. "Kihonteki" means "basic" or "fundamental," and "ni wa" is a grammatical structure that indicates the speaker's stance or viewpoint. This phrase is often used to frame an evaluation or opinion, providing context for the following statement. It serves as an extrasyntactic premodifier, setting the stage for the main clause. By saying "kihonteki ni wa," the speaker suggests that there may be some issues or nuances to consider. It can imply a nuance of not being entirely opposed to something, but rather providing a framework for discussion. The term "extrasyntactic premodifier" is not widely established in linguistic terminology, but it effectively describes the function of this phrase in immediate grammar.

\newpage

\section*{361. Take a detour}

\textbf{Date}: 2025.04.06 (Sun)\\
\textbf{Index}: yorimichishite\\
\textbf{Expression (JA)}: (寄:よ)り(道:みち)して\\
\textbf{Romanization}: Yorimichi shite\\
\textbf{Expression (EN)}: Take a detour

\textbf{Variants (JA)}: \\
\textbf{Variants (EN)}: 

\textbf{Intent (JA)}: 提案, 誘い\\
\textbf{Intent (EN)}: suggestion, invitation

\textbf{Tags (JA)}: 即時文法, 省略, 状況依存型, 提案, リアルタイム性\\
\textbf{Tags (EN)}: immediate grammar, omission, context-dependent, suggestion, strong real-time nature

\textbf{Adjusted Expression (JA)}:\\
(寄:よ)り(道:みち)して、(帰:かえ)りましょう。\\
\textbf{Adjusted Expression (EN)}:\\
Let's take a detour and go home.

\textbf{Example (JA)}:\\
A「(寄:よ)り(道:みち)して、(帰:かえ)ろ」B「いいね」\\
\textbf{Example (EN)}:\\
A: Let's take a detour and go home. B: Sounds good!

\textbf{Notes (JA)}:\\
「寄り道して帰ろう」は、帰る途中で寄り道をすることを提案する表現です。「寄り道」は「detour」を、「して」は動詞「する」（to do）の連用形、「帰ろ」は動詞「帰る」の意向形です。このフレーズは、カジュアルな会話で使われることが多く、計画の変更や旅を楽しむための提案として使われます。意向形を使うことで、提案や誘いのニュアンスが強まります。斎藤和義の「歩いて帰ろう」などの歌詞にも使われている表現です。

\textbf{Notes (EN)}:\\
"Yorimichi shite kaero" is a casual phrase used to suggest taking a detour on the way home. "Yorimichi" means "detour," "shite" is the te-form of the verb "suru" ("to do"), and "kaerou" is the volitional form of "kaeru" ("to go home"). This expression is commonly used in everyday conversation to propose a spontaneous change of plans or to make the journey more enjoyable. The volitional form adds a nuance of suggestion or gentle invitation. It's a familiar phrase even in pop culture, for example, in the lyrics of Kazuyoshi Saito's song "Aruite Kaerou."

\newpage

\section*{360. Let's go home}

\textbf{Date}: 2025.04.05 (Sat)\\
\textbf{Index}: kaero\\
\textbf{Expression (JA)}: (帰:かえ)ろ\\
\textbf{Romanization}: Kaero\\
\textbf{Expression (EN)}: Let's go home

\textbf{Variants (JA)}: \\
\textbf{Variants (EN)}: 

\textbf{Intent (JA)}: 帰宅, 提案\\
\textbf{Intent (EN)}: going home, suggestion

\textbf{Tags (JA)}: 即時文法, 省略, 状況依存型, 提案, リアルタイム性\\
\textbf{Tags (EN)}: immediate grammar, omission, context-dependent, suggestion, strong real-time nature

\textbf{Adjusted Expression (JA)}:\\
(帰:かえ)りましょう\\
\textbf{Adjusted Expression (EN)}:\\
Shall we go home?

\textbf{Example (JA)}:\\
A「いっしょに(帰:かえ)ろ」B「うん、(帰:かえ)ろう」\\
\textbf{Example (EN)}:\\
A: Let's go home together. B: Yeah, let's go.

\textbf{Notes (JA)}:\\
「帰ろ」は、帰ることを提案するカジュアルな表現です。「帰ろ」は動詞「帰る」の意向形で、正式には「帰ろう」ですが、カジュアルな会話では「帰ろ」と省略されることが多いです。このフレーズは、友達や家族との会話でよく使われます。意向形を使うことで、より直接的でカジュアルな印象を与えます。

\textbf{Notes (EN)}:\\
"Kaero" is a casual way to suggest going home. It is the volitional form of the verb "kaeru" (to go home), and in formal situations, it would be "kaerou." However, in casual conversations, it is often shortened to "kaero." This phrase is commonly used in conversations with friends or family. By using the volitional form, it gives a more direct and casual impression.

\newpage

\section*{359. Take it easy}

\textbf{Date}: 2025.04.04 (Fri)\\
\textbf{Index}: nonbiri\\
\textbf{Expression (JA)}: (今度:こんど)、のんびり、(コーヒー:こーひー)でもどう？\\
\textbf{Romanization}: Kondo, nonbiri, koohii demo dou?\\
\textbf{Expression (EN)}: How about taking it easy with a coffee next time?

\textbf{Variants (JA)}: \\
\textbf{Variants (EN)}: 

\textbf{Intent (JA)}: 提案, 誘い\\
\textbf{Intent (EN)}: suggestion, invitation

\textbf{Tags (JA)}: 即時文法, 省略, 状況依存型, 提案, リアルタイム性\\
\textbf{Tags (EN)}: immediate grammar, omission, context-dependent, suggestion, strong real-time nature

\textbf{Adjusted Expression (JA)}:\\
(今度:こんど)、のんびり、(コーヒー:こーひー)でも(飲:の)みませんか？\\
\textbf{Adjusted Expression (EN)}:\\
Would you like to have a relaxing coffee together sometime?

\textbf{Example (JA)}:\\
A「(今度:こんど)、のんびり、(コーヒー:こーひー)でもどう？」B「いいね」\\
\textbf{Example (EN)}:\\
A: How about taking it easy with a coffee next time? B: Sounds good!

\textbf{Notes (JA)}:\\
「のんびり」は、リラックスしたり、ゆったりとした状態を表現するためのフレーズです。「のんびり」は「take it easy」や「relax」を意味し、カジュアルな会話で使われることが多いです。このフレーズは、誰かに休憩を取ったり、一緒にゆったりとした時間を楽しむことを提案するときによく使われます。文末の「どう」は、提案や招待を示す表現です。本当に親しい人から、このように言われたときには、相手の気持ちを考えて、あまり強く断らない方が良いでしょう。たとえば、「今度、のんびり、コーヒーでもどう？」と聞かれたときに、「今度は無理」とだけ言うのは少し失礼です。相手の気持ちを考えて、「また誘ってね」と添えると、やわらかく断ることができます。また、こうした誘いは社交辞令であることも多く、必ずしも実行を前提にしていない場合もあります。そんなときは、「じゃ、いつにする？」などと詰めずに、「ええ、ぜひ」と返しておくのがスマートです。

\textbf{Notes (EN)}:\\
"Nonbiri" is a phrase used to express a relaxed or leisurely state. It means "take it easy" or "relax," and it is often used in casual conversations to suggest a laid-back attitude. The use of the word "dou" at the end of the sentence indicates that the speaker is making a suggestion or invitation. When a close friend says this, it is better not to refuse too strongly, considering the other person's feelings. For example, if someone asks, "How about taking it easy with a coffee next time?" saying "I can't this time" might be a bit rude. Instead, it would be better to say something like, "I can't this time, but please invite me again." Also, just because you're invited like this doesn't mean you have to go. The other person may not really want to go together. This is what we call 'shakoujirei' in Japanese, what might be called 'an empty promise disguised as a courteous remark' in English. In such cases, instead of asking, "So when should we do it?" just respond with "Sure, let's do it!"

\newpage

\section*{358. Is that so?}

\textbf{Date}: 2025.04.03 (Thu)\\
\textbf{Index}: sounano\\
\textbf{Expression (JA)}: そうなの？\\
\textbf{Romanization}: Sou nano?\\
\textbf{Expression (EN)}: Is that so?

\textbf{Variants (JA)}: \\
\textbf{Variants (EN)}: 

\textbf{Intent (JA)}: 確認, 疑問\\
\textbf{Intent (EN)}: confirmation, question

\textbf{Tags (JA)}: 即時文法, 省略, 状況依存型, 疑問文, リアルタイム性\\
\textbf{Tags (EN)}: immediate grammar, omission, Context-dependent expression, Question, Strong real-time nature

\textbf{Adjusted Expression (JA)}:\\
そうなんですか？(知:し)りませんでした。\\
\textbf{Adjusted Expression (EN)}:\\
Is that so? I didn't know.

\textbf{Example (JA)}:\\
A「(明日:あす)、(試験:しけん)だよね」B「そうなの？じゃないと(思:おも)うけど」\\
\textbf{Example (EN)}:\\
A: The exam is tomorrow, right? B: Is that so? I don't think so.

\textbf{Notes (JA)}:\\
「そうなの？」は、驚きや確認を表現するためのフレーズです。「そう」は「like that」や「is that so」を意味し、「なの」は疑問の助詞「の」の口語的な形です。このフレーズは、カジュアルな会話で情報を確認したり、信じられないことを表現するためによく使われます。疑問の助詞「の」を使うことで、問いかけや好奇心を感じさせる効果があります。「じゃないと思うけど」は、軽く否定する表現で、相手の発言を受けて「そうではない」と思うことを伝えています。このように、「じゃない」の前のことばを省略して短く表現することで、母語話者のような自然な発話になります。

\textbf{Notes (EN)}:\\
"Sou nano?" is a phrase used to express surprise or confirmation. "Sou" means "like that" or "is that so," and "nano" is a colloquial form of the question particle "no." This phrase is often used in casual conversations to confirm information or express disbelief. The use of the question particle "no" adds a sense of inquiry or curiosity. The phrase "ja nai to omou kedo" is a light denial expression, indicating that the speaker thinks it is not the case. This allows for a shorter expression by omitting the word before "ja nai," making it sound more natural and fluent, as native speakers often do.

\newpage

\section*{357. Finally, it's done}

\textbf{Date}: 2025.04.02 (Wed)\\
\textbf{Index}: yatto dayo\\
\textbf{Expression (JA)}: やっと、だよ\\
\textbf{Romanization}: Yatto dayo\\
\textbf{Expression (EN)}: Finally, it's done

\textbf{Variants (JA)}: \\
\textbf{Variants (EN)}: 

\textbf{Intent (JA)}: 達成, 完了\\
\textbf{Intent (EN)}: achievement, completion

\textbf{Tags (JA)}: 即時文法, 省略, 状況依存型, 感情表現, リアルタイム性\\
\textbf{Tags (EN)}: immediate grammar, omission, Context-dependent expression, Emotional expression, Strong real-time nature

\textbf{Adjusted Expression (JA)}:\\
やっと、(宿題:しゅくだい)が(終:お)わりました。\\
\textbf{Adjusted Expression (EN)}:\\
Finally, the homework is done.

\textbf{Example (JA)}:\\
A「(宿題:しゅくだい)、もう(終:お)わったの？」B「もう、じゃなくて、やっと、だよ」\\
\textbf{Example (EN)}:\\
A: You done with your homework already? B: Not already, finally!

\textbf{Notes (JA)}:\\
「やっとだよ」は、達成感や物事が完了したことによる安堵感を伝える表現です。「やっと」は「finally」、「だ」は「〜です」「〜だ」といった断定の表現、「よ」は相手にとって新規の情報を表す終助詞です。このフレーズは、長い時間や努力を要してようやく何かを終えたときによく使われます。「やっと」によって、話し手の安心感や満足感が強調されます。

\textbf{Notes (EN)}:\\
"Yatto da yo" is a phrase used to express a sense of achievement or relief upon completion. "Yatto" means "finally," "da" is the plain form of the verb "to be" (like "is"), and "yo" is a sentence-final particle that indicates the speaker is presenting new or important information to the listener. This expression is commonly used when something has taken a long time or much effort to complete. The use of "yatto" highlights the speaker's feeling of relief or satisfaction.

\newpage

\section*{356. Doing this and that}

\textbf{Date}: 2025.04.01 (Tue)\\
\textbf{Index}: taritari\\
\textbf{Expression (JA)}: (寝:ね)たり、(食:た)べたり\\
\textbf{Romanization}: Neta ri, tabe tari\\
\textbf{Expression (EN)}: Sleeping, eating, and so on

\textbf{Variants (JA)}: \\
\textbf{Variants (EN)}: 

\textbf{Intent (JA)}: 列挙, 網羅\\
\textbf{Intent (EN)}: enumeration, comprehensive

\textbf{Tags (JA)}: 即時文法, 省略, 状況依存型, 列挙, リアルタイム性\\
\textbf{Tags (EN)}: immediate grammar, omission, Context-dependent expression, Enumeration, Strong real-time nature

\textbf{Adjusted Expression (JA)}:\\
(寝:ね)たり、(食:た)べたりします。\\
\textbf{Adjusted Expression (EN)}:\\
I sleep, eat, and so on.

\textbf{Example (JA)}:\\
A「(休:やす)みの(日:ひ)は、(何:なに)してるの？」B「(寝:ね)たり、(食:た)べたりしてる」\\
\textbf{Example (EN)}:\\
A: What do you do on your days off? B: I sleep, eat, and so on.

\textbf{Notes (JA)}:\\
「V1-たり、V2-たりします」は、複数の行動や活動を列挙するためのフレーズです。V1の動詞が1つの行動を示し、V2の動詞が別の行動を示します。このフレーズは、網羅的でない一連の行動や活動を表現するためによく使われます。「...たり」の文形を作るためには、たとえば、動詞「寝る」の過去形「寝た」を「...たり」と組み合わせます。この表現は、いくつかの代表的な行動をとりあげて、それらのようなことをしていることを示します。すべての行動を網羅するのではなく、例として挙げた行動から、それ以外の同種の行動も行っていることをほのめかすような用法です。

\textbf{Notes (EN)}:\\
"Neta ri, tabe tari" is a phrase used to list multiple actions or activities. The verb in the te-form is repeated with "tari" to indicate that there are other similar actions or activities. This phrase is often used to express a non-exhaustive list of actions or activities. To form the "...tari" sentence structure, you combine the verb in the past tense with "tari." For example, you combine the past tense of the verb "neru" (to sleep) with "tari." This expression suggests that you are doing various representative actions and implies that you are doing other similar actions as well, without listing all of them.

\newpage

\section*{355. While doing something}

\textbf{Date}: 2025.03.31 (Mon)\\
\textbf{Index}: nagara\\
\textbf{Expression (JA)}: (音楽:おんがく)、(聞:き)きながら...\\
\textbf{Romanization}: ...nagara\\
\textbf{Expression (EN)}: While listening to music...

\textbf{Variants (JA)}: \\
\textbf{Variants (EN)}: 

\textbf{Intent (JA)}: 同時性, 並行\\
\textbf{Intent (EN)}: simultaneity, concurrency

\textbf{Tags (JA)}: 即時文法, 省略, 状況依存型, 同時性, 並行, リアルタイム性\\
\textbf{Tags (EN)}: immediate grammar, omission, Context-dependent expression, Simultaneity, Concurrency, Strong real-time nature

\textbf{Adjusted Expression (JA)}:\\
(音楽:おんがく)を(聞:き)ながら...\\
\textbf{Adjusted Expression (EN)}:\\
While listening to music...

\textbf{Example (JA)}:\\
A「(音楽:おんがく)、(聞:き)ながら、(勉強:べんきょう)してる」B「いいね」\\
\textbf{Example (EN)}:\\
A: I'm studying while listening to music.

\textbf{Notes (JA)}:\\
「V1-しながら、V2」は、2つの行動が同時に行われていることを表すフレーズです。V1の動詞が随伴的・背景的な行動を示し、V2の動詞が主要な行動を示します。このフレーズは、複数の行動や並行して行われている活動を表現するためによく使われます。「...しながら」の文形を作るためには、たとえば、動詞「聞く」の連用形「聞き」を「...しながら」と組み合わせます。

\textbf{Notes (EN)}:\\
"V1 nagara, V2" is a phrase used to express that two actions are happening simultaneously. The verb V1 indicates the accompanying or background action, and the verb V2 indicates the main action. This phrase is often used to express multiple actions or activities happening concurrently. To form the "...nagara" sentence structure, you combine the verb in the te-form with "nagara." For example, you combine the te-form of the verb "kiku" (to listen) with "nagara."

\newpage

\section*{354. Anything will be fine}

\textbf{Date}: 2025.03.30 (Sun)\\
\textbf{Index}: nandemo\\
\textbf{Expression (JA)}: なんでもいいよ\\
\textbf{Romanization}: Nandemoiiyo\\
\textbf{Expression (EN)}: Anything would be fine

\textbf{Variants (JA)}: \\
\textbf{Variants (EN)}: 

\textbf{Intent (JA)}: 選択, 柔軟\\
\textbf{Intent (EN)}: choice, flexibility

\textbf{Tags (JA)}: 即時文法, 省略, 状況依存型, 許可, リアルタイム性\\
\textbf{Tags (EN)}: immediate grammar, omission, Context-dependent expression, Permission, Strong real-time nature

\textbf{Adjusted Expression (JA)}:\\
どれでもかまいませんよ。\\
\textbf{Adjusted Expression (EN)}:\\
Any of them will be fine.

\textbf{Example (JA)}:\\
A「(何:なに)がいい？」B「なんでもいいよ」\\
\textbf{Example (EN)}:\\
A: What would you like? B: Anything is fine.

\textbf{Notes (JA)}:\\
「なんでもいいよ」は、何でも受け入れられることを表現するフレーズです。「なんでも」は「anything」を、「いい」は形容詞「いい」（good）を、「よ」は同意を求める文末の助詞です。このフレーズは、誰かが好みを尋ねたときに、柔軟性を示したり、許可を与えたりするときによく使われます。"なんでもいいよ"を使うことで、どの選択肢にも対応できることを伝えることができます。しかし、実際には、どの選択肢も受け入れられるわけではない場合もありますので、文脈によって使い分けることが重要です。特に危険なのは、ガールフレンドやボーイフレンドに対して「なんでもいいよ」と言った場合です。あなたが「なんでもいいよ」と言ったとき、相手は「本当になんでもいいの？」と聞いてくるかもしれません。あるいは、あなたが「なんでもいいよ」と言ったとき、相手はあなたが自分に興味がないと感じるかもしれません。これはもしかしたら、最悪の事態をひき起こすかもしれません。

\textbf{Notes (EN)}:\\
"Nandemo iiyo" is a phrase used to express that you are okay with anything. "Nandemo" means "anything," "ii" is the adjective "ii" (good), and "yo" is a sentence-final particle that seeks agreement. This phrase is often used to show flexibility or give permission when someone asks for your preference. By saying "nandemo iiyo," you can convey that you are open to any option. However, in reality, not all options may be acceptable, so it is important to use this phrase appropriately depending on the context. Particularly dangerous is when you say "nandemo iiyo" to your girlfriend or boyfriend. When you say "nandemo iiyo," the other person may ask, "Do you really mean anything?" Or when you say "nandemo iiyo," the other person may feel that you are not interested in them. This could lead to the worst-case scenario.

\newpage

\section*{353. Isn't this amazing?}

\textbf{Date}: 2025.03.29 (Sat)\\
\textbf{Index}: sugokunai?kore\\
\textbf{Expression (JA)}: (凄:すご)くない？これ。\\
\textbf{Romanization}: Sugoku nai? Kore.\\
\textbf{Expression (EN)}: Isn't this amazing?

\textbf{Variants (JA)}: \\
\textbf{Variants (EN)}: 

\textbf{Intent (JA)}: 感動, 驚き\\
\textbf{Intent (EN)}: impressed, surprised

\textbf{Tags (JA)}: 即時文法, 省略, 状況依存型, 感情表現, 疑問文, リアルタイム性\\
\textbf{Tags (EN)}: immediate grammar, omission, Context-dependent expression, Emotional expression, Question, Strong real-time nature

\textbf{Adjusted Expression (JA)}:\\
これは、すごいですね。\\
\textbf{Adjusted Expression (EN)}:\\
This is amazing, isn't it?

\textbf{Example (JA)}:\\
A「(世界新記録:せかいしんきろく)、(凄:すご)くない？これ」\\
\textbf{Example (EN)}:\\
A: It's a world record. Isn't this amazing?

\textbf{Notes (JA)}:\\
「すごくない？これ」は、驚きや感動を表現するための表現です。「すごく」は形容詞「すごい」の副詞形で「amazing」や「incredible」を、「ない」は動詞「ある」（to be）の否定形、「これ」は「this」を意味します。この表現は、「ない」を使っていますが、否定しているわけではありません。疑問文で使われる特徴を表わすことばの否定形は、否定を表すだけでなく、驚きや感動を表すときにも使われます。

\textbf{Notes (EN)}:\\
"Sugoku nai? Kore" is an expression used to express amazement or surprise. "Sugoku" is the adverbial form of the adjective "sugoi" (amazing), "nai" is the negative form of the verb "aru" (to be), and "kore" means "this." Although the phrase uses the negative form, it is not denying anything. The negative form of question words used in questions has the function of expressing surprise or amazement.

\newpage

\section*{352. I'm tired}

\textbf{Date}: 2025.03.28 (Fri)\\
\textbf{Index}: tsukareta,a-\\
\textbf{Expression (JA)}: あー、(疲:つか)れた\\
\textbf{Romanization}: Aa, tsukareta\\
\textbf{Expression (EN)}: I'm tired

\textbf{Variants (JA)}: \\
\textbf{Variants (EN)}: 

\textbf{Intent (JA)}: 疲労, 不満\\
\textbf{Intent (EN)}: fatigue, frustration

\textbf{Tags (JA)}: 即時文法, 省略, 状況依存型, 感情表現, 感動詞使用, リアルタイム性\\
\textbf{Tags (EN)}: immediate grammar, omission, Context-dependent expression, Emotional expression, interjectional usage, Strong real-time nature

\textbf{Adjusted Expression (JA)}:\\
ちょっと、(疲:つか)れてしまいました。\\
\textbf{Adjusted Expression (EN)}:\\
Ah, I'm tired.

\textbf{Example (JA)}:\\
A「(今日:きょう)は(忙:いそが)しかったね」B「あー、(疲:つか)れた」\\
\textbf{Example (EN)}:\\
A: You had a busy day today. B: I'm tired.

\textbf{Notes (JA)}:\\
「あー、疲れた」は、疲れたことを表現するためのフレーズです。「あー」は感情を表す単語の前に使用する感動詞で、「疲れた」は動詞「疲れる」（to get tired）の過去形です。このフレーズは、長い一日の終わりや激しい活動の後によく使われます。主語が省略されていますが、文脈から意味が明確に理解されます。

\textbf{Notes (EN)}:\\
"Aa, tsukareta" is a phrase used to express that you are tired. "Aa" is an interjection expressing emotions, and "tsukareta" is the past tense of the verb "tsukareru" (to get tired). This phrase is often used at the end of a long day or after intense activity. Although the subject is omitted, the meaning is clearly understood from the context.

\newpage

\section*{351. How's it going?}

\textbf{Date}: 2025.03.27 (Thu)\\
\textbf{Index}: saikin-dou\\
\textbf{Expression (JA)}: (最近:さいきん)、どう？\\
\textbf{Romanization}: Saikin, dou?\\
\textbf{Expression (EN)}: How's it going?

\textbf{Variants (JA)}: \\
\textbf{Variants (EN)}: 

\textbf{Intent (JA)}: 挨拶, 近況確認\\
\textbf{Intent (EN)}: greeting, check-in

\textbf{Tags (JA)}: 即時文法, 省略, 状況依存型, 疑問文, リアルタイム性\\
\textbf{Tags (EN)}: immediate grammar, omission, Context-dependent expression, Question, Strong real-time nature

\textbf{Adjusted Expression (JA)}:\\
最近、どうですか？\\
\textbf{Adjusted Expression (EN)}:\\
How have you been lately?

\textbf{Example (JA)}:\\
A「(最近:さいきん)、どう？」B「(元気:げんき)だよ、ありがとう」\\
\textbf{Example (EN)}:\\
A: How's it going? B: I'm good, thanks.

\textbf{Notes (JA)}:\\
「最近、どう？」は、誰かの近況を尋ねるカジュアルな表現です。「最近」は「lately」や「recently」を、「どう」は「how」を意味します。このフレーズは友達や家族とのカジュアルな会話でよく使われます。即時文法の典型的な例であり、動詞が省略されていますが、文脈から意味が明確に理解されます。

\textbf{Notes (EN)}:\\
"Saikin, dou?" is a casual way to ask how someone is doing. "Saikin" means "lately" or "recently," and "dou" is a question word that means "how." This phrase is often used in casual conversations among friends or family. It is a typical example of immediate grammar, as the verb is omitted but the meaning is clearly understood from context.

\newpage

\section*{350. Endure in silence}

\textbf{Date}: 2025.03.26 (Wed)\\
\textbf{Index}: nakineiridesune\\
\textbf{Expression (JA)}: それは、(結局:けっきょく)のところ、(泣:な)き(寝入:ねい)りですね\\
\textbf{Romanization}: Sore wa, kekkyoku no tokoro, naki-neiri desu ne\\
\textbf{Expression (EN)}: So in the end, you just have to endure it in silence, huh?

\textbf{Variants (JA)}: \\
\textbf{Variants (EN)}: 

\textbf{Intent (JA)}: \\
\textbf{Intent (EN)}: 

\textbf{Tags (JA)}: 即時文法, 状況によっては調整文法, 評価表現\\
\textbf{Tags (EN)}: immediate grammar, contextually adjustive, evaluation expression

\textbf{Adjusted Expression (JA)}:\\
それは、(:)結局、どこにも(:)訴えることはできません。\\
\textbf{Adjusted Expression (EN)}:\\
In the end, there was no legal recourse available.

\textbf{Example (JA)}:\\
A「(銀行:ぎんこう)(振:ふ)り(込:こ)みで(商品:しょうひん)が(送:おく)ってこなくて、(相手:あいて)の(住所:じゅうしょ)もわからないんです」B「それは、(結局:けっきょく)のところ、(泣:な)き(寝入:ねい)りですね」\\
\textbf{Example (EN)}:\\
A: I transferred the money to their account, but the product never came, and I don't even know where they are. B: Well... sounds like you're just going to have to suffer in silence.

\textbf{Notes (JA)}:\\
「泣き寝入り」は、黙って耐えるしかない状況を表す表現です。「泣き」は動詞「泣く」のマス形（「泣きます」の「ます」を取り除いた形で、連用形とも言う）、「寝」は「sleep」、「入り」は「enter」、「ですね」は「isn't it」を意味します。その人が文句を言わず、助けを求めずに状況に耐えるしかないことを暗示しています。

\textbf{Notes (EN)}:\\
"Naki-ne-iri desu ne" is an expression used to describe a situation where someone has no choice but to endure silently. "Naki" is the continuative form of the verb "naku" (to cry), "ne" is the noun form of "neru" (to sleep), "iri" comes from "iru" (to enter), and "desu ne" is a sentence-final particle that seeks agreement. The phrase implies resignation or powerlessness, often used when someone is left without recourse.

\newpage

\section*{349. Oh, it's ...ing!}

\textbf{Date}: 2025.03.25 (Tue)\\
\textbf{Index}: ...teru,a!\\
\textbf{Expression (JA)}: あ！(車:くるま)、(来:き)てる！\\
\textbf{Romanization}: A! Kuruma, kiteru!\\
\textbf{Expression (EN)}: Watch out! The car is coming!

\textbf{Variants (JA)}: \\
\textbf{Variants (EN)}: 

\textbf{Intent (JA)}: \\
\textbf{Intent (EN)}: 

\textbf{Tags (JA)}: 即時文法, 省略, 状況依存型, 警告, 縮約形, リアルタイム性\\
\textbf{Tags (EN)}: immediate grammar, Many omissions, Context-dependent expression, Warning, Contracted expressions, Strong real-time nature

\textbf{Adjusted Expression (JA)}:\\
あ！(車:くるま)が(来:き)ています！\\
\textbf{Adjusted Expression (EN)}:\\
Oh! The car is approaching.

\textbf{Example (JA)}:\\
A「あ！(車:くるま)、(来:き)てる！」B「(危:あぶ)ない！」\\
\textbf{Example (EN)}:\\
A: Watch out! The car is coming! B: Be careful!

\textbf{Notes (JA)}:\\
「あ！車、来てる！」は、危険をすぐに知らせるための表現です。「あ」は驚きや警告を表す間投詞で、「車」は「car」、「来てる」は「来ている」の縮約形です。このフレーズは、車が近づいてきたり、誰かが何かにぶつかりそうな状況など、すぐに誰かに警告する必要がある場面でよく使われます。動詞「来ている」を「来てる」に縮約することで、表現がより簡潔になり、リアルタイムのコミュニケーションに適しています。

\textbf{Notes (EN)}:\\
"A! Kuruma, kiteru!" is an expression used to warn someone immediately of danger. "A" is an interjection expressing surprise or warning, "kuruma" means "car," "kiteru" is the contracted form of "kite iru" (is coming). This phrase is used when you need to warn someone immediately, for example, when a car is approaching or someone is about to bump into something. By contracting the verb "kite iru" to "kiteru," the expression becomes more concise and suitable for real-time communication.

\newpage

\section*{348. When are you coming back?}

\textbf{Date}: 2025.03.24 (Mon)\\
\textbf{Index}: itsumodottekuruno\\
\textbf{Expression (JA)}: (何時:いつ)、(戻:もど)ってくるの\\
\textbf{Romanization}: Itsu, modotte kuru no\\
\textbf{Expression (EN)}: When are you coming back?

\textbf{Variants (JA)}: \\
\textbf{Variants (EN)}: 

\textbf{Intent (JA)}: \\
\textbf{Intent (EN)}: 

\textbf{Tags (JA)}: 即時文法, 状況依存型, 疑問「の」, リアルタイム性\\
\textbf{Tags (EN)}: immediate grammar, Context-dependent expression, Question particle "no", Strong real-time nature

\textbf{Adjusted Expression (JA)}:\\
いつ、(戻:もど)って(来:く)るのですか。\\
\textbf{Adjusted Expression (EN)}:\\
When will you be returning?

\textbf{Example (JA)}:\\
A「(何時:いつ)、(戻:もど)ってくるの？」B「(夕方:ゆうがた)には(戻:もど)るよ」\\
\textbf{Example (EN)}:\\
A: When are you coming back? B: I'll be back in the evening.

\textbf{Notes (JA)}:\\
「いつ、戻って来るの」は、いつ戻るか尋ねるためのカジュアルな表現です。「いつ」は「when」、「戻って」は動詞「戻る」（to return）のテ形、「来る」は動詞「来る」（to come）の現在形、「の」は疑問を示す助詞です。疑問を示す助詞「の」は、カジュアルな会話で情報を求めたり、理由を尋ねる際にも使われます。

\textbf{Notes (EN)}:\\
"Itsu, modotte kuru no" is a casual way to ask when someone will return. "Itsu" means "when," "modotte" is the te-form of the verb "modoru" (to return), "kuru" is the present form of the verb "kuru" (to come), and "no" is a question particle. The particle "no" adds an inquisitive tone, often used in casual speech to seek information or ask for reasons.

\newpage

\section*{347. I told you so}

\textbf{Date}: 2025.03.23 (Sun)\\
\textbf{Index}: dakaraittajan\\
\textbf{Expression (JA)}: だから(言:い)ったじゃん\\
\textbf{Romanization}: Dakara ittajan\\
\textbf{Expression (EN)}: I told you so

\textbf{Variants (JA)}: \\
\textbf{Variants (EN)}: 

\textbf{Intent (JA)}: \\
\textbf{Intent (EN)}: 

\textbf{Tags (JA)}: 即時文法, 省略, 状況依存型, くりかえし表現, 縮約表現, リアルタイム性\\
\textbf{Tags (EN)}: immediate grammar, Many omissions, Context-dependent expression, Repetitive expressions, Contracted expressions, Strong real-time nature

\textbf{Adjusted Expression (JA)}:\\
ですから、(私:わたし)はそのように(言:い)いました。\\
\textbf{Adjusted Expression (EN)}:\\
That is why I mentioned it before.

\textbf{Example (JA)}:\\
A「この道は通れないね」B「だから、(言:い)ったじゃん、(行:い)き(止:ど)まりだって」A「ごめんごめん」\\
\textbf{Example (EN)}:\\
A: We can't go through this road. B: I told you so. It's a dead end. A: Sorry, sorry.

\textbf{Notes (JA)}:\\
「だから言ったじゃん」は、以前に言ったことが正しかったことを表現するフレーズです。「だから」は「therefore」や「so」を、「言った」は動詞「言う」（to say）の過去形、「じゃん」は「です」の口語的な形です。このフレーズは、友達や家族とのカジュアルな会話でよく使われます。即時文法の典型的な例であり、話し手が以前に何を言ったのかは省略されていますが、文脈から意味が明確に理解されます。

\textbf{Notes (EN)}:\\
"Dakara ittajan" is a phrase used to express that you were right about something you said earlier. The word "dakara" means "therefore" or "so," "itta" is the past tense of the verb "iu" (to say), and "jan" is a colloquial contraction of "desu" (to be). This phrase confirms that what you said earlier was correct and is commonly used in casual conversations among friends or family. It is a typical example of immediate grammar, as the specific details of what was said earlier are omitted, but the meaning is clearly understood from context.

\newpage

\section*{346. Let me see.}

\textbf{Date}: 2025.03.22 (Sat)\\
\textbf{Index}: doredore\\
\textbf{Expression (JA)}: どれどれ\\
\textbf{Romanization}: Dore dore\\
\textbf{Expression (EN)}: Let me see.

\textbf{Variants (JA)}: \\
\textbf{Variants (EN)}: 

\textbf{Intent (JA)}: \\
\textbf{Intent (EN)}: 

\textbf{Tags (JA)}: 即時文法, 省略, 状況依存型, くりかえし表現, リアルタイム性\\
\textbf{Tags (EN)}: immediate grammar, Many omissions, Context-dependent expression, Repetitive expressions, Strong real-time nature

\textbf{Adjusted Expression (JA)}:\\
どちらでしょうか。\\
\textbf{Adjusted Expression (EN)}:\\
Which one would it be?

\textbf{Example (JA)}:\\
A「これ、(見:み)て」B「どれどれ」\\
\textbf{Example (EN)}:\\
A: Look at this. B: Let me see.

\textbf{Notes (JA)}:\\
「どれどれ」は、誰かが何かを見せているときに、それを確認しようとする即時的な反応を示す表現です。「どれ」は「which one」を意味し、これを繰り返すことで、より自然なリアルタイムの反応になります。この表現では動詞が省略されていますが、文脈によって「見せて」「確認しよう」という意図が明確に伝わります。例えば、子供が何かを見せようとしているときに親しみを持って使うこともできますが、大人同士の会話でも「どれどれ？」と言って何かを確認するときに使えます。

\textbf{Notes (EN)}:\\
"Dore dore" is an expression used when you want to take a look at something that someone is showing or referring to. The phrase consists of "dore" (which one) repeated twice, creating a more natural and immediate reaction. This phrase is often used when you approach a child looking at toys, but it can also be used in casual conversations when someone is showing you something and you want to take a closer look. It is a typical example of immediate grammar, as the verb is omitted but the meaning is clearly understood from context.

\newpage

\section*{345. Look at this}

\textbf{Date}: 2025.03.21 (Fri)\\
\textbf{Index}: mitemite, kore\\
\textbf{Expression (JA)}: (見:み)て(見:み)て、これ\\
\textbf{Romanization}: Mite mite, kore\\
\textbf{Expression (EN)}: Look at this

\textbf{Variants (JA)}: \\
\textbf{Variants (EN)}: 

\textbf{Intent (JA)}: \\
\textbf{Intent (EN)}: 

\textbf{Tags (JA)}: \\
\textbf{Tags (EN)}: 

\textbf{Adjusted Expression (JA)}:\\
\\
\textbf{Adjusted Expression (EN)}:\\


\textbf{Example (JA)}:\\
A「(見:み)て(見:み)て、これ」B「(何:なに)？」\\
\textbf{Example (EN)}:\\
A: Look at this. B: What?

\textbf{Notes (JA)}:\\
「見て見て、これ」は、誰かの注意を何かに向けさせるための表現です。「見て」は動詞「見る」（to see）のテ形で、「これ」は「this」を意味します。この表現は、誰かに何かを見せたいときや、興味深いものを共有したいときに使われます。日本語はSOV（主語-目的語-動詞）言語なので、動詞を文末に置くのが基本ですが、カジュアルな会話では「見て！これ！」のように語順を変えることで、臨場感や驚きを表現できます。言いたい語を先に言うことで、強調したい部分を目立たせることができます。

\textbf{Notes (EN)}:\\
"Mite mite, kore" is an expression used to draw someone's attention to something. "Mite" is the te-form of the verb "miru" (to see), and "kore" means "this." This expression is used when you want to show someone something or share something interesting. Japanese is an SOV language, meaning that verbs are usually placed at the end in formal speech. However, in casual conversation, you can say "Mite! Kore!" to create a sense of immediacy or excitement. You can say the key word first, depending on what you want to emphasize.

\newpage

\section*{344. Where are you?}

\textbf{Date}: 2025.03.20 (Thu)\\
\textbf{Index}: dokoniiruno\\
\textbf{Expression (JA)}: どこにいるの\\
\textbf{Romanization}: Doko ni iru no\\
\textbf{Expression (EN)}: Where are you?

\textbf{Variants (JA)}: \\
\textbf{Variants (EN)}: 

\textbf{Intent (JA)}: \\
\textbf{Intent (EN)}: 

\textbf{Tags (JA)}: \\
\textbf{Tags (EN)}: 

\textbf{Adjusted Expression (JA)}:\\
\\
\textbf{Adjusted Expression (EN)}:\\


\textbf{Example (JA)}:\\
A「(今:いま)、どこにいるの？」B「(図書館:としょかん)にいるよ」\\
\textbf{Example (EN)}:\\
A: Where are you now? B: I'm at the library.

\textbf{Notes (JA)}:\\
「どこにいるの」は、誰かがどこにいるか尋ねるためのフレーズです。「どこ」は「where」を、「に」は場所を示す助詞を、「いる」は現在形の動詞「いる」（to be）を、「の」は確認を求める助詞を意味します。これは電話で相手の場所を尋ねるときや、誰かを探しているときに使われます。

\textbf{Notes (EN)}:\\
"Doko ni iru no" is a phrase used to ask where someone is. The word "doko" means "where," "ni" is a particle that indicates the location, "iru" is the verb "iru" (to be) in the present tense, and "no" is a particle that seeks confirmation. This phrase is used when asking someone's location over the phone or when looking for someone.

\newpage

\section*{343. would've done}

\textbf{Date}: 2025.03.19 (Wed)\\
\textbf{Index}: noni\\
\textbf{Expression (JA)}: ..V-taのに\\
\textbf{Romanization}: Noni\\
\textbf{Expression (EN)}: would have done

\textbf{Variants (JA)}: \\
\textbf{Variants (EN)}: 

\textbf{Intent (JA)}: \\
\textbf{Intent (EN)}: 

\textbf{Tags (JA)}: \\
\textbf{Tags (EN)}: 

\textbf{Adjusted Expression (JA)}:\\
\\
\textbf{Adjusted Expression (EN)}:\\


\textbf{Example (JA)}:\\
A「え、(歩:ある)いて来たの？」B「うん」A「(電話:でんわ)してくれれば、(迎:むか)えに(行:い)ったのに」\\
\textbf{Example (EN)}:\\
A: Huh? You walked here? B: Yeah. A: If you had called me, I would've picked you up.

\textbf{Notes (JA)}:\\
「のに」は、起こらなかったことに対する後悔や失望を表現するための接続詞です。「のに」は「although」や「despite」を意味します。「のに」の前の動詞は過去形で、後の動詞は過去形や仮定形です。「～たのに」は日本語では たった一語「のに」 を付けるだけで、仮定のニュアンスを簡単に表現できます。一方で、英語では "would have done" という複雑な文法構造が必要になります。この点を考えると、日本語のほうが直感的で使いやすいとも言えますね。

\textbf{Notes (EN)}:\\
"Noni" is a conjunction used to express regret or disappointment about something that did not happen. The word "noni" means "although" or "despite." The verb before "noni" is in the past tense, and the verb after it is in the past or conditional form. In Japanese, you can easily express a hypothetical situation with just one word, "noni." On the other hand, in English, you need a complex grammar structure like "would have done." Considering this point, Japanese is more intuitive and easier to use.

\newpage

\section*{342. It can't be helped}

\textbf{Date}: 2025.03.18 (Tue)\\
\textbf{Index}: shikataganai\\
\textbf{Expression (JA)}: (仕方:しかた)がない\\
\textbf{Romanization}: Shikata ga nai\\
\textbf{Expression (EN)}: No way for it

\textbf{Variants (JA)}: \\
\textbf{Variants (EN)}: 

\textbf{Intent (JA)}: \\
\textbf{Intent (EN)}: 

\textbf{Tags (JA)}: \\
\textbf{Tags (EN)}: 

\textbf{Adjusted Expression (JA)}:\\
\\
\textbf{Adjusted Expression (EN)}:\\


\textbf{Example (JA)}:\\
A「(雨:あめ)、(降:ふ)ってきた」B「(仕方:しかた)がないね」\\
\textbf{Example (EN)}:\\
A: It's raining. B: No way for it.

\textbf{Notes (JA)}:\\
「仕方がない」は、状況を変える方法がないことや、どうしようもないことを表現するためのフレーズです。「仕方」は「way」や「method」を、「が」は文の主語を示す助詞、「ない」は動詞「ある」（to be）の否定形です。このフレーズは、状況を受け入れるしかないことを表現するために使われます。コントロールできないことが起こったとき、例えば悪天候や公共交通機関の遅延など、よく使われます。

\textbf{Notes (EN)}:\\
"Shikata ga nai" is a phrase used to express that there is no way to change a situation or that something cannot be helped. The word "shikata" means "way" or "method," "ga" is a particle that indicates the subject of the sentence, and "nai" is the negative form of the verb "aru" (to be). This phrase is used to express that you have no choice but to accept a situation as it is. It is often used when something happens that is beyond your control, such as bad weather or a delay in public transportation.

\newpage

\section*{341. I thought that you wouldn't come.}

\textbf{Date}: 2025.03.17 (Mon)\\
\textbf{Index}: konainokatoomottayo\\
\textbf{Expression (JA)}: (来:こ)ないのかと(思:おも)ったよ\\
\textbf{Romanization}: konai no ka to omotta yo\\
\textbf{Expression (EN)}: I thought that you wouldn't come.

\textbf{Variants (JA)}: \\
\textbf{Variants (EN)}: 

\textbf{Intent (JA)}: \\
\textbf{Intent (EN)}: 

\textbf{Tags (JA)}: \\
\textbf{Tags (EN)}: 

\textbf{Adjusted Expression (JA)}:\\
\\
\textbf{Adjusted Expression (EN)}:\\


\textbf{Example (JA)}:\\
A「(遅:おそ)くなって、(ごめん)」B「(来:こ)ないのかと(思:おも)ったよ」\\
\textbf{Example (EN)}:\\
A: I'm sorry I'm late. B: I thought that you wouldn't come.

\textbf{Notes (JA)}:\\
「来ないのかと思ったよ」は、何かが起こらないと思っていたが、実際に起こったことを表現するためのフレーズです。「来ない」は動詞「来る」（to come）の否定形、「の」は文の主語を示す助詞、「か」は疑問を示す助詞、「と」は接続詞で「and」を意味し、「思ったよ」は動詞「思う」（to think）の過去形です。この例のように思っていたことと違う結果になったときの時間に関する表現は、言語によって異なることがあるので、確認しておきましょう。

\textbf{Notes (EN)}:\\
"Konai no ka to omottayo" is a phrase used to express that you thought something would not happen, but it did. The word "konai" is the negative form of the verb "kuru" (to come), "no" is a particle that indicates the subject of the sentence, "ka" is a question particle, "to" is a conjunction that means "and," and "omottayo" is the past tense of the verb "omou" (to think). Expressions of time when something happened differently from what you thought may vary depending on the language, so be sure to check.

\newpage

\section*{340. You know what?}

\textbf{Date}: 2025.03.16 (Sun)\\
\textbf{Index}: anone\\
\textbf{Expression (JA)}: あのね\\
\textbf{Romanization}: Ano ne\\
\textbf{Expression (EN)}: You know what?

\textbf{Variants (JA)}: \\
\textbf{Variants (EN)}: 

\textbf{Intent (JA)}: \\
\textbf{Intent (EN)}: 

\textbf{Tags (JA)}: \\
\textbf{Tags (EN)}: 

\textbf{Adjusted Expression (JA)}:\\
\\
\textbf{Adjusted Expression (EN)}:\\


\textbf{Example (JA)}:\\
A「あのね、ママ、(大好:だいす)きだよ」B「ありがとう」\\
\textbf{Example (EN)}:\\
A: You know what, Mom, I love you so much. B: Thank you.

\textbf{Notes (JA)}:\\
「あのね」は、会話を始めるときや、言いたいことを強調するためのフレーズです。「あの」は「that」や「um」を、「ね」は同意や理解を求める助詞です。「好き」は「love」や「like」を、「大好き」は「love so much」を意味します。このフレーズは、言いたいことに注意を引くために使われます。大切なことを共有したり、気持ちを表現したりしたいときによく使われます。難しい話をしたいときにも使われます。

\textbf{Notes (EN)}:\\
"Ano ne" is a phrase used to start a conversation or emphasize what you want to say. "Ano" is similar to "that" or "um," and "ne" is a particle that seeks agreement or understanding. "suki" means "love" or "like." "daisuki" means "love so much." This phrase is used to draw attention to what you want to say. It is often used when you want to share something important or express your feelings. It is also used when you want to talk about something difficult.

\newpage

\section*{339. I'll leave it to you}

\textbf{Date}: 2025.03.15 (Sat)\\
\textbf{Index}: omakaseshimasu\\
\textbf{Expression (JA)}: お(任:まか)せします\\
\textbf{Romanization}: Omakase shimasu\\
\textbf{Expression (EN)}: I'll leave it to you

\textbf{Variants (JA)}: \\
\textbf{Variants (EN)}: 

\textbf{Intent (JA)}: \\
\textbf{Intent (EN)}: 

\textbf{Tags (JA)}: \\
\textbf{Tags (EN)}: 

\textbf{Adjusted Expression (JA)}:\\
\\
\textbf{Adjusted Expression (EN)}:\\


\textbf{Example (JA)}:\\
A「どこで(晩:ばん)ご(飯:はん)、(食:た)べる？」B「お(任:まか)せします」\\
\textbf{Example (EN)}:\\
A: Where do you want to eat dinner? B: I'll leave it to you.

\textbf{Notes (JA)}:\\
「お任せします」は、決定を他の人に委ねるときに使うフレーズです。「お任せ」は「entrust」や「leave it to you」という意味があり、「します」は動詞「する」（to do）の丁寧形です。このフレーズは決まり文句として、友だちや家族との関係でも「お休みなさい」や「いただきます」と同じように丁寧な表現のまま使います。「任せる」と言うと、やや乱暴な印象や上から目線のニュアンスがあります。使うとすれば、例えばサッカーの試合でペナルティーキックを任せる場面など、緊張感のある状況で「おまえに任せる」と言うことが考えられます。また、「お任せ」には別の意味もあります。寿司屋で料理長にメニューの選択を任せることを指します。

\textbf{Notes (EN)}:\\
"Omakase shimasu" is a phrase used to leave a decision or task to someone else. The word "omakase" means "entrust" or "leave it to you," and "shimasu" is the polite form of the verb "suru" (to do). This phrase is a set expression and is used politely even among friends and family, similar to phrases like "good night" and "itadakimasu." Saying "makaseru" sounds a bit rough or condescending. It might be used in tense situations, such as during a sports game when a player says, "I'll leave it to you" before a penalty kick in soccer. "Omakase" also has another meaning: it refers to leaving the menu selection to the chef at a sushi restaurant.

\newpage

\section*{338. It looks interesting}

\textbf{Date}: 2025.03.14 (Fri)\\
\textbf{Index}: omoshirosou\\
\textbf{Expression (JA)}: おもしろそう\\
\textbf{Romanization}: Omoshirosou\\
\textbf{Expression (EN)}: It looks interesting

\textbf{Variants (JA)}: \\
\textbf{Variants (EN)}: 

\textbf{Intent (JA)}: \\
\textbf{Intent (EN)}: 

\textbf{Tags (JA)}: \\
\textbf{Tags (EN)}: 

\textbf{Adjusted Expression (JA)}:\\
\\
\textbf{Adjusted Expression (EN)}:\\


\textbf{Example (JA)}:\\
この(映画:えいが)、おもしろそう。\\
\textbf{Example (EN)}:\\
This movie looks interesting.

\textbf{Notes (JA)}:\\
「おもしろそう」は、何かがおもしろそうだということを表現するためのフレーズです。「おもしろい」は「interesting」や「fun」を、「そう」は何かがあるように見えることを示す接尾辞です。このフレーズは、最初の印象に基づいて何かに興味を持っていることを表現するために使われます。映画や本、イベントなど、初めて見るものに対してよく使われます。

\textbf{Notes (EN)}:\\
"Omoshirosou" is a phrase used to express that something looks interesting or fun. The word "omoshiroi" means "interesting" or "fun," and "sou" is a suffix that indicates that something looks a certain way. This phrase is used to express that you are interested in something based on your first impression. It is often used when you see something for the first time, such as a movie, book, or event.

\newpage

\section*{337. I can't believe it}

\textbf{Date}: 2025.03.13 (Thu)\\
\textbf{Index}: shinjirarenai\\
\textbf{Expression (JA)}: (信:しん)じられない\\
\textbf{Romanization}: Shinjirarenai\\
\textbf{Expression (EN)}: I can't believe it

\textbf{Variants (JA)}: \\
\textbf{Variants (EN)}: 

\textbf{Intent (JA)}: \\
\textbf{Intent (EN)}: 

\textbf{Tags (JA)}: \\
\textbf{Tags (EN)}: 

\textbf{Adjusted Expression (JA)}:\\
\\
\textbf{Adjusted Expression (EN)}:\\


\textbf{Example (JA)}:\\
A「お(待:ま)たせ」B「(信:しん)じられない。1(時間:じかん)も(遅:おく)れて」\\
\textbf{Example (EN)}:\\
A: Sorry to keep you waiting. B: I can't believe it. You're an hour late.

\textbf{Notes (JA)}:\\
「信じられない」は、不信や驚きを表現するためのフレーズです。「信じられない」は動詞「信じる」（to believe）の否定可能形です。否定可能形は、何かができないか、難しいことを表現するために使われます。このフレーズは、信じがたいことに驚いたりショックを受けたりすることを表現するために使われます。

\textbf{Notes (EN)}:\\
"Shinjirarenai" is a phrase used to express disbelief or surprise. The word "shinjirarenai" is the negative potential form of the verb "shinjiru" (to believe). The negative potential form is used to express that something is impossible or difficult to do. This phrase is used to express that you are surprised or shocked by something that is hard to believe.

\newpage

\section*{336. I'm into it}

\textbf{Date}: 2025.03.12 (Wed)\\
\textbf{Index}: muchuudesu\\
\textbf{Expression (JA)}: ...に(夢中:むちゅう)です\\
\textbf{Romanization}: ...ni muchuu desu\\
\textbf{Expression (EN)}: I'm into it

\textbf{Variants (JA)}: \\
\textbf{Variants (EN)}: 

\textbf{Intent (JA)}: \\
\textbf{Intent (EN)}: 

\textbf{Tags (JA)}: \\
\textbf{Tags (EN)}: 

\textbf{Adjusted Expression (JA)}:\\
\\
\textbf{Adjusted Expression (EN)}:\\


\textbf{Example (JA)}:\\
A「(最近:さいきん)、(何:なに)に(夢中:むちゅう)？」B「(読書:どくしょ)に(夢中:むちゅう)です」\\
\textbf{Example (EN)}:\\
A: What are you into lately? B: I'm into reading.

\textbf{Notes (JA)}:\\
「夢中になっています」は、「absorbed in」や「obsessed with」を意味する名詞「夢中」に、「に」は動作の対象を示す助詞、「です」は丁寧形の判定詞「だ」（to be）です。このフレーズは、何かに夢中であり、非常に興味を持っていることを表現するために使われます。だいたい「読書に夢中です」と言っておけば、勉強熱心な印象を与えることができます。マンガも読書の一つです。

\textbf{Notes (EN)}:\\
"Muchuu" is a noun that means "absorption" or "obsession." The particle "ni" indicates the object of the action, and "desu" is the polite form of the copula "da" (to be). This phrase is used to express that you are absorbed in something and have a strong interest in it. Saying "I'm into reading" can give the impression that you are eager to learn. Manga is also one of the things you can read.

\newpage

\section*{335. The heart of a parent}

\textbf{Date}: 2025.03.11 (Tue)\\
\textbf{Index}: okinodokusama\\
\textbf{Expression (JA)}: お(気:き)の(毒:どく)さま\\
\textbf{Romanization}: O-ki no doku\\
\textbf{Expression (EN)}: I'm sorry to hear that

\textbf{Variants (JA)}: \\
\textbf{Variants (EN)}: 

\textbf{Intent (JA)}: \\
\textbf{Intent (EN)}: 

\textbf{Tags (JA)}: \\
\textbf{Tags (EN)}: 

\textbf{Adjusted Expression (JA)}:\\
\\
\textbf{Adjusted Expression (EN)}:\\


\textbf{Example (JA)}:\\
A「(昨日:きのう)、(メガネ:めがね)、なくした」B「それは、お(気:き)の(毒:どく)さま。で、(見:み)つかった？」\\
\textbf{Example (EN)}:\\
A: I lost my glasses yesterday. B: I'm sorry to hear that. Did you find them?

\textbf{Notes (JA)}:\\
「お気の毒さま」は、誰かが不運な出来事に遭ったときに同情を示すフレーズです。「お気」は丁寧な接頭辞で、「気」は「心」や「気持ち」を意味します。「の」は所有を示す助詞、「毒」は「毒」や「不運」を意味し、「さま」は敬意を示す接尾辞です。

\textbf{Notes (EN)}:\\
"O-ki no doku sama" is an expression used to show sympathy when someone has experienced an unfortunate event. "O-ki" is a polite prefix with "ki" meaning "heart" or "mind." "No" is a possessive particle, "doku" means "poison" or "misfortune," and "sama" is an honorific suffix.

\newpage

\section*{334. Kind a like a...}

\textbf{Date}: 2025.03.10 (Mon)\\
\textbf{Index}: mitaina\\
\textbf{Expression (JA)}: みたいな\\
\textbf{Romanization}: Mitai na\\
\textbf{Expression (EN)}: Kinda like a...

\textbf{Variants (JA)}: \\
\textbf{Variants (EN)}: 

\textbf{Intent (JA)}: \\
\textbf{Intent (EN)}: 

\textbf{Tags (JA)}: \\
\textbf{Tags (EN)}: 

\textbf{Adjusted Expression (JA)}:\\
\\
\textbf{Adjusted Expression (EN)}:\\


\textbf{Example (JA)}:\\
A「(懐:なつ)かしい(気持:きも)ちになるよね」B「うん、(懐かしい:なつかしい)けど、(ワクワク:わくわく)する、みたいな」\\
\textbf{Example (EN)}:\\
A: It makes me feel nostalgic. B: Yeah, it's kind a like... nostalgic but also exciting.

\textbf{Notes (JA)}:\\
「みたいな」は、何かが他の何かに似ていることを表現するフレーズですが、多くの場合、適切な語句が見つからないときや、直接的な表現を避けたいときに使われます。「みたい」は「like」や「similar to」を意味し、「な」は名詞と接続するための助詞です。このフレーズを使うことで、類似性を曖昧に表現することができ、はっきりと断定せずに伝えることができます。日本語の「...みたいな...」は、具体的な説明を避けつつ何かに似ていることを伝えるときによく使われ、やや曖昧なニュアンスを含んでいます。また、「～みたいな感じ」「～みたいなもの」などの表現を用いることで、語尾をぼかすこともできます。

\textbf{Notes (EN)}:\\
"Mitai na" is a phrase used when you want to express that something is similar to something else but cannot find the exact word to describe it. "Mitai" means "like" or "similar to," and "na" is a particle that connects it to a noun. This phrase allows you to express similarity in an informal and slightly ambiguous way, without making a direct statement. In Japanese, "...mitai na..." is often used when you want to convey resemblance without giving a precise explanation. It is also common to blur the sentence ending with phrases like "...mitai na kanji" (like this feeling) or "...mitai na mono" (like this thing).

\newpage

\section*{333. Younger brother}

\textbf{Date}: 2025.03.09 (Sun)\\
\textbf{Index}: otoutosan\\
\textbf{Expression (JA)}: おとうとさん\\
\textbf{Romanization}: Otoutosan\\
\textbf{Expression (EN)}: Younger brother

\textbf{Variants (JA)}: \\
\textbf{Variants (EN)}: 

\textbf{Intent (JA)}: \\
\textbf{Intent (EN)}: 

\textbf{Tags (JA)}: \\
\textbf{Tags (EN)}: 

\textbf{Adjusted Expression (JA)}:\\
\\
\textbf{Adjusted Expression (EN)}:\\


\textbf{Example (JA)}:\\
A「(弟:おとうと)さん、(確:たし)かお(医者:いしゃ)さんでしたよね？」B「ええ、(弟:おとうと)は(今:いま)、NPO(法人:ほうじん)で(働:はたら)いていますよ」\\
\textbf{Example (EN)}:\\
A: Your younger brother was a doctor, wasn't he? B: Yes, my younger brother is working for an NPO now.

\textbf{Notes (JA)}:\\
一般的に「弟さん」は、他人の弟を丁寧に指すことばです。自分の弟について話すときは「弟」を使います。また、自分の弟を直接呼ぶときには、「弟さん」や「弟」のどちらも使えません。名前で呼びましょう。この例文は、相手の弟さんが医者だということをかつて聞き、それを思い出したことを確認しています。「お医者さんでした」の「でした」は「です」の過去形ですが、これは過去に知ったことを確認するときに使う表現です。弟さんがもう医者ではないということではありません。

\textbf{Notes (EN)}:\\
"Otoutosan" is a polite way to refer to someone else's younger brother. When referring to your own younger brother in conversation, you should use "otouto." You cannot refer to your own younger brother as either "otoutosan" or just "otouto" when addressing him directly, you should call him by his name. This example confirms that the other person's younger brother was once a doctor and recalls that fact. "Deshita" is the past tense of "desu," but it is used to confirm something you knew in the past. It does not mean that the younger brother is no longer a doctor.

\newpage

\section*{332. Older sister}

\textbf{Date}: 2025.03.08 (Sat)\\
\textbf{Index}: oneesan\\
\textbf{Expression (JA)}: おねえさん\\
\textbf{Romanization}: Oneesan\\
\textbf{Expression (EN)}: Older sister

\textbf{Variants (JA)}: \\
\textbf{Variants (EN)}: 

\textbf{Intent (JA)}: \\
\textbf{Intent (EN)}: 

\textbf{Tags (JA)}: \\
\textbf{Tags (EN)}: 

\textbf{Adjusted Expression (JA)}:\\
\\
\textbf{Adjusted Expression (EN)}:\\


\textbf{Example (JA)}:\\
A「(お姉:ねえ)さん、(何:なに)、(飲:の)んでるの？」B「(コーヒー:こーひー)」\\
\textbf{Example (EN)}:\\
A: What are you drinking, older sister? B: Coffee.

\textbf{Notes (JA)}:\\
一般的に「お姉さん」は、他人の姉を丁寧に呼ぶ言葉です。自分の姉に話しかけるときに使います。あなたのお姉さんでもない女の人に呼びかけるときは、非常に下品なことば遣いになります。

\textbf{Notes (EN)}:\\
"Oneesan" is a polite way to refer to someone else's older sister. It is also used when talking to your own older sister. However, when you use it to address someone who is not your older sister, it becomes very rude.

\newpage

\section*{331. Father}

\textbf{Date}: 2025.03.07 (Fri)\\
\textbf{Index}: otousan\\
\textbf{Expression (JA)}: おとうさん\\
\textbf{Romanization}: Otousan\\
\textbf{Expression (EN)}: Father or Sir

\textbf{Variants (JA)}: \\
\textbf{Variants (EN)}: 

\textbf{Intent (JA)}: \\
\textbf{Intent (EN)}: 

\textbf{Tags (JA)}: \\
\textbf{Tags (EN)}: 

\textbf{Adjusted Expression (JA)}:\\
\\
\textbf{Adjusted Expression (EN)}:\\


\textbf{Example (JA)}:\\
A「ちょっと、そこの(お父:とう)さん、(節税:せつぜい)してますか？」B「ええ、してますよ」\\
\textbf{Example (EN)}:\\
A: Excuse me, sir, are you doing tax planning? B: Yes, I am.

\textbf{Notes (JA)}:\\
一般的に「おとうさん」は、他人の父親を丁寧に呼ぶ言葉です。自分の父親に話しかけるときにも使います。しかし、この場合は、通りがかりの中年男性に呼びかけるときのことばで、その人が本当に誰かの父親であるかどうかはわかりません。

\textbf{Notes (EN)}:\\
"Otousan" is a polite way to refer to someone else's father. It is also used when talking to your own father. However, in this case, it is a word used to address a middle-aged man passing by, and it is not known whether he is really someone's father.

\newpage

\section*{330. How come it's not good?}

\textbf{Date}: 2025.03.06 (Thu)\\
\textbf{Index}: doushitedamenano\\
\textbf{Expression (JA)}: どうしてだめなの\\
\textbf{Romanization}: Doushite dame na no\\
\textbf{Expression (EN)}: How come it's not good?

\textbf{Variants (JA)}: \\
\textbf{Variants (EN)}: 

\textbf{Intent (JA)}: \\
\textbf{Intent (EN)}: 

\textbf{Tags (JA)}: \\
\textbf{Tags (EN)}: 

\textbf{Adjusted Expression (JA)}:\\
\\
\textbf{Adjusted Expression (EN)}:\\


\textbf{Example (JA)}:\\
A「これしちゃだめ」B「どうしてだめなの？」\\
\textbf{Example (EN)}:\\
A: Don't do this. B: Why not?

\textbf{Notes (JA)}:\\
「どうしてだめなの？」は、何かがよくない理由や何かが許されない理由を尋ねるためのフレーズです。「どうして」は「why」を、「だめ」は「no good」や「not allowed」を、「な」は形容詞を示す助詞、「の」は理由を示す助詞です。このフレーズは、何かがよくない理由や何かが許されない理由を尋ねるために使われます。もちろん、「どうして？」だけでも使えます。

\textbf{Notes (EN)}:\\
"Doushite dame na no?" is a phrase used to ask why something is not good or why something is not allowed. The word "doushite" means "why," "dame" means "no good" or "not allowed," "na" is a particle that indicates an adjective, and "no" is a particle that indicates the reason. This phrase is used to ask why something is not good or why something is not allowed. Of course, you can also use just "doushite?".

\newpage

\section*{329. What happened?}

\textbf{Date}: 2025.03.05 (Wed)\\
\textbf{Index}: doukashitano\\
\textbf{Expression (JA)}: どうかしたの？\\
\textbf{Romanization}: Douka shita no?\\
\textbf{Expression (EN)}: What happened?

\textbf{Variants (JA)}: \\
\textbf{Variants (EN)}: 

\textbf{Intent (JA)}: \\
\textbf{Intent (EN)}: 

\textbf{Tags (JA)}: \\
\textbf{Tags (EN)}: 

\textbf{Adjusted Expression (JA)}:\\
\\
\textbf{Adjusted Expression (EN)}:\\


\textbf{Example (JA)}:\\
A「(急:きゅう)に(顔色:かおいろ)が(悪:わる)くなったね。どうかしたの？」B「(うん)、(頭:あたま)が(痛:いた)いの」\\
\textbf{Example (EN)}:\\
A: You suddenly look pale. What happened? B: Yes, I have a headache.

\textbf{Notes (JA)}:\\
「どうかしたの？」は、誰かの外見や行動の変化に気づいたときに、何が起こったのか尋ねるためのフレーズです。「どうかした」は「what happened」や「what's wrong」を、「した」は動詞「する」（to do）の過去形です。「の」はいろいろな意味があるのですが、ここでは理由を述べています。他にも文末で使うと疑問や理由を尋ねる文が表現できます。

\textbf{Notes (EN)}:\\
"Douka shita no?" is a phrase used to ask someone what happened when you notice a change in their appearance or behavior. The word "doukashita" means "what happened" or "what's wrong," "shita" is the past tense of the verb "suru" (to do). The particle "no" has various meanings, but here it is explaining the reason. It can also be used at the end of a sentence to ask a question or give a reason.

\newpage

\section*{328. I can't get through}

\textbf{Date}: 2025.03.04 (Tue)\\
\textbf{Index}: tsunagaranai\\
\textbf{Expression (JA)}: (繋:つな)がらない\\
\textbf{Romanization}: Tsunagaranai\\
\textbf{Expression (EN)}: I can't get through

\textbf{Variants (JA)}: \\
\textbf{Variants (EN)}: 

\textbf{Intent (JA)}: \\
\textbf{Intent (EN)}: 

\textbf{Tags (JA)}: \\
\textbf{Tags (EN)}: 

\textbf{Adjusted Expression (JA)}:\\
\\
\textbf{Adjusted Expression (EN)}:\\


\textbf{Example (JA)}:\\
A「(電話:でんわ)、(繋:つな)がらない」B「(携帯:けいたい)、(電池:でんち)、(切:き)れてるかも」\\
\textbf{Example (EN)}:\\
A: I can't get through on the phone. B: Your cell phone battery may be dead.

\textbf{Notes (JA)}:\\
「つながらない」は、電話やメールで誰かと連絡が取れないことを示すためのフレーズです。「電話」は「phone」を、「つながらない」は自動詞「つながる」（to connect）の否定形、「携帯」は「cell phone」や「mobile phone」を、「電池」は「battery」を、「切れてる」は自動詞「切れる」（to be cut off）の現在の状況を示す「...ている」の構文です。

\textbf{Notes (EN)}:\\
"Tsunagaranai" is a phrase used to indicate that you cannot connect to someone by phone or email. The word "denwa" means "phone," "tsunagaranai" is the negative form of the verb "tsunagaru" (to connect), "keitai" means "cell phone" or "mobile phone," "denchi" means "battery," and "kireteru" is the present state of the verb "kireru" (to be cut off). "...te iru" is a construction that indicates the current state of something.

\newpage

\section*{327. I have allergies}

\textbf{Date}: 2025.03.03 (Mon)\\
\textbf{Index}: arerugi\\
\textbf{Expression (JA)}: (アレルギー:あれるぎー)があります\\
\textbf{Romanization}: Arerugii ga arimasu\\
\textbf{Expression (EN)}: I have allergies

\textbf{Variants (JA)}: \\
\textbf{Variants (EN)}: 

\textbf{Intent (JA)}: \\
\textbf{Intent (EN)}: 

\textbf{Tags (JA)}: \\
\textbf{Tags (EN)}: 

\textbf{Adjusted Expression (JA)}:\\
\\
\textbf{Adjusted Expression (EN)}:\\


\textbf{Example (JA)}:\\
A「(花粉:かふん)、(大丈夫:だいじょうぶ)？」B「ちょっと、(アレルギー:あれるぎー)があって、(目:め)が(痒:かゆ)くて」\\
\textbf{Example (EN)}:\\
A: Are you okay with pollen? B: I have allergies, so my eyes are itchy.

\textbf{Notes (JA)}:\\
「花粉」は「pollen」を、「アレルギー」は「allergy」を意味します。「大丈夫」は「okay」や「fine」、「目」は「eye」、「かゆい」は「itchy」という意味です。「アレルギーがあります」は、アレルギーを持っていることを示すためのフレーズです。日本では春先になると杉の花粉が飛び交い、アレルギー症状が出る人が多いです。この症状のことを「花粉症」といいます。

\textbf{Notes (EN)}:\\
"Kafun" means "pollen" and "arerugii" means "allergy." The word "daijoubu" means "okay" or "fine." The word "me" means "eye" and "kayui" means "itchy." The phrase "arerugii ga arimasu" is used to indicate that someone has allergies. The word "ga" is a particle that indicates the subject of the sentence and "arimasu" is the polite form of the verb "aru" (to have). In Japan, when spring comes, cedar pollen flies around, and many people have allergy symptoms. This condition is called "kafunshou."

\newpage

\section*{326. Just now, it's done}

\textbf{Date}: 2025.03.02 (Sun)\\
\textbf{Index}: choudoimadekitatokoro\\
\textbf{Expression (JA)}: ちょうど(今:いま)、(出来:でき)たところ。\\
\textbf{Romanization}: Choudo ima, dekita tokoro.\\
\textbf{Expression (EN)}: Just now, it's done.

\textbf{Variants (JA)}: \\
\textbf{Variants (EN)}: 

\textbf{Intent (JA)}: \\
\textbf{Intent (EN)}: 

\textbf{Tags (JA)}: \\
\textbf{Tags (EN)}: 

\textbf{Adjusted Expression (JA)}:\\
\\
\textbf{Adjusted Expression (EN)}:\\


\textbf{Example (JA)}:\\
A「(何:なに)、(作:つく)ってるの？」B「ちょうど(今:いま)、(カレー:かれー)うどん、できたところ」\\
\textbf{Example (EN)}:\\
A: What are you making? B: Just now, I finished making curry udon.

\textbf{Notes (JA)}:\\
「ちょうど今、出来たところ」は、何かがちょうど完成したことを示すためのフレーズです。「ちょうど」は「just」や「exactly」を、「今」は「now」を、「出来た」は動詞「できる」（to complete cooking or making）の過去形、「ところ」は「place」や「point」を意味します。このフレーズは、何かがちょうど完成したことを示すために使われます。とてもよいタイミングであることを意味します。

\textbf{Notes (EN)}:\\
"Choudo ima, dekita tokoro" is a phrase used to indicate that something has just been completed. The word "choudo" means "just" or "exactly," "ima" means "now," "dekita" is the past tense of the verb "dekiru" (to complete cooking or making), and "tokoro" means "place" or "point." This phrase is used to indicate that something has just been completed. It means that it is a very good timing.

\newpage

\section*{325. From now on}

\textbf{Date}: 2025.03.01 (Sat)\\
\textbf{Index}: korekara\\
\textbf{Expression (JA)}: これから\\
\textbf{Romanization}: Kore kara\\
\textbf{Expression (EN)}: From now on

\textbf{Variants (JA)}: \\
\textbf{Variants (EN)}: 

\textbf{Intent (JA)}: \\
\textbf{Intent (EN)}: 

\textbf{Tags (JA)}: \\
\textbf{Tags (EN)}: 

\textbf{Adjusted Expression (JA)}:\\
\\
\textbf{Adjusted Expression (EN)}:\\


\textbf{Example (JA)}:\\
A「(仕事:しごと)、(終:お)わった？」B「うん、これから、(友達:ともだち)と(食事:しょくじ)」\\
\textbf{Example (EN)}:\\
A: Is work over? B: Yes, from now on, I'm going to have dinner with my friends.

\textbf{Notes (JA)}:\\
「これから」は、今後何かが起こることを示すためのフレーズです。「これ」は「this」を、「から」は行動や出来事の始まりを示す助詞です。「これ」は「this」の意味を知っている人は多いと思いますが、会話では「これから」が「今から」の意味で使われることが多いのです。文章でも「将来」という意味で使われることもあります。

\textbf{Notes (EN)}:\\
"Kore kara" means "from now on" and you can use it to indicate that something will happen in the future. The word "kore" means "this" and "kara" is a particle that indicates the beginning of an action or event. Many people know that "kore" means "this," but "kore" often means "now" in conversation. In writing, it can also be used to mean "in the future".

\newpage

\section*{324. Let me explain this in detail after this song.}

\textbf{Date}: 2025.02.28 (Fri)\\
\textbf{Index}: kuwashikuhakono\\
\textbf{Expression (JA)}: (詳:くわ)しくは、この(曲:きょく)のあと。\\
\textbf{Romanization}: Kuwashiku ha, kono kyoku no ato.\\
\textbf{Expression (EN)}: Details after this song.

\textbf{Variants (JA)}: \\
\textbf{Variants (EN)}: 

\textbf{Intent (JA)}: \\
\textbf{Intent (EN)}: 

\textbf{Tags (JA)}: 即時文法, 省略, 状況依存型の表現, リアルタイム性が強い\\
\textbf{Tags (EN)}: immediate grammar, Many omissions, Context-dependent expression, Strong real-time nature

\textbf{Adjusted Expression (JA)}:\\
詳しくは、この曲が終わった後にご説明いたします。\\
\textbf{Adjusted Expression (EN)}:\\
Let me explain this in detail after this song is over.

\textbf{Example (JA)}:\\
A「これ、(プレゼント:ぷれぜんと)？」B「ええ、(詳:くわ)しくは、この(曲:きょく)のあと」\\
\textbf{Example (EN)}:\\
A: Is this a present? B: Yes, let me explain this in detail after this song.

\textbf{Notes (JA)}:\\
「詳しくは、この曲の後」は、何かの後に詳しく説明することを示すために使われる表現です。この表現では、通常の文であれば「詳しくは、この曲が終わった後にご説明いたします。」のように動詞が入るはずですが、ここでは動詞が省略されています。それでも文脈から意味が伝わるのは、即時文法の特徴です。言語は、文法を厳密に守らなくても、話し手の意図を伝えることができます。

\textbf{Notes (EN)}:\\
"Kuwashiku ha, kono kyoku no ato" is used to indicate that you will explain something in detail after something else. The word "kuwashiku" means "in detail," "ha" is a particle that indicates the topic of the sentence, "kono" means "this," "kyoku" means "song," and "ato" means "after." This expression omits the verb at the end, making it concise yet understandable in context. Japanese allows flexible word order and omission in spoken language.

\newpage

\section*{323. That's good...and so much.}

\textbf{Date}: 2025.02.27 (Thu)\\
\textbf{Index}: iidesune...sugoku\\
\textbf{Expression (JA)}: いいですね、すごく。\\
\textbf{Romanization}: Ii desu ne, sugoku.\\
\textbf{Expression (EN)}: That's good...and so much.

\textbf{Variants (JA)}: \\
\textbf{Variants (EN)}: 

\textbf{Intent (JA)}: \\
\textbf{Intent (EN)}: 

\textbf{Tags (JA)}: \\
\textbf{Tags (EN)}: 

\textbf{Adjusted Expression (JA)}:\\
\\
\textbf{Adjusted Expression (EN)}:\\


\textbf{Example (JA)}:\\
A「この(絵:え)、どう」B「いいですね、すごく」\\
\textbf{Example (EN)}:\\
A: How about this painting? B: That's good...and so much.

\textbf{Notes (JA)}:\\
「いいですね、すごく」は、何かが良いか印象的であることを表現するためのフレーズです。「いい」は「good」や「nice」を、「です」は丁寧形の判定詞「だ」（to be）、「ね」は同意や確認を求める文末の助詞、「すごく」は副詞で「so much」や「very」を意味します。このフレーズは、話しているときに、まず「いい」ということを言って、その後、それが「すごい」ことを言う形式です。書くときにはこの語順は使いません。とても印象的で素晴しいことを強調して表現するときに適しています。

\textbf{Notes (EN)}:\\
"Ii desu ne, sugoku" is a phrase used to express that something is good or impressive. The word "ii" means "good" or "nice," "desu" is the polite form of the copula "da" (to be), "ne" is a sentence-ending particle that seeks agreement or confirmation, and "sugoku" is an adverb that means "so much" or "very." You should not use this word order when writing. This phrase is suitable for emphasizing something very impressive and wonderful.

\newpage

\section*{322. Even if you do it}

\textbf{Date}: 2025.02.26 (Wed)\\
\textbf{Index}: shitatositemo\\
\textbf{Expression (JA)}: したとしても\\
\textbf{Romanization}: Shita to shite mo\\
\textbf{Expression (EN)}: Even if you do it

\textbf{Variants (JA)}: \\
\textbf{Variants (EN)}: 

\textbf{Intent (JA)}: \\
\textbf{Intent (EN)}: 

\textbf{Tags (JA)}: \\
\textbf{Tags (EN)}: 

\textbf{Adjusted Expression (JA)}:\\
\\
\textbf{Adjusted Expression (EN)}:\\


\textbf{Example (JA)}:\\
A「(一生懸命:いっしょうけんめい)、(勉強:べんきょう)したとしても、(合格:ごうかく)できなんじゃ意味がないよ」B「やってもみないのにそんなことわからないよ」\\
\textbf{Example (EN)}:\\
A: Even if you study hard, it's meaningless if you can't pass. B: You don't know that if you don't try.

\textbf{Notes (JA)}:\\
「したとしても」は、仮定の状況や条件を表現するためのフレーズです。「した」は動詞「する」（to do）の過去形、「と」は条件を示す助詞、「しても」は「even if」を意味するフレーズです。このフレーズは、実現可能性が低いわけではないが、仮に起こり得ることを表現するために使われます。

\textbf{Notes (EN)}:\\
"Shita to shite mo" is a phrase used to express a hypothetical situation or condition. The word "shita" is the past tense of the verb "suru" (to do), "to" is a particle that indicates the condition, and "shite mo" is a phrase that means "even if." This phrase is used to express that something is possible, but unlikely to happen.

\newpage

\section*{321. I finally got used to it}

\textbf{Date}: 2025.02.25 (Tue)\\
\textbf{Index}: youyakunaremashita\\
\textbf{Expression (JA)}: ようやく(慣:な)れました。\\
\textbf{Romanization}: Youyaku naremashita\\
\textbf{Expression (EN)}: I finally got used to it

\textbf{Variants (JA)}: \\
\textbf{Variants (EN)}: 

\textbf{Intent (JA)}: \\
\textbf{Intent (EN)}: 

\textbf{Tags (JA)}: \\
\textbf{Tags (EN)}: 

\textbf{Adjusted Expression (JA)}:\\
\\
\textbf{Adjusted Expression (EN)}:\\


\textbf{Example (JA)}:\\
A「(日本:にほん)の(生活:せいかつ)、(慣:な)れましたか？」B「うん、(ようやく)、(慣:な)れました」\\
\textbf{Example (EN)}:\\
A: Have you gotten used to life in Japan? B: Yes, I finally got used to it.

\textbf{Notes (JA)}:\\
「ようやく慣れました」は、誰かがついに何かに慣れたことを表現するためのフレーズです。「ようやく」は副詞で、「it took long time but finally...」という意味です。「慣れました」は動詞「慣れる」（to be accustomed to）の過去形です。

\textbf{Notes (EN)}:\\
"Youyaku naremashita" is a phrase used to express that someone has finally gotten used to something. The word "youyaku" is an adverb that means "it took a long time but finally..." and "naremashita" is the past tense of the verb "nareru" (to be accustomed to).

\newpage

\section*{320. This time, it's special/exceptional}

\textbf{Date}: 2025.02.24 (Mon)\\
\textbf{Index}: konkaitokubetsune\\
\textbf{Expression (JA)}: (今回:こんかい)、(特別:とくべつ)ね。\\
\textbf{Romanization}: Konkai, tokubetsu ne.\\
\textbf{Expression (EN)}: This time, it's special/exceptional

\textbf{Variants (JA)}: \\
\textbf{Variants (EN)}: 

\textbf{Intent (JA)}: \\
\textbf{Intent (EN)}: 

\textbf{Tags (JA)}: \\
\textbf{Tags (EN)}: 

\textbf{Adjusted Expression (JA)}:\\
\\
\textbf{Adjusted Expression (EN)}:\\


\textbf{Example (JA)}:\\
A「ママ、もう一つ、食べてもいい？」B「(今回:こんかい)、(特別:とくべつ)ね」\\
\textbf{Example (EN)}:\\
A: Mom, can I have one more? B: This time, it's special.

\textbf{Notes (JA)}:\\
「今回、特別ね」は、今回のことが特別であることを示すためのフレーズです。「今回」は「this time」を、「特別」は「今回だけ」や「例外」を、「ね」は同意や確認を求める文末の助詞です。この表現は、今回限りの例外的措置であることを認める表現になります。

\textbf{Notes (EN)}:\\
"Konkai, tokubetsu ne" is a phrase used to indicate that something is special or exceptional this time. The word "konkai" means "this time," "tokubetsu" means "special" or "exceptional," and "ne" is a sentence-ending particle that seeks agreement or confirmation. This expression acknowledges that the current situation is an exceptional measure.

\newpage

\section*{319. Let me think about it.}

\textbf{Date}: 2025.02.23 (Sun)\\
\textbf{Index}: chottokangaesasetekudasai\\
\textbf{Expression (JA)}: ちょっと(考:かんが)えさせてください\\
\textbf{Romanization}: Chotto kangae sasete kudasai\\
\textbf{Expression (EN)}: Let me think about it for a moment.

\textbf{Variants (JA)}: \\
\textbf{Variants (EN)}: 

\textbf{Intent (JA)}: \\
\textbf{Intent (EN)}: 

\textbf{Tags (JA)}: \\
\textbf{Tags (EN)}: 

\textbf{Adjusted Expression (JA)}:\\
\\
\textbf{Adjusted Expression (EN)}:\\


\textbf{Example (JA)}:\\
A「(今:いま)、(決:き)めるのは(早:はや)いですか？」B「ちょっと(考:かんが)えさせてください」\\
\textbf{Example (EN)}:\\
A: Is it too early to decide now? B: Let me think about it for a moment.

\textbf{Notes (JA)}:\\
「ちょっと考えさせてください」は、誰かに何かを考える時間を与えてもらうように頼むためのフレーズです。「ちょっと」は「a little」や「a bit」を、「考えさせて」は動詞「考える」（to think）の使役形で、「ください」は動詞「くれる」（to give）の丁寧な依頼形です。自分が何かしたいときにお願いするときにはこのように使役の形を使います。多くの言語で同様の表現があると思います。

\textbf{Notes (EN)}:\\
"Chotto kangaesasete kudasai" is a phrase used to ask someone to give you a moment to think about something. The word "chotto" means "a little" or "a bit," "kangaesasete" is the causative form of the verb "kangaeru" (to think), and "kudasai" is a polite request form of the verb "kureru" (to give). When asking someone to give you time to think about something, use the causative form of the verb. This is similar to many languages.

\newpage

\section*{318. Listen carefully}

\textbf{Date}: 2025.02.22 (Sat)\\
\textbf{Index}: mimisumasete\\
\textbf{Expression (JA)}: (耳:みみ)を(澄:す)ませて\\
\textbf{Romanization}: Mimi wo sumasete\\
\textbf{Expression (EN)}: Listen carefully

\textbf{Variants (JA)}: \\
\textbf{Variants (EN)}: 

\textbf{Intent (JA)}: \\
\textbf{Intent (EN)}: 

\textbf{Tags (JA)}: \\
\textbf{Tags (EN)}: 

\textbf{Adjusted Expression (JA)}:\\
\\
\textbf{Adjusted Expression (EN)}:\\


\textbf{Example (JA)}:\\
(耳:みみ)を(澄:す)ませて、(鳥:とり)の声、(聞:き)こえる？\\
\textbf{Example (EN)}:\\
Can you hear the birds singing? Listen carefully.

\textbf{Notes (JA)}:\\
「耳を澄ませて」は、誰かに注意深く聞くように頼むためのフレーズです。「耳」は「ear」を、「を」は動詞の目的を示す助詞を、「澄ませて」は動詞「澄ます」（to clear）のて形で、「て」は要求や命令を示す助詞です。

\textbf{Notes (EN)}:\\
"Mimi wo sumasete" is a phrase used to ask someone to listen carefully or to pay attention to a sound or voice. The word "mimi" means "ear," "wo" is a particle that indicates the object of the verb, "sumasete" is the te-form of the verb "sumasu" (to clear), and "te" is a particle that indicates a request or command.

\newpage

\section*{317. Just so you know}

\textbf{Date}: 2025.02.21 (Fri)\\
\textbf{Index}: chinamini\\
\textbf{Expression (JA)}: ちなみに\\
\textbf{Romanization}: Chinamini\\
\textbf{Expression (EN)}: On the related note

\textbf{Variants (JA)}: \\
\textbf{Variants (EN)}: 

\textbf{Intent (JA)}: \\
\textbf{Intent (EN)}: 

\textbf{Tags (JA)}: \\
\textbf{Tags (EN)}: 

\textbf{Adjusted Expression (JA)}:\\
\\
\textbf{Adjusted Expression (EN)}:\\


\textbf{Example (JA)}:\\
(彼:かれ)は、(若手:わかて)で(一番:いちばん)(人気:にんき)のある(画家:がか)です。そうそう、ちなみに、(彼:かれ)の(父親:ちちおや)も(画家:がか)です。\\
\textbf{Example (EN)}:\\
He is one of the most popular young painters. By the way, his father is also a painter.

\textbf{Notes (JA)}:\\
「ちなみに」は、追加情報や関連事実を提供するためのフレーズです。関連するトピックを紹介したり、本題とは直接関係のない追加情報を提供するために使われます。「ちなみに」は「just so you know」や「on the related note」を意味します。

\textbf{Notes (EN)}:\\
"Chinamini" is a phrase used to provide additional information or related facts. It is used to introduce a related topic or to provide additional information that is not directly related to the main topic. The word "chinamini" means "just so you know" or "on the related note."

\newpage

\section*{316. I thought you weren't coming}

\textbf{Date}: 2025.02.20 (Thu)\\
\textbf{Index}: konainokatoomotta\\
\textbf{Expression (JA)}: (来:こ)ないのかと(思:おも)った\\
\textbf{Romanization}: Konai no ka to omotta\\
\textbf{Expression (EN)}: I thought you weren't coming

\textbf{Variants (JA)}: \\
\textbf{Variants (EN)}: 

\textbf{Intent (JA)}: \\
\textbf{Intent (EN)}: 

\textbf{Tags (JA)}: \\
\textbf{Tags (EN)}: 

\textbf{Adjusted Expression (JA)}:\\
\\
\textbf{Adjusted Expression (EN)}:\\


\textbf{Example (JA)}:\\
A「(遅:おそ)いね」B「ごめん、(道:みち)が(混:こ)んでたから、(遅:おく)れちゃった」A「(ああ)、(来:こ)ないのかと(思:おも)った」\\
\textbf{Example (EN)}:\\
A: You're late. B: Sorry, the road was congested, so I'm late. A: Oh, I thought you weren't coming.

\textbf{Notes (JA)}:\\
「来ないのかと思った」は、誰かが来ないかもしれないと思ったことを表現するためのフレーズです。「来ない」は動詞「来る」（to come）の否定形、「の」は文の主語を示す助詞、「か」は疑問を示す助詞、「と」は話者の思考を示す助詞、「思った」は動詞「思う」（to think）の過去形です。

\textbf{Notes (EN)}:\\
"Konai no ka to omotta" is a phrase used to express that someone thought someone else might not come. The word "konai" is the negative form of the verb "kuru" (to come), "no" is a particle that indicates the subject of the sentence, "ka" is a particle that indicates a question, "to" is a particle that indicates the speaker's thoughts, and "omotta" is the past tense of the verb "omou" (to think).

\newpage

\section*{315. Wait a moment}

\textbf{Date}: 2025.02.19 (Wed)\\
\textbf{Index}: chottomatte\\
\textbf{Expression (JA)}: ちょっと(待:ま)って\\
\textbf{Romanization}: Chotto matte\\
\textbf{Expression (EN)}: Hold on a second

\textbf{Variants (JA)}: \\
\textbf{Variants (EN)}: 

\textbf{Intent (JA)}: \\
\textbf{Intent (EN)}: 

\textbf{Tags (JA)}: \\
\textbf{Tags (EN)}: 

\textbf{Adjusted Expression (JA)}:\\
\\
\textbf{Adjusted Expression (EN)}:\\


\textbf{Example (JA)}:\\
(すみません)、(今:いま)、(電話:でんわ)してるんで、ちょっと(待:ま)ってもらえます？\\
\textbf{Example (EN)}:\\
Excuse me, I'm on the phone right now. Can you wait a moment?

\textbf{Notes (JA)}:\\
「ちょっと待って」は、誰かに少し待ってもらうように頼むためのフレーズです。「ちょっと」は「a little」や「a bit」を、「待って」は動詞「待つ」（to wait）のて形です。

\textbf{Notes (EN)}:\\
"Chotto matte" is a phrase used to ask someone to wait for a moment. The word "chotto" means "a little" or "a bit," and "matte" is the te-form of the verb "matsu" (to wait).

\newpage

\section*{314. Today, there's a lesson}

\textbf{Date}: 2025.02.18 (Tue)\\
\textbf{Index}: kyohajugyogaaru\\
\textbf{Expression (JA)}: (今日:きょう)は(授業:じゅぎょう)がある\\
\textbf{Romanization}: Kyou wa jugyou ga aru\\
\textbf{Expression (EN)}: I have a class today

\textbf{Variants (JA)}: \\
\textbf{Variants (EN)}: 

\textbf{Intent (JA)}: \\
\textbf{Intent (EN)}: 

\textbf{Tags (JA)}: \\
\textbf{Tags (EN)}: 

\textbf{Adjusted Expression (JA)}:\\
\\
\textbf{Adjusted Expression (EN)}:\\


\textbf{Example (JA)}:\\
A「(今日:きょう)、(授業:じゅぎょう)ある？」B「うん、(午後:ごご)から」\\
\textbf{Example (EN)}:\\
A: Do you have a class today? B: Yes, in the afternoon.

\textbf{Notes (JA)}:\\
「今日は授業がある」は、特定の日に誰かが授業やレッスンがあることを示すためのフレーズです。「今日」は「today」を、「は」はそれが特定であることを意味する助詞、「授業」は「class」や「lesson」を、「がある」は「to have」を意味するフレーズです。

\textbf{Notes (EN)}:\\
"Kyou wa jugyou ga aru" is a phrase used to indicate that someone has a class or lesson on a particular day. The word "kyou" means "today," "wa" is a particle that indicates the subject of the sentence, "jugyou" means "class" or "lesson," "ga" is a particle that something is specific, and "aru" is a phrase that means "to have."

\newpage

\section*{313. It could happen.}

\textbf{Date}: 2025.02.17 (Mon)\\
\textbf{Index}: sonnakotomoaruyo\\
\textbf{Expression (JA)}: そんなこともあるよ\\
\textbf{Romanization}: Sonna koto mo aru yo\\
\textbf{Expression (EN)}: It could happen.

\textbf{Variants (JA)}: \\
\textbf{Variants (EN)}: 

\textbf{Intent (JA)}: \\
\textbf{Intent (EN)}: 

\textbf{Tags (JA)}: \\
\textbf{Tags (EN)}: 

\textbf{Adjusted Expression (JA)}:\\
\\
\textbf{Adjusted Expression (EN)}:\\


\textbf{Example (JA)}:\\
A「(実験:じっけん)、(失敗:しっぱい)」B「そんなこともあるよ」\\
\textbf{Example (EN)}:\\
A: The experiment failed. B: It could happen.

\textbf{Notes (JA)}:\\
「そんなこともあるよ」は、日常生活で望ましくないことや不運なことが起こる可能性があることを表現し、誰かや自分の気持ちを変えるために使われるフレーズです。「そんな」は「such」を、「こと」は「thing」を、「も」は「also」を、「ある」は動詞「ある」（to exist）ですが、この時は「起こりうる」という意味、「よ」は新しい情報を提供する文末の助詞です。

\textbf{Notes (EN)}:\\
"Sonna koto mo aru yo" is a phrase used to express that something unwanted or unfortunate could happen in daily life and is used to change someone's or one's own feelings. The word "sonna" means "such," "koto" means "thing," "mo" means "also," "aru" means "to exist," but in this case, it means "could happen," and "yo" is a sentence-ending particle that provides new information.

\newpage

\section*{312. Can you fix this?}

\textbf{Date}: 2025.02.16 (Sun)\\
\textbf{Index}: korenaoseru\\
\textbf{Expression (JA)}: これ、(直:なお)せる？\\
\textbf{Romanization}: Kore, naoseru?\\
\textbf{Expression (EN)}: Can you fix this?

\textbf{Variants (JA)}: \\
\textbf{Variants (EN)}: 

\textbf{Intent (JA)}: \\
\textbf{Intent (EN)}: 

\textbf{Tags (JA)}: \\
\textbf{Tags (EN)}: 

\textbf{Adjusted Expression (JA)}:\\
\\
\textbf{Adjusted Expression (EN)}:\\


\textbf{Example (JA)}:\\
A「(携帯:けいたい)、(壊:こわ)れた。これ、(直:なお)せる？」B「(直:なせ)るよ」A「(本当:ほんとう)！」B「うん、お(金:かね)を(出:だ)せば」\\
\textbf{Example (EN)}:\\
A: My phone is broken. Is it possible to fix this? B: It can be fixed. A: Really? B: Yes, if you pay.

\textbf{Notes (JA)}:\\
「直せる」は「to repair」や「to fix」を意味する動詞です。「携帯」は「cell phone」や「mobile phone」を、「壊れた」は動詞「壊れる」（to break）の過去形です。「これ」は「this」を、「直せる」は動詞「直す」（to fix）の可能形です。「これ、直せる？」は、壊れたり損傷したものを修理できるかどうか尋ねるためのフレーズです。「出せば」は動詞「出す」（to pay）の条件形です。

\textbf{Notes (EN)}:\\
"Naoru" is a verb that means "to repair" or "to fix." The word "keitai" means "cell phone" or "mobile phone." The word "kowareta" is the past tense of the verb "kowareru" (to break). The word "kore" means "this," and the word "naoseru" is the potential form of the verb "naosu" (to fix). The phrase "kore, naoseru?" is used to ask if someone can fix something that is broken or damaged. The word "daseba" is the conditional form of the verb "dasu" (to pay).

\newpage

\section*{311. Let's try to see it.}

\textbf{Date}: 2025.02.15 (Sat)\\
\textbf{Index}: choottomitemiyouyo\\
\textbf{Expression (JA)}: ちょっと(見:み)てみようよ\\
\textbf{Romanization}: Chotto mite miyou yo\\
\textbf{Expression (EN)}: Let's try to see it.

\textbf{Variants (JA)}: \\
\textbf{Variants (EN)}: 

\textbf{Intent (JA)}: \\
\textbf{Intent (EN)}: 

\textbf{Tags (JA)}: \\
\textbf{Tags (EN)}: 

\textbf{Adjusted Expression (JA)}:\\
\\
\textbf{Adjusted Expression (EN)}:\\


\textbf{Example (JA)}:\\
A「これ、(何:なに)？」B「ちょっと(見:み)てみようよ」\\
\textbf{Example (EN)}:\\
A: What's this? B: Let's try to see it.

\textbf{Notes (JA)}:\\
「ちょっと見てみようよ」は、何かを試したり、何かを見たりすることを提案するためのフレーズです。「ちょっと」は「a little」や「a bit」を、「見て」は動詞「見る」（to see）のて形で、「見よう」は動詞「見る」（to see）の意向形です。助詞「よ」は文を強調する助詞です。

\textbf{Notes (EN)}:\\
"Chotto mite miyou yo" is a expression used to suggest trying something or looking at something. The word "chotto" means "a little" or "a bit," "mite" is the te-form of the verb "miru" (to see), and "miyou" is the volitional form of the verb "miru" (to see). The word "yo" is a particle that emphasizes the statement.

\newpage

\section*{310. If you have something you want me to help with, let me know}

\textbf{Date}: 2025.02.14 (Fri)\\
\textbf{Index}: moshitetsudattehoshiikotogaattaraittene\\
\textbf{Expression (JA)}: もし、(手伝:てつだ)ってほしいことがあったら(言:い)ってね\\
\textbf{Romanization}: Moshi, tetsudatte hoshii koto ga attara itte ne\\
\textbf{Expression (EN)}: If you need help, let me know

\textbf{Variants (JA)}: \\
\textbf{Variants (EN)}: 

\textbf{Intent (JA)}: \\
\textbf{Intent (EN)}: 

\textbf{Tags (JA)}: \\
\textbf{Tags (EN)}: 

\textbf{Adjusted Expression (JA)}:\\
\\
\textbf{Adjusted Expression (EN)}:\\


\textbf{Example (JA)}:\\
A「もし、(手伝:てつだ)ってほしいことがあったら(言:い)ってね」B「ありがとう、(今:いま)のところ、(大丈夫:だいじょうぶ)」\\
\textbf{Example (EN)}:\\
A: If you need help, let me know. B: Thank you, I'm good for now.

\textbf{Notes (JA)}:\\
「もし、手伝ってほしいことがあったら言ってね」は、誰かに手助けや援助を申し出るためのフレーズです。「もし」は「if」、「手伝って」は動詞「手伝う」（to help）のて形、「てほしい」は「動詞＋てほしい」でyou want me to doという意味です。

\textbf{Notes (EN)}:\\
"Moshi, tetsudatte hoshii koto ga attara itte ne" is a phrase used to offer help or assistance to someone. The word "moshi" means "if," "tetsudatte" is the te-form of the verb "tetsudau" (to help), "verb + te hoshii" means "want me to do,"

\newpage

\section*{309. Ending particles and omitting particles}

\textbf{Date}: 2025.02.13 (Thu)\\
\textbf{Index}: darewomatteiruno\\
\textbf{Expression (JA)}: (誰:だれ)を(待:ま)っているの？\\
\textbf{Romanization}: Dare o matte iru no?\\
\textbf{Expression (EN)}: Who are you waiting for?

\textbf{Variants (JA)}: \\
\textbf{Variants (EN)}: 

\textbf{Intent (JA)}: \\
\textbf{Intent (EN)}: 

\textbf{Tags (JA)}: \\
\textbf{Tags (EN)}: 

\textbf{Adjusted Expression (JA)}:\\
\\
\textbf{Adjusted Expression (EN)}:\\


\textbf{Example (JA)}:\\
A「(誰:だれ)を(待:ま)っているの？」B「(友:とも)だち」\\
\textbf{Example (EN)}:\\
A: Who are you waiting for? B: My friend.

\textbf{Notes (JA)}:\\
「誰を待っているの？」は、誰かが誰を待っているか尋ねるためのフレーズです。「誰」は「who」を、「を」は動詞の目的を示す助詞を、「待っている」は動詞「待つ」（to wait）の進行形で、「の」は所有ではなくて、柔らかく尋ねる時の終助詞です。会話でカジュアル体で「か」を使うと乱暴な問い掛けになるので使わないようにしましょう。カジュアルの会話で疑問文を作るときは「の」、フォーマルの会話で疑問文を作るときは「ですか」「ますか」が適切です。この文で「誰を」のように「を」を使う理由は、動作の主体「誰が」を尋ねる文と区別するためです。文脈からこの2つの関係がわかる場合には、助詞を使わなくても意味はわかります。

\textbf{Notes (EN)}:\\
"Dare o matte iru no?" is a phrase used to ask who someone is waiting for. The word "dare" means "who," "o" is a particle that indicates the object of the verb, "matte iru" is the progressive form of the verb "matsu" (to wait), and "no" is not possessive but a soft ending particle used when asking gently. In casual conversation, avoid using "ka" in the casual form, as it can sound rude. Use "no" in casual conversation to form a question, and "desu ka" or "masu ka" in formal conversation. The reason for using "wo" as in "dare wo" is to distinguish this question from asking who is waiting for whom. In context, these two relationships can be understood without using particles.

\newpage

\section*{308. What a heck!}

\textbf{Date}: 2025.02.12 (Wed)\\
\textbf{Index}: sorenani\\
\textbf{Expression (JA)}: それ(何:なに)!\\
\textbf{Romanization}: Sore nani!\\
\textbf{Expression (EN)}: What a heck!

\textbf{Variants (JA)}: \\
\textbf{Variants (EN)}: 

\textbf{Intent (JA)}: \\
\textbf{Intent (EN)}: 

\textbf{Tags (JA)}: \\
\textbf{Tags (EN)}: 

\textbf{Adjusted Expression (JA)}:\\
\\
\textbf{Adjusted Expression (EN)}:\\


\textbf{Example (JA)}:\\
A「(昨日:きのう)、(彼:かれ)、(突然:とつぜん)、(辞:や)めたんだって」B「それ(何:なに)！」\\
\textbf{Example (EN)}:\\
A: I heard he quit yesterday. B: What a heck!

\textbf{Notes (JA)}:\\
「それ何！」は、予期せぬことや異常なことに対する驚き、不信、またはショックを表現するためのフレーズです。「それ」は「that」を、「何」は「what」を意味し、感嘆符は話者の強い反応を強調します。もちろん、同じ表現で普通に「それは何ですか」という意味でもあります。

\textbf{Notes (EN)}:\\
"Sore nani!" is a phrase used to express surprise, disbelief, or shock about something unexpected or unusual. The word "sore" means "that," "nani" means "what," and the exclamation mark emphasizes the speaker's strong reaction. Of course, it also means "What is that?" in a normal context.

\newpage

\section*{307. You can do it anytime}

\textbf{Date}: 2025.02.11 (Tue)\\
\textbf{Index}: itsudemoiidesuyo\\
\textbf{Expression (JA)}: いつでもいいですよ\\
\textbf{Romanization}: Itsudemo ii desu yo\\
\textbf{Expression (EN)}: Anytime is fine

\textbf{Variants (JA)}: \\
\textbf{Variants (EN)}: 

\textbf{Intent (JA)}: \\
\textbf{Intent (EN)}: 

\textbf{Tags (JA)}: \\
\textbf{Tags (EN)}: 

\textbf{Adjusted Expression (JA)}:\\
\\
\textbf{Adjusted Expression (EN)}:\\


\textbf{Example (JA)}:\\
A「(明日:あした)、(会議:かいぎ)は(午後:ごご)でも(大丈夫:だいじょうぶ)ですか？」B「はい、いつでもいいですよ」\\
\textbf{Example (EN)}:\\
A: Is it okay to have the meeting in the afternoon tomorrow? B: Yes, anytime is fine.

\textbf{Notes (JA)}:\\
「いつでもいいですよ」は、何かがいつでも受け入れられるか、都合がいいことを示すためのフレーズです。「いつでも」は「anytime」や「whenever」を、「いい」は「good」や「fine」を、「です」は丁寧形の判定詞「だ」（to be）で、「よ」は文を強調する助詞です。

\textbf{Notes (EN)}:\\
"Itsudemo ii desu yo" is a phrase used to indicate that something is acceptable or convenient at any time. The word "itsudemo" means "anytime" or "whenever," "ii" means "good" or "fine," "desu" is a polite form of the copula "da" (to be), and "yo" is a particle that emphasizes the statement. 

\newpage

\section*{306. Let's hurry}

\textbf{Date}: 2025.02.10 (Mon)\\
\textbf{Index}: isogimashou\\
\textbf{Expression (JA)}: (急:いそ)ぎましょう\\
\textbf{Romanization}: Isogimashou\\
\textbf{Expression (EN)}: Let's hurry

\textbf{Variants (JA)}: \\
\textbf{Variants (EN)}: 

\textbf{Intent (JA)}: \\
\textbf{Intent (EN)}: 

\textbf{Tags (JA)}: \\
\textbf{Tags (EN)}: 

\textbf{Adjusted Expression (JA)}:\\
\\
\textbf{Adjusted Expression (EN)}:\\


\textbf{Example (JA)}:\\
(時間:じかん)が(掛:か)かってしまいましたね。(急:いそ)ぎましょう。\\
\textbf{Example (EN)}:\\
It took longer than expected. Let's hurry.

\textbf{Notes (JA)}:\\
「急ぎましょう」は、誰かに急いだり、速く動いたりするように提案するためのフレーズです。「急ぐ」は「to hurry」や「to rush」を、「ましょう」は「let's」や「shall we」を意味するフレーズです。

\textbf{Notes (EN)}:\\
"Isogimashou" is a phrase used to suggest that someone should hurry or move quickly. The word "isogu" means "to hurry" or "to rush," and "mashou" is a phrase that means "let's" or "shall we." 

\newpage

\section*{305. Don't make a noise}

\textbf{Date}: 2025.02.09 (Sun)\\
\textbf{Index}: ototatenaide\\
\textbf{Expression (JA)}: (音:おと)を(立:た)てないで\\
\textbf{Romanization}: Oto o tatenaide\\
\textbf{Expression (EN)}: Don't make a noise

\textbf{Variants (JA)}: \\
\textbf{Variants (EN)}: 

\textbf{Intent (JA)}: \\
\textbf{Intent (EN)}: 

\textbf{Tags (JA)}: \\
\textbf{Tags (EN)}: 

\textbf{Adjusted Expression (JA)}:\\
\\
\textbf{Adjusted Expression (EN)}:\\


\textbf{Example (JA)}:\\
(今:いま)、(録画:ろくが)しているから、(音:おと)を(立:た)てないで。\\
\textbf{Example (EN)}:\\
I'm recording now, so don't make a noise.

\textbf{Notes (JA)}:\\
「音を立てないで」は、誰かに音や騒音を立てないように頼むためのフレーズです。「音」は「sound」や「noise」を、「を」は動詞の目的を示す助詞を、「立て」は動詞「立てる」（to make）の否定形で、「ないで」は動詞「ない」（not to be）のテ形です。

\textbf{Notes (EN)}:\\
"Oto o tatenaide" is a phrase used to ask someone not to make a sound or noise. The word "oto" means "sound" or "noise," "o" is a particle that indicates the object of the verb, "tate" is the verb "tateru" (to make) in the negative form, and "naide" is the te-form of the verb "nai" (not to be). 

\newpage

\section*{304. Did you hear that?}

\textbf{Date}: 2025.02.08 (Sat)\\
\textbf{Index}: imanokikoeta?\\
\textbf{Expression (JA)}: (今:いま)の(聞:き)こえた？\\
\textbf{Romanization}: Ima no kikoeta?\\
\textbf{Expression (EN)}: Did you hear that?

\textbf{Variants (JA)}: \\
\textbf{Variants (EN)}: 

\textbf{Intent (JA)}: \\
\textbf{Intent (EN)}: 

\textbf{Tags (JA)}: \\
\textbf{Tags (EN)}: 

\textbf{Adjusted Expression (JA)}:\\
\\
\textbf{Adjusted Expression (EN)}:\\


\textbf{Example (JA)}:\\
A「(ドア:どあ)の(音:おと)、(今:いま)の(聞:き)こえた？」B「うん、(何:なに)か(人:ひと)が(入:はい)ってきたみたい」\\
\textbf{Example (EN)}:\\
A: Did you hear the sound of the door just now? B: Yeah, it seems like someone came in.

\textbf{Notes (JA)}:\\
「今の聞こえた？」は、ちょうど今起こった音や騒音を聞いたかどうか尋ねるためのフレーズです。「今」は「now」、「聞こえた」は動詞「聞こえる」（to hear）の過去形です。「聞く」と「聞こえる」の違いは、「聞く」は意識的に聞くことを指し、「聞こえる」は偶然聞こえることを指します。「聞こえる」は現象ですから、この動詞の可能形はありません。したがって、否定形は「聞こえない」のように使います。

\textbf{Notes (EN)}:\\
"Ima no kikoeta?" is a phrase used to ask if someone heard a sound or noise that occurred just now. The word "ima" means "now," "kikoeta" is the past tense of the verb "kikoeru" (to hear). The difference between "kiku" and "kikoeru" is that "kiku" refers to actively listening, while "kikoeru" refers to hearing something by chance. Since "kikoeru" is a phenomenon, it does not have a potential form. Therefore, the negative form is used as "kikoenai."

\newpage

\section*{303. Do you want a little?}

\textbf{Date}: 2025.02.07 (Fri)\\
\textbf{Index}: sukoshihoshii?\\
\textbf{Expression (JA)}: (少:すこ)しほしい？\\
\textbf{Romanization}: Sukoshi hoshii?\\
\textbf{Expression (EN)}: Do you want a little?

\textbf{Variants (JA)}: \\
\textbf{Variants (EN)}: 

\textbf{Intent (JA)}: \\
\textbf{Intent (EN)}: 

\textbf{Tags (JA)}: \\
\textbf{Tags (EN)}: 

\textbf{Adjusted Expression (JA)}:\\
\\
\textbf{Adjusted Expression (EN)}:\\


\textbf{Example (JA)}:\\
A「(コーヒー:こーひー)、(少:すこ)しほしい？」B「いいえ、(大丈夫:だいじょうぶ)です」\\
\textbf{Example (EN)}:\\
A: Do you want a little coffee? B: No, I'm good.

\textbf{Notes (JA)}:\\
「少しほしい？」は、食べ物、飲み物、または援助など、何かの少量を欲しいかどうか尋ねるためのフレーズです。「少し」は「a little」や「a few」を、「ほしい」は動詞「ほしい」（to want）の形容詞形です。

\textbf{Notes (EN)}:\\
"Sukoshi hoshii?" is a phrase used to ask if someone wants a small amount of something, such as food, drink, or assistance. The word "sukoshi" means "a little" or "a few," "hoshii" is the adjective form of the verb "hoshii" (to want).

\newpage

\section*{302. for now}

\textbf{Date}: 2025.02.06 (Thu)\\
\textbf{Index}: imanotokoro\\
\textbf{Expression (JA)}: (今:いま)のところ、\\
\textbf{Romanization}: Ima no tokoro\\
\textbf{Expression (EN)}: for now

\textbf{Variants (JA)}: \\
\textbf{Variants (EN)}: 

\textbf{Intent (JA)}: \\
\textbf{Intent (EN)}: 

\textbf{Tags (JA)}: \\
\textbf{Tags (EN)}: 

\textbf{Adjusted Expression (JA)}:\\
\\
\textbf{Adjusted Expression (EN)}:\\


\textbf{Example (JA)}:\\
(今:いま)のところ、(晴:は)れてるけど、(今朝:けさ)の(予報:よほう)では(雨:あめ)。\\
\textbf{Example (EN)}:\\
For now, it's sunny, but the forecast this morning said it would rain.

\textbf{Notes (JA)}:\\
「今のところ」は、特定の時点での現在の状況や状態を示し、将来的に変化や変更の可能性があることが含まれています。「今」は「now」を、「の」は所有や関連を示す助詞を、「ところ」は「point」や「place」を意味します。

\textbf{Notes (EN)}:\\
"Ima no tokoro" is a phrase used to indicate the current situation or status at a specific point in time. The word "ima" means "now," "no" is a particle that indicates possession or association, and "tokoro" means "point" or "place."

\newpage

\section*{301. Full of}

\textbf{Date}: 2025.02.05 (Wed)\\
\textbf{Index}: tappuri\\
\textbf{Expression (JA)}: たっぷり\\
\textbf{Romanization}: Tappuri\\
\textbf{Expression (EN)}: Full of

\textbf{Variants (JA)}: \\
\textbf{Variants (EN)}: 

\textbf{Intent (JA)}: \\
\textbf{Intent (EN)}: 

\textbf{Tags (JA)}: \\
\textbf{Tags (EN)}: 

\textbf{Adjusted Expression (JA)}:\\
\\
\textbf{Adjusted Expression (EN)}:\\


\textbf{Example (JA)}:\\
(今日:きょう)は、(自由:じゆう)な(時間:じかん)がたっぷりあってうれしい。\\
\textbf{Example (EN)}:\\
I'm happy because I have plenty of free time today.

\textbf{Notes (JA)}:\\
「たっぷり」は、時間や空間、資源などが豊富にあることを表す言葉です。「full of」や「plenty of」の意味に相当します。何かが十分にあり、ゆとりをもって使える状態を表すときによく使われます。「たくさん」との違いは、「たっぷり」には単に量が多いだけでなく、「余裕をもって十分に使える」というニュアンスが含まれる点です。

\textbf{Notes (EN)}:\\
"Tappuri" is a word used to express an abundant amount of something, such as time, space, or resources. It means "full of" or "plenty of." This expression is commonly used to indicate that something is available in a sufficient or generous amount. The difference between "tappuri" and "takusan" is that "tappuri" implies not just a large quantity, but also that the amount is sufficient and can be fully utilized.

\newpage

\section*{300. I just remembered}

\textbf{Date}: 2025.02.04 (Tue)\\
\textbf{Index}: imaomoidashita\\
\textbf{Expression (JA)}: (今:いま)、(思:おも)い(出:だ)した！\\
\textbf{Romanization}: Ima omoidashita\\
\textbf{Expression (EN)}: I just rememberedi!

\textbf{Variants (JA)}: \\
\textbf{Variants (EN)}: 

\textbf{Intent (JA)}: \\
\textbf{Intent (EN)}: 

\textbf{Tags (JA)}: \\
\textbf{Tags (EN)}: 

\textbf{Adjusted Expression (JA)}:\\
\\
\textbf{Adjusted Expression (EN)}:\\


\textbf{Example (JA)}:\\
(今:いま)(思:おも)い(出:だ)した！(明日:あした)、(会議:かいぎ)があるんだった。\\
\textbf{Example (EN)}:\\
Oh, I just remembered! I have a meeting tomorrow!

\textbf{Notes (JA)}:\\
「今思い出した」は、何かを思い出したことを表現するフレーズです。「今」は「now」、「思い出した」は動詞「思い出す」（to remember）の過去形ですが、この場合、単なる過去ではなく「ちょうど今、思い出した」という意味になります。このように、突然何かを思い出したり、気づいたりしたときによく使われます。同様に「会議があるんだった」の「だった」も、過去の出来事を指しているのではなく「今、気づいた」という意味を持ちます。未来の予定について過去形が使われるのは、日本語において、再認識や思い出しのニュアンスを表すための用法です。

\textbf{Notes (EN)}:\\
"Ima omoidashita" means "I just remembered!" "Ima" means "now," and "omoida-shita" is the past tense of "omoidasu" (to remember). Although it appears to be in the past tense, it expresses a realization happening at this very moment. This usage is common when someone suddenly recalls something they had forgotten. Similarly, in "arun-datta", even though the meeting is in the future (tomorrow), "datta" (past tense) is used. This past form does not indicate a past event but rather that the speaker's realization at the present moment.

\newpage

\section*{299. That's right}

\textbf{Date}: 2025.02.03 (Mon)\\
\textbf{Index}: sonotoori\\
\textbf{Expression (JA)}: その(通:とお)り\\
\textbf{Romanization}: Sono toori\\
\textbf{Expression (EN)}: That's right

\textbf{Variants (JA)}: \\
\textbf{Variants (EN)}: 

\textbf{Intent (JA)}: \\
\textbf{Intent (EN)}: 

\textbf{Tags (JA)}: \\
\textbf{Tags (EN)}: 

\textbf{Adjusted Expression (JA)}:\\
\\
\textbf{Adjusted Expression (EN)}:\\


\textbf{Example (JA)}:\\
A「つまり、(日本語:にほんご)では(主語:しゅご)を(省略:しょうりゃく)することが(多:おお)いんですね？」 B「その(通:とお)り！」\\
\textbf{Example (EN)}:\\
A: So, in Japanese, it's common to omit the subject, isn't that right? B: That's right!

\textbf{Notes (JA)}:\\
「その通り」は、単なる同意ではなく、相手が論理的に導き出した正しい推論を認める表現であり、短く明確に答えられるのが特徴です。そのため、単なる同意の「ええ、そうですね」とは意味が異なります。「その」は「that」を、「通り」は「道」や「やり方」という意味で、そこから「正しいこと」「正確であること」を示します。

\textbf{Notes (EN)}:\\
"Sono toori" is not just a simple agreement but an acknowledgment that the other person's logically derived inference is correct. It is a concise and clear way to affirm their reasoning. In this sense, it differs from a simple agreement like "Yes, that's right." The word "sono" means "that," and "toori" means "way" or "manner," implying correctness or accuracy.

\newpage

\section*{298. Do you want more?}

\textbf{Date}: 2025.02.02 (Sun)\\
\textbf{Index}: mouchottoirimasuka\\
\textbf{Expression (JA)}: もうちょっと(要:い)りますか。\\
\textbf{Romanization}: Mou chotto irimasuka\\
\textbf{Expression (EN)}: Do you want more?

\textbf{Variants (JA)}: \\
\textbf{Variants (EN)}: 

\textbf{Intent (JA)}: \\
\textbf{Intent (EN)}: 

\textbf{Tags (JA)}: \\
\textbf{Tags (EN)}: 

\textbf{Adjusted Expression (JA)}:\\
\\
\textbf{Adjusted Expression (EN)}:\\


\textbf{Example (JA)}:\\
A「(コーヒー:こーひー)、もうちょっと(要:い)りますか？」B「いいえ、どうも」\\
\textbf{Example (EN)}:\\
A: Do you want more coffee? B: No, I'm good.

\textbf{Notes (JA)}:\\
「もうちょっといりますか」は、食べ物、飲み物もう一つほしいかどうかを尋ねるためのフレーズです。「もう」は「more」や「another」を、「ちょっと」は「a little」や「a bit」を、「いりますか」は動詞「いる」（to need or want）の丁寧形です。

\textbf{Notes (EN)}:\\
"Mou chotto irimasuka" is a phrase used to ask if someone wants more of something, such as food, drink, or assistance. The word "mou" means "more" or "another," "chotto" means "a little" or "a bit," "irimasuka" is a polite form of the verb "iru" (to need or want). 

\newpage

\section*{297. We are on the same boat}

\textbf{Date}: 2025.02.01 (Sat)\\
\textbf{Index}: norikakattafune\\
\textbf{Expression (JA)}: (乗:の)りかかった(船:ふね)\\
\textbf{Romanization}: Norikakatta fune\\
\textbf{Expression (EN)}: We are on the same boat

\textbf{Variants (JA)}: \\
\textbf{Variants (EN)}: 

\textbf{Intent (JA)}: \\
\textbf{Intent (EN)}: 

\textbf{Tags (JA)}: \\
\textbf{Tags (EN)}: 

\textbf{Adjusted Expression (JA)}:\\
\\
\textbf{Adjusted Expression (EN)}:\\


\textbf{Example (JA)}:\\
(彼:かれ)に(頼:たの)まれて(プロジェクト:ぷろじぇくと)に(参加:さんか)したけれど、(途中:とちゅう)でやめるわけにはいかない。(乗:の)りかかった(船:ふね)だし、(最後:さいご)までやるしかない。\\
\textbf{Example (EN)}:\\
I was asked by him to join the project, but I can't quit halfway. We are on the same boat, so I have to finish it.

\textbf{Notes (JA)}:\\
日本語の慣用句「乗りかかった船」は、途中で関与してしまい、もう後戻りできない状況を指します。そのため、最後までやり遂げるしかないという意味を持ちます。「乗りかかった」は動詞「乗りかかる」から来ており、「乗る」は「何かに乗る」（例：船に乗る）、「かかる」は「未完了または途中の状態」を表します。「be on the same boat」と似ていますが、「乗りかかった船」は個人の責任感を強調する一方、「be on the same boat」は状況や困難を共有する連帯感を強調します。

\textbf{Notes (EN)}:\\
The Japanese idiom "norikakatta fune" means being halfway into something and unable to back out, so you must see it through. "Norikakatta" comes from the verb "norikakaru," where "noru" means "to get on" (e.g., a boat) and "kakaru" describes an incomplete or ongoing state. While similar to "be on the same boat," "norikakatta fune" emphasizes personal responsibility, whereas "be on the same boat" highlights shared circumstances or solidarity.

\newpage

\section*{296. I can't easily fall asleep}

\textbf{Date}: 2025.01.31 (Fri)\\
\textbf{Index}: nakanakanamerenai\\
\textbf{Expression (JA)}: なかなか(眠:ねむ)れない\\
\textbf{Romanization}: Nakanaka nemurenai\\
\textbf{Expression (EN)}: I can't sleep well

\textbf{Variants (JA)}: \\
\textbf{Variants (EN)}: 

\textbf{Intent (JA)}: \\
\textbf{Intent (EN)}: 

\textbf{Tags (JA)}: \\
\textbf{Tags (EN)}: 

\textbf{Adjusted Expression (JA)}:\\
\\
\textbf{Adjusted Expression (EN)}:\\


\textbf{Example (JA)}:\\
(最近:さいきん)、なかなか(眠:ねむ)れないんです。\\
\textbf{Example (EN)}:\\
Lately, I haven't been able to sleep well.

\textbf{Notes (JA)}:\\
「なかなか眠れない」は、寝つきが悪いか、よく眠れないことを表現するためのフレーズです。「なかなか」は「not easily」や「not readily」を、「眠れない」は動詞「眠る」（to fall asleep）の可能形を否定形にしたものです。

\textbf{Notes (EN)}:\\
"Nakanaka nemurenai" is a phrase used to express that one has difficulty falling asleep or sleeping well. The word "nakanaka" means "not easily" or "not readily," "nemurenai" is the negative potential form of the verb "nemuru" (to fall asleep).

\newpage

\section*{295. How was it?}

\textbf{Date}: 2025.01.30 (Thu)\\
\textbf{Index}: doudatta\\
\textbf{Expression (JA)}: どうだった\\
\textbf{Romanization}: Dou datta\\
\textbf{Expression (EN)}: How was it?

\textbf{Variants (JA)}: \\
\textbf{Variants (EN)}: 

\textbf{Intent (JA)}: \\
\textbf{Intent (EN)}: 

\textbf{Tags (JA)}: \\
\textbf{Tags (EN)}: 

\textbf{Adjusted Expression (JA)}:\\
\\
\textbf{Adjusted Expression (EN)}:\\


\textbf{Example (JA)}:\\
A「(映画:えいが)、どうだった？」B「(面白:おもしろ)かったよ」\\
\textbf{Example (EN)}:\\
A: How was the movie? B: It was interesting.

\textbf{Notes (JA)}:\\
「どうだった」は、誰かの経験、印象、または何かについての意見を尋ねるためのフレーズです。「どう」は「how」を、「だった」は動詞「だ」（to be）の過去形です。このように短かく簡単で一息で言えるフレーズをたくさん使ってみることが上達の定石です。

\textbf{Notes (EN)}:\\
"Dou datta" is a phrase used to ask about someone's experience, impression, or opinion of something. The word "dou" means "how," "datta" is the past form of the verb "da" (to be). Such as "dou datta" is one of the fundamental techniques to improve your Japanese skills. It is recommended to use many short, simple phrases that can be said in one breath.

\newpage

\section*{294. Whose is that?}

\textbf{Date}: 2025.01.29 (Wed)\\
\textbf{Index}: sorehadareno\\
\textbf{Expression (JA)}: それは(誰:だれ)の？\\
\textbf{Romanization}: Sore wa dare no?\\
\textbf{Expression (EN)}: Whose is that?

\textbf{Variants (JA)}: \\
\textbf{Variants (EN)}: 

\textbf{Intent (JA)}: \\
\textbf{Intent (EN)}: 

\textbf{Tags (JA)}: \\
\textbf{Tags (EN)}: 

\textbf{Adjusted Expression (JA)}:\\
\\
\textbf{Adjusted Expression (EN)}:\\


\textbf{Example (JA)}:\\
A「この(傘:かさ)、(誰:だれ)の？」B「あ、それ、(私:わたし)の」\\
\textbf{Example (EN)}:\\
A: Whose umbrella is this? B: Oh, that's mine.

\textbf{Notes (JA)}:\\
「それは誰の」は、何かが誰のものであるか、または誰が何かに責任を持っているかを尋ねるためのフレーズです。「それ」は「that」を、「は」は文の主題を示す助詞を、「誰」は「who」を、「の」は所有や関連を示す助詞です。

\textbf{Notes (EN)}:\\
"Sore wa dare no" is a phrase used to ask who something belongs to or who is responsible for something. The word "sore" means "that," "wa" is a particle that marks the topic of the sentence, "dare" means "who," and "no" is a particle that indicates possession or association. 

\newpage

\section*{293. I have no idea}

\textbf{Date}: 2025.01.28 (Tue)\\
\textbf{Index}: kentoumotsukimasen\\
\textbf{Expression (JA)}: (見当:けんとう)もつきません\\
\textbf{Romanization}: Kentou mo tsukimasen\\
\textbf{Expression (EN)}: I have no idea

\textbf{Variants (JA)}: \\
\textbf{Variants (EN)}: 

\textbf{Intent (JA)}: \\
\textbf{Intent (EN)}: 

\textbf{Tags (JA)}: \\
\textbf{Tags (EN)}: 

\textbf{Adjusted Expression (JA)}:\\
\\
\textbf{Adjusted Expression (EN)}:\\


\textbf{Example (JA)}:\\
(何:なん)でみんな(並:なら)んでいるのか、(見当:けんとう)もつきません。\\
\textbf{Example (EN)}:\\
I have no idea why everyone is lining up.

\textbf{Notes (JA)}:\\
「見当もつきません」は、何かについてのアイデアや推測が全くないことを表現するためのフレーズです。「見当」は「idea」や「guess」を、「も」は「even」を意味し、「つきません」は動詞「つく」（to attach or estimate）の否定形です。この表現は、特定のトピックや状況についての知識、理解、または洞察がないことを伝えるためによく使われます。

\textbf{Notes (EN)}:\\
"Kentou mo tsukimasen" is a phrase used to express that one has no idea, clue, or estimate about something. The word "kentou" means "idea" or "guess," "mo" is a particle that means "even," and "tsukimasen" is the negative form of the verb "tsuku" (to attach or estimate). This expression is commonly used to convey a lack of knowledge, understanding, or insight about a particular topic or situation.

\newpage

\section*{292. What do you think?}

\textbf{Date}: 2025.01.27 (Mon)\\
\textbf{Index}: douomou\\
\textbf{Expression (JA)}: どう(思:おも)いますか。\\
\textbf{Romanization}: Dou omoimasuka\\
\textbf{Expression (EN)}: What do you think?

\textbf{Variants (JA)}: \\
\textbf{Variants (EN)}: 

\textbf{Intent (JA)}: \\
\textbf{Intent (EN)}: 

\textbf{Tags (JA)}: \\
\textbf{Tags (EN)}: 

\textbf{Adjusted Expression (JA)}:\\
\\
\textbf{Adjusted Expression (EN)}:\\


\textbf{Example (JA)}:\\
A「この(問題:もんだい)、どう(思:おも)いますか？」B「(難:むずか)しいですね」\\
\textbf{Example (EN)}:\\
A: What do you think about this problem? B: It's difficult.

\textbf{Notes (JA)}:\\
「どう思う」は、特定のトピックや状況について誰かの意見、考え、または感情を尋ねるためのフレーズです。「どう」は「how」や「what」を、「思う」は「to think」や「to feel」を意味します。この表現は、特定の問題や問題に対する誰かの視点、視点、または反応について尋ねるためによく使われます。親しい人には「どう思う」と短かく尋ねることができます。

\textbf{Notes (EN)}:\\
"Dou omou" is a phrase used to ask for someone's opinion, thoughts, or feelings about a particular topic or situation. The word "dou" means "how" or "what," and "omou" means "to think" or "to feel." This expression is commonly used to inquire about someone's perspective, viewpoint, or reaction to a specific matter or issue. You can ask someone close to you with a short "dou omou."

\newpage

\section*{291. The majority of people}

\textbf{Date}: 2025.01.26 (Sun)\\
\textbf{Index}: ookunohito\\
\textbf{Expression (JA)}: (多:おお)くの(人:ひと)\\
\textbf{Romanization}: Ooku no hito\\
\textbf{Expression (EN)}: The majority of people

\textbf{Variants (JA)}: \\
\textbf{Variants (EN)}: 

\textbf{Intent (JA)}: \\
\textbf{Intent (EN)}: 

\textbf{Tags (JA)}: \\
\textbf{Tags (EN)}: 

\textbf{Adjusted Expression (JA)}:\\
\\
\textbf{Adjusted Expression (EN)}:\\


\textbf{Example (JA)}:\\
(多:おお)くの(人:ひと)はそう(思:おも)っている。\\
\textbf{Example (EN)}:\\
The majority of people think so.

\textbf{Notes (JA)}:\\
「多くの人」は、グループや人口の大多数またはほとんどの人々を指すためのフレーズです。「多く」は「many」や「a lot」を、「の」は所有や関連を示す助詞を、「人」は「people」や「person」を意味します。この表現は、ほとんどの人々が何を考え、信じ、または行うかについて一般化したり、述べたりするためによく使われます。似ていることばに「多い」があります。これはイ形容詞ですが、「多い人」という言い方はないので注意が必要です。

\textbf{Notes (EN)}:\\
"Ooku no hito" is a phrase used to refer to the majority or most people in a group or population. The word "ooku" means "many" or "a lot," "no" is a particle that indicates possession or association, and "hito" means "people" or "person." This expression is commonly used to make generalizations or statements about what most people think, believe, or do. A similar phrase is "ooi," which is an i-adjective, but be careful as there is no expression like "ooi hito."

\newpage

\section*{290. I do things my way}

\textbf{Date}: 2025.01.25 (Sat)\\
\textbf{Index}: bokuhabokunarini\\
\textbf{Expression (JA)}: ぼくはぼくなりに\\
\textbf{Romanization}: Boku wa boku nari ni\\
\textbf{Expression (EN)}: In my own way

\textbf{Variants (JA)}: \\
\textbf{Variants (EN)}: 

\textbf{Intent (JA)}: \\
\textbf{Intent (EN)}: 

\textbf{Tags (JA)}: \\
\textbf{Tags (EN)}: 

\textbf{Adjusted Expression (JA)}:\\
\\
\textbf{Adjusted Expression (EN)}:\\


\textbf{Example (JA)}:\\
A「(この問題:このもんだい)、(どう:どう)やって(解:と)く？」B「ぼくはぼくなりに(考:かんが)えてみるよ」\\
\textbf{Example (EN)}:\\
A: How do you solve this problem? B: I'll think about it in my own way.

\textbf{Notes (JA)}:\\
「僕は僕なりに」は、自分の信念、価値観、または好みに従って物事を行うことを表現するためのフレーズです。「僕」は「I」や「me」を、「は」は文の主題を示す助詞を、「なり」は「in one's own way」を意味し、「に」は何かを行う方法や方法を示す助詞です。この表現は、特定の問題や状況に対する独自のアプローチや視点を持っていることを伝えるためによく使われます。「ぼく」は男性が使う「I」や「me」の意味です。女性の場合は「わたし」を使います。

\textbf{Notes (EN)}:\\
"Boku wa boku nari ni" is a phrase used to express that one does things according to their own beliefs, values, or preferences. The word "boku" means "I" or "me," "wa" is a particle that marks the topic of the sentence, "nari" is a particle that means "in one's own way," and "ni" is a particle that indicates the method or manner of doing something. This expression is commonly used to convey that one has their own unique approach or perspective on a particular matter or situation. "Boku" is the meaning of "I" or "me" used. In the case of females, use "watashi".

\newpage

\section*{289. What do you think from your standpoints?}

\textbf{Date}: 2025.01.24 (Fri)\\
\textbf{Index}: santekiniwadoudesuka\\
\textbf{Expression (JA)}: (田中:たなか)さん(的:てき)にはどうですか。\\
\textbf{Romanization}: Tanaka-san teki ni wa dou desuka\\
\textbf{Expression (EN)}: What do you think from your standpoint?

\textbf{Variants (JA)}: \\
\textbf{Variants (EN)}: 

\textbf{Intent (JA)}: \\
\textbf{Intent (EN)}: 

\textbf{Tags (JA)}: \\
\textbf{Tags (EN)}: 

\textbf{Adjusted Expression (JA)}:\\
\\
\textbf{Adjusted Expression (EN)}:\\


\textbf{Example (JA)}:\\
A「この(提案:ていあん)、(賛成:さんせい)ですか？」B「(田中:たなか)さん(的:てき)にはどうですか」\\
\textbf{Example (EN)}:\\
A: Do you agree with this proposal? B: What do you think from your standpoint, Tanaka-san?

\textbf{Notes (JA)}:\\
「的にはどうですか」は、特定のトピックや状況についての誰かの意見や視点を尋ねるためのフレーズです。「的」は「-like」や「-ish」を、「に」は行動の対象や受取人を示す助詞で、「は」は文の主題を示す助詞で、「どうですか」は「how is it?」や「what do you think?」を意味するフレーズです。この表現は、個人の視点や立場に基づいて、誰かの考えや感情、反応について尋ねるためによく使われます。

\textbf{Notes (EN)}:\\
"Teki ni wa dou desuka" is a phrase used to ask for someone's opinion or perspective on a particular topic or situation. The word "teki" means "-like" or "-ish," "ni" is a particle that indicates the target or recipient of an action, "wa" is a particle that marks the topic of the sentence, and "dou desuka" is a phrase that means "how is it?" or "what do you think?" This expression is commonly used to inquire about someone's thoughts, feelings, or reactions based on their personal viewpoint or standpoint.

\newpage

\section*{288. I can't anymore}

\textbf{Date}: 2025.01.23 (Thu)\\
\textbf{Index}: もう(無理:むり)\\
\textbf{Expression (JA)}: もう(無理:むり)\\
\textbf{Romanization}: Mou muri\\
\textbf{Expression (EN)}: I can't anymore

\textbf{Variants (JA)}: \\
\textbf{Variants (EN)}: 

\textbf{Intent (JA)}: \\
\textbf{Intent (EN)}: 

\textbf{Tags (JA)}: \\
\textbf{Tags (EN)}: 

\textbf{Adjusted Expression (JA)}:\\
\\
\textbf{Adjusted Expression (EN)}:\\


\textbf{Example (JA)}:\\
A「もう(一杯:いっぱい)、いかがですか？」B「もう(無理:むり)、どうもありがとう。これ(以上:いじょう)(飲:の)んだら(明日:あした)が(大変:たいへん)だから」\\
\textbf{Example (EN)}:\\
A: How about one more drink? B: I can't anymore, thank you. If I drink any more, tomorrow's going to be rough.

\textbf{Notes (JA)}:\\
「もう無理」は、疲労、疲労、または困難により、限界に達したか、続けることができないことを表現するためのフレーズです。「もう」は「already」や「anymore」を、「無理」は「impossible」や「unbearable」を意味します。

\textbf{Notes (EN)}:\\
"Mou muri" is a phrase used to express that one has reached their limit or cannot continue due to exhaustion, fatigue, or difficulty. The word "mou" means "already" or "anymore," and "muri" means "impossible" or "unbearable."

\newpage

\section*{287. I don't know the reason}

\textbf{Date}: 2025.01.22 (Wed)\\
\textbf{Index}: wakewakaran\\
\textbf{Expression (JA)}: わけわからん\\
\textbf{Romanization}: Wakewakaran\\
\textbf{Expression (EN)}: I don't know the reason

\textbf{Variants (JA)}: \\
\textbf{Variants (EN)}: 

\textbf{Intent (JA)}: \\
\textbf{Intent (EN)}: 

\textbf{Tags (JA)}: \\
\textbf{Tags (EN)}: 

\textbf{Adjusted Expression (JA)}:\\
\\
\textbf{Adjusted Expression (EN)}:\\


\textbf{Example (JA)}:\\
この(数学:すうがく)の(問題:もんだい)、(公式:こうしき)がたくさん(出:で)てきてわけわからん！\\
\textbf{Example (EN)}:\\
There are so many formulas in this math problem that I don't know what's going on!

\textbf{Notes (JA)}:\\
「わけがわからない」は、何かの理由や原因が理解できないことを表現するためのフレーズです。「わけ」は「reason」や「cause」を、「が」は文の主語を示す助詞で、「わからない」は動詞「わかる」（to understand）の否定形です。

\textbf{Notes (EN)}:\\
"Wake ga wakaranai" is a phrase used to express that one doesn't understand the reason or cause of something. The word "wake" means "reason" or "cause," "ga" is a particle that marks the subject of the sentence, and "wakaranai" is the negative form of the verb "wakaru" (to understand). 

\newpage

\section*{286. I couldn't do anything}

\textbf{Date}: 2025.01.21 (Tue)\\
\textbf{Index}: nanimodekinakatta\\
\textbf{Expression (JA)}: (何:なに)もできなかった\\
\textbf{Romanization}: Nani mo dekinakatta\\
\textbf{Expression (EN)}: I couldn't do anything

\textbf{Variants (JA)}: \\
\textbf{Variants (EN)}: 

\textbf{Intent (JA)}: \\
\textbf{Intent (EN)}: 

\textbf{Tags (JA)}: \\
\textbf{Tags (EN)}: 

\textbf{Adjusted Expression (JA)}:\\
\\
\textbf{Adjusted Expression (EN)}:\\


\textbf{Example (JA)}:\\
(今日:きょう)(忙:いそが)しくて、(何:なに)もできなかった。\\
\textbf{Example (EN)}:\\
I was so busy today that I couldn't do anything.

\textbf{Notes (JA)}:\\
「何もできなかった」は、忙しかったり、能力がなかったり、障害に直面していたりするなど、さまざまな理由で何もできなかったことを表現するためのフレーズです。「何」は「what」を、「も」は「anything」を意味し、「できなかった」は動詞「できる」（to be able to do）の過去否定形です。

\textbf{Notes (EN)}:\\
"Nani mo dekinakatta" is a phrase used to express that one was unable to do anything due to various reasons such as being busy, lacking the ability, or facing obstacles. The word "nani" means "what," "mo" is a particle that means "anything," "dekinakatta" is the past negative form of the verb "dekiru" (to be able to do).

\newpage

\section*{285. It'll be all right}

\textbf{Date}: 2025.01.20 (Mon)\\
\textbf{Index}: kittodaijoubu\\
\textbf{Expression (JA)}: きっと(大丈夫:だいじょうぶ)\\
\textbf{Romanization}: Kitto daijoubu\\
\textbf{Expression (EN)}: It'll be all right

\textbf{Variants (JA)}: \\
\textbf{Variants (EN)}: 

\textbf{Intent (JA)}: \\
\textbf{Intent (EN)}: 

\textbf{Tags (JA)}: \\
\textbf{Tags (EN)}: 

\textbf{Adjusted Expression (JA)}:\\
\\
\textbf{Adjusted Expression (EN)}:\\


\textbf{Example (JA)}:\\
A「(大丈夫:だいじょうぶ)かな？」B「きっと(大丈夫:だいじょうぶ)だよ」\\
\textbf{Example (EN)}:\\
A: I wonder if it's okay? B: I'm sure it'll be fine.

\textbf{Notes (JA)}:\\
「大丈夫」は、すべてが大丈夫、問題ない、安全であることを表現するためにさまざまな状況で使われる日本語の汎用語です。誰かを安心させるために使われたり、何かが受け入れられることを確認したり、問題がないことを示したりするために使われます。

\textbf{Notes (EN)}:\\
"Daijoubu" is a versatile Japanese word that can be used in various situations to express that everything is okay, fine, or safe. It can be used to reassure someone, confirm that something is acceptable, or indicate that there is no problem. 

\newpage

\section*{284. If you think you can do it, give it a try}

\textbf{Date}: 2025.01.19 (Sun)\\
\textbf{Index}: yarerumononarayattemiro\\
\textbf{Expression (JA)}: やれるものなら、やってみろ。\\
\textbf{Romanization}: Yareru mono nara, yatte miro\\
\textbf{Expression (EN)}: If you think you can do it, give it a try

\textbf{Variants (JA)}: \\
\textbf{Variants (EN)}: 

\textbf{Intent (JA)}: \\
\textbf{Intent (EN)}: 

\textbf{Tags (JA)}: \\
\textbf{Tags (EN)}: 

\textbf{Adjusted Expression (JA)}:\\
\\
\textbf{Adjusted Expression (EN)}:\\


\textbf{Example (JA)}:\\
A「(今日:きょう)こそ(俺:おれ)たちが(勝:か)つからな！」B「ふん、やれるものなら、やってみろ！」\\
\textbf{Example (EN)}:\\
A: Today, we're going to win! B: Hmph, if you think you can do it, give it a try!

\textbf{Notes (JA)}:\\
「やれるものなら、やってみろ」は、自分がそれをできると思うなら、行動を起こしたり、試みたりすることを奨励するためのフレーズです。「やれる」は「to be able to do」を、「もの」は「thing」を、「なら」は「if」を意味し、「やってみろ」は「give it a try」を意味します。この表現は、誰かが自分の能力やスキルを試すために誰かを励ましたり、挑戦したりするためによく使われます。

\textbf{Notes (EN)}:\\
"Yareru mono nara, yatte miro" is a phrase used to encourage someone to take action or make an attempt if they believe they are capable of doing so. The word "yareru" means "to be able to do," "mono" means "thing," "nara" is a conditional form that means "if," and "yatte miro" is a phrase that means "give it a try." This expression is commonly used to motivate or challenge someone to test their abilities or skills.

\newpage

\section*{283. Is that on purpose?}

\textbf{Date}: 2025.01.18 (Sat)\\
\textbf{Index}: wazatodesuka\\
\textbf{Expression (JA)}: わざとですか。\\
\textbf{Romanization}: Wazato desuka\\
\textbf{Expression (EN)}: Is that on purpose?

\textbf{Variants (JA)}: \\
\textbf{Variants (EN)}: 

\textbf{Intent (JA)}: \\
\textbf{Intent (EN)}: 

\textbf{Tags (JA)}: \\
\textbf{Tags (EN)}: 

\textbf{Adjusted Expression (JA)}:\\
\\
\textbf{Adjusted Expression (EN)}:\\


\textbf{Example (JA)}:\\
A「プリン、(間違:まちが)えて(食:た)べちゃった」B「それ、わざとですか？」\\
\textbf{Example (EN)}:\\
A: I accidentally ate the pudding. B: Did you do that on purpose?

\textbf{Notes (JA)}:\\
「わざとですか」は、誰かが意図的に何かをしたかどうかを尋ねるためのフレーズです。「わざと」は「on purpose」や「intentionally」を、「ですか」は丁寧な終わり方で質問を形成します。この表現は、行動や言葉が怪しいか異常な場合に、相手の意図や動機を明確にするためによく使われます。こういう時は、だいたいわざとですね。自分で買ってもいないのに、普通、他人のプリンを間違えて食べません。

\textbf{Notes (EN)}:\\
"Wazato desuka" is a phrase used to ask if someone did something intentionally or deliberately. The word "wazato" means "on purpose" or "intentionally," and "desuka" is a polite ending that forms a question. This expression is commonly used to clarify someone's intentions or motives when their actions or words seem suspicious or unusual. In most cases, it's on purpose. Normally, you wouldn't eat someone else's pudding by mistake if you didn't buy it yourself.

\newpage

\section*{282. I'll leave it up to you}

\textbf{Date}: 2025.01.17 (Fri)\\
\textbf{Index}: omakaseshimasu\\
\textbf{Expression (JA)}: おまかせします。\\
\textbf{Romanization}: Omakase shimasu\\
\textbf{Expression (EN)}: I'll leave it up to you

\textbf{Variants (JA)}: \\
\textbf{Variants (EN)}: 

\textbf{Intent (JA)}: \\
\textbf{Intent (EN)}: 

\textbf{Tags (JA)}: \\
\textbf{Tags (EN)}: 

\textbf{Adjusted Expression (JA)}:\\
\\
\textbf{Adjusted Expression (EN)}:\\


\textbf{Example (JA)}:\\
A「(今日:きょう)の(晩御飯:ばんごはん)、(何:なに)にしようかな」B「おまかせします」\\
\textbf{Example (EN)}:\\
A: What should we have for dinner tonight? B: I'll leave it up to you.

\textbf{Notes (JA)}:\\
「おまかせします」は、他の人に決定やタスクを任せることを表現するためのフレーズです。「おまかせ」は「to entrust」や「to leave it up to someone」を意味し、「します」は動詞「する」（to do）の丁寧形です。この表現は、責任や意思決定を他の人に委ねるためによく使われ、自信がない場合や好みがない場合によく使われます。

\textbf{Notes (EN)}:\\
"Omakase shimasu" is a phrase used to express that one is entrusting a decision or task to someone else. The word "omakase" means "to entrust" or "to leave it up to someone," and "shimasu" is the polite form of the verb "suru" (to do). This expression is commonly used to delegate responsibility or decision-making to another person and is often used in situations where one is unsure or has no preference.

\newpage

\section*{281. Either is fine}

\textbf{Date}: 2025.01.16 (Thu)\\
\textbf{Index}: docchidemoii\\
\textbf{Expression (JA)}: どっちでもいい\\
\textbf{Romanization}: Docchi demo ii\\
\textbf{Expression (EN)}: Either is fine

\textbf{Variants (JA)}: \\
\textbf{Variants (EN)}: 

\textbf{Intent (JA)}: \\
\textbf{Intent (EN)}: 

\textbf{Tags (JA)}: \\
\textbf{Tags (EN)}: 

\textbf{Adjusted Expression (JA)}:\\
\\
\textbf{Adjusted Expression (EN)}:\\


\textbf{Example (JA)}:\\
A「(寿司:すし)にしますか、(中華:ちゅうか)にしますか」B「どっちでもいいですよ」\\
\textbf{Example (EN)}:\\
A: Would you like sushi or Chinese food? B: Either is fine.

\textbf{Notes (JA)}:\\
「どっちでもいい」は、どちらの選択肢も受け入れ可能または適していることを表現するためのフレーズです。「どっち」は「which」や「either」を、「でも」は「or」を意味し、「いい」は「good」や「fine」を意味します。この表現は、2つ以上の選択肢の間に好みがないことを示したり、両方の選択肢が同じくらい受け入れられることを示すためによく使われます。

\textbf{Notes (EN)}:\\
"Docchi demo ii" is a phrase used to express that either option is acceptable or suitable. The word "docchi" means "which" or "either," "demo" is a particle that means "or," and "ii" means "good" or "fine." This expression is commonly used to indicate that one has no preference between two or more options or that both options are equally acceptable.

\newpage

\section*{280. Let's give it a try}

\textbf{Date}: 2025.01.15 (Wed)\\
\textbf{Index}: toriaezuyattemimashou\\
\textbf{Expression (JA)}: とりあえずやってみましょう。\\
\textbf{Romanization}: Toriaezu yatte mimashou\\
\textbf{Expression (EN)}: Let's just try it for now

\textbf{Variants (JA)}: \\
\textbf{Variants (EN)}: 

\textbf{Intent (JA)}: \\
\textbf{Intent (EN)}: 

\textbf{Tags (JA)}: \\
\textbf{Tags (EN)}: 

\textbf{Adjusted Expression (JA)}:\\
\\
\textbf{Adjusted Expression (EN)}:\\


\textbf{Example (JA)}:\\
A「この(料理:りょうり)、(難:むずか)しそうですね」B「とりあえず、(作:つく)ってみましょうか」\\
\textbf{Example (EN)}:\\
A: This dish looks difficult. B: Let's just try making it for now.

\textbf{Notes (JA)}:\\
「とりあえずやってみましょう」は、何かを最初の一歩として、結果を気にせずに一時的に試してみることを提案するためのフレーズです。「とりあえず」は「for now」や「first of all」を、「やって」は動詞「やる」（to do）のテ形を、「みましょう」は動詞「見る」（to see or try）の意向形を意味します。この表現は、行動を起こすことやためらいや過度な考え込みをせずに試みることを奨励するためによく使われます。

\textbf{Notes (EN)}:\\
"Toriaezu yatte mimashou" is a phrase used to suggest trying something as a first step or temporarily without worrying about the outcome. The word "toriaezu" means "for now" or "first of all," "yatte" is the te-form of the verb "yaru" (to do), and "mimashou" is the volitional form of the verb "miru" (to see or try). This expression is commonly used to encourage taking action or making an attempt without overthinking or hesitating.

\newpage

\section*{279. Does it suit me?}

\textbf{Date}: 2025.01.14 (Tue)\\
\textbf{Index}: niatteimasuka\\
\textbf{Expression (JA)}: (似合:にあ)っていますか。\\
\textbf{Romanization}: Niatteimasuka\\
\textbf{Expression (EN)}: Does it suit me?

\textbf{Variants (JA)}: \\
\textbf{Variants (EN)}: 

\textbf{Intent (JA)}: \\
\textbf{Intent (EN)}: 

\textbf{Tags (JA)}: \\
\textbf{Tags (EN)}: 

\textbf{Adjusted Expression (JA)}:\\
\\
\textbf{Adjusted Expression (EN)}:\\


\textbf{Example (JA)}:\\
この(服:ふく)、(似合:にあ)っていますか。\\
\textbf{Example (EN)}:\\
Does this outfit suit me?

\textbf{Notes (JA)}:\\
「この服、似合っていますか」は、何かがよく見えるか、誰かに合っているかを尋ねるためのフレーズです。「似合って」は動詞「似合う」（似合う、よく見える）のテ形で、「いますか」は現在の状態を質問しています。この表現は、服、アクセサリー、髪型を試す際に、自分がどう見えるかについて他人の意見を尋ねるためによく使われます。

\textbf{Notes (EN)}:\\
"Niatte imasuka" is a phrase used to ask if something looks good or suits someone. "Niau" is the verb "to suit" or "to look good," and "imasuka" is a question form indicating the current state. This expression is commonly used to ask for others' opinions on how one looks when trying on clothes, accessories, or hairstyles.

\newpage

\section*{278. About how much is it?}

\textbf{Date}: 2025.01.13 (Mon)\\
\textbf{Index}: daitaiikuradesuka\\
\textbf{Expression (JA)}: だいたいいくらですか。\\
\textbf{Romanization}: Daitai ikura desuka\\
\textbf{Expression (EN)}: About how much is it?

\textbf{Variants (JA)}: \\
\textbf{Variants (EN)}: 

\textbf{Intent (JA)}: \\
\textbf{Intent (EN)}: 

\textbf{Tags (JA)}: \\
\textbf{Tags (EN)}: 

\textbf{Adjusted Expression (JA)}:\\
\\
\textbf{Adjusted Expression (EN)}:\\


\textbf{Example (JA)}:\\
(参加費:さんかひ)はだいたいいくらですか？\\
\textbf{Example (EN)}:\\
About how much is the participation fee?

\textbf{Notes (JA)}:\\
「だいたいいくらですか」は、何かのおおよその価格や費用を尋ねるためのフレーズです。「だいたい」は「about」や「approximately」を、「いくら」は「how much」を、「ですか」は丁寧な終わり方で質問を形成します。

\textbf{Notes (EN)}:\\
"Daitai ikura desuka" is a phrase used to ask for an approximate or general price or cost of something. The word "daitai" means "about" or "approximately," "ikura" means "how much," and "desuka" is a polite ending that forms a question.

\newpage

\section*{277. Sorry for the trouble}

\textbf{Date}: 2025.01.12 (Sun)\\
\textbf{Index}: gomeiwaku\\
\textbf{Expression (JA)}: ご(迷惑:めいわく)\\
\textbf{Romanization}: Gomeiwaku\\
\textbf{Expression (EN)}: Trouble

\textbf{Variants (JA)}: \\
\textbf{Variants (EN)}: 

\textbf{Intent (JA)}: \\
\textbf{Intent (EN)}: 

\textbf{Tags (JA)}: \\
\textbf{Tags (EN)}: 

\textbf{Adjusted Expression (JA)}:\\
\\
\textbf{Adjusted Expression (EN)}:\\


\textbf{Example (JA)}:\\
しばらく、ご(連絡:れんらく)、できなく、ごめいわくをおかけしました。\\
\textbf{Example (EN)}:\\
I'm sorry for the trouble I caused by not being able to contact you for a while.

\textbf{Notes (JA)}:\\
「ご迷惑」は、誰かに不便やトラブル、迷惑をかけたことに対する謝罪や後悔を表現するためのフレーズです。「ご」は尊敬やフォーマルな感覚を加える接頭辞で、「迷惑」は「trouble」や「annoyance」を意味します。

\textbf{Notes (EN)}:\\
"Gomeiwaku" is a phrase used to express an apology or regret for causing inconvenience, trouble, or annoyance to someone. The word "go" is a prefix that adds a sense of respect or formality, and "meiwaku" means "trouble" or "annoyance."

\newpage

\section*{276. Gets on my nerves}

\textbf{Date}: 2025.01.11 (Sat)\\
\textbf{Index}: gamandekinai\\
\textbf{Expression (JA)}: がまんできない\\
\textbf{Romanization}: Gamanndekinai\\
\textbf{Expression (EN)}: I can't stand it

\textbf{Variants (JA)}: \\
\textbf{Variants (EN)}: 

\textbf{Intent (JA)}: \\
\textbf{Intent (EN)}: 

\textbf{Tags (JA)}: \\
\textbf{Tags (EN)}: 

\textbf{Adjusted Expression (JA)}:\\
\\
\textbf{Adjusted Expression (EN)}:\\


\textbf{Example (JA)}:\\
あの(芸能人:げいのうじん)の(話:はな)し(方:かた)はがまんできない。\\
\textbf{Example (EN)}:\\
I can't stand the way that celebrity talks.

\textbf{Notes (JA)}:\\
「我慢できない」は、特定の状況、行動、または人を我慢できないことを表現するためのフレーズです。「我慢」は「patience」や「endurance」を意味し、「できない」は動詞「できる」（to be able to do）の否定形です。このフレーズは、耐えがたいものや我慢できないものに直面した際に、不快感、いらだち、あるいは憤りを表現するためによく使われます。

\textbf{Notes (EN)}:\\
"Gaman dekinai" is a phrase used to express that one cannot tolerate or endure a particular situation, behavior, or person. The word "gaman" means "patience" or "endurance," and "dekinai" is the negative form of the verb "dekiru" (to be able to do). This phrase is commonly used to convey a sense of frustration, annoyance, or exasperation when faced with something that is difficult to bear or put up with.

\newpage

\section*{275. There's nothing I can do}

\textbf{Date}: 2025.01.10 (Fri)\\
\textbf{Index}: doushiyounomonai\\
\textbf{Expression (JA)}: どうしようもない\\
\textbf{Romanization}: Doushiyou mo nai\\
\textbf{Expression (EN)}: Hopeless

\textbf{Variants (JA)}: \\
\textbf{Variants (EN)}: 

\textbf{Intent (JA)}: \\
\textbf{Intent (EN)}: 

\textbf{Tags (JA)}: \\
\textbf{Tags (EN)}: 

\textbf{Adjusted Expression (JA)}:\\
\\
\textbf{Adjusted Expression (EN)}:\\


\textbf{Example (JA)}:\\
(電車:でんしゃ)が(止:と)まっている(以上:いじょう)、どうしようもないですね。\\
\textbf{Example (EN)}:\\
Since the train has stopped, there's nothing we can do.

\textbf{Notes (JA)}:\\
「どうしようもない」は、解決策や治療法がない状況を表現するためのフレーズです。「どう」は「how」を、「しよう」は動詞「する」（to do）の意向形です。状況に使えるだけでなく「あの人はどうしようもない」のように人にも使えます。

\textbf{Notes (EN)}:\\
"Doushiyou mo nai" is a phrase used to express a situation where there is no solution or remedy available. The word "dou" means "how," "shiyou" is the volitional form of the verb "suru" (to do). The phrase can be used to describe a situation where there is no solution or treatment available, and it can also be used to describe a person as in "ano hito wa doushiyou mo nai" (That person is hopeless).

\newpage

\section*{274. I have no idea}

\textbf{Date}: 2025.01.09 (Thu)\\
\textbf{Index}: mattakuwakarimasen\\
\textbf{Expression (JA)}: まったくわかりません。\\
\textbf{Romanization}: Mattaku wakarimasen\\
\textbf{Expression (EN)}: I have no idea

\textbf{Variants (JA)}: \\
\textbf{Variants (EN)}: 

\textbf{Intent (JA)}: \\
\textbf{Intent (EN)}: 

\textbf{Tags (JA)}: \\
\textbf{Tags (EN)}: 

\textbf{Adjusted Expression (JA)}:\\
\\
\textbf{Adjusted Expression (EN)}:\\


\textbf{Example (JA)}:\\
A「(問題:もんだい)、(解:と)けそうですか？」B「まったくわかりません。どこから(始:はじ)めればいいのか」\\
\textbf{Example (EN)}:\\
A: Do you think you can solve the problem? B: I have no idea. I don't know where to start.

\textbf{Notes (JA)}:\\
「まったくわかりません」は、理解が全然できないがくないことを表現するフレーズです。「まったく」は「completely」や「totally」を意味し、「わかりません」は動詞「わかる」（理解する）の否定形で、「I don't understand」を意味します。

\textbf{Notes (EN)}:\\
"Mattaku wakarimasen" is a phrase used to express complete lack of understanding or knowledge about a particular topic or situation. The word "mattaku" means "completely" or "totally," "wakarimasen" is the negative form of the verb "wakaru" (to understand), and it means "I don't understand."

\newpage

\section*{273. Do you think so?}

\textbf{Date}: 2025.01.08 (Wed)\\
\textbf{Index}: souomouno\\
\textbf{Expression (JA)}: そう(思:おも)うの？\\
\textbf{Romanization}: Sou omou no?\\
\textbf{Expression (EN)}: Do you think so?

\textbf{Variants (JA)}: \\
\textbf{Variants (EN)}: 

\textbf{Intent (JA)}: \\
\textbf{Intent (EN)}: 

\textbf{Tags (JA)}: \\
\textbf{Tags (EN)}: 

\textbf{Adjusted Expression (JA)}:\\
\\
\textbf{Adjusted Expression (EN)}:\\


\textbf{Example (JA)}:\\
A「この(映画:えいが)、(結構:けっこう)(面白:おもしろ)かったよね」B「そう(思:おも)うの？ (私:わたし)はちょっと(退屈:たいくつ)だったな」\\
\textbf{Example (EN)}:\\
A: This movie was pretty interesting, wasn't it? B: Do you think so? I found it a bit boring.

\textbf{Notes (JA)}:\\
「そう思うの？」は、相手の意見や発言に確認を求めるためのフレーズです。「そう」は「that way」を、「思う」は「to think」を意味し、「の」は確認や同意を求める文末の助詞です。この表現は、異なる意見についてさらに議論を促したりするために、会話でよく使われます。

\textbf{Notes (EN)}:\\
"Sou omou no?" is a phrase used to ask for confirmation with someone's opinion or statement. The word "sou" means "that way." "omou" is the verb "to think," and "no" is a sentence-ending particle that seeks confirmation or agreement. This expression is commonly used in conversations to check if the listener shares the speaker's perspective or to prompt further discussion about differing opinions.

\newpage

\section*{272. A little too much}

\textbf{Date}: 2025.01.07 (Tue)\\
\textbf{Index}: chottoiisugi\\
\textbf{Expression (JA)}: ちょっと(言:い)いすぎ\\
\textbf{Romanization}: Chotto ii sugi\\
\textbf{Expression (EN)}: A little too much

\textbf{Variants (JA)}: \\
\textbf{Variants (EN)}: 

\textbf{Intent (JA)}: \\
\textbf{Intent (EN)}: 

\textbf{Tags (JA)}: \\
\textbf{Tags (EN)}: 

\textbf{Adjusted Expression (JA)}:\\
\\
\textbf{Adjusted Expression (EN)}:\\


\textbf{Example (JA)}:\\
A「あなたって、いつも何も考えてないよね」B「え、それはちょっと言いすぎじゃない？ちゃんと考えてるよ」\\
\textbf{Example (EN)}:\\
A: You never think about anything, do you? B: Huh, isn't that a little too much? I do think things through.

\textbf{Notes (JA)}:\\
「ちょっと言い過ぎ」は、発言や行動がやや誇張されているか、あるいは過度に過激であることを表現するためのフレーズです。「ちょっと」は"a little"や"slightly"を意味し、「言い過ぎ」は「言う」のマス形に「過ぎる」のマス形を付けた名詞で、"too much to say"という意味になります。この表現は悪い評価でも良い評価でも使えます。

\textbf{Notes (EN)}:\\
"Chotto ii sugi" is an expression used to indicate that a statement or action is slightly exaggerated or overly extreme. The word "chotto" means "a little" or "slightly," "iisugi" is the masu-form of the verb "iu (to say)" with the  masu-form of "sugiru (to exceed)," meaning "too much to say." It can be used in both negative and positive contexts.

\newpage

\section*{271. I don't mind at all}

\textbf{Date}: 2025.01.06 (Mon)\\
\textbf{Index}: mattakukininaranai\\
\textbf{Expression (JA)}: まったく(気:き)にならない\\
\textbf{Romanization}: Mattaku ki ni naranai\\
\textbf{Expression (EN)}: I don't mind at all

\textbf{Variants (JA)}: \\
\textbf{Variants (EN)}: 

\textbf{Intent (JA)}: \\
\textbf{Intent (EN)}: 

\textbf{Tags (JA)}: \\
\textbf{Tags (EN)}: 

\textbf{Adjusted Expression (JA)}:\\
\\
\textbf{Adjusted Expression (EN)}:\\


\textbf{Example (JA)}:\\
A「ちょっと(髪型:かみがた)(失敗:しっぱい)しちゃったかも」B「いやいや、(全然:ぜんぜん)(似合:にあ)ってるし、まったく(気:き)にならないよ」\\
\textbf{Example (EN)}:\\
A: I might have messed up my hairstyle a bit. B: No way, it looks great on you, I don't mind at all.

\textbf{Notes (JA)}:\\
「まったく気にならない」は、何かに気になったり、気にすることがないことを表現するためのフレーズです。「まったく」は「completely」や「totally」を意味し、「気にならない」は「not to mind」や「not to be concerned」を意味します。本当は目立つほど失敗作だけれども、言わないだけかもしれません。「髪型失敗しちゃったかも」と言わなければ、気がつかなかったのに、逆に言ったので、気づいてしまうこともあります。だいたいは、他人は自分が思っているほど、他人のことを気にしていないものです。

\textbf{Notes (EN)}:\\
"Mattaku ki ni naranai" is a phrase used to express that one doesn't mind or is not bothered by something. The word "mattaku" means "completely" or "totally," "ki ni naranai" means "not to mind" or "not to be concerned." It may be a complete failure, but if you don't say it, it may not be noticed. If you say, "I might have messed up my hairstyle," others might notice it, In general, people don't care about others as much as they think they do.

\newpage

\section*{270. Anytime}

\textbf{Date}: 2025.01.05 (Sun)\\
\textbf{Index}: itsudemodouzo\\
\textbf{Expression (JA)}: いつでもどうぞ\\
\textbf{Romanization}: Itsudemo douzo\\
\textbf{Expression (EN)}: Anytime

\textbf{Variants (JA)}: \\
\textbf{Variants (EN)}: 

\textbf{Intent (JA)}: \\
\textbf{Intent (EN)}: 

\textbf{Tags (JA)}: \\
\textbf{Tags (EN)}: 

\textbf{Adjusted Expression (JA)}:\\
\\
\textbf{Adjusted Expression (EN)}:\\


\textbf{Example (JA)}:\\
(授業:じゅぎょう)について(質問:しつもん)がある(方:かた)は、いつでもどうぞ\\
\textbf{Example (EN)}:\\
If you have any questions about the class, feel free to ask anytime.

\textbf{Notes (JA)}:\\
「いつでもどうぞ」は、相手にいつでも行動を起こしたり、要求をしたりすることを奨励するためのフレーズです。「いつでも」は「anytime」や「whenever」を意味し、「どうぞ」は何かをするための招待や許可を表します。とはいえ、本当にアポイントメントも入れずにいつでもいいというわけではありません

\textbf{Notes (EN)}:\\
"Itsu demo douzo" is a phrase used to encourage someone to take action or make a request at any time. "Itsu demo" means "anytime" or "whenever," and "douzo" is an invitation or permission to do something. However, it doesn't mean you can do it anytime without making an appointment.

\newpage

\section*{269. I do remember it}

\textbf{Date}: 2025.01.04 (Sat)\\
\textbf{Index}: wasureruwakenai\\
\textbf{Expression (JA)}: (忘:わす)れるわけない\\
\textbf{Romanization}: Wasureru wake nai\\
\textbf{Expression (EN)}: I do remember it

\textbf{Variants (JA)}: \\
\textbf{Variants (EN)}: 

\textbf{Intent (JA)}: \\
\textbf{Intent (EN)}: 

\textbf{Tags (JA)}: \\
\textbf{Tags (EN)}: 

\textbf{Adjusted Expression (JA)}:\\
\\
\textbf{Adjusted Expression (EN)}:\\


\textbf{Example (JA)}:\\
A「(明日:あす)、(会議:かいぎ)だよね？」B「わすれるわけないでしょ！っていうか、もう(準備:じゅんび)できてるし」\\
\textbf{Example (EN)}:\\
A: We have a meeting tomorrow, right? B: There's no way I'd forget! Besides, I'm already prepared.

\textbf{Notes (JA)}:\\
「忘れるわけない」は、何かをしっかり覚えていて、忘れていないことを意味します。「忘れる」は「to forget」を意味し、「わけない」は理由や原因がないことを示すフレーズです。

\textbf{Notes (EN)}:\\
"Wasureru wake nai" is a phrase used to express certainty or confidence that one has not forgotten something. The word "wasureru" means "to forget," and "wake nai" is a phrase that indicates the absence of a reason or cause.

\newpage

\section*{268. Is it okay later?}

\textbf{Date}: 2025.01.03 (Fri)\\
\textbf{Index}: atodemoii\\
\textbf{Expression (JA)}: あとでもいい？\\
\textbf{Romanization}: Ato demo ii?\\
\textbf{Expression (EN)}: Is it okay later?

\textbf{Variants (JA)}: \\
\textbf{Variants (EN)}: 

\textbf{Intent (JA)}: \\
\textbf{Intent (EN)}: 

\textbf{Tags (JA)}: \\
\textbf{Tags (EN)}: 

\textbf{Adjusted Expression (JA)}:\\
\\
\textbf{Adjusted Expression (EN)}:\\


\textbf{Example (JA)}:\\
A「(食器:しょっき)(洗:あら)ってくれた？」B「うーん、あとでもいい？(今:いま)、(犬:いぬ)と(遊:あそ)んでるところなんだ」\\
\textbf{Example (EN)}:\\
A: Did you wash the dishes? B: Hmm, is it okay later? I'm playing with the dog right now.

\textbf{Notes (JA)}:\\
「あとでもいい？」は、後でまた別の時間にできるかどうかを尋ねるためのフレーズです。「あと」は「later」や「after」を意味します。この表現は、タスクを遅らせたり予定を変更したりしても不便や混乱を引き起こさないかどうかを確認するために、カジュアルな会話でよく使われます。とはいえ、こういう人はだいたいあとになってもしないことが多いですね。

\textbf{Notes (EN)}:\\
"Ato demo ii?" is a phrase used to ask if something can be done at a different time. "Ato" means "later" or "after," and "demo" is a particle that indicates a conditional or alternative situation. This expression is commonly used in casual conversations to check if delaying a task or changing a schedule would cause inconvenience or confusion. However, people who say this often don't do it later.

\newpage

\section*{267. Better than nothing}

\textbf{Date}: 2025.01.02 (Thu)\\
\textbf{Index}: shinaiyorimashi\\
\textbf{Expression (JA)}: しないよりまし\\
\textbf{Romanization}: Shinai yori mashi\\
\textbf{Expression (EN)}: Better than nothing

\textbf{Variants (JA)}: \\
\textbf{Variants (EN)}: 

\textbf{Intent (JA)}: \\
\textbf{Intent (EN)}: 

\textbf{Tags (JA)}: \\
\textbf{Tags (EN)}: 

\textbf{Adjusted Expression (JA)}:\\
\\
\textbf{Adjusted Expression (EN)}:\\


\textbf{Example (JA)}:\\
A「(今:いま)から(試験:しけん)(勉強:べんきょう)ですか？」B「ええ、しないよりまし」\\
\textbf{Example (EN)}:\\
A: Are you going to study for the exam now? B: Yes, better than nothing.

\textbf{Notes (JA)}:\\
「しないよりまし」は、何もしないよりも何かをする方が良いということを表現するためのフレーズです。「しない」は動詞「する」（to do）の否定形で、「まし」は形容詞「いい」（good）の比較形です。大方において、こういう人はしっかり準備をしているのかもしれません。

\textbf{Notes (EN)}:\\
"Shinai yori mashi" is a phrase used to express that doing something is better than doing nothing at all. The word "shinai" is the negative form of the verb "suru" (to do), and "mashi" is the comparative form of the adjective "ii" (good). It conveys the idea that taking action is preferable to inaction. In many cases, such people may be well-prepared.

\newpage

\section*{266. My pleasure}

\textbf{Date}: 2025.01.01 (Wed)\\
\textbf{Index}: yorokonde\\
\textbf{Expression (JA)}: よろこんで\\
\textbf{Romanization}: Yorokonde\\
\textbf{Expression (EN)}: My pleasure

\textbf{Variants (JA)}: \\
\textbf{Variants (EN)}: 

\textbf{Intent (JA)}: \\
\textbf{Intent (EN)}: 

\textbf{Tags (JA)}: \\
\textbf{Tags (EN)}: 

\textbf{Adjusted Expression (JA)}:\\
\\
\textbf{Adjusted Expression (EN)}:\\


\textbf{Example (JA)}:\\
A「(手伝:てつだ)ってくれませんか？」B「よろこんで」\\
\textbf{Example (EN)}:\\
A: Can you help me? B: My pleasure.

\textbf{Notes (JA)}:\\
「よろこんで」は、誰かを手伝ったり助けたりすることに対する意欲や喜びを表現するためのフレーズです。「よろこんで」は動詞「よろこぶ」（喜ぶ、楽しむ）のテ形で、この文脈で使われると、他人を助けることに喜びや満足感を表現します。しかし、ちょっと忙しく手伝っていられない場合、喜んで手伝える場合ではなくてもこのように言うことがあります。一瞬手伝うことで終りにできるなら、少しぐらいの時間の貢献は止むなしと言ったところでしょうか。忙しく仕事をしているのに頼んでくる人がいるならその人の問題です。

\textbf{Notes (EN)}:\\
"Yorokonde" is a phrase used to express willingness or joy in helping or assisting someone. It is the te-form of the verb "yorokobu" (to be pleased or delighted) and is used in this context to convey a sense of joy or satisfaction in helping others. However, it can also be used in situations where one may not be able to help with joy but is willing to assist. If the help can be done in a moment and is unavoidable, one might say this. If someone is asking for help while you are busy working, that is their problem.

\newpage

\section*{265. I brought everything}

\textbf{Date}: 2024.12.31 (Tue)\\
\textbf{Index}: zenbumottara\\
\textbf{Expression (JA)}: ぜんぶもった。\\
\textbf{Romanization}: Zenbu motta\\
\textbf{Expression (EN)}: I brought everything

\textbf{Variants (JA)}: \\
\textbf{Variants (EN)}: 

\textbf{Intent (JA)}: \\
\textbf{Intent (EN)}: 

\textbf{Tags (JA)}: \\
\textbf{Tags (EN)}: 

\textbf{Adjusted Expression (JA)}:\\
\\
\textbf{Adjusted Expression (EN)}:\\


\textbf{Example (JA)}:\\
A「(忘:わす)れ(物:もの)、ない？」B「ぜんぶもった」\\
\textbf{Example (EN)}:\\
A: Got everything? B: I brought everything.

\textbf{Notes (JA)}:\\
「全部持った」は、必要なものや必要なものをすべて持ってきたことを示すためのフレーズです。「全部」は「everything」や「all」を意味し、「持った」は動詞「持ってる」（持つ、持っている）の過去形です。日本語では過去形を使っても現在の状況を表現することができます。

\textbf{Notes (EN)}:\\
"Zenbu motta" is a phrase used to indicate that one has brought or obtained everything that is needed or required. "Zenbu" means "everything" or "all," and "motta" is the past tense of the verb "motteru" (to have or to bring). In Japanese, the past tense can be used to describe a current state or situation. This expression is commonly used to convey that one has all the necessary items or has completed a task by bringing everything required.

\newpage

\section*{264. Is it a big deal?}

\textbf{Date}: 2024.12.30 (Mon)\\
\textbf{Index}: sonnanisawagukotodesuka\\
\textbf{Expression (JA)}: そんなにさわぐことですか。\\
\textbf{Romanization}: Sonnani sawagu koto desuka\\
\textbf{Expression (EN)}: Is it a big deal?

\textbf{Variants (JA)}: \\
\textbf{Variants (EN)}: 

\textbf{Intent (JA)}: \\
\textbf{Intent (EN)}: 

\textbf{Tags (JA)}: \\
\textbf{Tags (EN)}: 

\textbf{Adjusted Expression (JA)}:\\
\\
\textbf{Adjusted Expression (EN)}:\\


\textbf{Example (JA)}:\\
A「(何:なに)か(問題:もんだい)が(起:お)きたのかな」B「そんなにさわぐことですか。ただ(疲:つか)れているだけです」\\
\textbf{Example (EN)}:\\
A: I wonder if something happened. B: Is it a big deal? I'm just tired.

\textbf{Notes (JA)}:\\
「そんなに騒ぐことですか」は、何かについて大騒ぎすることの重要性や必要性を問うためのフレーズです。「そんなに」は「to that extent」や「so much」を意味し、「騒ぐ」は「to make noise」や「to cause a disturbance」を意味します。これらの要素を組み合わせることで、過度の非難や他者の主張の言い過ぎを疑問視している感覚を伝えます。

\textbf{Notes (EN)}:\\
"Sonnani sawagu koto desuka" is a phrase used to question the importance or necessity of making a big fuss about something. "Sonnani" means "to that extent" or "so much," and "sawagu" means "to make noise" or "to cause a disturbance." By combining these elements, the phrase conveys a sense of questioning the significance or need to make excessive noise or fuss about something. This expression is often used to express skepticism about excessive criticism or overblown claims made by others.

\newpage

\section*{263. That seems to be the case}

\textbf{Date}: 2024.12.29 (Sun)\\
\textbf{Index}: sonoyoudesune\\
\textbf{Expression (JA)}: そのようですね。\\
\textbf{Romanization}: Sonoyoudesune\\
\textbf{Expression (EN)}: That seems to be the case

\textbf{Variants (JA)}: \\
\textbf{Variants (EN)}: 

\textbf{Intent (JA)}: \\
\textbf{Intent (EN)}: 

\textbf{Tags (JA)}: \\
\textbf{Tags (EN)}: 

\textbf{Adjusted Expression (JA)}:\\
\\
\textbf{Adjusted Expression (EN)}:\\


\textbf{Example (JA)}:\\
A「この(カフェ:かふぇ)、(最近:さいきん)(人気:にんき)みたいですね」B「そのようですね、いつも(混:こ)んでいる(気:き)がします」\\
\textbf{Example (EN)}:\\
A: This cafe seems to be popular lately. B: That seems to be the case; it always seems crowded.

\textbf{Notes (JA)}:\\
「そのようですね」は、他の人が述べた発言を認めたり同意したりするためのフレーズです。話し手の観察、意見、評価に同意や確認を表現する際によく使われます。「そのよう」は「that way」や「like that」を意味し、「ですね」は同意や肯定の意味を加える丁寧な終わり方です。これらの要素を組み合わせることで、他の人が説明した状況を理解したり認識したりする感覚を伝えます。この表現は、会話で話し手との注意深さ、同意、共有された認識を示すためによく使われます。

\textbf{Notes (EN)}:\\
"Sonoyoudesune" is a phrase used to acknowledge or agree with a statement made by someone else. It is often used to express agreement or confirmation with the speaker's observation, opinion, or assessment. The phrase "sono you" means "that way" or "like that," and "desu ne" is a polite ending that adds a sense of agreement or affirmation. By combining these elements, the phrase conveys a sense of understanding or recognition of the situation described by the other person. This expression is commonly used in conversations to show attentiveness, agreement, or shared perception with the speaker.

\newpage

\section*{262. Indescribable}

\textbf{Date}: 2024.12.28 (Sat)\\
\textbf{Index}: nantomoienai\\
\textbf{Expression (JA)}: なんともいえない\\
\textbf{Romanization}: Nantomo ienai\\
\textbf{Expression (EN)}: Indescribable

\textbf{Variants (JA)}: \\
\textbf{Variants (EN)}: 

\textbf{Intent (JA)}: \\
\textbf{Intent (EN)}: 

\textbf{Tags (JA)}: \\
\textbf{Tags (EN)}: 

\textbf{Adjusted Expression (JA)}:\\
\\
\textbf{Adjusted Expression (EN)}:\\


\textbf{Example (JA)}:\\
その(映画:えいが)の(結末:けつまつ)はなんともいえない。\\
\textbf{Example (EN)}:\\
The ending of the movie is indescribable.

\textbf{Notes (JA)}:\\
「なんとも言えない」は、ことばで表現するのが難しいものや、説明できないものを表現するためのフレーズです。複雑な気持ちになっている時、意見を求められても言わない方がいい時に使えます。

\textbf{Notes (EN)}:\\
"Nantomo ienai" is a phrase used to describe something that is difficult to express in words or that defies explanation. It is often used when encountering complex emotions or situations that are hard to describe or when it is better not to say anything when asked for an opinion.

\newpage

\section*{261. Anxious}

\textbf{Date}: 2024.12.27 (Fri)\\
\textbf{Index}: soudatoiikedo\\
\textbf{Expression (JA)}: そうだといいけど、\\
\textbf{Romanization}: Souda to iikedo\\
\textbf{Expression (EN)}: Anxious

\textbf{Variants (JA)}: \\
\textbf{Variants (EN)}: 

\textbf{Intent (JA)}: \\
\textbf{Intent (EN)}: 

\textbf{Tags (JA)}: \\
\textbf{Tags (EN)}: 

\textbf{Adjusted Expression (JA)}:\\
\\
\textbf{Adjusted Expression (EN)}:\\


\textbf{Example (JA)}:\\
A「(試験:しけん)、うまくいくといいね」 B「そうだといいけど、(勉強:べんきょう)してないから(不安:ふあん)だよ」\\
\textbf{Example (EN)}:\\
A: I hope the exam goes well. B: I hope so too, but I'm anxious because I haven't studied.

\textbf{Notes (JA)}:\\
「そうだといいけど」は、ポジティブな結果に対する希望や願いを表現する一方で、不確実性や不安の可能性を認識するためのフレーズです。文字通り「そうだといい」という意味で、動詞「そうだ」（それはそうだ）の条件形「そうだといい」と「けど」（しかし）が結びついており、楽観主義と現実主義をバランスよく表現しています。このフレーズは、希望と懸念の両方を伝えるために会話でよく使われ、好ましい結果を望む一方で、疑念や不安の可能性を認識する話し手の気持ちを反映しています。

\textbf{Notes (EN)}:\\
"Souda to iikedo" is a phrase used to express hope or wish for a positive outcome while acknowledging the possibility of uncertainty or anxiety. Literally meaning "it would be good if that's the case," the phrase combines the conditional form "souda to ii" of the verb "sou da" (it is so) with "kedo" (but), creating a nuanced expression that balances optimism with realism. This phrase is often used in conversations to convey a mix of hope and concern, reflecting the speaker's desire for a favorable result while recognizing the potential for doubt or unease.

\newpage

\section*{260. Do over}

\textbf{Date}: 2024.12.26 (Thu)\\
\textbf{Index}: yarinaosu\\
\textbf{Expression (JA)}: (やり直:なお)す\\
\textbf{Romanization}: Yarinaosu\\
\textbf{Expression (EN)}: Do over

\textbf{Variants (JA)}: \\
\textbf{Variants (EN)}: 

\textbf{Intent (JA)}: \\
\textbf{Intent (EN)}: 

\textbf{Tags (JA)}: \\
\textbf{Tags (EN)}: 

\textbf{Adjusted Expression (JA)}:\\
\\
\textbf{Adjusted Expression (EN)}:\\


\textbf{Example (JA)}:\\
おでんを(作:つく)ってたら、なぜか(カレー:かれー)ができてた。(最初:さいしょ)からやり(直:なお)そう。\\
\textbf{Example (EN)}:\\
I was making oden, but somehow it turned into curry. Let's start over from scratch.

\textbf{Notes (JA)}:\\
「やり直す」は、動詞「やる」のます形「やり」と「直す」が結びついた複合動詞で、「to redo」や「to do over」を意味します。この語は、何かを初めからやり直したり、間違いを修正する際に用いられます。また、「やり直し」とすることで名詞化され、行為そのものを指すようになります。おでんは、大根や卵、こんにゃくなどを出汁で煮込んだ日本の伝統的な料理で、冬の定番として親しまれています。一方、カレーはスパイスを使った煮込み料理で、日本ではカレーライスとして広く食べられており、特有の甘辛いルウが特徴です。

\textbf{Notes (EN)}:\\
'Yarinaosu' is a compound verb formed by combining the masu-form 'yari' of the verb 'yaru' (to do) with 'naosu' (to fix), and it means 'to redo' or 'to do over.' It is used when starting something over from the beginning or correcting a mistake. Additionally, when nominalized as 'yarinaoshi,' it refers to the act of redoing itself. Oden is a traditional Japanese dish made by simmering ingredients like daikon radish, eggs, and konjac in a light dashi broth, and it is popular as a winter staple. On the other hand, curry is a spiced stew, widely enjoyed in Japan as curry rice, characterized by its uniquely sweet and savory roux. Thus, 'yarinaosu' conveys the idea of retrying after failure or making a fresh start.

\newpage

\section*{259. may not look it, but...}

\textbf{Date}: 2024.12.25 (Wed)\\
\textbf{Index}: aamiete\\
\textbf{Expression (JA)}: ああ(見:み)えて、...\\
\textbf{Romanization}: Aamiete\\
\textbf{Expression (EN)}: may not look it, but...

\textbf{Variants (JA)}: \\
\textbf{Variants (EN)}: 

\textbf{Intent (JA)}: \\
\textbf{Intent (EN)}: 

\textbf{Tags (JA)}: \\
\textbf{Tags (EN)}: 

\textbf{Adjusted Expression (JA)}:\\
\\
\textbf{Adjusted Expression (EN)}:\\


\textbf{Example (JA)}:\\
ああ(見:み)えて、(彼:かれ)は(会社:かいしゃ)で「Excelの(魔術師:まじゅつし)」って(呼:よ)ばれてるらしい。\\
\textbf{Example (EN)}:\\
He may not look it, but at work, they call him the Wizard of Excel.

\textbf{Notes (JA)}:\\
「ああ見えて」は、外見や行動からはすぐにはわからない、誰かや何かについての驚くべき事実や予想外の事実を紹介するためのフレーズです。「ああ見えて」は文字通り「見えないかもしれない」という意味で、最初の印象と矛盾するか、最初の印象に疑問を投げかけるような発言が続きます。

\textbf{Notes (EN)}:\\
"Aamiete" is a phrase used to introduce surprising or unexpected facts or revelations about someone or something that may not be immediately apparent from their appearance or behavior. Literally meaning "may not look like that," the phrase is followed by statements that contradict or cast doubt on the initial impression. This expression is often used to highlight hidden talents, unexpected skills, or surprising qualities that are not immediately obvious.

\newpage

\section*{258. What should I do?}

\textbf{Date}: 2024.12.24 (Tue)\\
\textbf{Index}: doushitara\\
\textbf{Expression (JA)}: どうしたらいいんでしょうかねぇ。\\
\textbf{Romanization}: Doushitaraiindeshoukanee.\\
\textbf{Expression (EN)}: What should I do?

\textbf{Variants (JA)}: \\
\textbf{Variants (EN)}: 

\textbf{Intent (JA)}: \\
\textbf{Intent (EN)}: 

\textbf{Tags (JA)}: \\
\textbf{Tags (EN)}: 

\textbf{Adjusted Expression (JA)}:\\
\\
\textbf{Adjusted Expression (EN)}:\\


\textbf{Example (JA)}:\\
友達の犬が、なぜか私にだけ敬語を使うんですけど、どうしたらいいんでしょうかねぇ。\\
\textbf{Example (EN)}:\\
My friend's dog, for some reason, only uses polite/honorific language with me. What should I do?

\textbf{Notes (JA)}:\\
「どうしたらいいんでしょうかねぇ」は、状況に対する不確実性を表現したり、どう対処すべきかについてアドバイスを求めたりするためのフレーズです。「どうしたら」は「どうすればいいか？」や「最善の行動は何か？」を意味し、「どうしたらいいんでしょうかねぇ」という形で使われます。特に、個人的な問題を話し合ったり、友人や家族からアドバイスを求めたりする際に、カジュアルな会話でよく使われます。

\textbf{Notes (EN)}:\\
"Doushitara ii n deshouka nee" is a phrase used to express uncertainty about a situation or to seek advice on how to handle it. "Doushitara" means "what should be done?" or "what is the best course of action?" When used in the full phrase "Doushitara ii n deshouka nee," it conveys a casual and slightly introspective tone. This expression is commonly used in informal conversations when discussing personal concerns or seeking advice from friends or family.

\newpage

\section*{257. Can't repay}

\textbf{Date}: 2024.12.23 (Mon)\\
\textbf{Index}: kaeshikirenai\\
\textbf{Expression (JA)}: (返:かえ)しきれない\\
\textbf{Romanization}: Kaeshikirenai\\
\textbf{Expression (EN)}: Can't repay

\textbf{Variants (JA)}: \\
\textbf{Variants (EN)}: 

\textbf{Intent (JA)}: \\
\textbf{Intent (EN)}: 

\textbf{Tags (JA)}: \\
\textbf{Tags (EN)}: 

\textbf{Adjusted Expression (JA)}:\\
\\
\textbf{Adjusted Expression (EN)}:\\


\textbf{Example (JA)}:\\
(高利貸:こうりが)しとは(知:し)らずに、お(金:かね)を借りて、(返:かえ)しても(返:かえ)しきれない(借金:しゃっきん)になってしまった。\\
\textbf{Example (EN)}:\\
I borrowed money without realizing it was from a loan shark, and now I'm stuck with debt I can never fully repay.

\textbf{Notes (JA)}:\\
「返しきれない」は、借金や義務を完全に返済できない状況を表現するフレーズです。この言葉は、「返す」（返済する、戻す）という動詞に、「切る」（完全に終える）の否定可能形「切れない」を組み合わせた複合動詞で、「返済を終えることができない」という意味を持ちます。このフレーズは、財政的や時間的な制約を暗示するほか、比喩的には、恩返しや償いができない状況を表現することもあります。たとえば、「この御恩は返しても返しきれません」のように、非常に手厚い恩を受けたときの返答として述べることができます。

\textbf{Notes (EN)}:\\
"Kaeshikirenai" describes a situation where one cannot fully repay a debt or obligation. The term combines "kaesu" (to return or repay) and the negative potential form of "kiru" (to cut or finish), conveying the inability to complete repayment. This phrase often implies financial or time-related constraints and can also be used metaphorically to express the inability to reciprocate kindness or make amends. For example, it can be used in responses to deep gratitude, such as "This kindness is something I can never fully repay."

\newpage

\section*{256. Somehow it will work out}

\textbf{Date}: 2024.12.22 (Sun)\\
\textbf{Index}: nankatokanaru\\
\textbf{Expression (JA)}: なんとかなる\\
\textbf{Romanization}: Nantoka naru\\
\textbf{Expression (EN)}: Somehow it will work out

\textbf{Variants (JA)}: \\
\textbf{Variants (EN)}: 

\textbf{Intent (JA)}: \\
\textbf{Intent (EN)}: 

\textbf{Tags (JA)}: \\
\textbf{Tags (EN)}: 

\textbf{Adjusted Expression (JA)}:\\
\\
\textbf{Adjusted Expression (EN)}:\\


\textbf{Example (JA)}:\\
(財布:さいふ)が(空:から)なのに「(明日:あす)はなんとかなる！」って(能天気:のうてんき)だな。\\
\textbf{Example (EN)}:\\
Even though his wallet is empty, he is so optimistic saying, Somehow it will work out tomorrow!

\textbf{Notes (JA)}:\\
「なんとかなる」は、困難な状況や不確実な状況でも、何とかうまくいくという楽観的な考えを表現するためのフレーズです。「なんとか」は「somehow」や「in some way」を意味し、「なる」は「to become」や「to work out」を意味します。これらの言葉を組み合わせることで、現在の課題や障害にもかかわらず、最終的には物事がうまくいくという自信を表現します。この表現は、自分自身や他者に、逆境や不確実性に直面しても前向きで希望を持ち続けるよう励ますためによく使われます。

\textbf{Notes (EN)}:\\
"Nantoka naru" is a phrase used to express optimism or hope that things will work out somehow, even in difficult or uncertain situations. The word "nantoka" means "somehow" or "in some way," and "naru" means "to become" or "to work out." By combining these words, the phrase conveys a sense of confidence that things will eventually fall into place, despite the current challenges or obstacles. This expression is often used to encourage oneself or others to stay positive and hopeful, even when facing adversity or uncertainty.

\newpage

\section*{255. If there is something}

\textbf{Date}: 2024.12.21 (Sat)\\
\textbf{Index}: nankaattara\\
\textbf{Expression (JA)}: なんかあったら\\
\textbf{Romanization}: Nanka attara\\
\textbf{Expression (EN)}: If there is something

\textbf{Variants (JA)}: \\
\textbf{Variants (EN)}: 

\textbf{Intent (JA)}: \\
\textbf{Intent (EN)}: 

\textbf{Tags (JA)}: \\
\textbf{Tags (EN)}: 

\textbf{Adjusted Expression (JA)}:\\
\\
\textbf{Adjusted Expression (EN)}:\\


\textbf{Example (JA)}:\\
なんかあったら、すぐに(連絡:れんらく)してください。\\
\textbf{Example (EN)}:\\
If there is something, please contact me immediately.

\textbf{Notes (JA)}:\\
「なんかあったら」は、相手が助けを必要とする場合に、手助けを申し出る際に使うカジュアルな表現です。「なんか」はsomethingやanythingを意味し、曖昧なものを指す際によく使われます。「あったら」を加えることで、必要が生じた場合に助ける意思を伝えます。友人を支えたり、緊急時にサポートを申し出たりする際に使われる便利なフレーズです。

\textbf{Notes (EN)}:\\
"Nanka attara" is a casual phrase used to offer help if the other person needs assistance. The word "nanka" means "something" or "anything" and is often used in informal contexts to refer to unspecified things. By adding "attara," the speaker expresses readiness to help if the need arises. This phrase is commonly used to show willingness to lend a hand in various situations, such as supporting a friend or responding to an emergency.

\newpage

\section*{254. Not particularly}

\textbf{Date}: 2024.12.20 (Fri)\\
\textbf{Index}: betsuni\\
\textbf{Expression (JA)}: (別:べつ)に(自慢:じまん)しているわけではないんですが、\\
\textbf{Romanization}: Betsuni jiman shite iru wake de wa nain desu ga,\\
\textbf{Expression (EN)}: Not particularly proud, but...

\textbf{Variants (JA)}: \\
\textbf{Variants (EN)}: 

\textbf{Intent (JA)}: \\
\textbf{Intent (EN)}: 

\textbf{Tags (JA)}: \\
\textbf{Tags (EN)}: 

\textbf{Adjusted Expression (JA)}:\\
\\
\textbf{Adjusted Expression (EN)}:\\


\textbf{Example (JA)}:\\
別に自慢しているわけではないんですが、猫と目が合う確率が高いんです。\\
\textbf{Example (EN)}:\\
Not particularly proud, but the probability of making eye contact with cats is quite high.

\textbf{Notes (JA)}:\\
「別に自慢しているわけではないんですが」は、自慢や自己賞賛に聞こえるかもしれない発言を和らげるためのフレーズです。話し手が面白いと思う個人的な特徴や経験を共有する際によく使われます。このフレーズは、発言が少し自己顕示的に見えることを控えめに認めつつ、軽いユーモアを交えてリラックスした雰囲気を作ります。こうすることで、話し手は自分についてのユニークなエピソードを共有しやすくなります。もっとも、心の中ではちょっと自慢したい気持ちがあるのかもしれませんが。

\textbf{Notes (EN)}:\\
"Betsuni jiman shite iru wake de wa nain desu ga" s a phrase used to soften statements that might sound like boasting or self-praise. It’s often used when sharing a personal trait or experience the speaker finds amusing or interesting. The phrase acknowledges the potentially self-centered nature of the remark while adding a touch of humor to create a relaxed atmosphere. By using this phrase, the speaker can comfortably share quirky or unique anecdotes about themselves without appearing overly arrogant. That said, deep down, they might actually feel a little proud.

\newpage

\section*{253. For that reason}

\textbf{Date}: 2024.12.19 (Thu)\\
\textbf{Index}: souiuwakede\\
\textbf{Expression (JA)}: そういう(訳:わけ)で\\
\textbf{Romanization}: Souiu wake de\\
\textbf{Expression (EN)}: For that reason

\textbf{Variants (JA)}: \\
\textbf{Variants (EN)}: 

\textbf{Intent (JA)}: \\
\textbf{Intent (EN)}: 

\textbf{Tags (JA)}: \\
\textbf{Tags (EN)}: 

\textbf{Adjusted Expression (JA)}:\\
\\
\textbf{Adjusted Expression (EN)}:\\


\textbf{Example (JA)}:\\
(猫:ねこ)が(キーボード:きーぼーど)の(上:うえ)を(歩:ある)き(回:まわ)り、(原稿:げんこう)が(消:き)えた。そういう(訳:わけ)で、(今日:きょう)は(宿題:しゅくだい)が(出:だ)せなくてすみません。\\
\textbf{Example (EN)}:\\
The cat walked across my keyboard and deleted my manuscript. So, I'm sorry I couldn't turn in my homework today.

\textbf{Notes (JA)}:\\
「そういう訳で」は、日常のハプニングや言い訳をまとめる際に使われる便利な表現である。この表現は、状況によって、不器用だがどこか憎めない「clumsy excuse」として働くこともあれば、明らかに説得力に欠ける「lame excuse」として受け取られることもある。「猫がキーボードの上を歩き回った」というような、予測不能な出来事に責任を転嫁する場面では、不自然ながらも微笑ましい印象を与え、「まあ仕方ないか」と許されやすい。一方で、「宿題を宇宙人に盗まれた」といった見え透いた理由に使えば、信じがたい言い訳にもあたかも理屈が通っているかのような印象を与える。このように「そういう訳で」は、言い訳の質や状況に応じて柔軟に使え、その語感がユーモアや愛嬌を生み出す要素となっている。

\textbf{Notes (EN)}:\\
"Souiu wake de" is a versatile expression in Japanese that is used to summarize reasons, excuses, or outcomes. Its meaning can shift depending on the situation, functioning as both a clumsy excuse and a lame excuse. In the case of a clumsy excuse, such as "The cat walked across my keyboard," it frames an unexpected event as the cause of a problem. While the explanation may seem a little awkward or silly, it often carries a humorous and endearing quality, inviting the listener to respond with, "Well, I guess it can't be helped." On the other hand, in the case of a lame excuse, like "Aliens stole my homework," it is used to package an obviously unconvincing or far-fetched explanation as if it were logical or believable. This flexibility allows "souiu wake de" to adapt to a wide range of situations, adding humor or charm depending on the context and the nature of the excuse.

\newpage

\section*{252. Soon/in the near future}

\textbf{Date}: 2024.12.18 (Wed)\\
\textbf{Index}: chikaiuchini\\
\textbf{Expression (JA)}: (近:ちか)いうちに\\
\textbf{Romanization}: chikaiuchini\\
\textbf{Expression (EN)}: Soon/in the near future

\textbf{Variants (JA)}: \\
\textbf{Variants (EN)}: 

\textbf{Intent (JA)}: \\
\textbf{Intent (EN)}: 

\textbf{Tags (JA)}: \\
\textbf{Tags (EN)}: 

\textbf{Adjusted Expression (JA)}:\\
\\
\textbf{Adjusted Expression (EN)}:\\


\textbf{Example (JA)}:\\
(近:ちか)いうちに、また(食事:しょくじ)でもしましょう。\\
\textbf{Example (EN)}:\\
Let's have a meal together soon.

\textbf{Notes (JA)}:\\
「近いうちに」は英語で "soon" や "in the near future" と訳され、直近の未来を表す際に使われる表現です。

\textbf{Notes (EN)}:\\
"Chikaiuchini" is an adverbial phrase used to indicate that something will happen soon or in the near future.

\newpage

\section*{251. Probably}

\textbf{Date}: 2024.12.17 (Tue)\\
\textbf{Index}: tabun\\
\textbf{Expression (JA)}: たぶん\\
\textbf{Romanization}: tabun\\
\textbf{Expression (EN)}: Probably

\textbf{Variants (JA)}: \\
\textbf{Variants (EN)}: 

\textbf{Intent (JA)}: \\
\textbf{Intent (EN)}: 

\textbf{Tags (JA)}: \\
\textbf{Tags (EN)}: 

\textbf{Adjusted Expression (JA)}:\\
\\
\textbf{Adjusted Expression (EN)}:\\


\textbf{Example (JA)}:\\
あの(人:ひと)は(天才:てんさい)だ。でも、(本人:ほんにん)はそれに(気:き)づいていない...たぶん。\\
\textbf{Example (EN)}:\\
That person is a genius. But he doesn't realize it...probably.

\textbf{Notes (JA)}:\\
「たぶん」は「probably」や「likely」の意味を持つ副詞で、不確実なことや可能性がある事柄について述べる際に使われます。話し手が完全に確信していない場合や、控えめに意見を伝えたい時に用いられます。たとえば、「たぶん彼は来ない」(He probably won't come.) のように使い、確信はないものの、ある程度の予測や信念を示す表現です。

\textbf{Notes (EN)}:\\
"Tabun" is an adverb that means "probably" or "likely" and is used when referring to uncertain or possible situations. It is often employed when the speaker is not completely sure about something or wants to express a prediction without sounding overly assertive. For example, in the sentence "Tabun kare wa konai" ("He probably won't come"), the speaker conveys a belief or expectation while acknowledging a degree of uncertainty. "Tabun" is also commonly used in Japanese to soften statements, reflecting a cultural preference for indirect or less definitive expressions.

\newpage

\section*{250. To be honest}

\textbf{Date}: 2024.12.16 (Mon)\\
\textbf{Index}: shoujikinatokoro\\
\textbf{Expression (JA)}: (正直:しょうじき)なところ\\
\textbf{Romanization}: shoujikinatokoro\\
\textbf{Expression (EN)}: To be honest

\textbf{Variants (JA)}: \\
\textbf{Variants (EN)}: 

\textbf{Intent (JA)}: \\
\textbf{Intent (EN)}: 

\textbf{Tags (JA)}: \\
\textbf{Tags (EN)}: 

\textbf{Adjusted Expression (JA)}:\\
\\
\textbf{Adjusted Expression (EN)}:\\


\textbf{Example (JA)}:\\
(正直:しょうじき)なところ、この(計画:けいかく)にはあまり(自信:じしん)がありません。\\
\textbf{Example (EN)}:\\
To be honest, I'm not very confident about this plan.

\textbf{Notes (JA)}:\\
「正直なところ」という表現は、文字通り「honestly speaking」や「to be honest」という意味です。この表現は、話し手が本当の気持ちや正直な考えを率直かつ素直に伝える際に使われます。通常、この言葉は、話し手の正直な意見や感情、飾らない視点を反映する発言の前に使われます。また、直接的な表現の影響を和らげ、誠実さや謙虚さを感じさせる効果もあります。

\textbf{Notes (EN)}:\\
The phrase "shoujikinatokoro" literally means "honestly speaking" or "to be honest." It is used to express the speaker's true feelings or genuine thoughts in a candid and straightforward manner. This phrase is typically used to introduce a statement that reflects the speaker's honest opinion, real emotions, or unfiltered perspective. It can soften the impact of a direct statement, making it sound more sincere or humble.

\newpage

\section*{249. at least}

\textbf{Date}: 2024.12.15 (Sun)\\
\textbf{Index}: sukunakutomo\\
\textbf{Expression (JA)}: (少:すく)なくとも\\
\textbf{Romanization}: Sukunakutomo\\
\textbf{Expression (EN)}: at least

\textbf{Variants (JA)}: \\
\textbf{Variants (EN)}: 

\textbf{Intent (JA)}: \\
\textbf{Intent (EN)}: 

\textbf{Tags (JA)}: \\
\textbf{Tags (EN)}: 

\textbf{Adjusted Expression (JA)}:\\
\\
\textbf{Adjusted Expression (EN)}:\\


\textbf{Example (JA)}:\\
(少:すく)なくとも(一:ひと)つ(問題:もんだい)が(解決:かいけつ)した。\\
\textbf{Example (EN)}:\\
At least one problem has been solved.

\textbf{Notes (JA)}:\\
「少なくとも」は、何かが最低限、最低限度であることを表現するために使用されます。

\textbf{Notes (EN)}:\\
"Sukunakutomo" is used to express that something is the minimum or the least. The word "sukunakutomo" is a conjunction that means "at least" or "at the very least."

\newpage

\section*{248. that reminds me}

\textbf{Date}: 2024.12.14 (Sat)\\
\textbf{Index}: souieba\\
\textbf{Expression (JA)}: そういえば\\
\textbf{Romanization}: Souieba\\
\textbf{Expression (EN)}: that reminds me

\textbf{Variants (JA)}: \\
\textbf{Variants (EN)}: 

\textbf{Intent (JA)}: \\
\textbf{Intent (EN)}: 

\textbf{Tags (JA)}: \\
\textbf{Tags (EN)}: 

\textbf{Adjusted Expression (JA)}:\\
\\
\textbf{Adjusted Expression (EN)}:\\


\textbf{Example (JA)}:\\
A「(最近:さいきん)、(寒:さむ)くなってきたね」 B「そういえば、(明日:あす)、(雪:ゆき)、(降:ふ)るって(天気予報:てんきよほう)、(言:い)ってたよ」\\
\textbf{Example (EN)}:\\
A: It's getting cold recently. B: That reminds me, the weather forecast says it will snow tomorrow.

\textbf{Notes (JA)}:\\
「そういえば」は、現在の話題に関連する考えや情報を導入する際に使われます。一方、「ところで」は、現在の会話とは無関係の、完全に新しい話題を導入するために使われます。この違いは、「そういえば」が前の話題と関連性を持たせるのに対し、「ところで」は会話の流れを別の話題に移す点にあります。

\textbf{Notes (EN)}:\\
"Souieba" means "that reminds me" or "speaking of which," and is used to introduce a thought or information related to the current topic. In contrast, "Tokorode" means "by the way" and is used to introduce a completely new topic that is unrelated to the ongoing conversation. The key difference is that "souieba" builds a connection to the previous topic, while "tokorode" shifts the conversation to an entirely different subject.

\newpage

\section*{247. Don't look}

\textbf{Date}: 2024.12.13 (Fri)\\
\textbf{Index}: minaide\\
\textbf{Expression (JA)}: 見ないで\\
\textbf{Romanization}: Minaide\\
\textbf{Expression (EN)}: Don't look

\textbf{Variants (JA)}: \\
\textbf{Variants (EN)}: 

\textbf{Intent (JA)}: \\
\textbf{Intent (EN)}: 

\textbf{Tags (JA)}: \\
\textbf{Tags (EN)}: 

\textbf{Adjusted Expression (JA)}:\\
\\
\textbf{Adjusted Expression (EN)}:\\


\textbf{Example (JA)}:\\
(暗:くら)いところで(テレビ:てれび)を(見:み)ないで。\\
\textbf{Example (EN)}:\\
Don't watch TV in the dark.

\textbf{Notes (JA)}:\\
「見ないで」は、動詞「見る」のやわらかい命令形です。

\textbf{Notes (EN)}:\\
"Minaide" is the soft imperative form of the verb "miru" (to see).

\newpage

\section*{246. Happy}

\textbf{Date}: 2024.12.12 (Thu)\\
\textbf{Index}: ureshii\\
\textbf{Expression (JA)}: うれしい\\
\textbf{Romanization}: Ureshii\\
\textbf{Expression (EN)}: Happy

\textbf{Variants (JA)}: \\
\textbf{Variants (EN)}: 

\textbf{Intent (JA)}: \\
\textbf{Intent (EN)}: 

\textbf{Tags (JA)}: \\
\textbf{Tags (EN)}: 

\textbf{Adjusted Expression (JA)}:\\
\\
\textbf{Adjusted Expression (EN)}:\\


\textbf{Example (JA)}:\\
(今日:きょう)は(旧友:きゅうゆう)に(会:あ)えて、うれしかった。\\
\textbf{Example (EN)}:\\
I was happy to see an old friend today.

\textbf{Notes (JA)}:\\
「うれしい」は、何か良いことが起きたり、意味のあるものを得たりして、そのことに喜びを感じるときに使います。このことばは「たのしい」に似ていますが、使う場面が異なります。うれしいは「出来事や結果による個人的な喜び」、たのしいは「体験や環境から得られる心地よい気持ち」を表します。

\textbf{Notes (EN)}:\\
"Ureshii" is an i-adjective that means "happy" or "glad." It is used to express personal joy when something positive happens or when you receive something meaningful. While similar to "tanoshii" (fun or enjoyable), they are used in different contexts. "Ureshii" describes joy from specific events or outcomes, whereas "tanoshii" refers to pleasant feelings derived from experiences or environments.

\newpage

\section*{245. Fun}

\textbf{Date}: 2024.12.11 (Wed)\\
\textbf{Index}: tanoshii\\
\textbf{Expression (JA)}: (楽:たの)しい\\
\textbf{Romanization}: Tanoshii\\
\textbf{Expression (EN)}: Fun

\textbf{Variants (JA)}: \\
\textbf{Variants (EN)}: 

\textbf{Intent (JA)}: \\
\textbf{Intent (EN)}: 

\textbf{Tags (JA)}: \\
\textbf{Tags (EN)}: 

\textbf{Adjusted Expression (JA)}:\\
\\
\textbf{Adjusted Expression (EN)}:\\


\textbf{Example (JA)}:\\
A「(昨日:きのう)の(パーティー:ぱーてぃ)、(楽:たの)しかったね」 B「うん、(楽:たの)しかった」\\
\textbf{Example (EN)}:\\
A: Yesterday's party was fun. B: Yeah, it was fun.

\textbf{Notes (JA)}:\\
「楽しい」は「楽しいこと」や「愉快なこと」を表すイ形容詞です。イ形容詞は基本的に自分の気持ちを表現することばなので、丁寧に言う必要はありません。そのため、文末に「です」を付けると幼稚に聞こえます。また、「楽しいですか」と人に尋ねることもほとんどありません。もし尋ねる場合は、「本当にこれ、楽しいですか？」というように、その人が本当に楽しんでいるのか疑問を持っている場面で使われます。この場合、「楽しくありませんよね」と確認することで、相手も楽しんでいないのではないかと問いかける意味になります。

\textbf{Notes (EN)}:\\
"Tanoshii" is an i-adjective meaning "fun" or "enjoyable." I-adjectives are primarily used to express one's own feelings, and it is unnecessary to make such expressions polite. For this reason, adding the polite auxiliary verb "desu" makes the expression sound childish. It is also uncommon to ask someone, "Tanoshii desu ka?" If it is asked, it is typically in a context where you are questioning, "Is this really fun?"---perhaps implying doubt about whether the person is genuinely having fun. In such cases, you might also confirm your assumption by saying, "It's not fun, is it?" to suggest that the other person might not be enjoying themselves either.

\newpage

\section*{244. Become}

\textbf{Date}: 2024.12.10 (Tue)\\
\textbf{Index}: mononinaru\\
\textbf{Expression (JA)}: (物:もの)になる\\
\textbf{Romanization}: Mononinaru\\
\textbf{Expression (EN)}: Become

\textbf{Variants (JA)}: \\
\textbf{Variants (EN)}: 

\textbf{Intent (JA)}: \\
\textbf{Intent (EN)}: 

\textbf{Tags (JA)}: \\
\textbf{Tags (EN)}: 

\textbf{Adjusted Expression (JA)}:\\
\\
\textbf{Adjusted Expression (EN)}:\\


\textbf{Example (JA)}:\\
A「(フランス語:ふらんすご)、(物:もの)になった？」 B「いや、ぜんぜん」\\
\textbf{Example (EN)}:\\
A: Did your french become something? B: No, not at all.

\textbf{Notes (JA)}:\\
"物になる"は、自分が鍛錬してきた技術、技能が十分実務的に使えるようになるという意味です。

\textbf{Notes (EN)}:\\
"Mononinaru" means that your skills or abilities have become practical and useful. The word "mononinaru" is a verb that means "become."

\newpage

\section*{243. Method}

\textbf{Date}: 2024.12.09 (Mon)\\
\textbf{Index}: houhou\\
\textbf{Expression (JA)}: (方法:ほうほう)\\
\textbf{Romanization}: Houhou\\
\textbf{Expression (EN)}: Method

\textbf{Variants (JA)}: \\
\textbf{Variants (EN)}: 

\textbf{Intent (JA)}: \\
\textbf{Intent (EN)}: 

\textbf{Tags (JA)}: \\
\textbf{Tags (EN)}: 

\textbf{Adjusted Expression (JA)}:\\
\\
\textbf{Adjusted Expression (EN)}:\\


\textbf{Example (JA)}:\\
この(方法:ほうほう)で(問題:もんだい)が(解決:かいけつ)できる。\\
\textbf{Example (EN)}:\\
This method can solve the problem.

\textbf{Notes (JA)}:\\
"方法"は、何かをするための方法や手段を表現するために使用されます。

\textbf{Notes (EN)}:\\
"Houhou" is used to express a way or method to do something. The word "houhou" is a noun that means "method" or "way." 

\newpage

\section*{242. Perfect}

\textbf{Date}: 2024.12.08 (Sun)\\
\textbf{Index}: kannpeki\\
\textbf{Expression (JA)}: (完璧:かんぺき)\\
\textbf{Romanization}: Kannpeki\\
\textbf{Expression (EN)}: Perfect

\textbf{Variants (JA)}: \\
\textbf{Variants (EN)}: 

\textbf{Intent (JA)}: \\
\textbf{Intent (EN)}: 

\textbf{Tags (JA)}: \\
\textbf{Tags (EN)}: 

\textbf{Adjusted Expression (JA)}:\\
\\
\textbf{Adjusted Expression (EN)}:\\


\textbf{Example (JA)}:\\
A「(今日:きょう)の(プレゼン:ぷれぜん)、(完璧:かんぺき)だったね」 B「ありがとう」\\
\textbf{Example (EN)}:\\
A: Today's presentation was perfect. B: Thank you.

\textbf{Notes (JA)}:\\
"完璧"は、何かが完璧であることを表現するために使用されます。漢字の「完」は「完全」、「璧」は「欠点がない」という意味です。何かが間違いや欠点がないことを表現するときに使われます。「璧」の字をよく観察すると、下に「玉」があります。これは「玉石」という言葉にも使われる「玉」で、玉石には、瑕ひとつない綺麗な表面をしていて、欠点がないという意味があります。「璧」に似ている漢字で「壁」がありますが、よく見ると下は「土」になっていますので、間違えないようにしましょう。完璧にならなくなりますよ。

\textbf{Notes (EN)}:\\
"Kannpeki" is used to express that something is perfect or flawless. The word "kannpeki" is an adjective that means "perfect" or "flawless." It is used to express that something is without any mistakes or flaws. If you look closely at the kanji for "kannpeki," you will see that the character "完" means "complete" and "璧" means "without flaws." When you observe the character "璧" closely, you will see that there is a "玉" at the bottom. This is the "玉" used in the word "gyokuseki" (jewel), which has a beautiful surface without any flaws. Be careful not to confuse the character "璧" with the character "壁" (kabe), which means "wall." The character "壁" has "土" at the bottom, so be careful not to make a mistake. You won't be perfect.

\newpage

\section*{241. Core of your mind}

\textbf{Date}: 2024.12.07 (Sat)\\
\textbf{Index}: honnshinn\\
\textbf{Expression (JA)}: ねぇ、(本心:ほんしん)は？\\
\textbf{Romanization}: Nee, honshinn ha?\\
\textbf{Expression (EN)}: What is your true intention?

\textbf{Variants (JA)}: \\
\textbf{Variants (EN)}: 

\textbf{Intent (JA)}: \\
\textbf{Intent (EN)}: 

\textbf{Tags (JA)}: \\
\textbf{Tags (EN)}: 

\textbf{Adjusted Expression (JA)}:\\
\\
\textbf{Adjusted Expression (EN)}:\\


\textbf{Example (JA)}:\\
A「(君:きみ)が(必要:ひつよう)です」 B「ねぇ、(本心:ほんしん)は？」\\
\textbf{Example (EN)}:\\
A: I need you. B: What is true intention? (You don't really think so, do you?)

\textbf{Notes (JA)}:\\
"本心は？"は、相手が本当に何を考えているのか知りたいときに使われます。「好きだ」とか「必要だ」と言っても、本当はそうではないかもしれません。「本心は？」という質問は、「利用したいだけでしょ」という意図を聞くことになりますね。

\textbf{Notes (EN)}:\\
"Honnshinn ha?" is used to ask what the other person really thinks. The word "honnshinn" is a noun that means "true feelings" or "real intention." Even if someone says "I love you" or "I need you," it may not always be sincere. This question, "Honnshinn ha?" asks for their genuine feelings, often suggesting doubt or skepticism.

\newpage

\section*{240. Oniisan}

\textbf{Date}: 2024.12.06 (Fri)\\
\textbf{Index}: oniisan\\
\textbf{Expression (JA)}: お(兄:にい)さん\\
\textbf{Romanization}: Oniisan\\
\textbf{Expression (EN)}: Big brother

\textbf{Variants (JA)}: \\
\textbf{Variants (EN)}: 

\textbf{Intent (JA)}: \\
\textbf{Intent (EN)}: 

\textbf{Tags (JA)}: \\
\textbf{Tags (EN)}: 

\textbf{Adjusted Expression (JA)}:\\
\\
\textbf{Adjusted Expression (EN)}:\\


\textbf{Example (JA)}:\\
お(兄:にい)さん、(手伝:てつだ)ってくれる？\\
\textbf{Example (EN)}:\\
Big brother, can you help me?

\textbf{Notes (JA)}:\\
"お兄さん"は、日本語で「お兄さん」という意味です。年上の兄弟や、兄弟でなくても年上の男性を指すときに使われます。あまり親しくない間や見ず知らずの人に対して、使うとあなたが下品な人として見做されるので、知らない人には使わない方がよい。もちろん、「お姉さん」も同様です。

\textbf{Notes (EN)}:\\
"Oniisan" means "big brother" in Japanese. It is used to refer to an older brother or a young man who is older than the speaker even if he is not a brother. It is better not to use it for someone you don't know, you are regarded as a vulgar person. Of course, "oneesan" is the same.

\newpage

\section*{239. Akita}

\textbf{Date}: 2024.12.05 (Thu)\\
\textbf{Index}: akita\\
\textbf{Expression (JA)}: 飽きた\\
\textbf{Romanization}: akita\\
\textbf{Expression (EN)}: I'm tired of it

\textbf{Variants (JA)}: \\
\textbf{Variants (EN)}: 

\textbf{Intent (JA)}: \\
\textbf{Intent (EN)}: 

\textbf{Tags (JA)}: \\
\textbf{Tags (EN)}: 

\textbf{Adjusted Expression (JA)}:\\
\\
\textbf{Adjusted Expression (EN)}:\\


\textbf{Example (JA)}:\\
A「(毎日:まいにち)、(同:おな)じこと、ばっかしていると、(飽:あ)きたよね」 B「うん、(新:あたら)しいこと、(始:はじ)めたい」\\
\textbf{Example (EN)}:\\
A: Doing the same thing every day gets boring. B: Yeah, I want to start something new.

\textbf{Notes (JA)}:\\
"飽きた"は、何かに飽きたり、飽き飽きしたりすることを表現するために使用されます。この表現は「食べ飽きる」「見飽きる」のように複合動詞の形で使われることも多い。

\textbf{Notes (EN)}:\\
"Akita" express that you are tired of something or bored with it. This expression is often used in the form of a compound verb such as "tabeakita"(get tired of eating) or "miakita"(get tired of see). 

\newpage

\section*{238. Oh, no, not that much!}

\textbf{Date}: 2024.12.04 (Wed)\\
\textbf{Index}: sorehododemo\\
\textbf{Expression (JA)}: それほどでも...\\
\textbf{Romanization}: Sore hodo demo...\\
\textbf{Expression (EN)}: Not that much...

\textbf{Variants (JA)}: \\
\textbf{Variants (EN)}: 

\textbf{Intent (JA)}: \\
\textbf{Intent (EN)}: 

\textbf{Tags (JA)}: \\
\textbf{Tags (EN)}: 

\textbf{Adjusted Expression (JA)}:\\
\\
\textbf{Adjusted Expression (EN)}:\\


\textbf{Example (JA)}:\\
A「(君:きみ)は(記憶力:きおくりょく)がいいねぇ」 B「いいえ、それほどでも」\\
\textbf{Example (EN)}:\\
A: You have a good memory. B: Not that much.

\textbf{Notes (JA)}:\\
「それほどでも」は、何かが期待されたり、言及されたほどではないことを表現するために使用されます。自分がそのレベルには至っていないことを言うことで、謙遜していることを意味します。しかし、このように先生に褒められた学生の内心は嬉しいに決まっています。「それほどでも」は、褒めてもらった時に謙遜して答える表現です。「上手ですね」と褒めてもらった時に、「ええ、そうですね」と自分で自分の優れていることを認めるのは図々しいので、謙遜のことばを使います。謙遜せずに、一度「ええ、そうですね」と答えてみてください。人気者になれるかもしれませんが、やはり、そう言わないほうがよいと思います。

\textbf{Notes (EN)}:\\
The phrase "sore hodo demo" is used to express that something is not as much as expected or mentioned. It is also used to modestly respond when you are praised. In Japanese, you cannot say "Yes, that's right" when you are praised, so you use a humble expression. However, the inner thoughts of student B who is praised by teacher A are certainly happy. If you think it's a lie, try answering "Yes, that's right" once. You may become popular, but I think it's better not to say that.

\newpage

\section*{237. sou iu imi de,...}

\textbf{Date}: 2024.12.03 (Tue)\\
\textbf{Index}: souiuimide\\
\textbf{Expression (JA)}: そういう(意味:いみ)で、...\\
\textbf{Romanization}: Sou iu imi de,...\\
\textbf{Expression (EN)}: In that sense, ...

\textbf{Variants (JA)}: \\
\textbf{Variants (EN)}: 

\textbf{Intent (JA)}: \\
\textbf{Intent (EN)}: 

\textbf{Tags (JA)}: \\
\textbf{Tags (EN)}: 

\textbf{Adjusted Expression (JA)}:\\
\\
\textbf{Adjusted Expression (EN)}:\\


\textbf{Example (JA)}:\\
わずかですが、とても(貴重:きちょう)な(意見:いけん)が(得:え)られました。そういう(意味:いみ)で、むしろ(会議:かいぎ)は(成功:せいこう)だったと(思:おも)います。\\
\textbf{Example (EN)}:\\
We received only a few comments, but they were very valuable. In that sense, I believe the meeting was actually a success.

\textbf{Notes (JA)}:\\
「そういう意味で」は、必ずしも理想的ではない状況の中で、良い点や肯定的な側面を強調するために使われます。この表現は、価値を見出した発言を、不満足な状況と結びつける役割を果たします。しかし、だいたいは失敗したプロジェクトの中に少しでもよいところを見つけて、強引によい成果があったと言いたいときに使われることも少なくないように思います。あるいは、プロジェクトを失敗して落ち込んでいる人を慰めるために、ほんのわずかの良いところを指摘して、励ますときにも使われますね。

\textbf{Notes (EN)}:\\
"Sou iu imi de" is used to highlight positive aspects of a situation that is less than ideal. It connects a statement of value to an otherwise unsatisfactory context. This phrase is often used to find something good, however small, in a project that may have otherwise been a failure, allowing one to present the results in a more positive light. It is also used to encourage someone who is feeling down after a failed project, by pointing out even the smallest positive outcome to lift their spirits.

\newpage

\section*{236. soko made iu?}

\textbf{Date}: 2024.12.02 (Mon)\\
\textbf{Index}: sokomadeiu\\
\textbf{Expression (JA)}: そこまで(言:い)う？\\
\textbf{Romanization}: Soko made iu?\\
\textbf{Expression (EN)}: Do you really mean that?

\textbf{Variants (JA)}: \\
\textbf{Variants (EN)}: 

\textbf{Intent (JA)}: \\
\textbf{Intent (EN)}: 

\textbf{Tags (JA)}: \\
\textbf{Tags (EN)}: 

\textbf{Adjusted Expression (JA)}:\\
\\
\textbf{Adjusted Expression (EN)}:\\


\textbf{Example (JA)}:\\
A「その(服:ふく)、(趣味:しゅみ)悪いよ」 B「そこまで(言:い)う？」\\
\textbf{Example (EN)}:\\
A: That outfit is in bad taste. B: Do you really mean that?

\textbf{Notes (JA)}:\\
「そこまで言う？」は、何かに驚いたり傷ついたりしたときに使われます。また、簡潔ですが、「驚き」や「軽い抗議」のニュアンスをしっかり伝えられる表現です。"まで"は、範囲や限度を表す助詞です。このように使うと、相手の言葉が自分にとっては予想外だったり、自分にとっては少し厳しい言葉だったりすること、すなわち、「言いすぎ」を表現します。

\textbf{Notes (EN)}:\\
"Soko made iu?" is used to express feelings of surprise or being hurt by something. It is a concise yet effective phrase for conveying mild astonishment or a light protest. The particle made indicates a range or limit. Used in this way, it suggests that the other person's words were unexpected or felt a bit harsh, effectively conveying the sense of "going too far."

\newpage

\section*{235. sorya naiyo}

\textbf{Date}: 2024.12.01 (Sun)\\
\textbf{Index}: soryanaiyo\\
\textbf{Expression (JA)}: そりゃないよ\\
\textbf{Romanization}: Sorya naiyo\\
\textbf{Expression (EN)}: Oh! no, that sucks

\textbf{Variants (JA)}: \\
\textbf{Variants (EN)}: 

\textbf{Intent (JA)}: \\
\textbf{Intent (EN)}: 

\textbf{Tags (JA)}: \\
\textbf{Tags (EN)}: 

\textbf{Adjusted Expression (JA)}:\\
\\
\textbf{Adjusted Expression (EN)}:\\


\textbf{Example (JA)}:\\
A「ごめん、(約束:やくそく)してたけど(今日:きょう)(行:い)けなくなった」 B「えー、そりゃないよ！(楽:たの)しみにしてたのに」\\
\textbf{Example (EN)}:\\
A: Sorry, I can't go today even though I promised. B: Oh! no, that sucks! I was looking forward to it.

\textbf{Notes (JA)}:\\
"そりゃないよ"は、何かにがっかりしたり驚いたりするときに使われます。

\textbf{Notes (EN)}:\\
"Sorya naiyo" is used to express that you are disappointed or surprised by something. The word "sorya" is a contraction of the phrase "sore wa," which means "that is." The word "naiyo" is a contraction of the phrase "nai yo," which means "no."

\newpage

\section*{234. soreha souda}

\textbf{Date}: 2024.11.30 (Sat)\\
\textbf{Index}: sorehasouda\\
\textbf{Expression (JA)}: それはそうだ\\
\textbf{Romanization}: Soreha souda\\
\textbf{Expression (EN)}: That's true

\textbf{Variants (JA)}: \\
\textbf{Variants (EN)}: 

\textbf{Intent (JA)}: \\
\textbf{Intent (EN)}: 

\textbf{Tags (JA)}: \\
\textbf{Tags (EN)}: 

\textbf{Adjusted Expression (JA)}:\\
\\
\textbf{Adjusted Expression (EN)}:\\


\textbf{Example (JA)}:\\
A「やっぱり健康が一番大事だと思うよ」 B「それはそうだ。健康がなければ何も始まらないよね」\\
\textbf{Example (EN)}:\\
A: I think health is the most important thing after all. B: That's true. Without health, nothing can start.

\textbf{Notes (JA)}:\\
"それはそうだ"は、相手の言ったことに全面的に同意するときに使われます。

\textbf{Notes (EN)}:\\
"Soreha souda" is used to express that you totally agree with what the other person said. The word "sore" is a pronoun that means "that." The word "ha" is a particle that is used to indicate the topic of a sentence. The word "souda" is a contraction of the phrase "sou da," which means "that's true."

\newpage

\section*{233. masaka}

\textbf{Date}: 2024.11.29 (Fri)\\
\textbf{Index}: masaka\\
\textbf{Expression (JA)}: まさか\\
\textbf{Romanization}: Masaka\\
\textbf{Expression (EN)}: No way

\textbf{Variants (JA)}: \\
\textbf{Variants (EN)}: 

\textbf{Intent (JA)}: \\
\textbf{Intent (EN)}: 

\textbf{Tags (JA)}: \\
\textbf{Tags (EN)}: 

\textbf{Adjusted Expression (JA)}:\\
\\
\textbf{Adjusted Expression (EN)}:\\


\textbf{Example (JA)}:\\
A「明日、雪が降るらしいよ」 B「まさかぁ」\\
\textbf{Example (EN)}:\\
A: It seems like it's going to snow tomorrow. B: No way.

\textbf{Notes (JA)}:\\
"まさか"は、何かに驚いたり、信じられないときに使われます。

\textbf{Notes (EN)}:\\
"Masaka" is used to express that you are surprised or don't believe something. The word "masaka" is an adverb that means "no way" or "impossible."

\newpage

\section*{232. nannde}

\textbf{Date}: 2024.11.28 (Thu)\\
\textbf{Index}: nannde\\
\textbf{Expression (JA)}: (何:なん)で\\
\textbf{Romanization}: Nannde\\
\textbf{Expression (EN)}: How come?

\textbf{Variants (JA)}: \\
\textbf{Variants (EN)}: 

\textbf{Intent (JA)}: \\
\textbf{Intent (EN)}: 

\textbf{Tags (JA)}: \\
\textbf{Tags (EN)}: 

\textbf{Adjusted Expression (JA)}:\\
\\
\textbf{Adjusted Expression (EN)}:\\


\textbf{Example (JA)}:\\
(今日:きょう)、(電車:でんしゃ)、(何:なん)で(遅:おく)れるの？\\
\textbf{Example (EN)}:\\
How come the train is late today?

\textbf{Notes (JA)}:\\
"何で"は、何かに驚いたり、疑問に思ったりするときに使われます。ときどき、良くないことを非難するときにも使われます。

\textbf{Notes (EN)}:\\
"Nannde" is used to express that you are surprised or curious about something. The word "nannde" is an interrogative pronoun that means "how come?" or "why?" Sometimes it is used to criticize something bad.

\newpage

\section*{231. tatta ima}

\textbf{Date}: 2024.11.27 (Wed)\\
\textbf{Index}: tattaima\\
\textbf{Expression (JA)}: たった(今:いま)\\
\textbf{Romanization}: Tatta ima\\
\textbf{Expression (EN)}: Just now

\textbf{Variants (JA)}: \\
\textbf{Variants (EN)}: 

\textbf{Intent (JA)}: \\
\textbf{Intent (EN)}: 

\textbf{Tags (JA)}: \\
\textbf{Tags (EN)}: 

\textbf{Adjusted Expression (JA)}:\\
\\
\textbf{Adjusted Expression (EN)}:\\


\textbf{Example (JA)}:\\
たった(今:いま)、(メール:めーる)を(送:おく)ったので(確認:かくにん)してください。\\
\textbf{Example (EN)}:\\
I just sent an email, so please check it.

\textbf{Notes (JA)}:\\
"たった今"は、今、ちょうどその瞬間に何かが起こったことを表現するために使用されます。

\textbf{Notes (EN)}:\\
"Tattaima" is used to express that something happened just now. The word "tatta" is an adverb that means "just" or "only." The word "ima" is a noun that means "now."

\newpage

\section*{230. -ta hou ga ii}

\textbf{Date}: 2024.11.26 (Tue)\\
\textbf{Index}: tahougaii\\
\textbf{Expression (JA)}: -た(方:ほう)がいい\\
\textbf{Romanization}: -ta hou ga ii\\
\textbf{Expression (EN)}: It's better to

\textbf{Variants (JA)}: \\
\textbf{Variants (EN)}: 

\textbf{Intent (JA)}: \\
\textbf{Intent (EN)}: 

\textbf{Tags (JA)}: \\
\textbf{Tags (EN)}: 

\textbf{Adjusted Expression (JA)}:\\
\\
\textbf{Adjusted Expression (EN)}:\\


\textbf{Example (JA)}:\\
(試験:しけん)の(前日:ぜんじつ)は、(勉強:べんきょう)しても(点:てん)が(良:よ)くなるわけではないので、(早:はや)く(寝:ね)たほうがいい。\\
\textbf{Example (EN)}:\\
You should go to bed early because studying the day before the exam doesn't necessarily improve your score.

\textbf{Notes (JA)}:\\
"-た方がいい"は、何かをする方が良いということを表現するために使用されます。

\textbf{Notes (EN)}:\\
"-ta hou ga ii" is used to express that it is better to do something. The word "hou" is a noun that means "way" or "direction." The word "ii" is an adjective that means "good."

\newpage

\section*{229. Kanari}

\textbf{Date}: 2024.11.25 (Mon)\\
\textbf{Index}: kanari\\
\textbf{Expression (JA)}: かなり\\
\textbf{Romanization}: Kanari\\
\textbf{Expression (EN)}: Quite

\textbf{Variants (JA)}: \\
\textbf{Variants (EN)}: 

\textbf{Intent (JA)}: \\
\textbf{Intent (EN)}: 

\textbf{Tags (JA)}: \\
\textbf{Tags (EN)}: 

\textbf{Adjusted Expression (JA)}:\\
\\
\textbf{Adjusted Expression (EN)}:\\


\textbf{Example (JA)}:\\
この(実験:じっけん)を(成功:せいこう)させるのは、かなり(難:むずか)しい。\\
\textbf{Example (EN)}:\\
It's quite difficult to make this experiment successful.

\textbf{Notes (JA)}:\\
"かなり"は、何かの程度が大きいことを意味します。

\textbf{Notes (EN)}:\\
"Kanari" is used to express that something is quite or fairly. The word "kanari" is an adverb that means "quite" or "fairly." It is used to express that something is more than expected or more than usual.

\newpage

\section*{228. Nasakenai}

\textbf{Date}: 2024.11.24 (Sun)\\
\textbf{Index}: nasakenai\\
\textbf{Expression (JA)}: (情:なさ)けない\\
\textbf{Romanization}: Nasakenai\\
\textbf{Expression (EN)}: Pitiful

\textbf{Variants (JA)}: \\
\textbf{Variants (EN)}: 

\textbf{Intent (JA)}: \\
\textbf{Intent (EN)}: 

\textbf{Tags (JA)}: \\
\textbf{Tags (EN)}: 

\textbf{Adjusted Expression (JA)}:\\
\\
\textbf{Adjusted Expression (EN)}:\\


\textbf{Example (JA)}:\\
こんなに(何:なに)もできないなんて、(情:なさ)けない。\\
\textbf{Example (EN)}:\\
It's pitiful that I can't do anything.

\textbf{Notes (JA)}:\\
"情けない"は、誰かや何かに同情することを表現するために使用されます。また、自分自身の腑甲斐無いことを表わすのにも使います。

\textbf{Notes (EN)}:\\
"Nasakenai" is used to express that you feel sorry for someone or something. The word "nasakenai" is an adjective that means "pitiful" or "pathetic." It is used to express sympathy or pity for someone or something. It is also used to express self-pity.

\newpage

\section*{227. Yokan}

\textbf{Date}: 2024.11.23 (Sat)\\
\textbf{Index}: yokan\\
\textbf{Expression (JA)}: (予感:よかん)\\
\textbf{Romanization}: Yokan\\
\textbf{Expression (EN)}: Premonition

\textbf{Variants (JA)}: \\
\textbf{Variants (EN)}: 

\textbf{Intent (JA)}: \\
\textbf{Intent (EN)}: 

\textbf{Tags (JA)}: \\
\textbf{Tags (EN)}: 

\textbf{Adjusted Expression (JA)}:\\
\\
\textbf{Adjusted Expression (EN)}:\\


\textbf{Example (JA)}:\\
(今日:きょう)は(何:なに)か(良:い)いことがありそうな(予感:よかん)がする。\\
\textbf{Example (EN)}:\\
I have a premonition that something good will happen today.

\textbf{Notes (JA)}:\\
"予感"は、将来起こる何かについての感情や直感を表現するために使用されます。"いい予感"、"悪い予感"のように使用されます。

\textbf{Notes (EN)}:\\
"Yokan" is used to express that you have a feeling or intuition about something that will happen in the future. The word "yokan" is a noun that means "premonition" or "hunch." "ii yokan" is used to express a good premonition, and "warui yokan" is used to express a bad premonition.

\newpage

\section*{226. yosasou}

\textbf{Date}: 2024.11.22 (Fri)\\
\textbf{Index}: yosasou\\
\textbf{Expression (JA)}: (良:よ)さそう\\
\textbf{Romanization}: Yosasou\\
\textbf{Expression (EN)}: Looks good

\textbf{Variants (JA)}: \\
\textbf{Variants (EN)}: 

\textbf{Intent (JA)}: \\
\textbf{Intent (EN)}: 

\textbf{Tags (JA)}: \\
\textbf{Tags (EN)}: 

\textbf{Adjusted Expression (JA)}:\\
\\
\textbf{Adjusted Expression (EN)}:\\


\textbf{Example (JA)}:\\
この(サイズ:さいず)、ちょうど(良:よ)さそう！\\
\textbf{Example (EN)}:\\
This size looks just right!

\textbf{Notes (JA)}:\\
"よさそう"は、何かが良さそうであることを表現するために使用されます。「いい」は「いさそう」では言いにくいので、「よい」を使って「よさそう」と言います。

\textbf{Notes (EN)}:\\
"Yosasou" is used to express that something looks good or seems to be good. The word "yosasou" is an adjective that means "looks good" or "seems good." "Yoi" is an alternative to "ii" and is used to avoid saying "iisou" which is difficult to say.

\newpage

\section*{225. toriaezu}

\textbf{Date}: 2024.11.21 (Thu)\\
\textbf{Index}: toriaezu\\
\textbf{Expression (JA)}: (取:と)りあえず\\
\textbf{Romanization}: Toriaezu\\
\textbf{Expression (EN)}: For now

\textbf{Variants (JA)}: \\
\textbf{Variants (EN)}: 

\textbf{Intent (JA)}: \\
\textbf{Intent (EN)}: 

\textbf{Tags (JA)}: \\
\textbf{Tags (EN)}: 

\textbf{Adjusted Expression (JA)}:\\
\\
\textbf{Adjusted Expression (EN)}:\\


\textbf{Example (JA)}:\\
とりあえず、(今日:きょう)はここまで。\\
\textbf{Example (EN)}:\\
For now, let's stop here today.

\textbf{Notes (JA)}:\\
"とりあえず"は、何かを一時的にしたり、当座の間だけしたりすることを表現するために使用されます。

\textbf{Notes (EN)}:\\
"Toriaezu" is used to express that you are doing something temporarily or for the time being. The word "toriaezu" is an adverb that means "for now" or "for the time being."

\newpage

\section*{224. That's true!}

\textbf{Date}: 2024.11.20 (Wed)\\
\textbf{Index}: tashikani\\
\textbf{Expression (JA)}: (確:たし)かに\\
\textbf{Romanization}: Tashika ni\\
\textbf{Expression (EN)}: That's true!

\textbf{Variants (JA)}: \\
\textbf{Variants (EN)}: 

\textbf{Intent (JA)}: \\
\textbf{Intent (EN)}: 

\textbf{Tags (JA)}: \\
\textbf{Tags (EN)}: 

\textbf{Adjusted Expression (JA)}:\\
\\
\textbf{Adjusted Expression (EN)}:\\


\textbf{Example (JA)}:\\
A「(彼:かれ)、(最近:さいきん)、(変:へん)だよね」 B「確かに」\\
\textbf{Example (EN)}:\\
A: He's been acting strange lately. B: You're right.

\textbf{Notes (JA)}:\\
"確かに" は、「確実に」という意味もありますが、日常会話では「言われてみればそうだね」や「今気がついた」というニュアンスで使われることが多いです。

\textbf{Notes (EN)}:\\
"Tashikani" means "certainly" or "surely." However, it is more commonly used to express agreement or realization, like "You're right" or "I see."

\newpage

\section*{223. nanteiuka,..}

\textbf{Date}: 2024.11.19 (Tue)\\
\textbf{Index}: nanteiuka\\
\textbf{Expression (JA)}: なんていうか\\
\textbf{Romanization}: Nanteiuka\\
\textbf{Expression (EN)}: How should I put it

\textbf{Variants (JA)}: \\
\textbf{Variants (EN)}: 

\textbf{Intent (JA)}: \\
\textbf{Intent (EN)}: 

\textbf{Tags (JA)}: \\
\textbf{Tags (EN)}: 

\textbf{Adjusted Expression (JA)}:\\
\\
\textbf{Adjusted Expression (EN)}:\\


\textbf{Example (JA)}:\\
A「(彼:かれ)、(最近:さいきん)、(変:へん)だよね」 B「なんていうか、(忙:いそが)しいのかな」\\
\textbf{Example (EN)}:\\
A: He's been acting strange lately. B: How should I put it, maybe he's busy.

\textbf{Notes (JA)}:\\
"なんていうか"は、何かを表現する正しい言葉を見つけようとしていることを表現するために使用されます。しかしながら、いろいろな文脈で使えるところから、ほとんど意味がないと考えてよいでしょう。

\textbf{Notes (EN)}:\\
"Nanteiuka" is used to express that you are trying to find the right words to describe something. The word "nante" is an interrogative pronoun that means "what kind of." The word "iu" is a verb that means "to say." The word "ka" is a particle that is used to indicate a question. However, it is used in various contexts, so it can be considered almost meaningless.

\newpage

\section*{222. Shoganaine}

\textbf{Date}: 2024.11.18 (Mon)\\
\textbf{Index}: shoganaine\\
\textbf{Expression (JA)}: しょうがないね\\
\textbf{Romanization}: Shoganaine\\
\textbf{Expression (EN)}: It can't be helped

\textbf{Variants (JA)}: \\
\textbf{Variants (EN)}: 

\textbf{Intent (JA)}: \\
\textbf{Intent (EN)}: 

\textbf{Tags (JA)}: \\
\textbf{Tags (EN)}: 

\textbf{Adjusted Expression (JA)}:\\
\\
\textbf{Adjusted Expression (EN)}:\\


\textbf{Example (JA)}:\\
(急:きゅう)に(予定:よてい)が(入:はい)った？しょうがないね、また(今度:こんど)にしよう。\\
\textbf{Example (EN)}:\\
Did you have a sudden appointment? It can't be helped, let's do it next time.

\textbf{Notes (JA)}:\\
"しょうがないね"は、状況に対して何もできないことを表現するために使用されます。状況を受け入れるときに使える表現です。

\textbf{Notes (EN)}:\\
"Shou ga nai ne" expresses acceptance of a situation that can't be changed. The word "shou" comes from "shiyou," meaning "way." This phrase is used not only when something is unavoidable, but also to convey acceptance.

\newpage

\section*{221. Not that reason}

\textbf{Date}: 2024.11.17 (Sun)\\
\textbf{Index}: sonnawakedemonakute\\
\textbf{Expression (JA)}: そんな(訳:わけ)でもなくて\\
\textbf{Romanization}: Sonna wake demo nakute\\
\textbf{Expression (EN)}: Not that reason

\textbf{Variants (JA)}: \\
\textbf{Variants (EN)}: 

\textbf{Intent (JA)}: \\
\textbf{Intent (EN)}: 

\textbf{Tags (JA)}: \\
\textbf{Tags (EN)}: 

\textbf{Adjusted Expression (JA)}:\\
\\
\textbf{Adjusted Expression (EN)}:\\


\textbf{Example (JA)}:\\
A「(昨日:きのう)、(遅:おそ)くまで(起:お)きていたの？」 B「そんな(訳:わけ)でもなくて、(映画:えいが)を(見:み)ていたら(時間:じかん)が(経:た)ってしまってた」\\
\textbf{Example (EN)}:\\
A: Were you up late last night? B: Not that reason, I was watching a movie and time flew by.

\textbf{Notes (JA)}:\\
"そんなわけでもなくて"は、何かの理由が相手が思っているものと違うことを表現するために使用されます。しかし、その理由が外れているわけでもなくて、確かにそういう理由なのだけれども、それは意図してそうしたのではないということも意味します。

\textbf{Notes (EN)}:\\
"Sonna wake demo nakute" is used to express that the reason for something is not what the other person thinks. The word "sonna" is a demonstrative pronoun that means "that kind of." The word "wake" is a noun that means "reason." The word "demo" is a conjunction that means "but" or "however." The word "nakute" is the 'te-form' of the verb "nai," which means "to not be." In some cases, the reason given by the other person isn't exactly wrong, but it also suggests that it wasn't intended to be that way.

\newpage

\section*{220. Ma, ikka}

\textbf{Date}: 2024.11.16 (Sat)\\
\textbf{Index}: maikka\\
\textbf{Expression (JA)}: ま、いっか\\
\textbf{Romanization}: Ma, ikka\\
\textbf{Expression (EN)}: Oh well

\textbf{Variants (JA)}: \\
\textbf{Variants (EN)}: 

\textbf{Intent (JA)}: \\
\textbf{Intent (EN)}: 

\textbf{Tags (JA)}: \\
\textbf{Tags (EN)}: 

\textbf{Adjusted Expression (JA)}:\\
\\
\textbf{Adjusted Expression (EN)}:\\


\textbf{Example (JA)}:\\
A「(今日:きょう)のお(弁当:べんとう)は、すべて(売:う)り(切:き)れました」 B「ま、いっか」\\
\textbf{Example (EN)}:\\
A: Today's bento boxes are all sold out. B: Oh well, never mind.

\textbf{Notes (JA)}:\\


\textbf{Notes (EN)}:\\
"Ma, ikka" is used to express that you don't care about something or that you are not bothered by something. The word "ma" is an interjection that is used to express hesitation or uncertainty. The word "ikka" is a contraction of the phrase "ii kara," which means "it's okay."

\newpage

\section*{219. Maane}

\textbf{Date}: 2024.11.15 (Fri)\\
\textbf{Index}: maane\\
\textbf{Expression (JA)}: まあね\\
\textbf{Romanization}: Maane\\
\textbf{Expression (EN)}: Well, you know

\textbf{Variants (JA)}: \\
\textbf{Variants (EN)}: 

\textbf{Intent (JA)}: \\
\textbf{Intent (EN)}: 

\textbf{Tags (JA)}: \\
\textbf{Tags (EN)}: 

\textbf{Adjusted Expression (JA)}:\\
\\
\textbf{Adjusted Expression (EN)}:\\


\textbf{Example (JA)}:\\
A「こんなことしたくないよね」 B「まあね」\\
\textbf{Example (EN)}:\\
A: I don't want to do this. B: Well, you know.

\textbf{Notes (JA)}:\\
「まあね」は、消極的な同意や理解を示す表現です。同意しているように見えますが、実際には同意も否定もしたくないときに便利に使われます。そのため、どちらにも踏み込まず曖昧な返事をしたい場面でよく使われます。この表現を聞いたときには、相手からの積極的な賛同をあまり期待しない方がよいでしょう。

\textbf{Notes (EN)}:\\
"Maane" is an expression of passive agreement or understanding. While it conveys a sense of agreement, it is often used when one does not want to fully commit to agreeing or disagreeing. This makes it convenient in situations where one wishes to remain non-committal. When you hear this expression, it is best not to expect too much enthusiasm or certainty from the speaker.

\newpage

\section*{218. Tamani}

\textbf{Date}: 2024.11.14 (Thu)\\
\textbf{Index}: tamani\\
\textbf{Expression (JA)}: たまに\\
\textbf{Romanization}: Tamani\\
\textbf{Expression (EN)}: Sometimes

\textbf{Variants (JA)}: \\
\textbf{Variants (EN)}: 

\textbf{Intent (JA)}: \\
\textbf{Intent (EN)}: 

\textbf{Tags (JA)}: \\
\textbf{Tags (EN)}: 

\textbf{Adjusted Expression (JA)}:\\
\\
\textbf{Adjusted Expression (EN)}:\\


\textbf{Example (JA)}:\\
(インスタントラーメン:いんすたんとらーめん)はなるたけ(食:た)べないようにしているが、たまに(食:た)べたくなることがある。\\
\textbf{Example (EN)}:\\
I try not to eat instant noodles as much as possible, but sometimes I feel like eating them.

\textbf{Notes (JA)}:\\


\textbf{Notes (EN)}:\\
"Tamani" means sometimes or occasionally. "narutake" means as much as possible. "tabetakunaru" means to feel like eating. "koto ga aru" is used to express that something happens occasionally.

\newpage

\section*{217. Having said that..}

\textbf{Date}: 2024.11.13 (Wed)\\
\textbf{Index}: tohaittemo\\
\textbf{Expression (JA)}: とは(言:い)っても\\
\textbf{Romanization}: To ha itte mo\\
\textbf{Expression (EN)}: However

\textbf{Variants (JA)}: \\
\textbf{Variants (EN)}: 

\textbf{Intent (JA)}: \\
\textbf{Intent (EN)}: 

\textbf{Tags (JA)}: \\
\textbf{Tags (EN)}: 

\textbf{Adjusted Expression (JA)}:\\
\\
\textbf{Adjusted Expression (EN)}:\\


\textbf{Example (JA)}:\\
(彼:かれ)は(忙:いそが)しい。とは(言:い)っても、(手伝:てつだ)ってくれるよ。\\
\textbf{Example (EN)}:\\
He's busy. However, he will help you.

\textbf{Notes (JA)}:\\
"とは言っても"は、対照や矛盾を表現するために使用されます。

\textbf{Notes (EN)}:\\
"To ha ittemo" is used to express a contrast or contradiction. The word "to" is a particle that is used to indicate a quotation. The word "ha" is a particle that is used to indicate the topic of a sentence. The word "itte" is the te-form of the verb "iu," which means "to say." The word "mo" is a particle that is used to indicate that the speaker is adding information.

\newpage

\section*{216. taimu sa-bisu}

\textbf{Date}: 2024.11.12 (Tue)\\
\textbf{Index}: taimusa-bisu\\
\textbf{Expression (JA)}: (タイムサービス:たいむさーびす)\\
\textbf{Romanization}: Taimu saabisu\\
\textbf{Expression (EN)}: Time-limited offer

\textbf{Variants (JA)}: \\
\textbf{Variants (EN)}: 

\textbf{Intent (JA)}: \\
\textbf{Intent (EN)}: 

\textbf{Tags (JA)}: \\
\textbf{Tags (EN)}: 

\textbf{Adjusted Expression (JA)}:\\
\\
\textbf{Adjusted Expression (EN)}:\\


\textbf{Example (JA)}:\\
この(商品:しょうひん)は、(今:いま)、(タイム:たいむ)(サービス:さーびす)で(半額:はんがく)です。\\
\textbf{Example (EN)}:\\
This product is half price with a time-limited offer.

\textbf{Notes (JA)}:\\
"タイムサービス"は、商品やサービスが期間限定で割引価格で提供されていることを表現するために使用されます。"タイムサービス"は、"タイムセール"とも言います。

\textbf{Notes (EN)}:\\
"Taimu saabisu" is used to express that a product or service is available at a discounted price for a limited time. The word "taimu" is a noun that means "time." The word "saabisu" is a noun that means "service." "Taimu seeru" is also used.

\newpage

\section*{215. Kamoku}

\textbf{Date}: 2024.11.11 (Mon)\\
\textbf{Index}: kamoku\\
\textbf{Expression (JA)}: (火木:かもく)\\
\textbf{Romanization}: Kamoku\\
\textbf{Expression (EN)}: Tuesday and Thursday

\textbf{Variants (JA)}: \\
\textbf{Variants (EN)}: 

\textbf{Intent (JA)}: \\
\textbf{Intent (EN)}: 

\textbf{Tags (JA)}: \\
\textbf{Tags (EN)}: 

\textbf{Adjusted Expression (JA)}:\\
\\
\textbf{Adjusted Expression (EN)}:\\


\textbf{Example (JA)}:\\
(火木:かもく)で(柔道:じゅうどう)の(稽古:けいこ)に行っている。\\
\textbf{Example (EN)}:\\
I go to judo practice on Tuesdays and Thursdays.

\textbf{Notes (JA)}:\\
"火木"は、「火曜日」と「木曜日」を指します。同様に、「土日」は、「土曜日」と「日曜日」を指します。

\textbf{Notes (EN)}:\\
"Kamoku" means "kayoubi" (Tuesday) and "mokuyoubi" (Thursday). In the same way, "donichi" means "doyoubi" (Saturday) and "nichiyoubi" (Sunday).

\newpage

\section*{214. Shuuichi}

\textbf{Date}: 2024.11.10 (Sun)\\
\textbf{Index}: shuuichi\\
\textbf{Expression (JA)}: (週1:しゅういち)\\
\textbf{Romanization}: Shuuichi\\
\textbf{Expression (EN)}: Once a week

\textbf{Variants (JA)}: \\
\textbf{Variants (EN)}: 

\textbf{Intent (JA)}: \\
\textbf{Intent (EN)}: 

\textbf{Tags (JA)}: \\
\textbf{Tags (EN)}: 

\textbf{Adjusted Expression (JA)}:\\
\\
\textbf{Adjusted Expression (EN)}:\\


\textbf{Example (JA)}:\\
(週1:しゅういち)で、(ジム:じむ)に(行:い)っている。\\
\textbf{Example (EN)}:\\
I go to the gym once a week.

\textbf{Notes (JA)}:\\
"週1"は、何かが週に1回行われることを表現するために使用されます。同様に、「週2」は「週に2回」、「週3」は「週に3回」という意味です。

\textbf{Notes (EN)}:\\
"Shuuichi" is used to express that something happens once a week. The word "shuu" is a noun that means "week." The word "ichi" is a number that means "one." The word "de" is a particle that is used to indicate the frequency of an action. To the same token, "shuuni" means "twice a week," "shuumi" means "three times a week," and so on.

\newpage

\section*{213. sore ha nai}

\textbf{Date}: 2024.11.09 (Sat)\\
\textbf{Index}: sorehanai\\
\textbf{Expression (JA)}: それはない\\
\textbf{Romanization}: Sore ha nai\\
\textbf{Expression (EN)}: That's not true/That's impossible

\textbf{Variants (JA)}: \\
\textbf{Variants (EN)}: 

\textbf{Intent (JA)}: \\
\textbf{Intent (EN)}: 

\textbf{Tags (JA)}: \\
\textbf{Tags (EN)}: 

\textbf{Adjusted Expression (JA)}:\\
\\
\textbf{Adjusted Expression (EN)}:\\


\textbf{Example (JA)}:\\
A「(彼:かれ)が(嘘:うそ)を(つ:つ)いていると(思:おも)う？」 B「それはない」\\
\textbf{Example (EN)}:\\
A: Do you think he's lying? B: That's not true.

\textbf{Notes (JA)}:\\
"それはない"は、何かが真実でないことや不可能であることを表現するために使用されます。

\textbf{Notes (EN)}:\\
"Sore ha nai" is used to express that something is not true or impossible. The word "sore" is a pronoun that means "that." The word "ha" is a particle that is used to indicate the topic of a sentence. The word "nai" is a verb that means "to not be."

\newpage

\section*{212. imasara}

\textbf{Date}: 2024.11.08 (Fri)\\
\textbf{Index}: imasara\\
\textbf{Expression (JA)}: (今更:いまさら)\\
\textbf{Romanization}: Imasara\\
\textbf{Expression (EN)}: At this point

\textbf{Variants (JA)}: \\
\textbf{Variants (EN)}: 

\textbf{Intent (JA)}: \\
\textbf{Intent (EN)}: 

\textbf{Tags (JA)}: \\
\textbf{Tags (EN)}: 

\textbf{Adjusted Expression (JA)}:\\
\\
\textbf{Adjusted Expression (EN)}:\\


\textbf{Example (JA)}:\\
(今更:いまさら)、(謝:あやま)っても(遅:おそ)いよ。\\
\textbf{Example (EN)}:\\
It's too late to apologize now.

\textbf{Notes (JA)}:\\


\textbf{Notes (EN)}:\\
"Imasara" is used to express that it is too late to do something. The word "imasara" is an adverb that means "now" or "at this point."

\newpage

\section*{211. Invincible}

\textbf{Date}: 2024.11.07 (Thu)\\
\textbf{Index}: saikyou\\
\textbf{Expression (JA)}: (最強:さいきょう)\\
\textbf{Romanization}: Saikyou\\
\textbf{Expression (EN)}: Invincible

\textbf{Variants (JA)}: \\
\textbf{Variants (EN)}: 

\textbf{Intent (JA)}: \\
\textbf{Intent (EN)}: 

\textbf{Tags (JA)}: \\
\textbf{Tags (EN)}: 

\textbf{Adjusted Expression (JA)}:\\
\\
\textbf{Adjusted Expression (EN)}:\\


\textbf{Example (JA)}:\\
(彼:かれ)は(最強:さいきょう)の(チーム:ちーむ)に(所属:しょぞく)している。\\
\textbf{Example (EN)}:\\
He belongs to the strongest team.

\textbf{Notes (JA)}:\\
"最強"は、何かが最も強力であることを表現するために使用されます。

\textbf{Notes (EN)}:\\
"Saikyou" is used to express that something is the strongest or most powerful. The word "saikyou" is a noun that means "the strongest."

\newpage

\section*{210. How can I do that?}

\textbf{Date}: 2024.11.06 (Wed)\\
\textbf{Index}: doushimashou\\
\textbf{Expression (JA)}: どうしよう/どうしましょう？\\
\textbf{Romanization}: Dou shimashou?\\
\textbf{Expression (EN)}: How should I do that?

\textbf{Variants (JA)}: \\
\textbf{Variants (EN)}: 

\textbf{Intent (JA)}: \\
\textbf{Intent (EN)}: 

\textbf{Tags (JA)}: \\
\textbf{Tags (EN)}: 

\textbf{Adjusted Expression (JA)}:\\
\\
\textbf{Adjusted Expression (EN)}:\\


\textbf{Example (JA)}:\\
あ、(困:こま)った。どうしよう？\\
\textbf{Example (EN)}:\\
Oh no, what should I do?

\textbf{Notes (JA)}:\\
"どうしよう/しましょう"は、困難な状況にあり、どうすればいいかわからないときに使われる独り言です。

\textbf{Notes (EN)}:\\
"Dou shiyou/shimashou?" is a soliloquy used when you are in a difficult situation and don't know what to do. The word "dou" is an interrogative pronoun that means "how." The word "shimashou" is a polite expression of the verb "suru," which means "to do." This is used when you are talking to yourself.

\newpage

\section*{209. Nani yattenda?}

\textbf{Date}: 2024.11.05 (Tue)\\
\textbf{Index}: naniyattemda\\
\textbf{Expression (JA)}: (何:なに)やってんだ？\\
\textbf{Romanization}: Nani yattenda?\\
\textbf{Expression (EN)}: What am I doing?

\textbf{Variants (JA)}: \\
\textbf{Variants (EN)}: 

\textbf{Intent (JA)}: \\
\textbf{Intent (EN)}: 

\textbf{Tags (JA)}: \\
\textbf{Tags (EN)}: 

\textbf{Adjusted Expression (JA)}:\\
\\
\textbf{Adjusted Expression (EN)}:\\


\textbf{Example (JA)}:\\
(何:なに)やってんだ？\\
\textbf{Example (EN)}:\\
What am I doing?

\textbf{Notes (JA)}:\\
"何やってんだ？"は、自分の行動にが混乱しているという表現です。いわゆる、独り言ですが、よく使われます。

\textbf{Notes (EN)}:\\
"Nani yattenda?" is used to express that you are surprised or confused about your own actions. The word "nani" is a pronoun that means "what." The word "yattenda" is a contraction of the phrase "yatte iru n da," which means "doing." It is so called "talking to oneself" expression,jbut it is often used.

\newpage

\section*{208. wasurete shimatta}

\textbf{Date}: 2024.11.04 (Mon)\\
\textbf{Index}: wasureteshimatta\\
\textbf{Expression (JA)}: (忘:わす)れてしまった\\
\textbf{Romanization}: Wasurete shimatta\\
\textbf{Expression (EN)}: Forgot

\textbf{Variants (JA)}: \\
\textbf{Variants (EN)}: 

\textbf{Intent (JA)}: \\
\textbf{Intent (EN)}: 

\textbf{Tags (JA)}: \\
\textbf{Tags (EN)}: 

\textbf{Adjusted Expression (JA)}:\\
\\
\textbf{Adjusted Expression (EN)}:\\


\textbf{Example (JA)}:\\
(昨日:きのう)、(妻:つま)の(誕生日:たんじょうび)を(忘:わす)れてしまった。\\
\textbf{Example (EN)}:\\
I forgot my wife's birthday yesterday.

\textbf{Notes (JA)}:\\
"忘れてしまった"は、何かを忘れたことを表現するために使用されます。だいたいは、失敗したり後悔したりしたときに使います。

\textbf{Notes (EN)}:\\
"Wasurete shimatta" is used to express that you forgot something. The word "wasurete" is the te-form of the verb "wasureru," which means "to forget." The word "shimatta" is a contraction of the phrase "shimaimashita," which means "to have done something." When you use "te shimatta," it implies that you did something unintentionally or regretfully.

\newpage

\section*{207. arubaito}

\textbf{Date}: 2024.11.03 (Sun)\\
\textbf{Index}: arubaito\\
\textbf{Expression (JA)}: (アルバイト:あるばいと)\\
\textbf{Romanization}: Arubaito\\
\textbf{Expression (EN)}: Part-time job

\textbf{Variants (JA)}: \\
\textbf{Variants (EN)}: 

\textbf{Intent (JA)}: \\
\textbf{Intent (EN)}: 

\textbf{Tags (JA)}: \\
\textbf{Tags (EN)}: 

\textbf{Adjusted Expression (JA)}:\\
\\
\textbf{Adjusted Expression (EN)}:\\


\textbf{Example (JA)}:\\
(彼:かれ)は(週:しゅう)に3(回:かい)、(カフェ:かふぇ)で(アルバイト:あるばいと)をしています。\\
\textbf{Example (EN)}:\\
He works part-time at a cafe three times a week.

\textbf{Notes (JA)}:\\
"アルバイト"は、ドイツ語の"arbeit"という言葉からの借用語で、「仕事」という意味です。日本語では、part-time jobを指すために使用されます。

\textbf{Notes (EN)}:\\
"Arubaito" is a loanword from the German word "arbeit," which means "work." The word "arubaito" is a noun that means "part-time job."

\newpage

\section*{206. tetori ashitori}

\textbf{Date}: 2024.11.02 (Sat)\\
\textbf{Index}: tetoriashitori\\
\textbf{Expression (JA)}: (手取:てと)り(足取:あしと)り\\
\textbf{Romanization}: Tetori ashitori\\
\textbf{Expression (EN)}: Hand in hand

\textbf{Variants (JA)}: \\
\textbf{Variants (EN)}: 

\textbf{Intent (JA)}: \\
\textbf{Intent (EN)}: 

\textbf{Tags (JA)}: \\
\textbf{Tags (EN)}: 

\textbf{Adjusted Expression (JA)}:\\
\\
\textbf{Adjusted Expression (EN)}:\\


\textbf{Example (JA)}:\\
(彼:かれ)は(手取:てと)り(足取:あしと)り(教:おし)えてくれた。\\
\textbf{Example (EN)}:\\
He provided step-by-step guidance.

\textbf{Notes (JA)}:\\
"手取り足取り"は、ステップバイステップで手取り足取り指導されることを表現するために使用されます。実際に、手を触ったり、足を取ったりはしないので、この表現は丁寧に親切に教えてくれたという意味になります。

\textbf{Notes (EN)}:\\
"Tetori ashitori" is used to express that you are guided step by step or hand in hand. In reality, your hands or feet are not touched, so this expression means that you are kindly and carefully taught.

\newpage

\section*{205. Nantoka shinakya}

\textbf{Date}: 2024.11.01 (Fri)\\
\textbf{Index}: nantokashinakya\\
\textbf{Expression (JA)}: なんとかしなきゃ\\
\textbf{Romanization}: Nantoka shinakya\\
\textbf{Expression (EN)}: Have to do something

\textbf{Variants (JA)}: \\
\textbf{Variants (EN)}: 

\textbf{Intent (JA)}: \\
\textbf{Intent (EN)}: 

\textbf{Tags (JA)}: \\
\textbf{Tags (EN)}: 

\textbf{Adjusted Expression (JA)}:\\
\\
\textbf{Adjusted Expression (EN)}:\\


\textbf{Example (JA)}:\\
この(問題:もんだい)、なんとかしなきゃ。\\
\textbf{Example (EN)}:\\
I have to do something about this problem.

\textbf{Notes (JA)}:\\
"なんとかしなきゃ"は、状況や問題について何かしなければならないことを表現するために使用されます。

\textbf{Notes (EN)}:\\
"Nantoka shinakya" is used to express that you have to do something about a situation or problem. The word "nantoka" is an adverb that means "somehow" or "in some way." The word "shinakya" is a contraction of the phrase "shinakereba naranai," which means "must do."

\newpage

\section*{204. mou konna jikann}

\textbf{Date}: 2024.10.31 (Thu)\\
\textbf{Index}: moukonnajikann\\
\textbf{Expression (JA)}: もうこんな(時間:じかん)\\
\textbf{Romanization}: Mou konna jikann\\
\textbf{Expression (EN)}: This time already

\textbf{Variants (JA)}: \\
\textbf{Variants (EN)}: 

\textbf{Intent (JA)}: \\
\textbf{Intent (EN)}: 

\textbf{Tags (JA)}: \\
\textbf{Tags (EN)}: 

\textbf{Adjusted Expression (JA)}:\\
\\
\textbf{Adjusted Expression (EN)}:\\


\textbf{Example (JA)}:\\
もうこんな(時間:じかん)、(帰:かえ)らないと(遅:おく)れるよ。\\
\textbf{Example (EN)}:\\
It's already this time, so if you don't go home now, you'll be late.

\textbf{Notes (JA)}:\\
"もうこんな時間"は、すでにある時間であることを表現するために使用されます。少し驚いて急ぐときに使われます。

\textbf{Notes (EN)}:\\
"Mou konna jikann" is used to express that it is already a certain time. It is used when you are a little surprised and in a hurry. The word "mou" is an adverb that means "already." The word "konna" is a demonstrative pronoun that means "this kind of." The word "jikann" is a noun that means "time."

\newpage

\section*{203. yoyuu}

\textbf{Date}: 2024.10.30 (Wed)\\
\textbf{Index}: yoyuu\\
\textbf{Expression (JA)}: (余裕:よゆう)\\
\textbf{Romanization}: Yoyuu\\
\textbf{Expression (EN)}: Room to spare

\textbf{Variants (JA)}: \\
\textbf{Variants (EN)}: 

\textbf{Intent (JA)}: \\
\textbf{Intent (EN)}: 

\textbf{Tags (JA)}: \\
\textbf{Tags (EN)}: 

\textbf{Adjusted Expression (JA)}:\\
\\
\textbf{Adjusted Expression (EN)}:\\


\textbf{Example (JA)}:\\
(余裕:よゆう)があると、(仕事:しごと)も(楽:らく)に(進:すす)められる。\\
\textbf{Example (EN)}:\\
When you have time, you can work easily.

\textbf{Notes (JA)}:\\
"余裕がある"は、何かをするための時間やスペースがあることを表現するために使用されます。

\textbf{Notes (EN)}:\\
"Yoyuu ga aru" is used to express that you have time or space to do something. The word "yoyuu" is a noun that means "leeway" or "margin." The word "aru" is a verb that means "to exist."

\newpage

\section*{202. doko itteta?}

\textbf{Date}: 2024.10.29 (Tue)\\
\textbf{Index}: dokoitteta\\
\textbf{Expression (JA)}: どこ(行:い)ってた？\\
\textbf{Romanization}: Doko itteta?\\
\textbf{Expression (EN)}: Where have you been?

\textbf{Variants (JA)}: \\
\textbf{Variants (EN)}: 

\textbf{Intent (JA)}: \\
\textbf{Intent (EN)}: 

\textbf{Tags (JA)}: \\
\textbf{Tags (EN)}: 

\textbf{Adjusted Expression (JA)}:\\
\\
\textbf{Adjusted Expression (EN)}:\\


\textbf{Example (JA)}:\\
A「どこ(行:い)ってた？」 B「ちょっと、そこまで」 A「(探:さが)してたよ」 \\
\textbf{Example (EN)}:\\
A: Where have you been? B: Just over there. A: I was looking for you.

\textbf{Notes (JA)}:\\
"どこ行ってた？"は、どこに行ったのかを尋ねるため、なぜ居なかったのかを伝えるために使用されます。

\textbf{Notes (EN)}:\\
"Doko itteta?" is used to ask where someone has been or to tell them that you were looking for them. The word "doko" is a noun that means "where." The word "itta" is the past tense of the verb "iku," which means "to go."

\newpage

\section*{201. doudemo ii}

\textbf{Date}: 2024.10.28 (Mon)\\
\textbf{Index}: doudemoii\\
\textbf{Expression (JA)}: どうでもいい\\
\textbf{Romanization}: Doudemoii\\
\textbf{Expression (EN)}: It doesn't matter

\textbf{Variants (JA)}: \\
\textbf{Variants (EN)}: 

\textbf{Intent (JA)}: \\
\textbf{Intent (EN)}: 

\textbf{Tags (JA)}: \\
\textbf{Tags (EN)}: 

\textbf{Adjusted Expression (JA)}:\\
\\
\textbf{Adjusted Expression (EN)}:\\


\textbf{Example (JA)}:\\
(ゾウ:ぞう)の(体重:たいじゅう)が(何:なん)(キロ:きろ)かなんてどうでもいい。\\
\textbf{Example (EN)}:\\
Who cares how much an elephant weighs?

\textbf{Notes (JA)}:\\
"どうでもいい"は、何かが重要でないことや、気にしないことを表現するために使用されます。

\textbf{Notes (EN)}:\\
"Doudemo ii" is used to express that something doesn't matter or that you don't care about it. The word "doudemo" is an adverb that means "anyway" or "in any case." The word "ii" is an adjective that means "good."

\newpage

\section*{200. Kannpai}

\textbf{Date}: 2024.10.27 (Sun)\\
\textbf{Index}: kannpai\\
\textbf{Expression (JA)}: (乾杯:かんぱい)\\
\textbf{Romanization}: Kanpai\\
\textbf{Expression (EN)}: Cheers

\textbf{Variants (JA)}: \\
\textbf{Variants (EN)}: 

\textbf{Intent (JA)}: \\
\textbf{Intent (EN)}: 

\textbf{Tags (JA)}: \\
\textbf{Tags (EN)}: 

\textbf{Adjusted Expression (JA)}:\\
\\
\textbf{Adjusted Expression (EN)}:\\


\textbf{Example (JA)}:\\
(乾杯:かんぱい)！\\
\textbf{Example (EN)}:\\
Cheers!

\textbf{Notes (JA)}:\\
"乾杯"は、乾杯をする前に乾杯をするために使用されます。

\textbf{Notes (EN)}:\\
"Kanpai" is used to express a toast or to say cheers before drinking. The word "kanpai" is a noun that means "toast" or "cheers." The literal translation of "kanpai" is "dry cup."

\newpage

\section*{199. hitotema kakaru}

\textbf{Date}: 2024.10.26 (Sat)\\
\textbf{Index}: hitotemakakeru\\
\textbf{Expression (JA)}: (一手間:ひとてま)かける\\
\textbf{Romanization}: Hitotema kakaru\\
\textbf{Expression (EN)}: Put in an extra effort

\textbf{Variants (JA)}: \\
\textbf{Variants (EN)}: 

\textbf{Intent (JA)}: \\
\textbf{Intent (EN)}: 

\textbf{Tags (JA)}: \\
\textbf{Tags (EN)}: 

\textbf{Adjusted Expression (JA)}:\\
\\
\textbf{Adjusted Expression (EN)}:\\


\textbf{Example (JA)}:\\
(料理:りょうり)に(一手間:ひとてま)かけると、(味:あじ)が(ぐっと)(良:よ)くなる。\\
\textbf{Example (EN)}:\\
Putting in an extra effort in cooking makes the taste much better.

\textbf{Notes (JA)}:\\
"一手間かける"は、何かを改善するために余分な努力を払ったり、余分なステップを踏んだりすることを表現するために使用されます。

\textbf{Notes (EN)}:\\
"Hitotema kakaru" is used to express that you put in an extra effort or take an extra step to improve something. The word "hitotema" is a noun that means "extra effort" or "extra step." The word "kakaru" is a verb that means "to hang" or "to put in."

\newpage

\section*{198. Ima sugu}

\textbf{Date}: 2024.10.25 (Fri)\\
\textbf{Index}: imasugu\\
\textbf{Expression (JA)}: (今:いま)すぐ\\
\textbf{Romanization}: Ima sugu\\
\textbf{Expression (EN)}: Right now

\textbf{Variants (JA)}: \\
\textbf{Variants (EN)}: 

\textbf{Intent (JA)}: \\
\textbf{Intent (EN)}: 

\textbf{Tags (JA)}: \\
\textbf{Tags (EN)}: 

\textbf{Adjusted Expression (JA)}:\\
\\
\textbf{Adjusted Expression (EN)}:\\


\textbf{Example (JA)}:\\
(今:いま)すぐ(行:い)かないと(遅:おく)れるよ。\\
\textbf{Example (EN)}:\\
We'll be late if we don't leave right now.

\textbf{Notes (JA)}:\\
"今すぐ"は、何かをすぐにやらなければならないことを表現するために使用されます。

\textbf{Notes (EN)}:\\
"Ima sugu" is used to express that you need to do something immediately or right away. The word "ima" is a noun that means "now." The word "sugu" is an adverb that means "immediately" or "right away."

\newpage

\section*{197. ki ni iru}

\textbf{Date}: 2024.10.24 (Thu)\\
\textbf{Index}: kiniru\\
\textbf{Expression (JA)}: (気:き)に(入:い)る\\
\textbf{Romanization}: Ki ni iru\\
\textbf{Expression (EN)}: To like

\textbf{Variants (JA)}: \\
\textbf{Variants (EN)}: 

\textbf{Intent (JA)}: \\
\textbf{Intent (EN)}: 

\textbf{Tags (JA)}: \\
\textbf{Tags (EN)}: 

\textbf{Adjusted Expression (JA)}:\\
\\
\textbf{Adjusted Expression (EN)}:\\


\textbf{Example (JA)}:\\
(今日:きょう)の(服:ふく)、(気:き)に(入:い)っている。\\
\textbf{Example (EN)}:\\
I like today's outfit.

\textbf{Notes (JA)}:\\
"気に入る"は、何かを好きだとか、何かが目に留まることを表現するために使用されます。

\textbf{Notes (EN)}:\\
"Ki ni iru" is used to express that you like something or that something catches your eye. The word "ki" is a noun that means "spirit" or "mind." The word "iru" is a verb that means "to be."

\newpage

\section*{196. Motto yarou}

\textbf{Date}: 2024.10.23 (Wed)\\
\textbf{Index}: mottoyarou\\
\textbf{Expression (JA)}: もっとやろう\\
\textbf{Romanization}: Motto yarou\\
\textbf{Expression (EN)}: Let's do more

\textbf{Variants (JA)}: \\
\textbf{Variants (EN)}: 

\textbf{Intent (JA)}: \\
\textbf{Intent (EN)}: 

\textbf{Tags (JA)}: \\
\textbf{Tags (EN)}: 

\textbf{Adjusted Expression (JA)}:\\
\\
\textbf{Adjusted Expression (EN)}:\\


\textbf{Example (JA)}:\\
まだ(時間:じかん)があるから、もっとやろう。\\
\textbf{Example (EN)}:\\
We still have time, so let's keep going.

\textbf{Notes (JA)}:\\
「もっとやろう」は、誰かにもっとやるように促すために使用されます。

\textbf{Notes (EN)}:\\
The phrase "motto yarou" is used to encourage someone to do more or to keep going. The word "motto" is an adverb that means "more." The word "yarou" is a verb that means "to do."

\newpage

\section*{195. That's tough}

\textbf{Date}: 2024.10.22 (Tue)\\
\textbf{Index}: sorewanai\\
\textbf{Expression (JA)}: それはたいへん\\
\textbf{Romanization}: Sore wa taihen\\
\textbf{Expression (EN)}: That's tough

\textbf{Variants (JA)}: \\
\textbf{Variants (EN)}: 

\textbf{Intent (JA)}: \\
\textbf{Intent (EN)}: 

\textbf{Tags (JA)}: \\
\textbf{Tags (EN)}: 

\textbf{Adjusted Expression (JA)}:\\
\\
\textbf{Adjusted Expression (EN)}:\\


\textbf{Example (JA)}:\\
A「(昨日:きのう)、(徹夜:てつや)で(仕事:しごと)を(終:お)わらせた」 B「それはたいへん」\\
\textbf{Example (EN)}:\\
A: I stayed up all night to finish my work yesterday. B: That's tough.

\textbf{Notes (JA)}:\\
「それは大変」は、誰かが困難な状況やチャレンジングな状況について話すときに、同情や理解を表現するために使用されます。

\textbf{Notes (EN)}:\\
The phrase "sore wa taihen" is used to express sympathy or understanding when someone tells you about a difficult or challenging situation. The word "sore" is a pronoun that means "that." The word "wa" is a particle that is used to indicate the topic of a sentence. The word "taihen" is a noun that means "tough" or "difficult."

\newpage

\section*{194. Bikkuri shita}

\textbf{Date}: 2024.10.21 (Mon)\\
\textbf{Index}: bikkurishita\\
\textbf{Expression (JA)}: びっくりした\\
\textbf{Romanization}: Bikkuri shita\\
\textbf{Expression (EN)}: Be surprised

\textbf{Variants (JA)}: \\
\textbf{Variants (EN)}: 

\textbf{Intent (JA)}: \\
\textbf{Intent (EN)}: 

\textbf{Tags (JA)}: \\
\textbf{Tags (EN)}: 

\textbf{Adjusted Expression (JA)}:\\
\\
\textbf{Adjusted Expression (EN)}:\\


\textbf{Example (JA)}:\\
うわっ、びっくりした！心臓止まるかと思った。\\
\textbf{Example (EN)}:\\
Wow, that scared me! I thought my heart would stop.

\textbf{Notes (JA)}:\\
「びっくりした」は、突然の出来事に驚いたときに使えます。

\textbf{Notes (EN)}:\\
The phrase "bikkuri shita" is used to express that you are surprised or startled by something. The word "bikkuri" is a noun that means "surprise" or "shock." The word "shita" is the past tense of the verb "suru," which means "to do."

\newpage

\section*{193. Get through it}

\textbf{Date}: 2024.10.20 (Sun)\\
\textbf{Index}: norikiru\\
\textbf{Expression (JA)}: (乗:の)り(切:き)る\\
\textbf{Romanization}: Norikiru\\
\textbf{Expression (EN)}: Get through

\textbf{Variants (JA)}: \\
\textbf{Variants (EN)}: 

\textbf{Intent (JA)}: \\
\textbf{Intent (EN)}: 

\textbf{Tags (JA)}: \\
\textbf{Tags (EN)}: 

\textbf{Adjusted Expression (JA)}:\\
\\
\textbf{Adjusted Expression (EN)}:\\


\textbf{Example (JA)}:\\
(大変:たいへん)な(仕事:しごと)だったけど、(何:なん)とか(乗:の)り(切:き)った。\\
\textbf{Example (EN)}:\\
It was a tough job, but I managed to get through it.

\textbf{Notes (JA)}:\\
「乗り切る」は、困難な仕事や状況を成功裏に完了したことを表現するために使用されます。

\textbf{Notes (EN)}:\\
The phrase "nori kiru" is used to express that you have successfully completed a difficult task or situation. The word "nori" is a noun that means "ride" or "get on." The word "kiru" is a verb that means "to cut" or "to get through."

\newpage

\section*{192. On the edge of a cliff}

\textbf{Date}: 2024.10.19 (Sat)\\
\textbf{Index}: gakeppuchi\\
\textbf{Expression (JA)}: (崖:がけ)っぷち\\
\textbf{Romanization}: Gakeppuchi\\
\textbf{Expression (EN)}: At the edge of a cliff

\textbf{Variants (JA)}: \\
\textbf{Variants (EN)}: 

\textbf{Intent (JA)}: \\
\textbf{Intent (EN)}: 

\textbf{Tags (JA)}: \\
\textbf{Tags (EN)}: 

\textbf{Adjusted Expression (JA)}:\\
\\
\textbf{Adjusted Expression (EN)}:\\


\textbf{Example (JA)}:\\
もうこれ(以上:いじょう)、(時間:じかん)もお(金:かね)もかけられない(状況:じょうきょう)に(追:お)い(込:こ)まれて、(彼:かれ)は、(仕事:しごと)で(崖:がけ)っぷちに(立:た)たされている」\\
\textbf{Example (EN)}:\\
He's been driven into a situation where he can't spend any more time or money, and he's on the edge of a cliff in his work.

\textbf{Notes (JA)}:\\
「崖っぷち」は、行動を取らざるを得ない困難な状況にあることを表現するために使用されます。

\textbf{Notes (EN)}:\\
The phrase "gakeppuchi" is used to express that you are in a difficult situation where you have no choice but to take action. The word "gake" is a noun that means "cliff." The word "ppuchi" is a noun that means "edge."

\newpage

\section*{191. Kokoro ga odoru}

\textbf{Date}: 2024.10.18 (Fri)\\
\textbf{Index}: kokorogaodoru\\
\textbf{Expression (JA)}: (心:こころ)が(躍:おど)る\\
\textbf{Romanization}: Kokoro ga odoru\\
\textbf{Expression (EN)}: Be excited

\textbf{Variants (JA)}: \\
\textbf{Variants (EN)}: 

\textbf{Intent (JA)}: \\
\textbf{Intent (EN)}: 

\textbf{Tags (JA)}: \\
\textbf{Tags (EN)}: 

\textbf{Adjusted Expression (JA)}:\\
\\
\textbf{Adjusted Expression (EN)}:\\


\textbf{Example (JA)}:\\
(彼:かれ)は、(新:あたら)しい(仕事:しごと)に、(心:こころ)が(躍:おど)った。\\
\textbf{Example (EN)}:\\
He was excited about his new job.

\textbf{Notes (JA)}:\\
「心が躍る」は、何かに興奮したり、わくわくしたりすることを表現するために使用されます。

\textbf{Notes (EN)}:\\
The phrase "kokoro ga odoru" is used to express that you are excited or thrilled about something. The word "kokoro" is a noun that means "heart" or "mind." The word "odoru" is a verb that means "to dance" or "to jump."

\newpage

\section*{190. The sandman is coming}

\textbf{Date}: 2024.10.17 (Thu)\\
\textbf{Index}: nemuke\\
\textbf{Expression (JA)}: (眠気:ねむけ)\\
\textbf{Romanization}: Nemuke\\
\textbf{Expression (EN)}: Sleepy

\textbf{Variants (JA)}: \\
\textbf{Variants (EN)}: 

\textbf{Intent (JA)}: \\
\textbf{Intent (EN)}: 

\textbf{Tags (JA)}: \\
\textbf{Tags (EN)}: 

\textbf{Adjusted Expression (JA)}:\\
\\
\textbf{Adjusted Expression (EN)}:\\


\textbf{Example (JA)}:\\
(昼:ひる)ご(飯:はん)を(食:た)べたあと、(眠気:ねむけ)が(来:き)た。\\
\textbf{Example (EN)}:\\
I was about to crash/felt sleepy after eating lunch.

\textbf{Notes (JA)}:\\
「眠気」は、眠い、あくびが出る、だるいなど、眠気を感じることを表現するために使用されます。

\textbf{Notes (EN)}:\\
The word "nemuke" is used to express that are sleepy or drowsy. The word "nemuke" is a noun that means "sleepiness."

\newpage

\section*{189. Itami ga hashiru}

\textbf{Date}: 2024.10.16 (Wed)\\
\textbf{Index}: itamigahashiru\\
\textbf{Expression (JA)}: (痛:いた)みが(走:はし)る\\
\textbf{Romanization}: Itami ga hashiru\\
\textbf{Expression (EN)}: Sharp pain

\textbf{Variants (JA)}: \\
\textbf{Variants (EN)}: 

\textbf{Intent (JA)}: \\
\textbf{Intent (EN)}: 

\textbf{Tags (JA)}: \\
\textbf{Tags (EN)}: 

\textbf{Adjusted Expression (JA)}:\\
\\
\textbf{Adjusted Expression (EN)}:\\


\textbf{Example (JA)}:\\
(歯:は)に(痛:いた)みが(走:はし)った。\\
\textbf{Example (EN)}:\\
I felt a sharp pain in my tooth.

\textbf{Notes (JA)}:\\
「痛みが走る」は、体の特定の部位に鋭い痛みを感じることを表現するために使用されます。

\textbf{Notes (EN)}:\\
The phrase "itami ga hashiru" is used to express that you feel a sharp pain in a specific part of your body. The word "itami" is a noun that means "pain." The word "hashiru" is a verb that means "to run."

\newpage

\section*{188. Kuuki wo yomu}

\textbf{Date}: 2024.10.15 (Tue)\\
\textbf{Index}: kuukiwoyomu\\
\textbf{Expression (JA)}: (空気:くうき)を(読:よ)む\\
\textbf{Romanization}: Kuuki wo yomu\\
\textbf{Expression (EN)}: Read the room

\textbf{Variants (JA)}: \\
\textbf{Variants (EN)}: 

\textbf{Intent (JA)}: \\
\textbf{Intent (EN)}: 

\textbf{Tags (JA)}: \\
\textbf{Tags (EN)}: 

\textbf{Adjusted Expression (JA)}:\\
\\
\textbf{Adjusted Expression (EN)}:\\


\textbf{Example (JA)}:\\
A「(今:いま)、(話:はなし)を(続:つづ)けても(大丈夫:だいじょうぶ)かな？」 B「(空気:くうき)を(読:よ)んで」\\
\textbf{Example (EN)}:\\
A: Is it okay to continue the conversation now? B: Read the room.

\textbf{Notes (JA)}:\\
「空気を読む」は、状況の雰囲気やムードを把握し、それに応じて行動することを表現するために使用されます。

\textbf{Notes (EN)}:\\
The phrase "kuuki wo yomu" is used to express that you are aware of the atmosphere or mood of a situation and act accordingly. The word "kuuki" is a noun that means "air" or "atmosphere." The word "yomu" is a verb that means "to read."

\newpage

\section*{187. netsu ga sameru}

\textbf{Date}: 2024.10.14 (Mon)\\
\textbf{Index}: netsugasameru\\
\textbf{Expression (JA)}: (熱:ねつ)が(冷:さ)める\\
\textbf{Romanization}: Netsu ga sameru\\
\textbf{Expression (EN)}: Lose enthusiasm

\textbf{Variants (JA)}: \\
\textbf{Variants (EN)}: 

\textbf{Intent (JA)}: \\
\textbf{Intent (EN)}: 

\textbf{Tags (JA)}: \\
\textbf{Tags (EN)}: 

\textbf{Adjusted Expression (JA)}:\\
\\
\textbf{Adjusted Expression (EN)}:\\


\textbf{Example (JA)}:\\
(新:あたら)しい(趣味:しゅみ)に(夢中:むちゅう)だったが、すぐに(熱:ねつ)が(冷:さ)めてしまった。\\
\textbf{Example (EN)}:\\
I was really into my new hobby, but my enthusiasm cooled off quickly.

\textbf{Notes (JA)}:\\
「熱が冷める」は、「熱心さがなくなる」の意味。趣味や仕事など▽はじめは熱心に取り組んでいたが、だんだん興味を失ってしまったときに使う。もちろん、本当に熱が冷めるときもあるが、多くは興味を失うときに使われる。では、風邪をひいて、熱が出たときは「熱が下がる」とか「熱はもうない」と言う。

\textbf{Notes (EN)}:\\
The phrase "netsu ga sameru" means to lose enthusiasm. It is used when you lose interest in a hobby or work. Of course, it can be used when you really lose enthusiasm, but it is often used when you lose interest. When you catch a cold and have a fever, you say "netsu ga sagaru" or "netsu wa mou nai."

\newpage

\section*{186. Te wo nuku}

\textbf{Date}: 2024.10.13 (Sun)\\
\textbf{Index}: tewonuku\\
\textbf{Expression (JA)}: (手:て)を(抜:ぬ)く\\
\textbf{Romanization}: Te wo nuku\\
\textbf{Expression (EN)}: Cut corners

\textbf{Variants (JA)}: \\
\textbf{Variants (EN)}: 

\textbf{Intent (JA)}: \\
\textbf{Intent (EN)}: 

\textbf{Tags (JA)}: \\
\textbf{Tags (EN)}: 

\textbf{Adjusted Expression (JA)}:\\
\\
\textbf{Adjusted Expression (EN)}:\\


\textbf{Example (JA)}:\\
(ビル:びる)(建設:けんせつ)で、(手:て)を(抜:ぬ)いた(仕事:しごと)をした(結果:けっか)、(大事故:だいじこ)が(起:おこ)った。\\
\textbf{Example (EN)}:\\
A major accident occurred in the construction of the building as a result of the work done carelessly.

\textbf{Notes (JA)}:\\
「手を抜く」は、何かを手抜きしていい加減にすることを表現するために使用されます。

\textbf{Notes (EN)}:\\
The phrase "te wo nuku" means to do something carelessly or to cut corners. The word "te" is a noun that means "hand." The word "nuku" is a verb that means "to pull out."

\newpage

\section*{185. Be helpful}

\textbf{Date}: 2024.10.12 (Sat)\\
\textbf{Index}: tewosashinoberu\\
\textbf{Expression (JA)}: (手:て)を(差:さ)し(伸:の)べる\\
\textbf{Romanization}: Te wo sashinoberu\\
\textbf{Expression (EN)}: Be helpful

\textbf{Variants (JA)}: \\
\textbf{Variants (EN)}: 

\textbf{Intent (JA)}: \\
\textbf{Intent (EN)}: 

\textbf{Tags (JA)}: \\
\textbf{Tags (EN)}: 

\textbf{Adjusted Expression (JA)}:\\
\\
\textbf{Adjusted Expression (EN)}:\\


\textbf{Example (JA)}:\\
(彼:かれ)は、(困:こま)っている(人:ひと)を(見:み)ると、(手:て)を(差:さ)し(伸:の)べる。\\
\textbf{Example (EN)}:\\
He offers help to people in need when he sees them.

\textbf{Notes (JA)}:\\
"手を差し伸べる"は、誰か困っている人の助けを提供することを意味する表現です。

\textbf{Notes (EN)}:\\
"Te wo sashinoberu" is used to express that you offer help or support to someone in need. The word "te" is a noun that means "hand." The word "sashinoberu" is a verb that means "to stretch out."

\newpage

\section*{184. Smile}

\textbf{Date}: 2024.10.11 (Fri)\\
\textbf{Index}: mewohosomeru\\
\textbf{Expression (JA)}: (目:め)を(細:ほそ)める\\
\textbf{Romanization}: Me wo hosomeru\\
\textbf{Expression (EN)}: To smile gently

\textbf{Variants (JA)}: \\
\textbf{Variants (EN)}: 

\textbf{Intent (JA)}: \\
\textbf{Intent (EN)}: 

\textbf{Tags (JA)}: \\
\textbf{Tags (EN)}: 

\textbf{Adjusted Expression (JA)}:\\
\\
\textbf{Adjusted Expression (EN)}:\\


\textbf{Example (JA)}:\\
(彼女:かのじょ)は(孫:まご)の(成長:せいちょう)を(喜:よろこ)んで、(優:やさ)しく(目:め)を(細:ほそ)めた。\\
\textbf{Example (EN)}:\\
She smiled warmly as she watched her grandchild grow.

\textbf{Notes (JA)}:\\
「目を細める」は、「微笑む」という意味で使用されます。

\textbf{Notes (EN)}:\\
"Me wo hosomeru" is used to express that one smiles gently or warmly while watching something or someone. The word "me" is a noun that means "eye." The word "hosomeru" is a verb that means "to narrow."

\newpage

\section*{183. To be in a hurry}

\textbf{Date}: 2024.10.10 (Thu)\\
\textbf{Index}: isogu\\
\textbf{Expression (JA)}: (急:いそ)ぐ\\
\textbf{Romanization}: Isogu\\
\textbf{Expression (EN)}: Hurry

\textbf{Variants (JA)}: \\
\textbf{Variants (EN)}: 

\textbf{Intent (JA)}: \\
\textbf{Intent (EN)}: 

\textbf{Tags (JA)}: \\
\textbf{Tags (EN)}: 

\textbf{Adjusted Expression (JA)}:\\
\\
\textbf{Adjusted Expression (EN)}:\\


\textbf{Example (JA)}:\\
(急:いそ)いで！(遅:おく)れるよ。\\
\textbf{Example (EN)}:\\
Hurry up! You'll be late.

\textbf{Notes (JA)}:\\
"急ぐ"は動詞です。

\textbf{Notes (EN)}:\\
"Isogu" is used to express that you are in a hurry or rushing to do something. The word "isogu" is a verb that means "to hurry" or "to rush."

\newpage

\section*{182. kokoro wo ubau}

\textbf{Date}: 2024.10.09 (Wed)\\
\textbf{Index}: kokorowoubau\\
\textbf{Expression (JA)}: (心:こころ)を(奪:うば)う\\
\textbf{Romanization}: Kokoro wo ubau\\
\textbf{Expression (EN)}: Capture one's heart

\textbf{Variants (JA)}: \\
\textbf{Variants (EN)}: 

\textbf{Intent (JA)}: \\
\textbf{Intent (EN)}: 

\textbf{Tags (JA)}: \\
\textbf{Tags (EN)}: 

\textbf{Adjusted Expression (JA)}:\\
\\
\textbf{Adjusted Expression (EN)}:\\


\textbf{Example (JA)}:\\
(彼女:かのじょ)の(歌:うた)は、(聴:き)く(人:ひと)の(心:こころ)を(奪:うば)う。\\
\textbf{Example (EN)}:\\
Her singing captures the hearts of those who listen.

\textbf{Notes (JA)}:\\
"心を奪う"は、何かが誰かの心を魅了したり、深く感動させたりすることを表現するために使用されます。

\textbf{Notes (EN)}:\\
"Kokoro wo ubau" is used to express that something captivates or deeply moves someone's heart. The word "kokoro" is a noun that means "heart" or "mind." The word "ubau" is a verb that means "to snatch" or "to capture."

\newpage

\section*{181. Timid}

\textbf{Date}: 2024.10.08 (Tue)\\
\textbf{Index}: kigachiisai\\
\textbf{Expression (JA)}: (気:き)が(小:ちい)さい\\
\textbf{Romanization}: Ki ga chiisai\\
\textbf{Expression (EN)}: Timid

\textbf{Variants (JA)}: \\
\textbf{Variants (EN)}: 

\textbf{Intent (JA)}: \\
\textbf{Intent (EN)}: 

\textbf{Tags (JA)}: \\
\textbf{Tags (EN)}: 

\textbf{Adjusted Expression (JA)}:\\
\\
\textbf{Adjusted Expression (EN)}:\\


\textbf{Example (JA)}:\\
(気:き)が(小:ちい)さい(人:ひと)は、(失敗:しっぱい)を(恐:おそ)れて(チャレンジ:ちゃれんじ)しないことが(多:おお)い。\\
\textbf{Example (EN)}:\\
Timid people often avoid challenges because they fear failure.

\textbf{Notes (JA)}:\\
"気が小さい"は、誰かが臆病で勇気がないことを表現するために使用されます。

\textbf{Notes (EN)}:\\
"Ki ga chiisai" is used to express that someone is timid or lacks courage. The word "ki" is a noun that means "spirit" or "mind." The word "chiisai" is an i-adjective that means "small."

\newpage

\section*{180. Inumerable}

\textbf{Date}: 2024.10.07 (Mon)\\
\textbf{Index}: kazoekirenai\\
\textbf{Expression (JA)}: (数:かぞ)え(切:き)れない\\
\textbf{Romanization}: Kazoekirenai\\
\textbf{Expression (EN)}: Countless

\textbf{Variants (JA)}: \\
\textbf{Variants (EN)}: 

\textbf{Intent (JA)}: \\
\textbf{Intent (EN)}: 

\textbf{Tags (JA)}: \\
\textbf{Tags (EN)}: 

\textbf{Adjusted Expression (JA)}:\\
\\
\textbf{Adjusted Expression (EN)}:\\


\textbf{Example (JA)}:\\
(数:かぞ)え(切:き)れない(星:ほし)が(夜空:よぞら)に(輝:かがや)いています。\\
\textbf{Example (EN)}:\\
Countless stars are shining in the night sky.

\textbf{Notes (JA)}:\\
"数え切れない"は、何かが数えられないほど多いことを表現するために使用されます。

\textbf{Notes (EN)}:\\
"Kazoekirenai" is used to express that something is so numerous that it cannot be counted. The word "kazoeru" is a verb that means "to count." The word "kirenai" is a negative form of the verb "kiru" that means "to be able to do."

\newpage

\section*{179. Iu made mo nai}

\textbf{Date}: 2024.10.06 (Sun)\\
\textbf{Index}: iumademonai\\
\textbf{Expression (JA)}: (言:い)うまでもなく\\
\textbf{Romanization}: Iumademonaku\\
\textbf{Expression (EN)}: Needless to say

\textbf{Variants (JA)}: \\
\textbf{Variants (EN)}: 

\textbf{Intent (JA)}: \\
\textbf{Intent (EN)}: 

\textbf{Tags (JA)}: \\
\textbf{Tags (EN)}: 

\textbf{Adjusted Expression (JA)}:\\
\\
\textbf{Adjusted Expression (EN)}:\\


\textbf{Example (JA)}:\\
(言:い)うまでもなく、その(結果:けっか)は(私:わたし)たちにとって(非常:ひじょう)に(重要:じゅうよう)です。\\
\textbf{Example (EN)}:\\
Needless to say, the result is extremely important for us.

\textbf{Notes (JA)}:\\
「言うまでもなく」は、何かが非常に明らかで、言う必要がないことを表現するために使用されます。

\textbf{Notes (EN)}:\\
"Iumademonaku" is used to express that something is so obvious that it does not need to be said. The word "iu" is a verb that means "to say." The word "mademo" is a conjunction that means "even if." The word "naku" is a negative suffix that is used to negate the verb that comes before it.

\newpage

\section*{178. Ki ga suru}

\textbf{Date}: 2024.10.05 (Sat)\\
\textbf{Index}: kigasuru\\
\textbf{Expression (JA)}: (気:き)が(する)\\
\textbf{Romanization}: Ki ga suru\\
\textbf{Expression (EN)}: Feel like

\textbf{Variants (JA)}: \\
\textbf{Variants (EN)}: 

\textbf{Intent (JA)}: \\
\textbf{Intent (EN)}: 

\textbf{Tags (JA)}: \\
\textbf{Tags (EN)}: 

\textbf{Adjusted Expression (JA)}:\\
\\
\textbf{Adjusted Expression (EN)}:\\


\textbf{Example (JA)}:\\
この(問題:もんだい)、どこかで(見:み)たことがある(気:き)がする。\\
\textbf{Example (EN)}:\\
I feel like I've seen this problem somewhere before.

\textbf{Notes (JA)}:\\
「気がする」は、確信がないものの、直感的に何かを予感したり、漠然と感じ取る場合に使用されます。確実ではないけれど、何となくそう思う、というニュアンスが含まれています。

\textbf{Notes (EN)}:\\
"Ki ga suru" is used when you intuitively sense or vaguely feel something, even though you are not certain about it. It carries a nuance of uncertainty, suggesting that you feel something might be true without having concrete evidence. The word "ki" is a noun that means "spirit" or "mind." The word "suru" is a verb that means "to do."

\newpage

\section*{177. Te ni ireru}

\textbf{Date}: 2024.10.04 (Fri)\\
\textbf{Index}: teniireru\\
\textbf{Expression (JA)}: (手:て)に(入:い)れる\\
\textbf{Romanization}: Te ni ireru\\
\textbf{Expression (EN)}: Get/obtain

\textbf{Variants (JA)}: \\
\textbf{Variants (EN)}: 

\textbf{Intent (JA)}: \\
\textbf{Intent (EN)}: 

\textbf{Tags (JA)}: \\
\textbf{Tags (EN)}: 

\textbf{Adjusted Expression (JA)}:\\
\\
\textbf{Adjusted Expression (EN)}:\\


\textbf{Example (JA)}:\\
(彼:かれ)は(長年:ながねん)の(努力:どりょく)の(末:すえ)、ようやく(夢:ゆめ)の(家:いえ)を(手:て)に(入:い)れた。\\
\textbf{Example (EN)}:\\
After years of hard work, he finally obtained his dream house.

\textbf{Notes (JA)}:\\
"手に入れる"は、探し求めたり、努力して自分のものにすることを表現するために使用されます。単純に物を買ったり、もらったりすることではありません。

\textbf{Notes (EN)}:\\
"Te ni ireru" is used to express obtaining something after searching for it or making an effort to make it one's own. It is not simply about buying or receiving something.

\newpage

\section*{176. Mimi ni suru}

\textbf{Date}: 2024.10.03 (Thu)\\
\textbf{Index}: miminisuru\\
\textbf{Expression (JA)}: (耳:みみ)に(する)\\
\textbf{Romanization}: Mimi ni suru\\
\textbf{Expression (EN)}: hear mentioned frequently

\textbf{Variants (JA)}: \\
\textbf{Variants (EN)}: 

\textbf{Intent (JA)}: \\
\textbf{Intent (EN)}: 

\textbf{Tags (JA)}: \\
\textbf{Tags (EN)}: 

\textbf{Adjusted Expression (JA)}:\\
\\
\textbf{Adjusted Expression (EN)}:\\


\textbf{Example (JA)}:\\
A「(最近:さいきん)、(その:その)(問題:もんだい)を(耳:みみ)にする」\\
\textbf{Example (EN)}:\\
A: I hear about that issue frequently these days.

\textbf{Notes (JA)}:\\
「耳にする」は、何かを頻繁に耳にすることを表現するために使用されます。

\textbf{Notes (EN)}:\\
The phrase "mimi ni suru" is used to express that you frequently hear about something. The word "mimi" is a noun that means "ear." The word "suru" is a verb that means "to do." This idiomatic phrase, "mimi ni suru," is to hear something mentioned or referred to on multiple occasions in everyday life, often indirectly or casually, through conversations, media, or public discourse.

\newpage

\section*{175. Kuchi ni suru}

\textbf{Date}: 2024.10.02 (Wed)\\
\textbf{Index}: kuchinisuru\\
\textbf{Expression (JA)}: (口:くち)に(する)\\
\textbf{Romanization}: Kuchi ni suru\\
\textbf{Expression (EN)}: mention

\textbf{Variants (JA)}: \\
\textbf{Variants (EN)}: 

\textbf{Intent (JA)}: \\
\textbf{Intent (EN)}: 

\textbf{Tags (JA)}: \\
\textbf{Tags (EN)}: 

\textbf{Adjusted Expression (JA)}:\\
\\
\textbf{Adjusted Expression (EN)}:\\


\textbf{Example (JA)}:\\
(彼:かれ)はその(事件:じけん)について(一切:いっさい)(口:くち)にしなかった。\\
\textbf{Example (EN)}:\\
He didn't mention the incident at all.

\textbf{Notes (JA)}:\\
「口にする」は、何かを話題にしたり、言及したりすることを表現しますが、否定文で使われることが多く、特に意図的に触れない、または言わないという意味を持つことが一般的です。また、食べ物や飲み物を口に入れる意味でも使われます。

\textbf{Notes (EN)}:\\
"Kuchi ni suru" refers to mentioning or bringing up a topic, often used in negative sentences to express deliberately avoiding a subject or not saying something. It can also mean putting food or drink into one's mouth. The word "kuchi" is a noun that means "mouth." The word "suru" is a verb that means "to do."

\newpage

\section*{174. Ki ga kiku}

\textbf{Date}: 2024.10.01 (Tue)\\
\textbf{Index}: kigakiku\\
\textbf{Expression (JA)}: (気:き)が(利:き)く\\
\textbf{Romanization}: Ki ga kiku\\
\textbf{Expression (EN)}: Be attentive

\textbf{Variants (JA)}: \\
\textbf{Variants (EN)}: 

\textbf{Intent (JA)}: \\
\textbf{Intent (EN)}: 

\textbf{Tags (JA)}: \\
\textbf{Tags (EN)}: 

\textbf{Adjusted Expression (JA)}:\\
\\
\textbf{Adjusted Expression (EN)}:\\


\textbf{Example (JA)}:\\
A「この(家:いえ)の(デザイン:でざいん)は、(使:つか)いやすくて、なかなか(気:き)が(利:き)いているね」\\
\textbf{Example (EN)}:\\
A: The design of this house is user-friendly and quite attentive.

\textbf{Notes (JA)}:\\
「気が利く」は、何かが気配りができることを表現するために使用されます。

\textbf{Notes (EN)}:\\
The phrase "ki ga kiku" is used to express that something is attentive or considerate. The word "ki" is a noun that means "spirit" or "mind." The word "kiku" is a verb that means "to be effective" or "to be useful."

\newpage

\section*{173. Itsunomanika}

\textbf{Date}: 2024.09.30 (Mon)\\
\textbf{Index}: itsunomanika\\
\textbf{Expression (JA)}: いつのまにか\\
\textbf{Romanization}: Itsunomanika\\
\textbf{Expression (EN)}: while one is unaware of it

\textbf{Variants (JA)}: \\
\textbf{Variants (EN)}: 

\textbf{Intent (JA)}: \\
\textbf{Intent (EN)}: 

\textbf{Tags (JA)}: \\
\textbf{Tags (EN)}: 

\textbf{Adjusted Expression (JA)}:\\
\\
\textbf{Adjusted Expression (EN)}:\\


\textbf{Example (JA)}:\\
A「(いつのまにか)、(夕方:ゆうがた)になっていた」\\
\textbf{Example (EN)}:\\
A: It was evening before I knew it.

\textbf{Notes (JA)}:\\
「いつのまにか」は、何かが気づかないうちに起こったことを表現するために使用されます。

\textbf{Notes (EN)}:\\
The phrase "itsunomanika" is used to express that something happened without you realizing it. The word "itsu" is a pronoun that means "when." The word "no" is a particle that is used to indicate possession or connection. The word "ma" is a noun that means "moment", and the word "ni" is an adverbial particle that is used to indicate time or place. The word "ka" is a question particle that is used to indicate a question.

\newpage

\section*{172. Effort}

\textbf{Date}: 2024.09.29 (Sun)\\
\textbf{Index}: chikarawoireru\\
\textbf{Expression (JA)}: (力:ちから)を(入:い)れる\\
\textbf{Romanization}: Chikara wo ireru\\
\textbf{Expression (EN)}: Put effort into

\textbf{Variants (JA)}: \\
\textbf{Variants (EN)}: 

\textbf{Intent (JA)}: \\
\textbf{Intent (EN)}: 

\textbf{Tags (JA)}: \\
\textbf{Tags (EN)}: 

\textbf{Adjusted Expression (JA)}:\\
\\
\textbf{Adjusted Expression (EN)}:\\


\textbf{Example (JA)}:\\
A「(当店:とうてん)の(主力商品:しゅりょくしょうひん)の(売り込み:うりこみ)に(力:ちから)を(入:い)れている」\\
\textbf{Example (EN)}:\\
A: We are putting effort into promoting our main products.

\textbf{Notes (JA)}:\\
「力を入れる」は、何かに努力を入れていることを表現するために使用されます。

\textbf{Notes (EN)}:\\
The phrase "chikara wo ireru" is used to express that you are putting effort into something. The word "chikara" is a noun that means "power" or "strength." The word "ireru" is a verb that means "to put in." Shuryoku shouhin is a noun that means "main product." Urikomi is a noun that means "promotion."

\newpage

\section*{171. Last hope}

\textbf{Date}: 2024.09.28 (Sat)\\
\textbf{Index}: tanominotsuna\\
\textbf{Expression (JA)}: (頼:たの)みの(綱:つな)\\
\textbf{Romanization}: Tanomi no tsuna\\
\textbf{Expression (EN)}: Last hope

\textbf{Variants (JA)}: \\
\textbf{Variants (EN)}: 

\textbf{Intent (JA)}: \\
\textbf{Intent (EN)}: 

\textbf{Tags (JA)}: \\
\textbf{Tags (EN)}: 

\textbf{Adjusted Expression (JA)}:\\
\\
\textbf{Adjusted Expression (EN)}:\\


\textbf{Example (JA)}:\\
A「(頼:たの)みの(綱:つな)だ」\\
\textbf{Example (EN)}:\\
A: It's our last hope.

\textbf{Notes (JA)}:\\
「頼みの綱」は、何かが最後の希望であることや問題を解決する唯一の方法であることを表現するために使用されます。

\textbf{Notes (EN)}:\\
The phrase "tanomi no tsuna" is used to express that something is the last hope or the only way to solve a problem. The word "tanomi" is a noun that means "request" or "hope." The word "tsuna" is a noun that means "rope."

\newpage

\section*{170. katadukeru}

\textbf{Date}: 2024.09.27 (Fri)\\
\textbf{Index}: katadukeru\\
\textbf{Expression (JA)}: (片付:かたづ)ける\\
\textbf{Romanization}: Katadukeru\\
\textbf{Expression (EN)}: Tidy up

\textbf{Variants (JA)}: \\
\textbf{Variants (EN)}: 

\textbf{Intent (JA)}: \\
\textbf{Intent (EN)}: 

\textbf{Tags (JA)}: \\
\textbf{Tags (EN)}: 

\textbf{Adjusted Expression (JA)}:\\
\\
\textbf{Adjusted Expression (EN)}:\\


\textbf{Example (JA)}:\\
A「(一:ひと)つ(一:ひと)つを(片付:かたづ)けていこう」\\
\textbf{Example (EN)}:\\
A: Let's tidy up one by one.

\textbf{Notes (JA)}:\\
「片付ける」は、物を整理したり、整頓したりすることを表現するために使用されます。しかし、「片付ける」は、物を整理するだけでなく、問題を解決したり、終わらせたりすることも意味します。

\textbf{Notes (EN)}:\\
The phrase "katadukeru" is used to express that you are tidying up or organizing things. However, "katazukeru" not only means organizing things but also means solving problems or finishing things.

\newpage

\section*{169. Ki wo tsukau}

\textbf{Date}: 2024.09.26 (Thu)\\
\textbf{Index}: kiwotsukau\\
\textbf{Expression (JA)}: (気:き)を(使:つか)う\\
\textbf{Romanization}: Ki wo tsukau\\
\textbf{Expression (EN)}: Be considerate of

\textbf{Variants (JA)}: \\
\textbf{Variants (EN)}: 

\textbf{Intent (JA)}: \\
\textbf{Intent (EN)}: 

\textbf{Tags (JA)}: \\
\textbf{Tags (EN)}: 

\textbf{Adjusted Expression (JA)}:\\
\\
\textbf{Adjusted Expression (EN)}:\\


\textbf{Example (JA)}:\\
A「(彼女:かのじょ)は、(他人:たにん)に(気:き)を(使:つか)う」\\
\textbf{Example (EN)}:\\
A: She is always considerate of others.

\textbf{Notes (JA)}:\\
「気をつかう」は、他人に気を遣うことを表現するために使用されます。

\textbf{Notes (EN)}:\\
The phrase "ki wo tsukau" is used to express that you are considerate of others. The word "ki" is a noun that means "spirit" or "mind." The word "tsukau" is a verb that means "to use."

\newpage

\section*{168. Show one's face}

\textbf{Date}: 2024.09.25 (Wed)\\
\textbf{Index}: kaowodasu\\
\textbf{Expression (JA)}: (顔:かお)を(出:だ)す\\
\textbf{Romanization}: Kao wo dasu\\
\textbf{Expression (EN)}: Show up

\textbf{Variants (JA)}: \\
\textbf{Variants (EN)}: 

\textbf{Intent (JA)}: \\
\textbf{Intent (EN)}: 

\textbf{Tags (JA)}: \\
\textbf{Tags (EN)}: 

\textbf{Adjusted Expression (JA)}:\\
\\
\textbf{Adjusted Expression (EN)}:\\


\textbf{Example (JA)}:\\
A「(彼:かれ)は、(最近:さいきん)、この(喫茶店:きっさてん)に(顔:かお)を(出:だ)していない」\\
\textbf{Example (EN)}:\\
A: He hasn't shown up at this cafe recently.

\textbf{Notes (JA)}:\\
「顔を出す」は、誰かがある場所に現れることを表現するために使用されます。

\textbf{Notes (EN)}:\\
The phrase "kao wo dasu" is used to express that someone appears or shows up at a place. The word "kao" is a noun that means "face." The word "dasu" is a verb that means "to put out" or "to show."

\newpage

\section*{167. Pull the eyes}

\textbf{Date}: 2024.09.24 (Tue)\\
\textbf{Index}: hitomewohiku\\
\textbf{Expression (JA)}: (人目:ひとめ)を(引:ひ)く\\
\textbf{Romanization}: Hitome wo hiku\\
\textbf{Expression (EN)}: Attract attention

\textbf{Variants (JA)}: \\
\textbf{Variants (EN)}: 

\textbf{Intent (JA)}: \\
\textbf{Intent (EN)}: 

\textbf{Tags (JA)}: \\
\textbf{Tags (EN)}: 

\textbf{Adjusted Expression (JA)}:\\
\\
\textbf{Adjusted Expression (EN)}:\\


\textbf{Example (JA)}:\\
A「この(ポスター:ぽすたー)は、(人目:ひとめ)を(引:ひ)きますね」\\
\textbf{Example (EN)}:\\
A: This poster attracts attention.

\textbf{Notes (JA)}:\\
「人目を引く」は、何かが注目を集めることを表現するために使用されます。

\textbf{Notes (EN)}:\\
The phrase "hitome wo hiku" is used to express that something attracts attention. The word "hitome" is a noun that means "eye" or "attention." The word "hiku" is a verb that means "to pull" or "to attract."

\newpage

\section*{166. Bake the hands}

\textbf{Date}: 2024.09.23 (Mon)\\
\textbf{Index}: tewoyaku\\
\textbf{Expression (JA)}: (手:て)を(焼:や)く\\
\textbf{Romanization}: Te wo yaku\\
\textbf{Expression (EN)}: Cannot take it anymore

\textbf{Variants (JA)}: \\
\textbf{Variants (EN)}: 

\textbf{Intent (JA)}: \\
\textbf{Intent (EN)}: 

\textbf{Tags (JA)}: \\
\textbf{Tags (EN)}: 

\textbf{Adjusted Expression (JA)}:\\
\\
\textbf{Adjusted Expression (EN)}:\\


\textbf{Example (JA)}:\\
A「(子供:こども)の(態度:たいど)に(手:て)を(焼:や)いている」\\
\textbf{Example (EN)}:\\
A: I can't take my child's attitude anymore.

\textbf{Notes (JA)}:\\
「手を焼く」は、誰かや何かと取り扱いに困難を感じていることを表現するために使用されます。

\textbf{Notes (EN)}:\\
The phrase "te wo yaku" is used to express that you are having a hard time dealing with someone or something. The word "te" is a noun that means "hand." The word "yaku" is a verb that means "to bake."

\newpage

\section*{165. Cut the hands}

\textbf{Date}: 2024.09.22 (Sun)\\
\textbf{Index}: tewokiru\\
\textbf{Expression (JA)}: (手:て)を(切:き)る\\
\textbf{Romanization}: Te wo kiru\\
\textbf{Expression (EN)}: Quit a relationship

\textbf{Variants (JA)}: \\
\textbf{Variants (EN)}: 

\textbf{Intent (JA)}: \\
\textbf{Intent (EN)}: 

\textbf{Tags (JA)}: \\
\textbf{Tags (EN)}: 

\textbf{Adjusted Expression (JA)}:\\
\\
\textbf{Adjusted Expression (EN)}:\\


\textbf{Example (JA)}:\\
A「この(会社:かいしゃ)とは(手:て)を(切:き)って、もう(取:と)り(引:ひ)きは(行:おこな)わない」\\
\textbf{Example (EN)}:\\
A: I'm going to quit this company and never deal with them again.

\textbf{Notes (JA)}:\\
「手を切る」は、誰かとの関係やつながりを絶つことを表現するために使用されます。

\textbf{Notes (EN)}:\\
The phrase "te wo kiru" is used to express that you are going to end a relationship or connection with someone or something. The word "te" is a noun that means "hand." The word "kiru" is a verb that means "to cut."

\newpage

\section*{164. A hard head}

\textbf{Date}: 2024.09.21 (Sat)\\
\textbf{Index}: atamagakatai\\
\textbf{Expression (JA)}: (頭:あたま)が(固:かた)い\\
\textbf{Romanization}: Atama ga katai\\
\textbf{Expression (EN)}: Be stubborn

\textbf{Variants (JA)}: \\
\textbf{Variants (EN)}: 

\textbf{Intent (JA)}: \\
\textbf{Intent (EN)}: 

\textbf{Tags (JA)}: \\
\textbf{Tags (EN)}: 

\textbf{Adjusted Expression (JA)}:\\
\\
\textbf{Adjusted Expression (EN)}:\\


\textbf{Example (JA)}:\\
A「(彼:かれ)は、(自分:じぶん)の(意見:いけん)を(曲:ま)げない」 B「(頭:あたま)が(固:かた)いね」\\
\textbf{Example (EN)}:\\
A: He doesn't change his opinion. B: He's stubborn.

\textbf{Notes (JA)}:\\
「頭が固い」は、誰かが頑固で柔軟性がないことを表現するために使用されます。

\textbf{Notes (EN)}:\\
The phrase "atama ga katai" is used to express that someone is stubborn or inflexible. The word "atama" is a noun that means "head." The word "katai" is an i-adjective that means "hard" or "stiff."

\newpage

\section*{163. Tears of sparrows}

\textbf{Date}: 2024.09.20 (Fri)\\
\textbf{Index}: suzumenonamida\\
\textbf{Expression (JA)}: (雀:すずめ)の(涙:なみだ)\\
\textbf{Romanization}: Suzume no namida\\
\textbf{Expression (EN)}: Tears of sparrows

\textbf{Variants (JA)}: \\
\textbf{Variants (EN)}: 

\textbf{Intent (JA)}: \\
\textbf{Intent (EN)}: 

\textbf{Tags (JA)}: \\
\textbf{Tags (EN)}: 

\textbf{Adjusted Expression (JA)}:\\
\\
\textbf{Adjusted Expression (EN)}:\\


\textbf{Example (JA)}:\\
A「ねぇ、(初任給:しょにんきゅう)、(出:で)た？」 B「うん、(雀:すずめ)の(涙:なみだ)ほどね」\\
\textbf{Example (EN)}:\\
A: Hey, did you get your first paycheck? B: Yeah, it's just a drop in the bucket.

\textbf{Notes (JA)}:\\
「雀の涙」は、何かが非常に小さく、取るに足らないことを表現するために使用されます。

\textbf{Notes (EN)}:\\
The phrase "suzume no namida" is used to express that something is very small or insignificant. The word "suzume" is a noun that means "sparrow." The word "namida" is a noun that means "tear."

\newpage

\section*{162. Concerns}

\textbf{Date}: 2024.09.19 (Thu)\\
\textbf{Index}: kininaru\\
\textbf{Expression (JA)}: (気:き)になる\\
\textbf{Romanization}: Ki ni naru\\
\textbf{Expression (EN)}: Be concerned

\textbf{Variants (JA)}: \\
\textbf{Variants (EN)}: 

\textbf{Intent (JA)}: \\
\textbf{Intent (EN)}: 

\textbf{Tags (JA)}: \\
\textbf{Tags (EN)}: 

\textbf{Adjusted Expression (JA)}:\\
\\
\textbf{Adjusted Expression (EN)}:\\


\textbf{Example (JA)}:\\
A「その(ドラマ:どらま)の(続:つづ)きがどうなるのか、(気:き)になって(眠:ねむ)れない」\\
\textbf{Example (EN)}:\\
A: I can't sleep because I'm concerned about how the drama will end.

\textbf{Notes (JA)}:\\
「気になる」は、何かに興味を持ったり、心配したりすることを表現するために使用されます。

\textbf{Notes (EN)}:\\
The phrase "ki ni naru" is used to express that you are concerned or interested in something. The word "ki" is a noun that means "mind" or "spirit." The word "ni" is a particle that is used to indicate the target of an action or the direction of a movement. The word "naru" is a verb that means "to become."

\newpage

\section*{161. I am into it}

\textbf{Date}: 2024.09.18 (Wed)\\
\textbf{Index}: hamaru\\
\textbf{Expression (JA)}: (ハマ:はま)る\\
\textbf{Romanization}: Hamaru\\
\textbf{Expression (EN)}: I am into it.

\textbf{Variants (JA)}: \\
\textbf{Variants (EN)}: 

\textbf{Intent (JA)}: \\
\textbf{Intent (EN)}: 

\textbf{Tags (JA)}: \\
\textbf{Tags (EN)}: 

\textbf{Adjusted Expression (JA)}:\\
\\
\textbf{Adjusted Expression (EN)}:\\


\textbf{Example (JA)}:\\
A「(最近:さいきん)、(読書:どくしょ)にはまっている」\\
\textbf{Example (EN)}:\\
A: I've been into reading books lately.

\textbf{Notes (JA)}:\\
「ハマる」は、何かに興味を持っていて、それに多くの時間を費やしていることを表現するために使用されます。

\textbf{Notes (EN)}:\\
The phrase "hamaru" is used to express that you are interested in something and have been spending a lot of time on it. The word "hamaru" is a verb that means "to be into" or "to be interested in."

\newpage

\section*{160. Can't wait}

\textbf{Date}: 2024.09.17 (Tue)\\
\textbf{Index}: machikirenai\\
\textbf{Expression (JA)}: (待:ま)ちきれない\\
\textbf{Romanization}: Machikirenai\\
\textbf{Expression (EN)}: Can't wait

\textbf{Variants (JA)}: \\
\textbf{Variants (EN)}: 

\textbf{Intent (JA)}: \\
\textbf{Intent (EN)}: 

\textbf{Tags (JA)}: \\
\textbf{Tags (EN)}: 

\textbf{Adjusted Expression (JA)}:\\
\\
\textbf{Adjusted Expression (EN)}:\\


\textbf{Example (JA)}:\\
A「(明日:あした)の(旅行:りょこう)、(楽:たの)しみだね」 B「うん、(待:ま)ちきれない」\\
\textbf{Example (EN)}:\\
A: I'm looking forward to the trip tomorrow. B: Yeah, I can't wait.

\textbf{Notes (JA)}:\\
「待ちきれない」は、何かを楽しみにしていて、それが起こるのを待ちきれないことを表現するために使用されます。「待ち」は動詞の「待つ」のマス型に「きる(我慢する、できる)」の否定形「きれない」がついたものです。

\textbf{Notes (EN)}:\\
The phrase "machikirenai" is used to express that you are looking forward to something and cannot wait for it to happen. The word "machi" is a masu-form of a verb "matsu (to wait)" such as "machi-masu." The word "kirenai" is a negative form of the verb "kiru" that means "to be able to endure."

\newpage

\section*{159. Various kinds of services}

\textbf{Date}: 2024.09.16 (Mon)\\
\textbf{Index}: deribari-sa-bisu\\
\textbf{Expression (JA)}: (デリバリーサービス:でりばりーさーびす)\\
\textbf{Romanization}: Deribari- sa-bisu\\
\textbf{Expression (EN)}: Delivery service

\textbf{Variants (JA)}: \\
\textbf{Variants (EN)}: 

\textbf{Intent (JA)}: \\
\textbf{Intent (EN)}: 

\textbf{Tags (JA)}: \\
\textbf{Tags (EN)}: 

\textbf{Adjusted Expression (JA)}:\\
\\
\textbf{Adjusted Expression (EN)}:\\


\textbf{Example (JA)}:\\
A「(デリバリーサービス:でりばりーさーびす)を(利用:りよう)すると、(家:いえ)に(居:い)ながら(食:た)べ(物:もの)を(注文:ちゅうもん)できる」\\
\textbf{Example (EN)}:\\
A: If you use a delivery service, you can order food while staying at home.

\textbf{Notes (JA)}:\\
「デリバリーサービス」は、食べ物や商品を家に配達するサービスを表現するために使用されます。

\textbf{Notes (EN)}:\\
The phrase "deribarii saabisu" is used to express a service that delivers food or goods to your home. The word "deribarii" is a loanword from English that means "delivery." The word "saabisu" is a loanword from English that means "service."

\newpage

\section*{158. Nostalgic}

\textbf{Date}: 2024.09.15 (Sun)\\
\textbf{Index}: natsukashii\\
\textbf{Expression (JA)}: (懐:なつ)かしい\\
\textbf{Romanization}: Natsukashii\\
\textbf{Expression (EN)}: Nostalgic

\textbf{Variants (JA)}: \\
\textbf{Variants (EN)}: 

\textbf{Intent (JA)}: \\
\textbf{Intent (EN)}: 

\textbf{Tags (JA)}: \\
\textbf{Tags (EN)}: 

\textbf{Adjusted Expression (JA)}:\\
\\
\textbf{Adjusted Expression (EN)}:\\


\textbf{Example (JA)}:\\
A「この(曲:きょく)を(聞:き)くと、(昔:むかし)のことを(思:おも)い(出:だ)す」 B「(懐:なつ)かしいね」\\
\textbf{Example (EN)}:\\


\textbf{Notes (JA)}:\\
「懐かしい」は、何かが過去を思い出させ、懐かしく感じさせることを表現するために使用されます。

\textbf{Notes (EN)}:\\
The phrase "natsu kashii" is used to express that something reminds you of the past and makes you feel nostalgic.

\newpage

\section*{157. Atomosphere}

\textbf{Date}: 2024.09.14 (Sat)\\
\textbf{Index}: funniki\\
\textbf{Expression (JA)}: (雰囲気:ふんいき)\\
\textbf{Romanization}: Fun'iki\\
\textbf{Expression (EN)}: Atomosphere

\textbf{Variants (JA)}: \\
\textbf{Variants (EN)}: 

\textbf{Intent (JA)}: \\
\textbf{Intent (EN)}: 

\textbf{Tags (JA)}: \\
\textbf{Tags (EN)}: 

\textbf{Adjusted Expression (JA)}:\\
\\
\textbf{Adjusted Expression (EN)}:\\


\textbf{Example (JA)}:\\
A「この(店:みせ)の(雰囲気:ふんいき)が(好:す)き」\\
\textbf{Example (EN)}:\\
A: I like the atmosphere of this store.

\textbf{Notes (JA)}:\\
「雰囲気」は、場所や状況の全体的な感じや雰囲気を表す名詞です。具体的にいい言い表せない時、このことばを使うと便利です。他にも「なんとなく」という言い方もあります。

\textbf{Notes (EN)}:\\
The word "fun'iki" is a noun that means the overall feel or atmosphere of a place or situation. When you cannot quite put your finger on what makes a place feel a certain way, using the word "fun'iki" can be helpful. Another way to express this idea is to say "nantonaku," which means "somehow" or "in some way."

\newpage

\section*{156. More/better than anything}

\textbf{Date}: 2024.09.13 (Fri)\\
\textbf{Index}: naniyori\\
\textbf{Expression (JA)}: なにより\\
\textbf{Romanization}: Nani yori\\
\textbf{Expression (EN)}: More than anything

\textbf{Variants (JA)}: \\
\textbf{Variants (EN)}: 

\textbf{Intent (JA)}: \\
\textbf{Intent (EN)}: 

\textbf{Tags (JA)}: \\
\textbf{Tags (EN)}: 

\textbf{Adjusted Expression (JA)}:\\
\\
\textbf{Adjusted Expression (EN)}:\\


\textbf{Example (JA)}:\\
この(プロジェクト:ぷろじぇくと)が(成功:せいこう)したことがなによりの(成果:せいか)です。\\
\textbf{Example (EN)}:\\
The success of this project is more than anything.

\textbf{Notes (JA)}:\\
「何より」は、何かが他の何よりも重要であることを表現するために使用されます。

\textbf{Notes (EN)}:\\
The phrase "nani yori" is used to express that something is more important than anything else. The word "nani" is a pronoun that means "what." The word "yori" is a particle that is used to indicate comparison or preference.

\newpage

\section*{155. An egg}

\textbf{Date}: 2024.09.12 (Thu)\\
\textbf{Index}: tamago\\
\textbf{Expression (JA)}: (卵:たまご)\\
\textbf{Romanization}: Tamago\\
\textbf{Expression (EN)}: An egg/A talented or promising individual

\textbf{Variants (JA)}: \\
\textbf{Variants (EN)}: 

\textbf{Intent (JA)}: \\
\textbf{Intent (EN)}: 

\textbf{Tags (JA)}: \\
\textbf{Tags (EN)}: 

\textbf{Adjusted Expression (JA)}:\\
\\
\textbf{Adjusted Expression (EN)}:\\


\textbf{Example (JA)}:\\
(君:きみ)はまだ(プログラマ:ぷろぐらま)の(卵:たまご)だが、(素質:そしつ)は(十分:じゅうぶん)ある。\\
\textbf{Example (EN)}:\\
You're still an aspiring programmer, but you've got a lot of potential.

\textbf{Notes (JA)}:\\
「卵」は、卵という意味の名詞です。この表現では、「卵」ということばは比喩的に使われ、特定の分野でまだ経験が浅いか、初心者であるが、成長し、スキルを磨いていく可能性がある人を指します。「...の卵」というフレーズは、成功する可能性がある人を表現するために使われます。

\textbf{Notes (EN)}:\\
The word "tamago" is a noun that means "egg." In this expression, the word "tamago" is used metaphorically to refer to someone who is still inexperienced or a beginner in a particular field, but has the potential to grow and develop their skills. The phrase "... no tamago" is often used to describe someone who is new to a profession or activity, but has the potential to become skilled or successful in the future.

\newpage

\section*{154. No eyes?}

\textbf{Date}: 2024.09.11 (Wed)\\
\textbf{Index}: meganai\\
\textbf{Expression (JA)}: (目:め)がない\\
\textbf{Romanization}: Me ga nai\\
\textbf{Expression (EN)}: Be blind

\textbf{Variants (JA)}: \\
\textbf{Variants (EN)}: 

\textbf{Intent (JA)}: \\
\textbf{Intent (EN)}: 

\textbf{Tags (JA)}: \\
\textbf{Tags (EN)}: 

\textbf{Adjusted Expression (JA)}:\\
\\
\textbf{Adjusted Expression (EN)}:\\


\textbf{Example (JA)}:\\
(彼:かれ)は、(甘:あま)いものに(目:め)がない。\\
\textbf{Example (EN)}:\\
He has a sweet tooth.

\textbf{Notes (JA)}:\\
「目がない」は、誰かが何かを非常に好きであることを意味します。

\textbf{Notes (EN)}:\\
The phrase "me ga nai" means that someone likes something very much. The word "me" is a noun that means "eye." The word "nai" is an i-adjective that means "not have."

\newpage

\section*{153. Big eyes}

\textbf{Date}: 2024.09.10 (Tue)\\
\textbf{Index}: megatobidiru\\
\textbf{Expression (JA)}: (目:め)が(飛:と)び(出:で)る\\
\textbf{Romanization}: Me ga tobidiru\\
\textbf{Expression (EN)}: Eyes pop out

\textbf{Variants (JA)}: \\
\textbf{Variants (EN)}: 

\textbf{Intent (JA)}: \\
\textbf{Intent (EN)}: 

\textbf{Tags (JA)}: \\
\textbf{Tags (EN)}: 

\textbf{Adjusted Expression (JA)}:\\
\\
\textbf{Adjusted Expression (EN)}:\\


\textbf{Example (JA)}:\\
A「その(値段:ねだん)を(見:み)て、(目:め)が(飛:と)び(出:で)るほど(驚:おどろ)いた」\\
\textbf{Example (EN)}:\\
A: Seeing the price, I was so surprised that my eyes nearly popped out.

\textbf{Notes (JA)}:\\
この日本語の表現「その値段を見て、目が飛び出るほど驚いた」は、その価格が非常に驚くべきものであったため、極端な反応を引き起こしたことを意味します。「目が飛び出るほど」というフレーズは、文字通り「目が飛び出る」という意味で、誰かが非常に驚いたり、ショックを受けたりする様子を描写する生々しい表現です。この表現は、価格を見たときの反応が非常に強烈で誇張されていることを伝えています。

\textbf{Notes (EN)}:\\
This Japanese expression means that the price was so astonishingly high that it caused an extreme reaction. The phrase "目が飛び出るほど" (me ga tobi deru hodo) literally translates to "eyes popping out," which is a vivid way of describing someone being extremely surprised or shocked. The expression conveys that the reaction to seeing the price was intense and exaggerated, similar to how one might react if they were extremely astonished.

\newpage

\section*{152. A black stomach guy}

\textbf{Date}: 2024.09.09 (Mon)\\
\textbf{Index}: haragakuroi\\
\textbf{Expression (JA)}: (腹:はら)が(黒:くろ)い\\
\textbf{Romanization}: Hara ga kuroi\\
\textbf{Expression (EN)}: Stomach is black

\textbf{Variants (JA)}: \\
\textbf{Variants (EN)}: 

\textbf{Intent (JA)}: \\
\textbf{Intent (EN)}: 

\textbf{Tags (JA)}: \\
\textbf{Tags (EN)}: 

\textbf{Adjusted Expression (JA)}:\\
\\
\textbf{Adjusted Expression (EN)}:\\


\textbf{Example (JA)}:\\
A「(彼:かれ)は、(腹:はら)が(黒:くろ)い」\\
\textbf{Example (EN)}:\\
A: He is malicious.

\textbf{Notes (JA)}:\\
「腹が黒い」は、誰かが悪意を持っていることを表現するために使用されます。本当に内蔵が黒いわけではありません。

\textbf{Notes (EN)}:\\
The phrase "hara ga kuroi" is used to express that someone has malice. It does not mean that the person's internal organs are actually black. The word "hara" is a noun that means "stomach." The word "kuroi" is an i-adjective that means "black."

\newpage

\section*{151. Break a bone}

\textbf{Date}: 2024.09.08 (Sun)\\
\textbf{Index}: honewooru\\
\textbf{Expression (JA)}: (骨:ほね)を(折:お)る\\
\textbf{Romanization}: Hone wo oru\\
\textbf{Expression (EN)}: Break a bone

\textbf{Variants (JA)}: \\
\textbf{Variants (EN)}: 

\textbf{Intent (JA)}: \\
\textbf{Intent (EN)}: 

\textbf{Tags (JA)}: \\
\textbf{Tags (EN)}: 

\textbf{Adjusted Expression (JA)}:\\
\\
\textbf{Adjusted Expression (EN)}:\\


\textbf{Example (JA)}:\\
A「(先生:せんせい)は、(特許取得:とっきょしゅとく)のために、いろいろ(骨:ほね)を(折:お)ってくださいました」\\
\textbf{Example (EN)}:\\
A: The teacher went to great lengths to obtain a patent.

\textbf{Notes (JA)}:\\
「努力をつける」は、何かに多くの努力を払ったことを表現するために使用されます。本当に体の骨を折ったわけではありません。

\textbf{Notes (EN)}:\\
The phrase "hone wo oru" is used to express that you put a lot of effort into something. It does not mean that you actually broke a bone. The word "hone" is a noun that means "bone." The word "oru" is a verb that means "to break."

\newpage

\section*{150. I don't wanna hear that.}

\textbf{Date}: 2024.09.07 (Sat)\\
\textbf{Index}: mimigaitai\\
\textbf{Expression (JA)}: (耳:みみ)が(痛:いた)い\\
\textbf{Romanization}: Mimi ga itai\\
\textbf{Expression (EN)}: I don't wanna hear that.

\textbf{Variants (JA)}: 耳が痛い / 耳が痛い話 / 耳の中が痛い\\
\textbf{Variants (EN)}: I don't wanna hear that / that hits too close to home / my ear hurts

\textbf{Intent (JA)}: 自分にとって都合の悪い指摘への反応, 軽い自己認識の表現, 気まずさの共有\\
\textbf{Intent (EN)}: reacting to an inconvenient truth, expressing self-awareness lightly, sharing discomfort

\textbf{Tags (JA)}: 即時文法, 慣用表現, 口語, 自己認識, 気まずさ\\
\textbf{Tags (EN)}: immediate grammar, idiomatic expression, colloquial, self-awareness, discomfort

\textbf{Adjusted Expression (JA)}:\\
(それ)は、(自分:じぶん)に(とって)都合の(悪:わる)い(指摘:してき)だ。\\
\textbf{Adjusted Expression (EN)}:\\
That is a point that is inconvenient for me.

\textbf{Example (JA)}:\\
A「それは、(耳:みみ)が(痛:いた)い(話:はなし)だ」\\
\textbf{Example (EN)}:\\
A: That's a painful story.

\textbf{Notes (JA)}:\\
「耳が痛い」は、耳が痛いことを表現するために使用されることもないわけではありませんが、その時には、「耳の中が痛い」と言います。一般的には、「耳が痛い」という表現は、他人の話を聞いて、その内容が自分にとって不都合であると感じたときに使われます。

\textbf{Notes (EN)}:\\
There are some cases where the phrase "mimi ga itai" is used to express that your ear hurts, but in that case, you say "mimi no naka ga itai (my ear hurts)." Generally, the expression "mimi ga itai" is used when you listen to someone's story and feel uncomfortable because the content is inconvenient for you.

\newpage

\section*{149. Have an eye}

\textbf{Date}: 2024.09.06 (Fri)\\
\textbf{Index}: mewomotteiru\\
\textbf{Expression (JA)}: (目:め)を(持:も)っている\\
\textbf{Romanization}: Me wo motteiru\\
\textbf{Expression (EN)}: Have an eye

\textbf{Variants (JA)}: 目を持っている / 見る目がある / 目が利く\\
\textbf{Variants (EN)}: have an eye / have an eye for / have a good eye

\textbf{Intent (JA)}: 能力や感性の評価, 鑑識眼の表現, 人物評価\\
\textbf{Intent (EN)}: evaluating ability or sensitivity, expressing discernment, evaluating a person

\textbf{Tags (JA)}: 即時文法, 慣用表現, 口語, 能力評価, 感性\\
\textbf{Tags (EN)}: immediate grammar, idiomatic expression, colloquial, ability evaluation, sensitivity

\textbf{Adjusted Expression (JA)}:\\
(彼:かれ)は、(絵画:かいが)を(適切:てきせつ)に(評価:ひょうか)する(能力:のうりょく)を(持:も)っている。\\
\textbf{Adjusted Expression (EN)}:\\
He has the ability to evaluate paintings appropriately.

\textbf{Example (JA)}:\\
A「(彼:かれ)は、(絵:え)を(見:み)る(目:め)を(持:も)っている」\\
\textbf{Example (EN)}:\\
A: He has an eye for pictures.

\textbf{Notes (JA)}:\\
「目を持っている」は、絵画の美しさや価値を理解し、評価する能力や感性を持っていることを意味します。それは、物理的な見る能力だけでなく、芸術を鑑賞するために必要な美的感覚と理解を持っていることを意味します。「目」は目という意味の名詞です。「持っている」は、「持つ」という動詞のて形です。

\textbf{Notes (EN)}:\\
The phrase "me wo motteiru (he has the eyes to appreciate paintings)" means that he possesses the ability or sensitivity to understand and evaluate the beauty and value of paintings. It implies that he not only has the physical ability to see but also has the aesthetic sense and understanding required to appreciate art. The word "me" is a noun that means "eye." The word "motteiru" is the te-form of the verb "motsu" that means "to have."

\newpage

\section*{148. Get fired}

\textbf{Date}: 2024.09.05 (Thu)\\
\textbf{Index}: kubininaru\\
\textbf{Expression (JA)}: (首:くび)になる\\
\textbf{Romanization}: Kubi ni naru\\
\textbf{Expression (EN)}: Get fired

\textbf{Variants (JA)}: 首になる / クビになる / 首にされる\\
\textbf{Variants (EN)}: get fired / be fired / be let go

\textbf{Intent (JA)}: 解雇の表現, 結果の報告, 失敗や不注意の帰結\\
\textbf{Intent (EN)}: expressing being fired, reporting a result, consequence of failure or carelessness

\textbf{Tags (JA)}: 即時文法, 慣用表現, 口語, 仕事場面, 解雇\\
\textbf{Tags (EN)}: immediate grammar, idiomatic expression, colloquial, work situation, dismissal

\textbf{Adjusted Expression (JA)}:\\
(彼:かれ)は、(会社:かいしゃ)から(解雇:かいこ)された。\\
\textbf{Adjusted Expression (EN)}:\\
He was dismissed from the company.

\textbf{Example (JA)}:\\
A「(何回:なんかい)も(遅刻:ちこく)して、(アルバイト:あるばいと)を(首:くび)になった」\\
\textbf{Example (EN)}:\\
A: I was fired from my part-time job because I was late many times.

\textbf{Notes (JA)}:\\
「首になる」は、仕事をクビにされたことを表現するために使用されます。「首」は首という意味の名詞です。「なる」は「become」という意味の動詞です。

\textbf{Notes (EN)}:\\
The phrase "kubi ni naru" is used to express that you are fired from your job. The word "kubi" is a noun that means "neck." The word "naru" is a verb that means "to become."

\newpage

\section*{147. Secret weapon}

\textbf{Date}: 2024.09.04 (Wed)\\
\textbf{Index}: okunote\\
\textbf{Expression (JA)}: 奥の手\\
\textbf{Romanization}: Oku no te\\
\textbf{Expression (EN)}: Secret weapon

\textbf{Variants (JA)}: 奥の手 / 奥の手を使う / 最後の切り札\\
\textbf{Variants (EN)}: secret weapon / use a secret weapon / last resort

\textbf{Intent (JA)}: 隠していた手段の提示, 切り札の使用の予告, 状況打開の表現\\
\textbf{Intent (EN)}: presenting a hidden method, announcing the use of a trump card, expressing a way to turn the situation around

\textbf{Tags (JA)}: 即時文法, 慣用表現, 名詞句, 戦略, 切り札\\
\textbf{Tags (EN)}: immediate grammar, idiomatic expression, noun phrase, strategy, trump card

\textbf{Adjusted Expression (JA)}:\\
(彼:かれ)は、(秘策:ひさく)を(用:もち)いて、(問題:もんだい)を(解決:かいけつ)した。\\
\textbf{Adjusted Expression (EN)}:\\
He solved the problem by using a hidden strategy.

\textbf{Example (JA)}:\\
A「(奥:おく)の(手:て)を(使:つか)うから、まだ(大丈夫:だいじょうぶ)」\\
\textbf{Example (EN)}:\\
A: I'll use my secret weapon, so I'm still okay.

\textbf{Notes (JA)}:\\
「奥の手」は、何かを隠しておいて、必要な時に使うことを表現するために使用されます。「奥」とは奥深いところ、秘密の意味があります。「手」は手のひら、手の意味以外に、手段、方法があります。

\textbf{Notes (EN)}:\\
The phrase "oku no te" is used to express that you have something hidden and use it when needed. The word "oku" means "deep" or "secret." The word "te" means "hand," and also extends to the meaning of "means" or "method."

\newpage

\section*{146. Goof off}

\textbf{Date}: 2024.09.03 (Tue)\\
\textbf{Index}: aburawouru\\
\textbf{Expression (JA)}: (油:あぶら)を(売:う)る\\
\textbf{Romanization}: Abura wo uru\\
\textbf{Expression (EN)}: Goof off

\textbf{Variants (JA)}: 油を売る / 油を売っている / 油を売っていました\\
\textbf{Variants (EN)}: goof off / be goofing off / was goofing off

\textbf{Intent (JA)}: サボっていることの表現, 仕事をしていないことの指摘, 軽い批判や冗談\\
\textbf{Intent (EN)}: expressing that someone is slacking off, pointing out that someone is not working, light criticism or joking

\textbf{Tags (JA)}: 即時文法, 慣用表現, 口語, サボる, 仕事場面, 軽い批判\\
\textbf{Tags (EN)}: immediate grammar, idiomatic expression, colloquial, slacking off, work situation, light criticism

\textbf{Adjusted Expression (JA)}:\\
(彼:かれ)は、(私語:しご)にふけり、(業務:ぎょうむ)を(怠:おこた)っています。\\
\textbf{Adjusted Expression (EN)}:\\
He is absorbed in private chatting and neglecting his duties.

\textbf{Example (JA)}:\\
A「(彼:かれ)は、(女:おんな)の(子:こ)とおしゃべりばかりして、(油:あぶら)を(売:う)っていました」\\
\textbf{Example (EN)}:\\
A: He was just chatting with girls and goofing off.

\textbf{Notes (JA)}:\\
「油を売る」は、本当に油を売るのではなく、仕事をしないでサボることです。もし、彼の職業が本当に油売りだったら、その時は仕事をしているのか、サボっているのか、わかりませんので、注意が必要です。

\textbf{Notes (EN)}:\\
The phrase "abura wo uru" means not to sell oil literally, but to goof off without working. The word "abura" is a noun that means "oil." The word "uru" is a verb that means "to sell." If his job was really selling oil, it would be difficult to tell whether he was working or goofing off.

\newpage

\section*{145. I could use all the help I can get}

\textbf{Date}: 2024.09.02 (Mon)\\
\textbf{Index}: nekonotemokaritai\\
\textbf{Expression (JA)}: (猫:ねこ)の(手:て)も(借:か)りたい\\
\textbf{Romanization}: Neko no te mo karitai\\
\textbf{Expression (EN)}: I could use all the help I can get

\textbf{Variants (JA)}: 猫の手も借りたい / 猫の手も借りたいほど / 猫の手も借りたいくらい\\
\textbf{Variants (EN)}: I could use all the help I can get / I'm so busy I could even borrow a cat's paws / I need all the help I can get

\textbf{Intent (JA)}: 忙しさの誇張表現, 人手不足の表現, 助けが欲しい気持ちの共有\\
\textbf{Intent (EN)}: exaggerating how busy one is, expressing a lack of helping hands, sharing the feeling of wanting help

\textbf{Tags (JA)}: 即時文法, 慣用表現, 口語, 忙しさ, 誇張表現, 人手不足\\
\textbf{Tags (EN)}: immediate grammar, idiomatic expression, colloquial, busyness, hyperbolic expression, lack of help

\textbf{Adjusted Expression (JA)}:\\
(忙:いそが)しくて、(猫:ねこ)の(手:て)も(借:か)りたいほどです。\\
\textbf{Adjusted Expression (EN)}:\\
I'm so busy that I could even borrow a cat's paws.

\textbf{Example (JA)}:\\
A「(何:なに)か(細:こま)かい(作業:さぎょう)を(手伝:てつだ)ってくれる(人:ひと)が(欲:ほ)しい」 B「(猫:ねこ)の(手:て)も(借:か)りたいね」\\
\textbf{Example (EN)}:\\
A: I want someone to help me with some small tasks. B: I could use all the help I can get.

\textbf{Notes (JA)}:\\
「猫の手も借りたい」は、とても忙しくて、少しでも助けが欲しいことを表現する言い方です。ふつうは役に立ちそうにない猫の手でも借りたい、という誇張を含んだ表現です。

\textbf{Notes (EN)}:\\
The phrase "neko no te mo karitai" is used to express that you are extremely busy and want as much help as possible. It is a hyperbolic expression: even the paws of a cat, which would not normally be very helpful, would be welcome.

\newpage

\section*{144. A guardian dog}

\textbf{Date}: 2024.09.01 (Sun)\\
\textbf{Index}: komainu\\
\textbf{Expression (JA)}: (狛犬:こまいぬ)\\
\textbf{Romanization}: Komainu\\
\textbf{Expression (EN)}: A guardian dog

\textbf{Variants (JA)}: 狛犬 / 狛犬が置かれている / 狛犬を見る\\
\textbf{Variants (EN)}: a guardian dog / a guardian dog is placed there / see a guardian dog

\textbf{Intent (JA)}: 物の説明, 神社の様子の描写, 文化的事物の言及\\
\textbf{Intent (EN)}: describing an object, describing a shrine scene, mentioning a cultural object

\textbf{Tags (JA)}: 即時文法, 名詞, 文化語彙, 神社, 置物\\
\textbf{Tags (EN)}: immediate grammar, noun, cultural vocabulary, shrine, ornament

\textbf{Adjusted Expression (JA)}:\\
(神社:じんじゃ)の(入口:いりぐち)に(狛犬:こまいぬ)が(置:お)かれています。\\
\textbf{Adjusted Expression (EN)}:\\
A guardian dog is placed at the entrance of the shrine.

\textbf{Example (JA)}:\\
A「(神社:じんじゃ)の(入口:いりぐち)に(狛犬:こまいぬ)が(置:お)かれている」\\
\textbf{Example (EN)}:\\
A: A guardian dog is placed at the entrance of the shrine.

\textbf{Notes (JA)}:\\
「狛犬」は、神社の入り口に置かれている置物で、その神社を守るために置かれています。

\textbf{Notes (EN)}:\\
The word "komainu" refers to a guardian dog placed at the entrance of a shrine. The word "komainu" is treated here as a noun referring to that statue or figure.

\newpage

\section*{143. A beckoning cat}

\textbf{Date}: 2024.08.31 (Sat)\\
\textbf{Index}: manekineko\\
\textbf{Expression (JA)}: (招:まね)き(猫:ねこ)\\
\textbf{Romanization}: Manekineko\\
\textbf{Expression (EN)}: A beckoning cat

\textbf{Variants (JA)}: 招き猫 / 招き猫が置かれている / 招き猫を置く\\
\textbf{Variants (EN)}: a beckoning cat / a beckoning cat is placed there / place a beckoning cat

\textbf{Intent (JA)}: 物の説明, 縁起物の言及, 店頭の様子の描写\\
\textbf{Intent (EN)}: describing an object, mentioning a lucky charm, describing a storefront scene

\textbf{Tags (JA)}: 即時文法, 名詞, 複合名詞, 文化語彙, 縁起物, 店頭表現\\
\textbf{Tags (EN)}: immediate grammar, noun, compound noun, cultural vocabulary, lucky charm, storefront expression

\textbf{Adjusted Expression (JA)}:\\
(店:みせ)の(入口:いりぐち)には、(商売繁盛:しょうばいはんじょう)を(願:ねが)って(招:まね)き(猫:ねこ)が(置:お)かれています。\\
\textbf{Adjusted Expression (EN)}:\\
A beckoning cat is placed at the entrance of the store to wish for prosperity in business.

\textbf{Example (JA)}:\\
A「(店:みせ)の(入口:いりぐち)には、(商売繁盛:しょうばいはんじょう)を(願:ねが)って(招:まね)き(猫:ねこ)が(置:お)かれている」\\
\textbf{Example (EN)}:\\
A: A beckoning cat is placed at the entrance of the store to wish for prosperity in business.

\textbf{Notes (JA)}:\\
「招き猫」は、商売繁盛を願って店の入り口に置かれる縁起物を表現するために使用されます。

\textbf{Notes (EN)}:\\
The phrase "manekineko" refers to a beckoning cat placed at the entrance of a store to wish for prosperity in business. The word "manekineko" is a compound noun composed of the verb "maneku" ("to beckon") and the noun "neko" ("cat").

\newpage

\section*{142. A coincidence}

\textbf{Date}: 2024.08.30 (Fri)\\
\textbf{Index}: guuzennoicchi\\
\textbf{Expression (JA)}: (偶然:ぐうぜん)の(一致:いっち)\\
\textbf{Romanization}: Guuzen no icchi\\
\textbf{Expression (EN)}: A coincidence

\textbf{Variants (JA)}: 偶然の一致 / 単なる偶然 / たまたま\\
\textbf{Variants (EN)}: a coincidence / just a coincidence / by chance

\textbf{Intent (JA)}: 偶然であることの説明, 一致の理由の否定, 状況の軽い説明\\
\textbf{Intent (EN)}: explaining something as coincidence, denying intentional matching, giving a light explanation

\textbf{Tags (JA)}: 即時文法, 名詞句, 説明表現, 偶然, 口語\\
\textbf{Tags (EN)}: immediate grammar, noun phrase, explanatory expression, coincidence, colloquial

\textbf{Adjusted Expression (JA)}:\\
単なる(偶然:ぐうぜん)の(一致:いっち)ですよ。\\
\textbf{Adjusted Expression (EN)}:\\
It's just a coincidence

\textbf{Example (JA)}:\\
A「(何:なん)で(同:おな)じ(服:ふく)(着:て)るの」 B「単なる(偶然:ぐうぜん)の(一致:いっち)」\\
\textbf{Example (EN)}:\\
A: Why are you wearing the same clothes? B: It's just a coincidence.

\textbf{Notes (JA)}:\\
「偶然の一致」は、何かが単なる偶然であることを表現するために使用されます。

\textbf{Notes (EN)}:\\
The phrase "guuzen no icchi" is used to express that something is a mere coincidence. The word "guuzen" is a noun that means "chance" or "accident." The word "icchi" is a noun that means "coincidence" or "agreement."

\newpage

\section*{141. Please help yourself.}

\textbf{Date}: 2024.08.29 (Thu)\\
\textbf{Index}: gojiyuuniotorikudasai\\
\textbf{Expression (JA)}: ご(自由:じゆう)に\\
\textbf{Romanization}: Go jiyuu ni\\
\textbf{Expression (EN)}: Please help yourself.

\textbf{Variants (JA)}: ご自由に / ご自由にどうぞ / ご自由にお取りください\\
\textbf{Variants (EN)}: please help yourself / please feel free / please take one freely

\textbf{Intent (JA)}: 許可の丁寧な提示, 相手への配慮を示す表現, 案内表現\\
\textbf{Intent (EN)}: politely giving permission, showing consideration to the other person, guiding expression

\textbf{Tags (JA)}: 即時文法, 定型表現, 待遇表現, 丁寧表現, 案内表現, 許可表現, サービス場面\\
\textbf{Tags (EN)}: immediate grammar, formulaic expression, honorific expression, polite expression, guiding expression, permission expression, service situation

\textbf{Adjusted Expression (JA)}:\\
ご(自由:じゆう)にお(取:と)りください\\
\textbf{Adjusted Expression (EN)}:\\
Please take one freely

\textbf{Example (JA)}:\\
A「(パンフレット:ぱんふれっと)、ひとつ、よろしいですか？」 B「ええ、どうぞ、ご(自由:じゆう)にお(取:と)りください」\\
\textbf{Example (EN)}:\\
A: Can I have a pamphlet? B: Yes, please help yourself.

\textbf{Notes (JA)}:\\
「自由」は、「freedom」や「liberty」という意味です。「ご...に」の構造は、「自由」をより丁寧にします。「お...ください」は、「取り」を丁寧にするために使われます。「取り」は「取る」から来ています。「取る」は「取ります」に変更され、その後「ます」が取り除かれて「取り」が形成されます。「ご自由にお取りください」は、何かを自由に取ることができるか、自分の判断で取ることができるかを丁寧に伝える方法です。

\textbf{Notes (EN)}:\\
"Jiyuu" means "freedom" or "liberty." The construction "go...ni" makes "jiyuu" more polite. Similarly, "o...kudasai" adds politeness to the verb "tori," which comes from "toru" ("to take"). "Toru" is first changed to "torimasu," and then "masu" is removed to form "tori." The phrase "go jiyuu ni otori kudasai" is a polite way to say that you can take something freely or at your own discretion.

\newpage

\section*{140. Let's keep going}

\textbf{Date}: 2024.08.28 (Wed)\\
\textbf{Index}: dondonsusumeyou\\
\textbf{Expression (JA)}: どんどん(進:すす)めよう\\
\textbf{Romanization}: Dondon susumeyou\\
\textbf{Expression (EN)}: Let's keep going

\textbf{Variants (JA)}: どんどん進めよう / どんどん進めていこう / さあ、どんどん進めよう\\
\textbf{Variants (EN)}: let's keep going / let's keep moving forward / come on, let's keep going

\textbf{Intent (JA)}: 前進の呼びかけ, 作業継続の促し, 勢いをつける表現\\
\textbf{Intent (EN)}: encouraging progress, urging continuation of work, adding momentum

\textbf{Tags (JA)}: 即時文法, 口語, 勧誘, 進行, 勢い\\
\textbf{Tags (EN)}: immediate grammar, colloquial, invitation, progress, momentum

\textbf{Adjusted Expression (JA)}:\\
さあ、どんどん(進:すす)めましょう\\
\textbf{Adjusted Expression (EN)}:\\
Come on, let's keep going

\textbf{Example (JA)}:\\
A「この(プロジェクト:ぷろじぇくと)の(条件:じょうけん)は、すべて(揃:そろ)った。さあ、どんどん(進:すす)めよう」\\
\textbf{Example (EN)}:\\
A: All the conditions for this project are in place. Let's keep going.

\textbf{Notes (JA)}:\\
「どんどん進めよう」は、止まらずに前に進むべきだということを表現するために使用されます。「どんどん」は、「ますます」や「着実に」という副詞です。「進めよう」は、「進める」という他動詞の意向形です。

\textbf{Notes (EN)}:\\
The phrase "dondon susumeyou" is used to express that you should continue to move forward without stopping. The word "dondon" is an adverb that means "more and more" or "steadily." The word "susumeyou" is the volitional form of the verb "susumeru," which means "to advance" or "to proceed."

\newpage

\section*{139. Suddenly}

\textbf{Date}: 2024.08.27 (Tue)\\
\textbf{Index}: totsuzen\\
\textbf{Expression (JA)}: (突然:とつぜん)\\
\textbf{Romanization}: Totsuzen\\
\textbf{Expression (EN)}: Suddenly

\textbf{Variants (JA)}: 突然 / 突然に / 突然ですが\\
\textbf{Variants (EN)}: suddenly / all of a sudden / suddenly, but ...

\textbf{Intent (JA)}: 出来事の突然性の表現, 驚きの共有, 予期しない展開の導入\\
\textbf{Intent (EN)}: expressing suddenness, sharing surprise, introducing an unexpected turn

\textbf{Tags (JA)}: 即時文法, 副詞, 口語, 驚き, 予期せぬ出来事\\
\textbf{Tags (EN)}: immediate grammar, adverb, colloquial, surprise, unexpected event

\textbf{Adjusted Expression (JA)}:\\
(突然:とつぜん)ではありますが、...\\
\textbf{Adjusted Expression (EN)}:\\
Suddenly, but ...

\textbf{Example (JA)}:\\
A「(突然:とつぜん)、(電話:でんわ)が(鳴:な)って、びっくりした」\\
\textbf{Example (EN)}:\\
A: Suddenly, the phone rang and I was surprised.

\textbf{Notes (JA)}:\\
「突然」は、突然、予期せずにという副詞です。

\textbf{Notes (EN)}:\\
The word "totsuzen" is an adverb that means "suddenly" or "unexpectedly."

\newpage

\section*{138. About ...}

\textbf{Date}: 2024.08.26 (Mon)\\
\textbf{Index}: nitsuite\\
\textbf{Expression (JA)}: ～について\\
\textbf{Romanization}: \\textasciitilde{}Ni tsuite\\
\textbf{Expression (EN)}: About ...

\textbf{Variants (JA)}: ～について / ～につきまして / ～については\\
\textbf{Variants (EN)}: about ... / concerning ... / regarding ...

\textbf{Intent (JA)}: 話題の提示, 説明の対象の明示, 情報を求める表現\\
\textbf{Intent (EN)}: presenting a topic, indicating the target of explanation, asking for information

\textbf{Tags (JA)}: 即時文法, 助詞表現, 話題提示, 説明, 質問\\
\textbf{Tags (EN)}: immediate grammar, particle expression, topic presentation, explanation, question

\textbf{Adjusted Expression (JA)}:\\
～についてですが、...\\
\textbf{Adjusted Expression (EN)}:\\
About ..., but ...

\textbf{Example (JA)}:\\
A「(書類:しょるい)の(書:か)き(方:かた)について(教:おし)えていただけませんか」\\
\textbf{Example (EN)}:\\
A: Could you tell me how to write the document?

\textbf{Notes (JA)}:\\
「～について」は、何かについての情報や指示を求めるために使用されます。

\textbf{Notes (EN)}:\\
The phrase "ni tsuite" is used to ask for information or instructions about something.

\newpage

\section*{137. By the way}

\textbf{Date}: 2024.08.25 (Sun)\\
\textbf{Index}: tokorode\\
\textbf{Expression (JA)}: ところで\\
\textbf{Romanization}: Tokorode\\
\textbf{Expression (EN)}: By the way

\textbf{Variants (JA)}: ところで / ところでさ / ところでね\\
\textbf{Variants (EN)}: by the way / by the way, / anyway, by the way

\textbf{Intent (JA)}: 話題転換, 質問への導入, 会話の流れの切り替え\\
\textbf{Intent (EN)}: changing the topic, introducing a question, shifting the flow of conversation

\textbf{Tags (JA)}: 即時文法, 口語, 談話標識, 話題転換, 会話運営\\
\textbf{Tags (EN)}: immediate grammar, colloquial, discourse marker, topic shift, conversation management

\textbf{Adjusted Expression (JA)}:\\
ちょっと、(話題:わだい)は(変:か)わるんですが、...\\
\textbf{Adjusted Expression (EN)}:\\
Well, this is a bit of a change of topic, but ...

\textbf{Example (JA)}:\\
A「(昨日:きのう)、(映画:えいが)を(見:み)たよ」 B「そうなんだ。ところで、(明日:あした)の(予定:よてい)は？」 A「(特:とく)にないよ」\\
\textbf{Example (EN)}:\\
A: I watched a movie yesterday. B: I see. By the way, what are your plans for tomorrow? A: Nothing special.

\textbf{Notes (JA)}:\\
「ところで」は、会話の話題を変えたり、質問をしたりするために使用されます。

\textbf{Notes (EN)}:\\
The phrase "tokorode" is used to change the subject of the conversation or to ask a question.

\newpage

\section*{136. Of course}

\textbf{Date}: 2024.08.24 (Sat)\\
\textbf{Index}: mochiron\\
\textbf{Expression (JA)}: もちろん\\
\textbf{Romanization}: Mochiron\\
\textbf{Expression (EN)}: Of course

\textbf{Variants (JA)}: もちろん / もちろんです / もちろんだよ\\
\textbf{Variants (EN)}: of course / of course it is / of course

\textbf{Intent (JA)}: 当然さの表明, 肯定の応答, 即答的な反応\\
\textbf{Intent (EN)}: expressing obviousness, giving an affirmative response, an immediate response

\textbf{Tags (JA)}: 即時文法, 口語, 応答表現, 肯定, 当然\\
\textbf{Tags (EN)}: immediate grammar, colloquial, response expression, affirmation, obviousness

\textbf{Adjusted Expression (JA)}:\\
もちろんいいですよ。\\
\textbf{Adjusted Expression (EN)}:\\
Of course, that's fine.

\textbf{Example (JA)}:\\
A「(明日:あした)、(会:あ)える？」 B「もちろん」\\
\textbf{Example (EN)}:\\
A: Can we meet tomorrow? B: Of course.

\textbf{Notes (JA)}:\\
「もちろん」は、何かが当然であることを表現するために使用されます。

\textbf{Notes (EN)}:\\
The phrase "mochiron" is used to express that something is natural or obvious.

\newpage

\section*{135. A habit}

\textbf{Date}: 2024.08.23 (Fri)\\
\textbf{Index}: shuukan\\
\textbf{Expression (JA)}: (習慣:しゅうかん)\\
\textbf{Romanization}: Shuukan\\
\textbf{Expression (EN)}: A habit

\textbf{Variants (JA)}: 習慣 / 習慣になる / 新しい習慣\\
\textbf{Variants (EN)}: a habit / become a habit / a new habit

\textbf{Intent (JA)}: 習慣の説明, 状態の変化の報告, 日常的行動の言及\\
\textbf{Intent (EN)}: describing a habit, reporting a change of state, mentioning everyday behavior

\textbf{Tags (JA)}: 即時文法, 名詞, 状態変化, 日常語彙\\
\textbf{Tags (EN)}: immediate grammar, noun, state change, everyday vocabulary

\textbf{Adjusted Expression (JA)}:\\
(新:あたら)しい(習慣:しゅうかん)ができました\\
\textbf{Adjusted Expression (EN)}:\\
I've made a new habit

\textbf{Example (JA)}:\\
A「(毎日:まいにち)、(ジョギング:じょぎんぐ)することが(習慣:しゅうかん)になった」\\
\textbf{Example (EN)}:\\
A: I've made it a habit to jog every day.

\textbf{Notes (JA)}:\\
「習慣」とは、毎日何かを繰り返し行うことを指します。たとえば、健康のために、早寝早起きをする、少しヨーグルトを毎日食べる、などです。

\textbf{Notes (EN)}:\\
The phrase "atarashii shuukan" is used to express that you have started a new habit. "Atarashii" is an i-adjective meaning "new," and "shuukan" is a noun meaning "habit." For example, sleeping early and waking up early for one's health, eating a little yogurt every day, etc.

\newpage

\section*{134. Made it in time}

\textbf{Date}: 2024.08.22 (Thu)\\
\textbf{Index}: maniatta\\
\textbf{Expression (JA)}: (間:ま)に(合:あ)った\\
\textbf{Romanization}: Ma ni atta\\
\textbf{Expression (EN)}: Made it in time

\textbf{Variants (JA)}: 間に合った / 間に合いました / ぎりぎり間に合った\\
\textbf{Variants (EN)}: made it in time / made it just in time / barely made it in time

\textbf{Intent (JA)}: 間に合ったことの報告, 安堵の共有, 時間に関する即時的な反応\\
\textbf{Intent (EN)}: reporting that something was in time, sharing relief, immediate reaction about timing

\textbf{Tags (JA)}: 即時文法, 口語, 時間表現, 結果報告, 安堵\\
\textbf{Tags (EN)}: immediate grammar, colloquial, time expression, reporting a result, relief

\textbf{Adjusted Expression (JA)}:\\
(間:ま)に(合:あ)いました\\
\textbf{Adjusted Expression (EN)}:\\
I made it in time

\textbf{Example (JA)}:\\
ひとつ(前:まえ)の(電車:でんしゃ)に(間:ま)に(合:あ)ったので、(予定:よてい)より(早:はや)く(着:つ)きました。\\
\textbf{Example (EN)}:\\
I made it to the earlier train, so I arrived earlier than planned.

\textbf{Notes (JA)}:\\
「間に合った」は、何かにちょうど間に合ったことを表現します。「間」は「時間」や「間隔」を意味し、「に」は目標や時間を示す助詞です。「合う」は「一致する」「適合する」または「間に合う」という意味の動詞です。実際のところ、人生では間に合わないことのほうが多いかもしれませんね。

\textbf{Notes (EN)}:\\
"Ma ni atta" expresses that you managed to make it in time. "Ma" means "time" or "interval," and "ni" is a particle indicating the target time or goal. "Au" is a verb meaning "to match," "to fit," or "to be in time." In real life, there are probably more situations where you don't make it in time than when you do.

\newpage

\section*{133. A driver's license}

\textbf{Date}: 2024.08.21 (Wed)\\
\textbf{Index}: untenmenkyo\\
\textbf{Expression (JA)}: (運転:うんてん)(免許:めんきょ)\\
\textbf{Romanization}: Unten menkyo\\
\textbf{Expression (EN)}: Driver's license

\textbf{Variants (JA)}: 運転免許 / 運転免許を取る\\
\textbf{Variants (EN)}: driver's license / get a driver's license

\textbf{Intent (JA)}: 資格の言及, 取得の表現, 日常語彙\\
\textbf{Intent (EN)}: mentioning a qualification, expressing obtaining it, everyday vocabulary

\textbf{Tags (JA)}: 即時文法, 名詞, 複合名詞, 資格, 日常語彙\\
\textbf{Tags (EN)}: immediate grammar, noun, compound noun, qualification, everyday vocabulary

\textbf{Adjusted Expression (JA)}:\\
(運転:うんてん)(免許:めんきょ)を(取:と)りました。\\
\textbf{Adjusted Expression (EN)}:\\
I got my driver's license.

\textbf{Example (JA)}:\\
A「(運転:うんてん)(免許:めんきょ)を(取:と)った」\\
\textbf{Example (EN)}:\\
A: I got my driver's license.

\textbf{Notes (JA)}:\\
「運転免許」は、2つの名詞「運転（driving）」と「免許（license）」からなる複合名詞です。

\textbf{Notes (EN)}:\\
The phrase "unten menkyo" is a compound noun composed of two nouns: "unten" (driving) and "menkyo" (license).

\newpage

\section*{132. I hope ...}

\textbf{Date}: 2024.08.20 (Tue)\\
\textbf{Index}: dearimasuyouni\\
\textbf{Expression (JA)}: ～でありますように\\
\textbf{Romanization}: \\textasciitilde{}de arimasu you ni\\
\textbf{Expression (EN)}: I hope ...

\textbf{Variants (JA)}: ～でありますように / ～でありますように。\\
\textbf{Variants (EN)}: I hope ... / May ...

\textbf{Intent (JA)}: 願いの表現, 希望の表現, 祈りの表現\\
\textbf{Intent (EN)}: expressing a wish, expressing hope, expressing a prayer

\textbf{Tags (JA)}: 即時文法, 口語, 願望表現, 希望表現, 祈願表現\\
\textbf{Tags (EN)}: immediate grammar, spoken, wish expression, hope expression, prayer expression

\textbf{Adjusted Expression (JA)}:\\
～であることを(願:ねが)います\\
\textbf{Adjusted Expression (EN)}:\\
I hope that ...

\textbf{Example (JA)}:\\
A「(世:よ)の(中:なか)、(平和:へいわ)でありますように」 B「(僕:ぼく)は、(試験:しけん)、(合格:ごうかく)できますように」\\
\textbf{Example (EN)}:\\
A: I hope there will be peace in the world. B: I hope I can pass the exam.

\textbf{Notes (JA)}:\\
「～でありますように」は、願いや希望を表現するために使用されます。この表現は、「（動詞＋できます）ように」や「（名詞＋であります）ように」といったパターンで使われます。

\textbf{Notes (EN)}:\\
The phrase "\\textasciitilde{}de arimasu you ni" is used to express a wish or hope. This expression is used in patterns such as "(verb) you ni" and "(noun + de arimasu) you ni."

\newpage

\section*{131. move your hands}

\textbf{Date}: 2024.08.19 (Mon)\\
\textbf{Index}: tewougokasu\\
\textbf{Expression (JA)}: (手:て)を(動:うご)かす\\
\textbf{Romanization}: Te wo ugokasu\\
\textbf{Expression (EN)}: move your hands

\textbf{Variants (JA)}: 手を動かす / 手を動かしてみる\\
\textbf{Variants (EN)}: move your hands / try doing it with your hands

\textbf{Intent (JA)}: 実践の促し, 理解の促進, 行動のうながし\\
\textbf{Intent (EN)}: encouraging practice, promoting understanding, prompting action

\textbf{Tags (JA)}: 即時文法, 口語, 実践表現, 理解促進表現, 行動促進表現\\
\textbf{Tags (EN)}: immediate grammar, spoken, practice expression, understanding-promoting expression, action-prompting expression

\textbf{Adjusted Expression (JA)}:\\
(実際:じっさい)に(手:て)を(動:うご)かしてください\\
\textbf{Adjusted Expression (EN)}:\\
Please try doing it yourself.

\textbf{Example (JA)}:\\
A「(考:かんが)えるだけでなく、(実際:じっさい)に(手:て)を(動:うご)かしてみると、よくわかるよ」\\
\textbf{Example (EN)}:\\
A: If you don't just think about it but actually try doing it with your hands, you'll understand it better.

\textbf{Notes (JA)}:\\
「手を動かす」は、何かを理解するために手を動かすべきだということを表現するために使用されます。

\textbf{Notes (EN)}:\\
The phrase "te wo ugokasu" is used to express that you should actually use your hands in order to understand something. The word "te" is a noun that means "hand." The word "ugokasu" is a verb that means "to move."

\newpage

\section*{130. In a nutshell}

\textbf{Date}: 2024.08.18 (Sun)\\
\textbf{Index}: yousuruni\\
\textbf{Expression (JA)}: (要:よう)するに...\\
\textbf{Romanization}: Yousuruni...\\
\textbf{Expression (EN)}: In a nutshell,...

\textbf{Variants (JA)}: 要するに... / 要するに\\
\textbf{Variants (EN)}: In a nutshell,... / To put it simply,...

\textbf{Intent (JA)}: 要約の表現, 要点提示の表現, 説明の整理\\
\textbf{Intent (EN)}: expressing a summary, presenting the main point, organizing an explanation

\textbf{Tags (JA)}: 即時文法, 口語, 要約表現, 説明表現, 整理表現\\
\textbf{Tags (EN)}: immediate grammar, spoken, summary expression, explanation expression, organizing expression

\textbf{Adjusted Expression (JA)}:\\
(要:よう)するに、そういうことです。\\
\textbf{Adjusted Expression (EN)}:\\
In short, that is the point.

\textbf{Example (JA)}:\\
A「(要:よう)するに、この(機械:きかい)は(誰:だれ)にでも(使:つか)えるということです」\\
\textbf{Example (EN)}:\\
A: In a nutshell, this machine can be used by anyone.

\textbf{Notes (JA)}:\\
「要するに」は、何かを要約したり、要点を説明したりすることを表現するために使用されます。

\textbf{Notes (EN)}:\\
The phrase "yousuruni" is used to express that something is being summarized or that the main point is being explained briefly.

\newpage

\section*{129. Continuation is power}

\textbf{Date}: 2024.08.17 (Sat)\\
\textbf{Index}: keizoku\\
\textbf{Expression (JA)}: (継続:けいぞく)は(力:ちから)なり\\
\textbf{Romanization}: Keizoku ha chikara nari\\
\textbf{Expression (EN)}: Continuation is power

\textbf{Variants (JA)}: 継続は力なり\\
\textbf{Variants (EN)}: Continuation is power / Perseverance is strength

\textbf{Intent (JA)}: 教訓の表現, 価値の提示, 励まし\\
\textbf{Intent (EN)}: expressing a moral, presenting a value, encouragement

\textbf{Tags (JA)}: 即時文法, 口語, 諺表現, 価値表現, 励まし表現\\
\textbf{Tags (EN)}: immediate grammar, spoken, proverb expression, value expression, encouragement expression

\textbf{Adjusted Expression (JA)}:\\
(継続:けいぞく)は(力:ちから)なりです\\
\textbf{Adjusted Expression (EN)}:\\
Continuing is powerful.

\textbf{Example (JA)}:\\
A「(毎日:まいにち)、(勉強:べんきょう)していると、(成績:せいせき)が(上:あ)がるよね」 B「(継続:けいぞく)は(力:ちから)なり」\\
\textbf{Example (EN)}:\\
A: If you study every day, your grades will improve. B: Continuation is power.

\textbf{Notes (JA)}:\\
「継続は力なり」は、何かを続けることが力であることを表現した諺です。「継続」は、「続けること」や「連続」を意味する名詞です。

\textbf{Notes (EN)}:\\
The phrase "keizoku ha chikara nari" is a proverb expressing that continuing something leads to strength or success. "Keizoku" means continuation, and "chikara" means power or strength.

\newpage

\section*{128. Can't eat it all}

\textbf{Date}: 2024.08.16 (Fri)\\
\textbf{Index}: tabekirenai\\
\textbf{Expression (JA)}: (食:た)べきれない\\
\textbf{Romanization}: Tabekirenai\\
\textbf{Expression (EN)}: Can't eat it all

\textbf{Variants (JA)}: 食べきれない / 食べ切れない\\
\textbf{Variants (EN)}: Can't eat it all / Can't finish eating it

\textbf{Intent (JA)}: 限界の表現, 完遂不能の表現, 状態の説明\\
\textbf{Intent (EN)}: expressing a limit, expressing inability to finish, explaining a state

\textbf{Tags (JA)}: 即時文法, 口語, 限界表現, 完遂不能表現, 状態表現\\
\textbf{Tags (EN)}: immediate grammar, spoken, limit expression, inability-to-finish expression, state expression

\textbf{Adjusted Expression (JA)}:\\
(食:た)べきれません\\
\textbf{Adjusted Expression (EN)}:\\
I cannot eat it all.

\textbf{Example (JA)}:\\
A「(美味:おい)しい(料理:りょうり)が(多:おお)すぎて、(食:た)べきれない」\\
\textbf{Example (EN)}:\\
A: There are too many delicious dishes and I can't eat them all.

\textbf{Notes (JA)}:\\
「食べきれない」は、何かをすべて食べることができないことを表現するために使用されます。可能の表現は、場合によって、いくつかの意味があります。「食べきれない」には、「もうお腹がいっぱいで、これ以上は食べられない」という意味になります。「不味くて食べられない」や「毒物が含まれていて食べられない」の場合には、「～きれない」の可能の表現は使えません。

\textbf{Notes (EN)}:\\
The phrase "tabekirenai" is used to express that you cannot finish eating all of something. Potential expressions can have different meanings depending on the situation. In the case of "tabekirenai," it means something like "I'm full and can't eat any more." In cases such as "mazukute taberarenai" (I cannot eat it because it tastes bad) or "dokubutsu ga fukumareteite taberarenai" (I cannot eat it because it contains poison), the pattern "-kirenai" cannot be used.

\newpage

\section*{127. I've done everything I can}

\textbf{Date}: 2024.08.15 (Thu)\\
\textbf{Index}: yaritsukushita\\
\textbf{Expression (JA)}: やり(尽:つ)くした\\
\textbf{Romanization}: Yaritsukushita\\
\textbf{Expression (EN)}: I've done everything I can

\textbf{Variants (JA)}: やり尽くした / やり切った\\
\textbf{Variants (EN)}: I've done everything I can / I've done it all

\textbf{Intent (JA)}: 完了の表現, 達成感の表現, 限界の提示\\
\textbf{Intent (EN)}: expressing completion, expressing a sense of accomplishment, showing one's limit

\textbf{Tags (JA)}: 即時文法, 口語, 完了表現, 達成表現, 限界表現\\
\textbf{Tags (EN)}: immediate grammar, spoken, completion expression, accomplishment expression, limit expression

\textbf{Adjusted Expression (JA)}:\\
やり(尽:つ)くしました\\
\textbf{Adjusted Expression (EN)}:\\
I have done everything I can.

\textbf{Example (JA)}:\\
A「(準備:じゅんび)としてはやり(尽:つ)くした」\\
\textbf{Example (EN)}:\\
A: I've done everything I can to prepare.

\textbf{Notes (JA)}:\\
「やり尽くした」は、自分ができることをすべてやったことを表現するために使用されます。「やり」は、「する」や「行う」という動詞のマス形です。「尽くす」は、「使い果たす」や「使い切る」という動詞です。

\textbf{Notes (EN)}:\\
The phrase "yaritsukushita" is used to express that you have done everything you can. "Yari" comes from the verb "yaru," and "tsukushita" is the past form of "tsukusu," which means to use up or exhaust completely.

\newpage

\section*{126. Pros and cons}

\textbf{Date}: 2024.08.14 (Wed)\\
\textbf{Index}: sannpiryouron\\
\textbf{Expression (JA)}: (賛否両論:さんぴりょうろん)\\
\textbf{Romanization}: Sanpi ryouronn\\
\textbf{Expression (EN)}: Pros and cons

\textbf{Variants (JA)}: 賛否両論\\
\textbf{Variants (EN)}: pros and cons / mixed opinions

\textbf{Intent (JA)}: 評価の対立の表現, 意見の分裂の表現, 議論の共有\\
\textbf{Intent (EN)}: expressing opposing evaluations, expressing divided opinions, sharing a discussion

\textbf{Tags (JA)}: 即時文法, 口語, 評価表現, 議論表現, 対立表現\\
\textbf{Tags (EN)}: immediate grammar, spoken, evaluation expression, discussion expression, contrast expression

\textbf{Adjusted Expression (JA)}:\\
(賛否両論:さんぴりょうろん)があります\\
\textbf{Adjusted Expression (EN)}:\\
There are pros and cons.

\textbf{Example (JA)}:\\
A「(新:あたら)しい(政策:せいさく)には、(賛否両論:さんぴりょうろん)があって、そう(簡単:かんたん)には(結論:けつろん)は(言:い)えませんよね」\\
\textbf{Example (EN)}:\\
A: There are pros and cons to the new policy and it's not easy to come to a conclusion.

\textbf{Notes (JA)}:\\
「賛否両論」は、何かについて肯定的な意見と否定的な意見があることを表現するために使用されます。

\textbf{Notes (EN)}:\\
The phrase "sanpi ryoron" expresses that there are both positive and negative opinions about something.

\newpage

\section*{125. more or less}

\textbf{Date}: 2024.08.13 (Tue)\\
\textbf{Index}: tashou\\
\textbf{Expression (JA)}: (多少:たしょう)\\
\textbf{Romanization}: Tashou\\
\textbf{Expression (EN)}: More or less

\textbf{Variants (JA)}: 多少 / 多少の / 多少は\\
\textbf{Variants (EN)}: more or less / some / to some extent

\textbf{Intent (JA)}: 程度の表現, 数量の表現, 控えめな言い方\\
\textbf{Intent (EN)}: expressing degree, expressing quantity, speaking in a restrained way

\textbf{Tags (JA)}: 即時文法, 口語, 程度表現, 数量表現, 控えめ表現\\
\textbf{Tags (EN)}: immediate grammar, spoken, degree expression, quantity expression, restrained expression

\textbf{Adjusted Expression (JA)}:\\
(多少:たしょう)あります\\
\textbf{Adjusted Expression (EN)}:\\
There is some.

\textbf{Example (JA)}:\\
A「(今日:きょう)は、(雨:あめ)が(多少:たしょう)(降:ふ)るかもしれない」\\
\textbf{Example (EN)}:\\
A: There may be some rain today.

\textbf{Notes (JA)}:\\
「多少」は、何かが少しあることを表現するために使用されます。これは、分量だけでなく、技術や能力などの程度を表すこともできます。

\textbf{Notes (EN)}:\\
The phrase "tashou" is used to express that there is some amount or degree of something. It can be used not only for quantity but also for the degree of something such as skill or ability.

\newpage

\section*{124. Take a deep breath}

\textbf{Date}: 2024.08.12 (Mon)\\
\textbf{Index}: shinnkokyuu\\
\textbf{Expression (JA)}: (深呼吸:しんこきゅう)して\\
\textbf{Romanization}: Shinkokyuu shite\\
\textbf{Expression (EN)}: Take a deep breath

\textbf{Variants (JA)}: 深呼吸して / 深呼吸してみて\\
\textbf{Variants (EN)}: Take a deep breath / Just take a deep breath

\textbf{Intent (JA)}: 落ち着きを促す表現, 気づかいの表現, 行動のうながし\\
\textbf{Intent (EN)}: encouraging calmness, expressing care, prompting an action

\textbf{Tags (JA)}: 即時文法, 口語, 気づかい表現, 行動促進表現, 落ち着き表現\\
\textbf{Tags (EN)}: immediate grammar, spoken, caring expression, action-prompting expression, calming expression

\textbf{Adjusted Expression (JA)}:\\
深呼吸してください\\
\textbf{Adjusted Expression (EN)}:\\
Please take a deep breath.

\textbf{Example (JA)}:\\
A「(緊張:きんちょう)しているときは、(深呼吸:しんこきゅう)して」\\
\textbf{Example (EN)}:\\
A: When you're nervous, take a deep breath.

\textbf{Notes (JA)}:\\
「深呼吸して」は、緊張したり不安なときに深呼吸をするようにということを表現するために使用されます。

\textbf{Notes (EN)}:\\
The phrase "shinkokyuu shite" is used to tell someone to take a deep breath when they are nervous or anxious. The word "shinkokyuu" is a noun meaning "deep breath."

\newpage

\section*{123. This and that}

\textbf{Date}: 2024.08.11 (Sun)\\
\textbf{Index}: arekore\\
\textbf{Expression (JA)}: あれこれ\\
\textbf{Romanization}: Arekore\\
\textbf{Expression (EN)}: This and that

\textbf{Variants (JA)}: あれこれ / あれやこれや\\
\textbf{Variants (EN)}: this and that / various things

\textbf{Intent (JA)}: 不特定のものの表現, 曖昧な列挙, 会話の省略\\
\textbf{Intent (EN)}: referring to unspecified things, vague listing, ellipsis in conversation

\textbf{Tags (JA)}: 即時文法, 口語, 曖昧表現, 列挙表現, 省略表現\\
\textbf{Tags (EN)}: immediate grammar, spoken, vague expression, listing expression, ellipsis expression

\textbf{Adjusted Expression (JA)}:\\
いろいろなもの\\
\textbf{Adjusted Expression (EN)}:\\
various things

\textbf{Example (JA)}:\\
A「(昨日:きのう)、(何:なに)(食:た)べた？」 B「あれこれ(食:た)べた」\\
\textbf{Example (EN)}:\\
A: What did you eat yesterday? B: I tried a bit of everything.

\textbf{Notes (JA)}:\\
「あれこれ」は、「いろいろな」という意味を表す表現です。会話では、何かを一つ一つ具体的に言わなくても、複数のものをまとめて表すことができます。特に即時的なやりとりでは、個別の名詞を並べる代わりに、雑多にいろいろあることを手早く表せる便利な表現です。

\textbf{Notes (EN)}:\\
The phrase "arekore" means various things. In conversation, it lets the speaker refer to multiple things without naming them one by one. Especially in immediate interaction, it is a convenient expression for quickly conveying a mixed or miscellaneous set of things.

\newpage

\section*{122. Exactly!}

\textbf{Date}: 2024.08.10 (Sat)\\
\textbf{Index}: sonotoori\\
\textbf{Expression (JA)}: そのとおり！\\
\textbf{Romanization}: Sono toori!\\
\textbf{Expression (EN)}: Exactly!

\textbf{Variants (JA)}: そのとおり / その通り\\
\textbf{Variants (EN)}: Exactly! / That's right!

\textbf{Intent (JA)}: 同意の表現, 発話内容の確認, 会話の応答\\
\textbf{Intent (EN)}: expressing agreement, confirming what was said, responding in conversation

\textbf{Tags (JA)}: 即時文法, 口語, 同意表現, 応答表現, 会話表現\\
\textbf{Tags (EN)}: immediate grammar, spoken, agreement expression, response expression, conversation expression

\textbf{Adjusted Expression (JA)}:\\
そのとおりです\\
\textbf{Adjusted Expression (EN)}:\\
That is correct.

\textbf{Example (JA)}:\\
A「(答:こたえ)は、ない、ということですか」 B「そのとおり！」\\
\textbf{Example (EN)}:\\
A: So you mean there is no answer? B: Exactly!

\textbf{Notes (JA)}:\\
「その通り」は、「正しい」という意味ではなく、単純に相手が言ったことが正しいことを表します。

\textbf{Notes (EN)}:\\
The phrase "sono toori" is used to express agreement with what the other person has said. It confirms that the statement matches the speaker's understanding, rather than making an independent judgment of correctness.

\newpage

\section*{121. It's fresh!}

\textbf{Date}: 2024.08.09 (Fri)\\
\textbf{Index}: shinsen\\
\textbf{Expression (JA)}: (新鮮:しんせん)\\
\textbf{Romanization}: Shinsen\\
\textbf{Expression (EN)}: Fresh

\textbf{Variants (JA)}: 新鮮 / 新鮮だ / 新鮮な\\
\textbf{Variants (EN)}: Fresh / It's fresh / fresh

\textbf{Intent (JA)}: 評価の表現, 状態の表現, 印象の共有\\
\textbf{Intent (EN)}: expressing an evaluation, expressing a state, sharing an impression

\textbf{Tags (JA)}: 即時文法, 口語, 評価表現, 状態表現, 印象表現\\
\textbf{Tags (EN)}: immediate grammar, spoken, evaluation expression, state expression, impression expression

\textbf{Adjusted Expression (JA)}:\\
(新鮮:しんせん)です\\
\textbf{Adjusted Expression (EN)}:\\
It is fresh.

\textbf{Example (JA)}:\\
A「(新鮮:しんせん)な(野菜:やさい)が(食:た)べたい」\\
\textbf{Example (EN)}:\\
A: I want to eat fresh vegetables.

\textbf{Notes (JA)}:\\
「新鮮」は、収穫されてまもないという意味のナ形容詞です。食べ物だけでなく、斬新なという意味で「新鮮なアイデア」のような使い方もできます。

\textbf{Notes (EN)}:\\
The word "shinsen" is a na-adjective that means freshly harvested. It can be used not only for food but also for fresh ideas, such as "fresh ideas."

\newpage

\section*{120. Looks good}

\textbf{Date}: 2024.08.08 (Thu)\\
\textbf{Index}: yosasoune\\
\textbf{Expression (JA)}: よさそうね\\
\textbf{Romanization}: Yosasou ne\\
\textbf{Expression (EN)}: Looks good

\textbf{Variants (JA)}: よさそうね / よさそう / よさげだね\\
\textbf{Variants (EN)}: Looks good / Seems good / Looks promising

\textbf{Intent (JA)}: 印象の表現, 推量の表現, 評価の共有\\
\textbf{Intent (EN)}: expressing an impression, expressing conjecture, sharing an evaluation

\textbf{Tags (JA)}: 即時文法, 口語, 印象表現, 推量表現, 評価表現\\
\textbf{Tags (EN)}: immediate grammar, spoken, impression expression, conjecture expression, evaluation expression

\textbf{Adjusted Expression (JA)}:\\
よさそうですね\\
\textbf{Adjusted Expression (EN)}:\\
It looks good.

\textbf{Example (JA)}:\\
A「これ、どう？」 B「よさそうね」\\
\textbf{Example (EN)}:\\
A: What do you think of this? B: Looks good.

\textbf{Notes (JA)}:\\
「よさそうね」は、何かが良さそうだと思われることを表現するために使用されます。「よさそう」は、「良い」という形容詞「よい」に「そう」が付いたもので、「良さそう」や「良いようだ」という意味です。

\textbf{Notes (EN)}:\\
The phrase "yosasou ne" is used to express that something looks good or seems good. "Yosasou" comes from the adjective "yoi" and the ending "sou," and it expresses the speaker's impression or guess based on what they see or sense.

\newpage

\section*{119. One shot exit}

\textbf{Date}: 2024.08.07 (Wed)\\
\textbf{Index}: ippatsutaijou\\
\textbf{Expression (JA)}: (一発:いっぱつ)、(退場:たいじょう)\\
\textbf{Romanization}: Ippatsu taijou\\
\textbf{Expression (EN)}: One shot exit

\textbf{Variants (JA)}: 一発、退場 / 一発退場\\
\textbf{Variants (EN)}: One shot exit

\textbf{Intent (JA)}: 厳しい処分の表現, 即時判断の表現, スポーツ場面の共有\\
\textbf{Intent (EN)}: expressing a harsh penalty, expressing an immediate decision, sharing a sports situation

\textbf{Tags (JA)}: 即時文法, 口語, スポーツ表現, 処分表現, 即時判断表現\\
\textbf{Tags (EN)}: immediate grammar, spoken, sports expression, penalty expression, immediate-decision expression

\textbf{Adjusted Expression (JA)}:\\
たった(一度:いちど)の(反則:はんそく)で(即退場:そくたいじょう)です\\
\textbf{Adjusted Expression (EN)}:\\
It's an immediate exit with just one foul.

\textbf{Example (JA)}:\\
A「(悪質:あくしつ)な(反則:はんそく)で、(一発:いっぱつ)、(退場:たいじょう)！」\\
\textbf{Example (EN)}:\\
A: One shot exit for a malicious foul.

\textbf{Notes (JA)}:\\
「一発退場」は、誰かが一発でゲームや競技から退場することを表現するために使用されます。「一発」は 「一発」や「一撃」という名詞で、「退場」は「出口」や「撤退」という名詞です。誰かが一発でゲームや競技から退場することを表現するために使われます。

\textbf{Notes (EN)}:\\
The phrase "ippatsu taijou" is used to express that someone is ejected from a game or competition with a single foul.

\newpage

\section*{118. Smell}

\textbf{Date}: 2024.08.06 (Tue)\\
\textbf{Index}: nioi\\
\textbf{Expression (JA)}: (匂:にお)い\\
\textbf{Romanization}: Nioi\\
\textbf{Expression (EN)}: Smell

\textbf{Variants (JA)}: 匂い / におい\\
\textbf{Variants (EN)}: smell / odor

\textbf{Intent (JA)}: 感覚の表現, 異変の察知, 印象の共有\\
\textbf{Intent (EN)}: expressing sensory perception, noticing something unusual, sharing an impression

\textbf{Tags (JA)}: 即時文法, 口語, 感覚表現, 嗅覚表現, 印象表現\\
\textbf{Tags (EN)}: immediate grammar, spoken, sensory expression, olfactory expression, impression expression

\textbf{Adjusted Expression (JA)}:\\
(匂:にお)いがします\\
\textbf{Adjusted Expression (EN)}:\\
There is a smell.

\textbf{Example (JA)}:\\
A「ちょっと(変:へん)な(匂:にお)いがする」\\
\textbf{Example (EN)}:\\
A: I smell something strange.

\textbf{Notes (JA)}:\\
「におい」は、何かの香りや芳香を指す名詞です。花や食べ物、香水などの香りを表現するために使われます。目に見えないが、鼻で感じることができるものの香りを表現するために使われます。「かおり」はいい匂いを意味し、「変な香り」とは言えません。

\textbf{Notes (EN)}:\\
The word "nioi" is a noun that refers to the smell of something. Unlike "kaori," which refers to a good smell, "nioi" can be used more broadly. It is used to express something that is not visible but can be sensed by the nose.

\newpage

\section*{117. Scent}

\textbf{Date}: 2024.08.05 (Mon)\\
\textbf{Index}: kaori\\
\textbf{Expression (JA)}: (香り:かおり)\\
\textbf{Romanization}: Kaori\\
\textbf{Expression (EN)}: Scent

\textbf{Variants (JA)}: 香り / 匂い\\
\textbf{Variants (EN)}: scent / fragrance

\textbf{Intent (JA)}: 感覚の表現, 好みの表現, 印象の共有\\
\textbf{Intent (EN)}: expressing sensory perception, expressing preference, sharing an impression

\textbf{Tags (JA)}: 即時文法, 口語, 感覚表現, 嗅覚表現, 印象表現\\
\textbf{Tags (EN)}: immediate grammar, spoken, sensory expression, olfactory expression, impression expression

\textbf{Adjusted Expression (JA)}:\\
香りです\\
\textbf{Adjusted Expression (EN)}:\\
It is a scent.

\textbf{Example (JA)}:\\
A「この(花:はな)の(香り:かおり)が(好:す)き」\\
\textbf{Example (EN)}:\\
A: I like the scent of this flower.

\textbf{Notes (JA)}:\\
「かおり」は、何かの香りや芳香を指す名詞です。花や食べ物、香水などの香りを表現するために使われます。目に見えないが、鼻で感じることができるものの香りを表現するために使われます。

\textbf{Notes (EN)}:\\
The word "kaori" is a noun that refers to the scent or fragrance of something, such as flowers, food, or perfume. It describes something that is not visible but can be perceived through smell.

\newpage

\section*{116. secret ingredient}

\textbf{Date}: 2024.08.04 (Sun)\\
\textbf{Index}: kakushiaji\\
\textbf{Expression (JA)}: (隠:かく)し(味:あじ)\\
\textbf{Romanization}: Kakushi aji\\
\textbf{Expression (EN)}: secret ingredient

\textbf{Variants (JA)}: 隠し味\\
\textbf{Variants (EN)}: secret ingredient / hidden flavor

\textbf{Intent (JA)}: 料理の工夫の提示, 驚きの共有, 知識の共有\\
\textbf{Intent (EN)}: pointing out a cooking trick, sharing a surprise, sharing knowledge

\textbf{Tags (JA)}: 即時文法, 口語, 料理表現, 工夫表現, 知識共有表現\\
\textbf{Tags (EN)}: immediate grammar, spoken, cooking expression, technique expression, knowledge-sharing expression

\textbf{Adjusted Expression (JA)}:\\
(隠:かく)し(味:あじ)です\\
\textbf{Adjusted Expression (EN)}:\\
It is a secret ingredient.

\textbf{Example (JA)}:\\
A「この(料理:りょうり)の(隠:かく)し(味:あじ)は、(チョコレート:ちょこれーと)ですね」\\
\textbf{Example (EN)}:\\
A: The secret ingredient in this dish is chocolate.

\textbf{Notes (JA)}:\\
「隠し味」は、食べ物に隠れた味があることを表現するために使用されます。「隠し」は、「隠す」という動詞の名詞形で、直後に名詞を付して、後ろの名詞を修飾することができます。「味」は、「味」や「風味」という名詞です。表面的には気がつかないが、おいしさを引き出すために使われる調味料や調理法を指します。

\textbf{Notes (EN)}:\\
The phrase "kakushi aji" refers to a hidden flavor or a secret ingredient that enhances a dish without being immediately noticeable. "Kakushi" comes from the verb "kakusu (to hide)" and functions as a modifier, while "aji" refers to taste or flavor.

\newpage

\section*{115. There is a secret}

\textbf{Date}: 2024.08.03 (Sat)\\
\textbf{Index}: hiketsugaaru\\
\textbf{Expression (JA)}: (秘訣:ひけつ)がある\\
\textbf{Romanization}: Hiketsu ga aru\\
\textbf{Expression (EN)}: There is a secret

\textbf{Variants (JA)}: 秘訣がある / コツがある\\
\textbf{Variants (EN)}: There is a secret / There is a trick to it

\textbf{Intent (JA)}: 方法の示唆, 成功要因の提示, 知識の共有\\
\textbf{Intent (EN)}: suggesting a method, pointing to a key to success, sharing knowledge

\textbf{Tags (JA)}: 即時文法, 口語, 方法表現, 成功表現, 知識共有表現\\
\textbf{Tags (EN)}: immediate grammar, spoken, method expression, success expression, knowledge-sharing expression

\textbf{Adjusted Expression (JA)}:\\
秘訣があります\\
\textbf{Adjusted Expression (EN)}:\\
There is a key to it.

\textbf{Example (JA)}:\\
A「(成功:せいこう)の(秘訣:ひけつ)は(何:なん)だと(思:おも)う？」 B「(努力:どりょく)と(根気:こんき)だと(思:おも)う」\\
\textbf{Example (EN)}:\\
A: What do you think is the secret to success? B: I think it's effort and patience.

\textbf{Notes (JA)}:\\
「秘訣」は、名詞です。技術的に何かができること、成功することの秘密、成功の鍵を指します。何かをするための秘密があることを表現するために使用されます。

\textbf{Notes (EN)}:\\
The phrase "hiketsu ga aru" is used to express that there is a secret, a key, or an effective way to do something successfully. "Hiketsu" is a noun that refers to a secret to doing something, a key to success, or an important point for achieving a result.

\newpage

\section*{114. I can't do anything about it}

\textbf{Date}: 2024.08.02 (Fri)\\
\textbf{Index}: dounimokounimo\\
\textbf{Expression (JA)}: どうにもこうにも...\\
\textbf{Romanization}: Dounimo kounimo...\\
\textbf{Expression (EN)}: I can't do anything about it

\textbf{Variants (JA)}: どうにもこうにもできない / どうにもならない\\
\textbf{Variants (EN)}: I can't do anything about it / There's nothing I can do

\textbf{Intent (JA)}: 無力感の表現, 行き詰まりの表現, 状況の共有\\
\textbf{Intent (EN)}: expressing helplessness, expressing being stuck, sharing a situation

\textbf{Tags (JA)}: 即時文法, 口語, 無力表現, 行き詰まり表現, 状況共有表現\\
\textbf{Tags (EN)}: immediate grammar, spoken, helplessness expression, stuck expression, situation-sharing expression

\textbf{Adjusted Expression (JA)}:\\
どうにもこうにもできません\\
\textbf{Adjusted Expression (EN)}:\\
There is nothing I can do about it.

\textbf{Example (JA)}:\\
A「この(問題:もんだい)はどうにもこうにも(解決:かいけつ)できない」\\
\textbf{Example (EN)}:\\
A: This problem can't be solved no matter what.

\textbf{Notes (JA)}:\\
「どうにもこうにも」は、状況に対して何もできないことを表現するために使用されます。

\textbf{Notes (EN)}:\\
The phrase "dounimo kounimo" is used to express that you cannot do anything about a situation, even after trying various ways.

\newpage

\section*{113. Right?}

\textbf{Date}: 2024.08.01 (Thu)\\
\textbf{Index}: deshou\\
\textbf{Expression (JA)}: ～でしょ？\\
\textbf{Romanization}: \\textasciitilde{}desho?\\
\textbf{Expression (EN)}: Right?

\textbf{Variants (JA)}: ～でしょ？ / ～でしょう？\\
\textbf{Variants (EN)}: Right? / Isn't it?

\textbf{Intent (JA)}: 同意の確認, 再確認の表現, 口語的な表現\\
\textbf{Intent (EN)}: confirming agreement, reconfirming what was said, colloquial expression

\textbf{Tags (JA)}: 即時文法, 口語, 同意確認表現, 再確認表現, 対人配慮表現\\
\textbf{Tags (EN)}: immediate grammar, spoken, agreement-confirming expression, reconfirming expression, interpersonal expression

\textbf{Adjusted Expression (JA)}:\\
～でしょうか。\\
\textbf{Adjusted Expression (EN)}:\\
Would you agree?

\textbf{Example (JA)}:\\
A「これ、おいしいね」 B「でしょ？」\\
\textbf{Example (EN)}:\\
A: This is delicious. B: Right?

\textbf{Notes (JA)}:\\
「～でしょう？」は、自分が言ったことを相手が認めたときに、再確認するときの表現です。このことばは友人の間で使えますが、目上の人に使うと失礼になりますので、注意が必要です。

\textbf{Notes (EN)}:\\
The phrase "\\textasciitilde{}desho?" is used when the listener acknowledges something the speaker has said, and the speaker then reconfirms it. It can be used among friends, but it may sound rude when used with superiors, so care is needed.

\newpage

\section*{112. I'm feeling much better}

\textbf{Date}: 2024.07.31 (Wed)\\
\textbf{Index}: daibuyokunarimashita\\
\textbf{Expression (JA)}: だいぶ(良:よ)くなりました\\
\textbf{Romanization}: Daibu yoku narimashita\\
\textbf{Expression (EN)}: I'm feeling much better

\textbf{Variants (JA)}: だいぶ良くなりました / だいぶ良くなっています / だいぶ良くなった\\
\textbf{Variants (EN)}: I'm feeling much better / I'm doing considerably better / I'm much better

\textbf{Intent (JA)}: 状態改善の表明, 回復の表現\\
\textbf{Intent (EN)}: expressing improvement

\textbf{Tags (JA)}: 即時文法, 口語, 状態表現, 回復表現\\
\textbf{Tags (EN)}: immediate grammar, spoken, state expression, recovery expression

\textbf{Adjusted Expression (JA)}:\\
だいぶ(良:よ)くなっています\\
\textbf{Adjusted Expression (EN)}:\\
I'm doing considerably better.

\textbf{Example (JA)}:\\
A「(風邪:かぜ)、(大丈夫:だいじょうぶ)？」 B「だいぶ(良:よ)くなったよ」\\
\textbf{Example (EN)}:\\
A: Are you okay with your cold? B: I'm much better.

\textbf{Notes (JA)}:\\
「だいぶ良くなりました」は、状況が改善されたことを表現するために使用されます。「だいぶ」は、「もっと」や「よりよく」という副詞です。

\textbf{Notes (EN)}:\\
The phrase "daibu yoku narimashita" is used to express that the situation has improved. The word "daibu" is an adverb that means "much" or "considerably." It is used to express that the situation has improved.

\newpage

\section*{111. I can't ask someone else to do it}

\textbf{Date}: 2024.07.30 (Tue)\\
\textbf{Index}: tanomuwakeniwai\\
\textbf{Expression (JA)}: (頼:たの)むわけにはいかない\\
\textbf{Romanization}: Tanomu wake ni wa ikanai\\
\textbf{Expression (EN)}: I can't ask someone else to do it

\textbf{Variants (JA)}: 頼むわけにはいかない / 頼めない\\
\textbf{Variants (EN)}: I can't ask someone else to do it / I can't ask anyone to do it

\textbf{Intent (JA)}: 規範的制約の表現, 責任の表明, 行動の抑制\\
\textbf{Intent (EN)}: expressing constraint, taking responsibility, restraining an action

\textbf{Tags (JA)}: 即時文法, 口語, 制約表現, 責任表現, 判断表現\\
\textbf{Tags (EN)}: immediate grammar, spoken, constraint expression, responsibility expression, judgment expression

\textbf{Adjusted Expression (JA)}:\\
(頼:たの)むわけにはいきません\\
\textbf{Adjusted Expression (EN)}:\\
I cannot ask someone else to do it.

\textbf{Example (JA)}:\\
(自分:じぶん)のことだから、(他人:たにん)に(頼:たの)むわけにはいかない。\\
\textbf{Example (EN)}:\\
It's my own responsibility, so I can't ask someone else to do it.

\textbf{Notes (JA)}:\\
「頼むわけにはいかない」は、誰かに何かを頼むことができないことを表現するために使用されます。自分の責任、自分の問題、自分の仕事など、他人にしてもらったのでは、意味がない場合や失礼にあたる場合に使われます。「～わけにはいかない」という表現は「できない」という意味を持ちますが、可能不可能という意味ではなく、そうすることが不都合である場合に使われます。

\textbf{Notes (EN)}:\\
The phrase "tanomu wake ni wa ikanai" is used to express that you cannot ask someone to do something. It is used when doing so would be inappropriate, such as when it is your own responsibility or problem. The pattern "\\textasciitilde{}wake ni wa ikanai" does not simply mean 'impossible'; it expresses that something cannot be done due to social, moral, or situational reasons.

\newpage

\section*{110. It defeats the purpose}

\textbf{Date}: 2024.07.29 (Mon)\\
\textbf{Index}: motomokomonai\\
\textbf{Expression (JA)}: 元も子もない\\
\textbf{Romanization}: Moto mo ko mo nai\\
\textbf{Expression (EN)}: It defeats the purpose

\textbf{Variants (JA)}: 元も子もない / それじゃ元も子もない\\
\textbf{Variants (EN)}: It defeats the purpose / That would defeat the purpose / That ruins everything

\textbf{Intent (JA)}: 無意味の指摘, 価値の喪失の表現, 行動の否定\\
\textbf{Intent (EN)}: pointing out pointlessness, expressing loss of value, rejecting an action

\textbf{Tags (JA)}: 即時文法, 口語, 評価表現, 否定評価表現, 結果評価表現\\
\textbf{Tags (EN)}: immediate grammar, spoken, evaluation expression, negative evaluation expression, result evaluation expression

\textbf{Adjusted Expression (JA)}:\\
それでは元も子もありません\\
\textbf{Adjusted Expression (EN)}:\\
That defeats the purpose.

\textbf{Example (JA)}:\\
無理なダイエットをして体調を崩しては、元も子もない。\\
\textbf{Example (EN)}:\\
If you ruin your health with an extreme diet, it defeats the purpose.

\textbf{Notes (JA)}:\\
「元も子もない」は、何かがその価値や意味を失い、続ける意味がなくなったことを表現する際に使います。「元」は「起源」や「源」を意味し、「子」は「子供」や「成果」を意味します。起源と成果の両方を失うことで、何も残らないことを比喩的に表現しています。

\textbf{Notes (EN)}:\\
The phrase "moto mo ko mo nai" is used to express that something loses all its value or meaning, making the effort pointless. The words "moto" (origin) and "ko" (result or outcome) together suggest losing both the basis and the result, leaving nothing of value.

\newpage

\section*{109. It's about time}

\textbf{Date}: 2024.07.28 (Sun)\\
\textbf{Index}: sorosoro\\
\textbf{Expression (JA)}: そろそろ\\
\textbf{Romanization}: Sorosoro\\
\textbf{Expression (EN)}: It's about time

\textbf{Variants (JA)}: そろそろ / もうそろそろ\\
\textbf{Variants (EN)}: It's about time / soon / It should be about time

\textbf{Intent (JA)}: 時間の接近, 行動の促し, 予測の表現\\
\textbf{Intent (EN)}: indicating time is near, prompting action, expressing expectation

\textbf{Tags (JA)}: 即時文法, 口語, 時間表現, 行動促進表現, 予測表現\\
\textbf{Tags (EN)}: immediate grammar, spoken, time expression, action-prompting expression, expectation expression

\textbf{Adjusted Expression (JA)}:\\
そろそろです\\
\textbf{Adjusted Expression (EN)}:\\
It is about time.

\textbf{Example (JA)}:\\
A「(電車:でんしゃ)、(来:こ)ないね」 B「そろそろ(来:く)るよ」\\
\textbf{Example (EN)}:\\
A: The train isn't coming. B: It should come soon.

\textbf{Notes (JA)}:\\
「そろそろ」は、もうすぐ何かが起こることを期待したり、何かを行うことを予定したりすることを表現するために使用されます。動詞の前に置いて、「そろそろ来るでしょう」「そろそろ行きましょう」「そろそろ寝ましょう」のように使えます。

\textbf{Notes (EN)}:\\
The phrase "sorosoro" is used to express that something is about to happen or that you are about to do something. It is often placed before a verb, as in "sorosoro kuru deshou" (it should come soon), "sorosoro ikimashou" (it's about time to go), and "sorosoro nemashou" (it's about time to go to bed).

\newpage

\section*{108. I see}

\textbf{Date}: 2024.07.27 (Sat)\\
\textbf{Index}: aasouiu\\
\textbf{Expression (JA)}: ああ、そういうことか！\\
\textbf{Romanization}: Aa, sou iu koto ka!\\
\textbf{Expression (EN)}: I see, that's what it is!

\textbf{Variants (JA)}: ああ、そういうことか！ / ああ、そうか！ / なるほど！\\
\textbf{Variants (EN)}: I see, that's what it is! / I see! / Oh, I see.

\textbf{Intent (JA)}: 理解の表現, 気づきの共有, 納得の表現\\
\textbf{Intent (EN)}: expressing understanding, sharing realization, expressing agreement

\textbf{Tags (JA)}: 即時文法, 口語, 理解表現, 気づき表現, 納得表現\\
\textbf{Tags (EN)}: immediate grammar, spoken, understanding expression, realization expression, agreement expression

\textbf{Adjusted Expression (JA)}:\\
ああ、そういうことですか\\
\textbf{Adjusted Expression (EN)}:\\
I see.

\textbf{Example (JA)}:\\
A「(エレベータ:えれべーた)が(来:こ)なくて、(不思議:ふしぎ)に(思:おも)って」 B「(点検中:てんけんちゅう)って(書:か)いてあるね」 A「ああ、そういうことか！」\\
\textbf{Example (EN)}:\\
A: The elevator didn't come, and I was wondering. B: It says 'under maintenance.' A: I see, that's what it is!

\textbf{Notes (JA)}:\\
「ああ、そういうことか！」は、相手の説明を理解したことを表現するために使用されます。相手の具体的な説明を聞いた後に使います。ネイティブスピーカーの日常会話でよく使われるフレーズです。特徴は、「ああ」「そう」「いう」「こと」「か」などそれぞれの単語は基本的なものばかりですが、使う場面でとても上手な言語使用がわかります。基本的な単語の組み合わせを上手に使ってみることも考えたほうがよいでしょう。

\textbf{Notes (EN)}:\\
The phrase "aa, sou iu koto ka!" is used to express that you have understood the other person's explanation. It is typically used after hearing a specific explanation. Although each word in the phrase is basic, their combination shows a natural and effective use of language in conversation.

\newpage

\section*{107. Can that really happen?}

\textbf{Date}: 2024.07.26 (Fri)\\
\textbf{Index}: sonnakotoaru\\
\textbf{Expression (JA)}: そんなことある？\\
\textbf{Romanization}: Sonna koto aru?\\
\textbf{Expression (EN)}: Can that really happen?

\textbf{Variants (JA)}: そんなことある？ / そんなことあるの？ / マジで？\\
\textbf{Variants (EN)}: Can that really happen? / Does that really happen? / Seriously?

\textbf{Intent (JA)}: 驚きの表現, 不信の表現, 可能性の確認\\
\textbf{Intent (EN)}: expressing surprise, expressing disbelief, checking possibility

\textbf{Tags (JA)}: 即時文法, 口語, 驚き表現, 不信表現, 確認表現\\
\textbf{Tags (EN)}: immediate grammar, spoken, surprise expression, disbelief expression, checking expression

\textbf{Adjusted Expression (JA)}:\\
そのようなことはありますか？\\
\textbf{Adjusted Expression (EN)}:\\
Is that really possible?

\textbf{Example (JA)}:\\
A「(昨日:きのう)、(一日:いちにち)で(三回:さんかい)も(車:くるま)が(故障:こしょう)したの！そんなことある？」\\
\textbf{Example (EN)}:\\
A: My car broke down three times in one day yesterday! Can that really happen?

\textbf{Notes (JA)}:\\
「そんなことある？」は、予期しないことが起こったことに対して、驚きや信じられない気持ちを表現しています。もちろん、この表現を使う場面ですべて驚きを表現するために使われるわけではありません。その可能性があるかどうかを訊く場面でも使えます。驚きや信じられないを示す場合には少し声を大きく言うとそういう意味になるのはおそらくどの言語でも同じかもしれません。

\textbf{Notes (EN)}:\\
The phrase "sonna koto aru?" is used to express surprise or disbelief when something unexpected happens. It can also be used to ask whether something is actually possible. In conversation, raising your voice slightly can add a sense of surprise or disbelief.

\newpage

\section*{106. It all comes down to that}

\textbf{Date}: 2024.07.25 (Thu)\\
\textbf{Index}: tsukiru\\
\textbf{Expression (JA)}: それに(尽:つ)きるよね\\
\textbf{Romanization}: Sore ni tsukiru yo ne\\
\textbf{Expression (EN)}: It all comes down to that

\textbf{Variants (JA)}: それに尽きるよね / それに尽きる / 結局それだよね\\
\textbf{Variants (EN)}: It all comes down to that / That's what it all comes down to, right? / That's really what it comes down to

\textbf{Intent (JA)}: 要点のまとめ, 評価の共有, 同意の誘い\\
\textbf{Intent (EN)}: summarizing a point, sharing an evaluation, seeking agreement

\textbf{Tags (JA)}: 即時文法, 口語, まとめ表現, 評価共有表現, 同意誘導表現\\
\textbf{Tags (EN)}: immediate grammar, spoken, summary expression, evaluation-sharing expression, agreement-seeking expression

\textbf{Adjusted Expression (JA)}:\\
それに尽きます\\
\textbf{Adjusted Expression (EN)}:\\
That is what it comes down to.

\textbf{Example (JA)}:\\
A「(成功:せいこう)するためには、(準備:じゅんび)と(タイミング:たいみんぐ)、それに(尽:つ)きるよね」\\
\textbf{Example (EN)}:\\
A: To succeed, it all comes down to preparation and timing, right?

\textbf{Notes (JA)}:\\
「それに尽きるよね」は、すべてが一つのことにかかっていることを表現するために使用されます。「尽きる」は、ここでは物事が最終的に一つの要点に帰着することを表します。また、「よね」が付くことで、その評価を共有し、相手の同意を求める働きがあります。

\textbf{Notes (EN)}:\\
The phrase "sore ni tsukiru yo ne" is used to express that everything can be reduced to one key point. "Tsukiru" here means that things come down to or boil down to a single factor. The ending "yo ne" adds a sense of shared understanding and invites agreement.

\newpage

\section*{105. Full of energy}

\textbf{Date}: 2024.07.24 (Wed)\\
\textbf{Index}: ikiiki\\
\textbf{Expression (JA)}: (生:い)き(生:い)きしてる\\
\textbf{Romanization}: Ikiiki shiteru\\
\textbf{Expression (EN)}: Full of energy

\textbf{Variants (JA)}: 生き生きしてる / 生き生きしている / 生き生きとしている\\
\textbf{Variants (EN)}: Full of energy / Lively / Full of life

\textbf{Intent (JA)}: 様子の描写, 肯定的評価, 口語的な表現\\
\textbf{Intent (EN)}: describing someone's state, positive evaluation, colloquial expression

\textbf{Tags (JA)}: 即時文法, 口語, 様子描写表現, 肯定評価表現, 口語的表現\\
\textbf{Tags (EN)}: immediate grammar, spoken, state-description expression, positive-evaluation expression, colloquial expression

\textbf{Adjusted Expression (JA)}:\\
(生:い)き(生:い)きしています\\
\textbf{Adjusted Expression (EN)}:\\
She is full of energy.

\textbf{Example (JA)}:\\
A「なんか、(生:い)き(生:い)きしてる」 B「(好:す)きな(仕事:しごと)をしてるからかな？」\\
\textbf{Example (EN)}:\\
A: She looks lively. B: Maybe it's because she's doing work she loves?

\textbf{Notes (JA)}:\\
「生き生きしてる」は、誰かが元気いっぱいであることを表現するために使用されます。

\textbf{Notes (EN)}:\\
The phrase "ikiiki shiteru" is used to express that someone is full of energy and lively.

\newpage

\section*{104. Exhausted}

\textbf{Date}: 2024.07.23 (Tue)\\
\textbf{Index}: kutakuta\\
\textbf{Expression (JA)}: くたくた\\
\textbf{Romanization}: Kutakuta\\
\textbf{Expression (EN)}: Exhausted

\textbf{Variants (JA)}: くたくた / くたくたになる\\
\textbf{Variants (EN)}: Exhausted / I'm exhausted

\textbf{Intent (JA)}: 状態の表現, 疲労の共有, 口語的な表現\\
\textbf{Intent (EN)}: expressing a state, sharing fatigue, colloquial expression

\textbf{Tags (JA)}: 即時文法, 口語, 状態表現, 疲労表現, 口語的表現\\
\textbf{Tags (EN)}: immediate grammar, spoken, state expression, fatigue expression, colloquial expression

\textbf{Adjusted Expression (JA)}:\\
くたくたです\\
\textbf{Adjusted Expression (EN)}:\\
I'm exhausted.

\textbf{Example (JA)}:\\
A「(今日:きょう)は、(仕事:しごと)で、くたくた」\\
\textbf{Example (EN)}:\\
A: I'm exhausted today from work.

\textbf{Notes (JA)}:\\
「くたくた」は、疲れていることを表現するために使用されます。体力的にも精神的にも疲れているときに使われます。

\textbf{Notes (EN)}:\\
The phrase "kutakuta" is used to express that you are exhausted or very tired. It can refer to both physical and mental fatigue.

\newpage

\section*{103. No matter what.}

\textbf{Date}: 2024.07.22 (Mon)\\
\textbf{Index}: doushitemo\\
\textbf{Expression (JA)}: どうしても\\
\textbf{Romanization}: Doushitemo.\\
\textbf{Expression (EN)}: No matter what.

\textbf{Variants (JA)}: どうしても / どうしても〜したい / どうしても〜したくて\\
\textbf{Variants (EN)}: No matter what. / I really want to... / I just have to...

\textbf{Intent (JA)}: 決意, 強調, 願望の強調\\
\textbf{Intent (EN)}: determination, emphasis, strong desire

\textbf{Tags (JA)}: 即時文法, 口語, 強調表現, 情熱表現\\
\textbf{Tags (EN)}: immediate grammar, spoken, expression of emphasis, expression of strong feeling

\textbf{Adjusted Expression (JA)}:\\
どうしても、お(会:あ)いしたいと(思:おも)います。\\
\textbf{Adjusted Expression (EN)}:\\
I really would like to meet you.

\textbf{Example (JA)}:\\
どうしても、(今晩:こんばん)は(君:きみ)に(会:あ)いたい。\\
\textbf{Example (EN)}:\\
I want to see you tonight no matter what.

\textbf{Notes (JA)}:\\
「どうしても」は、どんな理由があっても何かをしたいという強い決意を表す表現です。「何が何でも」に近い意味を持つ副詞です。テレビドラマでもよく使われますが、実際に使うとどうなるでしょうか。大成功になるか、あるいは大失敗になるかでしょうか。

\textbf{Notes (EN)}:\\
The phrase "doushitemo" expresses a strong determination to do something regardless of circumstances. It is an adverb similar to "no matter what" or "by all means." You may hear it often in TV dramas, but if you actually use it in real life, the result could be either a great success or a complete failure.

\newpage

\section*{102. That sounds tough.}

\textbf{Date}: 2024.07.21 (Sun)\\
\textbf{Index}: sorehakitsuidesune\\
\textbf{Expression (JA)}: それはきついですね。\\
\textbf{Romanization}: Sore wa kitsui desu ne.\\
\textbf{Expression (EN)}: That sounds tough.

\textbf{Variants (JA)}: それはきついですね。 / それはきついね。 / きついですね。\\
\textbf{Variants (EN)}: That sounds tough. / That's tough. / That must be tough.

\textbf{Intent (JA)}: 共感, 評価, 相手への配慮\\
\textbf{Intent (EN)}: showing empathy, evaluating a situation, showing consideration

\textbf{Tags (JA)}: 即時文法, 口語, 共感表現\\
\textbf{Tags (EN)}: immediate grammar, spoken, expression of empathy

\textbf{Adjusted Expression (JA)}:\\
それは大変ですね。\\
\textbf{Adjusted Expression (EN)}:\\
That sounds difficult.

\textbf{Example (JA)}:\\
A「(試験:しけん)が(近:ちか)づいていて、(毎日:まいにち)(勉強:べんきょう)(漬:づ)けです」 B「それはきついですね」\\
\textbf{Example (EN)}:\\
A: The exam is approaching, and I'm studying every day. B: That sounds tough.

\textbf{Notes (JA)}:\\
「それはきついですね」は、相手の状況が大変であることを受けて、共感や評価を表す表現です。「きつい」は、身体的・精神的に苦しい状況を表すイ形容詞で、「きつい仕事」「きつい状況」のように使われます。

\textbf{Notes (EN)}:\\
The phrase "sore wa kitsui desu ne" is used to express that a situation sounds difficult or tough, often showing empathy for the other person. The adjective "kitsui" describes something physically or mentally demanding. It can also be used before a noun, such as "kitsui shigoto" (a tough job) or "kitsui joukyou" (a difficult situation).

\newpage

\section*{101. It depends.}

\textbf{Date}: 2024.07.20 (Sat)\\
\textbf{Index}: baainiyoru\\
\textbf{Expression (JA)}: (場合:ばあい)による\\
\textbf{Romanization}: Baai ni yoru.\\
\textbf{Expression (EN)}: It depends.

\textbf{Variants (JA)}: 場合による / 場合によります / ケースバイケース\\
\textbf{Variants (EN)}: It depends. / That depends. / Case by case.

\textbf{Intent (JA)}: 一般化の回避, 条件の提示, あいまいな応答\\
\textbf{Intent (EN)}: avoiding generalization, indicating conditions, vague response

\textbf{Tags (JA)}: 即時文法, 口語, 条件表現, 応答表現\\
\textbf{Tags (EN)}: immediate grammar, spoken, conditional expression, response expression

\textbf{Adjusted Expression (JA)}:\\
(場合:ばあい)によります。\\
\textbf{Adjusted Expression (EN)}:\\
It depends on the situation.

\textbf{Example (JA)}:\\
A「いつでも、そうなりますか」 B「いや、(場合:ばあい)による」\\
\textbf{Example (EN)}:\\
A: Does it always happen? B: No, it depends.

\textbf{Notes (JA)}:\\
「場合による」は、状況によって結果が変わることを表す表現です。「場合によります」と言うこともできますが、一般的な解釈について述べることが多いため、必ずしも丁寧に言う必要はありません。実際の会話では、状況に応じて、ですます体とだ体が使い分けられます。

\textbf{Notes (EN)}:\\
The phrase "baai ni yoru" is used to express that something changes depending on the situation. It can also be said as "baai ni yorimasu," but since it often refers to a general interpretation, it does not always need to be expressed politely. In actual conversation, speakers switch between plain and polite forms depending on the situation.

\newpage

\section*{100. I'm back.}

\textbf{Date}: 2024.07.19 (Fri)\\
\textbf{Index}: tadaima\\
\textbf{Expression (JA)}: ただいま\\
\textbf{Romanization}: Tadaima.\\
\textbf{Expression (EN)}: I'm back.

\textbf{Variants (JA)}: ただいま / ただいまー\\
\textbf{Variants (EN)}: I'm back. / I'm home.

\textbf{Intent (JA)}: 帰宅のあいさつ, 帰着の報告\\
\textbf{Intent (EN)}: returning greeting, announcing one's return

\textbf{Tags (JA)}: 即時文法, 口語, あいさつ表現, 多義的表現\\
\textbf{Tags (EN)}: immediate grammar, spoken, greeting expression, polysemous expression

\textbf{Adjusted Expression (JA)}:\\
ただいま戻りました。\\
\textbf{Adjusted Expression (EN)}:\\
I have just returned.

\textbf{Example (JA)}:\\
A「ただいま」 B「おかえり」\\
\textbf{Example (EN)}:\\
A: I'm back. B: Welcome back.

\textbf{Notes (JA)}:\\
「ただいま」は、家や職場に帰ってきたときに使われます。ただし、「ただいま」には2つの意味があります。ひとつはこの例のような帰宅時のあいさつで、もうひとつは「今すぐ...します」のような意味です。イントネーションが少し違い、後者の意味のときには「だ」の音を少し高めに発音します。

\textbf{Notes (EN)}:\\
The phrase "tadaima" is used when you return home or to your workplace. However, "tadaima" has two meanings. One is the greeting in the example above, and the other is "right now" or "at this moment," as in "I will do it right now." The intonation is slightly different. In the latter meaning, the "da" is pronounced a little higher.

\newpage

\section*{99. I'm hooked.}

\textbf{Date}: 2024.07.18 (Thu)\\
\textbf{Index}: hamacchatta\\
\textbf{Expression (JA)}: はまっちゃった\\
\textbf{Romanization}: Hamacchatta.\\
\textbf{Expression (EN)}: I'm hooked.

\textbf{Variants (JA)}: はまっちゃった / 最近、はまっちゃった / すっかりはまっちゃった\\
\textbf{Variants (EN)}: I'm hooked. / I've gotten hooked on it. / I'm really into it.

\textbf{Intent (JA)}: 夢中の表現, 近況の共有, 感情の表現\\
\textbf{Intent (EN)}: expressing enthusiasm, sharing a recent state, expressing a feeling

\textbf{Tags (JA)}: 即時文法, 口語, 夢中表現\\
\textbf{Tags (EN)}: immediate grammar, spoken, expression of being into something

\textbf{Adjusted Expression (JA)}:\\
最近、それに夢中になっています。\\
\textbf{Adjusted Expression (EN)}:\\
I've recently become really interested in it.

\textbf{Example (JA)}:\\
A「(最近:さいきん)、(漫画:まんが)にはまっちゃった」\\
\textbf{Example (EN)}:\\
A: I've gotten hooked on manga recently.

\textbf{Notes (JA)}:\\
「はまっちゃった」は、趣味やゲーム、テレビ番組など、何かに夢中になってしまったことを表す表現です。くだけた、いきいきした感じがあります。

\textbf{Notes (EN)}:\\
The phrase "hamacchatta" is used to express that you have become really into something, such as a hobby, game, or TV show. It often has a casual and lively feeling.

\newpage

\section*{98. Take care}

\textbf{Date}: 2024.07.17 (Wed)\\
\textbf{Index}: kiotsukete\\
\textbf{Expression (JA)}: (気:き)をつけて\\
\textbf{Romanization}: Ki o tsukete.\\
\textbf{Expression (EN)}: Take care.

\textbf{Variants (JA)}: 気をつけて / 気をつけてね / 帰り道、気をつけてね\\
\textbf{Variants (EN)}: Take care. / Be careful. / Take care on your way home.

\textbf{Intent (JA)}: 注意の促し, 気づかい, 別れ際の表現\\
\textbf{Intent (EN)}: urging caution, showing care, parting expression

\textbf{Tags (JA)}: 即時文法, 口語, 気づかい表現\\
\textbf{Tags (EN)}: immediate grammar, spoken, expression of care

\textbf{Adjusted Expression (JA)}:\\
お(気:き)をつけてお(帰:かえ)りください。\\
\textbf{Adjusted Expression (EN)}:\\
Please take care on your way home.

\textbf{Example (JA)}:\\
A「(帰:かえ)り(道:みち)、(気:き)をつけてね」\\
\textbf{Example (EN)}:\\
A: Take care on your way home.

\textbf{Notes (JA)}:\\
「気をつけて」は、相手に注意してほしいという気持ちを表す表現です。相手を気づかう気持ちを含み、特に別れ際や、何か気をつけるべき場面でよく使われます。

\textbf{Notes (EN)}:\\
The phrase "ki o tsukete" is used to tell someone to be careful. It often expresses concern or care for the other person, especially when they are leaving or about to do something that may involve some risk.

\newpage

\section*{97. I see.}

\textbf{Date}: 2024.07.16 (Tue)\\
\textbf{Index}: naruhodo\\
\textbf{Expression (JA)}: なるほど\\
\textbf{Romanization}: Naruhodo.\\
\textbf{Expression (EN)}: I see.

\textbf{Variants (JA)}: なるほど / なるほどね / なるほどなるほど\\
\textbf{Variants (EN)}: I see. / That makes sense. / I see, I see.

\textbf{Intent (JA)}: 理解の表明, 納得, 会話の促進\\
\textbf{Intent (EN)}: showing understanding, showing agreement, promoting conversation

\textbf{Tags (JA)}: 即時文法, 口語, 相づち表現, 会話ストラテジー\\
\textbf{Tags (EN)}: immediate grammar, spoken, backchannel expression, conversation strategy

\textbf{Adjusted Expression (JA)}:\\
なるほど、そういうことですね。\\
\textbf{Adjusted Expression (EN)}:\\
I see, that's what you mean.

\textbf{Example (JA)}:\\
A「(最近:さいきん)、(趣味:しゅみ)で(絵:え)を(描:か)き(始:はじ)めたんだ」 B「なるほど、それでいつも(楽:たの)しそうなんだ」\\
\textbf{Example (EN)}:\\
A: I recently started drawing pictures as a hobby. B: I see, that's why you always look happy.

\textbf{Notes (JA)}:\\
「なるほど」は、相手の説明を理解したことを表現するために使用されます。相手の具体的な説明を聞いた後に使います。短い会話でも自然に使える表現です。ネイティブスピーカーの日常会話でよく使われるフレーズですが、この表現を使ったからといって、本当に理解しているかどうかはわかりません。私のかつての学生は、この表現を使っているのに、実際には理解していないことがありました。しかし、それが無意味であるかどうかは、言い切れません。十分に理解できていなくても、相手の発言を促すために、「なるほどなるほど」と言っていたかもしれないからです。スムーズな会話を促すためのストラテジーとして使っていたのですね。

\textbf{Notes (EN)}:\\
The phrase "naruhodo" is used to express that you understand the other person's explanation. It is used after hearing a specific explanation from the other person. It is a natural expression that can be used even in short conversations. However, just because someone uses this expression does not necessarily mean that they really understand. Some of my former students used this expression even when they did not actually understand. Still, that does not mean it was meaningless. Even without full understanding, they may have been using "naruhodo naruhodo" as a strategy to encourage the other person to continue speaking. In that sense, it may have functioned as a strategy for promoting smooth conversation.

\newpage

\section*{96. I knew it.}

\textbf{Date}: 2024.07.15 (Mon)\\
\textbf{Index}: yappari\\
\textbf{Expression (JA)}: やっぱり\\
\textbf{Romanization}: Yappari.\\
\textbf{Expression (EN)}: I knew it.

\textbf{Variants (JA)}: やっぱり / やっぱりね / やっぱりそうか\\
\textbf{Variants (EN)}: I knew it. / See? / Just as I thought.

\textbf{Intent (JA)}: 予想の確認, 納得, 軽い共感\\
\textbf{Intent (EN)}: confirming an expectation, showing realization, light agreement

\textbf{Tags (JA)}: 即時文法, 口語, 反応表現\\
\textbf{Tags (EN)}: immediate grammar, spoken, reaction expression

\textbf{Adjusted Expression (JA)}:\\
やはりそうですね。\\
\textbf{Adjusted Expression (EN)}:\\
That is as expected.

\textbf{Example (JA)}:\\
A「(雨:あめ)が(降:ふ)ってきた」 B「やっぱり」\\
\textbf{Example (EN)}:\\
A: It's raining. B: I knew it.

\textbf{Notes (JA)}:\\
「やっぱり」は、結果が予想通りであるときに使われ、予想していたことが実際に起こったときの納得や確認の気持ちを表します。

\textbf{Notes (EN)}:\\
The phrase "yappari" is used when something turns out as expected. It often expresses the speaker's feeling of confirmation or realization when what they anticipated actually happens.

\newpage

\section*{95. That's not true.}

\textbf{Date}: 2024.07.14 (Sun)\\
\textbf{Index}: sonnakotonai\\
\textbf{Expression (JA)}: そんなことない\\
\textbf{Romanization}: Sonna koto nai.\\
\textbf{Expression (EN)}: That's not true.

\textbf{Variants (JA)}: そんなことない / そんなことないよ / そんなことないって\\
\textbf{Variants (EN)}: That's not true. / No, that's not true. / That's not the case.

\textbf{Intent (JA)}: 否定, 励まし, 評価の打ち消し\\
\textbf{Intent (EN)}: denial, encouragement, rejecting an evaluation

\textbf{Tags (JA)}: 即時文法, 口語, 否定表現, 励まし表現\\
\textbf{Tags (EN)}: immediate grammar, spoken, expression of denial, expression of encouragement

\textbf{Adjusted Expression (JA)}:\\
そんなことはありません。\\
\textbf{Adjusted Expression (EN)}:\\
That's not the case.

\textbf{Example (JA)}:\\
A「(今日:きょう)のプレゼン、うまくいかなかった(気:き)がする。」 B「そんなことない、ちゃんと(伝:つた)わってたよ」\\
\textbf{Example (EN)}:\\
A: I don't think today's presentation went well. B: That's not true, it came across well.

\textbf{Notes (JA)}:\\
「そんなことない」は、相手の言ったことを否定し、特に相手が自分を控えめに評価したときに、それを打ち消して励ますときによく使われます。

\textbf{Notes (EN)}:\\
The phrase "sonna koto nai" is used to deny what the other person has said, often to encourage them. It is commonly used when someone downplays their own performance or ability, and the speaker responds by rejecting that negative evaluation.

\newpage

\section*{94. What the heck!}

\textbf{Date}: 2024.07.13 (Sat)\\
\textbf{Index}: nantekotta\\
\textbf{Expression (JA)}: なんてこった。\\
\textbf{Romanization}: Nante kotta.\\
\textbf{Expression (EN)}: What the heck!

\textbf{Variants (JA)}: なんてこった。 / なんてこったよ。 / なんてことだ。\\
\textbf{Variants (EN)}: What the heck! / What on earth! / What the hell!

\textbf{Intent (JA)}: 驚きの表現, 困惑の表現, 感嘆\\
\textbf{Intent (EN)}: expressing surprise, expressing confusion, exclamation

\textbf{Tags (JA)}: 即時文法, 口語, 感嘆表現\\
\textbf{Tags (EN)}: immediate grammar, spoken, exclamatory expression

\textbf{Adjusted Expression (JA)}:\\
それは大変ですね。\\
\textbf{Adjusted Expression (EN)}:\\
That sounds serious.

\textbf{Example (JA)}:\\
A「(昨日:きのう)、(急:きゅう)に(パソコン:ぱそこん)が(壊:こわ)れちゃって、(データ:でーた)が(全部:ぜんぶ)(消:き)えたんだ」 B「なんてこった！(修理:しゅうり)できるの？」\\
\textbf{Example (EN)}:\\
A: Yesterday, my computer suddenly broke down and all the data disappeared. B: What the heck! Can it be repaired?

\textbf{Notes (JA)}:\\
「なんてこった」は、驚きや困惑を表す自然な日本語の表現で、予期しない出来事やトラブルに対して使われることが多いです。良い意味の驚きにも使えますが、やや困った状況で使われることが多い表現です。

\textbf{Notes (EN)}:\\
The phrase "nante kotta" expresses surprise or confusion. It is used when something unexpected happens, especially in troublesome situations. It can sometimes be used for positive surprises as well, but it is more commonly associated with trouble or shock.

\newpage

\section*{93. Sort of...}

\textbf{Date}: 2024.07.12 (Fri)\\
\textbf{Index}: maanantonaku\\
\textbf{Expression (JA)}: まあ、なんとなく。\\
\textbf{Romanization}: Maa nantonaku.\\
\textbf{Expression (EN)}: Well, sort of.

\textbf{Variants (JA)}: まあ、なんとなく。 / なんとなく。 / まあ、なんかそんな感じ。\\
\textbf{Variants (EN)}: Well, sort of. / Sort of. / Well, kind of.

\textbf{Intent (JA)}: 理由の回避, あいまいな応答, 感覚の表現\\
\textbf{Intent (EN)}: avoiding giving a reason, vague response, expressing a feeling

\textbf{Tags (JA)}: 即時文法, 口語, あいまい表現, 応答表現\\
\textbf{Tags (EN)}: immediate grammar, spoken, vague expression, response expression

\textbf{Adjusted Expression (JA)}:\\
(特:とく)に(理由:りゆう)はないんですが、なんとなくそう(思:おも)います。\\
\textbf{Adjusted Expression (EN)}:\\
I don't have a clear reason, but I kind of feel that way.

\textbf{Example (JA)}:\\
A「(何:なに)か(食:た)べたいな」 B「じゃあ、(何:なに)か(食:た)べる？」 A「カレーライス？」 B「なぜ？」 A「まあ、なんとなく」\\
\textbf{Example (EN)}:\\
A: I want to eat something. B: Then, do you want to eat? A: Curry rice? B: Why? A: Well, sort of.

\textbf{Notes (JA)}:\\
「まあ、なんとなく」は、それをする理由が明確でないときに使えます。「どうして」と聞かれても理由が明確に言えないことはよくあります。「そんな感じがしました」というよりも、一言で「なんとなく」と言えると自然です。他にも「まあ、なんとなく分かります。」「まあ、なんとなくそう思います。」「まあ、なんとなくそんな気がする。」のような使い方があります。

\textbf{Notes (EN)}:\\
The phrase "maa nantonaku" is used when you do not have a clear reason for something. It is often used when you cannot explain why even if you are asked. It is more natural to say "nantonaku" briefly than to give a full explanation like "I feel that way." Other expressions include "maa, nantonaku wakarimasu" (Well, I kind of understand), "maa, nantonaku sou omoimasu" (Well, I kind of think so), and "maa, nantonaku sonna ki ga suru" (Well, I kind of feel that way).

\newpage

\section*{92. Crowded trains are tough.}

\textbf{Date}: 2024.07.11 (Thu)\\
\textbf{Index}: mannindensha\\
\textbf{Expression (JA)}: (満員電車:まんいんでんしゃ)は(辛:つら)い\\
\textbf{Romanization}: Man'in densha wa tsurai.\\
\textbf{Expression (EN)}: Crowded trains are tough.

\textbf{Variants (JA)}: 満員電車は辛い / 満員電車、辛いよね / 満員電車は本当に辛い\\
\textbf{Variants (EN)}: Crowded trains are tough. / Crowded trains are really tough. / Rush hour trains are tough.

\textbf{Intent (JA)}: 状況の評価, 共感の共有, 不満の表現\\
\textbf{Intent (EN)}: expressing evaluation, sharing empathy, expressing discomfort

\textbf{Tags (JA)}: 即時文法, 口語, 評価表現, 不満表現\\
\textbf{Tags (EN)}: immediate grammar, spoken, expression of evaluation, expression of discomfort

\textbf{Adjusted Expression (JA)}:\\
満員電車は大変ですね。\\
\textbf{Adjusted Expression (EN)}:\\
Crowded trains are difficult, aren't they?

\textbf{Example (JA)}:\\
A「(満員電車:まんいんでんしゃ)は(辛:つら)い」\\
\textbf{Example (EN)}:\\
A: Crowded trains are tough.

\textbf{Notes (JA)}:\\
「満員電車」は、「満員」と「電車」の2つの語を組み合わせた複合語で、混雑した電車を指します。「辛い」は、身体的・精神的に苦しい状況を表すイ形容詞で、「辛い仕事」「辛い時間」のように、耐えるのが大変な状況に使われます。

\textbf{Notes (EN)}:\\
The expression "man'in densha" is a compound made from "man'in" (full) and "densha" (train), referring to a crowded train. The adjective "tsurai" expresses a physically or emotionally difficult situation. It is often used to describe something that is hard to endure, such as "tsurai shigoto" (a tough job) or "tsurai jikan" (a difficult time).

\newpage

\section*{91. I want to relax.}

\textbf{Date}: 2024.07.10 (Wed)\\
\textbf{Index}: nonbirishitaine\\
\textbf{Expression (JA)}: のんびりしたいね。\\
\textbf{Romanization}: Nonbirishitai ne.\\
\textbf{Expression (EN)}: I want to relax.

\textbf{Variants (JA)}: のんびりしたいね。 / 今日はのんびりしたいね。 / 何もしないでのんびりしたいね。\\
\textbf{Variants (EN)}: I want to relax. / I want to take it easy today. / I just want to relax and do nothing today.

\textbf{Intent (JA)}: 気持ちの共有, 願望の共有, くつろぎの表現\\
\textbf{Intent (EN)}: sharing a feeling, sharing a wish, expression of relaxation

\textbf{Tags (JA)}: 即時文法, 口語, 願望表現, くつろぎ表現\\
\textbf{Tags (EN)}: immediate grammar, spoken, expression of desire, expression of relaxation

\textbf{Adjusted Expression (JA)}:\\
(少:すこ)しゆっくりしたいですね。\\
\textbf{Adjusted Expression (EN)}:\\
I'd like to relax a little.

\textbf{Example (JA)}:\\
A「(今日:きょう)は、(何:なに)もしないで、のんびりしたいね」\\
\textbf{Example (EN)}:\\
A: I just want to relax and do nothing today.

\textbf{Notes (JA)}:\\
「のんびり」は、動詞の前に置いて、ゆっくり、リラックスする様子を表す副詞です。「のんびりする」「のんびりしたい」「のんびりしよう」「のんびり暮らす」のように使えます。

\textbf{Notes (EN)}:\\
The word "nonbiri" is an adverb that expresses a relaxed and leisurely state. It is used before a verb to describe doing something in a relaxed way. It can be used in expressions such as "nonbiri suru (to relax)," "nonbiri shitai (want to relax)," "nonbiri shiyou (let's relax)," and "nonbiri kurasu (to live a leisurely life)."

\newpage

\section*{90. Don't push yourself}

\textbf{Date}: 2024.07.09 (Tue)\\
\textbf{Index}: muriwashinai\\
\textbf{Expression (JA)}: (無理:むり)はしない\\
\textbf{Romanization}: Muri wa shinai\\
\textbf{Expression (EN)}: Don't push yourself

\textbf{Variants (JA)}: 無理はしない / 無理しない / 無理はしないで\\
\textbf{Variants (EN)}: Don't push yourself / Don't overdo it / Take it easy

\textbf{Intent (JA)}: 相手を気づかう表現, 無理を止める表現, 体調や状況への配慮\\
\textbf{Intent (EN)}: showing concern for someone, telling someone not to overdo it, showing consideration for someone's condition or situation

\textbf{Tags (JA)}: 即時文法, 口語, 配慮表現, 助言, 健康, 気づかい\\
\textbf{Tags (EN)}: immediate grammar, colloquial, considerate expression, advice, health, concern

\textbf{Adjusted Expression (JA)}:\\
(体調:たいちょう)が(悪:わる)いのだから、(無理:むり)をしないほうがよい。\\
\textbf{Adjusted Expression (EN)}:\\
Since you are not feeling well, you should not overexert yourself.

\textbf{Example (JA)}:\\
A「(明日:あした)、(試験:しけん)があるけれども、(熱:ねつ)があるんだから、(無理:むり)はしない」\\
\textbf{Example (EN)}:\\
A: You have an exam tomorrow, but since you have a fever, don't push yourself.

\textbf{Notes (JA)}:\\
「無理」は、できないことをやろうとすることです。病気のときに、試験のために治療をしないで試験勉強をしようとすることや、お金を持っていないのに、女の子にごちそうすることなどです。

\textbf{Notes (EN)}:\\
The word "muri" refers to situations where you try to do something beyond your capacity. For example, studying for an exam without taking proper rest when you are sick, or trying to treat a girl when you do not have enough money.

\newpage

\section*{89. The first impression}

\textbf{Date}: 2024.07.08 (Mon)\\
\textbf{Index}: daiichiinshou\\
\textbf{Expression (JA)}: (第一印象:だいいちいんしょう)\\
\textbf{Romanization}: Daiichiinshou\\
\textbf{Expression (EN)}: The first impression

\textbf{Variants (JA)}: 第一印象 / 第一印象がいい / 第一印象では\\
\textbf{Variants (EN)}: the first impression / make a good first impression / at first impression

\textbf{Intent (JA)}: 初対面の印象の表現, 最初の評価の共有, 経験の出だしの記述\\
\textbf{Intent (EN)}: expressing an initial impression, sharing a first evaluation, describing the beginning of an experience

\textbf{Tags (JA)}: 即時文法, 名詞, 複合語, 評価表現, 経験の記述\\
\textbf{Tags (EN)}: immediate grammar, noun, compound word, evaluative expression, describing experience

\textbf{Adjusted Expression (JA)}:\\
(最初:さいしょ)に(受:う)けた(印象:いんしょう)は、(ゴミ箱:ごみばこ)がないということだった。\\
\textbf{Adjusted Expression (EN)}:\\
My initial impression was that there were no trash cans.

\textbf{Example (JA)}:\\
A「(日本:にほん)はどうですか」 B「(ゴミ箱:ごみばこ)がないというのが(第一印象:だいいちいんしょう)ですね」\\
\textbf{Example (EN)}:\\
A: How is Japan? B: My first impression was that there are no trash cans.

\textbf{Notes (JA)}:\\
「第一印象」は、「第一」と「印象」の2つの語を組み合わせて作られた複合語です。この場合は「第一の印象」と言ってもよいのですが、すべての2語の組合せが複合語として使われるわけではありません。「第二印象」ということばはなく、「二番目の印象」と言わなければなりません。文法は便利かもしれませんが、ルールがあるからといって、すべての場合に適用できるわけではありませんので、実際に使われたことばを使うのがよいでしょう。このことから、できるだけ多く日本語に触れることが大切です。

\textbf{Notes (EN)}:\\
The phrase "daiichiinshou" is a compound word formed from "daiichi" and "inshou." In this case, you can also say "daiichi no inshou," but not all combinations of words are used as compounds. For example, there is no commonly used word like "dainiinshou" ("second impression"), so you need to say "nibanme no inshou." Grammar can be helpful, but rules do not apply in all cases, so it is better to use expressions that are actually used.

\newpage

\section*{88. Don't mix them, it's dangerous!}

\textbf{Date}: 2024.07.07 (Sun)\\
\textbf{Index}: mazerunakiken\\
\textbf{Expression (JA)}: まぜるな、(危険:きけん)！\\
\textbf{Romanization}: Mazeruna, kiken!\\
\textbf{Expression (EN)}: Don't mix them, it's dangerous!

\textbf{Variants (JA)}: まぜるな危険 / 混ぜるな危険 / それ、まぜるな危険じゃん\\
\textbf{Variants (EN)}: Don't mix them, it's dangerous! / What a dangerous combination! / Don't mix those!

\textbf{Intent (JA)}: 危険の警告, 注意喚起, 冗談としての強調表現\\
\textbf{Intent (EN)}: warning of danger, calling attention, emphatic expression used humorously

\textbf{Tags (JA)}: 即時文法, 定型表現, 口語, 警告表現, 注意喚起, ユーモア\\
\textbf{Tags (EN)}: immediate grammar, formulaic expression, colloquial, warning expression, caution, humor

\textbf{Adjusted Expression (JA)}:\\
(それら)を(混合:こんごう)すると、(危険:きけん)である。\\
\textbf{Adjusted Expression (EN)}:\\
Mixing them is dangerous.

\textbf{Example (JA)}:\\
A「(漂白剤:ひょうはくざい)と(洗剤:せんざい)を(混:ま)ぜたら、(危険:きけん)ですよ」B「えっ、それ、まぜるな、(危険:きけん)じゃん」\\
\textbf{Example (EN)}:\\
A: If you mix bleach and detergent, it's dangerous. B: What? That's a dangerous combination! Don't mix them!

\textbf{Notes (JA)}:\\
「まぜるな、危険！」は、2つ以上のものを混ぜると危険であることを表現するために使用されます。混ぜてはいけないものを混ぜないように誰かに警告するときに使用されます。塩素系漂白剤のボトルには「まぜるな危険」と書かれていることがあります。また、このことばを利用して、2つのものが非常に特殊で危険な組み合わせであるときには、「何！その組み合わせ、まぜるな、危険じゃん」と言って笑い話にすることもあります。

\textbf{Notes (EN)}:\\
The phrase "mazeruna, kiken!" is used to warn that mixing certain things is dangerous. It is commonly used when telling someone not to combine things that should not be mixed. Bottles of chlorine bleach often have warnings like "Do not mix" on them. This phrase is also used humorously, for example, when two things make a striking or risky combination: "What? That combination is dangerous! Don't mix them!"

\newpage

\section*{87. It's muggy}

\textbf{Date}: 2024.07.06 (Sat)\\
\textbf{Index}: mushiatsui\\
\textbf{Expression (JA)}: むし(暑:あつ)い\\
\textbf{Romanization}: Mushiatsui\\
\textbf{Expression (EN)}: It's muggy

\textbf{Variants (JA)}: むし暑い / 蒸し暑い / 今日はむし暑いね\\
\textbf{Variants (EN)}: It's muggy. / It's hot and humid. / It's really muggy today.

\textbf{Intent (JA)}: 天気や体感の共有, 不快な暑さの表現, 季節の会話\\
\textbf{Intent (EN)}: sharing the weather or physical feeling, expressing uncomfortable heat, seasonal conversation

\textbf{Tags (JA)}: 即時文法, 口語, 天気表現, 体感表現, 複合語, 夏\\
\textbf{Tags (EN)}: immediate grammar, colloquial, weather expression, physical sensation expression, compound word, summer

\textbf{Adjusted Expression (JA)}:\\
(今日は)、(気温:きおん)が(高:たか)く、(湿度:しつど)も(高:たか)いですね。\\
\textbf{Adjusted Expression (EN)}:\\
It is hot and humid today.

\textbf{Example (JA)}:\\
A「むし(暑:あつ)いね」 B「うん、(暑:あつ)いね」\\
\textbf{Example (EN)}:\\
A: It's muggy. B: Yeah, it's hot.

\textbf{Notes (JA)}:\\
「蒸し暑い」は、暑くて湿度が高い日本でよく言われることばです。この表現は、「蒸し」と「暑い」の2語の組み合わせで、これは複合語と呼ばれています。「蒸し」は「湿気が多い」という意味で、「暑い」は「気温が高い」という意味です。複合語は使わなくても、日常的に生きていけるので、意識して使わないとなかなか習得できません。使えると場面に合ったわかりやすい上手な言い方になります。

\textbf{Notes (EN)}:\\
The phrase "mushiatsui" is a common expression in Japan, often used on hot and humid days. This expression is a compound of "mushi" and "atsui." "Mushi" refers to humidity, and "atsui" means "hot." Since people can manage daily life without using such compound expressions, they can be difficult to master. Unless you consciously try to use them, you may not acquire them easily. However, if you can use them, your speech becomes clearer and more suitable for the situation.

\newpage

\section*{86. Let's go!}

\textbf{Date}: 2024.07.05 (Fri)\\
\textbf{Index}: ikou\\
\textbf{Expression (JA)}: (行:い)こう！\\
\textbf{Romanization}: Ikou!\\
\textbf{Expression (EN)}: Let's go!

\textbf{Variants (JA)}: (行:い)こう！ / (食:た)べに(行:い)こう！ / (飲:の)みに(行:い)こう！\\
\textbf{Variants (EN)}: Let's go! / Let's go eat! / Let's go for a drink!

\textbf{Intent (JA)}: 誘い, 働きかけ\\
\textbf{Intent (EN)}: invitation, prompting

\textbf{Tags (JA)}: 即時文法, 口語, 誘い表現\\
\textbf{Tags (EN)}: immediate grammar, spoken, invitation expression

\textbf{Adjusted Expression (JA)}:\\
(行:い)きましょう\\
\textbf{Adjusted Expression (EN)}:\\
Shall we go?

\textbf{Example (JA)}:\\
A「(食:た)べに(行:い)こう！」 B「いいね！」\\
\textbf{Example (EN)}:\\
A: Let's go eat! B: Sounds good!

\textbf{Notes (JA)}:\\
「行こう！」は、どこかに行くことを誘う表現です。実際の発音では、「いこう」の「う」はあまり強く発音されません。したがって、「いこ」と聞こえるぐらいです。また、「たべに」の「に」も強く発音されませんので、全体的に「たべいこ」と聞こえるかもしれません。同様に「のみいこ」「みにこ」「しにこ」など、きわめて簡単に発音されることが多いので、馴れておくとよいかもしれません。確かに短く発音すると言いやすいですよね。

\textbf{Notes (EN)}:\\
The phrase "ikou!" is used to invite someone to go somewhere. In actual pronunciation, the "u" in "ikou" is not pronounced very strongly. Therefore, it may sound like "iko". Also, the "ni" in "tabeni" is not pronounced strongly, so it may sound like "tabeiko". Similarly, expressions such as "nomiiko (let's go for a drink)," "miiko (let's go see it)," and "shiiko (let's go do it)" are often pronounced very simply, so it may be good to get used to it. It is certainly easy to pronounce it shortly, isn't it?

\newpage

\section*{85. Wanna go see it?}

\textbf{Date}: 2024.07.04 (Thu)\\
\textbf{Index}: miniiku\\
\textbf{Expression (JA)}: (見:み)に(行:い)く？\\
\textbf{Romanization}: Mi ni iku?\\
\textbf{Expression (EN)}: Wanna go see it?

\textbf{Variants (JA)}: (見:み)に(行:い)く？ / (映画:えいが)、(見:み)に(行:い)く？ / (見:み)に(行:い)かない？\\
\textbf{Variants (EN)}: Wanna go see it? / Wanna go see a movie? / Do you want to go see it?

\textbf{Intent (JA)}: 誘い, 予定の確認, 働きかけ\\
\textbf{Intent (EN)}: invitation, checking plans, prompting

\textbf{Tags (JA)}: 即時文法, 口語, 誘い表現, 予定確認表現\\
\textbf{Tags (EN)}: immediate grammar, spoken, invitation expression, plan-checking expression

\textbf{Adjusted Expression (JA)}:\\
(見:み)に(行:い)きませんか？\\
\textbf{Adjusted Expression (EN)}:\\
Would you like to go see it?

\textbf{Example (JA)}:\\
A「(映画:えいが)、(見:み)に(行:い)く？」 B「うん、(行:い)く」\\
\textbf{Example (EN)}:\\
A: Wanna go see a movie? B: Yeah, I do.

\textbf{Notes (JA)}:\\
「見に行く？」は、助詞の「に」を使うことで、行動の目的を表すことができます。助詞「に」の前には、目的となる名詞あるいは動詞の連用形が使われます。この表現は誘う表現でもありますし、予定を尋ねる表現でもあります。明確な誘いの表現ではないので、便利です。もし、もっと丁寧に誰かを何かを見に行くように誘いたい場合は、「見に行きませんか？」という表現を使うことができますが、断られたときは気まずいので、よほど自信がある場合でなければ使わない方がよいでしょう。「見に行く？」のような尋ね方は、相手の反応によって次の言い方を変えられます。もし「行かない」と言われたら、「そう、僕は行くよ」と答えて、誘ったのではなく、単純に予定を尋ねた表現になります。

\textbf{Notes (EN)}:\\
The phrase "mi ni iku?" uses the particle "ni" to express the purpose of the action. Before the particle "ni," a noun or the continuative form of a verb is used. This expression can work both as an invitation and as a way of asking about someone's plans. Because it is not a fully explicit invitation, it is convenient. If you want to invite someone more politely, you can say "mi ni ikimasen ka?". However, if you are refused, it may feel awkward, so it is better not to use it unless you are fairly confident. An intermediate question like "mi ni iku?" lets you change your next move depending on the other person's reaction. If the other person says, "I'm not going," you can answer, "I am," and it becomes a question about plans rather than a direct invitation.

\newpage

\section*{84. Do you want something?}

\textbf{Date}: 2024.07.03 (Wed)\\
\textbf{Index}: nanikahoshiimono\\
\textbf{Expression (JA)}: (何:なに)か(欲:ほ)しいもの、ある？\\
\textbf{Romanization}: Nanika hoshii mono, aru?\\
\textbf{Expression (EN)}: Do you want something?

\textbf{Variants (JA)}: (何:なに)か(欲:ほ)しいもの、ある？ / 何かいる？ / (欲:ほ)しいものある？\\
\textbf{Variants (EN)}: Do you want something? / Do you need something? / Is there anything you want?

\textbf{Intent (JA)}: 申し出, 気づかい, 確認\\
\textbf{Intent (EN)}: offering, consideration, confirmation

\textbf{Tags (JA)}: 即時文法, 口語, 気づかい表現, 申し出表現\\
\textbf{Tags (EN)}: immediate grammar, spoken, considerate expression, offering expression

\textbf{Adjusted Expression (JA)}:\\
(何:なに)か(欲:ほ)しいものがありますか？\\
\textbf{Adjusted Expression (EN)}:\\
Is there anything you would like?

\textbf{Example (JA)}:\\
A「(何:なに)か(欲:ほ)しいもの、ある？」 B「いえ、(大丈夫:だいじょうぶ)です」\\
\textbf{Example (EN)}:\\
A: Do you want something? B: No, I'm fine.

\textbf{Notes (JA)}:\\
「何か欲しいものある？」は、食べ物や飲み物、プレゼントなど、相手に何か欲しいものがあるか尋ねたいときに使用されます。「何が」ではなく「何か」であることに注意してください。

\textbf{Notes (EN)}:\\
The phrase "nanika hoshii mono aru?" is used when you want to ask if someone wants something, such as food, drink, or a gift. Please note that it is "nanika" not "nani ga".

\newpage

\section*{83. One thing at a time}

\textbf{Date}: 2024.07.02 (Tue)\\
\textbf{Index}: hitotsuzutsu\\
\textbf{Expression (JA)}: ひとつずつ\\
\textbf{Romanization}: Hitotsuzutsu\\
\textbf{Expression (EN)}: One thing at a time

\textbf{Variants (JA)}: ひとつずつ / (一:ひと)つずつ / (一:ひと)つ(一:ひと)つ\\
\textbf{Variants (EN)}: One thing at a time / One by one / One at a time

\textbf{Intent (JA)}: 順序立てた処理, 安心の表明, 作業方針の表明\\
\textbf{Intent (EN)}: step-by-step handling, showing reassurance, stating a way of proceeding

\textbf{Tags (JA)}: 語彙表現, 口語, 順序表現, 作業表現\\
\textbf{Tags (EN)}: lexical expression, spoken, order expression, work expression

\textbf{Adjusted Expression (JA)}:\\
ひとつずつ(確認:かくにん)しながら、(片付:かたづ)けていきます。\\
\textbf{Adjusted Expression (EN)}:\\
I will deal with them one by one while checking each one carefully.

\textbf{Example (JA)}:\\
A「(何:なに)か(手伝:てつだ)いましょうか？」 B「いえ、(大丈夫:だいじょうぶ)です。ひとつずつ、(片付:かたづ)けていきます」\\
\textbf{Example (EN)}:\\
A: Can I help you with something? B: No, I'm fine. I will clean things up one thing at a time.

\textbf{Notes (JA)}:\\
「ひとつずつ」は、物事を一つずつ行うことを表現するために使用されます。部屋の片付けや荷物の整理、仕事の進め方など、物事を一つずつするというのは基本的な意味です。ひとつずつやると言うと、「ひとつひとつ確認しながら、確実にする」という意味も含まれるでしょう。

\textbf{Notes (EN)}:\\
The phrase "hitotsuzutsu" is used to express that you are doing things one by one. Doing things one by one is a basic idea in situations such as cleaning up a room, organizing your belongings, or carrying out work. When you say you do things one by one, it may also include the meaning of checking each one and doing it carefully and reliably.

\newpage

\section*{82. After a little while}

\textbf{Date}: 2024.07.01 (Mon)\\
\textbf{Index}: ochitsuitekara\\
\textbf{Expression (JA)}: ちょっと(落:お)ち(着:つ)いてから\\
\textbf{Romanization}: Chotto ochitsuite kara\\
\textbf{Expression (EN)}: After a little while

\textbf{Variants (JA)}: ちょっと(落:お)ち(着:つ)いてから / 少し(落:お)ち(着:つ)いてから / ちょっとしてから\\
\textbf{Variants (EN)}: After a little while / After I settle down a bit / After a short while

\textbf{Intent (JA)}: 時間の猶予, やわらかい先送り, 場面調整\\
\textbf{Intent (EN)}: asking for a little time, soft postponement, adjusting the situation

\textbf{Tags (JA)}: 即時文法, 口語, 先送り表現, 場面調整表現\\
\textbf{Tags (EN)}: immediate grammar, spoken, postponement expression, interaction-adjusting expression

\textbf{Adjusted Expression (JA)}:\\
(少:すこ)し(落:お)ち(着:つ)いてからにしましょう。\\
\textbf{Adjusted Expression (EN)}:\\
Let's do it after I settle down a little.

\textbf{Example (JA)}:\\
A「すぐに(食:た)べましょうか？」 B「いえ、ちょっと(落:お)ち(着:つ)いてから...」\\
\textbf{Example (EN)}:\\
A: Shall we eat right away? B: No, after I settle down a bit first...

\textbf{Notes (JA)}:\\
「ちょっと落ち着いてから」は、外出から帰ってきたときに、少し、着替えをしたり、荷物整理をしたり、休憩をしたりしてから、といういろいろな意味を含め、少し休憩をしたい場合に使えます。単純に「休憩したい」というと、会話相手の人はあなたに具体的に理由や時間を聞きたくなるかもしれません。

\textbf{Notes (EN)}:\\
The phrase "chotto ochitsuite kara" is used when you want to take a break after returning from going out. It can be used when you want to take a break after returning from going out, such as changing clothes, organizing your belongings, or taking a break. If you simply say "I want to take a break," the person you are talking to may want to know the reason or time you want to take a break.

\newpage

\section*{81. By chance}

\textbf{Date}: 2024.06.30 (Sun)\\
\textbf{Index}: tamatama\\
\textbf{Expression (JA)}: たまたまですよ。\\
\textbf{Romanization}: Tamatama desu yo.\\
\textbf{Expression (EN)}: It's just a coincidence.

\textbf{Variants (JA)}: たまたまですよ。 / たまたまです。 / たまたまうまくいっただけです。\\
\textbf{Variants (EN)}: It's just a coincidence. / It just happened that way. / I just happened to do well this time.

\textbf{Intent (JA)}: 偶然の説明, 謙遜, 賞賛への応答\\
\textbf{Intent (EN)}: explaining chance, humility, response to praise

\textbf{Tags (JA)}: 即時文法, 口語, 謙遜表現, 応答表現\\
\textbf{Tags (EN)}: immediate grammar, spoken, humility expression, response expression

\textbf{Adjusted Expression (JA)}:\\
たまたまうまくいっただけです。\\
\textbf{Adjusted Expression (EN)}:\\
I just happened to do well this time.

\textbf{Example (JA)}:\\
A「(試験:しけん)、いい(点数:てんすう)がとれましたね」 B「たまたまですよ」\\
\textbf{Example (EN)}:\\
A: You got a good score on the exam. B: It's just a coincidence.

\textbf{Notes (JA)}:\\
「たまたまですよ」は、何かが偶然に起こったことを表現するために使用されます。しかし、この例文のように、何かを褒めてもらったときに、実際には偶然ではなく、あなたの努力によるものであったとしても、それが偶然であったように言うことで、謙遜の態度を示すことができます。あなたの心の中で、それが自分の努力による実力だと思っていることには、まったく問題ありません。

\textbf{Notes (EN)}:\\
The phrase "tamatama desu yo" is used to express that something happened by chance. However, as in this example, when you are praised for something, you can show a humble attitude by saying that it happened by chance, even if it was actually the result of your effort. It is perfectly fine to think in your heart that the result really came from your own ability and effort.

\newpage

\section*{80. I came here from far away.}

\textbf{Date}: 2024.06.29 (Sat)\\
\textbf{Index}: harubaru\\
\textbf{Expression (JA)}: (遥々:はるばる)やって(来:き)ました。\\
\textbf{Romanization}: Harubaru yatte kimashita.\\
\textbf{Expression (EN)}: I came here from far away.

\textbf{Variants (JA)}: (遥々:はるばる)やって(来:き)ました。 / (遥々:はるばる)(来:き)ました。 / (遠:とお)くから(来:き)ました。\\
\textbf{Variants (EN)}: I came here from far away. / I came all the way here. / I came from a long way away.

\textbf{Intent (JA)}: 遠方からの到着, 距離の強調, 軽いあいさつ\\
\textbf{Intent (EN)}: arrival from far away, emphasis on distance, light greeting

\textbf{Tags (JA)}: 語彙表現, 口語, 到着表現, 強調表現\\
\textbf{Tags (EN)}: lexical expression, spoken, arrival expression, emphatic expression

\textbf{Adjusted Expression (JA)}:\\
(遠方:えんぽう)から(参:まい)りました。\\
\textbf{Adjusted Expression (EN)}:\\
I have come from far away.

\textbf{Example (JA)}:\\
A「(遥々:はるばる)やって(来:き)ました」\\
\textbf{Example (EN)}:\\
A: I came here from far away.

\textbf{Notes (JA)}:\\
「遥々やって来ました」は、遠くから来たことを表現するために使用されます。単なる「来ました」ではなく「やって来ました」のように「やって」を前置するのがポイントです。さほど遠くもないところから来たときに、この表現を使うと当然、皮肉だと思われるでしょう。

\textbf{Notes (EN)}:\\
The phrase "harubaru yatte kimashita" is used to express that you have come from a long distance. The point is to use "yatte" before "kimashita" to express that you have come from far away. When you use this expression after coming from a place that is not so far away, it will naturally be taken as ironic.

\newpage

\section*{79. Thank you for taking care of me!}

\textbf{Date}: 2024.06.28 (Fri)\\
\textbf{Index}: osewaninarimashita\\
\textbf{Expression (JA)}: お(世話:せわ)になりました。\\
\textbf{Romanization}: Osewa ni narimashita.\\
\textbf{Expression (EN)}: Thank you for taking care of me!

\textbf{Variants (JA)}: お(世話:せわ)になりました。 / いろいろお(世話:せわ)になりました。 / 大変お(世話:せわ)になりました。\\
\textbf{Variants (EN)}: Thank you for taking care of me. / Thank you for all your help. / Thank you very much for your support.

\textbf{Intent (JA)}: 感謝, 別れのあいさつ, 丁寧表現\\
\textbf{Intent (EN)}: gratitude, parting greeting, polite expression

\textbf{Tags (JA)}: 定型表現, 口語, 感謝表現, あいさつ表現\\
\textbf{Tags (EN)}: formulaic expression, spoken, gratitude expression, greeting expression

\textbf{Adjusted Expression (JA)}:\\
大変お(世話:せわ)になりました。\\
\textbf{Adjusted Expression (EN)}:\\
Thank you very much for all your support.

\textbf{Example (JA)}:\\
A「(今日:きょう)は、お(世話:せわ)になりました」\\
\textbf{Example (EN)}:\\
A: Thank you for taking care of me today.

\textbf{Notes (JA)}:\\
「お世話になりました」は、受けたお世話や支援に対して感謝の気持ちを表現するために使用される丁寧な表現です。

\textbf{Notes (EN)}:\\
The phrase "osewa ni narimashita" is a polite expression used to express gratitude for the care and support you have received.

\newpage

\section*{78. Long time no see!}

\textbf{Date}: 2024.06.27 (Thu)\\
\textbf{Index}: ohisashiburi\\
\textbf{Expression (JA)}: お(久:ひさ)しぶりです。\\
\textbf{Romanization}: Ohisashiburi desu.\\
\textbf{Expression (EN)}: Long time no see!

\textbf{Variants (JA)}: お(久:ひさ)しぶりです。 / お(久:ひさ)しぶりですね。 / ご(無沙汰:ぶさた)しております。\\
\textbf{Variants (EN)}: Long time no see! / It's been a while. / I haven't seen you in a long time.

\textbf{Intent (JA)}: 再会のあいさつ, 丁寧な呼びかけ, 親しみの表現\\
\textbf{Intent (EN)}: greeting on meeting again, polite opening, expression of familiarity

\textbf{Tags (JA)}: 定型表現, 口語, あいさつ表現, 再会表現\\
\textbf{Tags (EN)}: formulaic expression, spoken, greeting expression, reunion expression

\textbf{Adjusted Expression (JA)}:\\
ご(無沙汰:ぶさた)しております。\\
\textbf{Adjusted Expression (EN)}:\\
It has been a long time since we last met.

\textbf{Example (JA)}:\\
A「お(久:ひさ)しぶりですね」\\
\textbf{Example (EN)}:\\
A: Long time no see!

\textbf{Notes (JA)}:\\
「お久しぶりです」は、久しぶりに会った人に対して使用される丁寧な表現です。学校を卒業してから会っていない友達や、転職してから会っていない同僚など、久しぶりに会った人に対して使用できます。

\textbf{Notes (EN)}:\\
The phrase "ohisashiburi desu" is a polite expression used when you meet someone you have not seen in a long time. It can be used when you meet a friend you have not seen since you graduated from school, or a colleague you have not seen since you changed jobs.

\newpage

\section*{77. What a waste.}

\textbf{Date}: 2024.06.26 (Wed)\\
\textbf{Index}: mottainai\\
\textbf{Expression (JA)}: もったいない。\\
\textbf{Romanization}: mottainai.\\
\textbf{Expression (EN)}: What a waste.

\textbf{Variants (JA)}: もったいない。 / (時間:じかん)がもったいない。 / それはもったいない。\\
\textbf{Variants (EN)}: What a waste. / That's a waste of time. / It would be a waste.

\textbf{Intent (JA)}: 無駄への反応, 後悔, 戒め\\
\textbf{Intent (EN)}: reaction to waste, regret, warning

\textbf{Tags (JA)}: 即時文法, 口語, 反応表現, 価値判断表現\\
\textbf{Tags (EN)}: immediate grammar, spoken, reaction expression, value-judgment expression

\textbf{Adjusted Expression (JA)}:\\
それは(無駄:むだ)になってしまうので、もったいないです。\\
\textbf{Adjusted Expression (EN)}:\\
That would be a waste, so it would be better not to do that.

\textbf{Example (JA)}:\\
A「(携帯:けいたい)を(見:み)ていたのでは、(時間:じかん)がもったいないので、(宿題:しゅくだい)をしてしまいましょう」\\
\textbf{Example (EN)}:\\
A: If we keep looking at our phones, it will be a waste of time, so let's get our homework done.

\textbf{Notes (JA)}:\\
「もったいない」は、何かが無駄になってしまったことに対して、後悔や失望を表現するために使用されます。何かが無駄になったり、効果的に使われていないことを示すために使用されます。食べ物を捨てるのを見たときや、まだ使えるものを使わないのを見たときに使われることが多いです。また、誰かが自分の能力を効果的に使っていないのを見たときにも使われます。

\textbf{Notes (EN)}:\\
The phrase "mottainai" is used to express regret or disappointment that something is wasted. It is used to indicate that something is wasted or not used effectively. It is often used when you see someone throwing away food or when you see someone not using something that is still useful. It is also used when you see someone not using their abilities effectively.

\newpage

\section*{76. Better that way}

\textbf{Date}: 2024.06.25 (Tue)\\
\textbf{Index}: sonohouga\\
\textbf{Expression (JA)}: その(方:ほう)が...\\
\textbf{Romanization}: Sono hou ga...\\
\textbf{Expression (EN)}: That might be better...

\textbf{Variants (JA)}: その(方:ほう)が... / そうした(方:ほう)が... / その(方:ほう)がいいかもね。\\
\textbf{Variants (EN)}: That might be better... / It may be better that way... / Maybe that's the better option...

\textbf{Intent (JA)}: 提案, 婉曲な助言, 配慮\\
\textbf{Intent (EN)}: suggestion, indirect advice, consideration

\textbf{Tags (JA)}: 即時文法, 口語, 婉曲表現, 配慮表現, 省略表現\\
\textbf{Tags (EN)}: immediate grammar, spoken, indirect expression, considerate expression, elliptical expression

\textbf{Adjusted Expression (JA)}:\\
そうした方がよいと思います。\\
\textbf{Adjusted Expression (EN)}:\\
I think it would be better that way.

\textbf{Example (JA)}:\\
A「やっぱり、(全部:ぜんぶ)(手作業:てさぎょう)でやります」 B「ええ、たしかに(時間:じかん)はかかるけど、その(方:ほう)が...」\\
\textbf{Example (EN)}:\\
A: I'll do it all by hand after all. B: Yes, it certainly takes time, but that might be better...

\textbf{Notes (JA)}:\\
「その方が...」は、相手の決定が必ずしも最善ではないにしても、ベターであることを示唆する表現です。省略せずに言うと、「そうした方がいいかもしれません」という意味になりますが、文を途中で終わらせることで、相手の考えや意思決定を尊重する姿勢を表現できます。この曖昧さが丁寧さや配慮を生み出します。

\textbf{Notes (EN)}:\\
The phrase "sono hou ga..." is an expression that suggests something might be better, even if it is not the perfect choice. If you say it without leaving it unfinished, it would mean something like "It may be better to do it that way," but ending the sentence midway shows a willingness to respect the other person's decision without explicitly stating everything. This subtlety creates politeness and consideration.

\newpage

\section*{75. That's it!}

\textbf{Date}: 2024.06.24 (Mon)\\
\textbf{Index}: soredesu\\
\textbf{Expression (JA)}: ええ、それです。\\
\textbf{Romanization}: Ee, sore desu.\\
\textbf{Expression (EN)}: That's it!

\textbf{Variants (JA)}: ええ、それです。 / あ、それです。 / それそれ。\\
\textbf{Variants (EN)}: That's it! / Yes, that's it. / That's the one.

\textbf{Intent (JA)}: 確認, 同意, 思い出しの達成\\
\textbf{Intent (EN)}: confirmation, agreement, successful recall

\textbf{Tags (JA)}: 即時文法, 口語, 確認表現, 反応表現\\
\textbf{Tags (EN)}: immediate grammar, spoken, confirmation expression, reaction expression

\textbf{Adjusted Expression (JA)}:\\
ええ、それが(正解:せいかい)です。\\
\textbf{Adjusted Expression (EN)}:\\
Yes, that's the correct answer.

\textbf{Example (JA)}:\\
A「(昨日:きのう)(行:い)った(カフェ:かふぇ)の(名前:なまえ)、(何:なん)だったっけ？」 B「(ブルー:ぶるー)じゃなかった？」 A「ええ、それです！」\\
\textbf{Example (EN)}:\\
A: What was the name of the cafe we went to yesterday? B: Wasn't it Blue? A: That's it!

\textbf{Notes (JA)}:\\
「ええ、それです」は、自分が思い出せないものの答を話し相手が教えてくれたときに、同意するときの表現です。探していたものを誰かが見つけてくれて、確かにそれであることを認めるときにも使用できます。「ええ」「それ」「です」など簡単な語の組み合わせですが、よく日本人の会話を観察していないと、このような簡単な語で表現できることはなかなかわかりません。このような表現をイディオムと言います。イディオムが使えると自然な言い方が体得できます。

\textbf{Notes (EN)}:\\
The phrase "ee, sore desu" is used to agree when someone tells you the answer to something you cannot remember. It can also be used when someone finds what you were looking for and you acknowledge that it is indeed what you were looking for. It is a combination of simple words such as "ee," "sore," and "desu," but if you do not observe Japanese conversations well, you may not realize that you can express it with such simple words. Such expressions are called idioms. If you can use idioms, you can naturally learn how to speak.

\newpage

\section*{74. Is that so?}

\textbf{Date}: 2024.06.23 (Sun)\\
\textbf{Index}: sounandesuka\\
\textbf{Expression (JA)}: そうなんですか？\\
\textbf{Romanization}: Sounandesuka?\\
\textbf{Expression (EN)}: Is that so?

\textbf{Variants (JA)}: そうなんですか？ / そうなんですか。 / そうなんですね。\\
\textbf{Variants (EN)}: Is that so? / Really? / Oh, is that so?

\textbf{Intent (JA)}: 驚き, あいづち, 疑い\\
\textbf{Intent (EN)}: surprise, backchannel response, doubt

\textbf{Tags (JA)}: 即時文法, 口語, 反応表現, あいづち表現\\
\textbf{Tags (EN)}: immediate grammar, spoken, reaction expression, backchannel expression

\textbf{Adjusted Expression (JA)}:\\
そうなのですか？\\
\textbf{Adjusted Expression (EN)}:\\
Is that so?

\textbf{Example (JA)}:\\
A「うちの子、もう(今年:ことし)、(小学校:しょうがっこう)に(入学:にゅうがく)するんですよ」 B「そうなんですか？」\\
\textbf{Example (EN)}:\\
A: My child is entering elementary school this year. B: Is that so?

\textbf{Notes (JA)}:\\
「そうなんですか？」は驚きや疑いを表現するために使用される間投詞です。何かが本当かどうかを確認するためというより、ちょっと驚いてみせて、喜びを受け取りましたという感じを表現するために使用されます。一方、疑いを表現するときには、語尾を少し上げて長めに発音します。発音の仕方で意味が大きく変化します。

\textbf{Notes (EN)}:\\
The phrase "sounandesuka?" is used as an interjection to express surprise or doubt. It is used not so much to confirm whether something is true, but to show a little surprise and to express that you have received the information with interest or even joy. On the other hand, when expressing doubt, raise the end of the phrase and pronounce it a little longer. The meaning changes significantly depending on how you pronounce it.

\newpage

\section*{73. How about it?}

\textbf{Date}: 2024.06.22 (Sat)\\
\textbf{Index}: ikagadesuka\\
\textbf{Expression (JA)}: いかがですか？\\
\textbf{Romanization}: Ikaga desu ka?\\
\textbf{Expression (EN)}: How about it?

\textbf{Variants (JA)}: いかがですか？ / いかがでしょうか？\\
\textbf{Variants (EN)}: How about it? / Would you like it? / What do you think?

\textbf{Intent (JA)}: 提案, 勧め, 丁寧な問いかけ\\
\textbf{Intent (EN)}: suggestion, offering, polite inquiry

\textbf{Tags (JA)}: 定型表現, 口語, 提案表現, 丁寧表現\\
\textbf{Tags (EN)}: formulaic expression, spoken, suggestion expression, polite expression

\textbf{Adjusted Expression (JA)}:\\
いかがでしょうか？\\
\textbf{Adjusted Expression (EN)}:\\
Would you like it?

\textbf{Example (JA)}:\\
A「(田中:たなか)(先生:せんせい)、(私:わたし)は、これを(食:た)べてみたいのですが、(先生:せんせい)もいかがですか？」\\
\textbf{Example (EN)}:\\
A: Professor Tanaka, I would like to try this. How about you?

\textbf{Notes (JA)}:\\
「いかがですか？」は、何かを相手に提案したいときによく使われる丁寧な表現です。日本語では一般的に相手の気持ちや意向を質問する時には、敬語を使わないと失礼になります。しかし、敬語の語形変化はさまざまなので、覚えたり、発音したりするのは、面倒です。そこで、自分の意向をあらかじめ伝えてから、「いかがですか」と言うと、敬語を使わずに丁寧に伺うことができます。

\textbf{Notes (EN)}:\\
The phrase "ikagadesuka?" is used to suggest something to someone. It is a polite expression that is often used when you want to propose something to someone. In Japanese, it is generally considered rude to ask someone's feelings or intentions without using honorific language. However, since there are various forms of honorific language, it is troublesome to remember or pronounce them. Therefore, if you express your intentions in advance and then say "ikagadesuka?", you can ask politely without using honorific language.

\newpage

\section*{72. Hobbies}

\textbf{Date}: 2024.06.21 (Fri)\\
\textbf{Index}: shumi\\
\textbf{Expression (JA)}: (趣味:しゅみ)は(何:なん)ですか？\\
\textbf{Romanization}: Shumi wa nan desu ka?\\
\textbf{Expression (EN)}: What is your hobby?

\textbf{Variants (JA)}: 趣味は何ですか？ / ご趣味は何ですか？\\
\textbf{Variants (EN)}: What is your hobby? / What are your hobbies?

\textbf{Intent (JA)}: 質問, 会話の開始, 相手への関心\\
\textbf{Intent (EN)}: question, starting a conversation, interest in the other person

\textbf{Tags (JA)}: 定型表現, 口語, 質問表現, 会話表現\\
\textbf{Tags (EN)}: formulaic expression, spoken, question expression, conversation expression

\textbf{Adjusted Expression (JA)}:\\
ご(趣味:しゅみ)は(何:なん)ですか？\\
\textbf{Adjusted Expression (EN)}:\\
What are your hobbies?

\textbf{Example (JA)}:\\
A「(趣味:しゅみ)は(何:なん)ですか？」 B「(音楽:おんがく)を(聴:き)くことです」\\
\textbf{Example (EN)}:\\
A: What is your hobby? B: Listening to music.

\textbf{Notes (JA)}:\\
「趣味は何ですか？」は、相手の趣味を尋ねるために使用される表現です。

\textbf{Notes (EN)}:\\
The phrase "shumi wa nan desu ka?" is used to ask someone what their hobby is.

\newpage

\section*{71. Sorry I couldn't do it well.}

\textbf{Date}: 2024.06.20 (Thu)\\
\textbf{Index}: gomenne\\
\textbf{Expression (JA)}: ごめんね、うまくできなくて。\\
\textbf{Romanization}: Gomen ne, umaku dekinakute.\\
\textbf{Expression (EN)}: I'm sorry I couldn't do it well.

\textbf{Variants (JA)}: ごめんね、うまくできなくて。 / ごめん、うまくできなくて。 / ごめんなさい、うまくできなくて。\\
\textbf{Variants (EN)}: I'm sorry I couldn't do it well. / Sorry I messed it up. / I'm sorry it didn't turn out well.

\textbf{Intent (JA)}: 謝罪, 失敗への反応, やわらかい言い方\\
\textbf{Intent (EN)}: apology, reaction to failure, soft expression

\textbf{Tags (JA)}: 即時文法, 口語, 謝罪表現, 失敗表現\\
\textbf{Tags (EN)}: immediate grammar, spoken, apology expression, failure expression

\textbf{Adjusted Expression (JA)}:\\
(申:もう)し(訳:わけ)ありません、うまくできなくて。\\
\textbf{Adjusted Expression (EN)}:\\
I'm sorry I wasn't able to do it well.

\textbf{Example (JA)}:\\
A「(料理:りょうり)、(失敗:しっぱい)しちゃった。ごめんね、うまくできなくて」\\
\textbf{Example (EN)}:\\
A: I messed up the cooking. I'm sorry I couldn't do it well.

\textbf{Notes (JA)}:\\
「ごめんね、うまくできなくて」は、何かをうまくできなかったことに対して謝罪を表現するために使用されます。

\textbf{Notes (EN)}:\\
The phrase "gomen ne, umaku dekinakute" is used to express an apology for not being able to do something well.

\newpage

\section*{70. It helps me a lot!}

\textbf{Date}: 2024.06.19 (Wed)\\
\textbf{Index}: tasukarimasu\\
\textbf{Expression (JA)}: (本当:ほんとう)に(助:たす)かります！\\
\textbf{Romanization}: hontouni tasukarimasu!\\
\textbf{Expression (EN)}: It helps me a lot!

\textbf{Variants (JA)}: 助かります。 / 本当に助かりました。\\
\textbf{Variants (EN)}: That really helps me. / That helped me a lot.

\textbf{Intent (JA)}: 感謝, 安心, 助力への反応\\
\textbf{Intent (EN)}: gratitude, relief, reaction to help

\textbf{Tags (JA)}: 即時文法, 口語, 感謝表現, 反応表現\\
\textbf{Tags (EN)}: immediate grammar, spoken, gratitude expression, reaction expression

\textbf{Adjusted Expression (JA)}:\\
(大変:たいへん)(助:たす)かります。\\
\textbf{Adjusted Expression (EN)}:\\
That is very helpful.

\textbf{Example (JA)}:\\
A「(手伝:てつだ)いましょうか？」 B「(本当:ほんとう)に(助:たす)かります！」\\
\textbf{Example (EN)}:\\
A: Can I help you? B: That really helps me!

\textbf{Notes (JA)}:\\
「本当に助かります」は、何かをしてもらう、あるいはしてもらったときに感謝を伝える表現です。

\textbf{Notes (EN)}:\\
The phrase "hontouni tasukarimasu" is used to express gratitude when someone will do or has done something for you.

\newpage

\section*{69. That's not fair!}

\textbf{Date}: 2024.06.18 (Tue)\\
\textbf{Index}: sorehanaideshou\\
\textbf{Expression (JA)}: それはないでしょう！\\
\textbf{Romanization}: sorewanai deshou!\\
\textbf{Expression (EN)}: That's not fair!

\textbf{Variants (JA)}: それはないでしょう！ / それはないよ。 / それはひどいよ。\\
\textbf{Variants (EN)}: That's not fair! / Come on, that's not right. / That's too much!

\textbf{Intent (JA)}: 不満, 異議, 理不尽への反応\\
\textbf{Intent (EN)}: dissatisfaction, objection, reaction to unfairness

\textbf{Tags (JA)}: 即時文法, 口語, 不同意表現, 反応表現\\
\textbf{Tags (EN)}: immediate grammar, spoken, disagreement expression, reaction expression

\textbf{Adjusted Expression (JA)}:\\
それは(少:すこ)しひどいと(思:おも)います。\\
\textbf{Adjusted Expression (EN)}:\\
I think that's a bit unfair.

\textbf{Example (JA)}:\\
A「お(菓子:かし)、(全部:ぜんぶ)、(食:た)べちゃった」 B「それはないでしょう！」\\
\textbf{Example (EN)}:\\
A: I ate all the sweets. B: That's not fair!

\textbf{Notes (JA)}:\\
「それはないでしょう」は、不公平や理不尽を表現するために使用されます。不満や異議を表現するために声を上げることになりますので、使いすぎると、周囲の人々はいつも不満を持っている人だと思うかもしれません。

\textbf{Notes (EN)}:\\
The phrase "sorewanai deshou" is used to express unfairness or injustice. It is used to raise your voice to express dissatisfaction or objection. If you use this expression too much, people around you may think you are always dissatisfied.

\newpage

\section*{68. I did it!}

\textbf{Date}: 2024.06.17 (Mon)\\
\textbf{Index}: yattaa\\
\textbf{Expression (JA)}: やったぁ！\\
\textbf{Romanization}: Yattaa!\\
\textbf{Expression (EN)}: I did it!

\textbf{Variants (JA)}: やった！ / やりました！\\
\textbf{Variants (EN)}: I did it! / Yay! / We made it!

\textbf{Intent (JA)}: 達成, 喜び, 成功への反応\\
\textbf{Intent (EN)}: achievement, joy, reaction to success

\textbf{Tags (JA)}: 即時文法, 口語, 反応表現, 感情表現\\
\textbf{Tags (EN)}: immediate grammar, spoken, reaction expression, emotional expression

\textbf{Adjusted Expression (JA)}:\\
やりました。\\
\textbf{Adjusted Expression (EN)}:\\
I did it.

\textbf{Example (JA)}:\\
A「(試験:しけん)、(合格:ごうかく)したよ」 B「やったぁ！」\\
\textbf{Example (EN)}:\\
A: I passed the exam. B: I did it!

\textbf{Notes (JA)}:\\
「やったぁ」は、達成感や成功の喜びを表現するために使用されます。この表現を使うときには、声を上げて、達成感や成功の喜びを表現します。しかし、この表現をあまりにも頻繁に使うと、周囲の人々に「自己顕示欲が強い人」と思われるかもしれません。また、この表現を使うときに限って、喜べない状況にある人がいるものです。

\textbf{Notes (EN)}:\\
The phrase "yattaa" is used to express a sense of accomplishment or joy in success. When you use this expression, raise your voice to express a sense of accomplishment or joy in success. However, if you use this expression too much, people around you may think you are a person with a strong desire for self-expression. Also, when you use this expression, there may be people who are in a situation where they cannot be happy.

\newpage

\section*{67. What you say and what you think..}

\textbf{Date}: 2024.06.16 (Sun)\\
\textbf{Index}: soukanaa\\
\textbf{Expression (JA)}: そうかなぁ？【(納得:なっとく)いかないけど】\\
\textbf{Romanization}: soukanaa? (nattoku ikanai kedo)\\
\textbf{Expression (EN)}: I see... (but I'm not convinced)

\textbf{Variants (JA)}: そうかなぁ？ / そうですかね。 / そうかな。\\
\textbf{Variants (EN)}: I'm not so sure. / Really? / I wonder about that.

\textbf{Intent (JA)}: 疑い, 異論, やわらかい不同意\\
\textbf{Intent (EN)}: doubt, disagreement, soft disagreement

\textbf{Tags (JA)}: 即時文法, 口語, 不同意表現, 反応表現\\
\textbf{Tags (EN)}: immediate grammar, spoken, disagreement expression, reaction expression

\textbf{Adjusted Expression (JA)}:\\
そうでしょうか。(私:わたし)はそうではないと(思:おも)います。\\
\textbf{Adjusted Expression (EN)}:\\
Is that so? I do not think so.

\textbf{Example (JA)}:\\
A「これ、おいしいね」 B「そうかなぁ？ じゃないと(思:おも)うけど」\\
\textbf{Example (EN)}:\\
A: This is delicious. B: I'm not so sure. I don't think so.

\textbf{Notes (JA)}:\\
「そうかなぁ？」は疑いや異論を表現するために使用されます。他の人が言ったことに納得していないことを示すために使用されます。他人の意見に納得していないときに強く否定するのはあまり頭のいい人がすることではありません。しかし、嘘をついて「そうですね」というのも「何でも同意する人」と思われるので得策ではありません。そんなときは、小さな声で「そうかなぁ」と囁きましょう。

\textbf{Notes (EN)}:\\
The phrase "soukanaa" is used to express doubt or disagreement. It is used to indicate that you do not agree with what someone else has said. It is not very smart to strongly deny when you do not agree with what someone else has said. However, it is not a good idea to lie and say, "Yes, it is." When you are not convinced of someone else's opinion, whisper "soukanaa" in a very small voice.

\newpage

\section*{66. Are you sure?}

\textbf{Date}: 2024.06.15 (Sat)\\
\textbf{Index}: hontoudesukaa\\
\textbf{Expression (JA)}: 本当ですかぁ？\\
\textbf{Romanization}: hontoudesukaa?\\
\textbf{Expression (EN)}: Really?

\textbf{Variants (JA)}: 本当ですか？ / ほんとうですか？ / 本当なの？\\
\textbf{Variants (EN)}: Really? / Is that true? / Are you sure?

\textbf{Intent (JA)}: 驚き, 疑い, 確認\\
\textbf{Intent (EN)}: surprise, doubt, confirmation

\textbf{Tags (JA)}: 即時文法, 口語, 反応表現, 確認表現\\
\textbf{Tags (EN)}: immediate grammar, spoken, reaction expression, confirmation expression

\textbf{Adjusted Expression (JA)}:\\
本当でしょうか。\\
\textbf{Adjusted Expression (EN)}:\\
Could that really be true?

\textbf{Example (JA)}:\\
A「(明日:あした)、(雨:あめ)が(降:ふ)るって(言:い)ってたよ」 B「(本当:ほんとう)ですかぁ？こんなに(晴:は)れているのに」\\
\textbf{Example (EN)}:\\
A: It's said that it will rain tomorrow. B: Really? Even though it's so sunny today.

\textbf{Notes (JA)}:\\
「本当ですかぁ？」は、驚きや疑いを表現するために使用される間投詞です。何かが本当かどうかを確認するときに使用されます。語尾を上げて少々伸ばして発音してください。

\textbf{Notes (EN)}:\\
The phrase "hontoudesukaa" is used to express surprise or doubt. It is used to confirm whether something is true or not. When you use this phrase, raise the end of the word and pronounce it a little longer.

\newpage

\section*{65. I'm relieved.}

\textbf{Date}: 2024.06.14 (Fri)\\
\textbf{Index}: hottoshita\\
\textbf{Expression (JA)}: ほっとした。\\
\textbf{Romanization}: hotto shita.\\
\textbf{Expression (EN)}: I'm relieved.

\textbf{Variants (JA)}: ほっとした。 / ほっとしました。\\
\textbf{Variants (EN)}: I'm relieved. / I feel relieved.

\textbf{Intent (JA)}: 安心, 解放感, 出来事への反応\\
\textbf{Intent (EN)}: relief, sense of release, reaction to an event

\textbf{Tags (JA)}: 即時文法, 口語, 感情表現, 反応表現\\
\textbf{Tags (EN)}: immediate grammar, spoken, emotional expression, reaction expression

\textbf{Adjusted Expression (JA)}:\\
ひとまず安心しました。\\
\textbf{Adjusted Expression (EN)}:\\
For now, I feel relieved.

\textbf{Example (JA)}:\\
A「(試験:しけん)、(全部:ぜんぶ)、(終:お)わって、ほっとした」\\
\textbf{Example (EN)}:\\
A: All the exams are over, and I'm kinda relieved.

\textbf{Notes (JA)}:\\
「ほっとした」は安心を表現する表現です。何かが終わったことや何かが解決したことを示すために使用されます。しかし、本当にほっとすることというのはあるのでしょうか。試験が終わったら、また次の試験があります。気になっていた女の子が自分のことが好きだということがわかって、そのときは「ほっとする」かもしれませんが、お互いにお互いが好きだとわかってから、以前よりドキドキして愛おしく思うことは少なくなるかもしれません。

\textbf{Notes (EN)}:\\
The phrase "hotto shita" is used to express relief. It is used to indicate that something has ended or been resolved. However, is there really anything to be relieved about? When one exam is over, there is another exam. When you find out that a girl you were worried about likes you, you may feel relieved at that time, but after you find out that you both like each other, you may feel less relieved and less excited and affectionate than before.

\newpage

\section*{64. Adding "o" to "hazukashii" makes it polite.}

\textbf{Date}: 2024.06.13 (Thu)\\
\textbf{Index}: ohazukashii\\
\textbf{Expression (JA)}: 失礼しました。これはこれは、お(恥:はず)かしい。\\
\textbf{Romanization}: Shitsurei shimashita. Kore wa kore wa, ohazukashii.\\
\textbf{Expression (EN)}: I'm sorry. Well, this is embarrassing.

\textbf{Variants (JA)}: これはこれは、お(恥:はず)かしい。 / お(恥:はず)かしいです。\\
\textbf{Variants (EN)}: Well, this is embarrassing. / How embarrassing.

\textbf{Intent (JA)}: 軽い謝罪, 雰囲気の緩和, 失敗への反応\\
\textbf{Intent (EN)}: light apology, easing the atmosphere, reaction to a mistake

\textbf{Tags (JA)}: 即時文法, 口語, 丁寧表現, 場面調整表現\\
\textbf{Tags (EN)}: immediate grammar, spoken, polite expression, interaction-adjusting expression

\textbf{Adjusted Expression (JA)}:\\
失礼しました。お(恥:は)ずかしいところをお(見:み)せして。\\
\textbf{Adjusted Expression (EN)}:\\
I'm sorry. I'm showing you an embarrassing side of me.

\textbf{Example (JA)}:\\
A「(失礼:しつれい)しました。これはこれは、お(恥:は)ずかしい」\\
\textbf{Example (EN)}:\\
A: I'm sorry. Well, this is embarrassing.

\textbf{Notes (JA)}:\\
イ形容詞の「恥ずかしい」ということばを直接発することは稀ですが、自分のとるにたらぬ、つまらぬ間違いを犯したときに、軽い謝罪のことばに添えて自分のことを「恥ずかしい」ということはあります。そのときには「お恥ずかしい」のように「お」を付けて発話すると丁寧で、雰囲気を和ませる効果があります。そのようにきれいな返答ができるあなたを誰も「恥ずかしい人」とは思わず、非常に世間的な常識を持った人だと思うでしょう。

\textbf{Notes (EN)}:\\
It is rare to directly say the word "hazukashii," which means "embarrassing," in Japanese. However, when you make a trivial mistake or a minor error, you may say this phrase with a light apology such as "shitsurei shimashita," which means "I'm sorry." When you utter this phrase, it is polite and has a relaxing effect on the conversation. If you can respond beautifully like this, no one will think of you as an "embarrassing person," but rather as a person with very worldly common sense.

\newpage

\section*{63. Wow!}

\textbf{Date}: 2024.06.12 (Wed)\\
\textbf{Index}: waa\\
\textbf{Expression (JA)}: わあ！\\
\textbf{Romanization}: waa!\\
\textbf{Expression (EN)}: Wow!

\textbf{Variants (JA)}: わあ。 / わあ、すごい！\\
\textbf{Variants (EN)}: Wow! / Wow, that's amazing!

\textbf{Intent (JA)}: 驚き, 感嘆, 即時反応\\
\textbf{Intent (EN)}: surprise, admiration, immediate reaction

\textbf{Tags (JA)}: 即時文法, 口語, 反応表現, 感動詞\\
\textbf{Tags (EN)}: immediate grammar, spoken, reaction expression, interjection

\textbf{Adjusted Expression (JA)}:\\
わあ、すばらしいですね。\\
\textbf{Adjusted Expression (EN)}:\\
Wow, that is wonderful.

\textbf{Example (JA)}:\\
A「わあ、すごい！」\\
\textbf{Example (EN)}:\\
A: Wow, that's amazing!

\textbf{Notes (JA)}:\\
「わあ」は驚きや感嘆を表現するために使用される間投詞です。この表現はカジュアルな場面でもフォーマルな場面でも使用できます。基本的に自分の態度を思わず発する表現を無理に丁寧な形にする必要はありません。

\textbf{Notes (EN)}:\\
The word "waa" is an interjection that is used to express surprise or admiration. It is used to indicate that something is impressive or unexpected. This expression can be used in both casual and formal situations. Basically, you do not have to express your attitude in a polite way.

\newpage

\section*{62. Oops! and Otto!}

\textbf{Date}: 2024.06.11 (Tue)\\
\textbf{Index}: otto\\
\textbf{Expression (JA)}: おっと！\\
\textbf{Romanization}: Otto!\\
\textbf{Expression (EN)}: Oops!

\textbf{Variants (JA)}: おっと。 / おっとっと。\\
\textbf{Variants (EN)}: Oops! / Whoops!

\textbf{Intent (JA)}: 気づき, 失敗への反応, 即時反応\\
\textbf{Intent (EN)}: realization, reaction to a mistake, immediate reaction

\textbf{Tags (JA)}: 即時文法, 口語, 反応表現, 感動詞\\
\textbf{Tags (EN)}: immediate grammar, spoken, reaction expression, interjection

\textbf{Adjusted Expression (JA)}:\\
あ、(傘:かさ)を持ってくるのを(忘:わす)れていました。\\
\textbf{Adjusted Expression (EN)}:\\
Oh, I forgot to bring my umbrella.

\textbf{Example (JA)}:\\
A「おっと、(傘:かさ)、持ってくるの、(忘:わす)れてた」\\
\textbf{Example (EN)}:\\
A: Oops, I forgot to bring my umbrella.

\textbf{Notes (JA)}:\\
「おっと」は、予期せぬことが起こったことを表現するために使用されます。 何かを間違えたり、不利なことが起こったことに気づいたときに無意識に発せられる表現です。これは、意図的な表現よりも自然な発言です。これを使うと、他人にあなたがネイティブに近い話者であることを感じさせるかもしれません。

\textbf{Notes (EN)}:\\
The phrase "otto" is used to express that something happened that you unexpected. The phrase is uttered unconsciously when you notice that you have done something wrong or unfavrable. This is a natural utterance rather than an intentional expression. If you use this, it may give others that you are a near-native speaker.

\newpage

\section*{61. Either Way}

\textbf{Date}: 2024.06.10 (Mon)\\
\textbf{Index}: izureniseyo\\
\textbf{Expression (JA)}: いずれにせよ\\
\textbf{Romanization}: Izure ni seyo\\
\textbf{Expression (EN)}: either way

\textbf{Variants (JA)}: いずれにしても / どちらにしても / ともかく\\
\textbf{Variants (EN)}: either way / in any case / at any rate

\textbf{Intent (JA)}: 論点の整理, 話の前進, 結論への接続\\
\textbf{Intent (EN)}: organizing the point, moving the discussion forward, linking to a conclusion

\textbf{Tags (JA)}: 定型表現, 接続表現, 説明表現, 論述表現\\
\textbf{Tags (EN)}: formulaic expression, connective expression, explanatory expression, discoursal expression

\textbf{Adjusted Expression (JA)}:\\
どのような(場合:ばあい)であっても、(準備:じゅんび)をしておいたほうがよいと(思:おも)います。\\
\textbf{Adjusted Expression (EN)}:\\
In any case, I think it is better to be prepared.

\textbf{Example (JA)}:\\
A「いずれにせよ、(準備:じゅんび)をしておいたほうがいい」\\
\textbf{Example (EN)}:\\
A: Either way, it is better to be prepared.

\textbf{Notes (JA)}:\\
「いずれにせよ」は、いくつかの可能性や前提の違いがあっても、その先で共通して言えることを述べるときに使う表現です。話の枝分かれにとどまらず、「どちらにしても」「ともかく」と話を先へ進める働きを持つ接続的な言い方です。

\textbf{Notes (EN)}:\\
"Izure ni seyo" is used when the speaker wants to move past differing possibilities and state what remains true in any case. It often works as a connective expression meaning something like "either way," "in any case," or "whatever happens."

\newpage

\section*{60. Everyone's Doing Fine}

\textbf{Date}: 2024.06.09 (Sun)\\
\textbf{Index}: minnagenkishiteruyo\\
\textbf{Expression (JA)}: みんな、(元気:げんき)してるよ。\\
\textbf{Romanization}: Minna, genki shiteru yo.\\
\textbf{Expression (EN)}: Everyone's doing fine.

\textbf{Variants (JA)}: みんな、(元気:げんき)してるよ。 / みんな(元気:げんき)だよ。 / みんな(元気:げんき)にしてるよ。\\
\textbf{Variants (EN)}: Everyone's doing fine. / Everyone's fine. / They're all doing well.

\textbf{Intent (JA)}: 近況報告, 応答, 気づかいへの返答\\
\textbf{Intent (EN)}: reporting how people are doing, response, reply to concern

\textbf{Tags (JA)}: 口語, 近況表現, 応答表現, あいさつ表現\\
\textbf{Tags (EN)}: spoken, well-being expression, response expression, greeting-related expression

\textbf{Adjusted Expression (JA)}:\\
みんな(元気:げんき)にしていますよ。\\
\textbf{Adjusted Expression (EN)}:\\
Everyone is doing well.

\textbf{Example (JA)}:\\
A「みんな、(元気:げんき)してる？」 B「みんな、(元気:げんき)してるよ」\\
\textbf{Example (EN)}:\\
A: How's everyone doing? B: Everyone's doing fine.

\textbf{Notes (JA)}:\\
「みんな、元気してるよ」は、家族や友人など、身近な人たちの様子をたずねられたときに、「みんな元気にしている」と気軽に答える口語表現です。あいさつそのものというより、近況をたずねられたときの返答として使われやすい表現です。方言によっては、「してる」が「しとる」になり、「元気しとる？」「元気しとるよ」のようにも言われます。

\textbf{Notes (EN)}:\\
"Minna, genki shiteru yo" is a casual way to say that everyone is doing well. It often appears as a response when someone asks about family members, friends, or people back home. In some dialects, "shiteru" may appear as "shitoru," as in "genki shitoru?" or "genki shitoru yo."

\newpage

\section*{59. Either Way}

\textbf{Date}: 2024.06.08 (Sat)\\
\textbf{Index}: dochiranishitemo\\
\textbf{Expression (JA)}: どちらにしても\\
\textbf{Romanization}: Dochira ni shitemo\\
\textbf{Expression (EN)}: either way

\textbf{Variants (JA)}: どちらにしても / どっちにしても / いずれにしても\\
\textbf{Variants (EN)}: either way / in any case / whichever way it goes

\textbf{Intent (JA)}: 状況整理, 結論の提示, 議論の収束\\
\textbf{Intent (EN)}: sorting out the situation, stating a conclusion, bringing discussion to a close

\textbf{Tags (JA)}: 即時文法, 口語, 結論提示表現, 議論整理表現\\
\textbf{Tags (EN)}: immediate grammar, spoken, conclusion-marking expression, discussion-organizing expression

\textbf{Adjusted Expression (JA)}:\\
いずれにしても、確認が必要です。\\
\textbf{Adjusted Expression (EN)}:\\
In any case, confirmation is necessary.

\textbf{Example (JA)}:\\
「どちらにしても、(確認:かくにん)が(必要:ひつよう)です」\\
\textbf{Example (EN)}:\\
Either way, we need to confirm it.

\textbf{Notes (JA)}:\\
「どちらにしても」は、どちらを選んでも結局同じ点が成り立つことを示す表現です。話をまとめたり、論点を絞ったり、だらだら続くやり取りを切り上げたりするときに便利です。会議などでも、長くてつまらない議論を、結局何が必要なのかという点に引き戻す働きをすることがあります。

\textbf{Notes (EN)}:\\
"Dochira ni shitemo" is used when the speaker wants to say that the same point holds regardless of which option is chosen. It is useful for summarizing a discussion, narrowing the focus, or moving people away from unproductive back-and-forth. In meetings, for example, it can help cut through a long and boring discussion by shifting attention to what still needs to be done anyway.

\newpage

\section*{58. Ultimately / In the End}

\textbf{Date}: 2024.06.07 (Fri)\\
\textbf{Index}: saishuutekiniwa\\
\textbf{Expression (JA)}: (最終的:さいしゅうてき)には\\
\textbf{Romanization}: Saishuuteki ni wa\\
\textbf{Expression (EN)}: ultimately / in the end

\textbf{Variants (JA)}: (最終的:さいしゅうてき)には / (結局:けっきょく) / (最後:さいご)には\\
\textbf{Variants (EN)}: ultimately / in the end / after all

\textbf{Intent (JA)}: 結論の提示, 結果のまとめ, 決定の表現\\
\textbf{Intent (EN)}: stating a conclusion, summing up a result, expression of decision

\textbf{Tags (JA)}: 即時文法, 口語, 結論提示表現, 結果表現\\
\textbf{Tags (EN)}: immediate grammar, spoken, conclusion-marking expression, result expression

\textbf{Adjusted Expression (JA)}:\\
(最終的:さいしゅうてき)な(判断:はんだん)としては、そのように(決:き)まりました。\\
\textbf{Adjusted Expression (EN)}:\\
As the final decision, that is what was decided.

\textbf{Example (JA)}:\\
「(最終的:さいしゅうてき)には、みんな(賛成:さんせい)しました。」\\
\textbf{Example (EN)}:\\
In the end, everyone agreed.

\textbf{Notes (JA)}:\\
「最終的には」は、いろいろな経過や議論のあとで、最後にどうなったかを示す表現です。ただし、この表現だけで、全員が納得して賛成したことまでは意味しません。実際には、何らかの形で決着がついたことを示すだけで、発言力の強い人の判断が通っただけという場合もあります。

\textbf{Notes (EN)}:\\
"Saishuuteki ni wa" is used when the speaker presents the final outcome after a process, discussion, or series of events. However, this expression by itself does not guarantee that everyone agreed willingly. In actual situations, it may simply mean that a decision was reached, sometimes because someone with stronger influence pushed it through.

\newpage

\section*{57. There's Nothing We Can Do}

\textbf{Date}: 2024.06.06 (Thu)\\
\textbf{Index}: dounimonaranai\\
\textbf{Expression (JA)}: どうにもならない。\\
\textbf{Romanization}: Dou ni mo naranai.\\
\textbf{Expression (EN)}: There's nothing we can do.

\textbf{Variants (JA)}: どうにもならない。 / もうどうにもならない。 / どうしようもない。\\
\textbf{Variants (EN)}: There's nothing we can do. / There's nothing I can do. / It can't be helped.

\textbf{Intent (JA)}: 無力感, あきらめ, 状況受容\\
\textbf{Intent (EN)}: helplessness, resignation, accepting the situation

\textbf{Tags (JA)}: 即時文法, 口語, あきらめの表現, 無力感の表現\\
\textbf{Tags (EN)}: immediate grammar, spoken, expression of resignation, expression of helplessness

\textbf{Adjusted Expression (JA)}:\\
この状況では、今のところどうすることもできません。\\
\textbf{Adjusted Expression (EN)}:\\
In this situation, there is nothing we can do for now.

\textbf{Example (JA)}:\\
A「(雨:あめ)が(降:ふ)っているから、(洗濯:せんたく)が(干:ほ)せない」 B「こればかりは、どうにもならないね」\\
\textbf{Example (EN)}:\\
A: It's raining, so I can't hang out the laundry. B: In this case, there's nothing we can do.

\textbf{Notes (JA)}:\\
「どうにもならない」は、何をしても状況を変えられないと感じるときに使う表現です。無力感、あきらめ、自分の力では及ばないという気持ちを表します。会話では、困ったことが起きているが、今は受け入れるしかないという即時的な反応として出やすい表現です。「こればかりは」を文頭に付けると、どうにもならない状況であることを強調できます。

\textbf{Notes (EN)}:\\
"Dou ni mo naranai" is used when the speaker feels that the situation cannot be changed no matter what they do. It expresses helplessness, resignation, or the sense that things are beyond one's control. In conversation, it often appears as an immediate reaction to an unpleasant but unavoidable situation. "kore bakari wa" can be added at the beginning to emphasize that this is a situation where nothing can be done.

\newpage

\section*{56. In the End}

\textbf{Date}: 2024.06.05 (Wed)\\
\textbf{Index}: kekkyoku\\
\textbf{Expression (JA)}: (結局:けっきょく)\\
\textbf{Romanization}: Kekkyoku\\
\textbf{Expression (EN)}: in the end

\textbf{Variants (JA)}: (結局:けっきょく) / (結局:けっきょく)のところ / (最後:さいご)には\\
\textbf{Variants (EN)}: in the end / after all / eventually

\textbf{Intent (JA)}: 結論の提示, 結果のまとめ, 見通しの表現\\
\textbf{Intent (EN)}: stating a conclusion, summing up a result, expression of outcome

\textbf{Tags (JA)}: 即時文法, 口語, 結論提示表現, 結果表現\\
\textbf{Tags (EN)}: immediate grammar, spoken, conclusion-marking expression, result expression

\textbf{Adjusted Expression (JA)}:\\
(最終的:さいしゅうてき)には、彼が来ることになりました。\\
\textbf{Adjusted Expression (EN)}:\\
In the end, it was decided that he would come.

\textbf{Example (JA)}:\\
A「(結局:けっきょく)、(彼:かれ)が(来:く)ることになりました」\\
\textbf{Example (EN)}:\\
A: In the end, it was decided that he would come.

\textbf{Notes (JA)}:\\
「結局」は、いろいろな成り行きや可能性があったあとで、最終的な結果を示すときに使う表現です。単に「最後に」というだけでなく、「あれこれあったが、結局こうなった」というまとめや帰着の感じを帯びやすい表現です。

\textbf{Notes (EN)}:\\
"Kekkyoku" is used when the speaker presents the final outcome after various possibilities, expectations, or developments. It often carries the feeling that things turned out that way after all, sometimes with a nuance of summary, resignation, or practical conclusion.

\newpage

\section*{55. Anyway / In Any Case}

\textbf{Date}: 2024.06.04 (Tue)\\
\textbf{Index}: douyara\\
\textbf{Expression (JA)}: どちみち\\
\textbf{Romanization}: Dochimichi\\
\textbf{Expression (EN)}: anyway / in any case

\textbf{Variants (JA)}: どのみち / いずれにせよ / とにかく\\
\textbf{Variants (EN)}: anyway / in any case / either way

\textbf{Intent (JA)}: 結論の提示, あきらめ, 状況整理\\
\textbf{Intent (EN)}: stating the outcome, resignation, sorting out the situation

\textbf{Tags (JA)}: 即時文法, 口語, 結論表現, 見通し表現\\
\textbf{Tags (EN)}: immediate grammar, spoken, conclusion-marking expression, outlook expression

\textbf{Adjusted Expression (JA)}:\\
どのみち、結果は変わらないと思います。\\
\textbf{Adjusted Expression (EN)}:\\
In any case, I think the result will not change.

\textbf{Example (JA)}:\\
A「(僕:ぼく)は(宿題:しゅくだい)はしないつもりだ」 B「そうだね、どちみち、(単位:たんい)は(取:と)れないだろう」\\
\textbf{Example (EN)}:\\
A: I don't plan to do my homework. B: Right, either way, you probably won't get the credits.

\textbf{Notes (JA)}:\\
「どちみち」は、「どのみち」「いずれにせよ」「どちらにしても」に近い意味をもつ口語表現です。どちらへ進んでも結局は同じだ、結果は変わらない、という見通しを出すときに使われます。会話では、少しあきらめた感じや、どうせそうなるという実際的な判断を帯びることがあります。「どのみち」という形もよく使われ、こちらの方がやや標準的に感じられることもあります。

\textbf{Notes (EN)}:\\
"Dochimichi" means "anyway," "either way," or "in any case." It is used when the speaker sees the outcome as essentially unchanged no matter which course is taken. In conversation, it can carry a nuance of resignation, inevitability, or rough practical judgment. The form "donomichi" is also common and is often felt to be slightly more standard.

\newpage

\section*{54. Not Really}

\textbf{Date}: 2024.06.03 (Mon)\\
\textbf{Index}: soudemonai\\
\textbf{Expression (JA)}: そうでもない\\
\textbf{Romanization}: Sou demo nai\\
\textbf{Expression (EN)}: not really

\textbf{Variants (JA)}: そうでもない / そうでもないよ / べつにそうでもない\\
\textbf{Variants (EN)}: not really / not that much / it's not really like that

\textbf{Intent (JA)}: やわらかい否定, 控えめな応答, 状況の軽減\\
\textbf{Intent (EN)}: soft denial, modest response, downplaying a situation

\textbf{Tags (JA)}: 即時文法, 口語, 応答表現, やわらかい否定表現\\
\textbf{Tags (EN)}: immediate grammar, spoken, response expression, soft denial expression

\textbf{Adjusted Expression (JA)}:\\
それほど(大変:たいへん)ではありません。\\
\textbf{Adjusted Expression (EN)}:\\
It is not that tough.

\textbf{Example (JA)}:\\
A「(大変:たいへん)だね」 B「そうでもないよ」\\
\textbf{Example (EN)}:\\
A: That sounds tough. B: Not really.

\textbf{Notes (JA)}:\\
「そうでもない」は、相手が今言ったことを、そのまま強く否定するのではなく、やわらかく弱めたり、少し引いた形で受け直したりするときに使う表現です。大変だ、すごい、深刻だ、などと言われたことに対して、「そこまでではない」と返す感じが中心です。言い方や文脈によっては、控えめさ、気づかい、強がり、空元気のような響きを帯びることもあります。

\textbf{Notes (EN)}:\\
"Sou demo nai" is used when the speaker responds by softening, downplaying, or gently denying what has just been said. It often appears when someone says something was hard, serious, great, or extreme, and the speaker answers that it was not quite so. Depending on tone and context, it can sound modest, reassuring, reserved, or slightly bluffing.

\newpage

\section*{53. I'll Give It a Try}

\textbf{Date}: 2024.06.02 (Sun)\\
\textbf{Index}: yattemimasu\\
\textbf{Expression (JA)}: やってみます。\\
\textbf{Romanization}: Yatte mimasu.\\
\textbf{Expression (EN)}: I'll give it a try.

\textbf{Variants (JA)}: やってみます。 / ちょっとやってみます。 / まずはやってみます。\\
\textbf{Variants (EN)}: I'll give it a try. / I'll try it. / I'll go ahead and try.

\textbf{Intent (JA)}: 試行, 意欲, 前向きな応答\\
\textbf{Intent (EN)}: trying, willingness, positive response

\textbf{Tags (JA)}: 即時文法, 口語, 応答表現, 試行表現\\
\textbf{Tags (EN)}: immediate grammar, spoken, response expression, trial expression

\textbf{Adjusted Expression (JA)}:\\
まずは一度、試しにやってみます。\\
\textbf{Adjusted Expression (EN)}:\\
First, I will try doing it once and see how it goes.

\textbf{Example (JA)}:\\
A「この(料理:りょうり)、(作:つく)ってみる？」 B「ええ、やってみます」\\
\textbf{Example (EN)}:\\
A: Do you want to try making this dish? B: Yes, I'll give it a try.

\textbf{Notes (JA)}:\\
「やってみます」は、何かを試しにやってみようという気持ちを表す表現です。提案や誘い、ちょっとした挑戦に対して前向きに応じるときによく使われます。うまくいくことを断言するのではなく、まずはやってみようという姿勢を示す言い方です。

\textbf{Notes (EN)}:\\
"Yatte mimasu" is used when the speaker expresses a willingness to try doing something. It often appears as a positive response to a suggestion, invitation, or challenge. Rather than guaranteeing success, it conveys a readiness to attempt something and see how it goes.

\newpage

\section*{52. I Can't Make Sense of It}

\textbf{Date}: 2024.06.01 (Sat)\\
\textbf{Index}: naniganandaka\\
\textbf{Expression (JA)}: (何:なに)が(何:なん)だかわからない。\\
\textbf{Romanization}: Nani ga nan da ka wakaranai.\\
\textbf{Expression (EN)}: I can't make sense of what's what.

\textbf{Variants (JA)}: (何:なに)が(何:なん)だかわからない。 / もう(何:なに)が(何:なん)だかわからない。 / (何:なに)がどうなっているのかわからない。\\
\textbf{Variants (EN)}: I can't make sense of what's what. / I have no idea what's going on. / I can't tell what's what anymore.

\textbf{Intent (JA)}: 混乱, 理解困難, 状況把握の失敗\\
\textbf{Intent (EN)}: confusion, difficulty in understanding, failure to grasp the situation

\textbf{Tags (JA)}: 即時文法, 口語, 混乱表現, 理解困難表現\\
\textbf{Tags (EN)}: immediate grammar, spoken, confusion expression, difficulty-understanding expression

\textbf{Adjusted Expression (JA)}:\\
(状況:じょうきょう)が(複雑:ふくざつ)で、(何:なに)がどうなっているのか理解できません。\\
\textbf{Adjusted Expression (EN)}:\\
The situation is so complicated that I cannot understand what is going on.

\textbf{Example (JA)}:\\
(彼:かれ)が(説明:せつめい)してくれるけど、あまりにも(速:はや)すぎて、(何:なに)が(何:なん)だかわからない。\\
\textbf{Example (EN)}:\\
He explains it to me, but he speaks so fast that I can't make sense of what's what.

\textbf{Notes (JA)}:\\
「何が何だかわからない」は、物事が入り組んでいて、何がどうなっているのか頭の中で整理できないときに使う表現です。説明が速すぎるとき、情報が多すぎるとき、状況が複雑すぎるときなどに出やすく、単に理解していないというだけでなく、その場で混乱している感じを伴いやすい言い方です。

\textbf{Notes (EN)}:\\
"Nani ga nan da ka wakaranai" is used when the speaker feels confused and cannot sort out what is what in the situation. It often appears when too many things are happening at once, when an explanation is hard to follow, or when the speaker feels mentally overwhelmed. Rather than a calm statement of non-understanding, it often carries the feeling of immediate confusion.

\newpage

\section*{51. Oh, I Didn't Know That}

\textbf{Date}: 2024.05.31 (Fri)\\
\textbf{Index}: soudeshitaka\\
\textbf{Expression (JA)}: そうでしたか。\\
\textbf{Romanization}: Sou deshita ka.\\
\textbf{Expression (EN)}: Oh, I didn't know that.

\textbf{Variants (JA)}: そうだったんですか。 / そうですか。 / そうでしたか。\\
\textbf{Variants (EN)}: Oh, I didn't know that. / I see. / Oh, is that so?

\textbf{Intent (JA)}: 理解, 気づき, 軽い驚き, 丁寧な反応\\
\textbf{Intent (EN)}: understanding, realization, mild surprise, polite response

\textbf{Tags (JA)}: 即時文法, 口語, 反応表現, 丁寧表現\\
\textbf{Tags (EN)}: immediate grammar, spoken, reaction expression, polite expression

\textbf{Adjusted Expression (JA)}:\\
それは知りませんでした。\\
\textbf{Adjusted Expression (EN)}:\\
I did not know that.

\textbf{Example (JA)}:\\
A「(先生:せんせい)、(来週:らいしゅう)は(休講:きゅうこう)です。」 B「そうでしたか。」\\
\textbf{Example (EN)}:\\
A: The class will be canceled next week. B: Oh, I didn't know that.

\textbf{Notes (JA)}:\\
「そうでしたか」は、今はじめて知ったことに対して、理解や軽い驚きをこめて返す丁寧な反応表現です。過去形になっていることで、相手の説明を受けて、そういうことだったのかと受け止める感じが出ます。「そっか」のようなくだけた反応よりも、やわらかく丁寧な響きを持つ表現です。

\textbf{Notes (EN)}:\\
"Sou deshita ka" is a polite reaction expression used when the speaker has just learned something new and responds with understanding, mild surprise, or acknowledgment. The past form gives it a reflective nuance, as if the speaker is taking in what has just been explained. It is often softer and more polite than a casual response such as "sokka."

\newpage

\section*{50. Excuse Me for Leaving First}

\textbf{Date}: 2024.05.30 (Thu)\\
\textbf{Index}: osakini\\
\textbf{Expression (JA)}: お(先:さき)に(失礼:しつれい)します。\\
\textbf{Romanization}: Osaki ni shitsurei shimasu.\\
\textbf{Expression (EN)}: Excuse me for leaving before you.

\textbf{Variants (JA)}: お(先:さき)に(失礼:しつれい)します。 / お(先:さき)に(上:あ)がります。 / お(先:さき)に(失礼:しつれい)いたします。\\
\textbf{Variants (EN)}: Excuse me for leaving first. / I'll be leaving now. / Please excuse me for heading out before you.

\textbf{Intent (JA)}: 退出, あいさつ, 丁寧表現\\
\textbf{Intent (EN)}: leaving, greeting, polite expression

\textbf{Tags (JA)}: 定型表現, 口語, 職場表現, 別れの表現\\
\textbf{Tags (EN)}: formulaic expression, spoken, workplace expression, parting expression

\textbf{Adjusted Expression (JA)}:\\
それでは、お(先:さき)に(失礼:しつれい)いたします。\\
\textbf{Adjusted Expression (EN)}:\\
Well then, please excuse me for leaving before you.

\textbf{Example (JA)}:\\
それでは、お(先:さき)に(失礼:しつれい)します。\\
\textbf{Example (EN)}:\\
Well then, please excuse me for leaving before you.

\textbf{Notes (JA)}:\\
「お先に失礼します」は、ほかの人がまだ残っている場面で、自分が先に退出するときに使う丁寧な定型表現です。一般的な「さようなら」とは少し違い、「先に失礼します」という事情つきの別れのあいさつです。職場では特によく使われ、「お疲れさまです/でした」とは役割や含みが異なります。

\textbf{Notes (EN)}:\\
"Osaki ni shitsurei shimasu" is a polite fixed expression used when the speaker leaves while others are still staying, especially in workplaces. It is not a general equivalent of "goodbye," but a more specific expression meaning that the speaker is excusing themselves for leaving first. A related workplace expression is "otsukaresama desu/deshita," but that has a different function and nuance.

\newpage

\section*{49. Thanks to You}

\textbf{Date}: 2024.05.29 (Wed)\\
\textbf{Index}: okagesamade\\
\textbf{Expression (JA)}: おかげさまで\\
\textbf{Romanization}: Okage-sama de\\
\textbf{Expression (EN)}: Thanks to you, I'm doing well.

\textbf{Variants (JA)}: おかげさまで、(元気:げんき)です。 / おかげさまで、なんとかやっています。\\
\textbf{Variants (EN)}: Thanks to you, I'm doing well. / Thanks to your support, I'm doing fine. / I'm doing well, thanks to you.

\textbf{Intent (JA)}: 感謝, 近況応答, 丁寧表現\\
\textbf{Intent (EN)}: gratitude, responding about one's condition, polite expression

\textbf{Tags (JA)}: 定型表現, 口語, 感謝表現, あいさつ表現\\
\textbf{Tags (EN)}: formulaic expression, spoken, gratitude expression, greeting expression

\textbf{Adjusted Expression (JA)}:\\
皆さまのおかげで、元気に過ごしております。\\
\textbf{Adjusted Expression (EN)}:\\
Thanks to everyone's support, I have been doing well.

\textbf{Example (JA)}:\\
A「お(元気:げんき)ですか。」 B「おかげさまで、(元気:げんき)です。」\\
\textbf{Example (EN)}:\\
A: How have you been? B: Thanks to you, I'm doing well.

\textbf{Notes (JA)}:\\
「おかげさまで」は、相手への感謝をにじませながら、自分の近況を丁寧に伝えるときによく使う定型表現です。単に「ありがとう」というだけではなく、今の自分の状態が他者の支えや配慮によるものだという含みを持っています。そのため、英語の一つの決まった語にそのまま対応させにくい表現です。

\textbf{Notes (EN)}:\\
"Okage-sama de" is a polite formulaic expression used to respond with gratitude, often when someone asks how you are doing. It does not simply mean "thanks," but carries the sense that one's present well-being is due to the kindness, support, or consideration of others. Because of this, it does not map neatly onto a single fixed English phrase.

\newpage

\section*{48. Really?}

\textbf{Date}: 2024.05.28 (Tue)\\
\textbf{Index}: hontou\\
\textbf{Expression (JA)}: ほんとー\\
\textbf{Romanization}: Hontoo\\
\textbf{Expression (EN)}: Really?

\textbf{Variants (JA)}: ほんと？ / ほんとう？ / ほんとー？\\
\textbf{Variants (EN)}: Really? / Is that true? / For real?

\textbf{Intent (JA)}: 驚き, 関心, 軽い不信\\
\textbf{Intent (EN)}: surprise, interest, mild disbelief

\textbf{Tags (JA)}: 即時文法, 口語, 反応表現, あいづち表現\\
\textbf{Tags (EN)}: immediate grammar, spoken, reaction expression, backchannel expression

\textbf{Adjusted Expression (JA)}:\\
(本当:ほんとう)ですか。\\
\textbf{Adjusted Expression (EN)}:\\
Is that really true?

\textbf{Example (JA)}:\\
A「その人、独学で10か国語できるんだって」 B「ほんとー」\\
\textbf{Example (EN)}:\\
A: They say that person taught themselves ten languages. B: Really?

\textbf{Notes (JA)}:\\
「ほんとー」は、驚き、関心、軽い不信をやわらかく表す口語的な反応表現です。強い不信をぶつけるというより、相手の話を受けて「へえ、ほんとう？」と返すような、聞き手の即時的な応答として出やすい表現です。私が子どものころ、母はいつもこう言って私の話を聞いてくれた。それがうれしくて、私は図書館へ行って何か新しいことを勉強し、母に教えたものだった。

\textbf{Notes (EN)}:\\
"Hontoo" is a casual reaction used when the speaker shows surprise, interest, or mild disbelief. It is often lighter and softer than a strong expression of disbelief, and it can function as a quick listener response. When I was a child, my mother often said this to me, and I was happy to hear it, so I would go to the library, learn something new, and tell it to her.

\newpage

\section*{47. You're Kidding!}

\textbf{Date}: 2024.05.27 (Mon)\\
\textbf{Index}: usodaa\\
\textbf{Expression (JA)}: うそだぁ～\\
\textbf{Romanization}: Uso daa\\textasciitilde{}\\
\textbf{Expression (EN)}: You're kidding!

\textbf{Variants (JA)}: うそだあ。 / うそでしょ？ / まさか！\\
\textbf{Variants (EN)}: You're kidding! / No way! / You can't be serious!

\textbf{Intent (JA)}: 不信, 驚き, 反応\\
\textbf{Intent (EN)}: disbelief, surprise, reaction

\textbf{Tags (JA)}: 即時文法, 口語, 反応表現, 不信表現\\
\textbf{Tags (EN)}: immediate grammar, spoken, reaction expression, disbelief expression

\textbf{Adjusted Expression (JA)}:\\
それはうそでしょう。\\
\textbf{Adjusted Expression (EN)}:\\
That cannot be true.

\textbf{Example (JA)}:\\
A「(田中:たなか)くんのお(父:とう)さんは(ゴジラ:ごじら)なんだって」 B「うそだぁ～」\\
\textbf{Example (EN)}:\\
A: I heard that Tanaka's father is Godzilla. B: You're kidding!

\textbf{Notes (JA)}:\\
「うそだぁ～」は、何かを聞いたときに、信じられない気持ちや驚きをその場で表す口語的な反応表現です。落ち着いて真偽を判断するというより、まず感情が先に出る即時的な言い方です。

\textbf{Notes (EN)}:\\
"Uso daa\\textasciitilde{}" is a colloquial reaction used when the speaker hears something and responds with disbelief or surprise. Rather than calmly judging whether something is true, it comes out as an immediate emotional response.

\newpage

\section*{46. Two thumbs up}

\textbf{Date}: 2024.05.26 (Sun)\\
\textbf{Index}: iine\\
\textbf{Expression (JA)}: いいね\\
\textbf{Romanization}: Ii ne\\
\textbf{Expression (EN)}: Nice

\textbf{Variants (JA)}: いいね。 / いいですね。 / それ、いいね。\\
\textbf{Variants (EN)}: That's good. / Two thumbs up. / That sounds good.

\textbf{Intent (JA)}: 賛成, 好意的評価, 共感\\
\textbf{Intent (EN)}: agreement, positive evaluation, shared feeling

\textbf{Tags (JA)}: 即時文法, 口語, 評価表現, 同意表現\\
\textbf{Tags (EN)}: immediate grammar, spoken, evaluative expression, agreement expression

\textbf{Adjusted Expression (JA)}:\\
それはよいと(思:おも)います。\\
\textbf{Adjusted Expression (EN)}:\\
I think that is good.

\textbf{Example (JA)}:\\
A「この(店:みせ)、(行:い)ってみない？」 B「いいね」\\
\textbf{Example (EN)}:\\
A: Why don't we try this place? B: Nice.

\textbf{Notes (JA)}:\\
「いいね」は、その場で賛成したり、よいと感じたり、相手の案や発言に好意的に反応したりするときによく使う口語表現です。「いいね」は、勢いよく前向きな感情の「Two thumbs up」の訳として、使われることもあります。この表現は、何かが好きだったり、同意したりすることを表すために、さまざまな文脈で使用できる便利な表現です。

\textbf{Notes (EN)}:\\
"Ii ne" is a common spoken expression used to show approval, agreement, or a favorable reaction on the spot. It can also be used as a translation of "Two thumbs up," which conveys a strong positive feeling. The phrase is versatile and can be used in various contexts to express that you like something or agree with it.

\newpage

\section*{46. That's good}

\textbf{Date}: 2024.05.26 (Sun)\\
\textbf{Index}: iine\\
\textbf{Expression (JA)}: いいね\\
\textbf{Romanization}: Iine\\
\textbf{Expression (EN)}: That's good

\textbf{Variants (JA)}: いいね / いいねえ / それ、いいね\\
\textbf{Variants (EN)}: That's good. / Sounds good. / That's nice.

\textbf{Intent (JA)}: 賛成の表現, 肯定的な反応, 軽い評価の共有\\
\textbf{Intent (EN)}: expressing agreement, giving a positive reaction, sharing a light evaluation

\textbf{Tags (JA)}: 即時文法, 口語, 応答表現, 肯定, 評価表現\\
\textbf{Tags (EN)}: immediate grammar, colloquial, response expression, affirmation, evaluative expression

\textbf{Adjusted Expression (JA)}:\\
それは、いいですね。\\
\textbf{Adjusted Expression (EN)}:\\
That is good.

\textbf{Example (JA)}:\\
\\
\textbf{Example (EN)}:\\


\textbf{Notes (JA)}:\\
「いいね」は、賛成や同意を表現するための表現で、多くの状況で使用できます。

\textbf{Notes (EN)}:\\
The word "iine" is an expression of approval or agreement. It is a versatile word that can be used in many situations.

\newpage

\section*{45. Oh, That's It?}

\textbf{Date}: 2024.05.25 (Sat)\\
\textbf{Index}: naannda\\
\textbf{Expression (JA)}: なぁんだ\\
\textbf{Romanization}: Naanda\\
\textbf{Expression (EN)}: Oh, that's it?

\textbf{Variants (JA)}: なんだ。 / なーんだ。 / なんだ、そういうことか。\\
\textbf{Variants (EN)}: Oh, that's it? / Oh, I see. / So that was it.

\textbf{Intent (JA)}: 気づき, 拍子抜け, 反応\\
\textbf{Intent (EN)}: realization, anticlimax, reaction

\textbf{Tags (JA)}: 即時文法, 口語, 反応表現, 気づきの表現\\
\textbf{Tags (EN)}: immediate grammar, spoken, reaction expression, realization expression

\textbf{Adjusted Expression (JA)}:\\
そういうことだったのですね。\\
\textbf{Adjusted Expression (EN)}:\\
I see. So that was what it was.

\textbf{Example (JA)}:\\
A「(リモコン:りもこん)が(見:み)つからないと(思:おも)ったら、(ソファ:そふぁ)の(下:した)にあった」 B「なぁんだ、そんなところにあったのか」\\
\textbf{Example (EN)}:\\
A: I thought I had lost the remote, but it was under the sofa. B: Oh, that's it? So it was there.

\textbf{Notes (JA)}:\\
「なぁんだ」は、事情がわかって、思ったより大したことがなかった、案外そんなことだったのか、という気持ちをその場で表す間投詞です。単なる理解だけでなく、少し拍子抜けした感じや、ほっとした感じを伴いやすい表現です。「そっか」が納得に寄るのに対して、「なぁんだ」は「なんだ、その程度か」という軽い落差を帯びやすい言い方です。

\textbf{Notes (EN)}:\\
"Naanda" is an interjection used when something becomes clear and the speaker reacts with a sense of anticlimax, mild surprise, or relief. It often comes out when the answer turns out to be simpler, smaller, or less serious than expected. Compared with "sokka," it more easily carries the feeling of "Oh, it was only that."

\newpage

\section*{44. So-So}

\textbf{Date}: 2024.05.24 (Fri)\\
\textbf{Index}: maamaa\\
\textbf{Expression (JA)}: まあまあ\\
\textbf{Romanization}: Maamaa\\
\textbf{Expression (EN)}: So-so.

\textbf{Variants (JA)}: まあまあだね。 / まあ、(悪:わる)くないよ。 / そんなに(悪:わる)くなかったよ。\\
\textbf{Variants (EN)}: So-so. / Not bad. / It was okay.

\textbf{Intent (JA)}: 控えめな評価, 反応, やわらかい返答\\
\textbf{Intent (EN)}: moderate evaluation, reaction, softened reply

\textbf{Tags (JA)}: 即時文法, 口語, 評価表現, あいまいな返答\\
\textbf{Tags (EN)}: immediate grammar, spoken, evaluative expression, ambiguous reply

\textbf{Adjusted Expression (JA)}:\\
(特:とく)にすごく(良:よ)かったわけではありませんが、(悪:わる)くもありませんでした。\\
\textbf{Adjusted Expression (EN)}:\\
It was not especially good, but it was not bad either.

\textbf{Example (JA)}:\\
A「(試験:しけん)、どうだった？」 B「まあまあ、(難:むずか)しかったけど、なんとかできたよ。」\\
\textbf{Example (EN)}:\\
A: How was the exam? B: So-so. It was difficult, but I managed.

\textbf{Notes (JA)}:\\
「まあまあ」は、何かについて控えめで断定しすぎない評価をするときによく使う口語表現です。特にすごく良いわけでも悪いわけでもない、という中くらいの感じを表しやすい言い方です。口調や文脈によっては、やや肯定的にも、どこか歯切れの悪い返答にも聞こえることがあります。

\textbf{Notes (EN)}:\\
"Maamaa" is a common spoken expression used when the speaker gives a moderate or noncommittal evaluation. It often means that something was neither especially good nor especially bad. Depending on tone and context, it can sound mildly positive, lukewarm, or slightly withholding.

\newpage

\section*{43. Oh, I See!}

\textbf{Date}: 2024.05.23 (Thu)\\
\textbf{Index}: sokka\\
\textbf{Expression (JA)}: そっか！\\
\textbf{Romanization}: Sokka!\\
\textbf{Expression (EN)}: Oh, I see!

\textbf{Variants (JA)}: そうか。 / そっか。 / そっか、なるほど。\\
\textbf{Variants (EN)}: Oh, I see. / I see now. / Oh, right.

\textbf{Intent (JA)}: 理解, 気づき, 反応\\
\textbf{Intent (EN)}: understanding, realization, reaction

\textbf{Tags (JA)}: 即時文法, 口語, 反応表現, 気づきの表現\\
\textbf{Tags (EN)}: immediate grammar, spoken, reaction expression, realization expression

\textbf{Adjusted Expression (JA)}:\\
そういうことなのですね。\\
\textbf{Adjusted Expression (EN)}:\\
I see. So that is how it is.

\textbf{Example (JA)}:\\
A「(明日:あした)は(休:やす)みだよ」 B「そっか！」\\
\textbf{Example (EN)}:\\
A: Tomorrow is a holiday. B: Oh, I see!

\textbf{Notes (JA)}:\\
「そっか」は、「そうか」がくだけた形になった口語表現です。相手の話を聞いて、なるほどと理解したり、そういうことかと気づいたりしたときに使われます。「へえー」が驚きや関心の反応になりやすいのに対して、「そっか」は、話の内容が自分の中でつながって納得した感じを表しやすい表現です。

\textbf{Notes (EN)}:\\
"Sokka" is a colloquial contracted form of "sou ka." It is used when the speaker understands something, realizes something, or sees how a situation fits together. Compared with "hee," it shows not just surprise or interest, but a sense that the information has now made sense to the speaker.

\newpage

\section*{42. Hmm, I See}

\textbf{Date}: 2024.05.22 (Wed)\\
\textbf{Index}: fuun\\
\textbf{Expression (JA)}: ふーん\\
\textbf{Romanization}: Fuun\\
\textbf{Expression (EN)}: Hmm, I see.

\textbf{Variants (JA)}: ふーん。 / ふうん。 / ふーん、そうなんだ。\\
\textbf{Variants (EN)}: Hmm. / I see. / Hmm, is that so?

\textbf{Intent (JA)}: 反応, 関心, 受け流し\\
\textbf{Intent (EN)}: reaction, interest, brushing off

\textbf{Tags (JA)}: 即時文法, 口語, あいづち表現, 反応表現\\
\textbf{Tags (EN)}: immediate grammar, spoken, backchannel expression, reaction expression

\textbf{Adjusted Expression (JA)}:\\
そうなんですね。\\
\textbf{Adjusted Expression (EN)}:\\
I see. Is that so?

\textbf{Example (JA)}:\\
A「その人、毎朝4時に起きるんだって」 B「ふーん」\\
\textbf{Example (EN)}:\\
A: They say that person gets up at four every morning. B: Hmm, I see.

\textbf{Notes (JA)}:\\
「ふーん」は、聞き手がその場で返す即時的な反応表現です。軽い関心や了解を表すこともありますが、言い方や文脈によっては、あまり強く感心していない感じや、少し距離を置いた受け止め方にもなります。「へえー」よりも平らで、冷めた響きになることもある表現です。

\textbf{Notes (EN)}:\\
"Fuun" is an interjection used as an immediate listener response. It can show mild interest, acknowledgment, or emotional distance, depending on intonation and context. Unlike "hee," it often sounds flatter and may even suggest that the speaker is not especially impressed.

\newpage

\section*{41. Oh, Really?}

\textbf{Date}: 2024.05.21 (Tue)\\
\textbf{Index}: he-\\
\textbf{Expression (JA)}: へえー\\
\textbf{Romanization}: Hee\\
\textbf{Expression (EN)}: Oh, really?

\textbf{Variants (JA)}: へえ。 / へえー。 / へえ、そうなんだ。\\
\textbf{Variants (EN)}: Oh, really? / Wow. / Is that so?

\textbf{Intent (JA)}: 驚き, 関心, あいづち\\
\textbf{Intent (EN)}: surprise, interest, backchannel response

\textbf{Tags (JA)}: 即時文法, 口語, あいづち表現, 反応表現\\
\textbf{Tags (EN)}: immediate grammar, spoken, backchannel expression, reaction expression

\textbf{Adjusted Expression (JA)}:\\
そうなんですか。\\
\textbf{Adjusted Expression (EN)}:\\
Oh, is that so?

\textbf{Example (JA)}:\\
A「その人、20か国語話せるんだって」 B「へえー」\\
\textbf{Example (EN)}:\\
A: They say that person can speak twenty languages. B: Oh, really?

\textbf{Notes (JA)}:\\
「へえー」は、驚き、関心、軽い感心などを表す間投詞です。実際の会話では、強い感嘆というより、相手の話を受けるときの即時的な反応として使われることが多い表現です。「わあ」のような大きな驚きよりも、むしろ「そうなんだ」「へえ、そうなのか」に近い聞き手の反応として自然に出やすい言い方です。

\textbf{Notes (EN)}:\\
"Hee" is an interjection used to show surprise, interest, or mild admiration. In actual conversation, it often works as a quick listener response rather than a strong exclamation. Compared with "wow," it is usually softer and closer to "oh, really?" or "is that so?" depending on the context.

\newpage

\section*{40. No Way!}

\textbf{Date}: 2024.05.20 (Mon)\\
\textbf{Index}: arien\\
\textbf{Expression (JA)}: ありえん！\\
\textbf{Romanization}: Arien!\\
\textbf{Expression (EN)}: No way!

\textbf{Variants (JA)}: ありえない！ / そんなのありえん！ / まさか！\\
\textbf{Variants (EN)}: No way! / That's impossible! / Unbelievable!

\textbf{Intent (JA)}: 不信, 驚き, 強い否定\\
\textbf{Intent (EN)}: disbelief, surprise, strong denial

\textbf{Tags (JA)}: 即時文法, 口語, 感情表現, 不信表現\\
\textbf{Tags (EN)}: immediate grammar, spoken, emotional expression, disbelief expression

\textbf{Adjusted Expression (JA)}:\\
それは、ありえないと(思:おも)います。\\
\textbf{Adjusted Expression (EN)}:\\
I think that is impossible.

\textbf{Example (JA)}:\\
3(回:かい)(連続:れんぞく)で(ジャンケン:じゃんけん)で(勝:か)った？ ありえん！\\
\textbf{Example (EN)}:\\
You won rock-paper-scissors three times in a row? No way!

\textbf{Notes (JA)}:\\
「ありえん」は、「ありえない」のくだけた口語的な形で、強い不信、驚き、信じられない気持ちをその場で表す表現です。丁寧に判断を述べるというより、まず感情が先に出るような即時的な言い方として使われやすい表現です。

\textbf{Notes (EN)}:\\
"Arien" is a colloquial spoken form of "arienai." It is used to express strong disbelief, surprise, or the feeling that something is beyond belief. In immediate use, it often comes out as a quick emotional reaction rather than as a carefully stated judgment.

\newpage

\section*{39. Everything}

\textbf{Date}: 2024.05.19 (Sun)\\
\textbf{Index}: zennbu\\
\textbf{Expression (JA)}: ぜんぶ\\
\textbf{Romanization}: Zenbu\\
\textbf{Expression (EN)}: everything

\textbf{Variants (JA)}: (全部:ぜんぶ) / すべて / (何:なに)もかも\\
\textbf{Variants (EN)}: everything / all of it / the whole thing

\textbf{Intent (JA)}: 包括, 強調, 全体の指示\\
\textbf{Intent (EN)}: inclusion, emphasis, referring to the whole

\textbf{Tags (JA)}: 語彙表現, 包括表現, 口語\\
\textbf{Tags (EN)}: lexical expression, inclusive expression, spoken

\textbf{Adjusted Expression (JA)}:\\
(関係:かんけい)するものを、全部まとめて(含:ふく)みます。\\
\textbf{Adjusted Expression (EN)}:\\
It includes everything related, all together.

\textbf{Example (JA)}:\\
これ、ぜんぶ(食:た)べていいの？\\
\textbf{Example (EN)}:\\
Can I eat all of this?

\textbf{Notes (JA)}:\\
「ぜんぶ」は、「すべて」「全部」「まるごと」を表す、日常会話で非常によく使われる語です。何かの全体や残りなく含まれることを言うときに使います。「すべて」よりもくだけた感じがあり、会話ではこちらの方が自然に出やすいことが多いです。

\textbf{Notes (EN)}:\\
"Zenbu" means "everything," "all," or "the whole thing." It is a very common everyday word used to refer to the whole of something. Compared with "subete," it often sounds more casual and more natural in conversation.

\newpage

\section*{38. You Know What?}

\textbf{Date}: 2024.05.18 (Sat)\\
\textbf{Index}: anone\\
\textbf{Expression (JA)}: あのね？\\
\textbf{Romanization}: Ano ne?\\
\textbf{Expression (EN)}: You know what?

\textbf{Variants (JA)}: あのね、 / あのさ、 / ねえ、あのね。\\
\textbf{Variants (EN)}: You know what? / Hey, listen. / So, you know...

\textbf{Intent (JA)}: 呼びかけ, 話し始め, 注意を向ける表現\\
\textbf{Intent (EN)}: calling someone, starting to speak, directing attention

\textbf{Tags (JA)}: 即時文法, 口語, 話し始めの表現, 注意喚起表現\\
\textbf{Tags (EN)}: immediate grammar, spoken, speech-opening expression, attention-directing expression

\textbf{Adjusted Expression (JA)}:\\
あのですね、ちょっとお(話:はな)ししたいことがあります。\\
\textbf{Adjusted Expression (EN)}:\\
Well, there is something I would like to tell you.

\textbf{Example (JA)}:\\
あのね、きょう(学校:がっこう)でね、おもしろいことがあったんだ。\\
\textbf{Example (EN)}:\\
You know what? Something interesting happened at school today.

\textbf{Notes (JA)}:\\
「あのね」は、これから何かを言い出す前に、相手の注意をやわらかく引き寄せる口語的な表現です。単に同意を求めるというより、話を始めるきっかけとして使われることが多く、個人的なことを伝えたり、ちょっとした出来事を話したりするときに自然に出やすい言い方です。

\textbf{Notes (EN)}:\\
"Ano ne" is a casual expression used to draw the listener in before saying something. Rather than simply seeking agreement, it often works as a soft way to begin speaking, especially when the speaker wants attention, wants to share something personal, or is about to tell a small story.

\newpage

\section*{37. Lovely!}

\textbf{Date}: 2024.05.17 (Fri)\\
\textbf{Index}: suteki\\
\textbf{Expression (JA)}: すてき！\\
\textbf{Romanization}: Suteki!\\
\textbf{Expression (EN)}: Lovely!

\textbf{Variants (JA)}: すてきですね。 / ほんとうにすてき！ / すてきだね。\\
\textbf{Variants (EN)}: Lovely! / That's lovely! / How nice!

\textbf{Intent (JA)}: 賞賛, 感動, 好意的評価\\
\textbf{Intent (EN)}: praise, admiration, positive evaluation

\textbf{Tags (JA)}: 即時文法, 口語, 評価表現, 感情表現\\
\textbf{Tags (EN)}: immediate grammar, spoken, evaluative expression, emotional expression

\textbf{Adjusted Expression (JA)}:\\
とてもすてきだと(思:おも)います。\\
\textbf{Adjusted Expression (EN)}:\\
I think it is very lovely.

\textbf{Example (JA)}:\\
その(服:ふく)、すてき！\\
\textbf{Example (EN)}:\\
That outfit is lovely!

\textbf{Notes (JA)}:\\
「すてき」は、好ましい感じや軽い感動を、その場でやわらかく表す表現です。「すばらしい」よりも少し軽く、個人的で親しみのある響きを持ち、見た目、雰囲気、センス、ちょっとした魅力をほめるときによく使われます。

\textbf{Notes (EN)}:\\
"Suteki" is used to express warm admiration or a favorable impression. Compared with "subarashii," it often sounds lighter, softer, and more personal, and it is frequently used for appearance, atmosphere, taste, or small moments that feel charming.

\newpage

\section*{36. Excellent}

\textbf{Date}: 2024.05.16 (Thu)\\
\textbf{Index}: sugoi\\
\textbf{Expression (JA)}: すごい！\\
\textbf{Romanization}: sugoi!\\
\textbf{Expression (EN)}: Excellent!

\textbf{Variants (JA)}: すごい！ / すごいですね。 / すごいなあ。\\
\textbf{Variants (EN)}: Excellent! / Amazing! / Great!

\textbf{Intent (JA)}: 感嘆, 賞賛, 驚き\\
\textbf{Intent (EN)}: admiration, praise, surprise

\textbf{Tags (JA)}: 即時文法, 口語, 感情表現, 評価表現\\
\textbf{Tags (EN)}: immediate grammar, spoken, emotional expression, evaluative expression

\textbf{Adjusted Expression (JA)}:\\
それはたいへんすばらしいと(思:おも)います。\\
\textbf{Adjusted Expression (EN)}:\\
I think that is excellent.

\textbf{Example (JA)}:\\
A「これなら(満足:まんぞく)？」 B「うん！(本当:ほんとう)にすごいよ！」\\
\textbf{Example (EN)}:\\
A: Is this satisfactory? B: Yes! It's really excellent!

\textbf{Notes (JA)}:\\
「すごい」は「すばらしい」や「素晴らしい」を意味します。感嘆や驚きを表現するために使用されます。 多目的なことばで、多くの状況で使用できます。

\textbf{Notes (EN)}:\\
The word "sugoi" means "amazing," "excellent," or "great." It is used to express admiration or surprise. It is a versatile word that can be used in many situations.

\newpage

\section*{35. When you are called your name}

\textbf{Date}: 2024.05.15 (Wed)\\
\textbf{Index}: naani\\
\textbf{Expression (JA)}: なあに？\\
\textbf{Romanization}: naani?\\
\textbf{Expression (EN)}: What?

\textbf{Variants (JA)}: なに？ / なあに。 / なに？どうしたの？\\
\textbf{Variants (EN)}: What? / Yes? / What is it?

\textbf{Intent (JA)}: 呼びかけへの応答, 確認, 会話の開始\\
\textbf{Intent (EN)}: response to being called, checking, starting interaction

\textbf{Tags (JA)}: 即時文法, 口語, 呼びかけ応答表現\\
\textbf{Tags (EN)}: immediate grammar, spoken, response-to-being-called expression

\textbf{Adjusted Expression (JA)}:\\
(何:なに)かご(用:よう)ですか。\\
\textbf{Adjusted Expression (EN)}:\\
Do you need me for something?

\textbf{Example (JA)}:\\
A「(田中:たなか)くん！」 B「なあに？」\\
\textbf{Example (EN)}:\\
A: Tanaka! B: What?

\textbf{Notes (JA)}:\\
「なあに」は「何？」や「何か用事でもありますか？」を意味する口語表現です。「なあに」は、名前を呼ばれたときに使用できます。

\textbf{Notes (EN)}:\\
"Naani" is a colloquial expression that means "What?" or "Do you need me for something?" It can be used when you are called by your name.

\newpage

\section*{34. Shall We Have Some Tea?}

\textbf{Date}: 2024.05.14 (Tue)\\
\textbf{Index}: ocha\\
\textbf{Expression (JA)}: お(茶:ちゃ)しませんか。\\
\textbf{Romanization}: Ocha shimasen ka.\\
\textbf{Expression (EN)}: Shall we have some tea?

\textbf{Variants (JA)}: ちょっとお(茶:ちゃ)しませんか。 / お(茶:ちゃ)でもしませんか。 / お(茶:ちゃ)しない？\\
\textbf{Variants (EN)}: How about some tea? / Would you like to have tea? / How about a tea break?

\textbf{Intent (JA)}: 誘い, 休憩, 会話のきっかけ\\
\textbf{Intent (EN)}: invitation, taking a break, starting a conversation

\textbf{Tags (JA)}: 即時文法, 口語, 誘い表現, 休憩表現\\
\textbf{Tags (EN)}: immediate grammar, spoken, invitation expression, break-time expression

\textbf{Adjusted Expression (JA)}:\\
もしよろしければ、お(茶:ちゃ)でも(飲:の)みませんか。\\
\textbf{Adjusted Expression (EN)}:\\
If you like, would you like to have some tea or something to drink?

\textbf{Example (JA)}:\\
(少:すこ)し(休:やす)みたいですね。お(茶:ちゃ)しませんか。\\
\textbf{Example (EN)}:\\
It looks like we could use a little break. Shall we have some tea?

\textbf{Notes (JA)}:\\
「お茶」は例であって本当に「お茶」を意味するわけではありません。「する」という動詞の丁寧の否定形ですが、これは否定ではなく、誘いです。「か」は疑問を示す助詞ですが、使わなくてもかまいません。「お茶」だからといって「コーヒー」を飲んではいけないわけではありません。

\textbf{Notes (EN)}:\\
The word "ocha" means "tea," however, A tea is one of examples to drink when having a break. "Shimasen" is the negative form of the verb "suru" which does not mean negative but an invitation. The word "ka" is a particle that indicates a question, but it is not necessary. Just because it is "ocha," it does not mean that you cannot drink coffee.

\newpage

\section*{33. Like Two Peas in a Pod}

\textbf{Date}: 2024.05.13 (Mon)\\
\textbf{Index}: urifutatsu\\
\textbf{Expression (JA)}: (瓜:うり)(二:ふた)つ\\
\textbf{Romanization}: Uri futatsu\\
\textbf{Expression (EN)}: like two peas in a pod

\textbf{Variants (JA)}: そっくり / 生き(写:うつ)し / よく(似:に)ている\\
\textbf{Variants (EN)}: exactly alike / almost identical / look just the same

\textbf{Intent (JA)}: 類似の表現, 評価, 観察\\
\textbf{Intent (EN)}: expression of resemblance, evaluation, observation

\textbf{Tags (JA)}: 語彙表現, 類似表現, 慣用表現\\
\textbf{Tags (EN)}: lexical expression, resemblance expression, idiomatic expression

\textbf{Adjusted Expression (JA)}:\\
(二人:ふたり)は(見分:みわ)けがつかないほどよく(似:に)ています。\\
\textbf{Adjusted Expression (EN)}:\\
The two are so alike that it is hard to tell them apart.

\textbf{Example (JA)}:\\
(二人:ふたり)はまるで(瓜:うり)(二:ふた)つ。\\
\textbf{Example (EN)}:\\
The two are just like two peas in a pod.

\textbf{Notes (JA)}:\\
「瓜二つ」は「二人がとてもよく似ている」ことを意味します。「瓜」は「メロン」や「ひょうたん」を意味します。

\textbf{Notes (EN)}:\\
The expression "uri futatsu" means "two peas in a pod." The word "uri" means "melon" or "gourd." The word "futatsu" means "two." The word "uri futatsu" means that two people are very similar.

\newpage

\section*{32. Shall I Get One?}

\textbf{Date}: 2024.05.12 (Sun)\\
\textbf{Index}: katteokimashouka\\
\textbf{Expression (JA)}: じゃ、それ、ひとつ、(買:か)っておきましょうか。\\
\textbf{Romanization}: Ja, sore, hitotsu, katte okimashou ka.\\
\textbf{Expression (EN)}: Well then, shall I get one?

\textbf{Variants (JA)}: じゃ、それ、(買:か)っておきましょうか。 / (先:さき)に(買:か)っておきましょうか。\\
\textbf{Variants (EN)}: Well then, shall I buy one? / Shall I get one in advance? / Shall I pick one up for later?

\textbf{Intent (JA)}: 申し出, 準備, 提案\\
\textbf{Intent (EN)}: offering, preparation, suggestion

\textbf{Tags (JA)}: 即時文法, 口語, 申し出表現, 準備表現\\
\textbf{Tags (EN)}: immediate grammar, spoken, offering expression, preparation expression

\textbf{Adjusted Expression (JA)}:\\
では、それをひとつ(買:か)っておきましょうか。\\
\textbf{Adjusted Expression (EN)}:\\
Well then, shall I buy one in advance?

\textbf{Example (JA)}:\\
(後:あと)で(必要:ひつよう)かもしれないから、じゃ、それ、ひとつ、(買:か)っておきましょうか。\\
\textbf{Example (EN)}:\\
We might need it later, so well then, shall I get one?

\textbf{Notes (JA)}:\\
「買っておきましょうか」は、何かを買うことを提案するために使用される。「買って」は「買う」という動詞のて形。「おきましょうか」は、「おく」という動詞の丁寧形で、「する」や「準備する」という意味です。このフレーズは、何かを買って準備することを提案するために使用される。準備として既に買ったものを表現するときは、「てある」を使うのに対し、これから買って、何かに備える場合には、「ておく」を使うとわかりやすい。

\textbf{Notes (EN)}:\\
"Katte okimashou ka" is used to suggest buying something. "Katte" is the te-form of the verb "kau" which means "to buy." "Okimashou ka" is the polite form of the verb "oku" which means "to do" or "to prepare."  It is easy to understand that when expressing something that has already been bought as a preparation, "te aru" is used, and when buying something for preparation, "te oku" is used.

\newpage

\section*{31. It's Already Bought and Ready}

\textbf{Date}: 2024.05.11 (Sat)\\
\textbf{Index}: kattearu\\
\textbf{Expression (JA)}: それは、もう(買:か)ってある。\\
\textbf{Romanization}: Sore wa, mou katte aru.\\
\textbf{Expression (EN)}: It's already bought and ready.

\textbf{Variants (JA)}: もう(買:か)ってあるよ。 / それは先に(買:か)ってあります。\\
\textbf{Variants (EN)}: I've already bought it. / It's already been bought. / I've already got it ready.

\textbf{Intent (JA)}: 準備完了, 安心させる表現, 確認\\
\textbf{Intent (EN)}: completion of preparation, reassurance, confirmation

\textbf{Tags (JA)}: 即時文法, 口語, 準備完了表現, 結果状態表現\\
\textbf{Tags (EN)}: immediate grammar, spoken, preparation-completed expression, result-state expression

\textbf{Adjusted Expression (JA)}:\\
それは、すでに(購入:こうにゅう)してあります。\\
\textbf{Adjusted Expression (EN)}:\\
It has already been purchased and prepared.

\textbf{Example (JA)}:\\
(牛乳:ぎゅうにゅう)はもう(買:か)ってあるから、(今日:きょう)は(買:か)わなくていいよ。\\
\textbf{Example (EN)}:\\
The milk has already been bought, so you don't need to buy it today.

\textbf{Notes (JA)}:\\
「買ってある」は「買う」という動詞の結果ですでに何かが存在していることを意味します。したがって、「もう買わなくてもよい」とか「準備済み」という意味になります。

\textbf{Notes (EN)}:\\
The phrase "katte aru" indicates that something has already been done and is prepared. The word "katte" is the te-form of the verb "kau" which means "to buy." The word "aru" is the verb "aru" which means "to exist." As a result of the action of the verb, something already exists. Therefore, it means "I don't have to buy it anymore" or "it's already prepared."

\newpage

\section*{30. Wonderful!}

\textbf{Date}: 2024.05.10 (Fri)\\
\textbf{Index}: subarashii\\
\textbf{Expression (JA)}: すばらしい！\\
\textbf{Romanization}: Subarashii!\\
\textbf{Expression (EN)}: Wonderful!

\textbf{Variants (JA)}: すばらしいですね。 / ほんとうにすばらしい！ / みごとだ。\\
\textbf{Variants (EN)}: Wonderful! / Excellent! / That's amazing!

\textbf{Intent (JA)}: 賞賛, 感動, 評価\\
\textbf{Intent (EN)}: praise, admiration, evaluation

\textbf{Tags (JA)}: 即時文法, 口語, 評価表現, 感情表現\\
\textbf{Tags (EN)}: immediate grammar, spoken, evaluative expression, emotional expression

\textbf{Adjusted Expression (JA)}:\\
たいへんすばらしいと(思:おも)います。\\
\textbf{Adjusted Expression (EN)}:\\
I think it is truly wonderful.

\textbf{Example (JA)}:\\
すばらしい！ よくできました。\\
\textbf{Example (EN)}:\\
Wonderful! You did very well.

\textbf{Notes (JA)}:\\
「すばらしい」は、強い賞賛や感動を表す表現です。演技、作品、考え方、行動など幅広い対象に使うことができ、即時的な運用では、そのまま感嘆として発せられることが多い表現です。

\textbf{Notes (EN)}:\\
"Subarashii" is used to express strong praise or admiration. It can be used for a performance, an idea, a piece of work, or a person's action, and in immediate use it often comes out as a direct exclamation.

\newpage

\section*{29. That Was Rude of Me}

\textbf{Date}: 2024.05.09 (Thu)\\
\textbf{Index}: shitsureishimashita\\
\textbf{Expression (JA)}: (失礼:しつれい)しました。\\
\textbf{Romanization}: Shitsurei shimashita.\\
\textbf{Expression (EN)}: I'm sorry. / That was rude of me.

\textbf{Variants (JA)}: (失礼:しつれい)しました。 / (失礼:しつれい)いたしました。 / さっきは(失礼:しつれい)しました。\\
\textbf{Variants (EN)}: I'm sorry. / That was rude of me. / Sorry about that.

\textbf{Intent (JA)}: 謝罪, 礼儀, 会話の調整\\
\textbf{Intent (EN)}: apology, politeness, adjusting interaction

\textbf{Tags (JA)}: 即時文法, 定型表現, 丁寧表現, 謝罪表現\\
\textbf{Tags (EN)}: immediate grammar, formulaic expression, polite expression, apology expression

\textbf{Adjusted Expression (JA)}:\\
(先:さき)ほどは(失礼:しつれい)いたしました。\\
\textbf{Adjusted Expression (EN)}:\\
I apologize for my rudeness earlier.

\textbf{Example (JA)}:\\
さっきは(失礼:しつれい)しました。\\
\textbf{Example (EN)}:\\
I'm sorry about earlier.

\textbf{Notes (JA)}:\\
「失礼」は「無礼」や「無礼」を意味します。「しました」は「する」という動詞の過去形で、「する」という動詞の過去形です。これは、謝罪するための丁寧な方法ですが、簡易的なもので、日常的に使ってもよいものです。

\textbf{Notes (EN)}:\\
"Shitsurei shimashita" can be decomposed with a noun "shitsurei" which means "rude" or "impolite," and a verb "shimashita" which is the past form of the verb "suru" which means "to do." This is a polite way to apologize, but it is a simple one and can be used in daily life.

\newpage

\section*{28. Sorry to Keep You Waiting}

\textbf{Date}: 2024.05.08 (Wed)\\
\textbf{Index}: omataseshimashita\\
\textbf{Expression (JA)}: お(待:ま)たせしました。\\
\textbf{Romanization}: Omatase shimashita.\\
\textbf{Expression (EN)}: Sorry to keep you waiting.

\textbf{Variants (JA)}: お(待:ま)たせいたしました。 / お(待:ま)たせしてすみません。\\
\textbf{Variants (EN)}: Sorry to have kept you waiting. / Thank you for waiting.

\textbf{Intent (JA)}: 謝罪, 接客, 会話の再開\\
\textbf{Intent (EN)}: apology, customer service, resuming interaction

\textbf{Tags (JA)}: 即時文法, 定型表現, 丁寧表現, 謝罪表現\\
\textbf{Tags (EN)}: immediate grammar, formulaic expression, polite expression, apology expression

\textbf{Adjusted Expression (JA)}:\\
お(待:ま)たせして(申:もう)し(訳:わけ)ありません。\\
\textbf{Adjusted Expression (EN)}:\\
I apologize for keeping you waiting.

\textbf{Example (JA)}:\\
お(待:ま)たせしました。こちらへどうぞ。\\
\textbf{Example (EN)}:\\
Sorry to keep you waiting. This way, please.

\textbf{Notes (JA)}:\\
「お待たせしました」は、丁寧の接頭辞「お」、動詞「待つ」、使役の助動詞「せる/させる」、丁寧の助動詞「ます」、過去の助動詞「た」で構成される挨拶ことばです。一言で言ってしまえば、難しくありませんが、分析すると非常に難しい接続であることがわかります。このことから、ことばは分析して考えることが得策ではないことが、よくわかります。単に「お待たせしました」は、誰かを待たせてしまったことを謝罪するための表現であると理解すればよいのです。

\textbf{Notes (EN)}:\\
"Omatase shimashita" can be decomposed with a polite prefix "o," a verb "matsu" which means "to wait," a causative auxiliary verb "seru/saseru," a polite auxiliary verb "masu," and a past auxiliary verb "ta." If you say this phrase, it is not difficult, but if you analyze it, you will find that it is a very difficult connection. From this, you can see that it is not a good idea to analyze and think about words. Only understand that "omatase shimashita" is an expression to apologize for making someone wait.

\newpage

\section*{27. Good Night}

\textbf{Date}: 2024.05.07 (Tue)\\
\textbf{Index}: oyasuminasai\\
\textbf{Expression (JA)}: お(休:やす)みなさい。\\
\textbf{Romanization}: Oyasumi nasai.\\
\textbf{Expression (EN)}: Good night.

\textbf{Variants (JA)}: お(休:やす)み。 / では、お(休:やす)みなさい。\\
\textbf{Variants (EN)}: Good night. / Sleep well.

\textbf{Intent (JA)}: あいさつ, 就寝前の表現, 別れの表現\\
\textbf{Intent (EN)}: greeting, expression before sleep, parting expression

\textbf{Tags (JA)}: 定型表現, 口語, あいさつ表現, 別れの表現\\
\textbf{Tags (EN)}: formulaic expression, spoken, greeting expression, parting expression

\textbf{Adjusted Expression (JA)}:\\
それでは、どうぞお(休:やす)みください。\\
\textbf{Adjusted Expression (EN)}:\\
Well then, please have a good rest.

\textbf{Example (JA)}:\\
もう(遅:おそ)いので、お(休:やす)みなさい。\\
\textbf{Example (EN)}:\\
It's getting late, so good night.

\textbf{Notes (JA)}:\\
「おやすみ」は「おやすみなさい」という意味です。「なさい」は、要求や命令をするために使用される丁寧な接尾辞です。「おやすみなさい」は、寝る前に誰かにおやすみなさいと言うために使用されます。また、夜、人と別れるときにも使われます。

\textbf{Notes (EN)}:\\
The word "oyasumi" is a noun made of a verb "yasumu" which means "to rest" or "to take a break." The word "nasai" is a polite suffix used to make a request or a command. The word "oyasumi nasai" is used to say good night to someone before they go to bed. Also, it is used when you part with someone at night.

\newpage

\section*{26. Before a Meal}

\textbf{Date}: 2024.05.06 (Mon)\\
\textbf{Index}: itadakimasu\\
\textbf{Expression (JA)}: いただきます。\\
\textbf{Romanization}: Itadakimasu.\\
\textbf{Expression (EN)}: Thank you for the food.

\textbf{Variants (JA)}: それでは、いただきます。 / ありがたくいただきます。\\
\textbf{Variants (EN)}: Thank you for the meal. / Let's eat.

\textbf{Intent (JA)}: 食事の開始, 感謝, 定型的なあいさつ\\
\textbf{Intent (EN)}: starting a meal, gratitude, formulaic greeting

\textbf{Tags (JA)}: 即時文法, 口語, 定型表現, 食事表現\\
\textbf{Tags (EN)}: immediate grammar, spoken, formulaic expression, meal expression

\textbf{Adjusted Expression (JA)}:\\
それでは、ありがたくいただきます。\\
\textbf{Adjusted Expression (EN)}:\\
Well then, thank you for the food.

\textbf{Example (JA)}:\\
(食:た)べる(前:まえ)に「いただきます」と(言:い)う。\\
\textbf{Example (EN)}:\\
Before eating, people say, Itadakimasu.

\textbf{Notes (JA)}:\\
「いただきます」は、日本で食事の前に言う定型表現です。食べ物や食事の場に対する感謝と、ありがたく受ける気持ちを込めて使われ、単純な逐語訳だけでは捉えにくい表現です。もとの動詞は謙譲の「いただく」ですが、実際には食事開始のあいさつとして機能しています。

\textbf{Notes (EN)}:\\
"Itadakimasu" is a fixed expression said before eating in Japan. It is used to begin a meal with a sense of gratitude and respectful receiving, rather than as a simple literal statement. Although the verb behind it is the humble form "itadaku," in actual use it functions as a formulaic mealtime expression.

\newpage

\section*{25. Excuse Me / I'm Sorry}

\textbf{Date}: 2024.05.05 (Sun)\\
\textbf{Index}: sumimasen\\
\textbf{Expression (JA)}: すみません。\\
\textbf{Romanization}: Sumimasen.\\
\textbf{Expression (EN)}: Excuse me. / I'm sorry.

\textbf{Variants (JA)}: すいません。 / (申:もう)し(訳:わけ)ありません。 / (失礼:しつれい)します。\\
\textbf{Variants (EN)}: Excuse me. / I'm sorry. / Pardon me.

\textbf{Intent (JA)}: 謝罪, 呼びかけ, 依頼, 注意喚起\\
\textbf{Intent (EN)}: apology, getting attention, request, calling for attention

\textbf{Tags (JA)}: 即時文法, 口語, 定型表現, 丁寧表現\\
\textbf{Tags (EN)}: immediate grammar, spoken, formulaic expression, polite expression

\textbf{Adjusted Expression (JA)}:\\
(申:もう)し(訳:わけ)ありません。\\
\textbf{Adjusted Expression (EN)}:\\
I sincerely apologize.

\textbf{Example (JA)}:\\
すみません、ちょっと(教:おし)えてください。\\
\textbf{Example (EN)}:\\
Excuse me, could you tell me something?

\textbf{Notes (JA)}:\\
「すみません」は多くの状況で使用できる多目的なことばです。謝罪するために使用したり、助けを求めたり、誰かの注意を引いたり、お願いをするために使用できます。これは、誰かの注意を引くための丁寧でフォーマルな方法です。間違いを謝罪したり、助けを求めたりするためにも使用されます。

\textbf{Notes (EN)}:\\
"Sumimasen" is a versatile word that can be used in many situations. It can be used to apologize, to ask for help, to get someone's attention, or to ask for a favor. It is a polite and formal way to get someone's attention. It is also used to apologize for a mistake or to ask for help.

\newpage

\section*{24. Making Okonomiyaki}

\textbf{Date}: 2024.05.04 (Sat)\\
\textbf{Index}: okonomiyaki\\
\textbf{Expression (JA)}: お(好:この)み(焼:や)きを(作:つく)る\\
\textbf{Romanization}: Okonomiyaki wo tsukuru\\
\textbf{Expression (EN)}: to make okonomiyaki

\textbf{Variants (JA)}: お(好:この)み(焼:や)きを(作:つく)ります / お(好:この)み(焼:や)きを(焼:や)く\\
\textbf{Variants (EN)}: to make okonomiyaki / to cook okonomiyaki / to make a Japanese savory pancake

\textbf{Intent (JA)}: 料理, 動作の表現, 食べ物の表現\\
\textbf{Intent (EN)}: cooking, expression of action, expression of food

\textbf{Tags (JA)}: 語彙表現, 料理表現, 食べ物の表現\\
\textbf{Tags (EN)}: lexical expression, cooking expression, food expression

\textbf{Adjusted Expression (JA)}:\\
(小麦粉:こむぎこ)、(卵:たまご)、キャベツなどを(混:ま)ぜて、お(好:この)み(焼:や)きを(作:つく)ります。\\
\textbf{Adjusted Expression (EN)}:\\
We make okonomiyaki by mixing flour, eggs, cabbage, and other ingredients and cooking it on a hot plate.

\textbf{Example (JA)}:\\
A「(お好:この)み(焼:や)き、(食:た)べたい？」 B「うん、(大好:だいす)き！」\\
\textbf{Example (EN)}:\\
A: Do you want to eat okonomiyaki? B: Yes, I love it!

\textbf{Notes (JA)}:\\
「お好み焼き」は、日本で人気のある料理で、ピザに似ています。「お好み」は「好きなもの」や「欲しいもの」を意味します。「焼き」は「焼く」や「調理する」を意味します。「作る」は「作る」を意味します。「お好み焼き」は、小麦粉、卵、キャベツ、その他の材料を混ぜ合わせて、熱いプレートで焼いて作ります。いくつかの地域のバリエーションがあります。最も有名なものは広島焼と大阪焼です。子供の頃、私はお好み焼きが大好きでした。その頃のお好み焼きは、今日のように華やかではありませんでした。小麦粉の薄いパンケーキにキャベツの一つまみと豚肉の一切れが入っていました。焼きそばが含まれていた場合、それを買うことができるのはごくわずかな裕福な子供だけでした。お好み焼きは、仮設のガレージで主婦が調理していました。

\textbf{Notes (EN)}:\\
The word "okonomiyaki" is a popular Japanese food that is similar to pizza. The word "okonomi" means "what you like" or "what you want." The word "yaki" means "grilled" or "cooked." The word "tsukuru" means "to make." The word "okonomiyaki" is made by mixing flour, eggs, cabbage, and other ingredients together and grilling them on a hot plate. There are several regional variations of okonomiyaki in Japan. The most famous ones are Hiroshima style and Osaka style. When I was a child, I loved eating okonomiyaki. Okonomiyaki back then was not as glamorous as it is today. It was a simple thin pancake of flour with a pinch of cabbage and a slice of pork. When yakisoba was included, only a few well-to-do children could afford it. Okonomiyaki was cooked by housewives in a temporary garage.

\newpage

\section*{23. Making Omurice}

\textbf{Date}: 2024.05.03 (Fri)\\
\textbf{Index}: omuraisu\\
\textbf{Expression (JA)}: (オムライス:おむらいす)を(作:つく)る\\
\textbf{Romanization}: Omuraisu wo tsukuru\\
\textbf{Expression (EN)}: to make omurice

\textbf{Variants (JA)}: (オムライス:おむらいす)を(作:つく)ります / (卵:たまご)で(包:つつ)んだケチャップライスを(作:つく)る\\
\textbf{Variants (EN)}: to make omurice / to make ketchup rice wrapped in a thin omelet / to make an omelet-wrapped rice dish

\textbf{Intent (JA)}: 料理, 動作の表現, 食べ物の表現\\
\textbf{Intent (EN)}: cooking, expression of action, expression of food

\textbf{Tags (JA)}: 語彙表現, 料理表現, 食べ物の表現\\
\textbf{Tags (EN)}: lexical expression, cooking expression, food expression

\textbf{Adjusted Expression (JA)}:\\
(ケチャップ:けちゃっぷ)(ライス:らいす)を(薄:うす)い(卵:たまご)で(包:つつ)んで、(オムライス:おむらいす)を(作:つく)ります。\\
\textbf{Adjusted Expression (EN)}:\\
We make omurice by wrapping ketchup rice in a thin omelet.

\textbf{Example (JA)}:\\
A「(オムライス:おむらいす)、(食:た)べたい？」 B「うん、(大好:だいす)き！」\\
\textbf{Example (EN)}:\\
A: Do you want to eat omurice? B: Yes, I love it!

\textbf{Notes (JA)}:\\
「オムライス」は、英語の借用語です。「オム」は「オムレツ」の短縮形です。「ライス」は英語の借用語です。「ライス」は「米」を意味します。「作る」は「作る」を意味します。日本に初めて訪れる際、多くの留学生はオムライスが大好きです。ほんのり甘い薄い卵がケチャップライスを包みます。日本で非常に人気のある料理です。

\textbf{Notes (EN)}:\\
The word "Omu-raisu" is a loanword from English. The word "omu" is a short form of the word "omelet." The word "raisu" is a loanword from English. The word "raisu" means "rice." The word "tsukuru" means "to make." When visiting to Japan first time, many international students would love an omuraisu. A hint of sweet thin egg wraps the ketchup rice. It is a very popular dish in Japan.

\newpage

\section*{22. Making Rice Balls}

\textbf{Date}: 2024.05.02 (Thu)\\
\textbf{Index}: onigiri\\
\textbf{Expression (JA)}: おにぎりを(握:にぎ)る\\
\textbf{Romanization}: Onigiri wo nigiru\\
\textbf{Expression (EN)}: to make rice balls

\textbf{Variants (JA)}: おにぎりを(作:つく)る / おむすびを(握:にぎ)る\\
\textbf{Variants (EN)}: to make rice balls / to make onigiri / to shape rice into rice balls

\textbf{Intent (JA)}: 料理, 動作の表現, 食べ物の表現\\
\textbf{Intent (EN)}: cooking, expression of action, expression of food

\textbf{Tags (JA)}: 語彙表現, 料理表現, 食べ物の表現\\
\textbf{Tags (EN)}: lexical expression, cooking expression, food expression

\textbf{Adjusted Expression (JA)}:\\
(炊:た)いたご(飯:はん)を(握:にぎ)って、おにぎりを(作:つく)ります。\\
\textbf{Adjusted Expression (EN)}:\\
We make rice balls by shaping cooked rice with our hands.

\textbf{Example (JA)}:\\
A「(おにぎり:おにぎり)を(握:にぎ)るのは、(楽:たの)しいね！」 B「うん、(炊:た)いたご(飯:はん)を(握:にぎ)って、おにぎりを(作:つく)るのは、(面白:おもしろ)いよね！」\\
\textbf{Example (EN)}:\\
A: Making rice balls is fun! B: Yes, shaping cooked rice into rice balls is interesting!

\textbf{Notes (JA)}:\\
「おにぎり」は「おにぎり」を意味します。「握る」は「作る」や「持つ」を意味します。「握る」はおにぎりを作る行為を表すために使用されます。「おにぎり」は、米を球状に形作り、海苔で包んで作られる人気のある日本食です。

\textbf{Notes (EN)}:\\
The word "onigiri" means "rice ball." The word "nigiru" means "to make" or "to hold." The word "nigiru" is used to describe the action of making a rice ball. The word "onigiri" is a popular Japanese food that is made by shaping rice into a ball and wrapping it in seaweed.

\newpage

\section*{21. You're Speaking Too Fast}

\textbf{Date}: 2024.05.01 (Wed)\\
\textbf{Index}: hayakuchi\\
\textbf{Expression (JA)}: (早口:はやくち)で(聞:き)き(取:と)れない\\
\textbf{Romanization}: Hayakuchi de kikitorenai\\
\textbf{Expression (EN)}: I can't catch what you're saying because you're speaking too fast.

\textbf{Variants (JA)}: (早口:はやくち)でわからない / (早:はや)すぎて(聞:き)き(取:と)れない\\
\textbf{Variants (EN)}: I can't follow you because you're talking too fast. / You're speaking too fast for me to catch it.

\textbf{Intent (JA)}: 聞き返し, 理解困難, 会話の調整\\
\textbf{Intent (EN)}: asking for repetition, difficulty in understanding, adjusting conversation

\textbf{Tags (JA)}: 即時文法, 口語, 聞き取り表現, 会話調整表現\\
\textbf{Tags (EN)}: immediate grammar, spoken, listening-comprehension expression, conversation-adjustment expression

\textbf{Adjusted Expression (JA)}:\\
(少:すこ)し(早口:はやくち)なので、(聞:き)き(取:と)れませんでした。\\
\textbf{Adjusted Expression (EN)}:\\
I could not catch what you are saying because you are speaking a little too fast.

\textbf{Example (JA)}:\\
A「(昨日:きのう)の(映画:えいが)、どうだった？」 B「(早口:はやくち)で(聞:き)き(取:と)れなかったけど、(面白:おもしろ)かったよ。」\\
\textbf{Example (EN)}:\\
A: How was the movie yesterday? B: I couldn't catch it because you were speaking too fast, but it was interesting.

\textbf{Notes (JA)}:\\
「早口」は、話す速度が速いことを表す語です。「聞き取れない」は、相手の言っていることを音としてうまく捉えられないことを表します。全体として、相手の話し方が速すぎて内容をきちんと追えないときに使う表現です。

\textbf{Notes (EN)}:\\
"Hayakuchi" refers to speaking quickly or in rapid speech. "Kikitorenai" means that the listener cannot catch or make out what is being said. The whole expression is used when someone is speaking so fast that the listener cannot follow the speech clearly.

\newpage

\section*{20. Thanks, Crocodile!}

\textbf{Date}: 2024.04.30 (Tue)\\
\textbf{Index}: doumokurokodairu\\
\textbf{Expression (JA)}: どうも、(クロコダイル:くろこだいる)！\\
\textbf{Romanization}: Doumo, kurokodairu!\\
\textbf{Expression (EN)}: Thanks, crocodile!

\textbf{Variants (JA)}: どうも、(アリゲーター:ありげーたー)！ / どうも、(クロコダイル:くろこだいる)！\\
\textbf{Variants (EN)}: Thanks, alligator! / Thanks, crocodile!

\textbf{Intent (JA)}: 感謝, ユーモア, 語呂合わせ\\
\textbf{Intent (EN)}: gratitude, humor, sound-based association

\textbf{Tags (JA)}: ユーモア表現, 発音連想, 言いまちがい\\
\textbf{Tags (EN)}: humorous expression, sound association, speech mistake

\textbf{Adjusted Expression (JA)}:\\
「ありがとう」と言おうとして、まちがって「クロコダイル」と言ってしまいました。\\
\textbf{Adjusted Expression (EN)}:\\
I tried to say "arigatou" but mistakenly said "crocodile."

\textbf{Example (JA)}:\\
「ありがとう」と(言:い)おうとして、「(クロコダイル:くろこだいる)」って(言:い)っちゃった\\
\textbf{Example (EN)}:\\
I tried to say arigatou but ended up saying crocodile.

\textbf{Notes (JA)}:\\
これは、「ありがとう」の代わりに、ユーモラスに使われる言いまちがい的な表現です。もとの逸話では、「ありがとう」を思い出すために「alligator」のような音の似た英語を手がかりにしようとして、実際には「クロコダイル」と言ってしまった、というおかしさがあります。大事なのは正しさよりも、言語をまたいだ音の連想のおもしろさです。

\textbf{Notes (EN)}:\\
This expression is used humorously as a mistaken or playful substitute for "arigatou." In the anecdote behind it, someone tried to remember "arigatou" by linking it to a similar-sounding English word such as "alligator," but ended up saying "crocodile" instead. The point is not correctness but the amusing sound association across languages.

\newpage

\section*{19. Happy Birthday!}

\textbf{Date}: 2024.04.29 (Mon)\\
\textbf{Index}: otanjoubiomedetou\\
\textbf{Expression (JA)}: お(誕生日:たんじょうび)、おめでとう！\\
\textbf{Romanization}: O-tanjoubi omedetou!\\
\textbf{Expression (EN)}: Happy Birthday!

\textbf{Variants (JA)}: お誕生日、おめでとう！ / お誕生日おめでとう！ / 誕生日おめでとう！\\
\textbf{Variants (EN)}: Happy Birthday! / Happy birthday! / Happy birthday to you!

\textbf{Intent (JA)}: 誕生日のお祝い, 祝福の表現, 親しみのある呼びかけ\\
\textbf{Intent (EN)}: celebrating someone's birthday, expressing congratulations, a friendly greeting

\textbf{Tags (JA)}: 即時文法, 定型表現, 挨拶表現, 祝福表現, 口語\\
\textbf{Tags (EN)}: immediate grammar, formulaic expression, greeting expression, congratulatory expression, colloquial

\textbf{Adjusted Expression (JA)}:\\
(彼:かれ)の(誕生日:たんじょうび)を(祝:いわ)っています。\\
\textbf{Adjusted Expression (EN)}:\\
They are celebrating his birthday.

\textbf{Example (JA)}:\\
\\
\textbf{Example (EN)}:\\


\textbf{Notes (JA)}:\\
「お誕生日」は「birthday」を意味します。「お」は丁寧に言うための接頭辞です。「おめでとう」は、特別な機会に誰かを祝うために使われます。

\textbf{Notes (EN)}:\\
The word "o-tanjoubi" means "birthday." "O" is a polite prefix. The word "omedetou" means "congratulations" and is used to celebrate a special occasion.

\newpage

\section*{18. Rome Wasn't Built in a Day}

\textbf{Date}: 2024.04.28 (Sun)\\
\textbf{Index}: ro-ma\\
\textbf{Expression (JA)}: (ローマ:ろーま)は(一日:いちにち)にして(成:な)らず\\
\textbf{Romanization}: Roma wa ichinichi ni shite narazu\\
\textbf{Expression (EN)}: Rome wasn't built in a day.

\textbf{Variants (JA)}: (大:おお)きなことを(成:な)し(遂:と)げるには、(時間:じかん)がかかります。 / (千里:せんり)の(道:みち)も(一歩:いっぽ)から / (石:いし)の(上:うえ)にも(三年:さんねん)\\
\textbf{Variants (EN)}: Great things take time to achieve. / A journey of a thousand miles begins with a single step. / Perseverance prevails even on a stone for three years.

\textbf{Intent (JA)}: 教訓, 励まし, 時間感覚\\
\textbf{Intent (EN)}: lesson, encouragement, sense of time

\textbf{Tags (JA)}: ことわざ, 努力表現, 時間表現\\
\textbf{Tags (EN)}: proverb, expression of effort, expression of time

\textbf{Adjusted Expression (JA)}:\\
(大:おお)きなことを(成:な)し(遂:と)げるには、(時間:じかん)がかかります。\\
\textbf{Adjusted Expression (EN)}:\\
Great things take time to achieve.

\textbf{Example (JA)}:\\
(ローマ:ろーま)は(一日:いちにち)にして(成:な)らず、(努力:どりょく)と(忍耐:にんたい)が(必要:ひつよう)です。\\
\textbf{Example (EN)}:\\
Rome wasn't built in a day. It takes effort and patience.

\textbf{Notes (JA)}:\\
これは、立派なことや大きな仕事は短い時間では成し遂げられない、という意味のよく知られたことわざです。人の成長や仕事の完成には時間がかかることを言うときに使われます。

\textbf{Notes (EN)}:\\
This is a well-known proverb meaning that great things cannot be accomplished quickly. It is used to remind someone that important work, growth, or achievement takes time.

\newpage

\section*{17. All or Nothing}

\textbf{Date}: 2024.04.27 (Sat)\\
\textbf{Index}: ichikabachika\\
\textbf{Expression (JA)}: (一:いち)か(八:ばち)か\\
\textbf{Romanization}: Ichi ka bachi ka\\
\textbf{Expression (EN)}: all or nothing

\textbf{Variants (JA)}: (運:うん)を(天:てん)に(任:まか)せる / (賭:か)けに(出:で)る\\
\textbf{Variants (EN)}: take a gamble / leave it to chance / go all in

\textbf{Intent (JA)}: 決断, 危険, 挑戦\\
\textbf{Intent (EN)}: decision, risk, challenge

\textbf{Tags (JA)}: 即時文法, 口語, 賭けの表現, 決断表現\\
\textbf{Tags (EN)}: immediate grammar, spoken, gambling-related expression, decision expression

\textbf{Adjusted Expression (JA)}:\\
(成功:せいこう)する(保証:ほしょう)はありませんが、(思:おも)い(切:き)ってやってみましょう。\\
\textbf{Adjusted Expression (EN)}:\\
There is no guarantee of success, but let's take a chance and try it.

\textbf{Example (JA)}:\\
(一:いち)か(八:ばち)かやってみよう。\\
\textbf{Example (EN)}:\\
Let's take a chance and try it.

\textbf{Notes (JA)}:\\
「一か八か」は、成功するか大失敗するかわからないまま、危険を承知で思い切ってやってみるときの表現です。語感としては賭け事を思わせるところがあり、当たるか外れるかの覚悟をこめた言い方です。

\textbf{Notes (EN)}:\\
"Ichi ka bachi ka" is an expression used when someone decides to take a risky chance without knowing whether the result will be a great success or a total failure. It is often associated with gambling in tone or background, and it conveys a do-or-die attitude.

\newpage

\section*{16. My Top Pick}

\textbf{Date}: 2024.04.26 (Fri)\\
\textbf{Index}: ichioshi\\
\textbf{Expression (JA)}: (一推:いちお)し\\
\textbf{Romanization}: Ichi-oshi\\
\textbf{Expression (EN)}: top pick

\textbf{Variants (JA)}: おすすめ / 特におすすめ / (推:お)し\\
\textbf{Variants (EN)}: recommendation / top recommendation / favorite pick

\textbf{Intent (JA)}: 推薦, 評価, 強調\\
\textbf{Intent (EN)}: recommendation, evaluation, emphasis

\textbf{Tags (JA)}: 語彙表現, 推薦表現, 評価表現\\
\textbf{Tags (EN)}: lexical expression, recommendation expression, evaluative expression

\textbf{Adjusted Expression (JA)}:\\
(今年:ことし)、(最:もっと)もおすすめしたい(映画:えいが)はこれです。\\
\textbf{Adjusted Expression (EN)}:\\
This is the movie I most highly recommend this year.

\textbf{Example (JA)}:\\
(今年:ことし)(一推:いちお)しの(映画:えいが)はこれだ。\\
\textbf{Example (EN)}:\\
This is my top movie pick of the year.

\textbf{Notes (JA)}:\\
「一推し」は「最高の一つ」という意味です。「一」は「一」や「最初」を意味します。「推し」は「推す」という意味で、「推薦」や「支持」を意味します。「一推し」は、そのクラスやカテゴリーの中で最高のものの一つや最高の選択肢の一つであるものを表すために使用されます。

\textbf{Notes (EN)}:\\
The word "ichi-oshi" means "one of the best." The word "ichi" means "one" or "first." The word oshi" means "push" which is used to mean "recommend" or "support." The word "ichi-oshi" is used to describe something that is one of the best in its class or category or one of the best choices.

\newpage

\section*{15. A Once-in-a-Lifetime Encounter}

\textbf{Date}: 2024.04.25 (Thu)\\
\textbf{Index}: ichigoichie\\
\textbf{Expression (JA)}: (一期一会:いちごいちえ)\\
\textbf{Romanization}: Ichigo-ichie\\
\textbf{Expression (EN)}: a once-in-a-lifetime encounter

\textbf{Variants (JA)}: (一生:いっしょう)に(一度:いちど) / またとない(出会:であ)い\\
\textbf{Variants (EN)}: once in a lifetime / a unique encounter / a meeting that comes only once

\textbf{Intent (JA)}: 価値づけ, 大切さの表明, 出会いの表現\\
\textbf{Intent (EN)}: valuation, expressing preciousness, expression of encounter

\textbf{Tags (JA)}: 語彙表現, 価値表現, 出会いの表現\\
\textbf{Tags (EN)}: lexical expression, value expression, expression of encounter

\textbf{Adjusted Expression (JA)}:\\
(一生:いっしょう)に(一度:いちど)しかない(出会:であ)いとして、(大切:たいせつ)にしたいと(思:おも)います。\\
\textbf{Adjusted Expression (EN)}:\\
I would like to cherish it as a once-in-a-lifetime encounter.

\textbf{Example (JA)}:\\
(一期一会:いちごいちえ)を(大切:たいせつ)にしたいと(思:おも)います。\\
\textbf{Example (EN)}:\\
I would like to cherish this once-in-a-lifetime encounter.

\textbf{Notes (JA)}:\\
「一期一会」は「一生に一度」という意味です。「一期」は「一度」や「一期間」を意味します。「一会」は「出会い」や「出会い」を意味します。「一期一会」は、一生に一度だけの出会いを表すために使用されます。この言葉は、トム・ハンクスが主人公を演じる映画「フォレスト・ガンプ」の日本語のサブタイトルとして使用されています。

\textbf{Notes (EN)}:\\
The word "ichigo-ichie" means "once in a lifetime." The word "ichigo" means "one time" or "one period." The word "ichie" means "meeting" or "encounter." The word "ichigo-ichie" is used to as a sub title in Japanese of the movie "Forrest Gump." in which the main character meets many people only once in his life. Tom Hanks plays the main character.

\newpage

\section*{14. Everything, All of It}

\textbf{Date}: 2024.04.24 (Wed)\\
\textbf{Index}: issaigassai\\
\textbf{Expression (JA)}: (一切合財:いっさいがっさい)\\
\textbf{Romanization}: Issaigassai\\
\textbf{Expression (EN)}: everything, all of it

\textbf{Variants (JA)}: なんでもかんでも / ありとあらゆる / (何:なに)もかも\\
\textbf{Variants (EN)}: everything / all of it / the whole nine yards

\textbf{Intent (JA)}: 包括, 強調, 申し出\\
\textbf{Intent (EN)}: inclusion, emphasis, offering service

\textbf{Tags (JA)}: 語彙表現, 包括表現, 強調表現\\
\textbf{Tags (EN)}: lexical expression, inclusive expression, emphatic expression

\textbf{Adjusted Expression (JA)}:\\
とりあえず、(考:かんが)えられるものをすべて(取:と)りまとめて、(対応:たいおう)します。\\
\textbf{Adjusted Expression (EN)}:\\
For now, we will gather together everything that can be considered and deal with it as a whole.

\textbf{Example (JA)}:\\
(一切合財:いっさいがっさい)、まとめてこちらでお(引:ひ)き(受:う)けします。\\
\textbf{Example (EN)}:\\
We will handle everything here, all together.

\textbf{Notes (JA)}:\\
「一切合財」は、「何も残さず全部」「まるごとすべて」という意味の表現です。全体をひとまとめにして強めに言う感じがあり、日常の中立的な言い方よりは、やや言い回しとして目立つことがあります。「なんでもかんでも」「ありとあらゆる」「何もかも」なども近い表現ですが、語感や使われ方は完全には同じではありません。

\textbf{Notes (EN)}:\\
"Issai-gassai" means "everything" or "all of it," with nothing left out. It often has a slightly emphatic tone and can sound somewhat old-fashioned or stylistically marked compared with neutral everyday expressions. Related expressions such as "nandemo-kandemo" and "ari to arayuru" are similar, but their tone and usage are not exactly the same.

\newpage

\section*{13. Doing One's Best}

\textbf{Date}: 2024.04.23 (Tue)\\
\textbf{Index}: isshokemmei\\
\textbf{Expression (JA)}: (一所懸命:いっしょけんめい)\\
\textbf{Romanization}: Issho-kemmei\\
\textbf{Expression (EN)}: doing one's very best

\textbf{Variants (JA)}: (一所懸命:いっしょけんめい) / (全力:ぜんりょく)で / (精一杯:せいいっぱい)\\
\textbf{Variants (EN)}: doing one's very best / with all one's effort / to the best of one's ability

\textbf{Intent (JA)}: 物事に全力で取り組むことを表す表現, 努力や頑張りを表す表現, 仕事や勉強などに対する姿勢を表す表現\\
\textbf{Intent (EN)}: Expression to indicate doing something with all one's effort, Expression to indicate effort and hard work, Expression to indicate attitude towards work or study

\textbf{Tags (JA)}: 表現, 日常会話, 努力の表現, 頑張りの表現, 仕事や勉強の姿勢の表現\\
\textbf{Tags (EN)}: expression, daily conversation, expression of effort, expression of hard work, expression of attitude towards work or study

\textbf{Adjusted Expression (JA)}:\\
(一所懸命:いっしょけんめい)に(勉強:べんきょう)する\\
\textbf{Adjusted Expression (EN)}:\\
to study with all one's effort

\textbf{Example (JA)}:\\
(一所懸命:いっしょけんめい)、がんばります。\\
\textbf{Example (EN)}:\\
I will do my very best.

\textbf{Notes (JA)}:\\
「一所懸命」は、力を尽くして物事に取り組むことを表す表現です。もともとは一つの所領に命をかけるという文字どおりの意味をもっていましたが、現代日本語では「全力で取り組む」「精いっぱいがんばる」という意味で広く使われます。

\textbf{Notes (EN)}:\\
The expression "issho-kemmei" means to do something with all one's effort. Historically, it literally referred to staking one's life on a single place or holding, but in modern Japanese it is commonly used to mean "doing one's very best" or "working as hard as possible."

\newpage

\section*{12. Top-Class}

\textbf{Date}: 2024.04.22 (Mon)\\
\textbf{Index}: ichiryu\\
\textbf{Expression (JA)}: (一流:いちりゅう)\\
\textbf{Romanization}: Ichi-ryu\\
\textbf{Expression (EN)}: top-class

\textbf{Variants (JA)}: (一流:いちりゅう) / (最高位:さいこうい) / (トップクラス:とっぷくらす)\\
\textbf{Variants (EN)}: top-class / the highest level / first-rate

\textbf{Intent (JA)}: 人の評価, 技術の評価, 製品の評価, 仕事ぶりの評価\\
\textbf{Intent (EN)}: Evaluation of a person, Evaluation of skill, Evaluation of product, Evaluation of performance

\textbf{Tags (JA)}: 表現, 日常会話, 評価, トップクラス, 最上位\\
\textbf{Tags (EN)}: expression, daily conversation, evaluation, top-class, the highest level

\textbf{Adjusted Expression (JA)}:\\
(一流:いちりゅう)の(料理人:りょうりにん)\\
\textbf{Adjusted Expression (EN)}:\\
a top-class chef

\textbf{Example (JA)}:\\
(中村:なかむら)さんは(一流:いちりゅう)の(料理人:りょうりにん)です。\\
\textbf{Example (EN)}:\\
Mr. Nakamura is a top-class chef.

\textbf{Notes (JA)}:\\
「一流」は、その分野で特にすぐれていて、水準が高いことを表す語です。人、技術、製品、仕事ぶりなどについて、「最高級だ」「第一級だ」という評価を表すときに使われます。

\textbf{Notes (EN)}:\\
The word "ichi-ryu" means "top-class," "first-rate," or "of the highest level." It is used to describe a person, skill, product, or performance that is regarded as outstanding in its field.

\newpage

\section*{11. Verb-ta / Adj-i / Adj-na + koto ni}

\textbf{Date}: 2024.04.21 (Sun)\\
\textbf{Index}: odoroitakotoni\\
\textbf{Expression (JA)}: おどろいたことに、...\\
\textbf{Romanization}: Odoroita koto ni, ...\\
\textbf{Expression (EN)}: surprisingly, ...

\textbf{Variants (JA)}: おもしろいことに、... / うれしいことに、... / 残念なことに、...\\
\textbf{Variants (EN)}: Interestingly, ... / Fortunately, ... / Unfortunately, ...

\textbf{Intent (JA)}: 評価の表現, 見方の表現, 反応の表現\\
\textbf{Intent (EN)}: Expression of evaluation, Expression of viewpoint, Expression of reaction

\textbf{Tags (JA)}: 表現, 日常会話, 評価の表現, 見方の表現, 反応の表現\\
\textbf{Tags (EN)}: expression, daily conversation, expression of evaluation, expression of viewpoint, expression of reaction

\textbf{Adjusted Expression (JA)}:\\
(驚:おどろ)きの(事実:じじつ)とでも(言:い)いましょうか、\\
\textbf{Adjusted Expression (EN)}:\\
How can I say it? It's a surprising fact that ...

\textbf{Example (JA)}:\\
(驚:おどろ)いたことに、(彼:かれ)は(日本人:にほんじん)なみに(日本語:にほんご)を(話:はな)す。\\
\textbf{Example (EN)}:\\
Surprisingly, he speaks Japanese almost as well as a native Japanese speaker.

\textbf{Notes (JA)}:\\
「動詞タ形 / 形容詞イ形 / 形容動詞 + ことに」は、文頭で話し手の評価や見方を添えるための言い方です。「おもしろいことに」「おどろいたことに」のような形は、英語の「interestingly」や「surprisingly」に近い働きをします。ここでの「ことに」は、続く内容に対する話し手の反応や視点を前置きする表現です。

\textbf{Notes (EN)}:\\
The pattern "verb-ta / adjective-i / adjective-na + koto ni" is used to add the speaker's evaluative comment at the beginning of a sentence. Expressions such as "omoshiroi koto ni" and "odoroita koto ni" work like sentence adverbs such as "interestingly" or "surprisingly." Here, "koto" is a nominal element, and the whole pattern presents the speaker's reaction or viewpoint toward what follows.

\newpage

\section*{10. Choice or Sentakushi}

\textbf{Date}: 2024.04.20 (Sat)\\
\textbf{Index}: sentakushi\\
\textbf{Expression (JA)}: (チョイス:ちょいす) or (選択肢:せんたくし)\\
\textbf{Romanization}: Choisu or sentaku-shi\\
\textbf{Expression (EN)}: choice or option

\textbf{Variants (JA)}: (選択肢:せんたくし)が(いくつか:いくつか)あって、.. / そういう(選択肢:せんたくし)があって、.. / いくつか(選:えら)べて、..\\
\textbf{Variants (EN)}: There are several options for that, ... / There are such options, ... / There are several choices, ...

\textbf{Intent (JA)}: 選択肢の提示, 選択肢の説明, 選択肢の提案\\
\textbf{Intent (EN)}: Presenting options, Explaining options, Suggesting options

\textbf{Tags (JA)}: 表現, 日常会話, 選択肢の表現, 提案の表現\\
\textbf{Tags (EN)}: expression, daily conversation, option expression, suggestion expression

\textbf{Adjusted Expression (JA)}:\\
いくつか(選択肢:せんたくし)から、お(選:えら)びできます。\\
\textbf{Adjusted Expression (EN)}:\\
You can choose from several options.

\textbf{Example (JA)}:\\
ええ、それにはいくつか(選択肢:せんたくし)があって、..\\
\textbf{Example (EN)}:\\
Yes, there are several options for that, ...

\textbf{Notes (JA)}:\\
「チョイス」は「選択肢」や「選択」を意味するために日本語で使用されます。これは英語からの借用語です。日本語の「選択肢」や「選択」のための日本語の言葉は「選択」や「選択」です。日本語をマスターするため、または良い日本語のユーザーであるふりをするために、このような固定フレーズをいくつか用意しておきます。

\textbf{Notes (EN)}:\\
The word "choice" is used in Japanese to mean "option" or "selection." It is a loanword from English. The Japanese word for "choice" is "sentaku" or "sentaku-shi." To master Japanese, or to pretend to be a good Japanese user, you prepare some fixed phrases like this one.

\newpage

\section*{9. Don't Do That}

\textbf{Date}: 2024.04.19 (Fri)\\
\textbf{Index}: suruna\\
\textbf{Expression (JA)}: するな！\\
\textbf{Romanization}: Suru na!\\
\textbf{Expression (EN)}: Don't do that!

\textbf{Variants (JA)}: やるな！ / するなよ！ / やるなってば！\\
\textbf{Variants (EN)}: Don't do it! / Don't do that, okay? / I said, don't do that!

\textbf{Intent (JA)}: 禁止の説明, 禁止の警告, 禁止の強調\\
\textbf{Intent (EN)}: Explanation of prohibition, Warning of prohibition, Emphasis of prohibition

\textbf{Tags (JA)}: 表現, 日常会話, 禁止の表現, 命令形の表現\\
\textbf{Tags (EN)}: expression, daily conversation, prohibition expression, imperative expression

\textbf{Adjusted Expression (JA)}:\\
そのような(行為:こうい)は(禁止:きんし)されています。\\
\textbf{Adjusted Expression (EN)}:\\
Such behavior is prohibited.

\textbf{Example (JA)}:\\
するな！\\
\textbf{Example (EN)}:\\
Don't do that!

\textbf{Notes (JA)}:\\
「するな」は、「する」という動詞の命令形です。「な」は、「ない」という動詞の否定命令形です。「するな」は、誰かに何かをしないように伝えるために使用されます。それは、誰かに何かをしないように伝えるための強い、直接的な方法です。

\textbf{Notes (EN)}:\\
The word "suruna" is a command form of the verb "suru" which means "to do." The word "na" is a negative command form of the verb "nai" which means "not to do." The word "suruna" is used to tell someone not to do something. It is a strong and direct way to tell someone not to do something.

\newpage

\section*{8. No Verb, Still a Sentence}

\textbf{Date}: 2024.04.18 (Thu)\\
\textbf{Index}: onakaippai\\
\textbf{Expression (JA)}: お(腹:なか)、いっぱい。\\
\textbf{Romanization}: Onaka, ippai.\\
\textbf{Expression (EN)}: I'm full.

\textbf{Variants (JA)}: もう、いっぱい。 / (食:た)べられない。\\
\textbf{Variants (EN)}: I'm full already. / I can't eat any more.

\textbf{Intent (JA)}: お腹がいっぱいです, もう、食べられません, もう、いっぱいです\\
\textbf{Intent (EN)}: I'm full., I can't eat any more., I'm full already.

\textbf{Tags (JA)}: 表現, 日常会話, 食事の表現, 動詞なしの表現\\
\textbf{Tags (EN)}: expression, daily conversation, food expression, expression without verb

\textbf{Adjusted Expression (JA)}:\\
もうこれ(以上:いじょう)は(食:た)べられません。\\
\textbf{Adjusted Expression (EN)}:\\
I can't eat any more.

\textbf{Example (JA)}:\\
お(腹:なか)、いっぱい。\\
\textbf{Example (EN)}:\\
I'm full.

\textbf{Notes (JA)}:\\
「お腹、いっぱい」は「お腹」と「いっぱい」だけですが、発話として十分に成り立つ表現です。動詞がなくても文として理解されます。日本語にはこのような言い方が多く見られます。

\textbf{Notes (EN)}:\\
"Onaka, ippai" consists of just two parts, "onaka" and "ippai," and yet it functions as a complete utterance. Even without a verb, it is understood as a full sentence in context, and Japanese has many expressions of this kind.

\newpage

\section*{7. Taking the Best Parts}

\textbf{Date}: 2024.04.17 (Wed)\\
\textbf{Index}: iitokodori\\
\textbf{Expression (JA)}: いいとこどり\\
\textbf{Romanization}: Ii-toko-dori\\
\textbf{Expression (EN)}: taking only the best parts

\textbf{Variants (JA)}: (都会:とかい)と(田舎:いなか)のいいとこどり。 / (仕事:しごと)のいいとこどりをする。 / この(店:みせ)は(安:やす)いし、(美味:おい)しいし、(いいとこどり)だね。\\
\textbf{Variants (EN)}: It takes the best parts of both city life and country life. / I take only the best parts of my job. / This restaurant is cheap, delicious, and it takes only the best parts.

\textbf{Intent (JA)}: いいとこどりをする, いいとこどりの生活をする, いいとこどりの店だね\\
\textbf{Intent (EN)}: to take only the best parts, to live a life that takes only the best parts, This restaurant takes only the best parts.

\textbf{Tags (JA)}: 表現, 日常会話, 褒める表現, 批判的な表現\\
\textbf{Tags (EN)}: expression, daily conversation, praising expression, critical expression

\textbf{Adjusted Expression (JA)}:\\
(双方:そうほう)の(利点:りてん)を(同時:どうじ)に(取:と)り(入:い)れました。\\
\textbf{Adjusted Expression (EN)}:\\
I took the advantages of both at the same time.

\textbf{Example (JA)}:\\
(都会:とかい)と(田舎:いなか)のいいとこどり。\\
\textbf{Example (EN)}:\\
It takes the best parts of both city life and country life.

\textbf{Notes (JA)}:\\
これは、何かの中で都合のよい部分や魅力的な部分だけを取ることを表す表現です。長所をうまく組み合わせるという前向きな言い方にも、自分に都合のよいところだけ取るという批判的な言い方にも使われます。

\textbf{Notes (EN)}:\\
This expression refers to taking or enjoying only the most favorable parts of something. It can be used positively, as in combining advantages, or critically, as in taking only what suits oneself.

\newpage

\section*{6. Reservation}

\textbf{Date}: 2024.04.16 (Tue)\\
\textbf{Index}: goyoyaku\\
\textbf{Expression (JA)}: ご(予約:よやく)\\
\textbf{Romanization}: Go-yoyaku\\
\textbf{Expression (EN)}: reservation

\textbf{Variants (JA)}: ご(予約:よやく)願います。 / ご(予約:よやく)をお(願:おねが)いします。\\
\textbf{Variants (EN)}: I would like to make a reservation. / I would like to make a reservation, please.

\textbf{Intent (JA)}: 予約をしたい, 予約をお願いします, 予約をしたいのですが\\
\textbf{Intent (EN)}: I want to make a reservation., I would like to make a reservation, please., I would like to make a reservation.

\textbf{Tags (JA)}: 表現, 日常会話, サービス業の表現\\
\textbf{Tags (EN)}: expression, daily conversation, service industry expression

\textbf{Adjusted Expression (JA)}:\\
ご(予約:よやく)をお(願:ねが)いします。\\
\textbf{Adjusted Expression (EN)}:\\
I would like to make a reservation, please.

\textbf{Example (JA)}:\\
(本日:ほんじつ)はご(予約:よやく)のお(客様:きゃくさま)のみの(営業:えいぎょう)とさせていただきます。\\
\textbf{Example (EN)}:\\
Today, we are open only to customers with reservations.

\textbf{Notes (JA)}:\\
これは、レストラン業界でよく使用されるフレーズです。しかし、本当に歩いてきた顧客を受け入れないレストランがあるかどうかはわかりません。

\textbf{Notes (EN)}:\\
This is a phrase that is often used in the restaurant business. But I do not know the restaurant really do not accept walk-in customers.

\newpage

\section*{5. I Can't Handle Hot Food}

\textbf{Date}: 2024.04.15 (Mon)\\
\textbf{Index}: nekojita\\
\textbf{Expression (JA)}: (猫舌:ねこじた)。\\
\textbf{Romanization}: Nekojita\\
\textbf{Expression (EN)}: I can't handle hot food or drinks.

\textbf{Variants (JA)}: 熱いものは苦手だ。 / 少し冷めてからでないと無理だ。\\
\textbf{Variants (EN)}: I can't handle hot food. / I can't drink hot drinks. / I need it to cool down first.

\textbf{Intent (JA)}: 熱い食べ物や飲み物が苦手, 熱いものを口にできない, 熱いものは避けたい\\
\textbf{Intent (EN)}: dislike hot food or drinks, cannot comfortably eat or drink hot things, prefer to avoid hot things

\textbf{Tags (JA)}: 表現, 日常会話, 食べ物の表現\\
\textbf{Tags (EN)}: expression, daily conversation, food expression

\textbf{Adjusted Expression (JA)}:\\
(熱:あつ)いものは(苦手:にがて)だと(思:おも)います。\\
\textbf{Adjusted Expression (EN)}:\\
I think I can't handle hot food or drinks.

\textbf{Example (JA)}:\\
A「(猫舌:ねこじた)だから、(熱:あつ)いのはちょっと」 B「じゃ、(アイス:あいす)にする？」\\
\textbf{Example (EN)}:\\
A: I can't handle hot things, so hot food is a bit... B: Then, how about cold things?

\textbf{Notes (JA)}:\\
英語にはこのような言葉はないので、英語話者は単に「熱いものが好きではない」、「熱い食べ物が嫌い」、「熱いものを食べられない」などと言います。

\textbf{Notes (EN)}:\\
There is no such word in English, so English speakers simply say "I don't like hot things"; "I dislike hot foods"; or "I cannot eat/drink hot things."

\newpage

\section*{4. You Have Potential}

\textbf{Date}: 2024.04.14 (Sun)\\
\textbf{Index}: nobishiro\\
\textbf{Expression (JA)}: (伸:の)び(代:しろ)がある。\\
\textbf{Romanization}: Nobishiro ga aru.\\
\textbf{Expression (EN)}: There is room for growth.

\textbf{Variants (JA)}: まだ上達する。 / まだ成長する。 / もっと上手になる。\\
\textbf{Variants (EN)}: There is room for growth. / You can still improve. / You can get better.

\textbf{Intent (JA)}: 成長の可能性, 改善の余地, 励まし\\
\textbf{Intent (EN)}: potential for growth, room for improvement, encouragement

\textbf{Tags (JA)}: 表現, 日常会話, 励ましの表現\\
\textbf{Tags (EN)}: expression, daily conversation, encouraging expression

\textbf{Adjusted Expression (JA)}:\\
(伸:の)び(代:しろ)があると(思:おも)います。\\
\textbf{Adjusted Expression (EN)}:\\
I think there is room for growth.

\textbf{Example (JA)}:\\
A「この記録だったら、まだまだ(伸:の)び(代:しろ)がある」 B「そうだね」\\
\textbf{Example (EN)}:\\
A: With this record, there is still room for growth. B: Yeah, I know.

\textbf{Notes (JA)}:\\
これは、改善の可能性があることを意味します。誰かを励ますときにこのフレーズを使用できます。

\textbf{Notes (EN)}:\\
It means that there is a possibility of improvement. You can use this phrase when you want to encourage someone.

\newpage

\section*{3. What a Farce}

\textbf{Date}: 2024.04.13 (Sat)\\
\textbf{Index}: chabangeki\\
\textbf{Expression (JA)}: (茶番劇:ちゃばんげき)。\\
\textbf{Romanization}: Chaban geki.\\
\textbf{Expression (EN)}: What a farce.

\textbf{Variants (JA)}: 茶番劇だ。 / インチキだ。 / とんだ茶番だ。\\
\textbf{Variants (EN)}: What a farce. / It's a sham. / What a ridiculous farce.

\textbf{Intent (JA)}: 不条理な状況, ばかばかしい, しらじらしい\\
\textbf{Intent (EN)}: absurd situation, ridiculous, insincere

\textbf{Tags (JA)}: 表現, 日常会話, 感情表現\\
\textbf{Tags (EN)}: expression, daily conversation, emotional expression

\textbf{Adjusted Expression (JA)}:\\
(茶番劇:ちゃばんげき)だと(思:おも)います。\\
\textbf{Adjusted Expression (EN)}:\\
I think it's a farce.

\textbf{Example (JA)}:\\
A「(茶番劇:ちゃばんげき)だ。」 B「そうだね」\\
\textbf{Example (EN)}:\\
A: What a farce. B: Yeah, I know.

\textbf{Notes (JA)}:\\
馬鹿げたまたは不条理な状況の文脈でよく使用されます。

\textbf{Notes (EN)}:\\
It is often used in the context of ridiculous or absurd situations.

\newpage

\section*{2. I'm Fed Up}

\textbf{Date}: 2024.04.12 (Fri)\\
\textbf{Index}: mouunzari\\
\textbf{Expression (JA)}: もう、うんざり。\\
\textbf{Romanization}: Mou unzari.\\
\textbf{Expression (EN)}: I'm fed up already.

\textbf{Variants (JA)}: もう、うんざり。 / もう、うんざりだ。 / もう、うんざりです。\\
\textbf{Variants (EN)}: I'm fed up already. / I'm sick of it already. / I've had enough already.

\textbf{Intent (JA)}: 不満, 疲れ, 主張\\
\textbf{Intent (EN)}: dissatisfaction, fatigue, assertion

\textbf{Tags (JA)}: 感情表現, 日常会話\\
\textbf{Tags (EN)}: emotional expression, daily conversation

\textbf{Adjusted Expression (JA)}:\\
もう、うんざりだと(思:おも)います。\\
\textbf{Adjusted Expression (EN)}:\\
I think I'm fed up already.

\textbf{Example (JA)}:\\
A「もう、うんざり。」 B「そうだね」\\
\textbf{Example (EN)}:\\
A: I'm fed up already. B: Yeah, I know.

\textbf{Notes (JA)}:\\
何かや誰かにうんざりしていることを意味します。 不当に扱われると、うんざりするかもしれません。

\textbf{Notes (EN)}:\\
It means that you are tired of something or someone. If you are treated unfairly, you may feel unzari.

\newpage

\section*{1. Do It or You'll Miss Out}

\textbf{Date}: 2024.04.11 (Thu)\\
\textbf{Index}: yaranakya\\
\textbf{Expression (JA)}: やらなきゃ(損:そん)！\\
\textbf{Romanization}: Yaranakya son!\\
\textbf{Expression (EN)}: You'll miss out if you don't do it!

\textbf{Variants (JA)}: やらなきゃ損！ / やらなければ損！ / やらないと損！\\
\textbf{Variants (EN)}: You'll miss out if you don't do it! / If you don't do it, you'll miss out. / You should do it, or you'll miss out.

\textbf{Intent (JA)}: 必要, 損得, 主張\\
\textbf{Intent (EN)}: necessity, profit and loss, assertion

\textbf{Tags (JA)}: 即時文法, 日常会話\\
\textbf{Tags (EN)}: immediate grammar, daily conversation

\textbf{Adjusted Expression (JA)}:\\
やらないと(損:そん)だと(思:おも)います。\\
\textbf{Adjusted Expression (EN)}:\\
I think you'll miss out if you don't do it.

\textbf{Example (JA)}:\\
A「やらなきゃ損！」 B「そうだね」\\
\textbf{Example (EN)}:\\
A: You'll miss out if you don't do it! B: Yeah, that's true.

\textbf{Notes (JA)}:\\
いつか何か良い利益や利益が得られるので、それをやるべきだという意味です。

\textbf{Notes (EN)}:\\
It means that you should do it because you will have some good profit or benefit someday.

\newpage

\section*{0. AEAD}

\textbf{Date}: 2024.04.10 (Wed)\\
\textbf{Index}: aead\\
\textbf{Expression (JA)}: (AEAD:えいあど)\\
\textbf{Romanization}: AEAD\\
\textbf{Expression (EN)}: An expression a day

\textbf{Variants (JA)}: \\
\textbf{Variants (EN)}: 

\textbf{Intent (JA)}: \\
\textbf{Intent (EN)}: 

\textbf{Tags (JA)}: \\
\textbf{Tags (EN)}: 

\textbf{Adjusted Expression (JA)}:\\
\\
\textbf{Adjusted Expression (EN)}:\\


\textbf{Example (JA)}:\\
\\
\textbf{Example (EN)}:\\


\textbf{Notes (JA)}:\\


\textbf{Notes (EN)}:\\


\newpage

\section*{. }

\textbf{Date}:  ()\\
\textbf{Index}: \\
\textbf{Expression (JA)}: \\
\textbf{Romanization}: \\
\textbf{Expression (EN)}: 

\textbf{Variants (JA)}: \\
\textbf{Variants (EN)}: 

\textbf{Intent (JA)}: \\
\textbf{Intent (EN)}: 

\textbf{Tags (JA)}: \\
\textbf{Tags (EN)}: 

\textbf{Adjusted Expression (JA)}:\\
\\
\textbf{Adjusted Expression (EN)}:\\


\textbf{Example (JA)}:\\
\\
\textbf{Example (EN)}:\\


\textbf{Notes (JA)}:\\


\textbf{Notes (EN)}:\\


\newpage

\section*{. }

\textbf{Date}:  ()\\
\textbf{Index}: \\
\textbf{Expression (JA)}: \\
\textbf{Romanization}: \\
\textbf{Expression (EN)}: 

\textbf{Variants (JA)}: \\
\textbf{Variants (EN)}: 

\textbf{Intent (JA)}: \\
\textbf{Intent (EN)}: 

\textbf{Tags (JA)}: \\
\textbf{Tags (EN)}: 

\textbf{Adjusted Expression (JA)}:\\
\\
\textbf{Adjusted Expression (EN)}:\\


\textbf{Example (JA)}:\\
\\
\textbf{Example (EN)}:\\


\textbf{Notes (JA)}:\\


\textbf{Notes (EN)}:\\


\newpage

\section*{. }

\textbf{Date}:  ()\\
\textbf{Index}: \\
\textbf{Expression (JA)}: \\
\textbf{Romanization}: \\
\textbf{Expression (EN)}: 

\textbf{Variants (JA)}: \\
\textbf{Variants (EN)}: 

\textbf{Intent (JA)}: \\
\textbf{Intent (EN)}: 

\textbf{Tags (JA)}: \\
\textbf{Tags (EN)}: 

\textbf{Adjusted Expression (JA)}:\\
\\
\textbf{Adjusted Expression (EN)}:\\


\textbf{Example (JA)}:\\
\\
\textbf{Example (EN)}:\\


\textbf{Notes (JA)}:\\


\textbf{Notes (EN)}:\\


\newpage
\end{document}

